\documentclass[10pt,a4paper,twoside,bg=print,twocolumn,openright,nodeprecatedcode]{dndbook}
\usepackage[utf8]{inputenc}
\usepackage{tabulary}
\usepackage[normalem]{ulem}
\usepackage{color,soul}
\usepackage{float}
\usepackage{caption}
\usepackage{wrapfig}
\usepackage{tocloft}
\usepackage[toc]{multitoc}
\usepackage[hidelinks]{hyperref}
\usepackage{csquotes}
\raggedbottom
\extrafloats{2000}
\cftsetindents{section}{0em}{0em}
\cftsetindents{subsection}{0em}{0em}
\cftsetindents{subsubsection}{0.2em}{0em}
\setcounter{tocdepth}{1}
\renewcommand*{\multicolumntoc}{3}
\setlist[description]{nolistsep,listparindent=3pt,topsep=6pt}
\title{ \Huge \scshape \nodesto\fontsize{50}{40}\selectfont
 Encyclopedia Exandria \\
 \medskip
 \nodesto
 \medskip\Huge
 A compendium of the Critical Role world}
\date{}
\begin{document}
\frontmatter
\maketitle
\tableofcontents
\mainmatter
\part{Exandria}
\chapter{Exandria Overview}
\section{Myth of Exandria}
\DndDropCapLine{B}{efore the wonder of these} lands can be explored, their context in the history of Exandria itself must be thoroughly established. The world of Exandria is home to four major continents, dozens of nations, and untold tiny pockets of civilization amid a wondrous, dangerous wilderness. Every civilization, from the largest empire to the tiniest nomadic clan, has its own interpretation of where its story began. Even within the world of Exandria, different cultures have creation myths that eventually converge with history, but there is no known definitive story.


Even so, life ever seeks to understand its inception. In this report, the collected minds of the \textit{Cobalt Soul} and \textit{Emon's} \textit{Alabaster Lyceum} can say, with only the most trivial of doubts, that the city of Vasselheim in \textit{Othanzia} truly deserves its title of the "Dawn City." It is here, to our researchers' best understanding, where the stories of all mortal peoples begin. Vasselheim became Exandria's oldest surviving city by virtue of not just being built long ago, but by enduring the \textit{Calamity} (which shall be elucidated presently), a miraculous feat that few other charted settlements can claim to have accomplished.

In short, Vasselheim houses the earliest known temples to the gods, and the earliest known records of history that survived this catastrophe. It is these sources that this joint research initiative treats as its primary documents, though great effort has been made to corroborate their claims through later sources. It is in the chronicles of Vasselheim that the earliest account of the myth of the \textit{Founding} is recorded—and indeed, this history is the most commonly accepted origin of the world in the Republic of Tal'Dorei. Though countless regional variations muddy the historical record, the broad strokes of this mythic history are undeniable.

\subsection{The Founding}
Long ago, this world was one of tumultuous and chaotic forces, naught but unbridled fire, seething oceans, and churning, gnashing rock. Through the ashen skies of Primordial Creation, the gods arrived from an unknown realm located beyond the ether. These ambitious divinities—young and still formless—looked down upon this roiling realm and saw potential for great strength and outstanding beauty, and the chance to learn their own place in creation.

From the divine hands of the The Arch Heart sprung forth the First Children, the elves. It is thought that they were made in their creator's image, yet it is just as likely that the The Arch Heart, and all the gods who shared their ideas with them, actually adopted forms inspired by the image of their creations. The First Children walked the lush Green and gazed upon the heavenly Blue as the gods tamed Primordial Creation around them.

Not long after, in the scale of the cosmos, a second mortal creation was wrought by the divine The All-Hammer: the dwarves, a hearty people imbued with the desire to continue the gods' work—taming a primordial world filled with the craft and invention of the divinity that lay beyond the ashen void. In their turn, a third people were given life by the passionate lovers of Order and Chaos, the The Lawbearer and the The Wildmother. Their children were the humans, endowed with hearts of passion to bring meaning to the world as it assumed an orderly form. Yet in their hearts burned a spark of chaos, cursing them with short life just as it blessed them with the unquenchable ambitions, mirth, and curiosities of the gods that crafted them.

As a note from the High Curator of Emon's Cobalt Reserve, the chronicle of Vasselheim states that the gods were given form and name only when the peoples of Exandria began to worship them—yet many of the various races of the world ascribe their creation to a single god or group of gods. It would be remiss to not at least conjecture that while the Protean Gods lacked distinct form, they still had unique thoughts, emotions, and motivations that were later codified through the worship of the people.

As inspiration flowed from the Protean Gods, other creations followed, giving life to Exandria's many peoples. Here the mythic record of Vasselheim grows unclear, for while the elves, dwarves, and humans were undeniably Exandria's first races, the identities of the many Children of Creation that followed were left unrecorded by the Dawn City's chroniclers.

These myriad young races explored their freshly founded world. As their knowledge grew, they attempted to build. But the land was fierce and treacherous, and the lives of the children were extinguished mercilessly by the still-chaotic world. Sorrow is said to have filled the hearts of the gods as they watched their children struggle against a land that rebuked them. The gods ultimately lent their children tiny fragments of their own power, gifts with which they might shape the world around them.

This small mercy of the gods was the birth of magic, and with it the people of Exandria learned to bend the angry earth to their will, to temper the wrath of its explosive mountains, to tame its unyielding floods, and to make the hard earth soft and bountiful. Wandering peoples settled, language became commonplace, survival gave way to art, and governance replaced anarchy. The Protean Creators, the divinity beyond the ashen skies, saw progress and saw that it was good—but fragile, and in need of guardians.

Thus were born the First Protectors: the Dragons Metallic of Tal'Dorei. These Protectors pledged themselves to defend the weak against tyrants. It is here that the mythic tales connect to the oldest, most fragmented passages of recorded history as civilization's reach stretched across the world. The people praised the gods for their newfound safety, and their prayers gave them forms, names, and purpose.

\subsection{Wrath of the Primordials}
Yet the realm refused to be tamed. Quaking cliffs roared in defiance of mortalkind's magical gifts. Seas swelled and seethed. Flames erupted from below with greater intensity than any had seen before. Beneath the elements, unknown to the Creators beyond the ashen skies, lived ancient beings who were born in the primeval chaos of the world the gods had found: the Primordials. These elemental titans that once dwelt deep within the land now rose from their unseen domain to sunder it once more. The gods watched as their children were crushed between gnashing rocks and fed to formless terrors unleashed in the wake of the destruction. Thus were the seeds of Exandria's greatest cataclysm set into motion.

Some gods were so full of grief and anger that they wished to abandon this world for another. They tried to convince their divine kindred to join the Primordials in reclaiming the realm for chaos, so that they could move on to start anew. Others wished to prove they could tame the world for their beloved children. This caused a divide among the gods. The schismatic deities, known today as the \textit{Betrayer Gods}, joined their songs and swords to those of chaos and destruction, taking and twisting their children into the image of their intent.

The cosmos itself seemed to recognize the magnitude of the battle on Exandria, and beings of Order and Chaos spilled forth from the still-young multiverse. Creatures of chaos emerged from a formless void that Exandria's peoples came to call the Abyss.

These demons fed upon the suffering of the wounded and the dying, and grew cruel. Sentinels of order emerged from a distant realm of celestial harmony. The \textit{Prime Deities} beseeched these celestials to aid them. In turn, the \textit{Betrayer Gods} forced order upon the demons and made them devils, perfectly lawful soldiers who lived only to cause harm.

As the clash of gods and Primordials spiraled out of control, the desperate mortal races sought their own way to defend themselves from the rebellious elements. Seeing that their creations still suffered, the The Arch Heart taught them the secrets of making magic themselves. The most intelligent and inquisitive mortal minds followed the path the gods had revealed to them, and created incantations and formulae to create new spells on their own terms, without the aid of divine power. These people called themselves arcanists, for they were masters of the arcane secrets of magic.

With their potent new magics, these resourceful mortals subdued the Primordials long enough for the \textit{Prime Deities} to banish their traitorous kin to secluded prison planes. With the \textit{Betrayer Gods} locked away, the gods were overjoyed to see the resourcefulness of their children. The defeated Primordials were destroyed, and from their ever-enduring essences, the gods channeled their wild and destructive powers into the Elemental Planes that surround the material world of Exandria.

For the first time since the Creation, Exandria finally settled into an idyllic peace. From this peace was born the first enduring mortal civilization, the heart of which was a grand city called Vasselheim, the Cradle of Creation and the Dawn City.

Culture developed anew, the races ventured beyond to explore and discover their own lands, and great music filled the air to give name to this world once and for all: Exandria.

\subsection{Age of Arcanum}
\DndQuote%
{}
{\enquote{The foremost scholars of our day would make it seem an indisputable fact that the gods created Exandria. That they blessed us with the bounty that all peoples of the world now make use of. So too, they say, was magic their gift to us. But such zealots ignore the fact that there is another likelihood: that the gods did not create mortals, and their world—but the opposite. That the dreams of mortals mingled in the wild magic of the young world, and gave that magic form and order: divinity. Of course, the dogmatic zealots of our day will dismiss such ideas as heresy—but not you, reader! The avid mind feeds not on dogma, but on new ideas!}}
{Anonymous}

Ages passed, and society flourished. Great kingdoms sprung up. Castles were built in a day, accelerated by the arcanists' newfound power. Even though magic could be used to complete the most difficult tasks with hitherto unknown speed, magic-users strove always to innovate. As mages practiced and perfected their powers of creation, they soon unlocked the secrets of life itself, giving birth to wondrous, dangerous new forms of life and power.

Yet, with the benefit of hindsight, it is clear that the very magic which allowed mortalkind to prosper also instilled a deep rot within their civilizations. The arcanists of the world grew arrogant. They came to see their arcane gifts as proof that the gods held no sway over their fate, and that with a sophisticated enough command of magic, they could become as powerful as the gods themselves. Though this hurt and surprised the \textit{Prime Deities} who had imprisoned their kin for the sake of their children, they sought to sympathize with the willfulness of their creations, remaining out of love and hoping that the mortals would learn the error of their ways.

The advent of the arcane seemed to be the key to a bountiful age of plenty, but also proved to threaten it, as prosperity soon gave way to greed. Rather than universal goodwill and comfort for all, the \textit{Age of Arcanum} came to be defined by endless petty squabbles over wealth and power among the elite, while those without magic hunted for their scraps. Tantalizing rumors of a path to immortality slithered through the most decadent circles of magi and nobility. One mortal mage, her name lost and presumed to be struck from history, crafted now-forbidden rites to challenge the God of Death, felling him and taking his place among the pantheon, making of her the first and only living mortal to ascend. The name of this original death god was likewise lost to history. Only the title of his now-godly successor, the The Matron of Ravens, survives.

Her victory over divinity was a catalyst for many of the horrors of the \textit{Age of Arcanum}. An archmage named Vespin Chloras, renowned throughout ancient Vasselheim for his wealth, skill, and cruelty, was inspired by this display. He sought the guidance and power of the banished gods, rending open the gates of their prisons and releasing the Betrayers into the mortal world.

In their imprisonment, these Gods of Hatred and Despair had warped their prisons into reflections of their depravity. The Abyss, once a formless place of chaos, a byproduct of the gods' Creation, was twisted into a place of evil. Countless other planes of cruelty were born from the \textit{Betrayer Gods'} hatred, and those realms' evil depths endlessly churned forth horrors that lived only to transform peace into suffering, and righteousness into arrogance and greed.

Released unto the mortal world once more, the \textit{Betrayer Gods'} urge to ruin was supplanted by the desire to dominate, and they turned their sights first to the archmage who had freed them, making Vespin the first of their many mortal thralls. Records of diabolical texts safeguarded within the libraries of the \textit{Alabaster Lyceum} suggest that Vespin, now long dead, serves the \textit{Betrayer Gods} to this day—now as a devil at the left hand of the The Lord of the Hells. These corrupt gods sought out the shattered remnants of the devils they had crafted and the demons they conscripted, and secretly created with them a fearsome new kingdom on the far end of the world, the capital of which was named Ghor Dranas, the Gathering of Shadows.

In this land, where the twisted power of the \textit{Betrayer Gods'} corrupted planes seeped into the world, the lords of evil welcomed mortals whose heartlessness and lust for power made them susceptible to the gods' grand promises. The most fertile soil of all for these noxious seeds lay in mortal hearts obsessed with the infinite power of the arcane. With a legion of the damned behind them, the \textit{Betrayer Gods} soon made their presence known to the world with a surprise assault on Vasselheim itself.

Though much of the city was reduced to rubble, Vasselheim weathered the initial assault, saved by the intervention of their protectors, the metallic dragons, and even a number of the \textit{Prime Deities} themselves. The creators descended to trade blows with their former brethren. The battle, which pitted gods against mortals and heroes against demons, raged ceaselessly for twenty days and nights, until the dark forces were forced to retreat, their surprise attack thwarted.

Yet this victory was a tainted one. Evil had been repulsed momentarily, but the revelation of such a terrible foe incited an arcane arms race. Trust was shattered indefinitely: if mortals could fall under the sway of the \textit{Betrayer Gods}, who could be true allies? If ruin like this could be unleashed under the watchful eyes of divinity, what value did mortal lives hold? Fearing all powers but their own, the most self-interested and singular human arcanists warped their greatest creations, magical instruments of prosperity and joy, into arms and armor of horrific power. The \textit{dwarves'} fascination with rock and earth turned toward isolation as they burrowed further into the mountains, using their divine gifts to animate legions of autonomous golems to protect their ancestral halls. elves used their understanding of creation's beauty and intricacies to weave spells of unimaginable destructive force, the likes of which Exandria had never before seen.

For the first time since the awakening of the Primordials, the focus of magic was warfare. The gods themselves agreed to join their children on the field of battle, descending from the heavens to take up arms once more for the war now referred to as the \textit{Calamity}.

\section{The Calamity}
Few records remain of the terrible war that followed, but its effects are still felt today. Most tales that remain are of great champions who the \textit{Prime Deities} blessed with their power. Some deities created warriors of their own from light and holy fire. Others imbued fragments of their power into the weapons now called \textit{Vestiges of Divergence}. More still bestowed blessings upon mortals and called them champions. One tale even tells of a champion blessed by three of the gods in their times of greatest need. Now all but forgotten, this Apotheon was forced by fate into terrible battles across the world.

The sheer magnitude of the energies unleashed in the ensuing battles of gods and mortals frayed the boundaries holding back the elemental chaos, spilling unbridled destruction into the world and bringing utter annihilation. Even the ley lines that direct the flow of magic across Exandria like veins direct blood through the body show signs of being warped by the raw power that the Calamity unleashed. Some well-preserved ruins of the ancient civilizations remain, such as the Drowned City of Cael Morrow in the lands of \textit{Marquet}, but for the most part, all traces of the old world were erased from the face of Exandria.

So great was the loss of life during the war that even liberal estimates suggest no more than a third of Exandria's population survived it. This unthinkable devastation inspired some civilizations to flee Exandria entirely and seek refuge on other planes, scattering the Children of Creation across the multiverse. By the war's end, the only bastion of civilization left on the face of Exandria was the Dawn City itself: Vasselheim.

\section{The Divergence}
In the wake of the \textit{Calamity}, the victorious \textit{Prime Deities} once more banished their traitorous kin to their realms of deception and hate. The world entered a long, dark period of recovery, as history had to be recovered and purpose restored, and the threat of the \textit{Betrayer Gods} still loomed heavily upon the minds of all. Even the \textit{Prime Deities} felt guilt for their role in the conflict, for it was the unrestrained clash of divine power that had unleashed such horror upon the world.

The records of the Scalebearers assert that the The Platinum Dragon and the The Lawbearer descended upon Vasselheim and spoke a decision: the \textit{Prime Deities} would depart from the world and establish a Divine Gate that would forever prevent any god from ever acting directly upon Exandria again. This proclamation shocked the other \textit{Prime Deities} and the beleaguered survivors of the war in equal measure, but while mortals railed against this announcement, the other gods quickly realized that the decision of the two most dutiful and self-sacrificing of their number was unimpeachable and just.

Thus, in hopes of ensuring such ruin would not befall Exandria again, they left their children to rebuild civilization anew within the walls of Vasselheim and beyond. The Creators returned to their own realms, sealing all divine powers behind their newly constructed Divine Gate. Only in this way could they prevent their corrupted brethren from physically returning to the Material Plane.

Much time has passed since, and the world has been reborn once again. The gods still exhibit their influence and guidance from beyond the Divine Gate, bestowing their knowledge and power to their most devout worshipers, but the path of mortals is now their own to make. New cities, kingdoms, and cultures have retaken the world, building over the ashes of the old. New songs fill the air, and the hope of a brighter future drives people day after day, while buried ruins and forgotten relics remind all people of a darker time—and of mistakes that should never be repeated.

\begin{DndSidebar}[float=!b]{Names of the Divergence}
Of these many formative events, the tale of the \textit{Divergence} is the most oft-told among the many faiths and civilizations of the world. When seeking the truth of this event, the researchers of the \textit{Cobalt Soul} found that it was called by many titles around the world, for many different reasons.




\begin{itemize}
\item \textbf{The Second Spark.} This name is most often used by those who study the arcane, for it was the gods giving them a second chance to use magic for the betterment of the world.
\item \textbf{The Penance.} Priests, clerics, and religious scholars across Exandria refer to the departure of the gods by this name, for they see it as a self-imposed atonement for the gods' role in the ravaging of the world.
\item \textbf{The Divergence.} The most common name for this event, the \textit{Divergence}, is used by people of all classes and creeds. This descriptive name simply refers to the gods' final act of divergence from the mortal world.
\end{itemize}



\end{DndSidebar}

\section{Pantheon of Exandria}
During the \textit{Divergence}, the \textit{Prime Deities} created a powerful barrier between the Material Plane and the supernatural realms of the Outer Planes. This gate sealed all gods, including the evil \textit{Betrayer Gods}, within their respective domains. If the Divine Gate were to be destroyed by the unanimous effort of the \textit{Prime Deities}, all powers would be unleashed, threatening armageddon. Thus, the gods patiently watch their creations from beyond the veil, aiding their faithful with what small power they can send through the barrier.

Some lesser idols, such as archdevils, demon lords, fiends and celestials of near-divine power, as well as a number of unaccounted-for demigods reside among the planes. They are affected by the Divine Gate, preventing them from crossing it of their own volition, yet are not so bound by it that they cannot be summoned by mortal magic.

The following gods are recommended as the existing Exandrian pantheon, but they are only a recommendation. As long as the GM and the players agree, you are welcome to tailor, alter, or completely change the gods in your Exandria campaign to fit your needs. The domains listed are the likely choices for followers of that deity, but these are not the only options.

\subsection{Divine Domains}
Most clerics can gain their magic from a single god, or even from a number of gods that share domains or divine provinces. Ultimately, how your cleric serves the gods or their core religious philosophies is for you and your Game Master to discuss. Each god in this chapter has two or three Divine Domains from which many of their clerics are granted power.

Nevertheless, almost all deities can have clerics who gain the power of any domain, if you and your Game Master agree.

\subsection{Prime Deities}
The Prime Deities are the leaders and luminary creators that battled the Primordial Titans and instigated the \textit{Founding}, forging the mortal races of Exandria. They hold sway over powers that represent a spectrum of sunlight, truth, moonlight, benevolent concealment, protection, love, death, and all other facets of freedom and life in the world. While these gods may disagree and squabble, they are united in an alliance dedicated to ensuring the survival of their creations.

\begin{DndTable}[header=Prime Deities]{XlX}

Deity & Alignment & Province\\
The Changebringer & Chaotic good & Change, freedom, luck\\
The Platinum Dragon & Lawful good & Honor, justice\\
The Arch Heart & Chaotic good & Art, beauty, elves\\
The Lawbearer & Lawful neutral & Civilization, law, peace\\
The Knowing Mentor & Neutral & Knowledge, learning, teaching\\
The Stormlord & Chaotic neutral & Battle, competition, storms\\
The Wildmother & Neutral & Seas, wilderness\\
The All-Hammer & Lawful good & Craft, creation\\
The Dawnfather & Neutral good & Healing, sun\\
The Everlight & Neutral good & Atonement, compassion\\
The Matron of Ravens & Lawful neutral & Death, fate, winter\\
The Moonweaver & Chaotic good & Illusion, moonlight, night\\
\end{DndTable}

\subsubsection{The Changebringer}
Also known as "She Who Makes the Path," the Changebringer champions freedom, travel, trade, and adventure across the lands. Her will heralds open frontiers, and her call beckons her followers to discover that which awaits them beyond the known. There are few urban temples dedicated to the Changebringer, but her shrines are often found along well-traveled roads and on isolated wilderness hill-tops. Many merchants, adventurers, and free-spirited folk worship the Changebringer, and many a gala toast and tavern song celebrate her as a bringer of luck and fortune.

\subparagraph{Depiction}
She is often depicted as a young woman of dark complexion and long, light brown hair that cascades to form the road left behind her. Most art shows her in constant motion, ever leading into the unknown.

\subparagraph{Divine Domains}
Many of the Changebringer's clerics are granted powers of Moon, Nature, or Trickery.

\subparagraph{Holy Day}
The Changebringer's holy day, called New Dawn, takes place on the first day of the first month, when the old year gives way to the new. In \textit{Emon}, the day of New Dawn is celebrated with a grand midnight feast that commonly features a short play celebrating the changes of the past year.

\subparagraph{Symbol}
A holy symbol of the Changebringer that shows the profile of a beautiful woman on a gold coin. She bears an intricate and ornate headpiece.
\pagebreak
{\par \vspace { 9pt plus 3pt minus 3pt } \noindent  \bfseries \dndsanscaps{Commandments of the Changebringer} \nopagebreak }
\begin{itemize}
\item Luck favors the bold. Your fate is your own to grasp, and to do so is to have the Changebringer behind you.
\item Change is inevitable. The righteous can ensure that such change is for the better.
\item Rise up against tyranny. Fight for the freedom of yourself and others when you can, and inspire others to fight when you cannot.
\end{itemize}


\subsubsection{The Platinum Dragon}
The pillar of justice, protection, nobility, and honor, the Platinum Dragon is a beacon to paladins of order and good, and is revered by most metallic dragons as the first of their kind. The crest of the Platinum Dragon adorns many halls of high leadership and judgment, invoking his will in all matters of justice. To follow him is to look after those who cannot look after themselves.

\subparagraph{Depiction}
The Platinum Dragon is often seen emblazoned on shields and armor, both functional and decorative, in the form of a brilliant silver dragon head in profile. Temples and works of art depict a massive, glittering dragon with vibrant platinum scales and a wingspan that fades into etherealness.

\subparagraph{Divine Domains}
Many of the Platinum Dragon's clerics are granted powers of Life, Order, or War.

\subparagraph{Holy Day}
The Platinum Dragon's holy day is called Embertide and is celebrated on the fifth day of the eleventh month. This is a day of remembrance, solemnity, and respect for those who have fallen in the defense of others.

\subparagraph{Symbol}
A dark blue heraldry of the profile of a platinum dragon.

{\par \vspace { 9pt plus 3pt minus 3pt } \noindent  \bfseries \dndsanscaps{Commandments of the Platinum Dragon} \nopagebreak }

\begin{itemize}
\item Stand as a paragon of honor and justice.
\item Smite evil wherever it is found, yet show compassion to those who have strayed from righteousness.
\item Defend the weak, bring freedom to the persecuted, and protect the ideals of justice and order.
\end{itemize}


\subsubsection{The Arch Heart}
Guardian of spring, beauty, and the arts, the Arch Heart is the patron of arcane magic, elves, and the fey. The \textit{Founding} inspired them to wander the twisted lands, seeding them with the first arcane magics and raising the most ancient of forests. It was by the Arch Heart's hand that the first elves wandered from the Fey Realm, and for this reason, they are considered the Parent of all elves. Those who seek art in all of their work, whether magic or mundane, often worship at the altar of the Arch Heart. They loathe the The Spider Queen and her priestesses for leading the drow astray.

\subparagraph{Depiction}
Most modern tapestries and tomes depict the Arch Heart as a genderless elf of impossible grace and beauty, androgynous and alluring, framed by long, wavy, golden hair. The Arch Heart inspired many early elven art pieces, and elements of their visage or symbol are included in most elven architecture.

\subparagraph{Divine Domains}
Many of the Arch Heart's clerics are granted powers of Arcana, Nature, or Light.

\subparagraph{Holy Day}
The Arch Heart's holy day is called Elvendawn, or Midsummer. It's celebrated on the twentieth day of the sixth month, and commemorates the \textit{elves'} first emergence from the Fey Realm. In \textit{Syngorn}, the elves build magical wards, open small doorways into the Fey Realm, and celebrate with uncharacteristic vigor with the wild fey.

\subparagraph{Symbol}
A silver holy symbol of two, blue crescent moons atop a four-pointed star with intricate elvish designs.

{\par \vspace { 9pt plus 3pt minus 3pt } \noindent  \bfseries \dndsanscaps{Commandments of the Arch Heart} \nopagebreak }

\begin{itemize}
\item Create, inspire, and find beauty in all that you do.
\item Follow the echoes of lost magic, forgotten sites,and ancient art, for within these lie the ArchHeart's first works.
\item Combat the followers of the The Spider Queen wherever they may be.
\end{itemize}


\subsubsection{The Lawbearer}
The driving inspiration behind many great inventions, the creation of vast cities, and law and order within society, the Lawbearer claims dominion over civilization. Judges and lawful rulers pay respect at her temples, which are central structures in major cities around the world. Peace and order—through structure and law—guide the will of her devout followers. The Lawbearer has a tempestuous romance with the The Wildmother, a furious love that is only tempered when civilization and nature are in balance. Civil unrest and riots blossom within her beloved cities in these times of imbalance.

\subparagraph{Depiction}
Illustrations and statues show the Lawbearer as a hooded, armored woman with light brown skin and a strong, lithe body, sitting atop a throne of pillars. Her face is generally obscured or depicted without expression, giving her presence an impartial yet imposing nature.

\subparagraph{Divine Domains}
Many of the Lawbearer's clerics are granted powers of Knowledge or Order.

\subparagraph{Holy Day}
The Lawbearer's holy day is Civilization's Dawn, and is celebrated on the autumnal equinox, usually the twenty-second day of the ninth month. In \textit{Emon}, each neighborhood celebrates by dancing around great bonfires in the square and giving gifts to celebrate their community.

\subparagraph{Symbol}
A double axe with golden scales engraved onto the blade.

{\par \vspace { 9pt plus 3pt minus 3pt } \noindent  \bfseries \dndsanscaps{Commandments of the Lawbearer} \nopagebreak }

\begin{itemize}
\item Utilize the company and aid of others. The efforts of the individual often pale against the capabilities of community.
\item Strive to tame the wilds in the name of civilization,and defend the points of light and order against the chaos of darkness.
\item Uphold and revere the spirit of invention. Create new settlements, build where inspiration strikes,and expand the edicts of the Lawbearer.
\end{itemize}


\subsubsection{The Knowing Mentor}
Revered by seers, sages, and teachers of all walks of life, the Knowing Mentor guided the growth of civilization throughout the \textit{Age of Arcanum} like sunlight guides the branches of a tree. Her teachings were silenced when she was grievously wounded by the The Chained Oblivion during the \textit{Calamity}, and though she survived and went into hiding behind the Divine Gate, her followers are hunted by agents of her ancient foes to this day. Her devout now worship in private, spreading knowledge, philosophy, and lore anonymously through traceless channels.

The one exception are the scholars of the \textit{Cobalt Soul}, an international organization of knowledge-seekers, spies, and archivists that seek to reveal the truth at all costs. Most members of the \textit{Cobalt Soul} aren't zealous worshipers of the Knowing Mentor, but still pay homage to her in some form.

\subparagraph{Depiction}
Common representations show the Knowing Mentor as a greying, mature woman of welcoming, matronly smile swathed in billowing robes and scarves that fan into books and scrolls. Some call her the Knowing Mistress, revering her as a headmistress of knowledge among the gods.

\subparagraph{Divine Domains}
Many of the Knowing Mentor's clerics are granted powers of Arcana, Knowledge, or Twilight.

\subparagraph{Holy Day}
The Knowing Mentor has no public holy day, for her public worship was shattered during the \textit{Calamity}, and she has since fallen into half-remembered myth. Only in \textit{Emon} has the Knowing Mentor's faith been publicly resurrected—though her worship by the intellectuals of the city bears little resemblance to that of the knowledge-seekers of old.

\subparagraph{Symbol}
A holy symbol that is a mossy stone with faded runes in the shape of a face with three eyes.
{\par \vspace { 9pt plus 3pt minus 3pt } \noindent  \bfseries \dndsanscaps{Commandments of the Knowing Mentor} \nopagebreak }

\begin{itemize}
\item Unmask those who would destroy the Knowing Mentor. Learn their secrets and unveil them to the world.
\item Uphold and teach the importance of reason, perception, and truth in guiding one's emotions and path.
\item Condemn those who lie without moral cause, for evil folk gain power when their followers obscure the truth. Never stoop to the level of selfish liars.
\end{itemize}


\subsubsection{The Stormlord}
Where thunder cracks and conflict rises, prayers to the Stormlord are shouted into the maelstrom. Reveling in all tests of strength, the blessing of the Stormlord finds those that prove themselves on the battlefield. Worshiped by athletes and warriors all across Exandria, he exalts those whose force of spirit and passion for victory call his attention. He brings tumultuous storms over land and sea, and those who wish for clearer skies offer their praises and prayers to appease him.

\subparagraph{Depiction}
Within his temples and carvings, the Stormlord is shown as a quintessential warrior, often nude with a beard and short, curly hair. As the epitome of muscle and strength, most art depicts the Stormlord in a stance of dominance, typically asserting his power by wrestling a terrible beast.

\subparagraph{Divine Domains}
Many of the Stormlord's clerics are granted powers of Tempest or War.

\subparagraph{Holy Day}
The Stormlord's holy day is the Day of Challenging, and it is celebrated on the seventh day of the second month. The Day of Challenging is one of the most raucous holidays in \textit{Emon}, and thousands of spectators attend the annual \textit{Godsbrawl} held in the Temple of the Stormlord to root for their favored deity's champion, particularly the chosen of the Stormlord and the Platinum Dragon.

\subparagraph{Symbol}
A round holy symbol that appears like a shield with four blue strikes of lightning crossing under a red diamond in the centre.
{\par \vspace { 9pt plus 3pt minus 3pt } \noindent  \bfseries \dndsanscaps{Commandments of the Stormlord} \nopagebreak }

\begin{itemize}
\item Bravery above all. There is no glory in cowardice.
\item Strength is the path to greatness, but greatness is the responsible use of strength.
\item The glory of the Stormlord lives through your own glory on the battlefield.
\end{itemize}


\subsubsection{The Wildmother}
The realm of the Wildmother extends to wherever the seas shift and the land grows over. As the keeper of the wilderness, she represents the wild creatures of nature, the rush of the angry rapids, and the heat-heavy stillness of the desert. Hunters, druids, and nomads worship her, seeking her blessing to live in harmony with lands that still bear their wild fangs. Likewise, those seeking safe passage across dangerous waters pray to her to guide them.

The druids of the Wildmother and clerics of her lover, the The Lawbearer, work together to preserve the balance of nature and civilization. It is said that the two goddesses grow furious when this balance is upset, and this fury manifests as devastating natural disasters. Clerics of the The Matron of Ravens also pay homage to the Wildmother, for the work she begins is the work that their deity must respectfully end.

\subparagraph{Depiction}
Immortalized through wooden reliefs and carved idols in hidden, overgrown groves and rural shrines, the Wildmother is shown as a beautiful woman with green or deep brown skin nearly swallowed by a wild, tangled wreath of hair, leaves, and vines that envelop her strong, full-figured form.

\subparagraph{Divine Domains}
Many of the Wildmother's clerics are granted powers of Life, Nature, or Tempest.

\subparagraph{Holy Day}
The Wildmother's holy day is Wild's Grandeur, and is celebrated on the vernal equinox, usually the twentieth day of the third month. The people of the southern wilds honor the Wildmother by journeying to a place of great natural beauty, like the top of a mountain waterfall or the center of a desert. Wild's Grandeur is rarely celebrated in \textit{Emon}, but some folk will plant trees in observance of the holiday.

\subparagraph{Symbol}
A wreath of golden wheat with green leaves and red berries sprouting out of it. A wooden shepherd's staff is in the centre of the wreath, entwined with some vines.
{\par \vspace { 9pt plus 3pt minus 3pt } \noindent  \bfseries \dndsanscaps{Commandments of the Wildmother} \nopagebreak }

\begin{itemize}
\item Protect the untamed wilderness from exploitation and destruction.
\item Slay abominations and other mockeries of nature.
\item Embrace and respect that which you cannot control in this world. Exist in harmony with it.
\end{itemize}


\subsubsection{The All-Hammer}
The patron of craft and creation, the All-Hammer is worshipped by smiths, artisans, and miners alike, granting inspiration where respect and prayer are given. He shaped the mountains from the chaos of the \textit{Founding}, and stands as the patron protector of home and family. Devotion to the All-Hammer is strongest in dwarven communities, and many of his temples mark the center of a mighty dwarven stronghold.

\subparagraph{Depiction}
Many guildhalls and workshops contain images of the All-Hammer, a faceless, stout dwarf of immense strength, hunched over a flaming heart clasped within his massive hands. The All-Hammer's dwarven appearance is consistent across Exandria, with some exceptions. Other cultures that pride themselves on their metallurgy, like the fire giant enclave of \textit{Vulkanon}, depict the All-Hammer as one of their own—in this case, a faceless giant with a flaming mane.

\subparagraph{Divine Domains}
Many of the All-Hammer's clerics are granted powers of Forge, Knowledge, or War.

\subparagraph{Holy Day}
The All-Hammer's holy day is Deep Solace, celebrated on the eighteenth day of the fifth month. Especially devout followers of the All-Hammer spend the day in isolation, meditating on the meaning of family and how they may be better mothers, fathers, siblings, and children. Dwarven communities in \textit{Kraghammer} and abroad celebrate with a full day of feasting and drinking.

\subparagraph{Symbol}
A holy symbol shaped like two hammers, back to back, and in profile, with a rectangular ruby in the center between them. The face of each hammer is shaped to look like the head of a bearded dwarf.

{\par \vspace { 9pt plus 3pt minus 3pt } \noindent  \bfseries \dndsanscaps{Commandments of The All-Hammer}  }
\begin{itemize}
\item Remain stoic and tenacious in the face of catastrophe.
\item Uphold and promote loyalty to your family, loyalty to your clan, and loyalty to your people.
\item Legacy is paramount. To create something that lasts is to change the world for the better.
\end{itemize}

\subsubsection{The Dawnfather}
The Dawnfather rules over sun and summer, and his vigil encircles the ages as the keeper of time. As the lord of agriculture and harbinger of the harvest, he is worshiped by farmers and most common folk, and his priests are welcome in lands all over. Supporter of the needy and destroyer of evil, the Dawnfather is often the patron of paladins and rangers who follow a similar creed. The Dawnfather is mythically renowned for his defeat of the The Chained Oblivion and is revered by those who hunt aberrations.

\subparagraph{Depiction}
Tapestries of old match early texts describing the Dawnfather as a kindly patriarch clad in silver and gold armor, his head a beacon of light and fire so bright that his face can barely be seen. Many statues in holy places use the Dawnfather's head as a brazier, lit with each dawn and extinguished with the dusk.

\subparagraph{Divine Domains}
Many of the Dawnfather's clerics are granted powers of Life, Light, or Nature.

\subparagraph{Holy Day}
The Dawnfather's holy day is called Highsummer and takes place on the fifteenth day of the seventh month. In \textit{Emon}, the entire week is celebrated with gift-giving and feasting. Festivities begin on Highsummer day and end at midnight on the twenty-first, the day that \textit{Zan Tal'Dorei} dethroned \textit{Trist Drassig} at the \textit{Battle of the Umbra Hills}. In \textit{Whitestone}, the feast of Highsummer is celebrated with gifts and a festival of lights around the \textit{Sun Tree}. Most folk choose to spend Highsummer with their family, recounting the small things they are thankful for.

\subparagraph{Symbol}
A circular holy symbol made of stained glass depicting a yellow sun with eight silver rays emanating from it.
{\par \vspace { 9pt plus 3pt minus 3pt } \noindent  \bfseries \dndsanscaps{Commandments of the Dawnfather} \nopagebreak }

\begin{itemize}
\item Be ever vigilant for evil. People are quick to forget the lessons of the past.
\item Help relieve the suffering of the innocent.
\item Deliver the light of the Dawnfather where darkness dwells, and do so with kindness, compassion, and mercy.
\end{itemize}


\subsubsection{The Everlight}
The Everlight, or Raei, is an adapted form of Sarenrae, a Paizo Product Identity.

God of compassion and redemption, the Everlight bears a divine message of determination and compassion, even in the darkest of times. She believes that the corrupt can be redeemed, a mindset that led to a betrayal by the The Lord of the Hells, who decimated her followers during the \textit{Calamity}. Only recently has her faith been rediscovered and her temples returned to prominence. The Everlight's followers are often rural healers and community philosophers, offering voices of reason and empathy in angry and cynical times.

\subparagraph{Depiction}
Those who rise to bring the Everlight's words back to the light either unearth her image from ruined temples or create new art to inspire others with her message. She is represented as a beautiful, strong woman with dark skin and light hair, rising betwixt a set of angelic ivory wings.

\subparagraph{Divine Domains}
Many of the Everlight's clerics are granted powers of Life, Light, or Peace.

\subparagraph{Holy Day}
The Everlight's holy day has long been forgotten, and her followers have yet to decide when her festival should be held, or even what the festivities should be. The debate has gone on for years, and no great miracles have yet been performed to unify the squabbling clerics.

\subparagraph{Symbol}
A holy symbol of a feminine, silver phoenix atop a yellow flame, her wings encircling a red ruby overhead.
{\par \vspace { 9pt plus 3pt minus 3pt } \noindent  \bfseries \dndsanscaps{Commandments of the Everlight} \nopagebreak }

\begin{itemize}
\item Lead with mercy, patience, and compassion. Inspire others to unite in fellowship.
\item Aid those who are without guide. Heal those who are without hope.
\item Those who are beyond redemption, who revel in slaughter and remorseless evil, must be dispatched with swift justice.
\end{itemize}


\subsubsection{The Moonweaver}
The Moonweaver is the god of moonlight and the autumn season, as well as the patron of illusion and misdirection. Widely worshipped in halfling and elven culture, she is also considered to be the deity of love, shielding the trysts of lovers in shadows of her own making. Those who work in darkness and trickery often ask for her blessing.

\subparagraph{Depiction}
Depictions of the Moonweaver are as numerous as the myths and stories of her meddling within simple unions. However, she is most often painted as a young girl of light-blue skin and white hair, with a body and limbs that dissolve into silky strands of silver moonlight, caressing and creating the edges of the shadows.

\subparagraph{Divine Domains}
Many of the Moonweaver's clerics are granted powers of Arcana, \textit{Moon}, or Twilight.

\subparagraph{Holy Day}
The Moonweaver has no holy day, but is celebrated by most elves on the night of the decade's largest full moon. Elven astronomers track the moon's phases and how it grows closer and farther to Exandria, and can predict these days with great accuracy. Many high and haughty elves use this festival as an excuse to be sly and mischievous; some younger elves use costumes and illusions to prank their peers. The \textit{wyvern-riding skyswimmers of Lyrengorn}, however, venerate the Moonweaver in their own way: with the annual arrival of the \textit{Moonweaver's Ribbons}, a majestic aurora that illuminates the skies above the \textit{Elvenpeaks}.


\subparagraph{Symbol}
A circular holy symbol with a crescent moon outlining the bottom of it, it's face upwards towards a starry, purple and white-blue borealis within. Two arrows curve over the top and bottom of the circle.
\pagebreak
{\par \vspace { 9pt plus 3pt minus 3pt } \noindent  \bfseries \dndsanscaps{Commandments of the Moonweaver} \nopagebreak }

\begin{itemize}
\item Seize your own destiny and pursue your own passions.
\item Let the shadows protect you from the burning light of fanaticism and the absolute darkness of despair.
\item Walk unbridled and untethered, forging new memories and experiences.
\end{itemize}

\subsection{Betrayer Gods}
The \textit{Betrayer Gods} are the deities who strayed from the ideals of the \textit{Founding} and embraced the destructive chaos of the Primordials. The \textit{Betrayer Gods} rarely work together, for they see one another as threats to their own goals. This very weakness allowed the \textit{Prime Deities} to defeat and banish them, ending the \textit{Calamity}.

\begin{DndTable}[header=Betrayer Gods]{XlX}

Deity & Alignment & Province\\
The Lord of the Hells & Lawful evil & God of devils and the Hells\\
The Strife Emperor & Lawful evil & Conquest, tyranny\\
The Ruiner & Chaotic evil & Slaughter, warfare\\
The Spider Queen & Chaotic evil & Deceit, spiders\\
The Chained Oblivion & Chaotic evil & Darkness, destruction\\
The Scaled Tyrant & Lawful evil & Dragon god of evil\\
The Crawling King & Neutral evil & Enslavement, torture\\
The Whispered One & Neutral evil & Necromancy, secrets\\
The Cloaked Serpent & Chaotic evil & Assassins, poison, snakes\\
\end{DndTable}

\subsubsection{The Lord of the Hells}
The devil god of the Hells is the master of tyranny and domination, and was the first to instill his despotic will into celestials, creating the first devils. His words are honeyed and carefully crafted, soothing and corrupting the hearts of mortals and immortals alike. The Lord of the Hells rules his domain with an iron fist, and the punishments awaiting those that cross him are the stuff of nightmares. Evil entities pay him tribute right alongside his devils, and many warlocks are drawn to his power.

\subparagraph{Depiction}
A twisted image of the celestial blood that once bore him, the Lord of the Hells is portrayed in many tomes and murals as a handsome humanoid with deep red skin and long black hair. Two curling horns rise from his crowned brow, and his lips bear an eternal, knowing grin.

\subparagraph{Divine Domains}
Many of the Lord of the Hells' clerics are granted powers of Order, Trickery, or War.

\subparagraph{Enemies}
The Lord of the Hells' enemies are numerous, even among the \textit{Betrayer Gods}, many of whom only follow him for fear of his immense power. His arch-nemesis is the The Changebringer, whose mischief and cunning has vexed him throughout the eons. It was she who defeated him in the \textit{Calamity} by tricking his fiendish armies into attacking one another.

The Devil Lord's greatest triumph during the \textit{Calamity} was fooling and betraying the The Everlight and slaughtering all her followers in one fell stroke—a move that today has united the followers of the The Everlight and the The Changebringer in a bond of unbreakable fellowship.

\subparagraph{Symbol}
A black crown with four red horns ascending from the top of it and two black and red horns curving out on the sides. Five black gems are embossed into the crown, with one lone black horn rising from the center.
{\par \vspace { 9pt plus 3pt minus 3pt } \noindent \bfseries \dndsanscaps{Commandments of the Lord of the Hells} \nopagebreak }
\begin{itemize}
\item Assert dominance and power over others. Show your strength of will in the image of your Lord.
\item Repay cruelty done unto you with further evil. If others show you kindness, exploit it.
\item As you ascend to power, do not pity or show mercy to those you climb over to get there. The weak do not warrant your compassion—compassion itself is a weakness.
\end{itemize}


\subsubsection{The Strife Emperor}
Blood-drenched armies of brutal warriors oft crush their foes in the name of the Strife Emperor, the patron of warfare and conquest. To serve his will is to accept the call to conflict, seeking lesser people to break and subjugate. Warmongers worship the Strife Emperor as they strike out at the world to bend it beneath them.

The Strife Emperor twists all living things to his iron will, even forcing nature itself to bow to his whims. During the \textit{Calamity}, he grew his armies not by rallying demons, but by corrupting noble creatures into monsters and by destroying the wilds to fuel his engines of war.

\subparagraph{Depiction}
The Strife Emperor is often depicted as a brutish, ogre-like man clad from head to toe in jagged black armor. Heads dangle from his belt, and the shadows that obscure his helmeted face do not hide his unblinking yellow eyes.

Those who long for conquest often see their own people as superior to all others. Thus, somewhat appropriately, depictions of the Strife Emperor are contentious and myriad: he is usually portrayed as a member of whatever race currently bears his banner. In the Republic of Tal'Dorei, his face most often has a human cast, for \textit{King Drassig's} bloodline paid him tribute long ago.

\subparagraph{Divine Domains}
Many of the Strife Emperor's clerics are granted powers of the Forge, Order, or War.

\subparagraph{Enemies}
The Strife Emperor's greatest enemy is the The Wildmother, who defeated him in \textit{Rifenmist} during the \textit{Calamity}. In turn, he is also despised by her beloved, the The Lawbearer, for perverting the benevolent rule of law into the manacles of authoritarianism. The Strife Emperor and the The Wildmother clashed many times during the war before his final defeat at \textit{Beynsfal Plateau}. His massive iron armor is still scattered across Beynsfal, and his gigantic helmet has become the seat of a ruthless empire.

\subparagraph{Symbol}
A raised armoured fist holding a copper chain that wraps around it.
{\par \vspace { 9pt plus 3pt minus 3pt } \noindent  \bfseries \dndsanscaps{Commandments of the Strife Emperor} \nopagebreak }

\begin{itemize}
\item Fear is your ally. Conquer your fear, and inspire it in your foes.
\item Disorder and rebellion are to be punished severely.
\item To kill is life's greatest pleasure, and perfecting the art of slaughter is life's greatest pursuit.
\end{itemize}

The Strife Emperor is also known as the progenitor of the Curse of Strife.


\subsubsection{The Ruiner}
The Ruiner commands roving hordes of barbaric marauders across the lands to destroy, pillage, and slaughter for the joy of it. Orderless and without honor, the creeds of the evil lord goad all living creatures to shrug off the oppressive chains of courtesy and civilization and to devour the world around them, giving in to the chaotic and selfish nature of the predator.

As civilization rises, people of all cultures feel the draw of the uncomplicated, primeval world. Some find their way to the loving embrace of the The Wildmother and her sects, going on to live in harmony with nature. Others fall into the cruel thrall of the Ruiner and his bloody cults of slaughter. In Tal'Dorei, the worst of these cults are the \textit{Ravagers}, the scourge of the \textit{Dividing Plains}.

\subparagraph{Depiction}
Simple clay representations in communities of the Ruiner's followers show him as a hulking, bulbous behemoth of an orc. His missing eye has shifted, the prominent remaining eye now centered in his face like that of a nightmarish cyclops.

\subparagraph{Divine Domains}
Many of the Ruiner's clerics are granted powers of Death, Tempest, or War.

\subparagraph{Enemies}
The The Arch Heart shot out the Ruiner's right eye during the \textit{Calamity}, and the god of slaughter longs for the day he can return the favor twofold. Likewise, the Ruiner is a dark mirror of the The Wildmother, and his hatred for her kindness is surpassed only by his lust for vengeance against the The Arch Heart.

\subparagraph{Symbol}
A black heraldry torn and ragged at the bottom. In its center is one, lone bleeding eye with one large drop of blood emanating from the bottom.
{\par \vspace { 9pt plus 3pt minus 3pt } \noindent  \bfseries \dndsanscaps{Commandments of the Ruiner} \nopagebreak }

\begin{itemize}
\item Ruin. Ravage. Kill.
\item The weak exist to be crushed by the strong. Be the strong.
\item There are no emotions but fury and joy. The rest are weakness.
\end{itemize}

The Ruiner is also known as the progenitor of the Curse of Ruin.


\subsubsection{The Spider Queen}
The god of deceit, shadows, and spiders, the Spider Queen weaves a complicated web of schemes and treachery through her worshipers, deceiving allies and enemies alike to gain power. The Spider Queen's worship is entwined with the society of dark elves across much of Exandria, even as the drow fall ever deeper into the thrall of the The Chained Oblivion aberrant minions. It is said that the Spider Queen can see through the eyes of all spiders and that she is truly all-knowing.

A mysterious faction of elves from \textit{Wildemount}, the Kryn Dynasty, escaped this evil god's thrall in ancient times. This civilization's geographical isolation from Tal'Dorei keeps communication between them limited, but a small number of Kryn dark elves have traveled to Tal'Dorei to investigate the \textit{subterranean tunnels beneath the Stormcrest Mountains}.

\subparagraph{Depiction}
Ancient icons and idols dedicated to the Spider Queen show an alluring woman with deep purple skin and silver hair, her abdomen swelling into the terrifying body of a monstrous spider. Her depictions are highly polarized between her elite devotees and her wretched supplicants. They are either opulent gilded statues, or deranged cave paintings.

\subparagraph{Divine Domains}
Many of the Spider Queen's clerics are granted powers of Knowledge or Trickery.

\subparagraph{Enemies}
The Spider Queen holds a searing grudge against the The Stormlord, for with one throw of his mighty thunderspear, he impaled her against a cliffside, leaving her drow armies leaderless during the \textit{Calamity}. The Spider Queen and the The Ruiner also share a burning hatred for the The Arch Heart, whom she claims drove her children below the earth. Despite their shared enemy, the Spider Queen often manipulates the The Ruiner followers into attacking her foes, so that the drow may remain safe. The only being the Queen fears is the The Chained Oblivion; she dreads its alien mind and is furious about how its madness has gripped her children.

\subparagraph{Symbol}
A pendant of a blue spider with red markings on its black. Its abdomen is marked with a golden crown.
{\par \vspace { 9pt plus 3pt minus 3pt } \noindent  \bfseries \dndsanscaps{Commandments of the Spider Queen} \nopagebreak }

\begin{itemize}
\item It is better to be loved than feared, but you may certainly try to be both.
\item Misdirection, slander, and shadowed steps have more function than direct conflict.
\item Death to the elves who live under the sun, and death to all their allies!
\end{itemize}


\subsubsection{The Chained Oblivion}
It is darkness unending, less like a god and more like another world of malevolent, alien intelligence. Life and death do not exist within the Chained Oblivion; only the absolute end of all things. Few minds can comprehend the depth of oblivion. Most who do are so completely lost to nihilistic despair or denial that they devote the rest of their lives to hastening the end of this reality so that it may be swallowed up by the great inevitability: utter, tenebrous cold.

Even the other \textit{Betrayer Gods} treat the Oblivion with caution. In its endless imprisonment, it dreams nightmarish aberrations into existence, deep beneath Exandria. The Chained Oblivion's demented cultists work without word from their twisted patron, awaiting the Epoch of Ends, when its freedom will be attained and all beings shall be consumed in deathlessness unending.

\subparagraph{Depiction}
Few visual attempts to depict the Chained Oblivion exist, but texts describe a creature of roiling ink and hungry darkness, a spreading cloud of lightless destruction born from a thousand ravenous mouths. Current references show the nightmare constrained by chains of black and gold, barely keeping the dark at bay.

\subparagraph{Divine Domains}
Many of the Chained Oblivion's clerics are granted powers of Death, Grave, or Trickery.

\subparagraph{Enemies}
The few remaining followers of the The Knowing Mentor, who was gravely wounded by the Oblivion in ancient times, are dedicated to ensuring that the Chained Oblivion is never again unchained. Their goddess led the charge against the all-consuming god's reemergence, and her bravery allowed the other \textit{Prime Deities} to shackle their enemy, though the victory nearly cost the Mentor her immortal life. Some say the The Knowing Mentor can only be fully restored if the Chained Oblivion is fully destroyed—for good.

\subparagraph{Symbol}
A holy symbol carved into a stone tablet. The symbol is of a jagged spiral surrounded by a barbed, silver chain.
{\par \vspace { 9pt plus 3pt minus 3pt } \noindent  \bfseries \dndsanscaps{Commandments of the Chained Oblivion} \nopagebreak }

\begin{itemize}
\item Offer light and life into its all-consuming maw so that it may surpass all divinities and be freed.
\item Uncover, restore, and exalt forgotten shrines and relics in its honor.
\item Ruin and raze the realms to hasten the arrival of the Epoch of Ends.
\end{itemize}


\subsubsection{The Scaled Tyrant}
The evil queen of dragons is a fearsome god of greed, envy, and hoarded wealth. While chromatic dragons are her foremost worshipers, the Scaled Tyrant accepts the worship of any who crave wealth. All chromatic dragons have a fearful reverence for their tyrannical queen, but many dragons of near-deific power and ambition chafe under her rule. Though the Scaled Tyrant was pleased by the devastation that the \textit{Chroma Conclave} spread across Exandria, the prideful \textit{Conclave} never considered themselves subservient to the Dragon Queen.

\subparagraph{Depiction}
Most visual representations of the Scaled Tyrant exist as warnings within sanctuaries of the The Platinum Dragon. She is shown as a dragon of frightful size with five vicious dragon heads, each one a chromatic color of her evil children. Massive, leathery wings sprout from her back, spreading clouds of poisonous mist.

\subparagraph{Divine Domains}
Many of the Scaled Tyrant's clerics are granted powers of Order, Trickery, or War.

\subparagraph{Enemies}
The Scaled Tyrant's hatred for the The Platinum Dragon is as old as the \textit{Founding}, and her cults are ever-hunted by his justice. For centuries, the zealous paladins of the The Platinum Dragon have limited the Tyrant's worship to the chromatic dragons, but sects of humanoid dragon-worshipers appeared in the wake of the \textit{Chroma Conclave's} destruction of \textit{Emon}, and have been secretly keeping her faith alive in the generation since the Conclave's defeat. Some Tal'Dorei oligarchs have even turned to the Scaled Tyrant in secret, as their hoarded wealth begins to run dry, sapped by the immense cost of rebuilding what was lost to the \textit{Conclave}.

\subparagraph{Symbol}
A black circular holy symbol with red runes outlining its entire circumference. In the center lies a piece of parchment with black markings upon it, made to look like a dragon's talon.
{\par \vspace { 9pt plus 3pt minus 3pt } \noindent  \bfseries \dndsanscaps{Commandments of the Scaled Tyrant} \nopagebreak }

\begin{itemize}
\item Amass wealth, but spend little. The gold—and the power that comes with it—is sufficient reward in itself.
\item Do not forgive nor forget an indignity to yourself. Let no affront go unpunished.
\item Take what you covet. Those without the strength to defend their dominion are not worthy of it.
\end{itemize}

\subsubsection{The Crawling King}
The god of the endless tunnels and caverns beneath Exandria, the Crawling King is patron to torturers, slavers, and jailers across the realms. The furious, molten tears he wept upon his defeat carved a network of infinite subterranean pathways beneath Exandria—and in the vastness of the Outer Planes, his plane of imprisonment is a realm of logic-defying caves and grasping manacles from which few return alive.

Those who rob others of their freedom offer prayers to him in cellars and other subterranean domains, and other creatures who live in the darkness below worship him and seek his guidance.

\subparagraph{Depiction}
The Crawling King is rendered as a swollen, malformed worm that slithers through the dark below, with a screaming, hairless human head and three arms carving through the lightless rock.

\subparagraph{Divine Domains}
Many of the Crawling King's clerics are granted powers of Death or Trickery.

\subparagraph{Enemies}
The The Dawnfather and the The Everlight defeated the Crawling King during the \textit{Calamity} by luring their nemesis above ground. The The Dawnfather pierced his tenebrous body with ten thousand lances of sunlight, and the The Everlight imprisoned him beyond the boundaries of Exandria within her serene rays. The tears of pain and anger the Crawling King shed burned through Exandria, and his faithful fled into these tunnels to escape their enemies' holy light.

\subparagraph{Symbol}
A circular copper holy symbol with a black void in the centre and three twisted pale hands extending from the void.

{ \par \vspace { 9pt plus 3pt minus 3pt } \noindent  \bfseries \dndsanscaps{Commandments of the Crawling King} \nopagebreak }
\begin{itemize}
\item Seek and exalt places where no light touches.
\item Revel in the pain you inflict upon others, and relish the pain you suffer yourself as an offering to the Crawling King.
\item Imprison those who cannot resist you, and drag all life into the darkness.
\end{itemize}

\subsubsection{The Whispered One}
The pain of defeat is cruel indeed. Twenty-four years ago, a cult known as the \textit{Remnants} rose up against the world, led by an undead archmage known as the Whispered One. This lich had toiled in secret for years within the Plane of Shadow, constructing the towering citadel of Thar Amphala there, while secretly spreading his poison throughout the storied lands of \textit{Issylra}. His goal was simple: repeat the lost rites of ascension that the The Matron of Ravens used in the \textit{Age of Arcanum} to become a god—and this time, to be on the opposite side of the Divine Gate, so that he could rule the world.

His ascension was successful, but the new god was ultimately banished behind the Divine Gate by the intervention of Vox Machina. Their victory, though grand, was not complete. The \textit{Remnants}, though bloodied, still venerate their spited master with offerings of secrets and unwilling souls to gain his favor.

The Whispered One is a god of hidden things, and the occult remnants of his followers prey upon those who harbor dark secrets, stoking their fears and driving them into his skeletal embrace. He also is a patron of undeath, and upon those who lavish him with their worship, he bestows the keys to creating ever-greater monstrosities.

\subparagraph{Depiction}
The Whispered One is a picture of sepulchral self-mutilation. Those faithful who saw him at Thar Amphala describe him as an unnaturally tall, gaunt humanoid with necrotic skin stretched taut over his bones, wrapped in robes of shadow, and adorned with jeweled talismans.

\subparagraph{Divine Domains}
Many of the Whispered One's clerics are granted powers of Arcana, Death, or Knowledge.

\subparagraph{Enemies}
The Whispered One now despises all the gods who stood in his way by anointing members of Vox Machina as their champions. The The Dawnfather, the The Everlight, the The Knowing Mentor, and the The Matron of Ravens all contributed to his defeat, and his \textit{Remnants} are hell-bent on subverting their mortal followings. He also harbors resentment for the The Moonweaver, who benevolently blesses the keepers of benign secrets, robbing him of followers.

\subparagraph{Symbol}
A light gold circular holy symbol with an upside-down five-pointed star. Sitting atop the star is a skull with one lone green light smoking from its left eye.
{\par \vspace { 9pt plus 3pt minus 3pt } \noindent  \bfseries \dndsanscaps{Commandments of the Whispered One} \nopagebreak }

\begin{itemize}
\item Learn all you can, and keep hidden that which you know. Reveal what pieces you must, but never the whole.
\item Express and cultivate the evil within yourself, and in doing so, recognize it in others to exploit them for your own benefit.
\item Seed the ruin of all who worship other deities, until only those who kneel before the Whispered One remain.
\end{itemize}


\subsubsection{The Cloaked Serpent}
A wanderer in the shadows and the creator of snakes and serpentkin, the Cloaked Serpent is the evil god of poisons, assassins, and darkness. The ancient serpentkin worship him over all other deities, dragging screaming offerings to their temples in his honor. Most of the Cloaked Serpent's worshipers were annihilated during the \textit{Calamity}, and those that weren't are either suspended in self-induced stasis or hunted for sport by the servants of the The Spider Queen and the The Crawling King. But perhaps the Cloaked Serpent is merely biding his time, waiting for the proper moment to unleash his hidden armies upon the world once more.

In Tal'Dorei, most of the Cloaked Serpent's serpentkin followers live in arcane suspension on the mysterious \textit{Visa Isle}.

\subparagraph{Depiction}
Many forgotten temples were once built to worship the Cloaked Serpent, and his image was embedded within much of their architecture. He possesses a humanoid body with six arms and a fanged, serpentine head adorned with a threatening hood. Strands of thick, dark hair sprout from his scaled body, obscuring his form in layers of swirling shadow.

\subparagraph{Divine Domains}
Many of the Cloaked Serpent's clerics are granted powers of Nature or Trickery.

\subparagraph{Enemies}
The Cloaked Serpent loathes the The Lawbearer and the The Wildmother, for he despises life, order, and love above all things. His surviving worshipers relish in using poison and fire to consume nature and undermine civilization, hoping to infuriate both goddesses and cast Exandria into chaos.

\subparagraph{Symbol}
A circular holy symbol of black silver with golden filigree upon it. Composing the entirety of its center is a coiled, golden snake - its fangs bared, its eyes green.
{\par \vspace { 9pt plus 3pt minus 3pt } \noindent  \bfseries \dndsanscaps{Commandments of the Cloaked Serpent} \nopagebreak }

\begin{itemize}
\item Keep your acts obfuscated and secret. The night is your greatest ally.
\item Strike quickly and without reason. Blind the target with their own confusion.
\item Kill slowly. Agonizingly. Or worse, make them enjoy it.
\end{itemize}


\subsection{Betrayer Curses}
The malign powers of some \textit{Betrayer Gods} linger in Exandria, both as myth and reality.


\DndItemHeader{Curse of Strife}{Curse}
The insidious voice of the Strife Emperor, both damnable and damning, invades the minds of countless people across Exandria who have been afflicted by the Curse of Strife. The curse is especially prevalent in Rifenmist, where the Strife Emperor's faithful clerics bestow it upon their soldiery, transforming them into perfect weapons of war. The curse acts in unpredictable ways, turning some into chaotic berserkers, others into obedient legionnaires, and others still into selfish, isolationist thieves. Cultists of the Strife Emperor also use their power to covertly spread the curse throughout hospitable and kind communities across Exandria, aiming to destroy them from within and make them more susceptible to conquest.

\paragraph{Origin of Strife} Such transformations are the result of the Strife Emperor's torment, which whispers constantly in the victims' minds. The free goblinkin of Tal'Dorei regard those who spread the curse with special contempt—and those who suffer from it with immense pity.

For centuries, small-minded people in Tal'Dorei have feared goblins because they are the children of the Strife Emperor. But while the first goblins, hobgoblins, and bugbears were created in the Calamity as soldiers for this fell god, the curse does not linger in their blood—it spreads via those who pledge their souls to the Strife Emperor. Clerics of the Strife Emperor can use the bestow curse spell to inflict the Curse of Strife.

\paragraph{Curing the Curse} Creatures who experience kindness or suffer near-death experiences can have epiphanies that allow their true self to tear through the bloody shroud of the Strife Emperor's lies. A curse-afflicted creature who spends one week being treated with love and compassion instead of hatred and fear can make a DC 20 Wisdom saving throw, breaking the curse on a success. The remove curse spell also immediately ends the curse—and many reformed clerics of the Strife Emperor have united with clerics of other faiths to loosen the bloody god's grip on their families.

\DndItemHeader{Curse of Ruin}{Curse}
Countless legends invoke hgar'Gruum, the Curse of Ruin. Great orc warriors have been laid low when bloodlust overtook their reason. Mighty mages driven by a foolhardy ambition forsook caution and delved furiously into knowledge no mortal eyes should behold. The Curse of Ruin is said to be the birthright of all orcs—a deadly taint carried in the blood of the Ruiner that forged their people during the Calamity.

\subparagraph{Truth of Ruin} The truth is that the Curse of Ruin is a myth. The blood of the Ruiner granted orcs physical strength and superhuman endurance, but it imparted none of the Betrayer God's evil. The power of the myth, however, is more ruinous than any magic.

\subparagraph{Consequences of the Myth} Countless superstitious folk across Exandria—including many orcs themselves—believe in the Curse of Ruin. This belief doesn't make the curse any more real. The myth's existence is a convenient way for any one orc to ignore the real consequences of their mistakes, or of any act of thoughtless harm—the curse takes the blame. Even worse is when this "birthright" is leveled at wise, even-tempered orcs who have done nothing wrong, and their frustrated reactions are then taken as proof of the curse's existence.


\subsection{Lesser Idols}
Though the gods still make their wills known upon the world by granting power to mortal clerics and by delivering edicts through their celestial or fiendish servitors, the Divine Gate inhibits them from exercising their full powers over the fate of Exandria. The same is not true for the lesser idols, beings of demi-godly power that have, in the absence of true gods, used their splendid and profane forms to amass their own cult followings.

Dozens of these powers are worshiped in relative secret across Exandria. Some are supernatural beings like archfey or archfiends that have taken up near-permanent residence on the Material Plane. Others are unknown fragments of the power of primordial Creation, given form. Most are mortal champions that have gained near-divine power through magic, the blessings of the gods, or through mythic effort and training. In all of these cases, the worship of mortals elevates these beings to something greater than they were before, imbuing them with fragments of divine power—and allowing them to grant scraps of that power back to the mortals in whom they take an interest.

\begin{DndTable}[header=Lesser Idols]{Xll}

Deity & Alignment & Province\\
Azgrah & Lawful neutral & Undying\\
Galdric & Lawful good & Celestial, Undying\\
Graz'tchar & Chaotic evil & Fiend, Hexblade\\
Khedive Xundi & Chaotic neutral & Genie\\
The Lost Beacon of Unknown Light & Neutral & Celestial\\
The Observer & Neutral good & Archfey, Celestial\\
The Sightless & Chaotic evil & Great Old One\\
The Traveler & Chaotic neutral & Archfey\\
Vesh & Neutral evil & Archfey, Undying\\
\end{DndTable}

\subsubsection{Azgrah}
The underground citadel of \textit{Emberhold}, along with the subterranean territories surrounding it, pays fealty to the The Crawling King as master of the underground. However, for all but the King's most faithful, the homage they pay to him is a distant second to their veneration of Azgrah, Lord of the Deep Dwelling. This mighty champion is a duergar warrior who completed six heroic labors in his prime, all of which have been frequently enshrined in bardsong, from the wooing of Prince Lhebeno to the Conquest of the Amber Worm—in which he defeated a servant sent by the The Crawling King to punish the present Queen of the \textit{Emberhold} for her impiety.

Azgrah has enjoyed an incredibly long life, even for a duergar, and now lives in isolation in a temple deep within the Amber Worm's writhing tunnels. He is tended to only by his most trusted warlocks, and is venerated by all as the Lord of the Deep Dwelling, and as the dwarf who will emerge in the \textit{Emberhold's} hour of greatest need.


\subsubsection{Galdric}
Given purpose and power as the companion to Pur-van Suul, champion of the The Matron of Ravens during the \textit{Calamity}, the black wolf Galdric was sealed within a \textit{raven's slumber}, one of several such relics made in the Matron's honor, and interred with his master in a tomb under the Marrowglade Loch in the distant land of \textit{Othanzia}.

Over twenty years ago, Vaxil'dan of Vox Machina was named champion in Purvan's stead, and Galdric was released from his slumber and bestowed with the charge of guarding the \textit{Parchwood} surrounding \textit{Whitestone}. This large and cunning wolf now stalks the woods as the city's silent protector, with the people crafting new legends about their bestial sentinel. On moonlit nights, some say you can see Galdric wander through the Greyfield, headed to the The Matron of Ravens shrine for communion.


\subsubsection{Graz'tchar}
Centuries ago, \textit{Zan Tal'Dorei} defeated \textit{King Trist Drassig} in the \textit{Battle of Umbra Hills}. \textit{Drassig} died wielding a blade called Graz'tchar, gifted to him by a cunning demon among the ranks of the fiends granted to him by the The Strife Emperor. Zan let \textit{Drassig's} blade lie where it fell, along with his body.

This blade contains a sliver of the ego of the Demon Prince of Indulgence, and enacts his anarchic, hedonistic will upon the world. From up to one thousand miles away from its resting place in the \textit{Umbra Hills}, the voice of Graz'tchar, the Decadent End can be heard faintly in the minds of those it reaches out to—and, so it is said, only one who has heard the voice of the blade can find it amidst the blackened crags of that accursed battlefield.

Graz'tchar, the Decadent End is willing to grant power to those who accept it, even from afar. It prefers to test those who hear its call before leading them to find it. When the time comes that someone does wield Graz'tchar once more, its game statistics are Graz'tchar, the Decadent End.


\subsubsection{Khedive Xundi}
Only about a year ago, a genie of elemental water was freed from captivity beneath Nicodranas—a city in \textit{Wildemount}, on the other side of the \textit{Lucidian Ocean}. This marid, named Kedive Xundi, has now returned to his Palace of Silver Silt on the Elemental Plane of Water, and is happy to remain there—but he is itching for payback.

But that's a continent away, and Xundi has time for both payback and entertainment. While he manipulates mortals on \textit{Wildemount's} Menagerie Coast, he also has deposited a beautiful clay amphora in the port city of \textit{Stilben}. Through this item, he watches passersby, occasionally enticing foolish mortals to pick it up and give it a rub. From there, Xundi offers them power, with no strings attached—save for a request for them to go on the most thrilling adventures a little mortal can. And maybe, if they survive long enough and become powerful enough, the marid will invite that mortal and their friends to his Palace of Silver Silt for a grand banquet, all the pleasures they can imagine, and then a request for them to engage in the most thrilling adventure he can think of—and Xundi has had the length of a most interminable imprisonment to think of incredible adventures.


\subsubsection{The Lost Beacon of Unknown Light}
In a time before time, when Exandria was naught but darkness in the chaos before even the \textit{Founding}, a luminous being called the Luxon sought answers about their existence. They lit the first fires within Exandria with their unquenchable light, and then began to wait. Gods came. They created beings that spread across the planet. Those creatures died and procreated and died in an endless cycle, of which the Luxon grew curious. Surely, by watching these short-lived beings, they could learn the answers they sought.

So the Luxon divided their own light into beacons, dodecahedral crystals of eternal radiance. They cast them across the world and fell into a deep slumber, during which they hoped people would find the beacons and bond to them, allowing the mortals' souls to be reborn again and again so they could learn everything there is to learn about existence.

Many such beacons have been found across the world, mostly by a people who have learned the ways of the Luxon in a far-off land. However, many of the beacons remain unaccounted for. One such beacon was unearthed about ten years ago by Tal'Dorei settlers in the \textit{Cliffkeep Mountains}, far from any civilization. The people of this settlement, \textit{Reaching Bluff}, call it the Beacon of Unknown Light, and treat it as a mystical oracle that will guide their young settlement to prosperity.

Despite the settlers' best efforts, word of their discovery has escaped \textit{Reaching Bluff}. As adventurers and treasure-seekers begin to trickle in, the settlers protect their newfound oracle with cultish devotion—despite the fact that it has never once spoken to them, as they have beseeched it to countless times.


\subsubsection{The Observer}
Amidst the \textit{Age of Arcanum}, the people of the great \textit{Tetrarchy of Qoniira} were visited by a being of sheltering power and comforting light. The people of Qoniira were on the path to creating the most prosperous civilization that the \textit{Rifenmist Jungle}—and indeed, all of \textit{Gwessar}—had ever seen. In its peace and magical might, it exceeded even the grasp of the nearby elven domain of \textit{Syngorn}. Magic, entirely separate from the arcane rituals practiced by arcanists of the era, fueled their city's rise to splendor, and the being that came to them looked on in fascination.

It appeared before them in a form that at once seemed familiar and majestic. It had a body like one of the \textit{Catfolk}—a wiry humanoid frame with a feline face, covered in tawny fur and possessed of a tail and limbs that ended in large paws. Beyond this familiar appearance, four broad, avian wings burst from this ten-foot-tall avatar's back. It was surrounded by a golden light that seemed not to illuminate and saturate the world around it, but make it dimmer, ashy, and washed-out. Yet in this pallid light came a sharpness of perception, as if it were more real than the material world around it. The being appeared in the center of the Qoniira civilization, at the heart of the city of \textit{Niirdal-Poc}, and it spoke a single phrase: "I am the Observer, and I give unto you the Gift." Then, it vanished from sight.

All at once, the gray-gold light of the Observer surged throughout Qoniira. Though the citizens didn't notice a difference at first, as time went by, many people began to manifest new forms of magic. This magic was not the controlled, scientific power of far-off arcanists in their flying cities, but chaotic, unpredictable, and miraculous.

The once-sprawling civilization of Qoniira has shrunk over the centuries, but its people still draw upon the miracles of the Gift, and they are still watched from afar by the unseen Observer. This entity does not reign like a god, but watches impassively those whom it gave its unpredictable Gift. As those born in \textit{Niirdal-Poc} under the auspices of the Observer travel north into the Republic of Tal'Dorei, or south into the oppressive realm of the \textit{Iron Authority}, the gaze of the Observer is moved even farther afield—adding ever more variables to its mysterious experiment.


\subsubsection{The Sightless}
In the \textit{Age of Arcanum}, an archmage named Clemain Astural was the first to make contact with an otherworldly power of the Realms Beyond that called itself the Sightless. His legendary writings, the well-guarded \textit{Astural Scrolls}, say that the entity's words of greeting to him were:

"I am the Sightless. Will you not lend me your eyes?"

The \textit{Astural Scrolls} are vague about the powers of the Sightless—also sometimes called the Sightless One, the Borrower of Eyes, and the Lurker Behind One Thousand Eyes—but they are specific about several things. First, Astural willingly relinquished his own sight to the Sightless in exchange for his patron's magical powers. He gained all the power he ever wanted, but Exandria was as alien to the Sightless as the Sightless was to Exandria. Forced to rely upon this entity's horrific interpretations of the world around it, Astural quickly lost his grip on reality.

In time, Astural's unchecked devastation was stopped by great heroes of the time, and his scrolls were confiscated. They exist in the \textit{Westhall Academy in Westruun}, and the second scroll, though secured in a deep vault, still holds the secrets of the ritual Astural used to contact the Sightless all those years ago.


\subsubsection{The Traveler}
Over the past few years, rumors of a green-cloaked deity with a wiry frame and fiery red hair have spread across Tal'Dorei and the neighboring continent of \textit{Wildemount}. This fey trickster, though not a god, has certainly enjoyed pretending to be one to spread mirthful chaos across the world.

Hundreds of isolated individuals across Tal'Dorei have heard the voice of the Traveler—or they've at least experienced one of his wily pranks. Within the past year, a convention of the Traveler's devotees congregated on the tiny island of Rumblecusp in the \textit{Lucidian Ocean} to exult his name and deeds.

And then, suddenly, he vanished. The tales of a trickster deity that washed across Tal'Dorei trickled to a halt, and a few of his once-devoted worshipers hung up their cloaks in disappointment. Something happened on Rumblecusp that shook the faith of dozens of Traveler-worshipers. Some did not return at all.

But amazingly, even though the divine narrative of the Traveler has been abandoned by all, the Traveler's devoted have only grown in number and gleeful impishness. The current story is remarkably close to the truth: that the Traveler is a mighty lord of the fey. This archfey, known as Artagan to a select few, despises the relentless pressure of celebrity godhood, and just wants to bring joy and freedom from responsibility to all. He still has plenty to learn about mortals, and about when a bit of good, harmless fun becomes not-so-harmless after all.

\subparagraph{Symbol}
An arch covered in emerald green vines with purple and orange flowers. Through the arch is a winding dirt path leading through a valley into a sunrise on the horizon.

\subsubsection{Vesh}
Blood and pleasure follow in the wake of the siren Vesh. Long has she stalked the shadows of far-off \textit{Issylra}, amassing secretive orgies of people of all genders and ancestries. These grand events have a membership that numbers in the hundreds, all of whom are bound to utmost secrecy, and are thrown in obscure retreats that are hidden in cities and towns across the continent. Little do the foolish pleasure-seekers in attendance know that the siren's great fetes of debauchery can end with only two left alive: Vesh and her new mate, who will go on to form a mystical blood-bond.

\enquote{Beware the Bloody Siren}," sob those few who have escaped her orgy of blood, but their warnings have only fueled her legend. In Tal'Dorei, where Vesh's presence is known only as a legend of a foreign land, her name is most often invoked by youths who want to impress their peers. But who can say what would happen if an attractive and foolish mortal called too fervently or too frequently to the Bloody Siren? Perhaps she would appear, baring her subtly radiant gray skin, stroking her flowing violet hair, and gazing invitingly at the fool with her blank white eyes.


\begin{DndSidebar}[float=!b]{Other Secluded Champions}
A handful of other mortals exist in Tal'Dorei who have risen to nearly divine power through destiny or sheer will. Much like Azgrah, the vast majority of these champions live out their lives in seclusion—and perhaps it is their seclusion that grants them the ability to lend their powers to others, while they await the day that their full strength is at last needed to defend their people.



All have their own stories, but no matter their tales, whether they are honored or despised by their kinsfolk, they are all mythic iconoclasts. Other secluded champions of Tal'Dorei that you may wish to develop into idols of your own include:




\begin{itemize}
\item Flavia, Queen of Unending Storms (cloud giant)
\item Suthine, Eternal Princess of Vulkanon (fire giant)
\item Valdemar, Draugr-Jarl of Farborg (frost giant)
\item Ghaladron, Traitor to the Iron Crown (hobgoblin)
\item Tevrosk, Whose Axe Was Wreathed in Flowers (orc)
\item The First Favored Champion (stone giant)
\item Typhoe, the Dreaming Shark (storm giant)
\end{itemize}



\end{DndSidebar}

\section{Calendar and the Cosmos}
The world of Exandria is a planet with regular seasons and a wide variety of climes, ranging from frigid ice caps to arid deserts, temperate forests, windswept plains, sweltering jungles, and storm-tossed oceans. It orbits a single sun, is orbited in turn by two moons, and is surrounded by distant stars that twinkle in the evening skies.

The ancient people of Exandria told time by these celestial bodies, and the elves created a calendar based on their movements.

\subsection{Calendar}
The calendar year of Exandria runs a total of 328 days, grouped into seven-day weeks over the course of eleven months. These months are outlined below in the order of their arrival within the calendar year, along with their respective number of days and holidays common throughout Exandria (many of which are detailed within the "\textit{Pantheon of Exandria}" section in \textit{chapter 2}).

\subsubsection{Counting Years}
In Exandria, the years are most commonly counted by a reckoning known as "\textit{Post Divergence}," with year 0 being the fateful year that the \textit{Prime Deities} ended the \textit{Calamity} by sealing all the gods behind the Divine Gate. This book describes the world as it is in the year 836 P.D. If you're familiar with the \textit{Critical Role} campaigns, this year is twenty-four years after the end of Campaign 1 (Vox Machina) and toward the end of Campaign 2 (The Mighty Nein).

\subsubsection{Months and Seasons}
The names of the months and days of the week vary between cultures, but the elven calendar described here is the standard means of reckoning dates in Tal'Dorei—as well as in most nations around the world, for ease of international trade and diplomacy.

Just as in our world, there are four seasons, and the year ends and begins just after the winter solstice.

The first blush of spring in Tal'Dorei is felt as new flowers bloom on Wild's Grandeur, a holy day on the 20th of Dualahei. The official start of the season, however, is observed a week earlier and celebrated with games, music, and just-ripened food at the Renewal Festival. Summer brings hotter days toward the middle of Unndilar. The noonday sun of the 26th day is called the Zenith, and this solstice is considered the first true moment of the summer.

The colors change and winds cool as autumn begins in the early days of the month of Fessuran, marked by the Harvest's Close on the 3rd day. The chill of winter arrives to bring longer nights and cleansing snow on Barren Eve, the 2nd of Duscar, a nighttime celebration and remembrance of those who fell in battle.

The seven days of the week are named Miresen, Grissen, Whelsen, Conthsen, Folsen, Yulisen, and Da'leysen. Each day is 24 hours long, and most city-dwellers in the Republic of Tal'Dorei are expected to work eight hours of that day—except for Yulisen and Da'leysen, which comprise the weekend, and on holidays.

These social rules don't necessarily apply to cultures where sleeping patterns are different (such as in \textit{Syngorn}, where most of its people are elves who trance a mere four hours of the day), or to people outside of cities—such as farmers, who rise with the sun, and traveling merchants and adventurers, who work highly unusual hours.

\subsubsection{Using a Calendar}
The Exandrian calendar is intentionally quite different from our own. The creation of calendars is a messy and often illogical process, with strange artifacts of their making that linger on only because of tradition—and because it would be too much of a hassle to get the entire world to change their way of counting time and days.

You don't need to bother with the Exandrian calendar in your home game. It's perfectly reasonable for you to use a nice, simple seven-day week, thirty-day month, and twelve-month year. You may even use our names for the months and the days of the week so that when an NPC tells the party to meet them at noon on "Grissen, Misuthar 7th," they're basically talking about "Tuesday, February 7th."

However, if you want to immerse your players in the fantasy world of Exandria, using this unfamiliar calendar may help you on your quest.

\begin{DndTable}[header=Exandrian Calendar]{lcX}

Month & Day & Holidays\\
Horisal & 29 & New Dawn (1st) \newline Hillsgold (27th)\\
Misuthar & 30 & Day of Challenging (7th)\\
Dualahei & 30 & Renewal Festival (13th) \newline Wild's Grandeur (20th)\\
Thunsheer & 31 & Harvest's Rise (11th) \newline Merryfrond's Day (31st)\\
Unndilar & 28 & Deep Solace (8th) \newline Zenith (26th)\\
Brussendar & 31 & Artisan's Faire (15th) \newline Elvendawn or Midsummer (20th)\\
Sydenstar & 32 & Highsummer (7th) \newline Morn of Largesse (14th)\\
Fessuran & 29 & Harvest's Close (3rd)\\
Quen'pillar & 27 & The Hazel Festival (10th) \newline Civilization's Dawn (22nd)\\
Cuersaar & 29 & Night of Ascension (13th) \newline Zan's Cup (21st)\\
Duscar & 32 & Barren Eve (2nd) \newline Embertide (5th) \newline Winter's Crest (20th)\\
\end{DndTable}

\subsection{Catha and Ruidus}
When the people of Tal'Dorei gaze into the morning sky, they see the sun, a blazing disc the size of a gold piece, looking back at them. It is the domain of the The Dawnfather, a god of life and light honored politely by nearly all the realm's denizens.

At night, the skies are filled with countless distant stars, a gleaming silver-white moon called Catha, and occasionally, a dim, ruddy moon called Ruidus that is nearly half Catha's size. Both moons are the domain of the The Moonweaver, a capricious god of trickery and illusions, but there are sects throughout the world that believe only Catha represents the The Moonweaver cunning and grace. Catha's pearly glow is said to bless the just with cunning and caution, and to make hidden the goodhearted when they require stealth and subtlety.

Ruidus, on the other hand, is so surrounded by disquieting rumors and folkloric tales of misfortune that some believe another unknown god or power rules this small, reddish-brown moon. Cultures around the world tell countless legends of prideful rulers who made grand plans or attempted deeds under the moon's full light—when it shines a brilliant vermilion rather than its usual ruddy color—and were forced to watch in horror as their endeavors fell to unforeseen misfortune. It is said those who fall afoul of Ruidus failed to give it the deference it is due—and so superstitious folk rarely dare to make plans while the full light of Ruidus shines above, let alone enact them.

Worse yet, some tales forebode dark fortunes for those born under the light of a full Ruidus, a curse of ill luck that will follow them throughout their lives. Though the cruel practice of moon-sacrifice is no longer permitted in any of Tal'Dorei's cities, some far-flung settlements still secretly sacrifice children born under a full Ruidus to "save" them from a cursed life, and to appease the dark and unknowable appetites of the red moon.

Fortunately for the superstitious, Ruidus is rarely full. While its cousin, Catha, completes a full cycle approximately once a month, Ruidus's haunting, half-year orbit, combined with its eerie and unexplainable tendency to simply not appear in the sky on certain nights or glow with unknown light, creating an unexpected full moon, has only added to its mythic reputation as an omen of ill fortune.

\begin{DndSidebar}[float=b]{Prologue to the Song of Alyxian}
    \setlength{\parindent}{0pt}
\enquote{O Ruidus, grant humble chorus leave\\
To sing the song which hails the zenith of\\
Your accurséd, thrice-blessed Apotheon;\\
Remember'd best by deeds in war, and yet\\
Whose acts were driven oft by fate most foul.

\vspace{6pt}

His kindly brow bore gifts from gods of change,\\
And art, and moon, yet in his soul was pain,\\
The suffering of your vermilion light\\
Drove the Paragon to a desert realm\\
Bestrewn with blades and drenched in crimson blood.

\vspace{6pt}

So hear, O moon of curséd deeds and fates!\\
The song of he who rose above your great\\
And mighty pow'r, to save Exandria.\\
From flames of war fanned by the Ruiner's blade.}

\end{DndSidebar}

\subsection{Planes}
Though this book concerns itself mostly with the world of Exandria, and specifically the continent of Tal'Dorei, high-level adventurers may eventually gain magic that allows them to travel to other planes of existence beyond the material world in which Exandria resides.

The planes surrounding Exandria are similar to the planes presented in fifth edition lore. Below are brief descriptions of some of the planes that a party of powerful adventurers may travel to. These planes go by many poetic and fanciful names.

\subsubsection{Material Plane}
The Material Plane is a land where natural laws work as we expect. It is a place of all four elements, and of plants, animals, and minerals. Exandria is but one of many worlds on the Material Plane—though none but the most accomplished astronomers have become aware of this fact, much less devised a way of traveling between material worlds.

\subsubsection{Elemental Planes}
When the Protean Gods discovered Exandria, it was like a barren ship adrift in a sea of churning, chaotic elements. As a part of their creation, the elements were thrust away from the world. Some theorize that the Elemental Planes were created from this banished chaos, while others believe the elemental chaos seeped into the Material Plane from tears in the fabric of the planes.

The latter may be more likely, as rifts between the Elemental Planes can be found across Exandria. The largest of these rifts are tended to by the four Ashari civilizations, and any who gaze within find wild wastelands of pure elemental power: roiling seas of fire, towering mountains of gnashing stone, infinite skies filled with cloud-top settlements, and endless, stormy oceans of water.

\subsubsection{Faerie, Shadow, and Ether}
Parallel to Exandria are odd reflections of its reality. The Fey Realm is a plane of natural beauty and ferocity, inhabited by fey folk. The Plane of Shadow is a dour, sinister mirror of the world that is. And the Ethereal Plane is a hazy film across the material world, as if it were viewed through an old, cloudy mirror, inhabited not only by restless spirits, but also by predators that target those who cling to the world of the living.

\subsubsection{Astral Plane}
All the prior planes exist directly adjacent to the Material Plane, buttressing and occasionally co-mingling with its inhabitants. All the planes that follow are separated from the material world by a vast sea of starlight, barren asteroids, and misty nebulae of unpredictable magic. In the ancient past, arcanists had to craft vessels to sail the treacherous sea of stars to travel between their home plane and the mysterious Outer Planes.

These days, magic such as the plane shift spell has allowed planar travelers to bypass these methods, but there are still a number of strange creatures and forgotten, imprisoned demigods who float eternally through the Cosmic Sea.

\subsubsection{Outer Planes}
The \textit{Prime Deities} imprisoned the \textit{Betrayer Gods} in prison planes after the \textit{Founding}, and retired to planes of their own creation after the \textit{Divergence}. Over time, the divine beings on these planes reshaped their new homes in their own image, creating places where natural laws don't function as one native to the Material Plane might expect.

All of these planes are strongly aligned to one or more cosmic principles of good, evil, law, and chaos, based on the divine powers that reside there. Because of the Divine Gate established after the \textit{Divergence}, no divine powers can depart these planes, for the sake of all Creation. The gods' supernatural creations can pass through the Divine Gate, though the more powerful they grow—and thus, the closer they come to divinity—the harder it is for them to pass through. However, mortals and lesser supernatural beings who possess magic to travel between the planes can travel to and from Outer Planes like the Hells, the Abyss, Elysium, and Pandemonium unhindered by the Divine Gate.

\subsubsection{Realms Beyond}
Far beyond even the divine Outer Planes is a realm where all laws of nature break down completely. Mortal senses have no way of properly perceiving the nature of these far-distant realms, and most minds that try to comprehend them snap under the strain of so much unfamiliar stimulus. These maddening realities and the alien horrors that live within them are mercifully kept away from the known multiverse by the unimaginable distance of space.

Some creatures of the beyond have lurked on the Material Plane since ancient times, perhaps at the invitation of the The Chained Oblivion itself, or because the divine powers unleashed in the \textit{Calamity} punched holes in time and space that allowed such aberrations to slither into Exandria unimpeded.

Major known locations of aberrant activity in Tal'Dorei include \textit{Yug'Voril} beneath \textit{Kraghammer}, the \textit{Crystalfen Caverns} beneath \textit{Emon}, and the realm of \textit{Ruhn-Shak} beneath the \textit{Stormcrest Mountains}.

\subsubsection{Demiplanes}
Massive, infinite spaces such as the Material, Elemental, and Outer Planes are not the only types of reality. Through the use of powerful magic, arcanists can create stable demiplanes—finite pockets of extraplanar reality that they can shape to their whims. These "pocket dimensions" exist as bubbles of floating color within the Cosmic Sea, and can only be traced or penetrated by extremely powerful, lost magical arts.

\section{Other Lands of Exandria}
The land of \textit{Gwessar}, now known far and wide as the young nation of Tal'Dorei, is a sprawling, curious land full of adventure, danger, and glory. Though Tal'Dorei is a land rife with grand tales, the world of Exandria is far larger and more ancient than just the lands between the \textit{Lucidian Coast} and the \textit{Bladeshimmer Shoreline}. Many of the people who live here either descend from or are immigrants to this realm.

Some of these other lands have been explored by Vox Machina and the Mighty Nein in \textit{Critical Role}, but are beyond the scope of this book. You can populate these lands with dungeons, cities, characters, and stories of your own creation—or be inspired by the history of Exandria in \textit{chapter 1} of this book, and the stories in the campaigns of \textit{Critical Role}, to create a version of Exandria entirely your own.

\subsection{Issylra}
Northwest of Tal'Dorei, across the freezing waters of the Ozmit Sea, lies Issylra, home to Vasselheim, commonly called the oldest surviving city on Exandria. Issylra is considered by most to be the land where all civilization emerged and spread in the time of the \textit{Founding}, making it the religious core for many, and the focal point of historical study for others. It was also where most of the flying cities of the \textit{Age of Arcanum} first rose from the ground. Its people rose high—and fell hard from those dizzying heights, and to this day, many there still hate and fear the power of arcane magic.

The territory known as \textit{Othanzia} encompasses a majority of the continent's developed regions, keeping most domains protected and governed by the might of the theocracy of Vasselheim. Nestled among eastern taiga forests of the Vesper Timberland, the city of Vasselheim stands as the oldest known seed of civilization and religion, and the only city to withstand every catastrophe that has rocked Exandria since the gods brought life to it. As the spiritual center of most of the world, Vasselheim is ruled by a coalition of high priests of the many \textit{Prime Deities}. These hierophants are viewed by their people as the defenders of their virtuous society, and of all of \textit{Othanzian} culture.

Beyond Vasselheim's ancient walls are frigid, untamed wilds, some of which are untouched since the \textit{Founding} itself. These lands are populated by ancient monsters and dire beasts. Smaller camps, communities, and townships are scattered across the surrounding snowy plains, and those who call \textit{Othanzia} home have become hardy and self-sufficient.

Elsewhere in \textit{Othanzia}, the lands of Issylra are sparsely populated by scattered republics of varying morality. Traveling between these city-states is difficult, due to the lingering beasts of the \textit{Calamity} that stalk the lands beyond Vasselheim.

\subsection{The Shattered Teeth}
Few sail into the dangerous waters southeast of Tal'Dorei, past the ever-present fog bank called the Fool's Curtain, and into the Shattered Teeth. This cluster of broken lands is made up of forty-three islands, varying in size from small, mile-wide reef-toppers to the large, city-holding isles like Ruukva. Two societies dwell within the Shattered Teeth. The first is the Ossended Host, a vast fishing culture that values their independence—even to the point of isolationism—and self-empowerment through the worship of the power of dreams and nightmares. The second is the Wanderman Assembly, a centuries-old trade company that, upon being stranded across the Shattered Teeth after losing their way in a typhoon, grew into a brutal, ruthless society dedicated to gaining profit at all costs. It masquerades poorly under the banner of honor and brotherhood.

This tension has led to bloodshed in recent years, and each group is paranoid that the other is bent on destroying their way of life. What brims beneath the surface of this conflict is the nature and source of the dreams the Ossended Host have been harnessing for generations, and what that answer may mean for the struggles of the people of the Shattered Teeth.

\part{Tal'Dorei}
\chapter{Welcome to Tal'Dorei}
\DndDropCapLine{T}{al'Dorei} is a land of epic adventure, filled with heroes who do great deeds for love, for justice, and—maybe more often than not—for gold and their own egos. Villains are draped in shadow and draw upon the power of cruel gods—but they are just as motivated by love, their own twisted concepts of justice, and their hunger for gold, power, and glory.


The name Tal'Dorei refers to three things throughout this book. Most often, Tal'Dorei means the continent that is home to all of the people, locations, monsters, magic items, and more within these pages. It is one of four major continents on the world of Exandria. Tal'Dorei can also mean the Republic of Tal'Dorei, the largest nation on the continent. And finally, Tal'Dorei was first the surname of a mighty hero who lent her name to the land that she helped liberate from evil.

See "\textit{A History of Tal'Dorei}" later in this chapter for more details about the heroic \textit{Zan Tal'Dorei}, the land she founded, and her descendants.

\section{Nations of Tal'Dorei}
The continent of Tal'Dorei encompasses several nations and independent city-states, all of which are home to diverse peoples who view the world in ways influenced by their cultures and the lands they call home. The most significant nations of Tal'Dorei include:

\subparagraph{The Republic of Tal'Dorei}
From the great city of \textit{Emon} in the west, across the vast \textit{Dividing Plains}, and even as far as the isolated city-state of \textit{Whitestone} in the northeast, the Republic of Tal'Dorei spans the breadth of the continent that bears its name. It is ruled by a council of elected representatives and is cordially allied with all other nations on the continent.

\subparagraph{The Ashari}
The elementally attuned Ashari people are reclusive and dedicated to their mission of safeguarding elemental rifts across the world. Two of the greatest rifts exist in Tal'Dorei, and the Earth Ashari of \textit{Terrah} and the Air Ashari of \textit{Zephrah} keep watch over them from their secluded homes.

\subparagraph{The Deep Halls of Kraghammer}
Deep beneath the \textit{Cliffkeep Mountains} is \textit{Kraghammer}, the ancestral home of Tal'Dorei's dwarven people. Though \textit{Kraghammer} is home to people other than dwarves, its dark and claustrophobic halls act as a deterrent to non-dwarven residents—as does the dwarf clans' famous inhospitality.

\subparagraph{The Verdant Enclave of Syngorn}
In the heart of the trackless \textit{Verdant Expanse} is a city of elves who hid in the Fey Realm from a war that nearly annihilated all life on Exandria. \textit{Syngorn} is welcoming to all who brave the \textit{Verdant Expanse} to reach them, though there is an expectation that those who visit take pains not to flaunt the strict and ancient rules of \textit{Syngornian} society, lest their stay come to a swift and unceremonious end.

\subparagraph{The Iron Authority}
This imperial force rules the \textit{Beynsfal Plateau} in the southernmost reaches of the continent. There they train legions of soldiers in an endless campaign to conquer the entire \textit{Rifenmist Peninsula}. The other powers of Tal'Dorei are just beginning to take notice of the movements of the \textit{Iron Authority}, but most of their knowledge is based in rumor. The \textit{Tal'Dorei Council} has little interest in drawing the authority's ire while it is preoccupied with its conquest of \textit{Rifenmist}.

\section{Running a Tal'Dorei Campaign}
Tal'Dorei is more than just its history; do not treat this book as an encyclopedia, nor as shackles that bind you to portray its characters and environments with perfect fidelity. As you read this manual, consider why you want to run an RPG campaign set in Tal'Dorei, and in the larger world of Exandria. What about \textit{Critical Role} inspires you as a Game Master? You might cherish the bonds of family and love between its many heroic characters. Perhaps you are drawn to the various factions that vie for power over Tal'Dorei's peoples. Or perhaps it's the way that evils from Exandria's mythic past inevitably rear their ugly heads, always when circumstances are already at their most dire, that most captures your attention.

No matter the initial spark that led you here, hold tight to it as you read, and search for ways to feed your personal flame of imagination. And when you are done, create. Create boldly, using these words as a launching point for your own tales.

This section is for Game Masters' eyes only, though GMs are welcome to share tidbits with their players to set the stage for the campaign, or to inspire them to create characters that will dovetail with the GM's story ideas.

\subsection{Secrets of Tal'Dorei}
The world of Exandria is steeped in familiar fantasy themes, with all the elves, dwarves, and wizards you might expect from a fantasy roleplaying game. But it also holds secrets that will challenge your expectations, and you may enjoy putting your own twists on these great mysteries of the world. Each of these secrets can form the backbone of an entire campaign, or can serve as flavor to create the impression of a world that is larger than just the player characters' story.

\subsubsection{All-Powerful, Imprisoned Gods}
Tal'Dorei was nearly destroyed when the gods fought one another in the \textit{Calamity}. Now, the \textit{Betrayer Gods} are sealed in otherworldly prison planes, and the victorious \textit{Prime Deities} live in self-imposed exile. An unbreakable barrier stands between the mortal and immortal realms. At best, a deity can grant their most skilled devotees a small measure of divine power to enact their will upon Exandria.

However, there are a number of cults who yearn to circumvent this barrier and summon their dark patron onto the Material Plane. Others still seek to unravel the forgotten rites of ascension, just as the The Matron of Ravens did in ancient times—just as The Whispered One did in living memory. What would happen if an omnipotent deity somehow managed to arise on the mortal side of the Divine Gate? Or even simply managed to fool the gods into thinking they had done so? Would the gods of good be forced to unleash their powers as one to break the Gate, unleashing another \textit{Calamity} upon the world?

\subsubsection{Dawn of a Republic}
It was the dying decree of Sovereign Uriel Tal'Dorei II that the continent-spanning Tal'Dorei Empire be reborn as the Republic of Tal'Dorei. No more would a sovereign rule the myriad peoples of this land, but instead the wise \textit{Council of Tal'Dorei} would enact the will of the people and protect them from the dangers of the world. But things don't always go according to plan. The nascent Republic of Tal'Dorei is suffering growing pains, and unless some heroes are able to help, the realm may fall into anarchy and chaos.

While fantasy politics isn't every gamer's cup of tea, the struggles of the fledgling republic can serve several purposes. Players looking for fantasy intrigue can meddle in the internal politics of the city of \textit{Emon} and root out corruption within the council, which may have been planted by any number of this land's power-hungry factions, from the invasive crime syndicate known as the \textit{Myriad} to the shadowy mages within the \textit{Arcana Pansophical} who splintered off into a new schismatic faction, the \textit{League of Miracles}.

Game Masters can also use political elements to flavor or provide a backdrop for their otherwise action-focused games. On a smaller scale, corrupt margraves may hire bandits to pillage their own villages, requiring heroes to step in. On a broader scale, a threat to the stability of the Republic of Tal'Dorei itself may force the characters to stop a megalomaniac bent on conquering the republic from within.

\subsubsection{Criminal Shadow War}
A criminal syndicate known as the \textit{Clasp} has lurked in Tal'Dorei's underworld since the foundation of the old Tal'Dorei Empire. It's almost become a cliché, for people across the land know the \textit{Clasp} as the stereotypical face of organized crime. Novels romanticizing their allure have been adapted into stage plays and bard-song in all corners of the republic.

And perhaps the \textit{Clasp} has bought its own good publicity. After the \textit{Cinder King's} famous conquest of \textit{Emon} some twenty years ago, it was the \textit{Clasp} that saved countless innocent people and led them to safety through their extensive network of smuggling tunnels beneath the city. It was the organized action of criminals, just as much as Vox Machina's heroic power, which saved Tal'Dorei—and the \textit{Tal'Dorei Council} knows it. The \textit{Clasp} has gone mainstream.

However, the \textit{Clasp} is under threat from the \textit{Myriad}, another criminal syndicate from the nearby continent of \textit{Wildemount}. The \textit{Myriad} sees Tal'Dorei's slow legitimization of the \textit{Clasp} as an opportunity to seize power in the nation's criminal underground. Violence is on the rise in Tal'Dorei's cities as the \textit{Myriad} uses local gangs as pawns in their bid to dethrone the \textit{Clasp}.

\subsubsection{A Land Thrice-Destroyed}
First there was the \textit{Calamity}, and the palaces of the gods were devoured by the earth. Then came the wars of \textit{Drassig}, and his defeat at the hands of \textit{Zan Tal'Dorei}. And then came \textit{Thordak}. The \textit{Cinder King} and the \textit{Chroma Conclave} shook the very foundations of Tal'Dorei. \textit{Thordak} "merely" reduced the city of \textit{Emon} to slag, but through fear, panic, and malice, the flames of destruction were spread across the continent by mortal hands. Today, a generation after the dawning of the Republic of Tal'Dorei, the ruins of three civilizations lie buried beneath your feet, and some secrets long to be unearthed.

\subsubsection{Mother of Invention}
Percival Fredrickstein von Musel Klossowski de Rolo III—Percy, to his friends—wanted to take the secrets of his inventions to his grave, especially the deadly power of his firearms. Unfortunately, because of his now-deceased archenemy Anna Ripley, the secret of gunpowder has begun to spread across Tal'Dorei. Now, isolated groups of tinkerers, both noble and unscrupulous are working to develop new technologies from Percival's research, and other black powder technologies from \textit{Marquet}. Some of these modern marvels will no doubt help the people of Exandria—but others, if placed in the wrong hands, are sure to enable the deaths of thousands.

Renaissance technologies, especially firearms, are not common in Tal'Dorei. If you want to keep your fantasy purely medieval, it's a simple matter to ignore the tinkerers, gunsmiths, and budding industrialists scattered across the continent.

\subsubsection{Uncharted Territory}
The corners of Tal'Dorei's maps are not all filled in. North of the \textit{Cliffkeep Mountains}, the trackless \textit{Neverfields} push up against the northern edge of the map. South of the \textit{Verdant Expanse} is the \textit{Rifenmist Peninsula}, where vast jungles and the brutal forces of the \textit{Iron Authority} prevent any explorer from filling in the map's southern tips.

To help you create stories that feel at home in Tal'Dorei, \textit{chapter 3} of this book contains a number of story hooks throughout its description of Tal'Dorei's landscape and settlements. See "\textit{Creating Adventures}" for more information on how to turn these plot seeds into full adventures that you can run for your gaming group.

While this book gives insight to one version of Exandria, its lore exists to serve you, not the other way around. Your campaign is your own creation—if you delight in adhering to \textit{Critical Role} canon and delving into its lore, then by all means, dig deeper! If you like subverting expectations, revealing that things believed to be true were nothing but myth and rumor all along, tell that story! And, if you love characters and places from \textit{Critical Role}, but want your stories to go to brand-new lands of your own creation, take inspiration from this book and steer your campaign in whatever direction you see fit. Never forget that you are the creator of your own version of Tal'Dorei. No matter who is playing at your table, even if it's Matthew Mercer himself, no one can tell you that your vision of the world is inaccurate—because it's yours now.

\chapter{History of Tal'Dorei}


\begin{DndReadAloud}
    \DndQuote%
    {}
    {\enquote{...then, the vast empty heart of endless shadow befell the light of the Seeker. Cascading from the pits was the font of life. The first moon formed the endless oceans. The second sun brought the lush soil. The third wind carried breath of life. The fourth flame ignited the Heart of Exandria.}}
    {}
\end{DndReadAloud}

Where did Tal'Dorei come from? The origin of this land is a time-shrouded question that everyone has pondered in some way, from the most exalted, world-traveling hero to the common farmer who has never traveled three miles beyond their village. Everyone has their own ideas on where Tal'Dorei, and all of the wide world of Exandria, came from. These varied thoughts all have some things in common, since they are all born of jumbled misunderstandings of ancient myths told orally from bard to bard, then passed from father to daughter, and eventually codified into religious re-writings of history that favor the teachings of a given deity.

To speak a full and accurate truth of Tal'Dorei is an impossible task, yet Tal'Dorei is known for creating folk who regularly do the impossible. What follows is the Myth of Exandria, an account based on the investigations and best efforts of Tal'Dorei's foremost scholars—the homegrown \textit{Alabaster Lyceum} of \textit{Emon}, and the international knowledge-seekers and anti-propagandists of the \textit{Cobalt Soul}. These two organizations have assembled this history by digging beneath the layers of dust and decay that hide battles long past, unearthing texts long thought burned and censored, and exhuming historical treasures from tombs of heroes long forgotten. The details are often debated, for the question of the land's origin remains—consuming the curious, calling those hungry for purpose, and fueling the business of adventuring to delve into the tantalizing unknown places of the world.

\section{Origins of Tal'Dorei}
Following the creation, and subsequent razing, of Exandria, a \textit{post-Divergence} world was now left to rise from the ashes and begin a new era. While every region had its own rebirth following the terrible destruction of the \textit{Calamity}, the continent that now bears the name of this setting shall remain the focus. The modern calendar began in the year 0 P.D. (\textit{Post-Divergence}), over eight hundred years before the writing of this text. Though the modern nation-state of Tal'Dorei did not emerge for some time after year 0, the history of Tal'Dorei truly began with the founding of \textit{Gwessar}.

\subsection{Gwessar}
At the dawning of modern history, what is now known as the continent of Tal'Dorei housed the germinating seeds of civilization. It was the hardy, dependable dwarves who best weathered the war between gods and mortals, beneath the \textit{Cliffkeep Mountains}. Dozens of dwarven redoubts rose and fell in the tumult of the \textit{Calamity} and the uncertain era that followed it. For many years, dwarves lived in the tunnels beneath the \textit{Cliffkeep Mountains}, leaderless and chaotic, until the various clans unified to form the subterranean city of \textit{Kraghammer}. Proud of their grand new home, the dwarf clans were content to remain beneath the earth. They saw the war-scarred surface world as a source of naught but misery and death, and their new underground domain as a place of undiscovered prosperity. This city was dug underneath the mountains, and its deep roots, though young, survived where more ancient ancestral halls were annihilated. Clan Jaggenstrike were the architects of \textit{Kraghammer's} hardy, unassailable redoubts, and easily became the first ruling clan of the \textit{Kraghammer} dwarves.

While the dwarves busied themselves excavating the world below, a group of elves appeared suddenly in the south. At the outset of the \textit{Calamity}, a society of elves used a powerful, obscure ritual to transport many of their people to safety on another plane: the Fey Realm, a plane of primordial beauty where elven legend says the The Arch Heart lovingly fashioned their people before placing them upon the face of Exandria. They returned at the turning of the ages under the guidance of an elven sorceress named Yenlara Alderwreath.

The elves rallied around her both for her defiant strength and for her compassion in the face of adversity. It was under Yenlara's wise rule that elven society once again began to reform. Upon returning to Exandria, Yenlara led her people westward to the \textit{Verdant Expanse}, an untamed forest born from the residual elemental power that lingered after the end of the \textit{Calamity}.

So beautiful was the forest they settled in that Yenlara's people were the first elves known to bear the now-widespread name of syn'alfen—wood elves, in Common parlance. Likewise, the beauty of their forest home inspired the name of their first settlement: \textit{Syngorn}, the great elvenhome that stands proud within the \textit{Verdant Expanse} to this very day.

As Yenlara's people began to explore beyond the forest, they found vast fields of grass emerging from the ash-darkened ground, with tall, snowcapped mountains visible on the horizon. Overwhelmed by the beauty of the restored world, the syn'alfen called their new land \textit{Gwessar}, the Fields of Joy, a name by which elves (and those enamored with elven culture) still call the continent of Tal'Dorei to this day.

The dwarves and elves are long-lived people, and when they struggled to rebuild their civilizations, there were those among them who still remembered the world that was. Humanity was not so fortunate. Human histories, written by wasteland warlords in fading ink on waterlogged parchment and vellum, did not survive the years. Yet, somehow, humanity endured. Several centuries after the Jaggenstrike dwarves began this period of renewal, a clan of humans braved the angry Ozmit Sea and sailed to \textit{Gwessar's} western coast from the continent of \textit{Issylra}, though whether or not they came from Vasselheim itself is unknown.

These people had the sea in their blood, and leapt from island to island for generations, but something called them to \textit{Gwessar}. The ruins of their first settlement—the port city of \textit{O'Noa}—still stand today. From \textit{O'Noa} the seafarers expanded outward, until all of the western shores flew their banner. In the north, they found fertile fields protected from the salt of the sea, and an inlet unmarred by looming rocks and dangerous reefs, and they began to build. They did not know their city would become the heart of a great empire. They did not know the glory and sorrow that would surround the city of \textit{Emon}.

\subsection{The Iron Rule of Drassig}
The rise of human colonies irritated the elves of \textit{Syngorn}. Forests that had stood for centuries fell under the axes of rash, short-lived beings who sought to exploit, expand, and propagate thoughtlessly. These tensions did not build to war, for long were the memories of the elves, but humanity remained a thorn in the side of elfkind, chafing at their hard-won efforts to protect their society and threatening to encroach upon their sacred woodlands.

As the first human civilization on this part of the world after the \textit{Calamity}, \textit{Emon} gave rise to a handful of self-entitled noble houses who went on to establish civilization in the name of the The Lawbearer—but those who wrote the game stacked the deck in their favor. Corruption spread through the upper echelons of \textit{Emon}, and power-hungry politicians seized each new and valuable resource that was discovered in their bountiful new kingdom. They turned their citizens against each other, forcing them to fight for scraps while they hoarded the lion's share.

\textit{Emon} was a political war zone, and the greatest warrior of them all was a loudmouthed braggart and cunning oligarch named Warren Drassig. The chaos and mistrust in \textit{Emon} allowed Drassig and his agents to seize power and transform the realm into the Kingdom of Drassig, with Warren himself as its supreme monarch. Drassig was quick to sever any remaining connections with the elves of \textit{Syngorn} and make new alliances with the dwarves of \textit{Kraghammer}, marrying Drassig's autocratic power with the \textit{dwarves'} immense material wealth.

The elves were furious, but the ambassador from \textit{Syngorn} to \textit{Emon}, an idealistic grandson of the still-living Yenlara, hoped to resolve this diplomatically. Upon arrival, he was apprehended, tortured, and slain. This final act of treachery drove \textit{Syngorn} to arms, and the continent erupted into a terrible, protracted war between Yenlara's kin and Drassig's bloodline known as the \textit{Scattered War}.

\subsection{The Scattered War}
The Scattered War lasted for thirty-two years, and is recalled in bard-song as the Time of Shrouds. The war spanned the \textit{Cliffkeep Mountains} and parts of the \textit{Verdant Expanse}, with the human colonies spread throughout the soon-to-be-warring territories. \textit{King Drassig} encountered little resistance as he conquered each village, town, and city in turn, giving these isolated settlements no way of warning one another of the warlord's advance.

No common warmonger, \textit{Drassig} was a cunning, heartless master of psychological warfare. Weeks in advance of his armies, he seeded spies throughout the land. They traveled swiftly through the secret dwarf-tunnels that crisscrossed the \textit{Cliffkeep Mountains}, and infiltrated settlements throughout the region, turning the people of each settlement against their own leaders, eating at them from within, poisoning them against one another. By the time \textit{Drassig's} armies reached their targets, they were already poised to fall. The lights of human towns and elven groves were snuffed out with equal savagery.

Yet, a mere nine years after the war began, King Warren Drassig died.

It happened at \textit{Torthil}, an abandoned village in \textit{Gwessar's} heartlands. For years, rumors had grown that \textit{Torthil} was a haven for war refugees. The king and his soldiers fell upon the village with a vengeance, seeking survivors to slaughter. As they regrouped in the town square, \textit{Drassig} stood among them, blade raised, eyes wild, and voice booming, to rally their bloodthirsty spirits. It was at the height of his wicked speech that the first arrow struck, followed shortly by a volley that clouded the sky. After nine years of suffering, Yenlara's wood elves and the rebellious humans of the scattered colonies had joined forces and formed an alliance against their common enemy.

After Warren \textit{Drassig} fell, he was succeeded by his eldest son, Neminar \textit{Drassig}. Neminar shunned the crude, brutal methods of his father, instead finding his interests in more sinister powers. He became known as "Neminar the Black-Fingered," as his forays into necromancy left one arm withered and tainted. Nevertheless, his weaponization of forbidden magic elevated the threat of \textit{Drassig's} war machine. King Neminar, bent on bloody-minded vengeance, led his army back to \textit{Torthil}, where the corpses of the elven and human traitors were piled high upon the outer walls. Their pyres ignited the night as the entire city was reduced to rubble. And so the Scattered War continued, unabated by the death of Warren Drassig.

The rule of Neminar marked the darkest days of the war. He introduced tactics and magic that drove his soldiers beyond human limits, leaving their bodies altered and mutated by foul necromancy. Their warped minds craved bloodshed and domination, and their bodies needed no sustenance to fuel their bloodlust. \textit{Drassig's} forces became the perfect weapons of war, their ranks blessed by the divine touch of the The Strife Emperor—the first \textit{Betrayer God} to extend his power to mortals beyond the Divine Gate. All the gods soon learned that though they could not walk upon Exandria themselves, they could still grant magic to their faithful.

In the face of such terrible power, the alliance of humans and elves bolstered their ranks with vengeful orphans and eager heroes of all ancestries who suffered under \textit{Drassig's} tyranny. One such hero who came to be instrumental in the coming conflict was \textit{Zan Tal'Dorei}. A human who rose from the harried streets of \textit{Syngorn}, Zan quickly showed her mettle as both a warrior and an inspirational leader. Rallying the broken ranks of the resistance, Zan lured Neminar and his lead forces into the \textit{Verdant Expanse}.

Neminar's undoing was an arrogant attack upon the \textit{Shifting Keep}, a \textit{Syngornian} outpost protected by the slippery magic of the Fey Realm. The illusory fortress vanished, leaving the \textit{Drassig} army without a target and vulnerable to ambush. The forest itself seemed to lend its fury to the weapons of Zan's warriors, and Neminar and his blighted soldiers were crushed in one swift stroke.

Word of Neminar's defeat and death spread like wildfire across \textit{Gwessar}, and the banner of \textit{Zan Tal'Dorei's} rebel army became a symbol of hope for all of the oppressed peoples of the realm.

Even so, the last years of fighting were still ahead of them, for the late King Warren Drassig's youngest son still stood to take power, and swore to avenge his family.

\subsection{Battle of the Umbra Hills}
\textit{King Trist Drassig}, second-born son of the despotic Warren \textit{Drassig}, took the throne with unearned confidence. So enamored was he with arrogant visions of victory that he ignored the truth of his tenuous rule: he was neither as brilliant nor as charismatic as his father, his armies were stretched thin, and his people rioted. The rebels, led by the young warrior \textit{Zan Tal'Dorei}, won battle after battle, and before long, King Trist found his army backed against the imposing base of the \textit{Cliffkeep Mountains}.

In the wake of certain victory, Zan and her rebels, allied with the elves of \textit{Syngorn}, pursued \textit{Drassig} to this valley, but were met with fiends amidst the ranks of their enemy. King Trist had a secret weapon—through his family's dealings with the The Strife Emperor, the spawn of the Betrayers had returned to the world. They spilled into the battlefield like a river of nightmares, and in minutes, the surrounding hills ran dark with blood and ichor, and the bodies of humans and demons alike littered the battlefield. Yet, despite all odds, the hero Zan defeated King Trist, ending the \textit{Drassig} bloodline—and with it, the demonic pact the \textit{Drassigs} had made. It is said that the grass and flowers of the \textit{Umbra Hills} grow black and burnt as an echo of this battle, their sap coursing with the fiendish blood that was spilled that day.

\section{Tal'Dorei Ascendant}
Society had all but collapsed once more amidst the \textit{Scattered War}. Within the \textit{Verdant Expanse}, the triumphant rebels assembled a council of trusted and proven minds, but the people were accustomed to a singular leader, a king. In the hearts of the people, it was not the crown that had failed them, but the bloody-minded family that wore it. The council assented to the will of their people, and nominated the war hero Zan Tal'Dorei to take the throne. She humbly accepted the role, but refused to take the title of king or queen. Her name, she asserted, was Zan. After some debate, she eventually relented to being called Sovereign Tal'Dorei, if such formality was required.

Again in spite of Zan's protestations, the council unanimously agreed the realm should be renamed Tal'Dorei in her honor. Power was divided between the sovereign and the council, both of whom ruled from \textit{Emon}. From there, \textit{Emon's} leaders expunged any remnants of \textit{Drassig's} reign from their city, while the allied armies of Tal'Dorei and \textit{Syngorn} did the same throughout the realm. The leading clans of \textit{Kraghammer} pleaded for clemency, claiming they were coerced to serve \textit{Drassig}. They spent many years paying reparations and making amends, but the fractured trust between \textit{Gwessar's} peoples took generations to heal.

\begin{DndReadAloud}{}

\enquote{
    \small\itshape
    A king should not stand apart from their people. A ruler must know sacrifice, for how else are they to know those who till the fields, who tend to the cattle, and who raise our children? A ruler should know pain and sorrow, for war is bloody and cruel. A ruler who asks their people to fight is one who asks them to die, and no person should die for a selfish king or an unjust cause. A ruler should know justice, but also mercy, for not all crimes are committed with evil intent, but rather in grief and shame. I am no king. I am the peasant in the fields. I am the painter at the easel. I am the young soldier donning their father's rusted armor. You will not find me sitting upon a bejeweled throne disdaining my supplicants, but working beside your mothers and fathers, fighting beside your valorous children, and aiding those who have suffered in the name of unjust kings.
    
    Henceforth, Tal'Dorei shall not be just a name: it is a principle. From this day, Tal'Dorei is a principle that stands for honor, an everlasting honor that binds its people. A principle that is a beacon for the lost, the helpless, and the forgotten. Tal'Dorei is a principle that you should declare with honor as you stand side-by-side with the people of Tal'Dorei as kin. Tal'Dorei is not a name; it is to be proud, it is to have honor, it is to be loyal, and loving, and kind. Tal'Dorei is not a name; it is a people. And you are the people of Tal'Dorei! I refuse the title of king. If you wish me to rule, I shall do so as a sovereign of the people—and if you permit me, if you accept me as such, the Sovereign of Tal'Dorei! For the Honor of Tal'Dorei!
}
\end{DndReadAloud}

\subsection{The Icelost Years}
Not two years into Sovereign Zan's reign, a cataclysmic evil tested her mettle and the skill of the council. A creature from the Elemental Plane of Ice slunk into the Material Plane through a hitherto unknown rift between the planes. It appeared in a mystical forest of eternal winter known today as the \textit{Frostweald}. From its hibernal abode, this elemental spirit spied upon the untested realm of Tal'Dorei and saw an opportunity. The spirit dove back to its extraplanar home and informed its master of a land ripe for conquest. Errevon the Rimelord, a cruel elemental so mighty that he is speculated to have been a scion of the Primordials, one that escaped the gods' wrath in the \textit{Founding}, seized this opportunity.

On the next winter solstice, Errevon tore wide the \textit{Frostweald} rift and commanded the endless blizzards of his home plane to ravage the prosperous land around it. Tal'Dorei's southern reaches, thus weakened, were easy pickings for the Algid Legion, a mighty army of frost giants and animate blizzards led by the Rimelord himself. Their relentless advance covered all they touched in snow and death.

Countless brave warriors and civilians alike fell to the invaders' gelid weapons, for they were unprepared for such an onslaught so soon after the \textit{Drassigs'} protracted war. The storm of ice widened year by year, consuming the Tal'Dorei heartland as Errevon claimed all lands that fell beneath the ice as his own. North of the rift, he built a towering citadel of frost, forcing those he conquered to swear fealty in exchange for warmth and liquid water. The Rimelord's tyranny lasted for three long years as the nascent \textit{Council of Tal'Dorei} struggled to convince \textit{Syngorn} and \textit{Kraghammer} to ally for the sake of the realm. This monumental diplomatic mission completed, the combined might of Tal'Dorei, \textit{Syngorn}, and \textit{Kraghammer} assaulted Errevon's Citadel and forced the Rimelord back to the rift from whence he came. It was in this time of need that the druidic Ashari people first revealed themselves to the people of Tal'Dorei, and used their unrivaled elemental magic to seal the rift once and for all. The snow began to melt, the fortress toppled into the nearby \textit{Foramere Basin}, and the people celebrated their freedom from the oppressive cold. This victory is now celebrated annually throughout Tal'Dorei as the Winter's Crest festival.

\subsection{Thordak, the Cinder King}
\textit{Zan Tal'Dorei} passed peacefully after a long and happy rule, and was succeeded by a long line of her descendants. Many generations of peace ensued, and the Tal'Dorei Empire formed as dozens of tiny city-states across the \textit{Dividing Plains} flocked to the stable rule and heroic reputation of Tal'Dorei in the wake of the \textit{Icelost Years}.

The adroit statecraft of the \textit{Council of Tal'Dorei} maintained cordial alliances with both \textit{Syngorn} and \textit{Kraghammer}, but, despite their best efforts, couldn't keep the elves and dwarves from regressing to hostility over ancient quarrels. In the time of Sovereign Odellan Tal'Dorei, word of a shadow in the south came to the ears of the council. Contact with numerous outlying townships at the edges of the empire halted abruptly, and traders that made circuits from \textit{Emon} to the \textit{Rifenmist Peninsula} in the south vanished without a trace. The council was slow to act in defense of territories that routinely flaunted imperial law, but were finally driven to action when allied villages south of the \textit{Verdant Expanse} begged for protection against a "nightmare of fire and malice." The dour Sovereign Odellan blocked the council's relief efforts, claiming the economic cost of sending a regiment that far south to outsider communities was a misuse of resources, especially without proof of a threat. It wasn't until two years later, when reports of a powerful red dragon reached the ears of the sovereign, that the \textit{Council of Tal'Dorei} was able to act.

By then, much of the \textit{Mornset Countryside} had been reduced to ash by a self-obsessed and power-hungry red dragon named \textit{Thordak}, who preferred the self-styled moniker of "the Cinder King." According to \textit{Cobalt Soul} records from the Temple of the Mentor in the lands of \textit{Marquet}, the \textit{Cinder King} was presumed dead nearly two hundred years prior, killed by the brass dragon Devo'ssa—yet it could not be denied that \textit{Thordak} survived, and that he sought new peoples to enslave, ones with fewer defenses.

The mighty armies of Tal'Dorei marched finally south, but morale was low. They marched toward fiery, painful death, and soldiers grumbled of a heartless sovereign who luxuriated in his castle in \textit{Emon} while his people put their lives on the line for the realm. They finally clashed with \textit{Thordak} and his mercenary armies south of the \textit{Stormcrest Mountains}, but the forces of Tal'Dorei were soon forced into a panicked, disorganized retreat as the rocky, exposed terrain worked to the advantage of their airborne foe.

The valiant soldiers of Tal'Dorei would surely have been brought to utter ruin were it not for the heroic intervention of a band of young—yet able—adventurers. Led by \textit{Arcanist Allura Vysoren}, who sits today upon the \textit{Tal'Dorei Council}, they cornered the \textit{Cinder King} alone and defeated him—though his power was so great that they could not manage to kill the dragon outright. Instead, \textit{Allura} and her companions used powerful magic to imprison \textit{Thordak} within the Elemental Plane of Fire forever. The Fire Ashari of distant \textit{Issylra} vowed to watch over \textit{Thordak's} prison until age could consume even him.

\subsection{The Chroma Conclave}
Forever lasted a mere sixteen years. In that time, life returned to mundane quarrels over the price of bread and taxes. During this time of quietude, an old and unexpected ally of \textit{Thordak's}, a green dragon called Raishan the Diseased Deceiver, found \textit{Thordak} and conspired to release him from his prison so he could conquer the land that had once rebuked him. Raishan formed a plan and forged an unprecedented alliance with three other power-hungry chromatic dragons: Umbrasyl the Hope Devourer, Brimscythe the Iron Storm, and Vorugal the Frigid Doom.

After infiltrating the Fire Ashari of Pyrah for nearly five years, Raishan finally unlocked the seal between the planes that bound \textit{Thordak} and set him free. \textit{Thordak's} physical body was warped by the fiery energies of his prison. Now stronger than ever, he rallied his draconic allies, the Chroma Conclave, to destroy or conquer Tal'Dorei as they saw fit—starting with a retributive assault of catastrophic proportions on \textit{Emon}. Nearly three centuries of Tal'Dorei rule in the realm ended under the rule of Sovereign Uriel Tal'Dorei II, his life cut short when much of \textit{Emon} was reduced to rubble.

The attack not only killed the Sovereign Uriel, but also scattered the \textit{Council of Tal'Dorei}, throwing the realm into chaos. The remnants of the civilized lands were divided up and taken as trophies by the Conclave, while the capital \textit{Emon} was ruled by the \textit{Cinder King} himself, with Raishan secretly manipulating the course of his rule from the shadows.

The \textit{Alabaster Lyceum} need hardly introduce the heroes of this story. One by one, the members of the Conclave fell to the might and cleverness of a band of misfit warriors known as Vox Machina, their might bolstered by a number of the lost \textit{Vestiges of Divergence}, artifacts of incalculable power forged amidst the chaos of the \textit{Calamity}. After many adventures, these heroes gathered their allies and stormed the capital of \textit{Emon}, slaying \textit{Thordak the Cinder King} and freeing its people. Upon discovering the machinations of the real mastermind, Raishan the Diseased Deceiver, with whom Vox Machina had held a tenuous alliance, they gave chase and finally slew this last standing member of the Chroma Conclave, ending their reign over the land and restoring rule once more to the \textit{Council of Tal'Dorei}.

\section{A Fledgling Republic}
The final act of Sovereign Uriel Tal'Dorei II's life was to abdicate the throne and end the line of sovereigns—permanently. To analyze the historic Last Sovereign's final acts is the duty of biographers, not chroniclers, yet it can readily be surmised that he saw flaws inherent to his station. Uriel was the son of Odellan Tal'Dorei, the sovereign who singlehandedly prevented the council from swiftly handling the \textit{Cinder King} when first he rose in the wilderness south of the empire.

How could the Last Sovereign not feel that it was his office's immense power to overrule the council that ultimately cost thousands of loyal citizens their lives? How could he not compare the singular power of the sovereign to the abuses of power committed by \textit{King Drassig} in ancient times?

Regardless of his private feelings, Uriel II's last act as sovereign transformed his family's imperial dynasty into the Republic of Tal'Dorei. \textit{Zan Tal'Dorei} was reluctant to accept the singular role of rulership at the dawning of the nation, an ideal she had to compromise upon to ensure stability for her fledgling nation. Uriel's relinquishment of power is seen today as a fulfillment of that ideal. It will be the duty of future historians to judge if his decision stood the test of time.

\subsection{The Time of Regrowth}
The \textit{Tal'Dorei Council} has guided the fate of the nation since Vox Machina slew the \textit{Cinder King}. For most of the past two decades, the goal of the council has been to rebuild that which was lost—and to assure far-flung territories that the nascent republic is just as stable as the empire ever was.

Recovery has been slow in both Tal'Dorei's heartland and its most distant reaches. Even though cities like \textit{Kymal} and \textit{Stilben} were not direct targets of the dragons' rampage, the age of unrest that followed made these small cities easy targets for ransacking by roving bandits and corrupt magistrates alike. Even today, outlying settlements are slow to trust the good intentions of the \textit{Tal'Dorei Council}.

Major cities like \textit{Westruun} and \textit{Emon} have been the greatest beneficiaries of the Time of Regrowth, for even though they suffered the worst of the \textit{Conclave's} destruction, the dragons also kept their lairs close by, allowing vast amounts of plundered wealth from their treasure hoards to flow rapidly back into these cities' coffers. Over the last twenty years, these cities have hired legions of arcanists to magically restore much of what was lost.

However, perhaps the one city that can say it directly benefited from the \textit{Chroma Conclave's} reign of chaos is \textit{Whitestone}, a formerly independent city-state that is now a member of the republic. \textit{Whitestone} took in countless refugees from \textit{Emon} after the \textit{Conclave's} initial attack, and countless more in the aftermath. Sequestered in Tal'Dorei's cold northeastern reaches, \textit{Whitestone} is an isolated city, but it has rapidly become a powerhouse of political, cultural, artistic, and economic influence—to the chagrin of some traditionalists in \textit{Emon}.

\subsection{Apotheosis Thwarted}
The world was troubled by a disturbance in \textit{Issylra} a year after the \textit{Chroma Conclave's} defeat, as a lich attempted to achieve godhood and crush the Dawn City of Vasselheim in one fell swoop. His success would have meant catastrophe for the Material Plane, leaving him the sole god on this side of the world's impenetrable Divine Gate. However, beyond hearing news of Vox Machina's involvement in this conflict, few in Tal'Dorei were concerned by rumblings on the other side of the world.

Though the lich's defeat confined him to a remote demiplane, there is no way to strip him of his stolen divinity. Occult worshipers yet scurry throughout Tal'Dorei's darkest corners, seeking the favor of a lich-god known as The Whispered One.

\subsection{Seeds of Peril}
A generation has passed in relative peace. Yet all who have studied history know that each age's conflict sows the seeds of the next. The Republic of Tal'Dorei has proven it can prosper in peacetime, but it has yet to be tested by the pressures of war or a supernatural cataclysm.

After two decades of expensive reconstruction and princely investments in \textit{skyship} moorings, the Republic of Tal'Dorei is deeply indebted to arcanists across the land. Chief among them is the \textit{League of Miracles}, a consortium of thaumaturgists who have claimed untold fortunes from reconstruction bounties all over the country.

The simple fact is that the cost of rebuilding Tal'Dorei was far more than the council or local leaders could afford to pay. Far be it from the \textit{Alabaster Lyceum} to cast aspersions on so beneficial an organization, but there are fears among the Tal'Dorei people that the league has already begun to collect on what they are owed in forms other than gold, as more and more league mages receive preferential treatment from local margraves and even members of the \textit{Tal'Dorei Council} itself.

Beyond \textit{Emon}, the landscape of Tal'Dorei has been forever scarred by the presence of the \textit{Conclave}, especially the \textit{Cinder King}. His hellfire left marks on the land that will never heal—magical wounds that bleed fire and ooze chaotic magic. Smaller cities beyond the reach of the council have begun to succumb to sinister influences. Demons and devils crawl into the world, summoned by occult powers answering fell voices from beyond the Divine Gate. The \textit{Betrayer Gods} sense weakness within the world. The entropic fury of the The Chained Oblivion churns beneath \textit{Gatshadow}. The Graz'tchar, the decadent end calls from the \textit{Umbra Hills} for a new master. The The Cloaked Serpent raises armies south of \textit{Syngorn}. Cults of the now-divine The Whispered One spread like poison through the veins of Tal'Dorei's cities.

Yet despite these threats, the legendary heroes of Vox Machina have largely set down their blades, guns, and spells and taken up the mantle of leadership. More than twenty years have passed since their greatest adventures, and though the surviving members of Vox Machina are still canny rulers and cunning warriors, the \textit{Tal'Dorei Council} can no longer rely upon their skills as readily as they once did. Rumors of the heroes' permanent retirement swirl, inspiring fear and nostalgic sorrow in many of the people—but adventurers across the land see this as an opportunity to rise to the level of Vox Machina and beyond.

Tal'Dorei needs heroes—and this time, Vox Machina alone will not be enough.

\chapter{Factions and Societies}
The lands of Tal'Dorei are governed not just by their history, but by factions and faiths striving to achieve myriad goals, fighting for conflicting ideals, and vying for the power to come out on top. The gods of Exandria may be barred from walking upon the land in material form, but they can still put their thumbs on the scales of history through their mortal followers. Such followers include clerics with the power to cast mighty spells, paladins who wield blades alight with idealistic verve, and untold thousands of common folk whose voices can affect the decisions of even the most powerful leaders of the land.

A number of factions hold just as much sway over the daily lives of the people as the gods and their clergy—perhaps even more. These factions range in might and influence from the greatest of the great to the smallest of the small. The \textit{Tal'Dorei Council} institutes political policies that affect everyone in grand, sweeping, and often inscrutable ways, while the \textit{Golden Grin} is a group of folk heroes who seek to subvert evil across the land. Whether these factions are governments or just tiny organizations, they can add depth and conflict to your Tal'Dorei campaign.

While the dominant cultures of the continent have molded structure and law from the rough clay left beyond the \textit{Divergence}, these very cultures are comprised of a number of smaller factions that together create the whole. Whether they be a union of opportunists, ever-vigilant to reap the benefits of an exploitable populace, or bands of like-minded politicians, striving to maintain control in a land beset by chaos and danger, each faction seeks to drive and manipulate the social and political direction of Tal'Dorei. The alliances and tensions that can arise between the following societies can guide your adventurers' destiny across the spectrum of heroism and villainy.

\section{Tal'Dorei Council}
When the new reign of Tal'Dorei was established three centuries ago, \textit{Zan Tal'Dorei} deposed a brutal tyrant. Though she was beloved by the people, she did not wish to be given unchecked or un-counseled rule over the realm. Nevertheless, her inner circle obeyed the will of the people and crowned her Sovereign of Tal'Dorei. She took to her new role well and swiftly expanded her council by appointing some of her trusted allies—and even some dissenting voices—to form an assembly of scholars, generals, and philosophers to govern at her side. This \textit{Council of Tal'Dorei} worked alongside their sovereign to rebuild the realm following the end of \textit{Drassig's} line, and were instrumental in the following prosperity of the nascent Tal'Dorei Empire.

Twenty-six years ago, Sovereign Uriel Tal'Dorei II, the latest and last of the Tal'Dorei line, publicly ended the monarchy and passed power over to the council, formally divesting his line of power. Tragically, this ceremony was interrupted by the arrival of the four ancient dragons known as the \textit{Chroma Conclave}. The dragons destroyed the city of \textit{Emon}, and the man who had until moments ago been sovereign was one of the first casualties. Uriel Tal'Dorei is survived by his wife Salda, and three children—\textit{Odessa}, Illiya, and Gren—who have used their family's still considerable wealth and power to become well-known public figures in \textit{Emon}. \textit{Odessa} herself won a seat on the council in a landslide victory.

In the generation since the fall of \textit{Emon}, the council has risen to the challenge of rebuilding and governing their new republic. This state's power is highly decentralized; each individual city-state in Tal'Dorei largely acts according to its own interests, with the council in place to ensure that pursuing those interests benefits the common good. Some critics of the republic say that the \textit{Tal'Dorei Council} is little more than a glorified city council for \textit{Emon}, though it is just as often praised for governing lightly as a rule, and for using the spy network helmed by the council's Master of Law to root out corruption as needed.

In addition to the historical Six Masters whose positions have existed on the \textit{Tal'Dorei Council} since Zan's time, a variable number of delegates also cast votes in matters of state. Each settlement in Tal'Dorei whose population numbers at least five thousand is offered the option to elect delegates to the council if they wish to engage in national politics. The settlement is permitted one delegate per five thousand people, and the means of their election is left to the settlement's local government. Additionally, any settlement with a population of over twenty thousand can appoint an ambassador to \textit{Emon}; these ambassadors also sit on the council and have voting power. Once in \textit{Emon}, each township's delegation is expected to negotiate with one another and find agreement. If an accord can't be reached, they must go through the slow process of drafting proposals to be voted on by seated members of the council.

Though \textit{Kraghammer} and \textit{Syngorn} remain sovereign city-states, the invasion of the \textit{Chroma Conclave} convinced lawmakers that even tighter alliances between the peoples of Tal'Dorei were necessary. One ambassador each from both \textit{Kraghammer} and \textit{Syngorn} now also holds a seat with voting power on the council. Each city-state has likewise sent a delegation to \textit{Emon} to vouch for their interests.

\subsection{Goals}
Following \textit{Zan Tal'Dorei's} coronation, the \textit{Tal'Dorei Council} existed to aid and inform the sovereign on how best to execute their duties. The council has evolved to manage the affairs of the city of \textit{Emon} and provide a private forum for delegates from across the land to shape the future of Tal'Dorei. The council is made up of six masters, whose roles have existed since the time of \textit{Zan Tal'Dorei}, and a variable number of ambassadors from all city-states across the continent who wish to participate in the coalition of power that is the Republic of Tal'Dorei.

The Six Masters are elected by a popular vote of all city-states who have sent an ambassador to the council. These masters are each beholden to a domain of responsibility, and they either inherit the existing infrastructure of civil servants, or restructure with the approval of the rest of the council. The domains are divided as such: Development, Arcana, Law, Commerce, Information, and Defense. Additionally, the first business of the council upon delivering or receiving a declaration of war is to elect a Seventh Master: the Master of War, who is appointed to the council by a vote of all sitting masters, ambassadors, and delegates. The council members are responsible for these domains throughout the city of \textit{Emon} and work closely with community leaders in other cities to maintain order.

\textbf{The Master of Development} works with the masonry and carpentry guilds to oversee all major infrastructure projects within \textit{Emon}, as well as allocating resources to projects needed in settlements that have a delegation in the council. In the past two decades, the current Master of Development has added an unprecedented eight Lieutenant Masters of Development to his staff to serve the reconstruction needs of each city in the republic.

\textbf{The Master of Arcana} keeps vigil over the lawful use of magics within Tal'Dorei, informs the council on matters beyond the mundane, and occasionally works with the \textit{Alabaster Lyceum's} Headmaster. This council member is currently embroiled in debates over the danger that the \textit{League of Miracles} poses to organizations like the \textit{Arcana Pansophical}.

\textbf{The Master of Commerce} works with all major merchant guilds and trading networks within \textit{Emon} and beyond. They are to maintain healthy trade at home and abroad, especially between squabbling delegates, and establish the government's annual budget.

\textbf{The Master of Information} manages foreign diplomacy, and also serves as the head of a network of spies that, these days, is supplemented by the skilled eyes of the \textit{Clasp}. Their work as spymaster and chief diplomat seeks to root out threats to the republic from within and without.

\textbf{The Master of Defense} is head of the \textit{Emon} Guard, High Commander of \textit{Fort Daxio}, and manages the allocation of federal resources to city guard groups throughout Tal'Dorei. They work closely with the Master of Law to mobilize armed troops throughout the country if supernatural or nonmilitary threats arise during peacetime.

\textbf{The Master of War} is a temporary position only instated during times of open or impending warfare. As Tal'Dorei has only a small standing army, this newly appointed master then assembles a military from city guards and town militias across the nation, mobilizing the elite warriors of \textit{Fort Daxio} as needed. When peace is declared, the position is dissolved until it is called for again.

\subsection{Relationships}
The council maintains free travel, free trade, and free communication between its constituent settlements, and between the republic and its allies—both on Tal'Dorei and abroad. The incursion of the \textit{Chroma Conclave} has united the oft-feuding city-states of \textit{Syngorn} and \textit{Kraghammer} like never before, and all of Tal'Dorei now enjoys free trade and open borders between even its most remote enclaves. Thanks to the actions of Vox Machina, the previously solitary city-state of \textit{Whitestone} has joined the republic, to the benefit of both.

The \textit{Clasp} was once considered a serious threat to the council's rule, but times have changed. These days, \textit{Clasp} Spirelings in \textit{Emon} frequently make public appearances before the council on behalf of their own organization and their often unscrupulous business partners. This is to say nothing of the \textit{Clasp's} deep involvement in the Master of Information's spy ring.

The \textit{Tal'Dorei Council} has comfortable trade and diplomatic agreements between the realm of \textit{Othanzia} in \textit{Issylra} and the loosely federated cities of the Clovis Concord in \textit{Wildemount}. The recent conclusion of a bloody war between the Dwendalian Empire and the Kryn Dynasty in \textit{Wildemount} has given the \textit{Tal'Dorei Council} cause to send ambassadors to both nations to establish embassies, in hopes of creating an historic arrangement with all of the sovereign nations of that oft-maligned continent. On the continent of \textit{Marquet} to the southwest, the great city of \textit{Ank'Harel} has been friendly and open when it comes to trade and commerce, but has shown no interest in any political relationships beyond that.

\subsection{Figures of Interest}
These figures are major members of the \textit{Tal'Dorei Council} that you can use as NPCs in your game. You should also take inspiration from these figures to make up new NPCs for your own campaign.

\subsubsection{Guardian Tofor Brotoras}
\textit{\textit{Dragonblood} Master of Defense}

\noindent
Stern, honorable, and quick to anger, Tofor is proud to serve on the council. After losing many friends (and one of her own arms) in the struggles against the \textit{Chroma Conclave}, she took the vulnerability of the city to heart, swearing never to allow such a thing to happen again. She is a beloved hero of the war against the \textit{Chroma Conclave} and has served on the council since before the death of the Last Sovereign. Most folk in \textit{Emon} suspect she'll win reelection until the day she dies.

\subsubsection{Seeker Odessa Tal'Dorei}
\textit{Human Master of Information}

\noindent
When the eldest daughter of the Last Sovereign campaigned for the Master of Information's seat, its current holder was Seeker Assum Emring, a loyal servant of her father's. He bowed out of the race, and his blithe quote made headlines in \textit{Emon}: "I taught her everything she knows." She won in a landslide. Odessa wields the dual blades of Assum's spy network and the mercenary agents of the \textit{Clasp} with practiced finesse. Nevertheless, her name and the legacy of her father's martyrdom weigh heavily upon her. As Seeker, Odessa is calm, methodical, and laconic, but she carries a dignified sorrow.

\subsubsection{Arcanist Allura Vysoren}
\textit{Human Master of Arcana}

\noindent
\textit{Allura Vysoren} has held the position of Master of Arcana since before the \textit{Chroma Conclave's} attack, and is held in the esteem of Tal'Dorei's people—mages and common folk alike. As a member of the \textit{Arcana Pansophical}, she is also increasingly worried as more and more of Tal'Dorei's politicians fall in debt to the \textit{League of Miracles}. After decades of investigating supernatural abnormalities herself, age has finally forced her to limit her adventuresome excursions. She now is one of the council's foremost employers of adventurers.

\subsubsection{Arbiter Brom Goldhand}
\textit{Human Master of Law}

\noindent
A stone-faced, gray-haired, flame-scarred cleric of the The Knowing Mentor, Brom is as mysterious as he is proficient with law. One thing is clear: he is dedicated to keeping the realm just and orderly, and there are few who would dare challenge his judgement.

\subsubsection{Hearthmaker Edelbern Cleareyes}
\textit{Goliath Master of Development}

\noindent
As a goliath of the \textit{Rivermaw nomads}, no one expected Edelbern to have an interest in civic architecture. Yet, after growing up in \textit{Westruun} while it was occupied by the Herd of Storms some twenty years ago, he found that he preferred the beauty of stone towers to towering pines. He also possessed a remarkable gift of mathematical acuity, and used his skills to set \textit{Westruun}—and eventually all of Tal'Dorei—on the path of reconstruction. Edelbern is also idealistic to a fault, and has, in his dedication to restoration, been manipulated into handing numerous lucrative contracts to the \textit{League of Miracles}, allowing them to amass greater and greater power.

\subsubsection{Coinmistress Vex'ahlia de Rolo}
\textit{Half-elf Master of Commerce}

\noindent
As a member of Vox Machina and hero of the realm, Vex'ahlia could help Tal'Dorei in any number of ways. She already held a ceremonial title in \textit{Whitestone} as \textit{Grand Mistress of the Grey Hunt}. Yet when she and her husband pushed for \textit{Whitestone} to be formally recognized as a constituent city of the Republic of Tal'Dorei, all who knew her name (and her legendary financial savvy) clamored for her to be elected to the council as Coinmistress. She took to the duty with aplomb, and fights hard to ensure that \textit{Emon} never gets the bad end of a deal when gold is involved.

\subsubsection{Ambassador Syldor Vessa}
\textit{Elf ambassador from \textit{Syngorn}}

\noindent
Vex'ahlia's elven father has served as his city's ambassador to \textit{Emon} for well over a century, and requested to be appointed a seated ambassador as soon as the office was created. Though he is otherwise a staunch traditionalist, Syldor was inspired by his brave children to move past his old racial prejudices and drive his homeland to be less bigoted and xenophobic—learning from his own mistakes as a father to two half-elven children. His relationship with his elder daughter is still strained, and it is a newsworthy occasion when they vote in accord.

\subsubsection{Ambassador Deirdrik Greyspine}
\textit{Dwarf ambassador from \textit{Kraghammer}}

\noindent
The youngest son of the ruling House Greyspine of \textit{Kraghammer}, Deirdrik has little hope of being elected Ironkeeper. Instead, this ambitious young dwarf—a scant seventy years old—left his home to win political renown in the national arena. He is a romantic idealist at heart and always votes for bold new policies, but doesn't always think through the consequences of his actions.

\subsubsection{Ambassador Keyleth of the Ashari}
\textit{Half-elf ambassador from \textit{Zephrah}}

\noindent
A legendary archdruid and member of Vox Machina, Keyleth is both leader of the Air Ashari and ambassador for all Tal'Dorei's \textit{Ashari} peoples. She despises politics, and is focused entirely on improving the wellbeing of as many people as possible—as well as involving \textit{the Ashari} on the national stage. She spends most of her time in \textit{Zephrah}, with her people—including her father, Korrin, and her newly returned mother, Vilya—but she can and will arrive in \textit{Emon} at the drop of a hat, if Vex asks it of her.

\subsubsection{Delegate Kel'jaia Uloeh}
\textit{Gnome delegate from \textit{Syngorn}}

\noindent
Kel'jaia is a confident gnome with incredible strength and resilience—both physical and emotional. She is thought by all to be the council's first choice for Master of War, should the time ever call for it. She was born in the \textit{Rifenmist Jungle} and grew up amidst mercenaries who fought against the \textit{Iron Authority}. After her commander was killed in battle, she traveled north to \textit{Syngorn} to seek aid in her fight. She dreams of forging bonds of unity between her homeland and Tal'Dorei.

\section{Arcana Pansophical}
All of Tal'Dorei's greatest mages fought at the \textit{Battle of the Umbra Hills}. They saw the incursion of demons and devils firsthand, and were rightly terrified that the spawn of the \textit{Betrayer Gods} could ever return to the world. It was clear that the \textit{Divergence} had not wholly prevented supernatural evil from infringing upon the world, and that the mages of Exandria must prevent future "surprises" of this nature. The mistakes of the \textit{Age of Arcanum} must never be repeated. In light of this, \textit{Yurek Windkeeper}, high enchanter of \textit{Syngorn} and a well-respected arcanist, reached out to the realm's most trusted and capable mages, and together they formed a private society called the \textit{Arcana Pansophical}.

\subsection{Goals}
The chief priority of this fellowship of powerful arcane specialists is to maintain the Truth of Magic, a strict code enacted to never again allow the rampant misuse of magic seen in the \textit{Age of Arcanum}. General members are referred to as Guides, while the higher circle of members who deal in judgment and oversight are called Makers. At least one member is established at each place of arcane learning.

Membership is largely kept secret unless out of necessity, though the existence of the order is relatively common knowledge. Anyone found to be practicing magic is logged and routinely checked up on, as misuse of magic is grounds for arrest, judgment by the Makers, and punishment outside of the local law. Even so, the reach and sight of the \textit{Arcana Pansophical} are not without limits, and many dark and terrible dealings manage to escape their notice.

\begin{DndSidebar}[float=!b]{The Truth of Magic}
The \textit{Pansophical's} myriad rules on the use and misuse of magic are innumerable—detailed at length in centuries of amendments made to the full text of the Truth of Magic. Though the Truth's bizarre intricacies are left up to the GM, the Truth's preamble establishes three unbreakable edicts:

\begin{itemize}
    \item In the Eternal interest of the Preservation of Harmony within the realm of Exandria, we the Arcana Pansophical establish the Truth of Magic, and three Edicts from which none shall stray. Those Arcanists found in gross violation of the Truth will be punished by the full extent of this Order—outside of any local law.
    \item Though the school of Necromancy shall be restricted to none, in the interest of magical study and Understanding, the animation of the dead is a Violation of the Truth.
    \item The Arcane is a tool to be wielded for the good of the People. To use its Power in the pursuit of wanton Destruction or unjust Murder is a Violation of the Truth.
\end{itemize}

In addition, though the Jurisdiction of the Pansophical supersedes the power of local laws, a mage who willfully breaks the laws of the land is in Violation of the Truth.


\end{DndSidebar}

\subsection{Relationships}
The \textit{Arcana Pansophical} is respected, if somewhat feared, throughout the Republic of Tal'Dorei. \textit{Syngorn} and other elf-founded settlements see the \textit{Pansophical} as allies, as the roots of the society began within \textit{Syngorn} itself. The people of \textit{Kraghammer} are suspicious of the \textit{Pansophical's} legalistic Truth of Magic, and they fear that its members hold more loyalty to the \textit{Pansophical} than to their homeland.

Recently, a group of profit-driven arcanists splintered from the \textit{Arcana Pansophical} to form the self-styled \textit{League of Miracles}. Their involvement in the for-profit reconstruction of Tal'Dorei has made them popular among mages who see the \textit{Pansophical} as a stodgy, elitist institution. The \textit{Pansophical}, however, sees the \textit{League of Miracles} as a distasteful upstart that will soon buckle under its own ambition.

\subsection{Figures of Interest}
These figures are major members of the \textit{Arcana Pansophical} that you can use as NPCs in your game. You should also take inspiration from these figures to make up your own NPCs.

\subsubsection{Enchanter Yurek Windkeeper}
\textit{Elf wizard (enchanter) and Maker}

\noindent
Founder of the \textit{Arcana Pansophical}, the ancient archmage Yurek generally stays within \textit{Syngorn}, acting as the nexus of the order and a consultant to the \textit{Wardens of Syngorn}.

\subsubsection{Alchemist Oz Gruude}
\textit{Halfling alchemist and Maker}

\noindent
Gruude is an ornery old halfling who is more interested in traveling the world to study plants and minerals with latent magical properties than doing alchemical lab work. He is rarely seen in the halls of the \textit{Pansophical}.

\subsubsection{Seer Gloria Ios}
\textit{Human wizard (diviner) and Maker}

\noindent
Gloria is one of the great diviners of the age, and her duties have sent her to Vasselheim, where she acts as a foreign outreach for the \textit{Pansophical}.

\subsubsection{Magus Seanor Wiles}
\textit{Human wizard (enchanter), Maker, and tutor}

\noindent
A humble-looking traveler, Wiles delights in finding fresh and emerging talent and helping his new students refine their arcane gifts.

\subsubsection{Planerider Ryn}
\textit{Tiefling wizard (conjurer), Maker, and planar studies specialist}

\noindent
This former pupil of \textit{Yurek Windkeeper} has grown into a master of planar travel. She likewise has taken an interest in the flow of magic across the world by means of ley lines. She's constantly on the move, traveling from plane to plane to conduct her research.

\subsubsection{Tinkerer Quash Dentdruggle}
\textit{Gnome construct specialist and Guide}

\noindent
Obsessed with the development of autonomous and living arcane constructs, Quash is usually holed up in his laboratory within \textit{Kraghammer}.

\subsubsection{Arcanist Allura Vysoren}
\textit{Human wizard (abjurer) and Guide}

\noindent
As a member of the \textit{Tal'Dorei Council} and a retired adventurer, \textit{Allura} has helped maintain the political significance of the \textit{Arcana Pansophical} within \textit{Emon}. Even among her colleagues in the \textit{Pansophical}, her unpretentious strength commands great respect.

\subsubsection{Elementalist Drake Thunderbrand}
\textit{Dwarf wizard (evoker) and Guide}

\noindent
Driven to bind the physical elements of the chaotic planes to his will, Drake is a prodigy of new incantations in the field. He and \textit{Allura Vysoren} are lifelong friends, and together they're renowned for achieving the impossible.

\subsubsection{Realmseer Eskil Ryndarien}
\textit{Human archmage and Maker emeritus}

\noindent
Though he was a Maker for most of his life, Realmseer Ryndarien grew frustrated with the politics and limitations of the \textit{Pansophical}, eventually retiring to \textit{Westruun} to pursue his own interests. These days, he has found a way to rejuvenate his aged body and is now renowned as a handsome and aloof man about town—much to the annoyance of all who formerly knew him.

\section{The Ashari}
The Material Plane is surrounded on all sides by the churning chaos of the elements. This elemental power is invisible to most, but in some places, the veil between this world and the Elemental Planes grows thin—and sometimes is torn entirely asunder. These places where the Elemental Planes bleed into the world are called rifts, and they first appeared on a long-ago Celestial Solstice when the planes aligned. That day, elemental rifts of untold proportions burst open across Exandria. Wildfires, tsunamis, hurricanes, and earthquakes ravaged the lands, and the people of the world struggled to combat the elemental beings that poured forth from the rifts. The tide of rampaging elementals was eventually stemmed, but the world was gripped by fear—what would happen on the next solstice?

In the wake of the planar disturbance, a society of druids called the Ashari made a covenant to divide and form four separate orders to seal and watch over the sources of these unleashed energies. Most of the lesser rifts that opened in that time have been healed, but one rift to each of the four Elemental Planes remains in Exandria. The four orders of the Ashari settled near these rifts, watching over them and soothing them when they grow volatile. Of these orders, two are in Tal'Dorei, where they watch over the rift of earth and the rift of air. A dozen generations later, the Ashari still hold watch over these tumultuous locations.

\subsection{Ashari Culture}
Though the Ashari began as a single group of elementalists, their individual orders have each developed their own ways of living over the centuries.

The people of Pyrah value reactivity and coolness under pressure, for panic is death during an eruption. The Ashari of Vesrah have a close relationship with the marine life that surrounds their rift, and their actions always prioritize the health of Exandria's waters. \textit{Zephrah's} culture values the freedom of the open sky, and balances their sense of duty with a free-spirited sense of play. Finally, the society of \textit{Terrah} is rooted in planning for the future, and fixing what is broken with patience and humility.

\subsection{Goals}
Each Ashari order guards their respective rift between one of the elemental planes and Exandria. Here, they strive to understand, master, and ultimately subdue the ever-escaping elemental forces. They attempt to bring balance to the surrounding lands and heal the rift over time, though recurring cosmic events tend to reopen these wounds.

\textit{The Ashari} of Pyrah oversee the Fire Elemental Rift in the valley caldera of the Sunderpeak Mountain range in \textit{Othanzia}, upon the continent of \textit{Issylra}. \textit{The Ashari} of Vesrah watch over the Water Elemental Rift on the Islands of Anamn, within the Ozmit Sea. On the southwest edge of the \textit{Cliffkeep Mountains}, \textit{the Ashari} of \textit{Terrah} stand vigilant over the Earth Elemental Rift, while high atop the windy mountaintops of the \textit{Summit Peaks}, \textit{the Ashari} of \textit{Zephrah} keep the Air Elemental Rift. The Air Ashari also watch over the \textit{Frostweald}, where a being of elemental ice once broke through to cause untold destruction.

Due to their isolation, \textit{the Ashari} go to great lengths to maintain contact with their kin. Each order is led by a single elemental master, a title that is passed on by blood or honor for generations. In every new generation that comes of age, the young leader-to-be must embark upon an Aramente. This "noble odyssey" sends them across the world to meet with the other Ashari leaders and train until they have earned both the respect and an understanding of their hosts.

Upon returning, they ascend to the head of their order, while the elder retires. This journey often takes years, and should one such Ashari not return within a reasonable period of time, a next-of-kin or a suitable elemental adept is chosen to take up the Aramente.

\subsection{Relationships}
\textit{The Ashari} maintain peaceful, if sometimes strained, relationships between the orders, though inner conflict occasionally arises within them. Beyond this, much of life within \textit{the Ashari} is hidden from outsiders' prying eyes. In recent years, however, Keyleth of the \textit{Zephrah} has taken steps to unite \textit{the Ashari} with the rest of the people of Tal'Dorei. She has gone so far as to become an ambassador for \textit{the Ashari} people on the \textit{Tal'Dorei Council}. Part of this work has been expanding the crisis orbs system, which \textit{the Ashari} used to contact their distant kin in times of peril, to connect cities across Tal'Dorei in the event of another crisis on the level of the \textit{Chroma Conclave}.

\textit{The Ashari} of \textit{Terrah} nurture a seed of disrespect and mistrust for \textit{Kraghammer}, whose miners constantly disregard their warnings against reckless mining. This arrogance caused countless deaths in the most recent eruption of the \textit{Terrah} Rift—and led to the collapse of an entire mountain.

Each Ashari culture has its own favored deities. Many Ashari of Pyrah venerate the The Dawnfather, for his flames give life to Exandria. Many in Vesrah worship the The Moonweaver, for her work controls the tides. In \textit{Zephrah}, the faith of the The Wildmother is predominant, for her freedom is like that of the wind. And the people of \textit{Terrah} hold the The All-Hammer in great respect for his association with creation and the earth. There are many Ashari, as well, who prefer to revere the elements themselves, forgoing the worship of gods.

\subsection{Figures of Interest}
These figures are major members of \textit{the Ashari} that you can use as NPCs in your game. You should also take inspiration from these figures to make up new NPCs for your own campaign.

\subsubsection{Pa'tice, Heart of the Mountain}
\textit{Human leader of \textit{Terrah}}

\noindent
Elderly leader of the Earth Ashari, Pa'tice still shows vigor and strength to resist his age, having led his people through many very tough challenges. His stubbornness is legendary, and it's said even death has been forced to wait until Pa'tice sees it his proper time. Many believe that his druidic training will keep him alive until he himself chooses to let go of this mortal coil.

\subsubsection{Korrin of Zephrah}
\textit{Half-elf elder of \textit{Zephrah}}

\noindent
His wife Vilya was lost while she undertook the Aramente, and when the previous Voice of the Tempest died of illness, Korrin was called upon to act as Voice until his daughter Keyleth came of age and completed her Aramente. Korrin was shocked when Vilya suddenly returned, and now spends his days in disbelieving bliss, reunited with his long-lost wife. He is gentle with those who show honor, and merciless to those who threaten his people.

\subsubsection{Vilya, Who Escaped Death}
\textit{Half-elf elder of \textit{Zephrah}}

\noindent
Vilya disappeared on her Aramente nearly thirty years ago, and was presumed dead. Her daughter, Keyleth, took up her own Aramente ten years after Vilya left, and discovered from the Water Ashari of Vesrah that Vilya was killed by a \textbf{kraken}, and only her leg was able to be recovered. Less than a year ago, Vilya suddenly strode out of a tree in \textit{Zephrah}, her leg replaced with a druidic prosthetic of her own creation, shocking all—though none were more overwhelmed with emotion than her daughter and husband. Though Vilya never earned the title of Voice of the Tempest, she now aids her daughter in leading her people without the faintest trace of jealousy.

\subsubsection{Keyleth, Voice of the Tempest}
\textit{Half-elf leader of \textit{Zephrah}}

\noindent
Hero of Tal'Dorei and member of Vox Machina, Keyleth completed her Aramente and took leadership among her people following the destruction of the \textit{Chroma Conclave}. She seems to have not aged a day since Vox Machina's victory over The Whispered One. Since her beloved's death, she has kept busy to fill the void in her heart. From expanding the crisis orb network across the realm and building beautiful windmills in the \textit{Summit Peaks} with Percy's aid to training Air Ashari heroes at home, she has never allowed herself a free moment. She has rarely left \textit{Zephrah} in the months since her mother's emotional return.

\section{Brawler's League}
The brutal blood sports of the Brawler's League began, unsurprisingly, during the vicious \textit{Drassig} Dynasty. In that cruel era, many of \textit{Emon's} affluent and indolent social elite used their wealth to engage in entertainment of a more "visceral" variety. The kingdom's poor were already fighting each other for scraps; why not pay them to literally fight to the death for their amusement?

Some strong, charismatic ruffians used this opportunity to gain social power for themselves, and chief among them was a rough-and-tumble entrepreneur from \textit{Kraghammer} named Bruel Sunderchin. Reaching out to his network of lowlifes and cutthroats, Sunderchin became ringleader of an annual gladiatorial tournament that caters expressly to Exandria's wealthiest elite.

Those who have become invested members of the league's events are anonymously delivered a golden coin bearing their symbol in the weeks leading up to an event. These coins are required for entry into the events, and those who wish to join, whether as a patron or participant, require a coin of their own.

\subsection{Goals}
The \textit{Brawler's League} is always in search of hidden and well-guarded locations around Tal'Dorei (and sometimes other realms) to establish their semi-annual event. A combatant or team is required to have a known and cleared sponsor or manager to enter the event. There are two tiers of bouts; one tier has teams of six battling via single elimination until a victor is declared. The other features single combatants fighting one-on-one, single-elimination, until a victor is declared. Killing your opponents isn't encouraged, but isn't explicitly discouraged, either. Bets are placed in person at the event, or remotely (via arcane messages and \textit{sending stones}) from all around Exandria.

\subsection{Relationships}
The \textit{Brawler's League} is extremely illegal in all major cities across Tal'Dorei. Should law-keepers of any nation discover the location of the semiannual event, the venue would find itself swiftly strangled by the long arm of the law. As such, magical protections and sentinels are placed throughout the event to prevent any unapproved entry, or to wipe the memory of those that do infiltrate.

A common rumor states that numerous political figures who publicly decry the event secretly partake in the gambling, and that the location of the event is concealed not only by magic, but also impenetrable layers of bureaucracy. \textit{Kraghammer} has hosted the event a number of times, with most local guards and great houses turning a blind eye.

The \textit{Brawler's League} frequently hires local \textit{Clasp} agents (or \textit{Myriad} operatives, in a pinch) to spread word of their next event through the criminal underground. If all goes well, they get hired at the event itself to provide protection, run interference, and make money on the side by providing loans and taking bets.

\subsection{Figures of Interest}
These figures are major members of the \textit{Brawler's League} that you can use as NPCs in your game. You should also take inspiration from these figures to make up new NPCs for your own campaign.

\subsubsection{Oben "Broadstaff" Sunderchin}
\textit{Dwarf organizer and head of \textit{Brawler's League}}

\noindent
Descendant of Bruel and inheritor of the league's top position, Oben secretly goes by the name Broadstaff and doesn't directly appear at events, as he avoids public connection with the league.

\subsubsection{Kradin Grimthorne}
Dwarf fight manager

A previously disgraced dwarf, exiled from \textit{Kraghammer} for murder, Grimthorne is the foremost brawler talent scout in Tal'Dorei.

\section{Chamber of Whitestone}
Generations ago, a team of intrepid explorers and fortune-seekers sailed from Port Damali in \textit{Wildemount} to northeastern Tal'Dorei through the \textit{Shearing Channel}, but wrecked upon the \textit{Alabaster Sierras'} rocky cliffs. The expedition was led by the de Rolo family, and these brave humans led the survivors to safety after the wreck. They took refuge in a valley within the \textit{Alabaster Sierras}, and there they discovered the holy \textit{Sun Tree}, a radiant tree blessed by the The Dawnfather. A community grew around the \textit{Sun Tree}, and the respected and beloved de Rolos were asked by their shipwrecked allies to be the leaders of this new settlement. Pressed against the chalky cliffs of the \textit{Alabaster Sierras}, their new homeland was dubbed \textit{Whitestone}.

The sovereign city-state of \textit{Whitestone} enjoyed peace for many generations, secluded from the rest of Tal'Dorei and \textit{Wildemount}. This serenity was destroyed when Lord Sylas and Lady Delilah Briarwood arrived in \textit{Whitestone} and murdered most of the de Rolo family. Lady Delilah, a necromancer of The Whispered One and a disgraced member of the Cerberus Assembly in \textit{Wildemount}, and her husband, a power-hungry \textbf{vampire}, transformed the quiet city-state into a realm of death. They attempted to take control of all of Tal'Dorei by puppeteering the Last Sovereign, Uriel Tal'Dorei II, but were thwarted by Vox Machina.

Thirty years after \textit{Whitestone's} fall, the \textit{Briarwoods} and their sepulchral master are defeated, and \textit{Whitestone} has returned to the blessed peace of its past. Nevertheless, much has changed. The surviving de Rolo siblings, Percival and Cassandra, reclaimed the city with the aid of Vox Machina. With the tyrants dethroned, Percival established a council titled the Chamber of \textit{Whitestone} to rule in his frequent absence, and the city now sends a delegation to the \textit{Tal'Dorei Council} to help safeguard the republic against evil.

\subsection{Goals}
Under the watch of Lady Cassandra de Rolo, the chamber handles and surveys all aspects of political life in \textit{Whitestone}. Its members are appointed by the head of the chamber, who is likely to be a scion of the de Rolo family for generations to come. The head can also add or remove members from the council. This system is new and untested by a ruler less noble than Percival or Cassandra. In many ways, this system of government is another one of Percy's tinkering projects.

Members of the Chamber are granted distinct areas of responsibility to the realm.

\textbf{The Guardian of Woven Stone} is the title given to the central figure on the council, responsible for higher civil justice, community unification, and oversight of internal government.

\textbf{The Steward of Sunblessed Gifts} is responsible for local farming and goods, overseeing distribution between citizens and exportation.

\textbf{The Curator of Fortune's Bounty} is responsible for all major commerce and economic business, both locally and with allied and foreign entities.

\textbf{The Keeper of Divine Virtue} is responsible for religious organization and spiritual guidance regarding multiple faiths within the region, both as a religious leader and a protector of practices.

\textbf{The Commissioner of Wardship} is responsible for the enforcement of local law, the organization of the Paleguard, acting as judge on smaller matters of civil justice, and defense of the city during times of war.

\textbf{The Architect of Enlightened Progress} is responsible for the construction and infrastructure of the city and castle.

\textbf{The Grand Mistress of the Grey Hunt} is responsible for diplomatic relations and for rooting out dangers beyond the city's walls.

\subsection{Relationships}
Despite its benefits, \textit{Whitestone's} previous solitary existence left it without allies in its time of need. The Chamber of \textit{Whitestone} has sent a delegation to the \textit{Tal'Dorei Council} for the betterment of both. These days, trade caravans are a common sight along the \textit{Silvercut Roadway}, transporting goods from \textit{Whitestone} to \textit{Westruun}, \textit{Kymal}, and ultimately \textit{Emon}. The most expensive goods travel via skyship, which land at \textit{Whitestone's} brand-new skyport.

\subsection{Figures of Interest}
These figures are major members of the \textit{Chamber of Whitestone} that you can use as NPCs in your game. You should also take inspiration from these figures to make up new NPCs for your own campaign.

\subsubsection{Lady Cassandra Johanna de Rolo, Guardian of Woven Stone}
\textit{Human head of the \textit{Chamber of Whitestone}}

\noindent
Cassandra is Percival's younger sister, thought dead until he discovered her in the thrall of the \textit{Briarwoods}. She now manages \textit{Whitestone's} most pressing matters, with daily briefings from other members of the chamber. She is beloved throughout the city for her resilience and her decisive leadership.

\subsubsection{Lord Percival Fredrickstein von Musel Klossowski de Rolo III, Architect of Enlightened Progress}
\textit{Human member of the \textit{Chamber of Whitestone}}

\noindent
Percival is Cassandra's older brother, thought dead until Vox Machina discovered him in a prison cell. His life was filled with countless joys and regrets. He has taken well to retirement from adventuring, and is renowned across Tal'Dorei for his creations.

\subsubsection{Lady Vex'ahlia de Rolo, Baroness of the Third House of Whitestone and Grand Mistress of the Grey Hunt}
\textit{Half-Elf Member of the \textit{Chamber of Whitestone}}

\noindent
Vex'ahlia is the Grand Mistress of the Grey Hunt, a title she has held since before she married Percival. Over the past twenty-four years, she has become a talented and respected leader of the \textit{Grey Hunters}. And though the hunters use guns inspired by her husband's own designs, she still prefers to wield her trusty bow.

\subsubsection{Grog Strongjaw, Grand Poobah de Doink, In Charge of All This-and-That}
\textit{Goliath friend}

\noindent
The noble, if simple, warrior of Vox Machina, no one really knows if Grog's title is a real title at all, or if Percy was just joking when he gave it to him. Grog likes it, and that's good enough for most people.

\section{Claret Orders}
\begin{DndReadAloud}{}
\DndQuote%
{}
{\enquote{As blood flows through and invigorates the body, so faith flows through and invigorates the soul. Ever remember this: your blood and your faith are one and the same. Spill not your blood nor that of another heedlessly, but do not hesitate to spill either when the cause is just.}}
{From The Crimson Canon}
\end{DndReadAloud}

Hundreds of years ago, a fell power exerted its dark will over \textit{Wildemount}, corrupting its outlying townships and spreading chaos among its people. Undeath washed over gravesites like a plague, shadowed beasts stalked the midnight woods, and the influence of fiends befouled even \textit{Wildemount's} purest souls.

In the eleventh hour, a priest of the The Matron of Ravens named Trence Orman prayed for a way to protect his flock. The The Matron of Ravens inspiration came in the form of long-hidden knowledge: the secrets to blood magic. Taking his gifts, Trence trained his most trusted warriors in these techniques, giving a portion of their humanity in exchange for the power to defend their people. This marked the origin of the \textit{Claret Orders}, and the first \textit{Blood Clerics}, \textit{Blood Mages}, and Blood Hunters.

The \textit{Claret Orders} work in relative secrecy, as the nature of their abilities is largely misunderstood, which often leads to persecution. Their numbers are small, as the price is great and the \textit{Orders} only accept those who are strong of heart and have little to lose. Small sects of the \textit{Orders} have begun work far enough West to wander the lands of Tal'Dorei in recent years, keeping out of the public eye and remaining vigilant to the signs of familiar darkness undermining the region.

Despite the \textit{Claret Orders'} best attempts at secrecy, the esoteric art of blood magic has escaped their clutches. Though rare, rogue blood mages and blood hunters have spread across both \textit{Wildemount} and Tal'Dorei, unbound by the tenets and strictures of the \textit{Claret Orders}.

\subsection{Goals}
The Creed of the \textit{Claret Orders} asks that its followers commit their lives to the hunt of entities and creatures that threaten the sanctity of life and joy. While the sacrifice is great, the reward is the continued existence of purity and good in the world. The \textit{Orders} strive to protect those who cannot protect themselves from the shadows, taking no credit for their services beyond the means to live and travel. Some have strayed, focusing on building a fortune for themselves, but should word of such intent make it back to the heads of the \textit{Orders}, and their actions begin to threaten the innocents they've sworn to protect, they may find themselves among the hunted.

\subsection{Relationships}
The \textit{Claret Orders} have no alliance with—nor allegiance to—any standing government or group. Solitary in action and discreet in display, they've not drawn the attention of many factions in Tal'Dorei as of yet. However, it is only a matter of time before curious eyes begin to seek answers to questions about these dark guardians and their true nature.

\subsection{Figures of Interest}
These figures are major members of the \textit{Claret Orders} that you can use as NPCs in your game. You should also take inspiration from these figures to make up new NPCs for your own campaign.

\subsubsection{Claret Director Alasterre de Vitrevos}
\textit{Human Blood Domain cleric of Ravens}

\noindent
Calm and sanguine to a fault, \textit{Alasterre de Vitrevos} is the youngest-ever director of Tal'Dorei's \textit{Claret Order}. He sees his position as one he deserves without question, and is willing to use his command more boldly than his mentor, director emeritus Teresa Dulamar. He is frequently seen in the field, working with various adventuring parties and his own blood hunters to wipe out monsters wherever they emerge. The combination of his relaxed personality and his aggressive strategy may seem a contradiction at first glance, but those who know him see the truth of the immense self-assurance that lets him make hard decisions without hesitation. This quality is a strength when cool-headed leadership is needed in the thick of battle, but it can make him challenging to work with as an equal.

\subsubsection{Director Emeritus Teresa Dulamar}
\textit{Dwarf Blood Domain cleric of Ravens}

\noindent
The former head of Tal'Dorei's \textit{Claret Orders} retired ten years ago after suffering a traumatic death in combat. Her soul accepted resurrection, but she took her defeat as a sign to let younger, fresher blood take her place. She remains a part of the \textit{Claret Orders} to advise the current director on matters he is too young to have experienced himself. Though she displays a chilly demeanor, she simply keeps her compassionate side hidden from the untrusted riffraff.

\subsubsection{Hunter Jorick LaMensh}
\textit{Human blood hunter}

\noindent
Once among the greatest of all blood hunters, Jorick is now Director Emeritus Teresa Dulamar's personal bodyguard. He is still a capable member of the Order of the Profane Soul, and urges immediate, decisive action to destroy demons and devils whenever their occult activity is discovered.

\section{The Clasp}
The \textit{Clasp} is a well-established secret organization consisting of spies, thieves, fences, assassins, saboteurs, and smugglers. It was built upon the ruins of the Winksman Thieves' Guild, an organization that flourished during the reign of \textit{Drassig}, but crumbled under the rule of \textit{Zan Tal'Dorei}. Based largely in \textit{Emon}—and under it, via a network of subterranean tunnels—the \textit{Clasp} also runs smaller outfits in \textit{Westruun}, \textit{Stilben}, and \textit{Kymal}. Members are branded with the guild's symbol on their mid-back, signifying allegiance and commitment. The \textit{Clasp} operates like a business-minded family or mafia, celebrating loyalty and an ability to keep one's word. A member pays their debts, keeps their promises, and always seeks opportunity for the betterment of the \textit{Clasp}.

Leaders of a \textit{Clasp} sect bear the title of Spireling, and they oversee assigned elements of the organization's functions. Spirelings hold equal power within a sect, and there is always an odd number acting within that location to break any disagreements over business ventures and actions.

\subsection{Goals}
While most common folk assume that the \textit{The Clasp} is merely interested in profit at any cost (and that is a merit of the \textit{The Clasp}), they also wish to preserve the function and prosperity of civilization. It is commonly held that a thriving city is a profitable one, and an enemy of civilization is a foe of the \textit{The Clasp}. They largely consider themselves a necessary dark underbelly, willing to do the under-the-table deeds that more wealthy and powerful citizens cannot commit with their own hands. On the other hand, less scrupulous Spirelings also understand that discontent breeds profit, and have occasionally kindled the flames of unrest to keep things tense and profitable.

Were it not for the branch of the \textit{The Clasp} in \textit{Emon}, the \textit{Chroma Conclave's} attack would likely have destroyed the city's culture entirely. However, thanks to the \textit{Clasp's} role as a "shadow civilization," the people of \textit{Emon} were able to maintain their foothold at home, even under the watchful eye of \textit{Thordak, the Cinder King}.

\subsection{Relationships}
A generation ago, the \textit{Clasp} was considered by most to be an uncouth cesspool of cutthroats and criminals. These days, thousands of people know the \textit{Clasp} as the roguish heroes who saved their lives, or the lives of their parents, from the \textit{Chroma Conclave}. Nowhere is this truer than in \textit{Emon}, where its Spirelings make public appearances and even, on occasion, deal openly with the \textit{Tal'Dorei Council}. Nevertheless, \textit{Clasp} agents outside of \textit{Emon} keep their wits about them, for most local governments still view them as underhanded scoundrels—and not without reason. Most guards and town militias will arrest known \textit{Clasp} members on sight; however, most \textit{Clasp} members know well enough to have a friend on the force that will help them escape.

The \textit{Clasp} frowns upon dogmatic religion or cultish behavior, as it has conflicted with the organization's business in the past, and always ended messily.

Since the arrival of the \textit{Myriad} within Tal'Dorei from \textit{Wildemount}, the \textit{Clasp} has been on extremely high alert, seeking information on them and any means of banishing this new competition. Gang warfare has broken out sporadically between the \textit{Clasp} and the \textit{Myriad} in cities from \textit{Stilben} to \textit{Emon} over the past twenty years, and though the rival organizations are currently at peace, everyone knows the whole thing is just one poached client away from reigniting.

\subsection{Figures of Interest}
These figures are major members of the \textit{Clasp} that you can use as NPCs in your game. You should also take inspiration from these figures to make up new NPCs for your own campaign.

\subsubsection{Spireling Shenn}
\textit{Human Spireling of Secrets (\textit{Emon})}

\noindent
A powerful illusionist, Shenn runs a far-reaching network of spies and informants. In his younger days, he was a canny negotiator who sought to spread the \textit{Clasp} beyond Tal'Dorei. The arrival of the \textit{Myriad} forced the \textit{Clasp} onto the defensive, and he will never forgive them for souring his ambitions.

\subsubsection{Spireling Warren}
\textit{Half-orc Spireling of Blades (\textit{Emon})}

\noindent
Deadly duelist turned master assassin of the \textit{Clasp}, Warren handles all contracts that involve kidnapping, assassination, and intimidation. He hates public appearances.

\subsubsection{Spireling Zilloa}
\textit{Half-elf Spireling of Shadows (\textit{Emon})}

\noindent
An accomplished burglar, Zilloa has risen to the head of all thieving, smuggling, and black-market dealings within the \textit{Clasp} of \textit{Emon}. She plays a dangerous game by enjoying folk-hero status as the public liaison between the \textit{Clasp} and the \textit{Tal'Dorei Council}.

\subsubsection{Spireling Gholesh}
\textit{Human Spireling of Shadows (\textit{Westruun})}

\noindent
A crass and ambitious man, Gholesh's appearance is a mystery as his face is always hidden behind a leather mask that resembles an ogre. He has secretly thrown in his lot with the \textit{League of Miracles} and, with their backing, he aspires to take over the \textit{Clasp} and enjoy all its wealth and prestige.

\subsubsection{Spireling Melchior}
\textit{Tiefling Spireling of Blades (\textit{Westruun})}

\noindent
This clean-faced, curly-haired tiefling with pale blue skin and stubby horns may have a soft, babyish visage, but it belies their powerful physique. They are equally skilled with a sword, axe, or bow, and made their name in underground brawling rings before joining the \textit{Clasp}.

\subsubsection{Spireling Fetch}
\textit{Halfling Spireling of Secrets (\textit{Westruun})}

\noindent
Creepy and unkempt in appearance, Fetch's gaunt face only distracts from his cunning and skills of keen manipulation.

\section{Golden Grin}
In dark times, entertainment and good cheer are just as vital as resistance. During the despotic rule of King Warren \textit{Drassig} and his sons, a group of bards, dancers, and storytellers united to subvert his rule and inspire the people to rise against him. This band took the name "the \textit{Golden Grin}" and traveled across the realm to offer escapism and entertainment, all while planting the seeds of discontent, rebellion, and heroism.

Since the fall of \textit{Drassig}, the \textit{Golden Grin} has been ever-present, using song, comedy, and theatrics to inspire the people of Tal'Dorei to endure hardship and fight back against evil. These "Grinners" include among their ranks spiritual leaders, artists, musicians, innkeepers, and dozens of retired adventurers who still want to do good outside the public eye. Little is publicly known about the \textit{Golden Grin}, other than myths and rumors, and the Grinners gladly perpetuate these tall tales.

\subsection{Goals}
The \textit{Golden Grin} (or simply "the Grin") exists to enlighten and inspire, keeping ears and eyes out for the murmurs of tyranny and cruelty. The Grin's core philosophy is that each and every person is capable of great things alone, and even greater things together. Rather than taking direct action themselves, most Grinners use songs and tales of great heroes to galvanize people to action.

\subsection{Relationships}
The \textit{Golden Grin} claims no formal allies, nor do members reveal their secret affiliation to outsiders unless circumstances are dire indeed. They use their covert, grassroots power to empower local freedom fighters and activists. Some Grinners do occasionally find their way into government, granting them a higher vantage point and platform to disseminate the ideals behind their cause—but even these members of the Grin take a dim view of the slowly turning gears of bureaucracy.

Enemies of the \textit{Golden Grin} are those that oppress and abuse freedom of expression and living. The Grin likewise mistrusts people and institutions that control or regulate the populace too heavily, even if they seem to have the best of intentions.

\subsection{Figures of Interest}
These figures are major members of the \textit{Golden Grin} that you can use as NPCs in your game. You should also take inspiration from these figures to make up new NPCs for your own campaign.

\subsubsection{Theona Balmhand}
\textit{Gnome cleric of the Lawbearer}

\noindent
A wandering healer hailing from \textit{Kymal}, Theona lives by donation and perpetually travels from place to place, seeking to aid who she can and to learn whatever the people whisper.

\subsubsection{Doctor Dranzel}
\textit{Half-orc traveling bard}

\noindent
A jovial and lackadaisical violinist and frontman for Doctor Dranzel's Traveling Troupe, the renowned Doctor Dranzel exaggerates his stage persona for a number of reasons, not the least of which is to make him seem like a preening fool—the perfect cover for a master spy of the Grin.

\section{Houses of Kraghammer}
The fortified, subterranean city-state of \textit{Kraghammer} was carved out of the bedrock of the \textit{Cliffkeep Mountains} in the years following the \textit{Divergence}. It is the ancestral home of the native dwarves of Tal'Dorei, and the five major bloodlines that founded the city in ages past still rule to this day. The ruling families are Houses Greyspine, Zuurthom, Bronzegrip, Thunderbrand, and Glorendar. These Great Houses elect an official to the office of Ironkeeper every ten years to keep the alliances between houses just and healthy.

Though the culture of \textit{Kraghammer} is highly traditional and as stoic as the stone it was carved from, it is not without unrest. It seems every dwarven generation has its own great social movement where disenfranchised workers or greedy, lesser noble houses attempt to unseat the ruling five. Yet for good or for ill, the Great Houses have time and time again shown that they can put aside their differences to preserve their mutual rulership.

\subsection{Goals}
Each ruling house has a specialty and holds responsibility over that domain within \textit{Kraghammer}. House Greyspine maintains the largest mine under \textit{Kraghammer}, as well as the Pools of Solace. House Zuurthom trains the chief architects of the city, and most masonry and building is organized through them. House Bronzegrip funds the smithing guilds and maintains the Bronzegrip Metalworks, \textit{Kraghammer's} largest blast furnace and ore refinery. House Thunderbrand prides itself on training \textit{Kraghammer's} foremost scholars and premiere arcanists. House Glorendar maintains the law and delivers punishment.

\subsection{Relationships}
Almost every dwarf living in \textit{Kraghammer} can remember a time when, "Kraghammer for the dwarves!" was a common cheer at the pubs. And when they said, "for the dwarves," they meant just the dwarves born in \textit{Kraghammer}. Times have changed. The prevailing opinion is that a unified Tal'Dorei is best for all. Those dwarves whose traditionalism is just a veil for bigotry have found that they aren't welcome in \textit{Kraghammer}—they've subsequently retreated to make their own isolated enclaves in the \textit{Cliffkeeps}.

The \textit{Houses of Kraghammer} facilitate a healthy trade of metals, weapons, and armor with \textit{Emon} and \textit{Westruun}. These bonds of commerce and communication have only grown stronger since an ambassador from \textit{Kraghammer} has taken a seat on the \textit{Tal'Dorei Council}.

Many members of House Thunderbrand also have ties to the \textit{Arcana Pansophical}. However, lesser clans with aspirations of supplanting the renowned Thunderbrands have thrown their lot in with the \textit{League of Miracles}, hoping to become the new masters of the arcane. They have been wooed by to this upstart faction with promises of disrupting the old ways of commerce and prestige, establishing something new.

\begin{DndReadAloud}{}
\enquote{%
\small\itshape
Five the fingers that grasp hammer or blade,\\
Five Houses of Kraghammer rightly made:\\
Greyspine delves 'neath dwarven halls,\\
Zuurthom raises stony walls,\\
Bronzegrip masters of metal wrought,\\
Thunderbrand uncovers secrets sought.\\
And Glorenthal to judgements hear—\\
In Kraghammer, dwarf-folk need never fear.
}
\end{DndReadAloud}

\subsection{Figures of Interest}
These figures are major members of the \textit{Houses of Kraghammer} that you can use as NPCs in your game. You should also take inspiration from these figures to make up new NPCs for your own campaign.

\subsubsection{Ironkeeper Palluda Zuurthom}
\textit{Dwarf Ironkeeper of House Zuurthom}

\noindent
Even after weathering the \textit{Chroma Conclave}, uniting \textit{Kraghammer's} squabbling houses in support of the \textit{Tal'Dorei Council} was no mean feat. Ironkeeper Palluda of House Zuurthom proved up to the task. They made a name for themself by taking control of a trade deal that ultimately sold tens of thousands of gold pieces' worth of quarried stone to the \textit{League of Miracles}, which began the reconstruction of Tal'Dorei in earnest. Palluda became a hero of the working dwarf within \textit{Kraghammer}. They are a master at getting all sides to agree to a compromise, and do so with an iron-edged smile.

\subsubsection{Lord Nostoc Greyspine}
\textit{Dwarf overseer of House Greyspine}

\noindent
The dour executive of dozens of quarries across the \textit{Cliffkeep Mountains}, Nostoc makes up for his unpleasant personality by being a formidable businessman. In recent years, he has become fabulously wealthy as the reconstruction of cities across Tal'Dorei devoured thousands of tons of stone. House Thunderbrand and the \textit{League of Miracles} made it possible to transport it all with teleportation spells.

\subsubsection{Elementalist Drake Thunderbrand}
One of the eldest of his House, Drake is known for his surprising empathy and fierce temper. He is also a member of the \textit{Arcana Pansophical}. His friendship with \textit{Allura Vysoren} made him a popular choice for \textit{Kraghammer's} ambassador to the council, but he politely refused, claiming it was a younger dwarf's duty.

\section{League of Miracles}
"Magic isn't cheap, but miracles are priceless." Such is the credo of the \textit{Wonderworker}, the mysterious thaumaturge whose words inspire and guide the \textit{League of Miracles}. This company of mercenary magi is the reason why Tal'Dorei has recovered so rapidly from the near-apocalyptic destruction of the \textit{Chroma Conclave}. \textit{Westruun} and \textit{Emon} rose from rubble to the height of their former glory in a matter of months.

The people of Tal'Dorei adore the league. Their mighty spellwrought adranachs can be seen lifting thousand-pound bricks of limestone high into the sky to rebuild walls, towers, and even castles in a matter of days. Many of the league's spellwrights have become folk heroes and celebrities for their miraculous acts, which range from healing those maimed in the \textit{Conclave's} attack to transmuting raw timber, mud, and straw into full-fledged houses in minutes.

For most of Tal'Dorei's elite, the \textit{League of Miracles} is likewise a great ally. Trade and travel have resumed stronger than before, for the league has eagerly subcontracted their talented mages out as guards for caravans and persons of interest.

Yet, not everyone sees the league as an ally. In the early aftermath of \textit{Thordak's} defeat, the new and untested \textit{Tal'Dorei Council} put out thousands of bounties on reconstruction projects to get their towns and cities back on their feet, and local mayors and margraves across the nation agreed to these subsidies to reduce the cost of reconstruction for their already devastated populace. They expected Tal'Dorei's rebirth to be a multi-generational undertaking—yet the \textit{League of Miracles} snapped up these bounties like candy, and the league does not work their miracles cheaply. Bills for hundreds of thousands of gold pieces' worth of bounties found their way to the desks of local leaders and to the \textit{Tal'Dorei Council} by the end of the year. Such vast wealth for thirty years' worth of rebuilding in three years' time simply did not exist.

But this was all part of the \textit{Wonderworker's} plan, and their spellwrights executed it flawlessly. The league needs money to purchase their expensive spell components and to pay their talented mercenary membership, of course, but their true goal is the priceless currency of power, in the form of political favors from leaders ranging from small-town mayors, to city-ruling margraves, and even delegates to the \textit{Tal'Dorei Council}. Rumors even swirl that one or more of the members of the council itself are in the league's pocket.

The league's members are delineated into a simple hierarchy.

\subparagraph{Mercenary Mage}
Spellcasters drawn from across Tal'Dorei form the backbone of the \textit{League of Miracles}. The league works as an agency, handing mercenaries lucrative jobs, collecting the payment themselves, and paying out in a timely fashion—minus twenty percent. These mercenaries come from all walks of life. Some are aspiring adventurers seeking a less life-threatening first gig, while others are skilled arcanists with shady pasts, including Remnant cultists and double-dipping \textit{Myriad} agents with a talent for magic.

\subparagraph{Spellwright}
The most talented (or best-connected) members of the league are awarded the title of spellwright. Only two dozen such masters of the arcane exist across Tal'Dorei. These spellwrights know the secrets of how to create adranachs, mighty constructs of magical energy that were used publicly to rebuild Tal'Dorei's ravaged infrastructure—and are also used in secret as hunter-destroyers of the league's enemies. These spellwrights sign a blood pact that connects their minds to the mysterious communications of the \textit{Wonderworker}, which all spellwrights are blood-bound to interpret and enact.

\subparagraph{Grand Thaumaturge}
The official title of the league's shadowy master is the Grand Thaumaturge, though most colloquially refer to them as the \textit{Wonderworker}. None know their true identity, or their ultimate goal. The spellwrights understand that the Grand Thaumaturge has a master plan that must be followed. Only these spellwrights have the privilege of hearing their master's psychic messages. Over the years, they have interpreted them to mean that the \textit{League of Miracles} must become the secret puppet-masters of all Tal'Dorei.

\subsection{Goals}
The \textit{League of Miracles} is comprised mostly of mercenary magi pulled from all corners of Tal'Dorei, with some even coming from abroad as news of the league's reputation spreads across Exandria. For most of these arcanists, their goal is simple: to make money. The goals of the league's spellwrights are more elaborate. They listen for the cryptic words of the mysterious Grand Thaumaturge and dedicate themselves to deciphering and enacting their grand design. The spellwrights' current aim is one of total control—once all of Tal'Dorei's social and political elite are in their pockets, they can shape the land as they see fit, all to the cheering adoration of the fawning people of Tal'Dorei.

\subsection{Relationships}
The \textit{League of Miracles} tries to keep cordial, if distant, relationships with every major power on Tal'Dorei. Their public image is shiny and unimpeachable. Many members of the \textit{Tal'Dorei Council} are already in their pocket and cast votes that further the league's self-interest. The \textit{Arcana Pansophical}—especially \textit{Allura Vysoren}, who also sits on the council—are deeply distrustful of this unregulated, uncommonly powerful organization.

The \textit{League of Miracles} has a tense relationship with the \textit{Chamber of Whitestone}, as they have long sought access to the \textit{residuum} unique to that territory. The spellwrights fear outright hostility with the legendary heroes of Vox Machina, and have supposedly abandoned their quest to mine \textit{Whitestone's} naturally occurring \textit{residuum}.

As "no questions asked" employers of mages, the league is in the good graces of many less scrupulous arcane factions, including the \textit{Remnants}. Necromancers, blood mages, Ruidus cultists, and other unsavory magi now exist under the league's protection.

\subsection{Figures of Interest}
These figures are major members of the \textit{League of Miracles} that you can use as NPCs in your game. You should also take inspiration from these figures to make up new NPCs for your own campaign.

\subsubsection{The Wonderworker}
\textit{Mysterious enactor of true miracles}

\noindent
All spellwrights hear the authoritative—yet enigmatically distorted—voice of the \textit{Wonderworker}. The \textit{Wonderworker} has supposedly appeared in the flesh only three times throughout the league's twenty-four-year history, as a humanoid draped in robes of starlight and wearing a mithral mask. In those moments, they have produced true miracles—raising a castle from the ground in a matter of seconds, halting the rampaging silver dragon Saldarthoryn in her tracks and healing her twisted mind, and stopping the Center Slab of \textit{Kraghammer} from collapsing into the Pits in a catastrophic mining accident that would have cost countless lives. These legendary acts inspired utmost awe and faith in all who witnessed them, leading the spellwrights to create adranachs in the \textit{Wonderworker's} image and enact their cryptic will.

\subsubsection{Blaine Kraverrogg}
\textit{Half-halfling mercenary necromancer}

\noindent
Short and brimming with a malign lust for power, Blaine appears to be a short, skeletally skinny human who hides his lack of stature in billowing black-and-gold robes. He is a member of the \textit{Remnants}, but pays his bills by doing undercover grave-robbing for the league. He longs to reach the bottom of \textit{Shadebarrow}, a mysterious crypt somewhere on the \textit{Dividing Plains}. His obsession with \textit{Shadebarrow} stems from a hunger for the power of the Demon Prince of Undeath.

\subsubsection{Honor Kinnabari}
\textit{Tiefling blood mage and novice spellwright}

\noindent
As the newest spellwright in the \textit{League of Miracles}, \textit{Honor Kinnabari} has a lot to prove. She was granted a chance to prove her quality in the eyes of the \textit{Wonderworker} when the league was contracted to rebuild the Market Ward of \textit{Westruun} after it was devastated by a herd of magma landsharks that descended from the \textit{Cliffkeep Mountains}. Her red eyes and sheet-white skin and horns complement a traditional red-and-white mage's robe from the reclusive dwarven clans of the \textit{Alabaster Sierras}. She is fully dedicated to the league, but she is as honorable as her name implies; her fellow spellwrights fear that she could turn traitor, should the unsavory truth of the league be laid plain before her.

\section{Library of the Cobalt Soul}
The Cobalt Reserve in \textit{Westruun} is home to a monastic order so old, enduring, and tirelessly dedicated to the preservation of knowledge that most people in Tal'Dorei simply take it for granted. Just as birds have always flown and grass has always grown, the Library of the Cobalt Soul has always sent its agents across the world, taking minutes at public trials, archiving broadsheets, exploring ancient ruins, and copying scrolls. To all but the most dedicated scholars, the work of the Cobalt Soul seems like dreary, banal stuff.

Nothing could be further from the truth. Beneath the veneer of being a dopey order of bookish monks, the Cobalt Soul is one of the most dogged espionage organizations in the world. The secret agents of the Cobalt Soul, its expositors, are spies who wield the truth as a weapon more potent than any blade. The Cobalt Soul was born in \textit{Wildemount}, just after the end of the \textit{Calamity}, when the The Knowing Mentor was maimed by an evil god. Her most loyal followers went into hiding and vowed to continue her divine work in knowing all there is to be known—and most importantly, to use the truth as a holy weapon against lies and propaganda.

Each archive of the Cobalt Soul has a common hierarchy, and higher-ranking members of the Cobalt Soul can command their lower-ranking comrades from another archive.

\subparagraph{Monk}
The monks of the Cobalt Soul are their most visible members, for they travel the land in search of knowledge both modern and archaic. Not all monks are engaged in fieldwork, especially new recruits. Those who work entirely within an archive aid in the organization and preservation of information. Few know that these monks are dedicated to the The Knowing Mentor; the common opinion in Tal'Dorei is that they are ascetics, obsessively dedicated to the chronicling of knowledge, not devotees of an obscure god.

\subparagraph{Expositor}
Trained in the art of espionage, the expositors of the Cobalt Soul use their infiltration expertise to uncover the secrets of governments, world leaders, and obscure cults. In times of great crisis, undercover expositors have even incited revolutions by sharing the knowledge they discovered with people toiling under evil regimes. This true mission of the expositors is a secret even to low-ranking monks, who see them simply as well-trained field agents.

\subparagraph{Archivist}
Though they are not involved in field work, archivists play the vital role of organizing the overwhelming flow of information that the monks and expositors return to their archive. They also train new monks and perform administrative tasks, like negotiating with sellers of artifacts and with government officials to ensure the stability of their archive.

\subparagraph{Curator}
A small group of three to five curators manage the top-level operations of each archive. One of their number is given the special designation of High Curator, to whom all other members of their archive must ultimately answer.

\subsection{Goals}
The public mission statement of the Cobalt Soul is to archive knowledge for the betterment of all peoples of Exandria. The secret mission of the Cobalt Soul is to use their accumulated knowledge to root out tyranny and ensure the freedom of all people.

As instructed by their reclusive god, the The Knowing Mentor, the Cobalt Soul in Tal'Dorei makes their archives available to everyone, so that knowledge can be used to improve and enrich the lives of all. Despite this, the curators of the Cobalt Reserve privately determine what knowledge is too volatile to be freely disseminated. Typically, this privileged information is magical in nature, but it also includes secrets recovered by expositors that must be strategically employed as a weapon against unjust leaders, cults, and power-hungry factions.

The Cobalt Soul itself is not immune to corruption and complacency. Nowhere is this truer than in the Cobalt Soul's birthplace of \textit{Wildemount}, where it must work hand-in-hand with the authoritarian Dwendalian Empire. Its highest officials say that the Cobalt Soul has outlasted every tyrant on Exandria, and that they must be content with small acts of insurrection—subtly subverting propaganda and sharing the truth with any who will listen—while waiting for the right moment to strike a fatal blow. Nevertheless, this mantra of patience has caused some of its more radical members to abandon the order entirely. Typically, deserters of the Cobalt Soul are expositors, who see the ugliest side of the world firsthand again and again.

\subsection{Relationships}
The Cobalt Soul is an international organization with archives in the Dwendalian Empire and Clovis Concord of \textit{Wildemount}, the Republic of Tal'Dorei, the Dawn City of Vasselheim on \textit{Issylra}, and the great city of \textit{Ank'Harel} on \textit{Marquet}. Each of these archives acts independently under the orders of their commanding High Curator, though they must ultimately answer to the joint leadership of the Cobalt Soul: the High Curators of the Rexxentrum and Vasselheim archives.

Though Tal'Dorei is a freer society than the Dwendalian Empire, the Cobalt Soul knows better than to think that means it is free of corruption. From the Cobalt Reserve, their archive in the city of \textit{Westruun}, the Cobalt Soul keeps close tabs on members of the \textit{Tal'Dorei Council}, the leaders of \textit{Syngorn} and \textit{Kraghammer}, and other powerful individuals across the land.

The Cobalt Soul actively opposes the \textit{Remnants}—followers of The Whispered One, a tyrannical god of secrets who is diametrically opposed to the goals of the The Knowing Mentor. They are also concerned by the meteoric rise of a guild of arcanists called the \textit{League of Miracles}, whose political maneuverings have made puppets of dozens of delegates to the \textit{Tal'Dorei Council} and political leaders across the land. This worry has only deepened as rumors reach them of the league filling their ranks with Remnant cultists who seek to dominate Tal'Dorei.

\subsection{Figures of Interest}
These figures are major members of the \textit{Library of the Cobalt Soul} that you can use as NPCs in your game. You should also take inspiration from these figures to make up new NPCs for your own campaign.

\subsubsection{High Curator Greisalda Cassios}
\textit{Half-\textit{dragonblood} leader of the Cobalt Reserve}

\noindent
The High Curator of the \textit{Westruun} archive of the \textit{Cobalt Soul} is a high-cheekboned, bald woman who would appear fully human if not for her slitted, silver eyes and the silver scales that run across her head, neck, and upper arms. She is regarded as one of the fastest minds in Tal'Dorei, for she has the rare gift of perfect recall. Monks tremble when her clear, unpretentious voice rings through an archive chamber, for Greisalda remembers acts of justice and personal slights with equal intensity, making her a powerful ally and a furious enemy. She is concerned that corruption is spreading throughout the \textit{Tal'Dorei Council}.

\subsubsection{Cressida Holt}
\textit{Orc expositor of the Cobalt Reserve}

\noindent
The name \textit{Cressida Holt} is synonymous with style and effortless charisma—and also a blunt, merciless coldness. \textit{Cressida} is nonbinary, and uses both she and they pronouns interchangeably. Though she's fresh blood to the expositors of the Cobalt Reserve, she has already taken advantage of her license to kill on her espionage missions. They are slim for an orc, but still more muscular than many of their human colleagues. Their assignments take them across Tal'Dorei in search of political corruption. She is currently stationed in \textit{Emon} to observe a person of interest on the \textit{Tal'Dorei Council}, where she alternates between lurking on rooftops and striding through ballrooms in fine suits.

\subsubsection{Archivist Salomane Rothtree}
\textit{Half-elf archivist of the Cobalt Reserve}

\noindent
A haggard, overworked administrator, Archivist Rothtree is known throughout the Cobalt Reserve as someone who tries too hard to prove himself and constantly bites off more than he can chew. Half-orc and half-elf, he has said that he hopes to accomplish as much in his lifetime as an elf could in theirs. He often is tasked with helping adventurers find information. Though he resents the task, he may slowly open up to heroes who treat him well.

\section{The Myriad}
\textit{Wildemount's} greatest organized crime ring is now becoming known throughout Tal'Dorei, but it came from humble beginnings. The Myriad began in the Dwendalian city of Yrrosa, as a small shipping company specializing in exotic goods and services, exploiting and selling curiosities from the distant continents of \textit{Othanzia} and \textit{Ank'Harel}. Over time, as legal trading became more difficult, the Myriad's illegal occupations became their best business. Through years establishing underworld contacts and tactical blackmail, the Myriad evolved into a powerful, well-veiled criminal faction. Over the past two decades, the Myriad's aspirations have spread beyond \textit{Wildemount}, and they are now deeply entrenched in Tal'Dorei's criminal underworld.

Known for displays of opulence and social grace, this misdirection enables the Myriad to seduce and control the lives and wills of whomever is deemed valuable to their plans. Dozens of truly amoral cutthroats and thieves noticed when the \textit{Clasp} began to work hand-in-hand with the \textit{Tal'Dorei Council}. They felt the dagger of accountability pressed against their backs, and broke with their former guild. Now, the Myriad has filled the vacuum of truly lawless criminality that the \textit{Clasp} left in its wake.

The Myriad operates as a loose network of gangs whose bosses run their own localized sects, without direct oversight from the mysterious heads of the society. Tithes and information are expected to be delivered to the leadership at a steady rate from each satellite of the organization, and should the success slow, threats of enslavement or destruction soon follow.

\subsection{Goals}
The \textit{Myriad} seeks wealth in hideous excess. The means are unimportant, as long as they're lucrative. No trade is too abhorrent, from bloody theft to kidnapping and slavery. And though they lack the political capital to bribe and puppeteer politicians the way they did in \textit{Wildemount}, the \textit{Myriad} are confident that they can get even the supposedly noble \textit{Tal'Dorei Council} in their pocket.

Most fronts for the \textit{Myriad} are businesses that publicly deal in antiquities and exotic textiles. Behind these masks is a ruthlessly efficient black market of slaves, magical beasts, weapons, and illicit substances. Overt violence tends to draw unwanted attention, and the \textit{Myriad} avoid it when possible, preferring to discredit, blackmail, or frame those who cross them. Should that fail, however, kidnapping or assassination can be a quick and traceless final alternative.

\subsection{Relationships}
Though they are relative newcomers to the criminal underworld of Tal'Dorei, the \textit{Myriad} is well-versed in the covert art of making money by any means necessary. The \textit{Clasp} is well aware of their presence in the broadest sense, and the two thieves' guilds have spilled blood in the streets of every city in Tal'Dorei at some point in the past twenty years. These days, the bloodshed has subsided into a cold war of espionage and counter-espionage, as each organization tries to sabotage the other.

The \textit{Myriad} has only been routed from cities a number of times—they no longer have any presence in \textit{Whitestone} or \textit{Syngorn}—and each time, it has been by a concentrated citizens' uprising. In both instances, when they sought the source, \textit{Myriad} leaders were furious to discover an empty safe house containing only a single gold coin, marked by a grinning face: the calling card of the \textit{Golden Grin}. Even across the \textit{Lucidian Ocean}, the old enmity between the \textit{Myriad} and the \textit{Golden Grin} remains today.

\subsection{Figures of Interest}
These figures are major members of the \textit{Myriad} that you can use as NPCs in your game. You should also take inspiration from these figures to make up new NPCs for your own campaign.

\subsubsection{Lustran "the Backstabber" Zeth}
\textit{Human boss of \textit{Stilben}}

\noindent
Lustran was one of the first \textit{Clasp} members to turn double agent. By the time anyone realized his betrayal, he had granted the \textit{Myriad} their foothold in Tal'Dorei in the swampy port city of \textit{Stilben}. His lean, hungry face and graying hair give him a handsome, intimidating appearance. All \textit{Myriad} agents in Tal'Dorei respect him, and all members of the \textit{Clasp} hate him. He revels in his bad reputation as the renowned Backstabber of \textit{Stilben}.

\subsubsection{Tully "the Conclave of One" Dezzeram}
\textit{Gold \textit{dragonblood} boss of \textit{Emon}}

\noindent
Tully's distasteful nickname delights her inner circle in \textit{Emon's} underground. The \textit{Myriad} boss of \textit{Emon} is a position that usually changes hands by the season, thanks to the diligent efforts of the \textit{Clasp}, but Tully has evaded capture for three years, and grown rich and powerful in her reign of smuggling, trafficking, and grand burglary.

\subsubsection{Sylker "the Millionaire" Uttolot}
\textit{Half-elf boss of \textit{Kymal}}

\noindent
This gold-eyed half-elf, a former underling of the crime families of Shady Creek Run in \textit{Wildemount}, joined the \textit{Myriad} and became the first individual on that continent to hoard one million gold pieces. He has a masterful mind for money, and spearheaded the \textit{Myriad's} infiltration of \textit{Kymal}. "Gold is a disease," he often says, "and I'm bloody sick."

\section{The Remnants}
In a time long past, an archmage of unparalleled power known as The Whispered One attempted to ascend to godhood by conducting the Ritual of Seeding from high atop the citadel of Thar Ampala within the Plane of Shadow. The Whispered One surrounded himself with a cult of fanatical mages and empowered murderers to protect him during his apotheosis. Despite these protections, he failed. His cult, however, survived, and helped him return to power centuries later, after the defeat of the \textit{Chroma Conclave}. He succeeded in his bid for godhood, but was kept from utterly dominating the Material Plane at the last second by Vox Machina, who banished him beyond the Divine Gate, where he was shackled by all the same limitations as the other gods.

His cult, the Remnants, went into hiding once again. Their fanatical ranks, bolstered by his true apotheosis, now seethe in secret, convinced that there is a way to grant their god the ultimate power he sought. Many of them have wormed their way into positions of power within Tal'Dorei's political institutions, churches, arcane organizations, and criminal outfits. Many more have joined up with the dubiously moral, profit-driven arcanists of the \textit{League of Miracles}. They exist in shadows, lurking, waiting, longing for a new dark prophet to lead them to glory.

\subsection{Goals}
The Whispered One achieved an incomplete victory over twenty years ago, and now his cult is building a new generation of faithful to make his megalomaniacal dreams a reality. The cult's leadership believes that reuniting these malign artifacts will give them greater power to commune with their god and let them begin their quest in earnest to reach him through the Divine Gate. This leadership is known as the Circle of the Dismembered. Each cult leader in the circle commands a branch of the cult, named for a part of The Whispered One body. These factions are:

\subsubsection{The Unveiled Sight}
Members of this cult are tasked with scouring the planes for arcane lore that could restore their master to the Material Plane.

\subsubsection{The Mordant Voice}
Members of this cult serve their master by infiltrating other factions across Tal'Dorei and subverting their purposes with subtle propaganda.

\subsubsection{The Devouring Blood}
Members of this cult specialize in the art of silent death. They train in stealth and necromancy to become truly horrific assassins.

\subsubsection{The Unseen Touch}
Members of this cult are trained in magic that corrupts the mind and perverts the will, making them vile puppet-masters that prepare the world for their god's return.

\subsubsection{The Insatiable Heart}
Members of this cult do the vital ground-level work of seeking out new blood. They are trained to exploit the emotional weaknesses of their targets to convert them to the cult, without the need of magic.

\subsection{Relationships}
The \textit{Remnants} remain largely unknown, even after their master's grand ascension, and wish to keep it that way. Those who are aware of their presence are either cult members or people who have uncovered their insidious existence and wish to stamp it out wherever it may appear. Thus, there are rarely any alliances made by or with the \textit{Remnants}.

\subsection{Figures of Interest}
These figures are major members of the \textit{Remnants} that you can use as NPCs in your game. You should also take inspiration from these figures to make up new NPCs for your own campaign.

\subsubsection{Quinton Puck}
\textit{Human jeweler, Master of the Unveiled Sight}

\noindent
While he runs a middling jewelry business within \textit{Emon}, Quinton is secretly the leader of the Unveiled Sight. He is a master of creating and piercing illusions, deciphering ancient script, and all things to do with the magical arts of obscurement and unveiling. He can be recognized by his blood-red spectacles.

\subsubsection{Ixrattu Khar}
\textit{Tiefling vampire, Master of the Devouring Blood}

\noindent
This escapee from the \textit{Black Bastille} prison in \textit{Emon} is Tal'Dorei's most infamous mass murderer, who once slaughtered and consumed hundreds of innocents on a deadly killing spree from \textit{Kymal} to \textit{Emon}. During her imprisonment, Ixrattu was plagued by shadowy visions of a woman beckoning to her from beyond the bars. The visions pointed her towards weaknesses within the \textit{Black Bastille's} guard, and helped her achieve the unthinkable: escape. Once free, she vanished for a year as she scoured the countryside, puzzling over the cryptic whispers of her continuing visions.

At the end of her quest was a fateful discovery: within a lightless underground passage, she found the venomous crossbow \textit{Condemner}, a \textit{Vestige of Divergence}. Now she lurks in \textit{Kraghammer}, where she trains a network of elite assassins who gleefully dispatch anyone the Unveiled Sight identifies as an enemy of the \textit{Remnants}.

\section{Wardens of Syngorn}
The elven city-state of \textit{Syngorn} was established in the \textit{Verdant Expanse} by the sorceress Yenlara, the wood \textit{elves'} first leader after the \textit{Divergence}. Today, the physical and cultural wellbeing of \textit{Syngorn} is safeguarded by the three offices of elders called \textit{Wardens}, united by the High Warden. Their primary duty is to protect the city and the surrounding lands from intrusion.

\subsection{Goals}
The \textit{Wardens} believe that preserving the ancient culture of \textit{Syngorn} is of utmost importance. Even as their fate grows ever more closely tied to the fate of all Tal'Dorei, the traditions that make \textit{Syngorn} unique must be conserved. The three offices of the \textit{Wardens} are the Verdant Lord, the \textit{Voice of Memory}, and the Guildrunner. The office of Verdant Lord commands \textit{Syngorn's} formidable military, which is maintained and scattered throughout the \textit{Verdant Expanse}, ready to be called to \textit{Syngorn's} defense at a moment's notice. As preservers of culture, the \textit{Wardens} also encourage public respect for the arts, including the art of magic. It's not uncommon for a \textit{Syngornian} to rise in social status after being publicly praised by a warden for a masterful artistic creation or performance.

\subsection{Relationships}
\textit{Syngorn} has long been an isolationist society. As a mostly elven enclave, they have often been criticized as xenophobic, or even racist, toward outsiders. There is truth to these claims. Nevertheless, the rise and fall of the \textit{Chroma Conclave} has goaded even some of \textit{Syngorn's} most conservative voices into strengthening the bonds of trade, travel, and kinship with the other free peoples of Tal'Dorei. Even the people of \textit{Kraghammer}, who many elves still remember siding with the tyrant \textit{Drassig} in centuries past, have been welcomed as friends and allies for the loyalty and courage they displayed in the battle against \textit{Thordak}.

\subsection{Figures of Interest}
These figures are major members of the \textit{Wardens of Syngorn} that you can use as NPCs in your game. You should also take inspiration from these figures to make up new NPCs for your own campaign.

\subsubsection{High Warden Tirelda}
\textit{Elf Warden of \textit{Syngorn}}

\noindent
Though she is in her twilight years, even by elven standards, High Warden Tirelda is a trusted descendent of Yenlara, founder of \textit{Syngorn}. She is skilled with magic and blessed with calm wisdom. Thus, she was entrusted with the highest level of responsibility amongst the \textit{Wardens}—to maintain law, guide society, and preside over all other \textit{Wardens}.

\subsubsection{Verdant Lord Celindar}
\textit{Elf Warden of \textit{Syngorn}}

\noindent
Head of the Verdant Guard of \textit{Syngorn}, and an accomplished warrior in their own right, Celindar handles all military strategy and instruction within the forest's bounds. They only return to \textit{Syngorn} when called by their fellow wardens, and can otherwise constantly be found prowling the \textit{Verdant Expanse}. They use any pronouns, and have relinquished any sort of gender presentation after receiving a vision from the The Arch Heart.

\subsubsection{Ouestra, the Voice of Memory}
\textit{Elf Warden of \textit{Syngorn}}

\noindent
A renowned historian and a beloved vocalist, Ouestra is charged with the preservation of \textit{Syngorn's} history and culture. It's said her most magical songs can control moonlight and bring light even to the Abyss itself. These days, she has grown somewhat unpopular by also singing songs of \textit{Syngorn's} darker days, for she says true greatness can only be inspired by learning from failure.

\subsubsection{Guildrunner Rawndel}
\textit{Elf Warden of \textit{Syngorn}}

\noindent
One of the oldest living elves within \textit{Syngorn}, Rawndel is more cunning and insightful than most, allowing him to be a very effective head of commerce, business, and the treasury.

\chapter{Tal'Dorei Gazetteer}
\begin{DndReadAloud}{}
\DndQuote%
{}
{\enquote{I've wandered over many a shore and plain in my journeys, and my mind still finds itself drawn back to Tal'Dorei. Ain't no place in this world more vibrant; the land feels young. Unscarred. Mind you, that sort of calm causes my old bones to itch, too. Folk in that land are kind and honest, and its vistas are second to none, yet...well, I'll just say that cities are cities, and the wilds are the wilds. Tal'Dorei's got plenty of both. Keep a sword or a wand with you wherever you go, adventurer.}}
{Norad Firth}
\end{DndReadAloud}

This chapter contains a wealth of information about the lands of Tal'Dorei, the adventures to be found therein, and the many centers of civilization to be found across its vast wilds. Consider this chapter a compilation of a wide variety of assorted notes, research, and scattered publications of historians of the \textit{Library of the Cobalt Soul} and the \textit{Alabaster Lyceum}.

Boundless possibilities await eager adventurers of all kinds who explore Tal'Dorei, from the vast fields of the \textit{Dividing Plains} to the treacherous crags of Tal'Dorei's many mountains, from the storm-wracked \textit{Lucidian Coast} to the immemorial depths of the \textit{Rifenmist Jungle}. Countless nomadic civilizations exist within these realms of magic and monsters, as do dozens of tiny villages, thriving towns, and bustling cities. The people of Tal'Dorei do what they can to survive, and share the tales of treasures and adventure with any who will listen.  Consider this gazetteer a compilation of tales and fragmented histories, passed down from bards and scholars. Your campaign can follow its descriptions to the letter, or diverge as wildly as you see fit.

\section{Settlements}
Every permanent settlement in Tal'Dorei has a few basic statistics that impart a top-level impression of that place. This includes what type of settlement it is—typically a village, town, or city—its population, and its rough demographics. Don't feel restricted by these demographics if you want your character to come from a specific location but you don't see their race listed there. People of all races can be found in just about every community in Tal'Dorei, and your character can come from anywhere.

Additionally, cities and large towns often have a number of small rural communities outside their fortified walls and population center. The population numbers of these larger settlements include the people in the outlying areas. On the other hand, the official populations of a village or other small settlement just includes the people living in the village itself.

In fact, though there are a few villages and small settlements listed in this gazetteer, Tal'Dorei is home to hundreds, if not thousands, of tiny settlements too small to detail or place on the map. The Game Master can make up any number of villages or other communities wherever they wish in Tal'Dorei without changing any of the lore presented in this gazetteer. And of course, the GM can also create large settlements unique to their version of Tal'Dorei—the only difference is that they should be prepared to think about how their version of Tal'Dorei is different, now that they've added another hub of population, trade, politics, and culture.

\section{Story Hooks}
This chapter contains a variety of story hooks that you can use as a starting point for Tal'Dorei adventures of your own creation. \textit{Chapter 5} of this book contains advice for expanding these hooks into full adventures, and even more advice for creating RPG stories can be found in the fifth edition core rules.

These plot hooks are tied to specific locations, and all have a level range. These level recommendations aren't strictly mandated, but have generally been chosen either because of the monsters involved, because of the narrative weight of the location or characters involved, or to indicate how epic and world-affecting the adventure should be.


\begin{itemize}
\item For low-level characters: 1st to 4th level
\item For mid-level characters: 5th to 10th level
\item For high-level characters: 11th to 16th level
\item For epic-level characters: 17th to 20th level
\end{itemize}

Additionally, some plot hooks are marked as "any level." These hooks present stories broad enough to suit characters of any level, anywhere in their party's story. These hooks generally don't suggest specific monsters to use as villains. When creating adventures based off of these hooks, use the encounter-building advice in the fifth edition core rules to choose appropriate foes for your players.

\section{Lucidian Coast}
The stormy Lucidian Ocean crashes furiously against the rocky coasts of eastern Tal'Dorei, chewing inlets into the coastline and spurring vast banks of impenetrable fog across the waters. Most ships rightly avoid the volatile northern waters, instead docking at southern ports like \textit{Stilben}. Adventurous folk leave their seafaring vessels behind and brave the \textit{Mooren River Run} to \textit{Drynna}. Yet even there, the Lucidian Coast is not without its dangers; vicious wildlife stalks the wilderness, and flocks of harpies have harried the shores since time immemorial, drawing lost sailors to their doom upon the rocks.

\begin{description}
\item[Majority Faiths:]
The Changebringer

\item[Minority Faiths:]
The Lawbearer

\item[Imports:]
Livestock, lumber

\item[Exports:]
Fish, ships, grain, iron, silver, gold, trade goods from \textit{Wildemount}
\end{description}

\subsection{Drynna}

\begin{description}
\item[Town:]
Population 2,765 (75\% humans, 8\% gnomes, 8\% halflings, 4\% elves, 5\% other races)

\end{description}

\noindent Along the shore of \textit{Mooren Lake} rests the sleepy town of Drynna. The community thrives on the bountiful wildlife around the lake, and the Drynna fisherfolk are responsible for most of eastern Tal'Dorei's freshwater fish supply. The social guidance of the city revolves around the Sunrise Lodge (or simply "the Lodge"), a membership of community leaders who meet to discuss and vote on important social and legal matters. The Lodge has also sent a variety of increasingly frustrated petitions to the \textit{Tal'Dorei Council}, since their community's small size means they have traditionally lacked a delegate of their own on the council. However, the recent surge in prospectors has increased their population enough for a member of the Lodge to attend council sessions, and the Lodge is doing everything it can to make sure the boom lasts as long as possible before it falls into bust.

\subsubsection{Boomtown}
\textit{Drynna} has a small militia to protect the town from wild beasts and roving monsters, and to capture horse thieves and other petty criminals. This militia is presently overwhelmed by hundreds of prospectors who have swept into town over the past three years after hearing rumors that chunks of \textit{whitestone} have been found in \textit{Mooren Lake}, washed down the river from the \textit{Alabaster Sierras}.

Given the high price of \textit{residuum}, this new occurrence has made \textit{Drynna} into something of a boomtown. Thousands of small homesteads now surround the town center, and burglary and banditry are both on the rise.

\subsubsection{Drynna Adventures}
Game Masters who set their adventure in \textit{Drynna} can use these plot hooks for inspiration.

\subparagraph{Unfriendly Waters (mid-level)}
Recently, \textit{Drynna} has suffered a number of losses, with prospectors and fisherfolk being attacked. Worse still, the supply of fish itself seems to be declining rapidly; without this local food staple, villages across eastern Tal'Dorei will certainly starve. The party hears the call for mercenaries reaching as far as \textit{Westruun} and \textit{Stilben}. The culprit is a \textbf{hydra} within the lake that has gained intelligence by living for so long in the \textit{residuum}-infused waters. It has an Intelligence score of 12 (+1), can speak Common, and commands a tribe of \textbf{lizardfolk}.

\subparagraph{Bleak Inheritance (any level)}
Last year, Hareth Val Bardo, a younger member of the Sunrise Lodge, made a pact with the Moon Mistress (see "\textit{Mooren Lake}," below). Hareth's elder brother, Pauvren, died in his sleep shortly thereafter, leaving the entire Val Bardo family fortune to Hareth, the only surviving heir. However, nightmares now haunt Hareth every night, and violent urges surge through his mind. The Lodge, having noticed Hareth's paranoid eyes and twitchy hands, calls upon outsiders to investigate.

\subsection{K'Tawl Swamp}
The sprawling swamplands of \textit{K'Tawl} cover much of the southern \textit{Lucidian Coast}, enveloping \textit{Stilben} and the K'Tawl Bay. The saltwater swamp is thick with cypress and tupelo trees reaching up from muddy waters that range, in places, from one to ten feet deep—or more. The hot, muggy air and thick, sludge-covered landscape host all manner of terrible beasts, including venomous flying snakes and man-eating giant frogs.

Bands of frog-headed humanoids harry the edges of \textit{Stilben} from their deep-marsh encampments, while a vicious school of \textbf{sahuagin} called the Ghormauth command the shoreline of \textit{K'Tawl}. The bloodthirsty \textit{Ravagers} use the forested edges of the swamp to ambush prey—as do countless bandits and petty thieves. As such, caravans and messengers must remain on constant alert for the possibility of an ambush. With such dangers lurking in the swamp, travel to and from \textit{Stilben} is rarely made without an armed escort.

\subsubsection{K'Tawl Swamp Adventures}
Game Masters who set their adventure in \textit{K'Tawl Swamp} can use these plot hooks for inspiration.

\subparagraph{They Shamble (mid-level)}
People have been disappearing from the isolated swamp-burghs within the \textit{K'Tawl}. The characters, perhaps en route to \textit{Stilben}, spend the night in the village of Dunghill. After hearing tales of the disappearances, the characters awaken the next morning to screams—as a terrified human is carried off within the rotting orifice of a \textbf{shambling mound} of swamp vegetation.

\subparagraph{Monolith of the Hells (high-level)}
Why have members of the Torchraiser herd been found dead, their hearts exploded and their teeth plucked from their mouths? Why have the stagnant \textit{K'Tawl} waters suddenly begun to flow toward the dismal center of the swamp? A divine messenger of the The Everlight visits the characters, its celestial form torn and wavering, and dies before it can impart its request. On the angel's lips is a single dying phrase: "Destroy K'Tawl's Heart."

\subsection{Mooren Lake}
North of the \textit{Summit Peaks} is a shimmering lake fed by the snowmelt from the \textit{Alabaster Sierras} by the \textit{Mooren River Run}. \textit{Mooren Lake} is surrounded both by bountiful wildlife and the frightful creatures attracted to such a thriving ecosystem, including roving \textbf{lizardfolk}, flocks of harpies, and several hyrdas that lurk in the lake's depths.

Centuries of \textit{whitestone} runoff from the \textit{Alabaster Sierras} have allowed deposits of \textit{residuum} to collect at the bottom of the lake—and the \textit{residuum} has formed into a thrumming organ of raw magical energy, beginning to gain sentience. The twin islands in the center of the lake are blanketed by an ever-present mist, and folk tales tell of a beautiful, ageless woman, known only as the Moon Mistress, who watches over the island shores. It's said that any mortal who seeks an audience is judged by her, and is either granted a favor, a boon of unveiled knowledge, or a swift and terrible end.

\subsubsection{Mooren Lake Adventures}
Game Masters who set their adventure in \textit{Mooren Lake} can use this plot hook for inspiration.

\subparagraph{Renew the Mists (mid-level)}
The Moon Mistress, a \textbf{night hag} who takes the guise of a water nymph, uses the fading mists to shroud her location from a blood hunter named Perron Brill, a woman dedicated to hunting the Mistress down out of vengeance. The mists require a blasphemous mortal sacrifice to renew, and the Mistress is willing to reward those who would aid her in retrieving a victim—willing or otherwise. Perron Brill resides in nearby \textit{Drynna}, and though she is no saint herself, more honorable characters may think it best to help the hunter find her prey in \textit{Rootgarden Marsh}.

\subsection{Rootgarden Marsh}
The shallows of \textit{Mooren Lake} leak into the tender soil of the lake-bed and spread eastward, creating the cold and desolate \textit{Rootgarden Marsh}. Tall, thin trees lift five or more feet above the muck upon stilt-like roots, like freshwater mangroves. Their sparse placement gives them the appearance of a village of scattered huts. The fresh water mingles with salt here, leaving only adapted vegetation. The smell of moss and rot fills the air, while strange, feral creatures leap from tree to tree, or patiently wait from within their dome-like tree-root dens to snatch wary travelers. An organized group of gnoll marauders called the Moonsteeth keep their cavern hideout here, striking out at easy targets outside of \textit{Drynna} and \textit{Mooren Lake}.

Since the region is populated almost entirely by monsters, most humanoids avoid it whenever possible. Those few who explore the deadly marsh do so only out of necessity, either because they need to unlock Rootgarden's many secrets, or because "civilization" is even more dangerous to them than the primal dangers found within.

\subsubsection{Rootgarden Marsh Adventures}
Game Masters who set their adventure in \textit{Rootgarden Marsh} can use these plot hooks for inspiration.

\subparagraph{The Ley-Knot (mid-level)}
Without warning, arcane magic stops functioning within a hundred-mile radius of \textit{Rootgarden Marsh}. As panic spreads across northeastern Tal'Dorei, the characters or a patron of theirs discover that a major ley line, the source of arcane power in this part of the world, has been twisted and tied into a knot at the center of the marsh. This knot has been tied by High Moonspeaker Luneskon of the Moonsteeth gnolls. Her ambitions are small, but whatever method she used to tangle the ley lines must be erased from history—knowledge of its existence could bring disaster to all of Tal'Dorei.

\subparagraph{Death's Garden (any level)}
Rootgarden's strange ecology has produced many distinctive and powerful forms of plant life. During the course of a campaign, a group of adventurers may need to create a unique potion requiring the essence of manawort, a sentient mushroom found only within Rootgarden. The manawort mushroom is fearful of being harvested, and has made close friends with a black dragon named Thardraxxus, who keeps the mushroom and its sporelings safe. This dragon can be of any age, depending on what level the characters are.

\subsection{Stilben}

\begin{description}
\item[Large Town:]
Population 10,660 (67\% humans, 11\% elves, 11\% halflings, 5\% half-elves, 6\% other races)

\end{description}

\noindent Stilben stinks to high heaven. Most sailors who pass through it—be they honest traders of the Clovis Concord in \textit{Wildemount} or pirates seeking their next conquest—call this dismal burgh the Rotted Lot. The stink and the crime that run through this town make the name an apt one, but things aren't all bad in Stilben. The largest port on this coast of Tal'Dorei, Stilben is the focal point of international trade with \textit{Wildemount}. Even though Stilben's permanent population is a scant ten thousand, its more transient population of sailors who stay for a few days or weeks at a time is just as large, and fully triples Stilben's total population in the summer months.

Predictably, the vast majority of Stilben's business caters to folk who won't stay for long. There are brothels aplenty, at least one inn and one tavern each on every street, and countless other comforts for the weary traveler. The town is nestled in the sticky depths of the \textit{K'Tawl Swamp}, and its putrescent stink mingles with the humidity and the ever-present buzz of insects to produce an air of misery around Stilben's outer districts. Pushed out of Stilben's cozy, cosmopolitan heart, its poor and disenfranchised citizens—derisively called muckdwellers by the town's bigshots—struggle to eke out a living in the \textit{K'Tawl's} murky waters.

Central Stilben houses most of its businesses, residences, and attractions. Those living in Stilben's heart like to present a clean, "civilized" view of their town, maintaining it as a thriving commercial hub. The cross-continental \textit{Silvercut Roadway} leads from Stilben all the way to \textit{Emon}, and its robust trade of rare goods, textiles, and spices has given rise to a formidable merchant ruling class. Though the aging Margrave Wendle Truss is Stilben's ruler in name, it is the guilds that hold true power. Their vast wealth and sizeable "protection" force have impelled the bombastic-yet-ineffectual margrave to acquiesce to the guilds' will in order to maintain peace, as well as to preserve his own station within the city.

Order is kept by the Waterwatch, Stilben's official constabulary, but corruption is common within their ranks. The past few decades have been hard for the \textit{Clasp} in Stilben. Though the local Spireling has kept a tight grip on the margrave and the Waterwatch, the \textit{Myriad} has used bribes, intimidation, and every dirty trick in the book to winnow away the \textit{Clasp's} power. Now, Stilben is little more than a crucible for the power struggle between the two criminal organizations.

\begin{DndSidebar}[float=!b]{Origins of Vox Machina}
The citizens of \textit{Stilben} do not know that some of Exandria's greatest heroes, Vox Machina, first met in their humble city and went on their first adventure in the muck of the \textit{K'Tawl Swamp}. Back then, even the exalted Vox Machina were just two-copper adventurers, quickly forgotten by \textit{Stilben's} haughty elite. One dwarf named Heinrich Runescribe, a junior member of the \textit{Alabaster Lyceum's} revered Lorekeeper Society, practically worships the heroes that defeated the \textit{Chroma Conclave}, and has made it his mission to record all of their travels. He has found a wealth of evidence that the heroes of Vox Machina first gathered in \textit{Stilben} under inauspicious circumstances and were dragged into an unexpected battle against the \textit{Myriad}.



\end{DndSidebar}

\subsubsection{Stilben Adventures}
Game Masters who set their adventure in \textit{Stilben} can use these plot hooks for inspiration.

\subparagraph{Lost Below (low-level)}
A \textit{Stilben} guild known as the Relic Seekers partnered with the \textit{Clasp} to smuggle a hoard of stolen \textit{Syngornian} artifacts through the heart of the \textit{K'Tawl Swamp} to the port of \textit{Stilben}. The smugglers were last seen by a \textit{Clasp} agent watching a bridge in the western swamp, but two weeks have passed since their last contact. They should have arrived by now. The characters are approached by one of several groups: an agent of the Relic Seekers in search of the relics, one of the \textit{Clasp} seeking to find their smugglers, or a \textit{Myriad} taskmaster who seeks to undermine the \textit{Clasp}. Their first contact: the bridge guard, a hulking tiefling \textbf{Clasp enforcer} named L'Arkhelle.

\subparagraph{Ashari Agitator (mid-level)}
Gang warfare between the \textit{Clasp} and the \textit{Myriad} has been a daily affair in \textit{Stilben} for the past thirty years, but something is changing. A cruel \textbf{Ashari waverider} named Galessi has joined the \textit{Myriad} and is using her elemental magic to turn the tide of this gang war. The \textit{Clasp}, the Waterwatch, and even the Air Ashari of \textit{Zephrah} want to see Galessi laid low, but she has a secret base within the swamp, protected by a charmed \textbf{giant crocodile} and putrid water elementals.

\subparagraph{Demonblood Vengeance (high-level)}
A \textbf{marilith} demon using the guise of a blonde human named Iselda has been lurking on the outskirts of \textit{Stilben} since her defeat at the hands of Vox Machina decades ago. Her demonic ichor has spawned a brood of demonfeed spiders that now view her as their mother. Iselda, sensing an opportunity, dedicated herself to the The Spider Queen and attracted a small legion of dark elf fanatics to her side. Aided by all these evil creatures, she plans on conquering \textit{Stilben}, not for the \textit{Myriad}, but for herself and her evil god.

\subsection{Summit Peaks}
The Summit Peaks loom over the southeastern plains of Tal'Dorei, like great spires of rock reaching for the sky. The tall, thin mountains burst so aggressively from the surrounding fields and swampland that even travelers from distant lands recognize them as a landmark of eastern Tal'Dorei. They are a welcome sight for any ship's crew making their way to shore. The mountain valleys recess into marshy gorges that house all manner of creatures migrating from the \textit{K'Tawl}. Reclusive hill giants and stone giants also make their homes in the marshy valleys. As the mountains rise higher, strong winds blow across the battered rock, leaving strangely smooth peaks that harbor \textbf{griffon} lairs, \textbf{harpy} dens, and convocations of giant eagles hunting the lower lands.

\subsection{Zephrah}

\begin{description}
\item[Small Town:]
Population 1,310 (62\% half-elves, 12\% elemental ancestry, 10\% humans, 9\% halflings, 5\% tieflings, 2\% other races)

\end{description}

\noindent High atop a mountain aerie within the \textit{Summit Peaks} sits Zephrah, the tribal home of the Air Ashari people, keeping vigil over a gate into the Elemental Plane of Air. \textit{The Ashari} people worship and live by the balance of the natural world, and the people of Zephrah have adapted to life among the clouds. Most of the tribe consists of skilled hunters and climbers who hunt mountain goats, large birds, and other cliff-dwelling creatures, along with subsisting on high-elevation vegetation. Utilizing the proximity of powerful elemental magic, the druidic masters of the Air Ashari craft skysail to ride the ever-present winds, allowing their people to traverse the sheer cliffs of the mountains. The leader of Zephrah is called the Voice of the Tempest; they oversee all travel through the rift, all major druidic rituals, and any visiting outsiders.

Though it's still a small settlement, Zephrah has grown greatly in the past several years. Many people with an innate connection to the element of air have petitioned the current Voice of the Tempest for her permission to live here among \textit{the Ashari}. Likewise, a number of druids and other spellcasters who wish to learn the ways of the Air Ashari have come to live in relative isolation here—after they have proven their good intentions and respect for the sanctity of the skies.

\subsubsection{Diplomacy and Developments}
\textit{Zephrah} has grown in more than just population. The close friendship shared between Keyleth, Voice of the Tempest, and the genius inventor \textbf{Percival de Rolo} of \textit{Whitestone} has led to a number of modern technologies being adopted in \textit{Zephrah}. Windmills now dot the cliffsides, milling grain purchased from farms on the \textit{Dividing Plains}. Keyleth continues to work with Percival to improve the lives of her people through technology and innovation. The operation of two of these mills is also kept private from the rest of \textit{Zephrah}, and they're used to refine \textit{whitestone} into \textit{residuum} for the construction of magic items to arm the Ashari skydancers.

Keyleth's confident leadership has also guided her to make further diplomatic ties between her people and the rest of Tal'Dorei. These attempts at breaking down the traditional isolationism of \textit{the Ashari} have been met with skepticism from within, slowing the process to a crawl. Allowing any outsiders into \textit{Zephrah} at all has been seen as a huge step. Beyond this, Keyleth has pushed for greater economic and political involvement, so that the needs of \textit{Zephrah} can be heard in the \textit{Tal'Dorei Council}. Resistance toward political integration has mostly been internal within \textit{the Ashari}, but Keyleth considers any progress at all a good start.

\subsubsection{Defense}
\textit{Zephrah} has never maintained a formal army, though past Voices of the Tempest have trained martial forces and druidic warriors to defend their home against the monsters of the \textit{Summit Peaks}. As Keyleth attempts to build bonds between her people and the rest of Tal'Dorei, she has expanded the training of new skydancers and created a unit of druidic warriors capable of responding to any emergency across Tal'Dorei. If a threat like the \textit{Chroma Conclave} were to rise again in Tal'Dorei, the skydancers of \textit{Zephrah} would be there to help save the nation's people.

Thanks to \textit{Zephrah's} mountainous terrain, \textit{the Ashari} enclave has never fallen under siege. Because of the dangerous beasts that lurk in the mountains, most Zephran adults have learned how to wield a \textit{spear} or a \textit{staff} to defend themselves. Roughly one-in-ten people have undergone a small amount of elemental training—enough to be able to cast the \textit{gust of wind} spell a number of times per day equal to their Wisdom modifier (minimum of once).

Were \textit{Zephrah} to fall under attack, the skydancers would defend the town and its elemental rift for as long as possible while the rest of its people evacuated to hidden retreats in the mountain caves.

\subsubsection{Zephrah Adventures}
Game Masters who set their adventure in \textit{Zephrah} can use these plot hooks for inspiration.

\subparagraph{Poached Eggs (low-level)}
The local Ashari have noticed a terrible trend of poachers climbing the peaks to steal entire \textbf{griffon} nests full of eggs, leaving the parents slain. The efficiency, brutality, and growing rate of these attacks suggest the hunters are dangerous professionals, and \textit{the Ashari} are seeking aid in finding and punishing those responsible before the locally endangered creatures are all dead or taken. The culprits, human poachers (use \textbf{scout} statistics) led by a \textbf{quasit} demon named Ryzzix, are fencing the stolen eggs on the \textit{Stilben} black market.

\subparagraph{A Welcome Overstayed (mid-level)}
Shaktia, an overbearing \textbf{djinni} soldier from the Plane of Air, has recently discovered and come through the elemental rift into \textit{Zephrah}. However, the \textbf{djinni} refuses to leave, instead choosing to bother the people with pranks, debates, and invitations to become his well-paid servants. The Voice of the Tempest, at the end of her patience, seeks adventurers to capture Shaktia and return him through the rift to the Elemental Plane of Air. However, when the adventurers learn what caused the \textbf{djinni} to leave his home plane in the first place, they may find it hard to force him to leave.

\section{Alabaster Sierras}
In the northeastern reaches of the continent, where the waters of the \textit{Lucidian Ocean} break between the rocky shores of Tal'Dorei, a lush treeline fills the valleys between the chalk-white bluffs of the Alabaster Sierras. This gelid northern territory has for decades been beyond the reach of the \textit{Tal'Dorei Council}, and fallen under the rule of the \textit{Chamber of Whitestone}. However, over twenty years ago, thousands of \textit{Emonian} refugees fled to Whitestone to ride out the devastation of the \textit{Chroma Conclave}. These refugees included the bereaved family of the Last Sovereign. Since those dark days, the \textit{Chamber of Whitestone} has enjoyed a close relationship with the \textit{Tal'Dorei Council}. Lady Vex'ahlia de Rolo, a member of the \textit{Chamber of Whitestone}, also sits on the \textit{Tal'Dorei Council} itself, further strengthening the bonds of unity between this once-isolated land and the rest of the realm.

Beyond \textit{Whitestone}, the region's vibrant-yet-gothic urban center, the Alabaster Sierras are framed by dangerously windy shoreside cliffs and dry, desolate mountains swarming with harpies, cold-adapted wyverns, and other hungry monstrosities. Beneath the peaks, the valley opens into the thickly wooded pine forest known as the \textit{Parchwood}. This dense, temperate rainforest is beautiful and lush, but many shadowed threats loom in the dark of the wood.

The northern center of the valley marks a single clearing, where the humble city of \textit{Whitestone} stands as the source of rule and law in this land. The people here have been hardened through strife and exalted through freedom, enough to know honor and appreciate loyalty like few other civilizations do.


\begin{description}
\item[Majority Faiths:]
The Dawnfather
\item[Minority Faiths:]
The Wildmother

\item[Imports:]
Industrial and precious metals, cloth, livestock, grain

\item[Exports:]
Granite, limestone, whitestone

\end{description}


\subsection{Mooren River Run}
While many smaller streams and rivers run through the \textit{Parchwood Timberland}, gathering rain and snow runoff from the nearby mountains, they nearly all empty into the massive Mooren River Run. A powerful river that tears through the heavy forests of the \textit{Alabaster Sierras}, the Mooren River Run flows fiercely southward, out of the woods and into \textit{Mooren Lake}, near the \textit{Lucidian Coast}. Freshwater fish are in ready supply along the river's path, and \textit{Whitestone} maintains a well-guarded dock for their fishermen to feed their people.

The small town surrounding the dock is simply called Mooren, and seasonally is home to some five thousand people of various ancestries in the spring and summer. When conditions grow too cold, mages from \textit{Whitestone} are called to place \textit{arcane lock} upon all establishments and to place the empty town in lockdown until next season.

Due to the plentiful nature of the river, many denizens of the \textit{Parchwood} wander its banks to reap its bounty, or to seek bounty from those who frequent it.

\subsubsection{Mooren River Run Adventures}
Game Masters who set their adventure in \textit{Mooren River Run} can use these plot hooks for inspiration.

\subparagraph{Bear of the White River (low-level)}
While traveling near \textit{Whitestone} or through the \textit{Parchwood Timberland}, the characters find the corpse of half-orc forester; a successful DC 10 Intelligence (Nature) check reveals that he was savaged by a large bear. This half-orc's forester clothing marks her as a cleric of the The Stormlord and a member of the \textit{Rivermaw herd}, a long way from the \textit{Dividing Plains}. Investigating the scene shows that bear tracks in the mud lead westward, toward the \textit{Mooren River Run}. These tracks, when followed, soon fade into the bare footprints of a humanoid—the killer was a \textbf{werebear}, but who could it be?

\subparagraph{Off-Season Suspicion (mid-level)}
The cold nights of autumn and winter force even the most stalwart fisherfolk away from Mooren. Some Mooren residents who have returned to \textit{Whitestone} for the off-season have seen lights in the abandoned town, and are spreading rumors that some of the mages contracted to "lock up" after everyone has left town are actually agents of the \textit{Myriad}. The truth is that these mages are mercenaries hired by the \textit{League of Miracles} to secretly conduct experiments on cold snap spirits to transform them into loyal elemental warriors.

\subsection{Parchwood Timberlands}
Dark, dense, and often assailed by snow and cold winds, the Parchwood Timberlands consume most of the valleys within the \textit{Alabaster Sierras} with pine trees and heavily overgrown pathways. The grim and forbidding interior of the Parchwood is home to many grim inhabitants, from the spectral remains of unlucky travelers past to wandering blights spawned by corrupt trees that consume the life around them. Howls can be heard in the darkest of nights, giving rumor to \textbf{werewolf} hunting parties that stalk the wood. A handful of nearby farmers swear they've encountered a community of territorial \textit{centaurs} in the southern region of the Parchwood, though such tales are largely met with derision. It's common knowledge that if you intend to travel through the thick fog of the Parchwood Timberlands, you stick to the roads, you bring some muscle, and you pray to whatever gods bring you comfort.

\subsubsection{Parchwood Timberlands Adventures}
Game Masters who set their adventure in the \textit{Parchwood Timberlands} can use these plot hooks for inspiration.

\subparagraph{Headless Horse-Man (mid-level)}
The \textit{centaurs} of the \textit{Parchwood} are no myth, and many are succumbing to a curse of undeath. One among them, a headless stallion-lord named Ichabarr, now leads the Herd of the Damned across the valleys of the \textit{Parchwood}. Several small settlements outside \textit{Whitestone} have told the \textit{Chamber} that these skeletal monsters have been spiriting away eligible young bachelors, and soon no men will be left in the Timberlands. These marauders are centaur skeletons. Their leader Ichabarr uses the same statistics, and can also use the optional Spirit Away feature.

\subparagraph{They Seek Only Rest (high-level)}
Those who travel the \textit{Parchwood} sometimes hear thunder roll through the valley, accompanied by the cracking of timbers, but no storm follows. Over twenty years ago, the \textit{Briarwoods} slaughtered and reanimated an entire clan of giants native to the \textit{Parchwood}. Some of these zombie giants survived the fall of the \textit{Briarwoods}. Most of these giants have been destroyed by the \textit{Grey Hunters}, but a few still remain. However, their mindless actions are not without purpose: each giant is tormented by being denied the halls of their ancestors, and quietly weeps for their lost afterlife as they roam.

Zombie giants can be represented by changing any giant's type to undead, decreasing its AC to 8, giving it resistance to bludgeoning damage from nonmagical attacks, and giving it immunity to poison damage and the poisoned condition. This does not affect its challenge rating.

\subsection{Salted Bluffs}
The tumultuous waters of the \textit{Lucidian Ocean} that encompass the northeast edge of Tal'Dorei end at the towering seaside cliffs called the Salted Bluffs. These cliffs mark the coastal boundary of the \textit{Alabaster Sierras} and are home to seemingly infinite cliffside caverns that overlook the ocean, sometimes hundreds of feet up. The waters along the base are rocky and dangerous to traverse by boat, and only the most skilled of captains even consider sailing through these treacherous waters. These factors make the Salted Bluffs an ideal hideaway for pirates, who have taken to these bluffs to hide away their riches or build a base of operations.

However, these outlaws must also deal with the colonies of harpies native to the region. They can never be too sure if a cave they're holing up in for the evening is abandoned—or if it's home to a beast of alluring song and shredding talons. Even empty caves house the remains of these harpies' brutal kills, for they sweep up unsuspecting creatures and dash them against the cliffs and rocks, recovering the broken corpses to feed upon safely.

\subsubsection{Salted Bluffs Adventures}
Game Masters who set their adventure on the Salted Bluffs can use these plot hooks for inspiration.

\subparagraph{Crow's Nest (low-level)}
The Ashwing \textbf{harpy} tribe that roost among the Salted Bluffs have recently commandeered the galley owned by the group of penitent corsairs and started raiding coastal villages. They are led by Captain Silvercomb, an arrogant, dashing, and well-spoken harpy grandmother. The other harpies are just pillaging for fun, but Silvercomb seems to be searching for something, and she always wears a pendant with a black feather around her neck. What could it be?

\subparagraph{Start Believing in Ghost Stories (low-level)}
Pirates are one thing. Undead pirates are another. A group of swashbucklers have willingly turned themselves in to the \textit{Chamber of Whitestone}, begging for absolution from the The Dawnfather for their crimes. They refuse to return to the cove where they kept their stolen treasure since it's guarded by the zombies and skeletons of their risen comrades—as well as the \textbf{wight} that killed them. Perhaps a group of adventurers can save the day?

\subsection{Shearing Channel}
Long ago, a land bridge once connected the landmasses of Tal'Dorei and \textit{Wildemount}. However, the catastrophic forces of the \textit{Calamity} blasted it apart, leaving in its place a deathly strait of jagged rock and gale-force winds. Numerous ships have been ruined against the tides of the channel, and sea creatures that thrive on wreckage have taken up residence beneath the waters. It's said that no honest sailor knows a safe passage through the channel, but rumors abound of hidden coves and unknown routes that the Revelry pirates of \textit{Wildemount's} Menagerie Coast use to safely traverse these deadly waters.

\begin{DndReadAloud}{}
\enquote{\itshape Amidst fluted white bones, serenaded by siren's call.\\
'Twas there did I bury my treasures, one and all.\\
But deaf and blind must the finder be,\\
Lest a briny death be that which ye truly seek.}
\end{DndReadAloud}

\subsubsection{Shearing Channel Adventures}
Game Masters who set their adventure in the \textit{Shearing Channel} can use this plot hook for inspiration.

\subparagraph{Scion of Oblivion (epic-level)}
The The Chained Oblivion has no physical form, but some residue of its power was cast into the ocean when the The Dawnfather pursued it from the \textit{Alabaster Sierras} to \textit{Gatshadow}. This seed of evil has had eons to gestate undisturbed in the tempestuous channel—and now it rises. A \textbf{kraken}, mutated by the The Chained Oblivion aberrant power, explodes from the water. Its mere presence causes the flesh and minds of sailors to fuse to their ships, sea creatures to hideously mutate into humanoid abominations, and the clouds above to rain thick, ichorous black droplets. A storm now lingers permanently over the \textit{Shearing Channel}, and forces from \textit{Whitestone} and the Menagerie Coast of western \textit{Wildemount} relay frantic messages, seeking to somehow halt this evil. It seems to be building an army here for some larger attack. What is it planning?

\subsection{Whitestone}

\begin{description}
\item[Small City:]
Population 17,710 (78\% humans, 7\% halflings, 6\% dwarves, 4\% tieflings, 5\% other races)

\end{description}

\noindent Over two hundred years ago, the de Rolo family helmed an expedition from \textit{Wildemount} to the shores of Tal'Dorei. When their vessel wrecked upon the stones of the \textit{Shearing Channel}, the de Rolos and their remaining crew survived for weeks in the \textit{Alabaster Sierras}, despite the tempestuous weather and slavering beasts. After a time, the survivors found a clearing containing a single tree that glowed with the radiance of the sun itself. They called it the \textit{Sun Tree}—a literal blessing of the The Dawnfather himself—and built their home around it.

Since then, the de Rolos' pioneering settlement has grown into a robust, if isolated, city. The sights of the city range from the gothic majesty of \textit{Castle Whitestone} to the radiant glory of the \textit{Sun Tree}, and from the quaint townhouses of the city center to the humble homesteads of its outlying townships. Because Whitestone is limited in fertile farmland, locally grown produce is scarce, leading to the necessity to develop a good trade relationship with the rural communities of the \textit{Turst Fields}.

Though the region's scarcity of vital resources made for a harrowing first few years, the de Rolo family stumbled upon a miraculous mineral they called \textit{whitestone}. This geological marvel both gave their city its name and quickly became its most important and most sought-after export. Before long, the de Rolos' sale of \textit{whitestone} allowed them to quickly amass a significant fortune and construct a princely keep within their city.

The city is constantly improved (and unsettled) by the brilliant tinkering of \textbf{Percival de Rolo}. His greatest creation, the \textit{Heart of Whitestone}, overlooks the city. Most prominent these days, however, is a rudimentary electric light grid. Using acid pits in the relatively unexplored caverns beneath the city as a sort of battery, six dozen electric lamps throughout the city are given just enough power to gently illuminate the streets, the watch posts on the castle, and the city walls with a soft, greenish glow all through the night.

\subsubsection{Scars of the Briarwoods}
Lady Delilah and Lord Sylas Briarwood, former cult leaders of The Whispered One, were ousted from \textit{Whitestone} over twenty years ago. Though they are no more, the evil they committed still lingers. The nearly six years they ruled \textit{Whitestone} sapped life and hope from its people. Though homes have been rebuilt and the Briarwoods' undead monstrosities have been eradicated, the trauma of the vampires' reign of terror will linger for generations. Some people of \textit{Whitestone} were so deeply affected that they turned to necromancy and vampirism themselves, and have fled into the \textit{Alabaster Sierras} to learn the secrets of the Briarwoods' foul magic—and to emulate their takeover of \textit{Whitestone}. The enigmatic \textit{Grey Hunters}, the secretive elite warriors of the \textit{Chamber of Whitestone}, have worked tirelessly to pull out the \textit{Remnants} by their roots, but permanently eliminating them seems an impossible task.

\subsubsection{An Enchanting Export}
One of the climactic battles of the \textit{Calamity} was a clash between the stalwart The Dawnfather and the engine of callous destruction known as the The Chained Oblivion. This battle came to a dramatic and explosive conclusion in the \textit{Alabaster Sierras}. The unstable arcane forces unleashed in their battle created the valley that now houses the \textit{Parchwood}, thrusting the surrounding mountains upward and infusing them with incredible amounts of \textit{residuum}—the physical residue of arcane energy.

The sheer amount of leftover arcane power within the \textit{whitestone} of the mountains is incredibly receptive to enchantment. Magic items that incorporate at least an ounce of \textit{whitestone} into their construction require only one-quarter the creation time of other magic items. Some alchemists in \textit{Emon} have discovered that dissolving \textit{whitestone} with specific acids can leave behind pure \textit{residuum}, which can substitute for expensive spell components when such materials are not readily available.

If a spell has a component cost that consumes the component, its caster can substitute the required component for an amount of \textit{residuum} of an equal value. One pound of pure \textit{residuum} is worth approximately 500 gp in most markets, though it becomes extremely expensive and difficult to find outside of \textit{Emon} or the city of \textit{Whitestone} itself.

\textit{Residuum} can also be heated and blown into glass, forming a shimmering, greenish surface. Residuum slates or orbs can be used in stationary enchanting tables. Archmages willing to pay exorbitant rates—easily 20,000 gp or more—may acquire an Residuum enchanting slate that can be used to expedite the creation of all magic items. Using such a table allows anyone crafting a magic item to do so in one-quarter of the usual time.

\subsubsection{Government}
For many years, the de Rolo family ruled \textit{Whitestone} with dignity and justice. Their noble rule was broken by Delilah Briarwood, an exiled necromancer of the Cerberus Assembly from \textit{Wildemount}, and her \textbf{vampire} husband, Sylas. Together, they slaughtered the family, took over \textit{Whitestone}, and plunged the city into despair.

Nearly six years later, Vox Machina returned with Percival, the only known de Rolo heir, and discovered that his sister Cassandra also still lived. The heroes defeated the \textit{Briarwoods} and expunged their evil from \textit{Whitestone}, and the de Rolo siblings founded a new system of governance: the \textit{Chamber of Whitestone}.

This council consists of members appointed by the Guardian of Woven Stone, who is to remain a member of the de Rolo bloodline as long as it exists, or until the rest of the Chamber unanimously agrees there are no suitable de Rolos to accept the position.

\subsubsection{Defense}
\textit{Whitestone's} first line of defense is made up of the obelisk wards. Hastily constructed during the time of the \textit{Chroma Conclave} to avoid detection by their roving \textbf{wyvern} armies, five \textit{whitestone} obelisks mark the perimeter of the city and castle grounds. These obelisks form a dome of illusory magic when powered by an arcane spellcaster, hiding the city under the guise of continued, undeveloped \textit{Parchwood}, and shielding everything within from arcane detection. The obelisks haven't been used since the attack of the \textit{Chroma Conclave} and are beginning to be covered by graffiti and overgrowth.

For each spell slot that is syphoned into an obelisk, two hours of illusion is powered per spell level sacrificed, and the cover of the illusion barrier is incomplete unless all five obelisks are powered simultaneously. This siphoning is a taxing experience, and for every 30 spell levels sacrificed within a 24-hour period, the caster gains 1 level of exhaustion.

In addition to this rarely used magical ward, three groups of defensive personnel protect \textit{Whitestone} from both within and without.

\subparagraph{Pale Guard}
The people of \textit{Whitestone} rely upon the Pale Guard to resolve any (admittedly rare) crimes and disturbances of the peace. The Pale Guard are a force of sentinels who answer to the Pale Lord of Wardship. Since criminal activity is fairly low these days, most of the Pale Guard are given other civic tasks to help \textit{Whitestone's} people.

These tasks range from helping newcomers navigate the city streets to escorting visiting dignitaries to diplomatic functions. While performing a duty that requires it, members of the Pale Guard are permitted to carry simple, three-barrel pepperbox revolvers.

\subparagraph{Whitestone Rifle Corps}
The defense of the city from outside invaders and threats from the darkened boughs of the \textit{Parchwood Timberlands} is attended to by the \textit{Whitestone Rifle Corps}, who are trained in the deadly long-range rifles created by \textbf{Percival de Rolo}. They are only permitted to wield the weapons while patrolling the city walls or castle battlements—except in times of emergency, or when granted special permission by the \textit{Chamber of Whitestone}.

\subparagraph{Grey Hunters}
An elite force called the Grey Hunters makes up the most secret of \textit{Whitestone's} defenses. Trained in the arts of hand-to-hand combat, marksmanship, espionage, and etiquette, the Hunters have a reach that extends across all Tal'Dorei, but they exist only as a rumor in the eyes of the public—and they answer only to the \textit{Grand Mistress of the Grey Hunt}.

The firearms wielded by members of the Grey Hunters are custom-made and named by their owner. Their rifles' stocks are made of \textit{Parchwood} inlaid with \textit{whitestone} and enchanted against wear with \textit{residuum}. The lock mechanism bears the \textit{Whitestone} crest. Hunters are trained to meticulously care for their weapons, and many learn enough to make personal modifications, or even build one from scratch. A Hunter must never lose their weapon. If they cannot prove it was destroyed, they aren't permitted to return to \textit{Whitestone} until it is recovered—although they may request help from other Hunters in retrieving it. If they come across a weapon of \textit{Whitestone} in the wrong hands, they are honor-bound to return it to the city.

This covert organization is headquartered in the ruins of a failed city expansion deep beneath the eastern \textit{Parchwood}. This base is used for training, for weapon storage, and as a repository for their mission logs, so that a complete record of every Hunter and their weapons can forever be kept. As with many of his inventions, \textbf{Percival de Rolo} fears that the Hunters are a powerful—and even necessary—tool in the hands of the righteous, but a terrible weapon, should they be perverted to suit evil.

\subsubsection{Crime}
The community of \textit{Whitestone} has endured much hardship, from its shipwrecked beginnings to the bloody reign of the \textit{Briarwoods}. Its hardy people are generally more concerned with personal duty, small comforts, and the creation of new art and technologies than with the pursuit of easy profit. The \textit{Myriad} has tried several times to gain a foothold in \textit{Whitestone}, but they've been foiled time and again by the consistent counterespionage efforts of the \textit{Grey Hunters}. Nevertheless, they continue to try, and their presence has inspired petty theft throughout the city.

The altruistic Lord Percival and Lady Vex'ahlia of \textit{Whitestone} prefer to address crime at the source—poverty and suffering—rather than punishing criminals. Likewise, most petty criminals are seen in \textit{Whitestone} as worthy of rehabilitation, not punishment.

\subsubsection{Geography}
\textit{Whitestone} was originally constructed around the newly discovered \textit{Sun Tree}, and has slowly expanded outward ever since. \textit{Castle Whitestone} was built atop a small hill overlooking the city to the north. Over the years, farmers and homesteaders have cleared more and more of the \textit{Parchwood} to the south of the city. Even today, as \textit{Whitestone's} population has swelled, the \textit{Parchwood} remains defiant, dark, and threatening. It is in little danger of being overrun by mortals, and its eerie creatures strike back whenever the people of \textit{Whitestone} press against it too greedily. The \textit{Grey Hunt} was formed to keep the deadly denizens of the woods at bay, though now their purview reaches farther than just the \textit{Parchwood}.

Dawnfather Square is the central district of \textit{Whitestone}. The square surrounds a small hill, and the golden glow of the revered \textit{Sun Tree} that sits atop it can be seen throughout the city. Many of \textit{Whitestone's} original settlers already worshiped the The Dawnfather, and the discovery of the \textit{Sun Tree} was seen as a sign. Today, Dawnfather Square is filled with quaint, bustling shops and small town-homes.

\subsubsection{Points of Interest}
The City of \textit{Whitestone} has a number of famous landmarks. They are keyed to the map of \textit{Whitestone}.

\subparagraph{1. Castle Whitestone}
This gothic fortress was erected not long after the city's founding, and has long been the seat of power for the sovereign de Rolo family. Percival and Cassandra de Rolo have opted to end their absolute power over the city by sharing it with the \textit{Chamber of Whitestone}—on which they each hold a seat.

The castle's interior has a certain stark regality to it. While the hunting trophies and tapestries that decorate its walls are clear signs of the de Rolos' ancestral wealth, there is a solemn, beautiful hollowness to its corridors. Cassandra de Rolo says that something of her brother's melancholy seeped into its \textit{whitestone} walls, and the castle will always carry the weight of his sorrow—and of the \textit{Briarwoods'} evil.

The castle itself is generally off-limits to the public, but its grounds are more accessible. The \textit{Whitestone Rifle Corps} train in a firing yard behind the castle, and many people of the town flock to watch—only to be politely rebuffed by the \textit{Pale Guard}. Open to all, however, are the gardens—except for Percival de Rolo's latest project, the Widow's Garden. This walled garden was named for Lady Melanie de Rolo, who originally cultivated its deadly blooms to slowly murder her husband. The Widow's Garden used to be a macabre menagerie of poisonous plants from across \textit{Wildemount}. It fell fallow during the reign of the \textit{Briarwoods}, but Percival has resurrected it with Keyleth's horticultural aid.

\subparagraph{2. The Sun Tree}
After the The Chained Oblivion nearly consumed the Knowing Mistress during the \textit{Calamity}, a furious and vengeful The Dawnfather pursued the fleeing The Chained Oblivion across the mountaintops and over the sea, to the peak of \textit{Gatshadow}, where the mad god was eventually defeated and imprisoned once again. The land wept with pain as the two wrathful gods laid waste to Tal'Dorei.

After his victory, the The Dawnfather returned to \textit{Whitestone} and placed a single seed into the ground where the The Knowing Mentor was nearly destroyed. The god wept, and from the seed grew a symbol of protection—a tree with radiant orange foliage and a golden glow that surrounds it when in bloom. The Sun Tree became a beacon of light within the dark forest that enclosed it. While the extent of the tree's divine power remains a mystery, it stands as an emblem of unity and hope, both within the city of \textit{Whitestone}, and upon its official crest.

\subparagraph{3. Temples and the Greyfield}
A patch of land outside the eastern wall of the city sees generations of \textit{Whitestone's} dead interred within a graveyard known as the Greyfield. Custody of this cemetery is shared by the Temple of the Dawn and the Altar of the Raven, sanctuaries maintained by the clergy of the The Dawnfather and the The Matron of Ravens, respectively. Since most people in \textit{Whitestone} consider the The Dawnfather their patron deity, the Temple of the Dawn is the largest of these shrines. The smaller Altar of the Raven is right next door, and most well-to-do devotees of the The Dawnfather look upon the dour, dark-robed clerics of the The Matron of Ravens with a mixture of awe and unease.

The temples' grounds also house a number of mausoleums containing the remains of long-forgotten noble families, though the last generation of lost de Rolos were not recovered following the rule of the \textit{Briarwoods}. Traces of the \textit{Briarwoods'} defiling magic still linger within the Greyfield, and the \textit{Pale Guard} are posted here at night to quell the restless dead that occasionally claw their way to the surface.

\subparagraph{4. The Heart of Whitestone}
A majestic clock tower called the Heart of \textit{Whitestone} stands proud and tall over the city. It is Percival de Rolo's favorite invention, one he dedicated much of the past two decades to constructing to honor his loyal friends, their sacrifices, and the trials they endured together. The area around it is known as Forlington Square, named for one of the first families to settle in \textit{Whitestone}. This bustling courtyard is lined with shops and street performers that attract locals and tourists alike.

The clock tower's construction was completed five years ago, but Percival continues to tinker with it, adding both functionality and self-indulgent embellishments. The tower has individual movements for \textit{most major holidays}, including depictions of famous moments and persons from \textit{Whitestone} history. The clock tracks the movements of the moons, stars, and spheres. The mechanics of the clock are only truly understood by Percy himself, though small clubs of scholars and esotericists take great joy in unraveling its mysteries. Further, it has long been rumored that Percival secreted a vault within the tower's innards to safeguard his most dangerous inventions.

\subparagraph{5. Snowdrop Memorial Trail}
In the east of the city, winding between the \textit{Grey Hunt Estate} and the temple to the The Matron of Ravens, is a private trail through the woods. Only locals know of this trail, and even among them, few know that the path was made by Keyleth of \textit{Zephrah} on one of her many visits to the city. Countless snowdrop flowers line this secluded track, and they are tended to by Lady Vex'ahlia when she has time away from her duties to both the \textit{Chamber of Whitestone} and the \textit{Tal'Dorei Council}. The significance of this memorial garden isn't listed on any plaque, but most in \textit{Whitestone} have an inkling that it's dedicated to a beloved champion of the god of death.

\subparagraph{6. Grey Hunt Manor}
The estate of the \textit{Grand Mistress of the Grey Hunt} is rarely used by the current holder of that title these days, as all her nights not spent away on \textit{Tal'Dorei Council} business are passed with her family in \textit{Castle Whitestone}. Members of the \textit{Grey Hunt} still use it as a formal meeting place—though meetings of true import take place in their secret lair in the eastern ruins (11).

\subparagraph{7. House of the Lawbearer}
This temple is both a place of worship and a stately courthouse. All who enter are sworn to tell the truth by the vigilant clerics of the The Lawbearer—for the entire building is under the protection of a vast, permanent \textit{zone of truth} spell. This protection is fallible, but it sends a clear message to all who enter.

\subparagraph{8. Slayer's Cake Bakery}
Founded by members of Vox Machina, the Slayer's Cake has quickly become one of the most beloved bakeries in Tal'Dorei, and it has a satellite shop in \textit{Emon}. It ships bread and pastries to Nicodranas in \textit{Wildemount} via \textit{teleportation circle} weekly, and the bakery's owners are thinking of expanding.

\subparagraph{9. Lord Trinket's Public Park}
A section of the \textit{Parchwood} to the west of the city was cleared in the aftermath of The Whispered One defeat in anticipation of a major expansion. However, more of the \textit{Emonian} refugees left the city than anticipated, leaving these plans untenable, but the forest already cleared. At the request of Vex'ahlia, Grand Mistress of the \textit{Grey Hunt}, the space was transformed into a public park that she uses to play with her now-venerable brown bear, \textbf{Trinket}—and a lifelike marble statue of that same bear, clad in magnificent armor, stands in its center.

Portions of this city expansion did eventually come to pass, and now several manor houses have been built near the edge of the park to house visiting dignitaries.

\subparagraph{10. Galdric's Gate}
The new, northwestern gate of \textit{Whitestone} opens onto the Timberland Trail, which grants hunters easy access to the city when they return from expeditions into the \textit{Parchwood}.

\subparagraph{11. Eastern Ruins}
\textit{Whitestone's} first expansion in recent times happened just after the fall of The Whispered One, to accommodate the city's large refugee population. However, a mix of social unrest and monstrous incursions from the \textit{Parchwood} stymied the project, and the bones of the expansion were left to fall into ruin. No attempt has been made to reclaim the ruins, but the \textit{Grey Hunt} have made a secret base of operations in the basement of a magnificent temple not yet dedicated to any god.

\subparagraph{12. Hearth of the Everlight}
Established by \textbf{Pike Trickfoot} in the years following Vox Machina's victory over The Whispered One, this humble temple stands as a bastion of welcoming light and shelter to those seeking solace and direction. Within its stone walls, the faithful of the The Everlight study her philosophy of compassion and fierce protection. Still, the temple is far from tranquil—the congregation also hosts monthly fêtes and revels as a part of their worship in the image of their exuberant founder.

\subsubsection{Whitestone Adventures}
Game Masters who set their adventure in \textit{Whitestone} can use these plot hooks for inspiration.

\subparagraph{Dad Scientist (any level)}
\textbf{Percival de Rolo} doesn't usually involve outsiders in his creative process. However, when his youngest daughter Gwendolyn takes one of his latest creations and disappears into the dungeons of \textit{Castle Whitestone}, he has no choice but to rally all adventurers in \textit{Whitestone} to join him, Vex'ahlia, and the city's elite warriors in scouring the dungeons for his missing girl.

\subparagraph{Matter of Some Importance (low-level)}
The \textit{Castle Whitestone} library recently acquired a new book containing unprecedented knowledge about the Outer Planes, but opening the tome also opens an infernal portal that summons a steady stream of imps. The characters are hired by Head Librarian JB Trickfoot to keep the imps out of her hair—and out of the library, please—while she scours the tome for information.

\subparagraph{Grey Hunt (high-level)}
The characters have been recommended to the \textit{Grand Mistress of the Grey Hunt} as formidable adventurers by an unknown ally—or adversary. Vex'ahlia de Rolo considers them cautiously, but offers them the chance to prove themselves worthy of the title of \textit{Grey Hunter}. She sets them a task to be completed within the next three months: track down a legendary \textbf{ember roc} in the \textit{Cliffkeep Mountains} and claim its head as a trophy. Of course, she knows that this quest is more than just a hunt; the question is, will the characters realize it takes more than just skill to join the hunters?

\section{Dividing Plains}
Tal'Dorei's vast heartland, framed by the \textit{Stormcrest Mountains} in the southwest, the \textit{Cliffkeeps} in the north, and the \textit{Summit Peaks} in the south, contains miles and miles of rolling hills, tall-grass prairies, sky-blue rivers, and fertile farmlands. Within these plains, the hardworking people of the realm tend to their fields and businesses while adventurers and village militias battle back the wild beasts and wandering bandits of the region. The safest way to avoid being beset by monsters is to travel the \textit{Silvercut Roadway}, which slices through the endless plains, connects the coasts, and shepherds caravans between them. Of course, the roads are coursing with highwaymen and bandits, so which route is safer is entirely up to the traveler to decide. Cults of the \textit{Betrayer Gods} also prey upon the people of the plains. The most fearsome of these death cults are the \textit{Ravagers}, a bloody band of warriors exiled from cities and nomadic clans alike. All people of the plains fear the \textit{Ravagers}, and even the smallest villages train people old enough to hold a spear how to defend their homes. Beyond the grassy plains, owlbears and other monstrosities stalk the forests and foothills that dot the plains and border it on many sides. Packs of \textit{centaurs} also gallop across the open plains, delivering the The Arch Heart justice as they see fit.


\begin{description}
\item[Majority Faiths:]
The Wildmother

\item[Minority Faiths:]
The Stormlord
\item[Imports:]
Fish, gold, diamond, spell components

\item[Exports:]
Lumber, grain, produce, cobalt, iron, silver, livestock

\end{description}


\subsection{The Bramblewood}
Surrounding the base of the southernmost stretch of the \textit{Cliffkeep Mountains} is the Bramblewood, a dark and tangled forest named for its jagged, thorny-barked trees. The Bramblewood's impenetrable canopy and thick, gnarled roots give it an ominous appearance, but its eerie boughs are populated with countless deadly beasts. It's a grim warning to all who would dare approach \textit{Gatshadow}, the mountain of despair: come no closer if you value your life.

Yet despite this, without the abundant game and plentiful timber of the Bramblewood, the nearby city of \textit{Westruun} would likely not exist. A trade city with nothing to trade rarely lasts long, and the Bramblewood ensures that local carpenters, artisans, and tradesfolk always have plentiful resources. For those brave enough to venture into the Bramblewood, a pair of well-tread paths have resisted its overgrowth; one leads to the iron-rich Murdoon Mines at the base of the \textit{Cliffkeeps}, while the other winds upward to \textit{Gatshadow's} storm-wreathed peak.

\subsubsection{Bramblewood Adventures}
Game Masters who set their adventure in the \textit{Bramblewood} can use these plot hooks for inspiration.

\subparagraph{Haunted Road (mid-level)}
Local folklore speaks of a spectral young woman sometimes appearing along the road through the forest, sobbing and leading men deeper into the \textit{Bramblewood} to never be seen again. While these tales are generally considered childish ghost stories, a mining crew was recently found dead near the fabled spot, unwounded and with looks of terror stretched across their lifeless faces. A \textbf{banshee}, aided by the specters of those she has slain, hungers for the sweet taste of another adventurer's dying breaths.

\subparagraph{Blood of the Ruiner (high-level)}
The Shields, \textit{Westruun's} beleaguered soldiery, have tracked a cell of \textit{Ravagers} to a base in the \textit{Bramblewood}. They have conducted a foul ritual to the The Ruiner to turn their cult master, an elf named Kalydria Darkeye, into a \textbf{Ravager slaughter lord}. They have also transformed a number of the giant spiders endemic to the \textit{Bramblewood} into hideous demonfeed spiders through the same bloody ritual. The \textit{Ravagers} have created a makeshift fortification, and the Shields need to know what their next move is—an infiltration mission only powerful adventurers are capable of completing.

\subsection{Foramere Basin}
This great lake was once the ice fortress of \textit{Errevon the Rimelord}, the elemental behemoth that nearly toppled the fledgling realm of Tal'Dorei just two years after its founding. It was here that the combined forces of \textit{Kraghammer}, \textit{Syngorn}, Tal'Dorei, and \textit{the Ashari} hurled the Rimelord back into the Elemental Plane of Ice and brought an end to the \textit{Icelost Years}. The massive lake left behind by \textit{Errevon's'} melted citadel is now the largest source of fresh water in central Tal'Dorei.

Foramere Basin's shores are peppered with tiny huts and fishing communities that trade with \textit{Kymal} and \textit{Westruun}. The great lake is fed by three rivers that give life to the \textit{Dividing Plains}: the Byhills, the Tundrun, and the Wildpath. After especially hard winters, these rivers run hard and fast, and the overflows spill into the Foramere Waterway, which flows into Owlset Bay, near \textit{Stilben}.

Some of these fishing villages are home to the descendants of the fighters that warred against the Rimelord's tyranny. Even in times when \textit{Syngorn} and \textit{Kraghammer} were bitter rivals, the elves and dwarves here took no part in their homelands' hatred.

\subsubsection{Foramere Basin Adventures}
Game Masters who set their adventure in \textit{Foramere Basin} can use this plot hook for inspiration.

\subparagraph{Nightmare Before Winter's Crest (mid-level)}
The festival of Winter's Crest celebrates the anniversary of the Rimelord's defeat and the fall of Foramere Citadel. The characters, far from home on the holiday, are not only witness to the fishing communities on the banks of the basin gathering to celebrate Winter's Crest, but to the grim sight of Foramere's waters surging upward from the basin. The water has begun to freeze into towers of ice—someone or something is recreating the domain of \textit{Errevon the Rimelord} and will not stop until all of Tal'Dorei is buried in ice.

\subsection{Ironseat Ridge}
The gnolls and hillsfolk of the \textit{Dividing Plains} say the crooked, angular mountain that juts out of the plains and foothills was once the throne of a titan. From a distance, \textit{Ironseat Ridge} does resemble a crooked chair. Whether or not this legend is true hardly matters; the tales of \textit{Ironseat Ridge} permeate the folklore of the plains.

The mountains surrounding Ironseat itself have been hollowed out over the centuries, and other legends say that the ruins of an abandoned dwarven redoubt called Inverspire can be found by explorers savvy enough to navigate the labyrinthine tunnels—and that it's surely filled with monsters and vast riches. Legends aside, countless abandoned mineshafts do snake through the stone like the veins of a dead titan. When the winds grow strong, you can sometimes hear a terrible, growling moan from the air pressing through the empty channels. People from \textit{Kymal} rarely venture close enough to the mountains to hear the noise, but sometimes on clear, silent nights, the monstrous sound echoes through \textit{Kymal's} brightly lit streets.

\subsubsection{Ironseat Ridge Adventures}
Game Masters who set their adventure in \textit{Ironseat Ridge} can use these plot hooks for inspiration.

\subparagraph{The Beast Below (any level)}
It's one of those still nights where the howls of the Ironseat echo across the plains. The next time the characters spend a night in town, whether in \textit{Kymal} or any small village of the plains, they learn of disappearances—explorers that went into the caves, never to return—and of the hefty reward their grieving families are offering for their return. What could be making that awful wailing in the caves? Werewolves? Kidnappers using the echoing caverns?

\subparagraph{Dead Fires of Inverspire (epic-level)}
The mythic redoubt of Inverspire contains a legendary creation of the The All-Hammer: a \textbf{forge guardian}. It was placed here by the The All-Hammer during the \textit{Calamity} to watch over a forge with the power to create arms and armor fit only for a war between gods and mortals. Over twenty years ago, the elemental power of the \textit{Cinder King} somehow ignited this divine creation's ashen heart and filled it with elemental life. It has turned to the forge it was made to guard, and has begun creating beings of steel and flame that will tear Tal'Dorei apart stone by stone in order to find its missing father, the The All-Hammer. Only definitive proof of the The All-Hammer imprisonment beyond the Divine Gate can quench the guardian's anguished flames.

\subsection{Ivyheart Thicket}
Clustered at the western side of the \textit{Ironseat Ridge}, this thicket contains the ruins of an ancient civilization that remains under study by the \textit{Alabaster Lyceum}. Among scattered relics and the rubble of forgotten villages lie references to the \textit{Founding}: the first age of the world. The excavation pushes onward under the watch of the \textit{Lyceum}, with the promise of discovery looming, but very little understanding of what may be found beneath.

The western edges of these ruins are the most dangerous—and the most alluring—part of the Ivyheart Thicket. A vast, well decorated with ivy-crept marble pillars, burns with eternal fire. Explorers in the employ of the \textit{Alabaster Lyceum} say that this place has become a hotbed of activity. \textbf{Azer} have flocked to it from the Elemental Plane of Fire, and have made it their home, forging beautiful bronze weapons, armor, and \textit{objets d'art}. The \textit{Lyceum} reports that elemental activity within the nearby redoubt of Inverspire is somehow related to this elemental incursion—and that it may be connected to the mythic fire giant stronghold of \textit{Vulkanon}.

\subsection{Kymal}

\begin{description}
\item[Small City:]
Population 29,440 (63\% humans, 24\% dwarves, 13\% other races)

\end{description}

\noindent Glitzy. Gilded. Grimy. The city of Kymal stands as the grumbling ghost of a gold-rush boomtown. These days, it's a haven for gamblers and criminals: a city of gold built atop a mountain of mud.

Over two centuries before, the discovery of rich gold veins within the \textit{Ironseat Ridge} sent a rush of eager folk to its base. Kymal began life as their base camp, and over the next century it grew into a lively, booming city. gold and minerals came out of Kymal, and all the pleasures of home followed. Music, theater, sex, food, liquor, drugs, trinkets—if you could want it, you could find it in Kymal. Anything you could spend money on was there, and people spent money like the world was ending. The rich binged anything they could get their hands on to pretend they were important, and poor folk indulged in whatever comfort they could find from their wretched existence in a town where money meant everything.

In time, of course, the veins ran dry. Ten years after Kymal boomed into being, it went bust. The miners packed up and left. Not many of them had struck it rich; in fact, most folk were worse off than they had been before. In the intervening centuries, Kymal has gone through too many cycles of boom and bust to count. New veins of gold or mithral or adamantine have been found in the mountains, and people rush in to claim it. Or a cunning entrepreneur comes up with a scheme to put Kymal on the map—and it works, for a few years, before the con runs its course and people are right back to where they were before.

Most folk here make a modest living in rural villages surrounding Kymal's city center, or trawling the nearby \textit{Foramere Basin} for fish, and are trying to build interest as a trade center like \textit{Westruun}, though without much luck. Kymal's most recent change of fortune, however, has yet to run dry. Some seventy years ago, an entrepreneur from \textit{Ank'Harel} named Calis Krishtan opened a small casino called "The Maiden's Wish." It started something of a gold rush of its own, one that continues to this day.

The city has grown over the past seventy years, becoming a hub of entertainment and games of chance in Tal'Dorei. So great was Calis's revival of Kymal that the \textit{Tal'Dorei Council} appointed his son (and heir to the Maiden's Wish), Jaktur Krishtan, the Margrave of Kymal.

\subsubsection{Government}
Now entering old age, Margrave Jaktur Krishtan is desperately trying to maintain the boom that his father created, while grooming his son to take over his gambling empire. Krishtan has recently made headlines across Tal'Dorei by relinquishing his position as margrave. It's an open secret that Jaktur Krishtan has never had a head for governance, simply enjoying the lavish, gold-plated life his casinos' profits bring him. He could easily have allowed his advisors to continue running \textit{Kymal} for him while he continued to throw magnificent galas, make nightly appearances at the Maiden's Wish, and put his stamp of approval on glitzy parades down the city's main street.

Investigators have flocked to \textit{Kymal} to uncover the truth of Margrave Krishtan's sudden change of heart. All suspect that the truth is a salacious secret, one which would shake \textit{Kymal} to its foundations. Krishtan himself seems unperturbed. He continues to live large as he always has, his detractors don't disparage him any more than usual, and the public that loves him hasn't changed their tune.

Jaktur's former majordomo, an austere tiefling woman named Khavasta "Prudence" Almiji, has taken up the mantle of Acting Margrave with quiet satisfaction. It is, after all, the title that she should have had for the past thirty years, and she hopes that the \textit{Tal'Dorei Council} will see fit to appoint her margrave permanently in very short order.

\subsubsection{Defense}
The acting margrave commands a force of about one thousand armed guards called the Public Defense. Despite the name, they are most often seen stationed around the casinos owned by ex-Margrave Krishtan, or working as bodyguards for \textit{Kymal's} politicians and wealthy cardsharps. Gold can buy anything in \textit{Kymal}, even the personal protection of the Public Defense.

\subsubsection{Crime}
\textit{Kymal} is no longer the den of crime it once was, but there is still scum beneath its glitzy and imperfect façade. Small-time cutpurses lurk in the shadows, waiting to make off with the small fortunes won by the casinos' luckiest patrons. The war between the \textit{Clasp} and the \textit{Myriad} smolders within \textit{Kymal} as well, but Acting Margrave Almiji has played her cards carefully and intelligently. Through an intricate web of bribes, promises, and threats, she has managed to suppress even the faintest trace of gang violence within her city.

Of course, crime still occurs in \textit{Kymal}. It's just not public and bloody. The \textit{Clasp} and \textit{Myriad} both own casinos in town—named the Wishing Well and the Dragon's Hoard, respectively. Not many know of these casinos' less-than-savory proprietorship. Even those who do fail to realize just how deeply rigged its tables are—or how many pickpockets and swindlers hide in plain sight in their gilded halls.

Moreover, the \textit{Clasp} and \textit{Myriad} alike have immense illicit drug operations within \textit{Kymal}. The \textit{Public Defense} is technically bound to arrest anyone caught selling or possessing these substances, but the bulk of these arrests fall on small-time crooks working for these criminal organizations, almost never the high-ranking members.

\subsubsection{Geography}
\textit{Kymal} is built on the rocky foothills of the \textit{Ironseat Ridge}. Exploring beyond the glamorous city leads one to nothing but a desolate expanse of boulders and scrub-brush.

\subsection{Ruins of Torthil}
While traveling south of \textit{Kymal} along the \textit{Wildwood Byway}, just outside of the Ivyheart Thicket, travelers can spot a mile-wide cluster of broken stone, salted earth, and a single obsidian monument to the death of Warren \textit{Drassig}, one of the most infamous and terrible figures in Tal'Dorei history. This scarred patch of land marks the ruins of the village of Torthil, which was the site of a pivotal battle of the \textit{Scattered War}—one in which King Warren \textit{Drassig} was ambushed and slain by a rebel army.

The city was destroyed in an act of vengeance by Warren's eldest son Neminar over two centuries ago. Now, its ruins harbor only memories—and occasionally outlaws and cultists who use the ghost stories that surround it to avoid suspicion. The now long-weathered wreckage and broken foundations of the town harbor only ghosts and unsavory creatures drawn to places of such dark history.

\subsubsection{Ruins of Torthil Adventures}
Game Masters who set their adventure in the Ruins of \textit{Torthil} can use this plot hook for inspiration.

\subparagraph{Forgotten War Machine (high-level)}
The Ruins of \textit{Torthil} hold one lost secret. Buried beneath a landslide in one of the nearby hills is a \textbf{platinum golem}. Any crests that may have indicated who it fought for—Tal'Dorei's rebels, the elves of \textit{Syngorn}, or \textit{Drassig's} legions—have been scoured away by time. The characters are drawn to \textit{Torthil} when exiles and criminals living in the ruins flood into \textit{Kymal}, saying that a greedy treasure hunter awakened a metallic monster that laid waste to the already ruined town—and is now making its way steadily toward \textit{Kymal}!

\subsection{Shadebarrow}
No one who has visited the \textit{Shadebarrow} has ever returned alive—but neither do they come back entirely dead. The \textit{Shadebarrow} was once a monolithic henge used as a ritual site for sun-worshiping druids known as the Dawn Circle, but their shrine became their tomb after \textit{Trist Drassig} slew them for aiding \textit{Zan Tal'Dorei} in the \textit{Scattered War}, centuries ago.

After \textit{Trist Drassig's} death, the trade guilds of \textit{Westruun} laid claim to the unowned land, and from there, the Dawn Circle's abandoned and treasure-less burial tunnels were purchased from the unions at great cost by an eccentric \textit{Westruunian} baron named Sevil Howthess. His obsession with its history led him to be interred there upon his passing, the tomb outfitted with protections and, supposedly, the remnants of his fortune.

The exact location of the \textit{Shadebarrow} has been forgotten by all but the Baron Howthess's few surviving grandchildren, and they deny any involvement with the forgotten crypt. The surviving descendants of Sevil Howthess, Camilla Tenver and Buddleia Austan, live in \textit{Westruun}, but have changed their names through marriage. Only searching through the public records in \textit{Westruun's} \textit{Hall of Reason} or seeking lore about the \textit{Shadebarrow} in the \textit{Cobalt Reserve} will turn up their names.

Despite the barrow's location being lost to time, treasure hunters and historians still seek out—and occasionally find—the infamous crypt. If they ever return to civilization, it is as a wailing spirit, cursed to eternally torment whoever they thought of in their final moments, or as a shambling corpse focused only on murdering their loved ones.

\subsubsection{Shadebarrow Adventures}
The \textit{Shadebarrow} is a crypt that you can turn into a dungeon of any size in your campaign. It may have only a single layer filled with ghosts, and perhaps a plot-important item. Or it could be a multi-level dungeon with interconnected passages and a mighty reward at the bottom, like a \textit{Vestige of Divergence} guarded by a terrible monster.

Game Masters who set their adventure in the \textit{Shadebarrow} can also use these plot hooks for inspiration.

\subparagraph{Forsaken Monolith (low-level)}
While traveling across the \textit{Dividing Plains}, the characters stumble upon the \textit{Shadebarrow}. The ancient druidic henges stand regally atop the grassy hill, and all seems peaceful—until the specters of dead explorers rise from the ground. They wail an incessant curse as they try to kill or scare intruders away: "Defiler! Howthess!"

\subparagraph{Undeath Unbound (epic-level)}
The Demon Prince of Undeath is one of the foul princes of the Abyss and a mighty servant of the \textit{Betrayer Gods}. The cruelty shown at the Dawn Circle so long ago drew his loathsome gaze, and for centuries he has bided his time, waiting for a mortal to disturb this nest of hate and undeath. Ever since Baron Howthess disturbed the \textit{Shadebarrow's} seal, this prince of evil has rallied an army of undeath in the accursed depths of this crypt. Now, the time has come for him to spill his plague of undeath across the \textit{Dividing Plains}. The undead are spreading to all corners of Tal'Dorei, and the characters must find the \textit{Shadebarrow} and reach its deepest sanctum to halt their spread for good.

\subsection{Silvercut Roadway}
Named for the shimmer of endless fields of windswept, sunlit grasses that accompany the way, the \textit{Silvercut Roadway} connects the heart of Tal'Dorei—\textit{Emon}—to its distant extremities, affording traders and travelers some measure of safety against the wilds. Parts of this ancient road predate the \textit{Tal'Dorei Council}, and it now runs from the \textit{Bladeshimmer Shoreline} to the \textit{Lucidian Coast}, connecting most major hubs of civilization across the country.

\subsubsection{Silvercut Crossroads}
The \textit{Silvercut Roadway} and the \textit{Wildwood Byway} meet in the shadow of the \textit{Ironseat Ridge}. During the rule of \textit{Drassig}, many innocents and war heroes were executed by hanging at a massive gallows that once stood at the crossroads, with bodies hung from signs as a warning to dissenters. The gallows were destroyed by \textit{Zan Tal'Dorei's} rebel army, and the rebels created a stone monument with an inscription to commemorate all who lost their lives to tyranny:

\DndQuote%
{}
{\enquote{The standard-bearers of revolution honor the innocent dead. Vengeance will not restore the lives of the lost, nor will it appease their spirits. We vow you will receive justice, for justice will ensure that no more guiltless blood shall be shed.}}
{}
Most travelers choose not to camp too closely to the site, for there are countless urban legends about hauntings that surround this grim monument.

\subsubsection{Silvercut Roadway Adventures}
Game Masters who set their adventure on the \textit{Silvercut Roadway} can use these plot hooks for inspiration.

\subparagraph{The Road Warriors (low-level)}
While traveling across the Silvercut, the characters come across the smoking wreckage of a massive wagon train, and investigation shows that it has been ransacked after a fierce battle. A team of mercenaries called the Road Warriors spring upon the characters, blaming them for the attack. Investigating the wreckage turns up unholy symbols of the The Ruiner, suggesting that the \textit{Ravagers} were responsible for the massacre, allowing the characters to easily clear their names.

But is it really that simple? Further investigation reveals that this convoy had a number of mercenary mages in their midst, and many of them had contracts from the \textit{League of Miracles}.

\subparagraph{Shard of the Shadebarrow (mid-level)}
A \textbf{wraith} calling herself the spirit of Archdruid Avandros has animated an army of ghouls and ghasts from the graveyard. The undead have taken over the \textit{Silvercut Crossroads}, and they must be dealt with before a national trade crisis emerges. But what is the spirit of an archdruid of the Dawn Circle doing in this mass grave? Sneaking into the graveyard reveals a human skull with sinister amethyst growths protruding from the bone, and Avandros's spirit is desperate to protect it.

\subsection{Throne of the Arch Heart}
In the foothills of the \textit{Ironseat Ridge} is a small rise known as the Throne of the The Arch Heart: a gathering place for the many nomadic \textit{centaurs} who ride across the \textit{Dividing Plains}. The \textit{centaurs} were first born when the The Arch Heart showed mercy to a tribe of elven horse riders who gave their lives serving their god during the \textit{Calamity}. For centuries, the \textit{centaurs} have all returned to the Throne of the The Arch Heart once per decade to gather in a Herdsmeet to discuss the state of the herds, and to pay homage to their deity on the very hill where the first \textit{centaurs} were created. This event causes great chaos in the eastern plains; the massive migration of the centaur herds disrupts trade and stirs the \textit{Ravagers} to acts of terrible vengeance against the creations of the The Arch Heart—the The Ruiner divine nemesis.

\subsubsection{Throne of the Arch Heart Adventures}
Game Masters who set their adventure in the Throne of the Arch Heart can use this plot hook for inspiration.

\subparagraph{Fist of the Ruiner (mid-level)}
Many orcs call this rise the Fist of the The Ruiner, for they too were created here in the image of their god. Legend tells that orcs and \textit{centaurs} were both created at this rise in the aftermath of a battle between the The Arch Heart and the The Ruiner. orcs with a deep appreciation for ancient myths from across Tal'Dorei make a pilgrimage at the time of the Herdsmeet. Some are long-time friends of the \textit{centaurs}, and they pay homage to the The Arch Heart together in thanks for spilling the blood that created them. Other orcs join with the \textit{Ravagers} to lay siege to the Herdsmeet—and the \textit{Ravagers'} bloodthirsty power is so great that orcs and \textit{centaurs} alike have begun hiring humanoid mercenaries to defend the herds and pilgrims as they attempt to worship the The Arch Heart in peace.

\subsection{Torian Forest}
Cradling the southwestern edge of the Cliffkeep Mountain range, the Torian Forest supplies much of \textit{Emon's} timber, as well as \textit{Kraghammer's} firewood. Though its locals proudly proclaim the Torian to be the "most welcoming" and "least haunted" of Tal'Dorei's forests, this is more of a myth than people realize. The forest is welcoming and bright at its edges, with a beautiful wetland at its northern end, but within its deep and unexplored heart are spirits not understood since the \textit{Age of Arcanum}—and a cruel and hungry \textbf{wraithroot tree} guarding a magical artifact likewise not seen since time immemorial.

\subsubsection{Rivermaw Herd}
Born from the wild peoples of the \textit{Dividing Plains}, free folk who shunned the constraints of society, the Rivermaw is the greatest of the \textit{Dividing Plains'} herds. It now boasts a membership of nearly ten thousand, along with a number of satellite herds that owe their allegiance to the Rivermaw.

Though the Rivermaw is now a mighty force for good on the \textit{Dividing Plains}, they had humble beginnings. Just two generations ago, they numbered little more than three hundred half-giants, orcs, gnolls, humans, and nomads of other races that made their home in the Torian Forest.

It wasn't until they were conquered by the Herd of Storms, helmed by the tyrannical goliath Kevdak, that a more violent way of life was forced upon them for a number of years. Once Kevdak was slain, the remaining members of both the Rivermaw and the Herd of Storms recanted his brutal ways and become the \textit{Rivermaw herd} known across Tal'Dorei today. It is led by two \textit{goliaths}: Zanroar, son of Kevdak, and Skyhaal, daughter of Zanroar.

\subsection{Turst Fields}

\begin{description}
\item[Village:]
Population 2,420 (40\% humans, 38\% halflings, 15\% gnolls, 8\% other races)

\end{description}

\noindent To the northwest of \textit{Drynna}, not far from the Dawnmist Pines, the farming community of Turst Fields toils over extremely fertile land fed by the \textit{Mooren River Run}. Providing much of Tal'Dorei and \textit{Whitestone's} grown food and produce, the community must import most all other commodities outside of wood, making it a lucrative stop for trade caravans traveling along the Parchwood Way to \textit{Drynna}—and an easy target for swindlers. As such, nearly two hundred Shields from \textit{Westruun} are permanently stationed here to keep the peace and prevent unlawful activity. Nearly as common as the arrival of a trader's wagon is an armored carriage transporting brigands, grifters, and horse thieves back to \textit{Westruun} for trial.

\subsubsection{Gnoll Mercy}
\textit{Turst Fields} has a significant \textbf{gnoll} population. As mutant beings created by demonic rituals, gnolls are shunned as monsters in most humanoid settlements on the \textit{Dividing Plains}. Nevertheless, without the Dustpaw gnolls, Turst would not have survived its first summer. When the village was founded, it was easy prey for the \textit{Ravagers}. Their leader sought blood tribute from the people of Turst, but the villagers refused to offer a sacrifice. This cell of \textit{Ravagers} had recently sought the favor of the Demon Lord of Slaughter, transforming a number of the plains-dwelling hyenas into gnolls in that pursuit.

When the villagers of Turst refused to give up their own, the \textit{Ravagers} unleashed the gnolls upon them. But something unusual happened. After the gnolls first drew blood, one of them began to weep. Then another, until all of the gnolls were howling in sorrow and fury. These beings, though infused with demon ichor, turned against their creators to defend their helpless prey.

To this day, the Dustpaw gnolls have lingered in Turst, slowly becoming trusted members of their community. Their unsettling fangs and canine faces upset many who travel to this village, but any who deride the gnolls do so at their own risk. Not at the risk of mauling from a \textbf{gnoll}—though any humanoid being will fight back if pushed too far—but because hateful agitators will soon find themselves surrounded by the spears of the Turst militia.

\subsubsection{Turst Fields Adventures}
Game Masters who set their adventure in \textit{Turst Fields} can use this plot hook for inspiration.

\subparagraph{Herd Mentality (mid-level)}
A half-elf, half-\textit{dragonblood} \textbf{Ravager slaughter lord} calling herself Grud the Great has taken control of an arm of the \textit{Ravagers}. Her highly organized regiment worships the The Strife Emperor. Now, they march on Turst to take an easy foothold in the northern plains. The army is too large to be stopped by just the characters, but a Dustpaw \textbf{gnoll} warrior named Akaksa knows of a number of nomadic groups wandering nearby that could help defend the town, if properly persuaded. The characters can seek out the Dawnmist Herd, led by a goliath \textbf{Rivermaw stormborn} named Gardessk, and convince him to aid the people of Turst.

\subsection{Westruun}

\begin{description}
\item[City:]
Population 51,980 (61\% humans, 10\% gnomes, 7\% tieflings, 5\% orcs, 4\% half-giants, 1\% gnolls, 12\% other races)

\end{description}

\noindent The city of Westruun is the center of Tal'Dorei, a crossroads of culture along the \textit{Silvercut Roadway} that welcomes adventurers, traders, and vagabonds of all kinds. Despite its name, Westruun is more east than west, and was in fact named after Palest Westruun, the trader and philosopher who founded the city. Westruun suffered under the rule of the goliath tyrant Kevdak's Herd of Storms. Worse, its outlying towns and great stone buildings were ravaged by the acidic breath of Umbrasyl the Hope Devourer, an \textbf{ancient black dragon} of the \textit{Chroma Conclave} who exacted tribute from his lair in the eerie mountain of \textit{Gatshadow}. Despite the damage inflicted by its evil rulers, Westruun was one of the first cities in Tal'Dorei to recover from the \textit{Conclave}. Though its government is deeply in debt to both the \textit{League of Miracles} and a number of adventuring parties, the city has long since reclaimed its place of importance as an economic, academic, and political power within Tal'Dorei.

\subsubsection{Government}
\textit{Westruun} underwent a political revolution in the wake of the \textit{Chroma Conclave's} defeat. The former margrave, Brandon Zimmerset, was a tyrant who claimed absolute power for himself following Kevdak's defeat at the hands of Vox Machina. The title of margrave is a holdover from the authoritarian rule of Warren \textit{Drassig} centuries ago, and those who hold that office are appointed for life by the council. A group of local idealists called the Retrievers banded with adventurers and mercenaries to overthrow Zimmerset a decade after he took power, and instated their leader as the lord mayor instead. Elections for the position of Lord Mayor of \textit{Westruun} are held every five years, with no term limits. All adult citizens of the city who have lived there for at least five years have a vote, though only a fraction of the people actually use it.

The Retrievers continue to be a political force in \textit{Westruun}, and several of \textit{Westruun's} delegates to the \textit{Tal'Dorei Council} are members of this faction. They are campaigning to abolish the office of margrave from all settlements in Tal'Dorei. Though they have broad support in the council, the council doesn't have the authority to unilaterally demand its constituent city-states to change their system of government. The revolution-minded Retrievers are somewhat stifled by the need to compromise with other cities.

Lord Mayor Lysandra Kallos is a middle-aged goliath and former leader of the Retrievers. She has never lost a vote for re-election, and holds executive power over the city's defenses, along with the ability to dissolve any business or guild, and to exile any person or group of up to ten people from the city at her discretion. She rules \textit{Westruun} with the aid of a council of ten advisors.

\subsubsection{Defense}
The lord mayor commands a city guard of about 1,200 trained soldiers known as the Shields of the Plain, though most folk just call them the Shields. \textit{Westruun's} soldiery is spread thin across the walls of the city itself, stationed at the gates of the dozens of outlying settlements and conducting raids against small cells of \textit{Ravagers} across the plains.

Many people in town would like to gather their forces and put an end to the \textit{Ravagers'} death cult once and for all, but the Shields are spread too thin to mount an effective offensive. Even in their current state, dozens—if not hundreds—of mercenaries, bounty hunters, and adventurers pass through \textit{Westruun} every year in search of bloody work that the Shields are too beleaguered to handle themselves.

\subsubsection{Debt}
Before he was ousted, Margrave Brandon Zimmerset authorized hundreds of reconstruction bounties—offers of gold to anyone who could rebuild the immense damage the city suffered during Umbrasyl's attack and the Herd of Storms' occupation of the city. Bounties were placed on the city's siege-breaking stone walls, majestic towers and parapets, and even its iconic post-and-beam houses, many of which were razed completely during the attack.

The margrave expected these bounties to be slowly completed over decades, but they were all snapped up and completed in a few short years by mercenary mages serving a group called the \textit{League of Miracles}. The margrave was saddled with a crushing debt, one that he was both unwilling and unable to pay before being ousted by the Retrievers. As his successor, Lord Mayor Lysandra Kallos inherited Zimmerset's immense debt to the \textit{League}, and is constantly seeking to find ways to alleviate their pressure—especially as the league continues to use the debt as leverage to buy the allegiance of her advisory council.

\subsubsection{Crime}
The \textit{Clasp} has become a major and public force in \textit{Westruun} since the criminal organization's dubiously virtuous makeover following the \textit{Cinder King's} defeat in \textit{Emon}. Many \textit{Clasp} agents enjoy the reputation of folk heroes across the \textit{Dividing Plains}, and make public appearances in taverns and town squares throughout \textit{Westruun}.

Nevertheless, the fact remains that the \textit{Clasp} is still a criminal organization. They have controlled and efficiently run their shadowy trades for decades right under the Shields' noses. Their network of fences runs from coast to coast, \textit{Stilben} to \textit{Emon}. Their hideout changes every few years, when the Shields manage to track them down. Their current hideout is secreted away behind the storefront of the Saddled Plainscow, a popular tavern owned and operated by \textit{Clasp} agents.

Folk in search of real folk heroes often seek the covert agents of the \textit{Golden Grin}, who tend to wait and watch in local taverns. The Grin sometimes gets a bad reputation in \textit{Westruun}, since people need help every day, but the Grinners prefer to wait for just the right moment to reveal themselves and strike against evil.

\subsubsection{Geography}
Miles of farmland and fields surround \textit{Westruun's} towering stone walls—except to the west, where the city is bordered by the \textit{Bramblewood} and dwarfed by the mist-shrouded peak of \textit{Gatshadow}. Nearly two dozen villages and towns surround the city's walls.

The city's interior is divided into six wards:

\subparagraph{Scholar Ward}
The city's small southern ward is replete with major temples to the \textit{Prime Deities} and homes for their clergy. The most important gods of the \textit{Dividing Plains} are the The Lawbearer, the The Wildmother, and the The Platinum Dragon, and it's made obvious by their huge, resplendent sanctuaries in the Temple Ward. The greatest temples here are the First Bastion, the Wild Arbor, and the The Platinum Dragon Rest.

\subparagraph{Opal Ward}
\textit{Westruun's} central district is also its municipal heart. The ward has no residential housing, but is home to the Lord Mayor's Residence, the \textit{Hall of Reason}, and the public square.

\subparagraph{Residential Ward}
The northeastern section of \textit{Westruun} is the city's most densely populated ward, housing over two-thirds of its population. There are no major landmarks in this ward, though every neighborhood in the ward has its own unique culture. Even though tourists might not visit the landmarks here, there are countless local bakeries, shops, and pubs beloved by its residents.

This ward is rife with crime. \textit{Westruun's} poorest citizens are pressed together in tight quarters, sometimes as little as a single city block away from the city's wealthiest manor houses. The \textit{Clasp} and the \textit{Myriad} prey upon the poor and affluent alike, often stoking the tensions affecting the city's downtrodden to create chaos that masks their crimes. The Shields are stretched too thin to properly deal with these agitators. Unlike in the Temple Ward, free expression of religion is not a guaranteed right here, even within the privacy of one's own home. Though the The Lawbearer, the The Wildmother, and the The Platinum Dragon are all benevolent deities, many of their \textit{Westruunian} followers are not particularly tolerant of other faiths.

\subparagraph{Market Ward}
\textit{Westruun's} southwestern ward contains the majority of its businesses, trade stands, and production warehouses. Some of these businesses are historic fixtures of the city, rebuilt by the \textit{League of Miracles} in years past, but many sell overpriced wares to cater to tourists, adventurers, and the city's wealthy elite.

It's currently in fashion for store owners to live in a second story above their shops, though less scrupulous businesspeople have turned their residences into gambling dens, black markets, and illegal brothels. The most significant building in this ward is the Exandrian Exchange, a marble-walled auction house where hundreds of vendors hawk rare and unusual items from across Tal'Dorei. A satellite location of \textit{Gilmore's Glorious Goods}, a famed magic item shop in \textit{Emon}, is a beloved attraction of this district.

\subparagraph{Underwalk Ward}
The city's newest ward is underground. It's an extension of the residential sprawl built from repurposed sections of the sewers, underground military bunkers, and the wine cellars of mansions annihilated by Umbrasyl and the Herd of Storms.

The Underwalk was frantically excavated during the invasion of the \textit{Chroma Conclave} in hopes of keeping the populace hidden. It has been largely abandoned for the past twenty-four years, save for those too poor to return even to the Residential Ward, but an effort has been made in the past two years to create a unique and vibrant ecosystem in the Underwalk. The optimistic entrepreneurs trying to popularize Underwalk housing hire adventurers in great quantities, for every day, new and unsettling rumors of many-headed monstrosities prowling the Underwalk's shadowed roads spread to the surface. The best-known entrance to this tangle of tunnels is the Underwalk Gates, a forty-foot-wide staircase in the Residential Ward that descends into a gaping opening lit only by torchlight. Another entrance known mostly to the \textit{Clasp} is a single unassuming manhole in the Opal Ward, topped with a metal cover emblazoned with the crest of \textit{Westruun}.

\subsubsection{Points of Interest}
The City of \textit{Westruun} has a number of noteworthy landmarks. They are keyed to the map of \textit{Westruun}.

\subparagraph{1. Cobalt Reserve}
This vast repository looks more like a house of worship than a library. Nevertheless, this landmark of the Scholar Ward is Tal'Dorei's largest house of knowledge, and it is maintained and fiercely defended by the monks of the \textit{Library of the Cobalt Soul}.

A character who spends 1 hour in the public section of the Cobalt Reserve can make a DC 15 Intelligence (History) check, discovering any piece of common knowledge in Tal'Dorei. On a failure, this check can be repeated after another hour of study. The GM can decide that the knowledge the character seeks can only be found in the restricted archives, which only members of the \textit{Cobalt Soul}, or characters in the \textit{Cobalt Soul's} good graces, can access.

\subparagraph{2. Westhall Academy}
No university in Tal'Dorei can hope to compete with the prestige of \textit{Emon's} \textit{Alabaster Lyceum}, but the Westhall Academy does its best. Its sand-colored stone walls and copper-patina dome are an unmistakable landmark of the Scholar Ward. They imply a stately and refined institution with an unpretentious, workaday attitude. Most of the Westhall Academy's students come from humble backgrounds, and the lord mayor does her level best to ensure that the academy is funded enough to let anyone from the plains who shows great promise attend, free of cost. Mages, architects, poets, artists, agriculturalists, and people of all disciplines can find an education at Westhall.

\subparagraph{3. Yuminor Observatory}
\textit{Westruun} is a fairly flat city, and one tower on the edge of the Scholar Ward looms high above the nearby \textit{Westhall Academy in Westruun}. It looks forever skyward, the arcanists inside constantly charting the movements of the heavens. A quasi-religious order of diviners and astromancers named the Scions of Yuminor study here, supported by the academy's headmistress, Estella Ladimar, who is herself a blood descendant of the heroic wizard Atz Yuminor. The observatory is open to all, and the diviners here will gladly give a star reading for 10 gp, once per day. Receiving this reading allows you to roll 1d20. Note the result—you can exchange it for any d20 roll made by a creature you can see. This divination fades after 24 hours. Despite the Scions of Yuminor's affability, something seems odd about them, like they have stared into the stars just a bit too long and seen something they wish they could forget.

\subparagraph{4. First Bastion}
The most magnificent of the temples within \textit{Westruun's} Temple Ward is the First Bastion. Dedicated to the The Lawbearer, this opulent sanctuary is a common attraction for visiting folk of all faiths. Its magnificent, even decadent, gold-and-marble exterior is a sight most people can't imagine exists upon the Material Plane—it seems like nothing less than a building from the realm of the gods themselves. The First Bastion was a secret home for refugees during Kevdak's occupation, and many \textit{Westruun} survivors have turned to the faith of the The Lawbearer after being protected by her grace.

\subparagraph{5. Lord Mayor's Residence}
The Lord Mayor's Residence is a squat manor in the Opal Ward. It was originally named the Margrave's Keep, and as such, it is a defensible and highly fortified compound buttressed against military barracks. The interior has been livened up significantly by the lord mayor, but the exterior is still a spare, military edifice.

\subparagraph{6. Hall of Reason}
This beautiful courthouse in the Opal Ward is constructed in sweeping Lyrengornian style and bears a marble statue of the The Platinum Dragon above its steps. The square before the courtroom steps is a forum for public gatherings and bears a fountain depicting Palest Westruun, the founder of the city, atop an onyx horse.

\subparagraph{7. The Black King}
Palest Westruun was not a king, but his new statue in the Opal Ward is often called the \textbf{Black King}, a subtle jab at the city's previous "Black King," the tyrannical dragon Umbrasyl. The high priest of the First Bastion blessed both the \textbf{Black King} and \textit{Westruun} itself when the statue was unveiled on the anniversary of \textit{Westruun's} reclamation, saying, "May the The Lawbearer grant us all strength of spirit when chaos next threatens our peaceful civilization. When that time comes, may the \textbf{Black King} unsheathe his sword to rally our people and lead \textit{Westruun} to victory."

When the safety of all of \textit{Westruun} is threatened, the \textbf{Black King} will animate and defend its people, becoming a stone golem with a speed of 60 feet, Charisma and Intelligence scores of 18 (+4), and the ability to speak Common. Its personality is dour but determined—regal and single-minded in its desire to protect its city and its people.

\subparagraph{8. Survivors' Legacy}
The names of all those in \textit{Westruun} who survived the reign of Kevdak and Umbrasyl are carved on a wall in the first chamber of the Underwalk. Legend has it that the memorial was created by \textbf{Percival de Rolo} when the people of \textit{Westruun} first hid here from the wrath of the \textit{Chroma Conclave} after Umbrasyl's death, but it has been expanded in the years since, now including the names of all who perished in the struggle as well.

The Survivors' Legacy is a symbol of hope, warding away the sorrow and evil of the Underwalk beneath, and to this day common, folk make trips here to honor their lost loved ones. Some of the more generous folk give aid to the homeless living here as well, and some more adventurous travelers make prayers to the The Dawnfather and the The Matron of Ravens here before delving into the Underwalk's deeper reaches.

Once per week, a character related to one of the people named on the Survivors' Legacy can pray there and gain Advantage and Disadvantage on death saving throws for the next 24 hours.

\subsubsection{Westruun Adventures}
Game Masters who set their adventure in \textit{Westruun} can use these plot hooks for inspiration.

\subparagraph{They Can't Have Flown Away (low-level)}
Over the past week, over four dozen plainscows, exceptional livestock and beasts of burden, have disappeared from \textit{Westruun's} outlying villages. Left in their place were nothing but a series of five-foot holes in the ground and traces of acid scarring. Farmers and caravanners alike are frightened and confused; these recurring attacks threaten to destroy \textit{Westruun's} agriculture and disrupt commerce. Jeremiah Sook, head of the regional Rancher's Guild, is willing to pay a hefty bounty for the culprits.

\subparagraph{An Assistant's Plight (mid-level)}
Just outside \textit{Westruun} is Slaterock Tower, the abode of renowned archmage \textit{Realmseer Eskil Ryndarien}. Normally, none are invited to the tower, but the characters, wherever they may be, receive a missive from the Realmseer's assistant, Jekt. The letter claims the Realmseer has gone missing within the tower, and the building itself now bars Jekt's entry—and apparently, the Realmseer's exit. Upon arrival, and confirmation of Jekt's request, the tower seems to constantly shift its interior structure, almost acting like a sentient entity. Challenge after deadly challenge is thrown at the party as they traverse the labyrinthine paths of the tower, only to find that the Realmseer himself is actually trying to keep something far more dangerous imprisoned within.

\subparagraph{The Stuff of Nightmares (mid-level)}
Reports of missing children rise once again within \textit{Westruun}, with a seemingly unrelated scourge of nightmares plaguing people old and young alike. Mystic expertise is sought to acquire information regarding either issue, but the sadistic \textbf{oni}, secretly responsible for both, is deft at erasing its tracks. Nyx and Vryll, twin tiefling mages with only a single horn each, have also recently opened a dream-reading stall, and their business is booming. Are they responsible, or is their exploitative business just a red herring?

\subparagraph{Down Below (high-level)}
A massive sinkhole has emerged beneath the First Bastion, taking a portion of the temple into the darkness of the Underwalk Ward. With the depth seemingly bottomless, and initial investigators going missing, high priest Umentu Helmdoff seeks the assistance of brave and able-bodied adventurers to find answers and save the people still in the temple when it vanished. Only one thing is certain when descending into the unknown below—the terror only grows as the light from above fades.

\subsection{Wildwood Byway}
Known as the oldest road in Tal'Dorei, this path linked the adolescent cities of \textit{Kraghammer} and \textit{Syngorn} before the landfall of humankind. During the reign of \textit{Drassig}, the \textit{Silvercut Roadway} intersected with this ancient road. Warren \textit{Drassig} built countless fortifications along this road to defend his holdings against \textit{Syngorn}, a move that made him many allies within \textit{Kraghammer}.

While traveling the Wildwood in broad daylight isn't any more dangerous than riding along the Silvercut, it's common knowledge that the forest at the northern end of this road is fraught with danger. Wild beasts, cultists, and bandits make the Wildwood their home.

\section{Cliffkeep Mountains}
Spanning the northern reaches of Tal'Dorei, the \textit{Cliffkeep Mountains} are an impenetrable barrier between the Republic of Tal'Dorei and the frozen \textit{Neverfields}. The southern foothills are dotted with forests and pockets of civilization, standing resolute against the unwelcoming topography. Deep within the earth itself, the dwarven stronghold of \textit{Kraghammer} flourishes and expands, just a few scant leagues above ancient caverns that plunge into a domain of unknowable horrors.

The \textit{Cliffkeeps} are the largest mountain range in all of Tal'Dorei, and innumerable types of beings make their homes upon the crags, in mountainside caves, and in the miles upon miles of tunnels that descend into the earth below them. dwarves, half-giants, goblins, and kobolds are the best-known peoples of the \textit{Cliffkeeps}, but the mountains are also home to the fire giants of \textit{Vulkanon}, wandering ettins and hill giants, hungry cyclopes, \textbf{frost giant} marauders from the \textit{Neverfields}, and cave-dwelling trolls.

Near the earth elemental rift and \textit{the Ashari} village of \textit{Terrah}, the mountains themselves seem to rise up in defense as earth elementals spring from the living stones. Northward, peaks reach higher and winds grow colder as snow and ice take the lands to become the frozen waste of the \textit{Neverfields}. These challenges keep the common folk at bay, the secrets and spoils held within these crags calling to the brave, the clever, and the foolhardy.

\begin{description}
\item[Majority Faiths:]
The All-Hammer

\item[Minority Faiths:]
The Platinum Dragon

\item[Imports:]
Lumber, spices, fish, grain, livestock

\item[Exports:]
Precious gems, industrial metals, quartz, granite, gold, silver, platnium, mithral, cobalt, jewelry, fungal produce, weapons, armor

\end{description}


\subsection{Fort Daxio}
Nestled within the \textit{Othendin Pass}, north of \textit{Emon}, \textit{Fort Daxio} is a massive stronghold and series of military structures that, under the rule of the \textit{Tal'Dorei Council}, trains Tal'Dorei's most elite warriors. The fortress was constructed following the \textit{Scattered War} as a strategic bulwark against opportunistic humanoid and giant bandit forces in the \textit{Cliffkeep Mountains}, and as a hidden reserve of military might in the event that \textit{Emon} was ever threatened.

Though Tal'Dorei doesn't maintain a standing army—rather, it assembles one in times of strife from mercenaries, adventurers, and local militias—the forces of \textit{Fort Daxio} are an exception. The legendary Daxio Outriders are a corps of cavaliers some four thousand strong. Half their forces remain stationed in \textit{Fort Daxio} and its surroundings at any given time, serving as a historic barrier between \textit{Emon} and the bandits and creatures of the \textit{Cliffkeep Mountains}. The other two thousand Outriders are spread across the continent, typically in battalions of up to one hundred riders, safeguarding villages, patrolling the \textit{Silvercut Roadway}, putting pressure on the \textit{Ravagers} of the \textit{Dividing Plains}, or protecting traveling diplomats.

\textit{Fort Daxio} itself is led by General Elle Gorgofon, a dour but diplomatic tactician. She served faithfully under the fort's previous leader, Warmaster Mikael Daxio, for nearly thirty-five years before he was killed at the hands of a \textbf{vampire} hoping to throw the fort into chaos. Daxio was resurrected, but is enjoying an extended holiday in \textit{Syngorn} rather than returning to arms.

The Daxio Outriders are divided into four different regiments: The Dusk Regiment, the Gale Regiment, the Aegis Regiment, and the Tide Regiment. Each regiment is led by their own assigned general, all answering directly to General Gorgofon (who in turn answers to the Master of War on the \textit{Tal'Dorei Council} in times of conflict).

The Dusk Regiment is responsible for covert operations across the country and is led by General Fei Yujian, a cheerful, affable man. The Gale Regiment is trained in skyborne combat in addition to horse-riding; after the \textit{Chroma Conclave} was defeated, \textit{Fort Daxio} trained a variety of griffons, hippogriffs, and wyverns to ward against airborne threats. It is led by General Jillian Sylph, a calm, tactical human with air elemental ancestry. The Aegis Regiment is charged with remaining in \textit{Fort Daxio} to protect \textit{Emon}, and is led by General Elle Gorgofon herself. The Tide Regiment is trained in spellcasting and ship-to-ship combat, in case \textit{Emon} is ever threatened from the western seas. This final regiment is led by General Kay Clearsight, a half-elf trained in the arcane arts at the renowned Soltryce Academy in the Dwendalian Empire of \textit{Wildemount}.

\subsubsection{Fort Daxio Adventures}
Game Masters who set their adventure in \textit{Fort Daxio} can use this plot hook for inspiration.

\subparagraph{Vampire Hunters (high-level)}
Years ago, Warmaster Mikael Daxio was killed by a \textbf{vampire} named Vrylaska Fellbranch. She disappeared without a trace, and the Daxio Outriders have been searching fruitlessly for her for years. While the characters are in \textit{Fort Daxio}, a scout returns with breathless news of vampire attacks in remote settlements in the \textit{Cliffkeep Mountains}. General Elle Gorgofon suspects a trap, but has no choice but to mobilize the Aegis Regiment—and the characters, if they accept her offer of 5,000 gp—to hunt the monster down.

Little do they know that, while the Aegis Regiment is away, the vampire's minions hidden in \textit{Fort Daxio} will come out to play.

\subsection{Gatshadow}
The mighty mountain known as \textit{Gatshadow} is a spire of death that looms over the \textit{Dividing Plains}. \textit{Westruun} sees it always upon the horizon, and the most vicious monsters of the \textit{Bramblewood} grow cruel in its fell shadow. Folkloric tales about the evil powers within \textit{Gatshadow} are told in all cultures of the \textit{Dividing Plains}: stories of giant spiders in the forest, of beings of pure shadow that dwell within the mountain's tunnels, and of the tunnels themselves, which are so labyrinthine and shifting that no one who enters will ever escape. These days, most of these stories are told to keep children from wandering into the forest alone—but they are far from simply children's tales.

\textit{Gatshadow} was transformed into a place of evil in the \textit{Age of Arcanum}. Acek Orattim, a priest of the Betrayer God known as the The Chained Oblivion, made the mountain his seat of power. The The Chained Oblivion inexorable might corrupted the mountain, stretching its stone-like putty until it was a jagged needle that towered over all other mountains, and etching a mind-bending labyrinth of twisting passages into the mountain's heart.

Acek spread the The Chained Oblivion nihilistic despair across the \textit{Cliffkeep Mountains} and the \textit{Dividing Plains} below. All bowed before their influence their minds lost to unfathomable contemplation of the end of all things. Untold thousands died at the hands of Acek and his god, willingly entering \textit{Gatshadow}—which then became their tomb.

The bitter cold of oblivion threatened to consume the land, but it was repulsed by the The Dawnfather and his warriors of light. The The Dawnfather struck a telling blow against the Oblivion, and doggedly pursued as it retreated to \textit{Gatshadow}. He shackled and banished the god of the end within \textit{Gatshadow's} depths. It was said that Acek's mind and spirit were so saturated with his master's essence that the god's banishment tore his body asunder. Today, only the empty maze of deadly chambers and traps remain beneath the mountain, forgotten and awaiting discovery.

The high bluffs of \textit{Gatshadow} have in recent years been the home of many a villain, most recently Umbrasyl the Hope Devourer, the black dragon of the \textit{Chroma Conclave}. Legend says that \textit{Gatshadow's} interior tunnels spiral infinitely downwards, into a void of utter darkness. Immortal, white-eyed devotees of the The Chained Oblivion pace endlessly into \textit{Gatshadow's} tunnels, carrying lightless torches, whispering that the secret to freeing their god from the Divine Gate is sealed in the depths below. The stone doors that bar entrance to \textit{Gatshadow's} lowest reaches are sealed by magic so great, even the most inquisitive of the \textit{Arcana Pansophical} can't breach them—and it is well that they cannot, for who knows what terror would be unleashed if they were opened.

\subsubsection{Gatshadow Adventures}
Game Masters who set their adventure in \textit{Gatshadow} can use these plot hooks for inspiration.

\subparagraph{Blood of the Dragon (high-level)}
The body of the \textbf{ancient black dragon} Umbrasyl was never recovered. Some think that's a waste of a perfectly good business opportunity—after all, a black dragon's bones are a potent alchemical reagent, and its scales stronger than any steel.

However, there's a reason no one has dared to recover the Hope Devourer's corpse. When Umbrasyl died, the caustic acid and magical blood within his body seeped through the caverns of Upper \textit{Gatshadow}, devouring stone and transforming the mountain's upper levels into a vitriolic maze teeming with oozes, half-black dragon warriors, and the restless spirits of \textit{goliaths} from the Herd of Storms that failed to kill the dragon.

\subparagraph{No Lights in the Darkness (epic-level)}
The \textit{Betrayer Gods} were banished to their planes of eternal imprisonment, but something has upset the laws of the planes. A being called the Nameless—supposedly a cultist of The Whispered One—has stolen her master's secret knowledge and fled to \textit{Gatshadow} in search of the tortured spirit of Acek Orattim. Agents of the \textit{Arcana Pansophical} have sensed that the lower gates of \textit{Gatshadow} have been breached, and frantically call for Tal'Dorei's greatest heroes to rush to the evil mountain. As the characters hurry to the bottom of the mind-sundering labyrinth within, they discover that they race against not only the new Herald of Oblivion, but also the \textit{Remnants} she betrayed.

\subsection{Jorenn Village}

\begin{description}
\item[Large Village:]
Population 2,480 (64\% humans, 9\% dwarves, 7\% halflings, 7\% gnomes, 13\% other races)

\end{description}

\noindent For many years, Jorenn Village was just a tiny pioneer town trying to profit from the shadegrass of the \textit{Umbra Hills}. Over the past century, the village has exploded in size and population, after plentiful silver veins were discovered within the nearby rocky hills.

Hardworking common folk and opportunistic prospectors alike flocked to Jorenn, despite its unsettling proximity to the \textit{Umbra Hills}. Jorenn is far enough away from \textit{Emon} that the influence of the \textit{Tal'Dorei Council} rarely touches it, giving opportunists and outlaws freedom to subvert and control both the village government and the local economy. Thirty years ago, a small circle of vigilantes and mercenaries declared themselves protectors of Jorenn, and took the name Shadewatch, hunting and slaying creeping dangers in the night and returning to great fanfare. Over the past three decades, others seeking fame and glory joined the ranks of the Shadewatch, and the group ballooned into the unofficial town law enforcement.

The Shadewatch is deeply corrupt, and it has been almost since its first year. The vigilantes personally stoke the fires of fear, all while demanding increasingly galling tribute for their services—and breaking kneecaps if they don't get it. Today, the people of Jorenn are just as afraid of the Shadewatch as they are of the creatures the watch keeps at bay.

The Shadewatch is a minor organization in the grand scheme of things, but Shademaster Arhanna Lewyn has grand ambitions. She has sent messages to every faction in Tal'Dorei, from the \textit{Myriad} to the \textit{League of Miracles}, with an interest in cruel subterfuge. It's clear that she's planning something: every night, she consults a ball of polished white quartz, and every week she sends a trio of her agents to the \textit{Umbra Hills}—or to the base of \textit{Gatshadow}.

\subsubsection{Jorenn Village Adventures}
Game Masters who set their adventure in Jorenn Village can use this plot hook for inspiration.

\subparagraph{An Offer You Can't Refuse (low-level)}
The characters are attacked by a group of worgs, and more just keep coming. Just as they're about to be completely overwhelmed, they're saved by a sudden volley of black-feathered bolts! Six thugs join the battle and "rescue" the characters by driving off the remaining worgs. These Shadewatch "officers" ask for the characters' names first, then brusquely declare that they like their moxie. They offer the characters a position in their force back in Jorenn, with unbelievable pay: 50 gp a day each.

If this deal sounds too good to be true, that's because it is: the Shademaster needs strong-willed sacrifices if she's to open the gates of \textit{Gatshadow}, a feat she can accomplish by teleporting up to twelve creatures into the depths of the mountain with her whispering quartz orb.

\subsection{Kraghammer}

\begin{description}
\item[City:]
Population 73,550 (73\% dwarves, 8\% gnomes, 5\% humans, 14\% other races)

\end{description}

\noindent Dozens of dwarven redoubts rose and fell in the tunnels beneath the \textit{Cliffkeep Mountains} during the \textit{Calamity}, and in the uncertain years that followed it. Legend holds that a great civilization called Uthtor once reigned supreme in the \textit{Cliffkeeps}, and that all of the deep settlements are fragments of Uthtor's glory, fractured by the \textit{Calamity}. Over the centuries, these fragments divided and coalesced dozens of times until, finally, those who remained came together to create Kraghammer.

Clan Jaggenstrike united thousands of dwarves, gnomes, and other dwellers of the deep halls in the creation of Kraghammer. Having quietly died out centuries ago, the clan now holds a position of mythic reverence within the city. Kraghammer's founding was dedicated to the The All-Hammer, divine creator of the dwarves and patron of creation and ingenuity. This god's industrious spirit filled the people of Kraghammer and blessed them with rich veins of platnium, mithral, adamantine, and other precious metals that have sustained them since the city's founding.

An ominous walk of a thousand stairs climbs the side of the mountain to the gates of the city—carved from the mountain itself, just wide enough for two dwarves to walk abreast. With the path too steep and treacherous for carts or common beasts of burden to traverse, traders and trappers use giant goats to transport their goods to and from market, while wealthy dwarven merchant lords employ serfs and laborers to carry their goods. Kraghammer's impenetrable adamantine gates stand at the end of a perilous bridge overlooking a hundred-foot chasm, guarded by three watchtowers carved directly into the mountain.

Within those gates, the massive stronghold is made up of three cylindrical levels called "slabs" that descend deep beneath the mountains. Kraghammer's granite walls are lit by roaring fires, enchanted torches, and shimmering candlerock, which glows with internal fire. Even when the evening fires are extinguished, the amber light of the Bronzegrip Metalworks at the base of the city bathes Kraghammer in its warm glow.

\subsubsection{Government}
\textit{Kraghammer} law is written by a consortium of the five most powerful houses within the city, each appointing a delegate to serve on their behalf. That tenuous union of houses is kept by the city's executive officer, the Ironkeeper. The Ironkeeper is elected by the houses and carries a term of ten years, with no term limit, though most Ironkeepers contain their rule to three to five terms.

The current sitting Ironkeeper is \textit{Palluda of House Zuurthom}. She rose to prominence advocating for the rights of workers in the wake of the \textit{Chroma Conclave}, and is a strong advocate for uniting with the Republic of Tal'Dorei to strengthen their might against future supernatural threats. She is sitting comfortably in the midst of her second term and enjoys strong support from members of \textit{Kraghammer's} working class—and staunch opposition from traditionalists, who cling to the city's fading culture of \textit{isolationism}.

The laws of the city are enforced by a military class of elite warriors called Carvers that act as guards, soldiers, and jailers under the guidance of the Ironkeeper. Becoming a Carver is a lengthy, rigorous, and taxing pursuit, and many who choose to join fail to complete their training.

Most who join the Carvers are dwarves from families with centuries of history within the Carvers. This isn't just because of the pressure to uphold their family's honor. Most hopefuls are also those who value \textit{Kraghammer's} ancient traditions and rigid social order—a set of beliefs largely held by dwarven clans with long family histories. Once officially named, a Carver is given a homestead built into the center slab of the city, a set of immaculate armor and weapons, and a steady income.

In the years following the defeat of the \textit{Chroma Conclave}, the famously \textit{isolationist} and xenophobic houses conceded to formally make an alliance with the Republic of Tal'Dorei. An ambassador from \textit{Kraghammer} now sits on the \textit{Tal'Dorei Council}, and rumors are beginning to circulate within the city that the heads of the Great \textit{Houses of Kraghammer} are in talks with the \textit{Tal'Dorei Council} to make \textit{Kraghammer} an official constituent city-state of the republic.

\subsubsection{Isolation of Kraghammer}
It's not easy being an outsider. \textit{Kraghammer} has long been a city of dwarves, for dwarves, and anyone else within the city is an interloper. Some traditionalists within the city still think that surface-born folk should spend as little time in \textit{Kraghammer} as possible, so as not to dilute the ancient splendor of their culture.

Yet times are changing. A coalition of people, ranging from dwarves to gnomes to traditionally surface-dwelling people who have made \textit{Kraghammer} their home, have demanded that the Great Houses amend the city's most xenophobic laws. Despite the angry protestations of the city's most bigoted dwarves, the Great Houses have acquiesced to the requests of their people. This is due in no small part to these protestors having the support of the current Ironkeeper—\textit{Palluda Zuurthrom} has long been a proponent of \textit{Kraghammer} entering the Republic of Tal'Dorei and ending their centuries of relative isolation.

No non-dwarf has ever served as a delegate of the houses, let alone the Ironkeeper. Progress is slow, especially in \textit{Kraghammer}, but it is happening. The most despicable of the city's xenophobes have left of their own volition, creating other dwarven settlements throughout the \textit{Cliffkeeps} where they can keep to themselves and "their own kind" in peace.

\subsubsection{Crime}
Citizens of \textit{Emon} who visit \textit{Kraghammer} often marvel at how lawful and organized the dwarven enclave is. Crime is ostensibly low within \textit{Kraghammer} thanks to the Carvers' protection, but this peace comes at a cost. Crime is low because the courts of \textit{Kraghammer} almost always reach a guilty verdict. As a long-lived and durable people, most dwarves, particularly those with a family reputation to uphold, would rather plead guilty and endure years, even decades, of hard labor than disturb the fabric of their society by fighting conviction.

Beyond this, corruption is rife within the Carvers, and they are able to brutalize even innocent people with relative impunity. This corruption is shielded from public scrutiny by the vast wealth and political influence of House Glorendar, their patron house. Because of Glorendar's close ties with the Carvers (and the mutual benefits thereof), corruption has been hard to uncover within either organization.

Over twenty years ago, however, a gnomish journalist named Ida Mudrake uncovered something damning within the vaults of House Glorendar, bringing some hope that change may yet move steadily forward. She has been forced to flee the city and now lives with a security detail in \textit{Emon}, but she works tirelessly to bring the crimes of House Glorendar to light. Many in \textit{Kraghammer} hope that this evidence, whatever it is, will shake both Glorendar and the Carvers to their foundations. Many more fear that even true justice is not worth the chaos that bringing Glorendar's misdeeds to light will cause.

\subsubsection{Geography and Climate}
Though the \textit{Cliffkeep Mountains} themselves are blanketed in thick snow, the bone-piercing chill does not extend into \textit{Kraghammer} itself. Heat from the Bronzegrip Metalworks and the other blast furnaces in the Bottom Slab rises and permeates the entire city, though the Top Slab still grows cold in the winter. The city itself is carved out of the granite core of Mount \textit{Kraghammer}; though the dwarves rarely see the night sky, the inner walls still sparkle like starlight in the amber glow of the furnaces.

\subparagraph{Top Slab}
Also called the Arch, this uppermost tier is the entry level of the city, built on a massive ring. Like a hub with many spokes, countless neighborhoods and tunnels spread out from the Top Slab. A fine layer of soot tends to coat much of this slab, due to imperfect ventilation from the industrious center slab; despite this, this level is where most people make their homes, with thousands of stone-built abodes dotting the walls, while popular taverns like the trendy Firebrook Inn and the historic but austere Ironhearth Tavern fill alcoves pushed into the rock. For hundreds of years, most non-dwarves within \textit{Kraghammer} were relegated to living in a Top Slab neighborhood called the "Otherwalk." These days, non-dwarves live all over the city, and the Otherwalk has been renamed the Cavernwalk. It's still mostly inhabited by non-dwarves, and it has a vibrant culture all its own.

As members of the \textit{Kraghammer} working class, "Topslabbers" either make the long commute to labor in the Bottom Slab or work in the Toppers' farms above the mountain. Others are indentured servants to wealthy dwarves in the Middle Slab, and only return to their families on the Top Slab on their monthly day of rest.

\subparagraph{Center Slab}
Also referred to as the Heart, this middle tier is the widest and most varied rung on the journey through \textit{Kraghammer}. It is stratified between the residences of \textit{Kraghammer's} old money (its long-established noble houses), and the new (the wealthy merchant class). At the exact center of the slab is the Pyrethrone, seat of the Ironkeeper. The geometric symbolism of the Pyrethrone being at the dead center of the city is a fact most visitors to \textit{Kraghammer} hear ad nauseum. Radiating out from the Pyrethrone are the various fortress-manors of the dwarven noble houses, with several Carver barracks in easy reach.

Below, the outer rings of the Center Slab are occupied by most of the city's non-mining businesses, from smithies, breweries, and jewelers, to tailors, butchers, and tinkerers like the \textit{Cracksackle Union}. The Hunter's Club is also well known as the chief provider for non-imported meats to the city, sending parties into the mountains or the nearby \textit{Torian Forest} for wild game. The extravagant marble temple known as The All-Hammer stands on the center slab of the city, the grandest of all shrines to the most worshipped deity under the mountain. Enormous in size and impeccable in both detail and architectural design, this hall calls to artisans and crafters throughout the city for inspiration.

Of the five noble houses, only two conduct their affairs entirely in the Center Slab. The magically talented Lord Steddos Thunderbrand and his family are known to rarely leave their mansion, except on business with the Ironkeeper. Wallera Glorendar used to oversee the combat training of all Carvers within \textit{Kraghammer}, but hasn't been seen in public for twenty years. Her daughter, Kazya, runs all her affairs within the public eye.

\subparagraph{Bottom Slab}
Best known as the Pit, this tier is the industrious powerhouse of \textit{Kraghammer}, boasting the most forges per square mile of any location in Tal'Dorei, courtesy of the Bronzegrip Metalworks and their competitors. Most of these forges are potent, highly industrialized, and mechanized blast furnaces. The expansive \textit{Keenstone Quarry} is chief among the numerous mining operations that have tunneled through the depths of the mountains—depths whose true dangers are obscured by centuries of legend, and therefore not fully understood.

A network of fungus-farming paths called the Glowgrove works its way around this level, passing through all sorts of caverns. The mushrooms are harvested by the Toppers, the dwarven agricultural league that got their name from the small topside farms they tend above the city.

Three of the five dwarven houses have business interests in the Bottom Slab, though they direct its operations from the next tier up. \textit{Nostoc Greyspine} is the grim overseer of \textit{Keenstone Quarry}. Haddi Bronzegrip is the foredwarf of the Bronzegrip Metalworks; she claims to be the wealthiest dwarf in the North, a claim no one has ever been able to successfully dispute. Blenton Zuurthom is an architectural savant, and one of the older and kinder noble dwarves; he tries his best to oversee every major construction project in the Bottom Slab.

\subsubsection{Points of Interest}
\textit{Kraghammer} has a number of famous landmarks, as well as a few mythic locations said to exist in the caverns beneath the city, as depicted on its map.

\subparagraph{The Starshrine (Top Slab)}
Most dwarves worship the The All-Hammer, and his teachings of community and ancestral piety have been all but completely subsumed into \textit{Kraghammer's} culture at large. For outlanders with different faiths, however, the monolithic religious culture can feel inescapable. Shrines to all the \textit{Prime Deities} have spread throughout \textit{Kraghammer} since the fall of the \textit{Chroma Conclave}, but the Starshrine was the first.

When the world trembled at the might of \textit{Thordak the Cinder King}, tremors rocked the upper layers of \textit{Kraghammer}, causing certain caverns to collapse—and in some cases, to reveal places of great beauty. The Starshrine is a vast grotto whose granite walls mysteriously sparkle with a perfect replica of the stars in the night sky. It can be found within the Top Slab's culturally diverse Cavernwalk district. Once per year, when a character prays to one of the \textit{Prime Deities} here, there is a Your prayer goes unanswered. chance per character level that their prayer is answered, as the cleric's Divine Intervention class feature.

\subparagraph{Cracksackle Guildhall (Middle Slab)}
The headquarters of the Cracksackle Union looks unlike anything in \textit{Kraghammer}. It's made of a dozen glistening steel domes, like sleek, metal igloos, and the light of unusual experiments flicker through its slitted windows at all hours of the night. Most within \textit{Kraghammer} consider the building an eyesore, but all interested in technological tinkering and ingenuity love it for its uniqueness. If nothing else, it plainly displays the oft-baffling ingenuity of its tinkerers.

The Cracksackle Union was founded by gnomish refugees from \textit{Wittebak} nearly four hundred years ago. While intrinsically tied to the Bronzegrip Metalworks for resources, contacts, and distribution, the Cracksacklers have made themselves indispensable to the dwarven elite with their inventive contributions to mining and masonry technology.

On any working day, a dozen tinkerers, many of them gnomes, will set up shop in the courtyard, selling strange and untested devices they have made in their workshops. The GM determines which non-magical curiosities can be purchased here, but some include:

\begin{DndTable}[header=Cracksackle Wares]{lcX}

Item & Cost & Description\\
\textit{Dynamite} & 100gp & This bundle of explosives can be lit and thrown up to 30 feet as an action. It explodes at the start of your next turn. Creatures and objects within 20 feet of it must make a DC 12 Dexterity saving throws, taking 6d6 fire damage on a failed save, or half as much damage on a successful one.\\
\textit{Gluebomb} & 50gp & This unique explosive can be primed and thrown up to 30 feet as an action. It explodes on impact. Creatures within 20 feet of it must make a DC 12 Dexterity saving throws or be restrained for 1 minute. A restrained creature can make a DC 12 Strength check as an action, escaping on a success.\\
\textit{Stink bomb} & 25gp & This smelly explosive can be primed and thrown up to 30 feet as an action. It explodes at the start of your next turn. Creatures within 20 feet of it must make a DC 12 Constitution saving throws or be poisoned for 1 minute. A creature can repeat this save at the end of each of its turns, ending the effect on itself on a success.\\
\end{DndTable}

\subparagraph{Keenstone Quarry (Bottom Slab)}
Despite all of \textit{Kraghammer's} wondrous crafts and goods, its most significant export is raw stone—most of which is excavated from Keenstone Quarry, making the quarry's overseer, \textit{Nostoc Greyspine}, a very rich dwarf indeed. As the legend of Vox Machina spread across Tal'Dorei, \textit{Overseer Greyspine} got an idea: open the quarry to adventurers. Vox Machina famously entered the vast network of dangerous tunnels beneath \textit{Kraghammer} from a fissure at the depths of the quarry, and countless imitators have tried to follow in their footsteps—some emerging with fabulous treasure, but most never returning at all.

To take advantage of this opportunity, \textit{Greyspine} has advertised the quarry far and wide as the "Gateway to the Depths," and posted an entrance fee of 25 gp per adventurer. Of course, neither \textit{Greyspine} nor the staff of the quarry bears any liability for death, dismemberment, or emotional trauma sustained with the depths.

\subparagraph{Hall of Burning Mushrooms}
The Bronzegrip Metalworks just carved deeper into the mountain to expand their furnaces, only to discover the most magnificent cavern: a mile-wide cave filled with bioluminescent purple mushrooms. Industry had to move in, and the Bronzegrips hired Wyrmhide Thunderbrand and his pyromancer brigade to torch the cavern. The mushrooms never stopped burning, and the myconids that lived in the fungal forest never stopped fighting. Characters known within \textit{Kraghammer} may be requested by House Bronzegrip or Thunderbrand to cut through the inferno and destroy the myconid monarch.

\subsection{Emberhold}

\begin{description}
\item[City:]
Population 21,180 (73\% dwarves, 15\% elves, 12\% other races)

\end{description}

\noindent A passage from the \textit{Keenstone Quarry} in \textit{Kraghammer} leads below the earth. Deep. Deeper. Until the light of day is entirely gone, and then deeper still. Across chasms miles-deep and full of magma, beneath great hanging gardens of glowing \textit{malachite}, lies the \textit{Emberhold}. It stares imperiously over an underground kingdom of the \textbf{duergar}, dwarves who long ago broke away from the rule of the Great \textit{Houses of Kraghammer}.

Millennia of fear and superstition have turned the dwarves of \textit{Kraghammer} into bitter, fearful enemies of their estranged kin. Untold dozens of duergar villages span the glassy lava-plains that the \textit{Emberhold} is built upon, with mighty armies constantly coming and going, sometimes returning with just a few bloody survivors, other times returning with captives in chains—and other times returning not at all.

The dwarves of the \textit{Emberhold} rarely rise to the surface; they have lived so long in the darkest depths of the earth that the light of the sun sears their eyes and prickles their skin. Few ever arise even to \textit{Kraghammer}, for they are taught that only death awaits them there. Nevertheless, many duergar doubt the narrative that has been fed to them since birth, and make attempts to escape to the surface. Those duergar that have left the ways of the \textit{Emberhold} behind and learned of the outside world have remarked how unbelievable it is that their people were so misled by lies about the barbaric cruelty of surface-dwellers.

\subsubsection{Government}
The \textit{Emberhold} and the towns that serve as its vassals are ruled by an absolute monarch. The current King of Embers is the wrathful King Thangrul II. He rules through fear and conquest, spreading his armies to the farthest corners of Tal'Dorei's underground caverns in the name of the The Crawling King.

Thangrul's rule is supported by a contingent of dark elves from their decadent, alien domain of \textit{Ruhn-Shak}. Most in Thangrul's court know well enough to mistrust these ambassadors of the far-flung caverns, for their minds have been twisted by the influence of not just the The Spider Queen, but also the aberrations with whom they have made strange alliances in order to sustain their power. Thangrul himself cares little for his advisors' mistrust of his elven allies and openly declares that he will make use of them as long as they are useful.

Vox Machina's assault of the \textit{Emberhold} over twenty years ago cost the duergar their king and queen, but rumors of a lost princess of Ember have begun to circulate through the realm, and the survivors of Vox Machina's raid have rebuilt their hierarchy and continued their endless war of conquest throughout the depths.

\subsubsection{Geography}
The \textit{Emberhold} itself is a fortress built of black stone and \textit{obsidian}, framed by rivers of molten rock and surrounded by a massive city of ramshackle dwellings and towers that mark the boundaries of its peoples' domain.

\subsubsection{Emberhold Adventures}
Game Masters who set their adventure in the \textit{Emberhold} can use this plot hook for inspiration.

\subparagraph{The Lost Heir (mid-level)}
King Murghol of the \textit{Emberhold} was killed in combat by \textbf{Scanlan Shorthalt}, gnome bard of Vox Machina, and his wife, Queen Ulara, was slain by the aberrant tyrant K'Varn. All of the \textit{Emberhold} knows that the King and Queen had a young daughter—only six years old at the time of their death—but none seem interested in seeking \textit{Emberhold's} missing heir. The new king, Thangrul II, grows particularly fiery whenever anyone mentions her name in his presence.

\subsection{Yug'Voril}
Where the dwarves of \textit{Kraghammer} whisper fearful words of the \textbf{duergar}, so too do the \textbf{duergar} whisper of Yug'Voril. Beneath the deep \textit{dwarves'} fiery kingdom, ever deeper into the heart of Exandria, there is a miles-wide grotto coated in glowing crystal. A lake fills the grotto, its surface so smooth and placid it seems like a \textit{jade} mirror, perfectly reflecting the light of the crystals around it. In the middle of the lake is an island, atop which loom the grand, austere towers of Yug'Voril.

The clean lines and beautiful facade of the city belie its current, abominable denizens.

Thousands of years old, the builders and previous occupants have long vanished, leaving the unspoiled ruins to become home to a host of psychic, mind-eating creatures from a plane beyond the boundaries of reality. The tyrannical society of these alien creatures flourishes here, far from their cosmic pursuers, giving them the freedom to build a legion of thralls: mind-slaves whose delicious brains are allowed to remain in their bodies...until those bodies fall apart from toil.

Yug'Voril and its inhabitants fell recently under the control of the extraplanar aberration called K'Varn, before he was slain by Vox Machina. Now that the alien mind-eaters have been freed from K'Varn's tyranny, they have reclaimed Yug'Voril as their own and have continued work on their strange and terrible plots.

\subsubsection{Yug'Voril Adventures}
Game Masters who set their adventure in \textit{Yug'Voril} can use this plot hook for inspiration.

\subparagraph{A Mind Is a Terrible Thing to Waste (high-level)}
Four of the foremost minds at \textit{Emon's} \textit{Alabaster Lyceum} have disappeared without a trace. All vanished on the same night, and there have been no confirmed sightings since. Investigations have been underway for several months, and some dwarves of \textit{Kraghammer} have recently reported seeing four figures that could be the missing scholars entering the lower passages of the \textit{Keenstone Mine} under cover of night.

Something has mentally dominated these scholars. The alien inhabitants of \textit{Yug'Voril} hunger for a specific morsel of powerful knowledge within their brilliant minds.

\subsection{Lyrengorn, the Elvenpeaks}

\begin{description}
\item[Town:]
Population 5,430 (93\% elves, 3\% dwarves, 2\% humans, 2\% other races)

\end{description}

\noindent Every year, dozens of travelers from across Tal'Dorei travel past the far northern reaches of the \textit{Cliffkeep Mountains} to see the The Moonweaver Ribbons twist around the Elvenpeaks, pulled through the sky by those mysterious figures on wyvernback. The Elvenpeaks are a single mountain split perfectly in the center, with ice-sheeted forests topping both peaks. While the exterior of these forest canopies appears frozen and without life, the interior is instead a warm and humid wood, blooming with flowers of all colors and roamed by living plants and wild game. In the boughs of these thick tropical trees also lives a society of wood elves, some of those wild folk who broke away from Yenlara's pact and the city of \textit{Syngorn}. Their city of Lyrengorn—\textit{lyren'gorn} literally meaning "mountain city" in Elvish—is the last bastion of safety before crossing into the endless ice of the \textit{Neverfields} to the north.

These lyren'alfen—"high elves"—call themselves friends of none and enemies of fewer; they welcome all who manage to make it to their home, provided they come in peace, and take no sides in any conflict. In ages past, they were highly sought-after as mercenaries because their warriors had learned the rare art of \textbf{wyvern} riding. When Warren \textit{Drassig's} legions marched to conquer Tal'Dorei, the elves of Lyrengorn refused to die for anyone else's war, retreating instead to their mountaintops and lush forests.

Today, few of the lyren'alfen are actually warriors, though many among them are hunters. In the intervening centuries, \textbf{wyvern} riding has transformed from a martial art into a ceremonial one; great athletes called skyswimmers train year-round for the annual coming of the The Moonweaver Ribbons, great streaks of multicolored light that illuminate the northern sky on the winter solstice. The skyswimmers and their \textbf{wyvern} Mounts and Vehicles seize the strands of light on their spears and paint radiant, esoteric portraits in the wintry sky, drawing intrepid outsiders from across Tal'Dorei to camp on remote mountains and marvel at their visual poetry.

Beyond the ceremony's artistic merit, it's also a ritual to the goddess known as the The Moonweaver, chief deity of these elves. The high priestess of the moon claims that the annual rite purifies Lyrengorn and gives the wild elves long life.

\subsubsection{Lyrengorn Adventures}
Game Masters who set their adventure in Lyrengorn can use these plot hooks for inspiration.

\subparagraph{The Impossible Art (low-level)}
While traveling through the \textit{Cliffkeep Mountains}, the cold becomes too much for the characters to bear, and they pass out in the snow. When they awaken, they find themselves being tended to by a trio of Lyrengorn \textbf{wyvern} riders, but their roles are suddenly reversed when a massive band of raiders, most of them goblins and human thugs, pincushions the riders in mid-flight. When the characters crash in the snowfields, they find their saviors dead, and must try to learn the impossible art of \textbf{wyvern} riding before they are claimed by the blizzard.

\subparagraph{The Incorruptible Light (high-level)}
While the characters are in Lyrengorn, they discover a small community of dark elf refugees. They fled here from a monstrous, nameless source of aberrant corruption that has been devouring their society for centuries. The dark elves hope that the skyswimmers' ceremony will cleanse them of the corruption and fits of sadistic cruelty that still plague them from their days under the abberrations' terrible thrall. Strange accidents conspire to kill the skyswimmers and the priestesses, and on the day of the solstice, a pair of many-eyed aberrations rise from the depths of the earth and try to bend the The Moonweaver Ribbons to their evil purposes.

\subsection{The Neverfields}
The farther north one goes into the \textit{Cliffkeep Mountains}, the taller the peaks grow and the colder the weather becomes, until the snowstorms are so thick that no human unprepared for their chill can survive. The spiky peaks grow shallow, trading treacherous rock for glacial ice and the untempered fury of endless blizzards. No bastions of living society exist this far north, though some brave, foolish warriors endure the elements here to hunt rare and powerful game or to seek lost relics from before the \textit{Calamity}.

Legends tell of a time when this region was green and temperate, when trees and joyous elven enclaves covered the land. Moreover, these tales say that relics of this elysian era are still buried deep beneath the icy wastes. Many that go seeking answers to these rumors never return. The few that do, however, occasionally emerge with artifacts of intense power and mystery, reigniting interest in the region and sending further generations to their doom among the ancient, ice-rimed ruins and plains of endless glacial wasteland.

\subsubsection{Neverfields Adventures}
Game Masters who set their adventure in the \textit{Neverfields} can use this plot hook for inspiration.

\subparagraph{Seed of Life (mid-level)}
Heroes of Tal'Dorei who have made themselves known to \textit{Syngorn} may suddenly find themselves contacted by \textit{Verdant Lord Celindar}, leader of \textit{Syngorn's} armies. His report is grim: the warding trees of the \textit{Verdant Expanse} no longer protect the city; in fact, some have turned against it. The treants that once defended the elves have grown pale and sickly, and are now attacking the elvenguard. \textit{Celindar} beseeches the characters to meet with his counterpart in Lyrengorn, Skyblade Hali Instralidar, who has been surveying the ice fields of the north for centuries. Why? Because elven legend speaks of the Seed of Life, the last bloom of an ancient \textbf{treant} named Cedargaunt, which may restore the withered treants' senses. The Seed of Life is worn on the neck of a \textbf{frost giant} jarl's pet \textbf{remorhaz} within the \textit{Neverfields}.

\subsection{Othendin Pass}
This massive valley carved into the southwest section of the \textit{Cliffkeep Mountains} was, in ancient times, the territory of \textbf{stone giant} clans that warred with the dwarves of \textit{Kraghammer}. After the arrival of humanity on Tal'Dorei's shores, the people of \textit{Kraghammer} and the new settlers of \textit{O'Noa} banded together and pushed deep into the mineral-rich valley.

The stone giants fought back, but they were overwhelmed by the might of their combined foe. They retreated to their caves higher into the mountains to the northeast, and have rarely been seen by humanoids in the centuries since. This allowed the forces of \textit{Emon} to claim the resources and land of the Othendin Pass for its people, establishing \textit{Fort Daxio} to defend it.

However, the population of deadly landsharks—creatures called bulettes—native to the region has been on the rise. Some say it's because the giants who hunted them were driven away, but since the giants have been gone for a century, few can understand why the landsharks are only taking over now. The miners, gatherers, and passing travelers do their best to cull the \textbf{bulette} population both on their own and with the help of hired adventurers, but they can only ever hope to keep the deadly monsters at bay.

\subsubsection{Othendin Pass Adventures}
Game Masters who set their adventure in \textit{Othendin Pass} can use this plot hook for inspiration.

\subparagraph{Bite the Bulette (mid-level)}
A druidic order called the Stoneherders, made up mostly of stone giants and \textit{goliaths}, have taken up residence in an important choke point of the \textit{Othendin Pass}. They claim that they are the Stoneherders, a great druidic order that keeps the balance between civilization and nature in the north. \textit{Fort Daxio's} military presence in the pass has kept them away for several generations, but an unusual rise in the \textbf{bulette} population has forced them to return.

The Stoneherders' plan to thin the \textbf{bulette} herd requires the characters' aid in rounding up the deadly creatures; once they have been assembled, the Stoneherders will tame as many as they can and move them elsewhere in the \textit{Cliffkeep Mountains}. Those that can't be tamed will be relocated to other parts of the mountains, where they can be hunted again.

\subsection{Palace of Wonder}
In the desolate and uninhabited northwestern reaches of the \textit{Cliffkeep Mountains} stands a ruined citadel. All that remains of it are crumbling walls, half-beaten to dust by the unrelenting winds of the mountains, and one forlorn tower, its turret torn asunder to reveal the mildewed, rain-fattened wooden beams within. This ruin, a relic of the early days of \textit{Issylran} settlement and a panoply of forgotten wars, is all but forgotten by the people of Tal'Dorei. Few explorers travel these wind-scarred peaks, or the storm-beaten northwestern coast of Tal'Dorei. The citadel itself is wracked by the elements, while the shoreline and the mountains beyond it are prowled by territorial elementals that brook no intrusions.

\subsubsection{Secrets of the Citadel}
The ruined citadel that overlooks this treacherous landscape is far more than meets the eye. In truth, it is the headquarters of the \textit{League of Miracles}. The power-hungry spellwrights who make up the league's upper echelons call this castle the Palace of Wonder. A crumbling but entirely illusory edifice has replaced the derelict fortress, maintained by the league to deter curiosity and avoid notice.

Beneath this arcane façade, the Palace of Wonder truly lives up to its name. The palace is made up of five wings, each of which has a dedicated \textit{teleportation circle}:

\subparagraph{Halls of Splendor}
The \textit{League of Miracles} is in the business of good first impressions. Though the Palace of Wonder is a secret to the people of Tal'Dorei, the spellwrights occasionally teleport particularly powerful mercenary mages to their home base in order to dazzle them with the castle's splendors. The palace's ground floor is made of a number of interconnected grand halls, awash in a tide of splendid, vibrant colors—bordering on garish—that crashes over all who enter.

These many grand dining, entertaining, and performance halls are gleaming spectacles of white marble, decorated by sparkling globes of magical light, and festooned with self-indulgent tapestries of the spellwrights' many dubious achievements. Any intruders here without the express permission of at least one spellwright will find themselves beset by a small army of animated armors and other constructs—including a \textbf{platinum golem} that watches over the front entrance.

\subparagraph{Bowers of Recovery}
The second floor of the palace is a verdant paradise, with vibrant plants curling around the castle's stately stone architecture. Most spellwrights have their quarters in this wing, which is protected by powerful enchantments that give life to the plants. All who enter here without a spellwright escort are viewed by the living plants as a tasty snack.

This wing of the palace also includes a three-tiered conference chamber where all twenty-four spellwrights can gather to discuss matters of import. In the center of the chamber is a pedestal, upon which sits a humanoid mask of gleaming mithral. When the \textit{Wonderworker}, the league's mysterious leader, wishes to deliver a proclamation to the spellwrights, their voice emanates from this mask.

\subparagraph{Towers of Farsight}
The pinnacle of the palace is made up of four sky-scraping towers that gaze out over the \textit{Cliffkeep Mountains}—and all of Exandria. These private studies are the personal quarters of four of the league's chief spellwrights, and they are equipped with all manner of \textit{scrying} crystals and tools of clairvoyance, clairaudience, and far-speech.

\subparagraph{Dungeons of Despair}
Woe betide the prisoners of the \textit{League of Miracles}, who are left to rot in the cells beneath their palace. The warden of this dungeon is an infallible \textbf{mage hunter golem}, and each cell is equipped with shackles that inhibit the spellcasting capabilities of anyone in their grasp. Instruments of unspeakable arcane torture litter these dungeons, capable of instilling endless pain, and even extracting the very soul from one's body. The shackled spirits of those who have died by this last vile torment serve eternally as the guardians of the palace dungeons.

\subparagraph{Cradle of the Adranach}
Winding about the palace dungeons are a network of twisting passages illuminated by glaring, magical lights. These forking tunnels shift daily, and the secret to navigating them is a mark visible only under the power of a \textit{true seeing} spell. Within these passages are a dozen laboratories stocked with the tools needed for the spellwrights to dream up new potions and spells. Some they sell at extravagant cost to the governments and militias of towns and cities across Tal'Dorei—but the best they keep for themselves, to wield against the factions they manipulate from the shadows, should they ever grow unruly.

At the deepest point of the tunnels is the Cradle of the Adranach, a unique arcane laboratory where all spellwrights are initiated. Here, they craft an \textbf{adranach}, a construct of mithral and astral energy that serves as a loyal servant—and that unmistakably marks its creator as a spellwright of the \textit{League of Miracles}.

\subsubsection{Gathering the Spellwrights}
The \textit{Palace of Wonder} is almost never empty, for at least one of its twenty-four spellwrights must always remain within its walls to protect it. Many more spellwrights frequently return by means of the palace's \textit{teleportation circle}, or upon the back of their personal adranachs, to take advantage of its magical equipment and nigh-impenetrable defenses.

At times, all spellwrights are requested to return to the palace at the behest of one of the four chief spellwrights, but unless one's presence is demanded by name, these summons are not mandatory. The one summon which cannot be refused, however, is the annual calling of the \textit{Wonderworker}. The nameless, faceless leader of the \textit{League of Miracles} holds a near-prophetic sway over the spellwrights that cannot be understated. When all twenty-four spellwrights are gathered within the conference chamber, the \textit{Wonderworker} intones a slew of cryptic yet unquestionable commandments that all of their followers must interpret and obey to the best of their ability.

\subsection{Pools of Wittebak}
Wittebak was once a peaceful city of three thousand or more rock gnomes, a safe haven for the small ones that was built into the mountain rock and heated by a series of geothermal pools. The Wittebak gnomes weathered great wars by not getting involved, investing in contraptions and laughter rather than wars and empires. They found joy in inventing complicated ways to simplify their daily lives. Yet, for all their efforts, trouble eventually came to them.

Around four hundred years ago, a displaced band of hill giants found and luxuriated within Wittebak's hot springs and made a homestead on their banks. It is said that Wittebak was founded underneath the full light of Ruidus, for by pure misfortune, a flaw in the stonework of the gnomish city cracked under the force of the giants frolicking within the pools. Worse still, the cave-in did more than simply crush a few buildings beneath. The seismic force of thousands of tons of rock collapsed into the caverns, cracking open buried pockets of volcanic gas.

Over a thousand gnomes suffocated that night, or were incinerated in isolated gas explosions. The survivors fled into the mountains, assailed by hungry giants as they ran. The gnomes have found unlikely refuge in \textit{Kraghammer}, but Wittebak has been abandoned entirely, choked by caustic gas and locked in a terrible moment in time. The ten-times great-grandchildren of the giants that caused the gnomes to flee now call the hot springs their ancestral home. To them, it is named Gundergrunch, and they flourish within their unchallenged territory.

Beneath them, the ruins of Wittebak are stratified into three distinct tiers of death and despair.

\subparagraph{Tinkerer's Tier}
This level is relatively unharmed, and the full glory of \textit{Wittebak's} tinkering shops can be seen here; some even include marvels of technology unseen for centuries. Creatures that need to breathe can survive the gases in the widest, uppermost level for a number of hours equal to their Constitution modifier, and take 3 (1d6) poison damage each minute after that time. The hill giants that live within have adapted to the gas and do not fear it.

\subparagraph{Homesteads}
The air in the middle level is heavy with gas, and the quaint gnomish homes on this level are filled with the decomposed remains of those too slow to escape. The entire area is lightly obscured with clouds of gas, and tiny streams of hot water leak from the walls, tumbling into the level below.

\subparagraph{Engine Core}
The bottommost tier is a mechanical marvel, a titanic wheel of rusted iron and shining steel with no known purpose. This level is completely flooded and filled with merrow. It is also the lair of a \textbf{sea hag} named Mauvlettir.

\subsubsection{Pools of Wittebak Adventures}
Game Masters who set their adventure in \textit{Wittebak} can use these plot hooks for inspiration.

\subparagraph{Getting into Gundergrunch (mid-level)}
Just getting into \textit{Wittebak} is an adventure unto itself. It's nearly impossible to fight through the hundreds of hill giants in the hot springs and the Tinkerer's Tier—and even if it weren't, such a fight would be a callous slaughter of countless innocents. Stealth is an option, for the giants take little heed of folk smaller than them. Diplomacy is another option; Gondalorr, prince of the giants, has a tender heart and loves the idea of repatriating \textit{Wittebak's} gnomes. His sister Gondabrau, on the other hand, is cruel, and would like nothing more than to pop wriggling gnomes into her mouth as a snack. Cunning adventurers could sneak in with Gondalorr's aid, or even cause chaos in the giant town by playing the siblings against each other.

\subparagraph{The Truth of the Disaster (mid-level)}
Was it truly just by chance that \textit{Wittebak} was destroyed? Perhaps there is truth to the rumors that Ruidus, the Moon of Ill Omen, shone brightly above \textit{Wittebak} on the night of its founding. The \textbf{merrow} that lurk in the Engine Core know nothing, but the \textbf{sea hag} who lairs there has an inkling. Delve deeper, and the \textbf{aboleth} and its aberrant servitors know more still—a secret that dates back to the \textit{Founding} itself.

\subsection{Reaching Bluff}

\begin{description}
\item[Village:]
Population 324 (53\% humans, 27\% dwarves, 3\% dark elves, 3\% earth elemental ancestry, 14\% other races)

\end{description}

\noindent In a small, fertile valley in the northwest of the \textit{Cliffkeep Mountains} is a small village founded by Tal'Dorei homesteaders just ten years ago. This simple settlement is inhabited by hardy folk seeking new places to mine gold and \textit{iron} from the earth. Many unpleasant exiles from \textit{Kraghammer} have joined them, so while this village has a charming, rustic atmosphere, few travelers stay for long. The hostility toward outsiders is palpable to all.

This insular attitude has only grown more hostile following the discovery of a gleaming dodecahedron about the size of a human head. One of the miners who unearthed it found that her pickaxe barely scratched its surface—and when struck, it began emanating a radiant, internal golden glow. This mystical object was returned to the village foredwarf, and before the day was out, all the villagers had learned of it. The people of Reaching Bluff call it the The Lost Beacon of Unknown Light, and worship it as a god.

The settlers protect their mysterious beacon with cultish fervor, and Foredwarf Daz, a deceitful dwarf \textbf{veteran} with earth elemental ancestry, is willing to murder anyone who tries to probe into their business.

\subsection{Serpent's Head}
\textit{Thordak} only flew to the \textit{Cliffkeep Mountains} once during his occupation of \textit{Emon}. No one living knows why the nearly invincible \textit{Cinder King} left his prized city to attack a single, isolated mountaintop village, but the ruins of \textit{Serpent's Head} prove that he did.

\textit{Serpent's Head} is a squat mountain in the southern reaches of the \textit{Cliffkeep Mountains}. Its people were mostly dwarves and humans, miners and farmers. They had no militia to speak of, and most had already evacuated to \textit{Kraghammer} when news of \textit{Emon's} fall reached them—but not everyone.

The village that surrounds \textit{Serpent's Head} Mountain is an ashen ruin. The immolated corpses of its people still lay buried beneath a thick layer of ash, and the mountain itself was transformed into a small volcano by \textit{Thordak's} magic.

\subsubsection{Serpent's Head Adventures}
Game Masters who set their adventure in \textit{Serpent's Head} can use this plot hook for inspiration.

\subparagraph{Factions of Flame (any level)}
Characters who wish to plumb the depths of \textit{Serpent's Head} for treasure—or perhaps to find an item or character connected to the greater story of your campaign—must contend with the creatures that have made it their home. Mindless cinderslag elementals prowl the shattered battlements of the fortress that once stood atop the mountain. Earth tremors have cracked the keep's foundation, revealing deep magma caverns filled with diamonds—precious and infused with chaotic magic. Occasionally, tiny rifts to the Elemental Plane of Fire flash open within the depths, allowing groups of \textbf{azer} and salamanders to wriggle through. The two factions have made camp in the lower caverns, and both struggle for a foothold in the Material Plane.

\subparagraph{Hot-Headed Wanderers (high-level)}
The \textit{Cinder King's} incursion into Tal'Dorei attracted elementally infused monsters from deep beneath the \textit{Cliffkeep Mountains} to the surface. Magma landsharks and their young are the worst of these monsters, and even decades after the \textit{Cinder King's} defeat, they continue to prowl the mountains from their nests beneath \textit{Serpent's Head}, devouring entire villages unless brave adventurers stand in their way.

\subsection{Terrah}

\begin{description}
\item[Village:]
Population 838 (61\% humans, 19\% dwarves, 5\% gnomes, 15\% other races)

\end{description}

\noindent Sheltered within a cauldron-like valley, miles deep into the northern mountain range, lies the home of the Earth \textit{Ashari}. The people of \textit{Terrah} have long stood watch over the exposed rift to the Elemental Plane of Earth, protecting the ever-quaking ground of this crumbling valley. Those without druidic power dig long trenches around their village to protect their people against rockslides that roll off the mountains, while a handful of mighty druids use their magic to reshape the valley and rebuild damaged structures.

Living off hunted game and scavenged fungi, the \textit{Terrah} people are stocky warriors of stubborn mind and immense pride and loyalty. Jewelry made of precious gemstones is regarded as a sign of station and respect. The Earth \textit{Ashari} supply \textit{Emon} with precious and semi-precious stones from the rift in exchange for essential goods and, occasionally, protection from the creatures of the mountains.

The peaks that surround the rift frequently seem to shake with rage, and from their angry slopes rise creatures of living stone, imbued with life by elemental spirits from the rift. The people of \textit{Terrah} sleep in shifts to keep watch for any sudden changes in seismic activity, lest the fury of the mountains undo all that \textit{the Ashari} have created.

\subsubsection{Government}
\textit{Terrah} is governed by an \textit{Ashari} archdruid given the honored title of Heart of the Mountain. This archdruid has completed their \textit{Aramente} and earned the respect of the four elemental enclaves of \textit{the Ashari}. Pa'tice, the current Heart of the Mountain, is a human man well over one hundred years old. Though tenacious and ornery, he is a venerable warrior with boundless appreciation and support for those who prove themselves honorable and self-reliant. Rumor has it that Pa'tice is seeking a worthy druid who can undertake the Aramente and succeed him as Heart of the Mountain.

\subsubsection{Terrah Adventures}
Game Masters who set their adventure in \textit{Terrah} can use these plot hooks for inspiration.

\subparagraph{Pestering Pa'tice (any level)}
Rumor has spread throughout the \textit{Cliffkeeps} that the Earth \textit{Ashari} are seeking a new Heart of the Mountain. In response, dozens of upstart druids with delusions of grandeur have flocked to \textit{Terrah}, seeking to impress the current leader of the Earth \textit{Ashari}. Pa'tice can't be bothered with this pageantry, and quietly hires a group of adventurers—the characters—to disperse these druids. If there is a druid in the party, perhaps their actions will impress the Heart of the Mountain and set them on the path of the Aramente.

\subparagraph{Dependable as Stone (high-level)}
The \textit{Terrah} people are insular and dutiful, and rarely take sides in any conflict. However, sometimes times of great peril require that even \textit{the Ashari} go to war. Whatever the conflict, the characters must convince the \textit{Terrah} to join their cause. Pa'tice meets with them and grimly states that the \textit{Terrah} cannot aid them because a terrible new threat has come through the portal. The Stonesight Council, four medusas aided by two gorgons, now command the earth elementals descending from the peaks around \textit{Terrah}. Pa'tice is a practical leader—he cannot lead his people to fight a threat to the outside world while danger lurks on their doorstep. "Eliminate this threat to Terrah," Pa'tice says, "and perhaps we can join your cause."

\subsection{Umbra Hills}
The \textit{Battle of the Umbra Hills}, the climactic battle of the \textit{Scattered War} and the death knell of the Kingdom of \textit{Drassig}, transformed the floral, sun-dappled Emerald Highlands into a blasted wasteland. A tide of demonic blood spilled across the highlands like fire, cursing the land forever. Generations have passed, but the heather that once grew on these hills is as black and burnt as the day of the battle. No animals live here, and the only plant that grows in the \textit{Umbra Hills} now is shadegrass. Its ash-gray stalks are dry and unfilling, but its unique, acrid flavor has drawn the interest of spice traders worldwide, propelling wealthy merchants to hire armed escorts for their spice-pickers.

Though no animals live here, danger still lurks within the \textit{Umbra Hills}. Undead soldiers rise from the grass when the silvery moon Catha is high, and the ruins of ancient \textit{Drassig} war camps seem to draw demons of all types, still bound by the contract between \textit{Trist Drassig} and the Demon Prince of Indulgence. Pathetic dretches gambol about in the ruins' long shadows, and rumors say that, on nights when the moon Ruidus is full, even mighty \textbf{balor} demons prowl the hills in search of mighty souls to corrupt.

\subsubsection{Umbra Hills Adventures}
Game Masters who set their adventure in the \textit{Umbra Hills} can use this plot hook for inspiration.

\subparagraph{Sword of the Demon Prince (high-level)}
In a dream, a skilled warrior within the party continually sees visions of a noble, kingly sword embedded in the ground, surrounded by shadows and demons. "Free me," it whispers. "I need your help...I am in the \textit{Umbra Hills}." The sword is not as it appears; though the Demon Prince of Indulgence was defeated at the \textit{Battle of the Umbra Hills} and banished to the Abyss, his accursed greatsword still lingers. The last shard of his malevolence in this world, \textit{Graz'tchar, the Decadent End} rests in the center of the valley, its tip embedded in the earth. The alluring, androgynous \textbf{marilith} demons that guard it pretend to be the ghosts of great knights, protecting the blade of a lost Sovereign of Tal'Dorei.

\subsection{Grey Valley}
North of the \textit{Umbra Hills} lies a wide valley of ill omen and dark history, named for the forest of petrified trees that fill the majority of the valley's forest floor. During the \textit{Scattered War}, \textit{Trist Drassig} obsessed over the possibility of losing and pledged his service to the Demon Prince of Indulgence in exchange for power to smite his foes.

The bargain accepted, the demon prince sent his forces to this valley to aid Trist in the \textit{Battle of the Umbra Hills}, but ultimately both fell to the incredible strength of the rebelling forces. The release of terrible energies exploded across the battlefield. Now left cursed and lifeless, the valley remains a warning against dealing with demons—and a dangerous locus of darker entities that are drawn to the ruined place.

\subsubsection{Grey Valley Adventures}
Game Masters who set their adventure in the \textit{Grey Valley} can use this plot hook for inspiration.

\subparagraph{Heart of the Wraithroot (high-level)}
At the dead center of the \textit{Grey Valley} is a petrified tree unlike any other. Knots in its stony bark gaze like eyes at its withered domain. This \textbf{wraithroot tree} is the eternal guardian of a dark treasure. Rumors abound about what could be secreted away behind this vile guardian—a relic of a demon prince, the crown of \textit{Drassig}, or an artifact of the \textit{Betrayer Gods} unknown to mortal minds? The true nature of this artifact is up to the Game Master, but whatever it is, rumors of its existence have lured countless adventurers to their deaths.

\section{Stormcrest Mountains}
Spanning Tal'Dorei's southern reaches, buttressed against the \textit{Verdant Expanse}, stand the massive \textit{Stormcrest Mountains}. The winds of the \textit{Lucidian Ocean} wail through the storm-battered peaks, and few dare make these treacherous mountains their home. The few who do live amidst the Stormcrests are ferocious monstrosities and the \textit{giantkin} strong enough to defeat them—and those few peoples cunning enough to evade storm, monster, and giant alike. The only friendly towns anywhere near these forsaken peaks lie in the \textit{Mornset Countryside} pushing along the Stormpoint range, far beyond the influence of the \textit{Tal'Dorei Council}. These rugged folk are used to fending for themselves and scoff at the thought of returning to so-called "civilization," preferring life among the swamps and shadowed forests of their homeland.

\begin{description}
\item[Majority Faiths:]
The Arch heart

\item[Minority Faiths:]
The Moonweaver

\item[Imports:]
Livestock, lumber, grain, silver, spell components

\item[Exports:]
Lumber, gold, grain, granite, quartz, iron

\end{description}


\subsection{Ashen Gorge}
One mountain stands out among the Stormcrests. Its twin peaks rumble with fury, breathing fire and smoke into the sky to mingle with the storm clouds that surround the mountains. Lightning tears through the ashen clouds with giddy excitement, making it clear to all who behold it that this is a place of supernatural power. Cartographers in \textit{Emon} call this crater the \textit{Ashen Gorge}, but the people of the low valleys call it the Dragon's Throne, for they know what creature once lived atop the mountain. They remember when they too were under the dominion of \textit{Thordak the Cinder King}.

Today, several tenuous alliances of \textbf{lizardfolk}, \textit{dragonblood}, kobolds, and even a few young dragons trying to expand their treasure hoards battle for dominion over this ruined vale. When the mountain was the lair of the \textit{Cinder King} during his first incursion into Tal'Dorei, the \textbf{lizardfolk} and kobolds flocked to his side, worshiping his might and cruelty. Though \textit{Thordak} is slain, his influence yet lingers, and his legend draws treasure hunters and would-be tyrants to this place. Gold, jewels, and magic items from all ages of the world lie in wait within the dozen or more secret vaults of the \textit{Cinder King's} lair. When the king of dragons still held dominion here, he had thousands of loyal scaled minions at his beck and call, moving invaluable treasures into deeper and deeper holds. Beyond the main cavern of the lair, even the \textit{Cinder King} himself had to magically take humanoid form to traverse the winding tunnels his servants carved through the rock.

The descendants of \textit{Thordak's} servitors now fight an endless war over his treasure. Every year, it seems a petty new "Cinder King" rises to power within the \textit{Ashen Gorge}, but they are always dethroned before their minions can enter the ancient vaults, still locked after all this time. To steal \textit{Thordak's} greatest treasures, an adventuring party would have to either broker a peace between the warring clans or somehow sneak into the mountain undetected.

The following major factions fight for control over the \textit{Ashen Gorge}. They make and break their alliances with one faction or another, and then break their new alliances just as readily. You can use these factions to inspire adventures set in the \textit{Ashen Gorge}.

\subparagraph{Scions of Flame}
This cult of black-scaled \textbf{lizardfolk} claim to have been \textit{Thordak's} elite guard in ancient times, and are the current rulers of the gorge. They're led by a five-hundred-year-old druid named Burning Oak, and while his awesome magic keeps the other tribes at bay, his frailty makes him vulnerable to attack—and the other tribes know it.

\subparagraph{Tinysoot}
Every ten years brings a new generation of kobolds, and while the Tinysoot tribe's warriors are not the mightiest, their overwhelming numbers have allowed them to overwhelm even the greatest of the Flame Scions' champions. The Tinysoots' leader, a stubborn young Queen of Soot named Yabber Tinysoot XIX, wears around her neck the key to Everflame Crevasse, \textit{Thordak's} deepest treasure vault.

\subparagraph{Black Snow}
When \textit{Thordak's} ally, Umbrasyl, was slain atop \textit{Gatshadow}, his caustic blood seeped into the corpses of a dozen goliath scouts from the Herd of Storms. The black dragons' magic, combined with the foul power of \textit{Gatshadow} itself, transformed the \textit{goliaths} into a squadron of undead black \textit{dragonblood} warriors. Skeletal wings sprouted from their backs, and the Black Snow tribe flew instinctively to the \textit{Ashen Gorge}. They were drawn to the Ruby of Oblivion, a pitch-black ruby of immense unholy power—now stored in the Obsidian Geode, a treasure vault guarded by a legion of golems and elementals. This squadron of hateful revenants is small, but each one of them wields the strength of twenty lesser warriors.

\subparagraph{Tatterwing Kin}
Toxoshuul, Chloroxodon, and Snagglefang, a trio of young green dragons, long to plunder \textit{Thordak's} lost vaults for their own hoards, and demand tribute from the other factions of the \textit{Ashen Gorge} in exchange for their aid in battle. All three are the children of Tatterwing, a feared dragon of the Kirmont Valley, and even if they find themselves on opposite sides of a skirmish, they never fight one another, only annihilating the pathetic humanoids beneath them with their toxic breath and razor-sharp talons.

\subsection{Dreamseep Marshlands}
East of the Stormcrest Range, within the Kirmont Valley, is the sprawling, fetid Dreamseep Marsh. The perpetual rain rolling off the Stormcrests' enchanted slopes has transformed this once-lush forest into a morass of rotting vegetation, sulfurous mud, and gnarled, weeping trees. So many cruel murders have been committed within the Dreamseep that the very land is cursed, forsaken by the gods. Negative energy pools there like water, and the dead drink deep of it. The marsh is populated by more than just the walking zombies of killers and their victims—awful amalgamations of dozens of bodies, both humanoid and bestial, are birthed within the Dreamseep's putrid womb.

Somewhere in the middle of this accursed realm is a sinkhole that plunges deep beneath the earth. At its deepest point is the Tomb of Udah, a fabled necropolis of countless chambers, littered with traps and treasures that have claimed the lives and imaginations of untold hundreds. Every death within Udah's accursed walls only adds to the legions at its dread master's command; characters who explore its tunnels face not only undead warriors in armor from a bygone millennium, but steely warriors in armor from every epoch in Tal'Dorei history.

The small community of Bronbog holds the only lights of civilization within the shadows here, and those that continue to thrive against the oppressive swamp are hearty and stubborn folk..

\subsubsection{Dreamseep Marshlands Adventures}
Game Masters who set their adventure in the Dreamseep Marshlands can use this plot hook for inspiration.

\subparagraph{Skull of Udah (epic-level)}
Two rumors persist of the Tomb of Udah. First, it is said that within the tomb's deepest sanctum is a three-eyed golden skull, and that in each socket is inlaid a ruby that can grant any wish. Second, the locals say that, unlike the mindless undead of the \textit{Dreamseep} above, the abominations in the necropolis below are commanded by a hideous master. Little do these rumor-spreaders know that the skull and the master are one and the same. Udah the Undying is a \textbf{lich} of tremendous power, devouring the souls of all who die within the \textit{Dreamseep}. His disembodied spirit controls the dungeon itself, and only retreats into the skull—the last remnant of his physical form—when it is in danger.

\subsection{Bronbog}

\begin{description}
\item[Village:]
Population 683 (53\% humans, 18\% half-elves, 11\% goblinkin, 9\% halflings, 9\% other races)

\end{description}

\noindent If there ever was a civilization that called the \textit{Dreamseep} home, Bronbog is all that remains of it. Its buildings are made from planks of waterlogged bogwood, and the only sign of its previous greatness is the ring of stone pillars that once supported the Temple of the The Dawnfather. Even though the temple to their patron god has crumbled, the people of Bronbog keep the faith; their belief in the sun god's afterlife is their greatest comfort against their meager lives. Yet, for all the gloom that surrounds the \textit{Dreamseep Marshlands}, the Bronboggi maintain a determinedly optimistic disposition. There's a saying in the village: "If you don't feel the storm, you can't know there ain't sun."

The only reason anyone north of the \textit{Stormcrest Mountains} knows of Bronbog is because of queenscap, a rare swamp fungus that serves as a reagent in creating potent potions of superior healing. It's been harvested almost to extinction in the more accessible \textit{K'Tawl Swamp}, and Tal'Dorei's alchemists will pay good money for a shipment.

The biggest building in Bronbog is its tavern, the Queen's Knight, which has about six beds in drafty rooms to accommodate the town's very few visitors.

\subsubsection{Bronbog Adventures}
Game Masters who set their adventure in Bronbog can use this plot hook for inspiration.

\subparagraph{Morning Glory (mid-level)}
On a queenscap-gathering expedition to Bronbog, an unusual storm forces the characters to stay the night. As they spend time in the swamp town, vine-choked stone obelisks slowly emerge from the murky waters, emanating an unearthly light. Some villagers immediately fall under the obelisks' sway, claiming they are "heralds of the Dawnfather, the rising sun," but all others are disturbed by their neighbors' feverish worship of the stones.

The true power behind these obelisks is a misguided \textbf{deva} that calls himself Eclipse. He seeks to create a brainwashed, perfectly loyal Army of the Dawn to fight against the The Dawnfather enemies. He is posing as Father Larkylai, the town's affable priest of the The Dawnfather, and conducts services from within the town's tavern.

\subsection{The Frostweald}
Along the northern base of the Stormcrests is a forest locked in perpetual winter, cursed to eternal cold by the invasion of \textit{Errevon the Rimelord} during the \textit{Icelost Years}. Within the forest, ponds magically freeze into perfect mirrors, reflecting the snowy sky above, and the snow-draped trees shelter families of fey hiding from enemies in the Fey Realm. At first, the \textit{Frostweald} seems a wonderland of crisp snow and aromatic pines and firs, yet the serene landscape belies sinister danger. Herds of basilisks roam the woods, and travelers who encounter mysterious snow-covered statuaries or copses of petrified trees must think fast or join this grisly statuary—permanently.

\subsubsection{Frostweald Inhabitants}
The only humanoid peoples who live within the \textit{Frostweald} are hunter-gatherers called the Shivergut. Many of their people are orcs who have lived within the \textit{Frostweald} for generations, but their ranks also include humans, elves, and dwarves who adopted their ways centuries ago during the \textit{Icelost Years}. Most are helpful to travelers, if reserved. Others, however, are dour, selfish isolationists who have been apart from other folk for so long that they care more for the thrill of bloodshed than a warm conversation over a crackling fire.

The \textit{Frostweald} is also home to a host of benevolent fey, most of them pixies or dryads, using the forest as a safe haven away from the Fey Realm. If pressed for answers on why they left the Fey Realm, their answers are cryptic and dissatisfactory, but always ominously suggest a "Great Shadow" has descended upon their lands—perhaps even a host of warriors from the Plane of Shadow? A nymph named Arethusa watches over a cluster of three mirror-like pools that form a pathway into the Fey Realm, each to a different Archfey's forest. She is suspicious of all mortals, and both her trust and a favor are required to gain passage.

\subsubsection{Ruins of the Frostweald}
Many forgotten obelisks of the The Knowing Mentor rest secretly beneath the ice and snow. These lost beacons were created to guide the faithful of the The Knowing Mentor to the \textit{Cavern of Axiom}, a sacred place guarded since time immemorial by a sphinx dedicated to the The Knowing Mentor. If any living creature knows why these ancient monoliths reside here when no other ruins of the The Knowing Mentor have been discovered, they have kept the knowledge secret from the other scholars of the world. Over thirty years ago, the half-orc \textit{Emonian} archaeologist Jorlund Vohr discovered a cache of Ioun stone buried here and brought them to the \textit{Alabaster Lyceum} of \textit{Emon}, but his research and findings were stolen the week after he returned. Vohr has since stated he "got over it," and is back to work on new research on \textit{Visa Isle}.

\begin{DndSidebar}[float=!b]{The Silver Tablet}
\begin{blockquote}
{    \small\itshape
    Ye who seek my boundless wisdom,\\
    Pass into the white trees' kingdom.\\
    Return to the shrine Mentor's eyes,\\
    And follow the light in the skies.\\
    I am the Beast of the Heavens;\\
    Thou shalt kneel before my presence,\\
    Lest the wisdom thou seekest\\
    Be lost to the cowardice of your weakest.
}
\end{blockquote}

This poem was discovered by Jorlund Vohr on his first expedition to the \textit{Frostweald}, on a silver tablet written in flowing Celestial script. Upon examination, scholars at the \textit{Alabaster Lyceum} discovered the language of the heavens is so sublime that its poetry rhymes even when translated into any language. The tablet was stolen later that week, along with the rest of Vohr's findings, but Jorlund keeps a parchment-and-charcoal copy of the tablet in a chest beneath his bed in \textit{Emon's} \textit{Erudite Quarter}.

\end{DndSidebar}

\subsection{Cavern of Axiom}
The Cavern of Axiom is a lost shrine to the The Knowing Mentor, spoken of only in riddles. Hidden by ever-shifting illusions, the entrance to this cavern opens only to those who are expected by fate. Its entrance looks different to every group that finds it, from the imperious to the humble, but it appears simply as a heavy snowdrift to those not fated to open its doors. Within its shifting facade are dangerous trials designed to test the will and mettle of those who seek the infinite knowledge of the The Knowing Mentor. The chambers plunge deeper into the rock beneath the mountain, while half-eroded murals and strange puzzles evaluate any wanderer who seeks an audience with the keeper of the cavern, an ageless \textbf{androsphinx} named Kamaljiori. The final challenges presented by the sphinx alter from subject to subject, and failure banishes the petitioners from the cavern, barred from ever returning.

\subsection{Ruhn-Shak}

\begin{description}
\item[Small City:]
Population 6,670 (80\% elves, 12\% dwarves, 5\% gnomes, 3\% other races)

\end{description}

\noindent If you find carved arches and steel gates in the mountain slopes, do not rest there, no matter the cold. There are no grand dwarf-halls among these forsaken peaks.

Deep below the surface world, countless caverns and tunnels wind into regions where light finds no purchase. It is here, beneath the \textit{Stormcrest Mountains}, that the largest subterranean society of the The Spider Queen devotees maintains its tyranny. Hidden entrances and tunnels riddle not just the mountain range, but the edges of the \textit{Dreamseep}, as well as darker regions of the \textit{Verdant Expanse}, granting raiding parties quick and easy access to the surface under cover of night—slaughtering many, capturing the rest, and taking the spoils as gifts to the The Spider Queen.

These tunnels, in their twisted network, are easily collapsed and reopened through the use of "Pit Witches," druids who master the art of rock and dirt manipulation, making it near impossible to give chase. The \textit{Wardens of Syngorn} are ever seeking a way to find and destroy this heart of the The Spider Queen power in Tal'Dorei, but the sly and cunning elves of Ruhn-Shak have yet to meet their match.

\subsubsection{Ruhn-Shak Adventures}
Game Masters who set their adventure in Ruhn-Shak can use these plot hooks for inspiration.

\subparagraph{Dynastic Intervention (any level)}
While in \textit{Syngorn}, the characters are called by \textit{Ouestra, Voice of Memory}, to meet with guests from \textit{Wildemount}. They are emissaries of the Kryn Dynasty, a society founded by dark elves who fled the The Spider Queen grip and now live on the surface. These Kryn envoys seek to undertake an expedition into \textit{Ruhn-Shak}, to try to use the radiant power of their mysterious god, the The Lost Beacon of Unknown Light, to free their kin from the The Spider Queen evil influence. The characters are offered a fine reward by the \textit{Wardens of Syngorn} to protect these Kryn priests on their mission.

\subparagraph{A Plea from the Dark (mid-level)}
The characters are drawn into the underground tunnels around \textit{Ruhn-Shak} by a psychic message: "I have displeased the Spider Queen; there is nothing I can do but beg for salvation—can anyone hear me?" The message comes from Ry'denyan, a \textbf{drow} of \textit{Ruhn-Shak} granted telepathic powers by an unknown aberration. This unexpected gift allowed him to break free of his society's brainwashing. He tried to escape, and was exiled into the tunnels, where he will surely be devoured by driders unless the characters save him.

\subsection{Wrettis}
A relic of the \textit{Age of Arcanum}, Wrettis is the ruined tower of a powerful mage, one who was driven mad by the seductive whispers of beings from the beyond. What few legends survive of the mage of Wrettis say he was known as Clemain Astural, the Sight Shephard, and that he was a powerful arcanist who peered into a realm beyond the planes in search of power to end the war that ravaged his world. He found it, and thought it would serve him.

He was wrong. Astural's sanity crumbled, but his power only grew, fueled by an entity he called the The Sightless. As devastation crept across southeastern Tal'Dorei, heroes of the land rode to Wrettis to end Astural's chaotic reign, destroy his tower, and bury his corpse in the rubble. Wrettis is now a moss-coated ruin, but some chambers within and below the tower still hold secrets of the mad mage, as well as some of his creations.

\begin{DndSidebar}[float=!b]{The Astural Scrolls}
Clemain Astural covered thousands of pages of parchment with his writings, ranging from alchemical formulas to powerful spells to esoteric star charts...but his later notes descend into paranoid, half-crazed ramblings. One of the heroes that killed Astural, a wizard named Atz Yuminor, took as many of the scrolls as he could find and brought them back to his house in the \textit{Dividing Plains}. Yuminor's bloodline continues in Estella Ladimar, headmistress of the \textit{Westhall Academy in Westruun} in \textit{Westruun}, and many of the Astural Scrolls are kept secretly by the \textit{Scions of Yuminor in the academy's observatory}.





\end{DndSidebar}

\section{Rifenmist Peninsula}
South of the \textit{Stormcrest Mountains} and the \textit{Verdant Expanse} are a number of autonomous communities that live outside the influence of the Republic of Tal'Dorei and its constituent city-states. A number of plans to expand into \textit{Rifenmist} have been brought to the \textit{Tal'Dorei Council}, but strong opposition from the \textit{Syngornian} delegation and the Master of Defense have consistently buried these plans before they can come to fruition.

A number of communities exist on the \textit{Rifenmist Peninsula}. Some are lost remnants of the \textit{Scattered War} that chose to flee \textit{Drassig's} rule to the isolated and unpopulated \textit{Mornset Countryside} rather than fight against his tyranny. These are the communities that some \textit{Emonian} merchants wish to bring into the fold of Tal'Dorei proper. Farther to the south (and occasionally roaming north into Mornset) are the \textit{Orroyen elves}, a collection of nomadic tribes. The \textit{Orroyen} and the elves of \textit{Syngorn} see one another as family, and the \textit{Syngornian} delegation refuses to allow the council to interfere with their cousins' way of life.

\begin{description}
\item[Majority Faiths:]
The Changebringer

\item[Minority Faiths:]
The Wildmother

\item[Imports:]
Arms, armor, stone, masonry, spell components

\item[Exports:]
Spell components, produce, rare metals
\end{description}

Beyond the Mornset people and the \textit{Orroyen} are a number of other small, individualistic settlements that live in fear of the imperial power that commands Rifenmist: the \textit{Iron Authority}. Tal'Dorei's Master of Defense warns the council of provoking this slumbering empire, for they are preoccupied with their unholy conquest of the \textit{Rifenmist Jungle} in the name of their patron deity, the The Strife Emperor. A number of \textit{Orroyen elves} and other people of \textit{Rifenmist} long for \textit{Emon's} aid in breaking the grip of the \textit{Iron Authority}, but fear that allowing the Republic of Tal'Dorei to gain a military foothold in \textit{Rifenmist} would simply be trading one imperial master for another.

Beyond the political struggles of this realm are terrors only true heroes can survive. Enormous, ancient beasts stalk the vine-twisted paths of the jungle that consumes the majority of the peninsula's coastline. An expanding morass of fungus spreads from a mysterious, corrupt source along the eastern depths of the underbrush, while a secret society of naga-worshippers bring bloody offerings to their snake-queen.

\subsection{Beynsfal Plateau}
Where below there is life and freedom, above there is naught but death and tyranny. In the high, southern plateaus of \textit{Rifenmist}, verdant jungle gives way to a cracked, volcanic landscape, its basalt fields interrupted only by towering structures of black iron. This hellish place is the home of the \textit{Iron Authority}, and one of the last battlegrounds of the \textit{Calamity}.

It is here that the Betrayer God known as the The Strife Emperor, aided by his legions of brainwashed goblinkin, clashed with the The Wildmother and her Free Children. Though the The Wildmother was victorious and the goblinkin were freed, the site of the The Strife Emperor defeat was reduced to rock and ash, a place where plants would never again grow.

Most of the goblinkin who survived the battle vanished into the jungle and traveled north into the rest of Tal'Dorei, leaving the shame of their connection to the The Strife Emperor behind. Centuries later, however, a number of goblinkin still remained among the towering iron of the The Strife Emperor fallen armor. One platoon of hobgoblins, the Iron Regiment, devoted themselves wholly to the The Strife Emperor, and pledged to revive his glorious legacy of conquest. They press-ganged as many of their kin as they could into armies and created the seed of an empire—one which they watered with blood.

When the Iron Regiment brought all of the Beyensfal Plateau under their rule and transformed themselves into the \textit{Iron Authority}, their empire only continued to grow. The Iron Tide never broke, and within a decade, it had engulfed all the other kingdoms along \textit{Rifenmist's} southern coast. They were united under the cruel banner of Tz'Jarr, the Iron Emperor.

While the entirety of Beynsfal is firmly controlled by the \textit{Iron Authority}, the empire is still reaching, now locked in a decades-long war to subjugate the native inhabitants of the jungles below: the \textit{Orroyen}.

Five major city-states comprise the core of the \textit{Iron Authority}, though all swear fealty to the Iron Emperor in its capital of \textit{Tz'Arrm}, the southernmost city of the plateau. From south to north, the other imperial city-states are Rybad-Kol, city of the Forge Lords; Hdar-Fye, the necropolis of Prince Hdar; Ezordam-Haar, eyrie of the red pegasi; and Ortem-Vellak, the petrified elventree.

\subsubsection{Beynsfal Plateau Adventures}
Game Masters who set their adventure on the \textit{Beynsfal Plateau} and the surrounding mountains can use this plot hook for inspiration.

\subparagraph{The Wayward Bloodline (mid-level)}
Tz'Jarr has a single heir to his throne: his \textbf{hobgoblin} child Zuun'dak, nearly of age and currently being groomed to become a great general. However, Zuun'dak has grown up with trepidation regarding the brutality of his society, and after being briefly captured by an outlying \textit{Orroyen} scouting party and learning of the gentler world abroad, he's begun to long for escape and reformation. When the party stumbles upon a disguised and fleeing Zuun'dak, they must decide whether to turn over such a valuable bargaining tool to the \textit{Orroyen}, or aid and groom the defecting heir as a rebellious force against his father.

\subsection{Mornset Countryside}
After a long journey south from \textit{Syngorn}, or north from the \textit{Rifenmist Jungle}, thick foliage gives way to rolling hills and winding rivers. No great cities are to be seen here, only an untouched wilderness pockmarked by fortified towns. The Mornset Countryside is home to many people whose ancestors fled the rule of \textit{Drassig}. Most of these people value their privacy—especially from \textit{Emonian} landscape painters who enrich themselves by romanticizing and exoticizing their idyllic way of life.

The Free Folk of Mornset, as they call themselves, include a wide variety of people who have spurned the ways of their forebears. The bulk of these people are the descendants of refugees from a long-ago war, but a number of Tal'Dorei people who long for freedom from the pressure of modern life have uprooted their lives and journeyed to Mornset to start anew. Likewise, some \textit{Orroyen elves} who dislike the nomadic lifestyle find homes in the small communities of Mornset.

Beyond those who follow their hearts to this realm, Mornset is also home to a number of criminals fleeing persecution. Most Free Folk don't ask questions about newcomers' pasts, and welcome them regardless. However, if these particular outlanders don't leave their past behind and continue to harm others, or bring bounty hunters (or worse) in their wake, every village in Mornset has a well-trained militia that will eagerly remove dangerous folk from their midst.

Bound by the necessity of survival, the people of this land share a common code of honor called the Fundaments. This code holds all to a baseline of respect for all travelers and outcasts, and codifies ancient principles of hospitality. Each denizen of the countryside looks out for their neighbor, and if a person in need arrives upon one's doorstep, the hospitality of food and shelter must be provided.

It is understood that the Fundaments may be taken advantage of by ruffians and power-hungry upstarts. Nevertheless, this code of honor has served the Free Folk well for centuries, and anyone who scorns it is looked upon with disdain and mistrust. Most of the Mornset Free Folk are humanoids, but many free settlements have alliances with the local hill giants, cyclopes, and ogres that help protect them against the \textit{Iron Authority} to the south.

\subsubsection{Mornset Adventures}
Game Masters who set their adventure in Mornset can use this plot hook for inspiration.

\subparagraph{Anger Everburning (low-level)}
A number of residents of the uncharted village of Heldenfaire have been unexpectedly—and suddenly—flying into a rage both mindless and destructive. Many within the village are confused, fearing a plague of violent madness, and they seek adventurers (especially parties with clerics or skilled healers) to save them. Investigation by the party reveals that each individual driven to these fits is connected to the town alchemist's dark past and a failed herbalism experiment from ten years ago, when she lived in \textit{Kymal}. If the characters wish to find an antidote, they must explore the alchemist's sinister laboratory in the hills.

\subsection{Byroden}

\begin{description}
\item[Town:]
Population 8,050 (45\% humans, 21\% half-elves, 12\% goblinkin, 8\% halflings, 8\% gnomes, 6\% other races)

\end{description}

\noindent The northernmost and largest town in the \textit{Mornset Countryside}, \textit{Byroden} borders the \textit{Gladepools} and marks the first stop for any traveler on their way to the \textit{Rifenmist Peninsula}. It's also a hub of trade for the entire \textit{Mornset Countryside}, and even some trappers and artisans of the \textit{Verdant Expanse}. The town has a dense town center surrounded by palisades and tall wooden walls, but most of the community lives in homesteads and farmhouses spread across the surrounding fields. \textit{Byroden} is mostly self-sufficient, and has a wealth of farmers, fishers, artisans, and miners, giving them everything they need to survive. Trade provides the opportunity for some luxuries, with raw materials serving as their primary currency.

While pleasant and welcoming on the outside, this lifestyle is protected by a fierce and intense drive to defend the community's way of life at the first sign of threat or danger. Most citizens are armed at all times, trained to rise as a communal militia at a moment's notice from the war ringer, a sentinel charged with the town's protection.

Despite (or perhaps because of) the ever-present dangers that threaten \textit{Byroden}, the town holds frequent festivals, feasts, and fairs, including the highly anticipated Gem of \textit{Byroden} pageant and the midsummer Peachiest Pie bake-off. On festival days, small children roam freely within the protective town walls as their parents enjoy the chance to let loose and celebrate life in all its ephemeral glory.

\subsubsection{Hillmaw Chasm}
A few miles north of the center of \textit{Byroden} is a gaping chasm. An earthquake thirty years ago split open Deercrest Hill. The chasm is referred to officially as Deercrest Ravine, but is known to most superstitious folk in town as Hillmaw Chasm. After the earthquake, locals stayed as far away from the chasm as they could—until four years later, when a group of adventurers found vast veins of precious gold and sturdy \textit{iron} within. Soon, all people of \textit{Byroden} strong enough to swing a pick or lift a shovel made their way into the chasm.

Most miners hire an armed escort because the chasm is known to be home to grumpy earth elementals, as well as unpleasant, slimy creatures of the dark. Most serious beliefs about the chasm being haunted have dissipated over the past thirty years. Still, rumors persist, and as folklore becomes myth, some say that myth can become reality.

\subsubsection{Advent of the Cinder King}
About forty years ago, \textit{Byroden} was brought to its knees by \textit{Thordak the Cinder King}. This incursion was about two decades before the \textit{Cinder King's} destruction of \textit{Emon}, and it was stopped by an adventuring party led by \textit{Allura Vysoren}, now a senior member of the \textit{Tal'Dorei Council}. Many in \textit{Byroden} still see \textit{Allura}, Lady Kima of Vord, and \textit{Drake Thunderbrand} as heroes. The graves of the companions they lost in that fight, Sirus Kaldrem, Dohla Lorian, and Ghenn Talevesh, hold places of honor in the town cemetery.

While \textit{Byroden} has long-since been built anew, the terror of that day still lives on in its elders and adults. All but the youngest generation of \textit{Byroden} takes great caution against another such monstrous attack—as well as against any attack by the \textit{Iron Authority} to the south.

\subsubsection{Byroden Adventures}
Game Masters who set their adventure in \textit{Byroden} can use these plot hooks for inspiration.

\subparagraph{Ironbound Scouts (low-level)}
All are awoken in \textit{Byroden} by the ringing of the town's alarm bell, including the characters. As the town militia is armed for battle, the characters overhear that \textbf{goblin} and \textbf{hobgoblin} scouts of the \textit{Iron Authority} have been spotted twenty miles south of town—closer than they've been seen in nearly two decades.

Worried that this scouting party is an omen of invasion, the town's war ringer deputizes the characters and offers them a handsome sum of money to investigate the scouting camp. Their goal is simple: learn the interlopers' mission, either by interrogating their commander or by finding a written command note. And, if the characters can capture or eliminate the scouts before the town's upcoming festival, the locals will celebrate their success with all the free food and drink they can stomach.

\subparagraph{The Dreamgate (mid-level)}
After a mining expedition goes missing deep within the Hillmaw Chasm, locals begin to report vivid, communal dreams involving members of the missing party pleading for help. A half dozen mercenaries have been hired to investigate—only for them to disappear as well. Just a day before the characters are asked for help, the townsfolk of \textit{Byroden} discovered with horror that the vanished mercenaries have also appeared in their dreams. The adventurers are asked to delve into this unearthed pathway and to find both answers and the missing people.

\subsection{Niirdal-Poc}

\begin{description}
\item[City:]
Population 15,442 (18\% humans, 17\% goblinkin, 12\% elves, 12\% Loxodon, 10\% half-elves, 7\% elemental ancestry, 24\% other races)

\end{description}

\noindent Long before the \textit{Calamity}, a great civilization called the Qoniira made \textit{Rifenmist} their home. During the \textit{Age of Arcanum}, the Qoniira Tetrarchy grew from a loose coalition of towns led by arcanists, clerics of the The Wildmother, and druids who spoke with the jungle around them. These towns were inhabited by humans and the animal-bodied humanoids of the jungle. The first stones of the Qoniiran heart-city Niirdal-Poc were laid by the broad, towering Loxodon, and places for new settlements were scouted out by \textit{Catfolk} seers whose eyes glistened with starlight.

\subsubsection{Surviving the Calamity}
In time, the Qoniira civilization grew into a Tetrarchy, with four great satellite cities and the central heart-city of \textit{Niirdal-Poc}. The tetrarchs established their cities with mathematical precision. A constellation of the night sky guided them—Surrac's Shield, a four-starred diamond with a fifth star in its center. The constellation still shines over \textit{Rifenmist} and southern Exandria today. The four outer cities of Qoniira match the points of Surrac's shield precisely, and \textit{Niirdal-Poc} is aligned with its center star.

The \textit{Calamity} struck at the peak of Qoniiran civilization, yet rather than take up arms against gods and mortals, the tetrarchs opted to remove themselves from the world—not unlike the elves of \textit{Syngorn} to the north. Something happened in \textit{Niirdal-Poc}, some sort of magic that even the Qoniirans of today do not understand.

\subsubsection{Geography}
Today, only the heart-city of \textit{Niirdal-Poc} survives. Its satellites are little more than crumbling ruins. Something shifted in the magic of the jungle when the Qoniirans chose not to fight in the \textit{Calamity}. The \textit{Rifenmist Jungle} closed its boughs against the world, shielding the city. When the The Strife Emperor died on the \textit{Beynsfal Plateau} to the south, the Qoniirans knew nothing of it. All the tumult of the world passed them by.

The \textit{Rifenmist Jungle} sunk \textit{Niirdal-Poc} into a deep basin, and its thick canopies shield it from above, while its intertwining boughs hold the heart-city tightly in its embrace. Few are permitted by the protective, ancient magic of the Qoniirans to enter the city. Its people live contentedly within \textit{Niirdal-Poc}, only occasionally traveling outside their basin to trade with \textit{Orroyen nomads} or, in mythically rare instances, with the town of \textit{Byroden}.

The four ruins which surround \textit{Niirdal-Poc} are:

\subparagraph{Niirdal-Gan}
This outer city to the north of \textit{Niirdal-Poc} was the agricultural hub of Qoniira. Some farmers have opted to live outside the basin and tend to these still-fertile lands. They are few, for magic is able to supply food for most people of \textit{Niirdal-Poc}.

\subparagraph{Niirdal-Teek}
This outer city, east of \textit{Niirdal-Poc}, was an academic center of the Qoniiran civilization. Most of the city is one vast university that once was populated by over 30,000 people. Some explorers venture into these ruins to find lost knowledge.

\subparagraph{Niirdal-Sarqet}
Located south of \textit{Niirdal-Poc}, this outer city was a center of faith for the Qoniirans. All that remains of the sacred city is a vast, floating cube called the Hexahedron of Sarqet. It is adorned with statues of the \textit{Prime Deities} and the \textit{Betrayer Gods} alike, and none living know what is inside its unbreachable walls. Even today, modern Qoniiran priests teach that balance is the true way of the world, telling their people to invoke the names of the \textit{Prime Deities} for blessings, and of the \textit{Betrayer Gods} for curses.

\subparagraph{Niirdal-Hup}
Precious little remains of this fortress-city, the heart of the Qoniiran military. A common parable reminds those who hear it to not fight fate—for of all the ruined cities, only the city of warriors was reduced to dust.

All of these ruins are filled with statues from the height of the Qoniiran Tetrarchy. Some statues have been known to move under the light of the stars, illuminating hidden runes upon their bodies as they patrol like sentinels of antiquity.

\subsubsection{Government}
The Tetrarchy of Qoniira, in ages past, was ruled by four tetrarchs, each of whom ruled one of the civilization's four outer cities. In times where the fate of Qoniira was at a crossroads, they gathered within the heart-city to guide the course of their civilization together.

The relative autonomy of the four outer cities and the looseness of the coalition that bound them seems at odds with the perfect precision with which the cities were built, so precisely matching the stars in the sky. The Qoniirans have always been an easygoing people who typically follow the flow of fate, rather than fight against it.

Today, the tetrarchy survives, even though the four outer cities are in ruins. They gather daily in the Hakredic Dome, an architecturally magnificent granite rotunda adorned with an outer layer of grayish-purple porphyry stone and lifelike, reddish porphyry statues. There, the tetrarchs and their advisors work to keep life simple and content for their people.

The tetrarchy is a hereditary office, but citywide elections are held between eligible heirs to determine which heir should take up their parent's mantle. Elections are also held when a tetrarch's line ends without heirs—often giving advisors a chance to gain power by virtue of their service to the city.

\subsubsection{Magic Beyond Reckoning}
Magic suffuses the city of \textit{Niirdal-Poc}. Its magic is something older than the gods, or the \textit{Calamity}, or anything wielded by mortals today. However, the mortal mages of \textit{Niirdal-Poc} wield the same arcane force gifted by the gods, and its druids and clerics touch upon the power of the gods behind the Divine Gate, just as Exandria's other mortal peoples do.

Both ancient and modern magic mingle in \textit{Niirdal-Poc}, for the patterns of the ancient powers are mysterious, and its methods are unknowable. Some call it the Light of Fate, others the Blessing of the Stars. The most common name is simply "the Gift." It was the Gift that protected \textit{Niirdal-Poc} in the \textit{Calamity}. It is the Gift that causes miracles to occur, like a person surviving a fall from the heights of Vedrim's Pinnacle to the streets of the city. It is the Gift that causes the Hexahedron of Sarqet to hover mysteriously; it is the Gift that surrounds the Statue of The Observer, a four-winged feline humanoid figure in the city's center, with a broken halo of ashen-gold light—and that causes that same light to flicker within the lamps of the city.

\subsection{Orroyen Tribes}

\begin{description}
\item[Nomadic Tribes:]
Population 11,850 (64\% elves, 12\% half-elves, 10\% humans, 14\% other races)

\end{description}

\begin{DndReadAloud}{}
\DndQuote%
{}
{\enquote{Storms bring death, but rain brings life. Storms are power, the power of the distant gods, on display for all to see. Climb into the Stormpoints, young one, and shout your fury into the relentless wind and booming thunder. Your prayers may not be answered by the gods, but you will feel them echo in your heart.}}
{Cloudwriter Ix}
\end{DndReadAloud}

\noindent Born from a wandering colony fleeing the \textit{Calamity}, these wood elves have not just survived within the jungles of \textit{Rifenmist}—they have flourished. Learning the dangers of the surrounding lands, refining their hunting techniques, and maintaining a healthy arms-length relationship with outsiders, the Orroyen call the jungle their domain and protect it furiously.

The Orroyen tribes have no central government, and their membership shifts as different clans move through the region and encounter other tribes, sometimes trading members as they go. Elders are given the most oversight and respect among a traveling tribe, the eldest given the title of dura, and the duty of leading their tribe. Each nomadic group maintains a population of anywhere between a hundred and a thousand people. Temporary lodging is constructed at each resting point, called a \textit{tomenda}.

Most Orroyen groups remain at a tomenda they've constructed for about two years, before depletion of resources or the arrival of dangerous creatures or soldiers forces them to move on. Countless abandoned tomenda can be found throughout the jungle, and Orroyen clans will frequently make camp in old tomenda made by a different tribe—it is considered a gift from the jungle.

To come into prominence within the tribe, or request admittance into the tribes as an ally or member, a hopeful entrant must complete a series of ceremonial trials. These trials change from dura to dura, but are designed to test strength of body and mind. It is not uncommon for those who attempt these trials to return maimed, or to not return at all, and many elders bear the marks and scars of their trials long past. Few outlanders from the north, known as orfindes to the Orroyen, have been allowed into any of the dozens of Orroyen tribes, and the elders have little intention of changing this, for they have heard distant tales of the ancient tyranny of \textit{Drassig} and the supernatural horrors the people of northern \textit{Gwessar} invite upon themselves.

Nevertheless, the Orroyen are practical, not prideful. They trade often with the Free Folk of Mornset, and sometimes even travel as far north as \textit{Byroden} to trade rare goods. There, some Orroyen meet people from \textit{Syngorn} or beyond and share stories of their lands. Love can bloom in these meetings of separate worlds, and sometime results in an orfinde being brought into an Orroyen clan. And sometimes curious Orroyen leave simply to see the world beyond the jungle.

\subsection{Rifenmist Jungle}
The jungles that cover Tal'Dorei's southern expanses are a vast wilderness. Many elves call this lush landscape their home, particularly those of the nomadic \textit{Orroyen Tribes}. There is no doubt that this jungle is as magical and mysterious as Tal'Dorei's temperate forests, for the jungle's greenery is home to countless mystical spirits, ferocious monsters, and arcane mysteries. The vine-strangled floors of the jungle are clustered with swollen trunks, spine-laden ferns, and carnivorous plant life that waits for the unknowing wanderer to make an incorrect step.

The sweltering heat of the sun-baked peninsula grows even more oppressive beneath the jungle's humid canopy. Clusters of venomous insects buzz through the mist-shrouded trees, alongside large lizards, and wild beasts in search of carrion. Sunken valleys lead to pockets of marshland or quicksand pit traps, while others are wrought with a dire, spreading fungus. Those who have lived all their lives in this environment know its dangers well and have hundreds of years of precautions to take against them. On the other hand, those who venture into the jungle from northern Tal'Dorei rarely return, and those that do warn others to never make the same mistake.

The farther south one travels into \textit{Rifenmist}, the denser the jungle grows. Though some elves and humans of northern Tal'Dorei have attempted to colonize the jungles for their vast resources in decades past, they were quickly rebuffed by the \textit{Orroyen elves} of the jungle. Their message is clear: no cities are to be built south of the jungle's edge—and if the northern jungles are cleared or burned, the trees' vengeance will be swift and just. The \textit{Tal'Dorei Council} reacted quickly as well—the elves of the jungle are to be heeded. The few settlers who have joined the Free Folk of the \textit{Mornset Countryside} are deeply superstitious about the southern jungles, and rarely travel deeper than they need to in order to parley with their \textit{Orroyen} allies.

Even the \textit{Iron Authority}, who has been at war with the trees of the jungle itself for decades, has not managed to destroy their foe. The \textit{Orroyen nomads} fear the conquerors of \textit{Beynsfal Plateau} and the massive armies supplied to them by their authoritarian leaders. Their raids ebb and flow as imperial leaders rise and fall. Each new era begins with a bold and cruel warlord rallying their people against their age-old enemy. War is waged, and countless innocents and Authority soldiers perish in a senseless fight. And then resources run low, morale dips, and the legions return to the \textit{Beynsfal Plateau} to lick their wounds.

Legend holds that the jungle will never die, so long as even one of the \textit{Orroyen} lives within it. So far, history has only proved the jungle's enduring strength.

\subsubsection{Rifenmist Jungle Adventures}
Game Masters who set their adventure in the \textit{Rifenmist Jungle} can use these plot hooks for inspiration.

\subparagraph{The Insatiable Sanctuary (mid-level)}
Long have rumors stirred regarding a vine-obscured temple from an epoch long past, but reports have mapped it at many different locations, with returning expeditions finding it missing. This temple has seemed to reveal itself once more in the northern region of the \textit{Rifenmist}, its shadowed passages and secret bounties ripe for the plucking. It isn't until a venturing band delves deeper into the structure that signs begin to reveal the shrine itself to be a living entity, and the treasure a lure to draw in prey.

\subparagraph{Seed of Immolation (high-level)}
Deep beneath the twisted floors of the jungle lies a long-slumbering relic from before the \textit{Founding}: an unborn primordial titan seed, its energy responsible for the fertile landscape and extreme overgrowth of the jungle. The \textit{Iron Authority} and their power-hungry clerics of the The Strife Emperor have tampered with the seed and unleashed powers beyond their control. The titan has begun to develop and wake from its eons-long slumber. If the process isn't halted, it will fully awaken and terrorize the continent. However, if the seed is destroyed, vast swathes of jungle may wither and die without its magical essence.

\subsection{Stormpoint Mountains}
The southern offshoot of the \textit{Stormcrest Mountains}, the Stormpoint Mountains push deep into \textit{Rifenmist Jungle} itself. These steep peaks are shrouded in low-sitting clouds and fog for most of the year, their rocky surface hidden from view by lush overgrowth that's fed by the near-perpetual rain.

This region is beautiful, and in spring, thousands of waterfalls spray mist and cast gorgeous rainbows across the crags. However, due to its distance from most of Tal'Dorei's population centers, few but adventurers and treasure hunters come here—seeking wealth, not beauty. Some Free Folk from the \textit{Mornset Countryside} and pilgrims of the The Wildmother do take great pains to travel here and revel in the Stormpoint Mountains' natural splendor.

\subsubsection{Stormpoint Mountains Adventures}
Game Masters who set their adventure in the Stormpoint Mountains can use this plot hook for inspiration.

\subparagraph{Allies of the Storm (high-level)}
The people of the \textit{Mornset Countryside} and the free people of the \textit{Rifenmist Jungle} long for an end to the tyranny of the \textit{Iron Authority}. A group of freedom fighters who have heard of the characters' legendary feats of heroism contact them with a request: help. They lack the funds to request the characters fight the Authority themselves, but they wish for the characters to recruit a clan of cyclopes that live in the Stormpoint Mountains, including the mighty cyclops stormcallers that lead their people. Once there, the characters learn that their clan is beset by demons who have possessed their stormcallers and are using their mighty powers to conjure a deadly storm to keep the rest of the cyclopes in submission. Only exorcising the demons and saving the stormcallers will earn the characters their aid.

\subsection{Tz'Arrm, Helm of the Emperor}

\begin{description}
\item[City:]
Population 47,400 (70\% hobgoblins, 10\% goblins, 20\% other races)

\end{description}

\begin{DndSidebar}[float=!b]{Content Warning: Slavery}
The \textit{Iron Authority} is a despicable society that embodies all the evils of ruthless authoritarianism. This includes the subjugation of innocents, racial supremacy, poverty, and slavery. People at your gaming table may be more sensitive to this type of realistic evil than villains who simply wish to destroy the world. It's worth discussing with your players in a \textit{session zero} if everyone will still have fun (and be comfortable) if the game contains this type of evil before introducing the \textit{Iron Authority} as a villain in your campaign.



If you and your players want to use the \textit{Iron Authority} in your game, but don't feel comfortable with the depth of its evil, you can strip out these elements. What remains is a classic "evil empire" that wants to conquer Tal'Dorei. Its conquest can serve as a backdrop to the other adventures that your players want to go on.



\end{DndSidebar}

\noindent When the The Strife Emperor fought the The Wildmother, he took the form of a giant of unimaginable size, clad from head to toe in armor of pitted iron. When he was defeated and cast back into his planar prison, his immense, divine armor remained and crashed to the ground. Centuries have passed, and the thirty-foot tall Helm of Strife is now the palace of the Iron Emperor. The rest of their imprisoned god's armor fuels the unstoppable war machine of the \textit{Iron Authority}, melted down to create walls, weapons, and worse.

\subsubsection{Government}
The armies of the \textit{Iron Authority} are strictly regimented, and its society is no less ruthlessly structured. Every \textbf{goblin}, \textbf{hobgoblin}, and \textbf{bugbear} knows their place within the imperial hierarchy. And every being of other races, including a number of elves who have betrayed the \textit{Orroyen} in hopes of gaining power within the Authority, know their place beneath the goblinkin. Iron Emperor Tz'Jarr is the supreme military commander of the empire, loyally served by the four royal generals that rule the empire's individual city-states. Tz'Jarr is unimaginably old, and it is said that his clerics have used blood magic granted to them by their god to sustain him beyond his years.

The military government of the \textit{Iron Authority} is headquartered in Tz'Arrm. Its cultish generals are doggedly focused on maintaining the Authority's war machine—and their own personal power. Beneath these generals are innumerable propagandists, military officials, slave-owning nobles, and their cronies, all of whom ravenously seek a way to climb the ladder of power. No elections are held; positions are only vacated in the event of death or promotion, and all leaders are appointed by the emperor or an immediate superior.

\subsubsection{Society}
As the capital of an empire with an insatiable appetite for expansion and conquest, Tz'Arrm is awash with propaganda and crawling with secret police. In addition to being the headquarters of the military, Tz'Arrm is also the center of operations for the Gauntlet, the empire's intelligence and espionage division. As the children of the The Strife Emperor, loyal goblinkin are seen as the supreme citizens of the empire. The upper echelons of the \textit{Iron Authority} are filled almost exclusively by hobgoblins, for the Iron Emperor has declared them the strongest and most worthy children of their god. By the same token, goblinkin who betray the empire are seen as the ultimate traitors, and suffer a fate worse than death.

Those with wealth and status within the empire are constantly paranoid about falling out of favor. Even great generals see all those around them as hungry wolves waiting to gobble them up. At the bottom rungs of society are the serving class, who eke out a meager existence as indentured servants for the elite, or as merchants who are bound to whatever master owns their shop. Beneath even them are slaves, most of whom are prisoners of war.

\subsubsection{Slavery}
The armies of the \textit{Iron Authority} have captured tens of thousands of slaves during their wars of conquest across \textit{Rifenmist}. The slaves of Tz'Arrm are mostly used as laborers, sent to chisel iron scraps from the The Strife Emperor fallen armor, or to venture back into the lower jungles to forage for food. The elite of the \textit{Iron Authority} sometimes play at benevolence by granting comfortable jobs indoors to slaves skilled in the arts—such as personal portrait-painters, musicians, or sculptors. Such kindnesses are rarely a sign of true mercy; they are more likely a simple (and meager) expression of guilt.

Groups of slaves that pass in and out of the city to labor in the armor yards or furnaces are escorted by a military detachment whose primary purpose is to discourage escape. The powerful people of the city prize their elven slaves above all others, passing them down like heirlooms from generation to generation. elves are treasured possessions by the Authority's despicable elite, not just because of their beauty and long life, but because of how difficult it is to procure slaves from the sheltering boughs of the jungle.

\subsubsection{Geography and Climate}
Tz'Arrm is hellishly hot. The black stone of the \textit{Beynsfal Plateau} and the iron walls of the city trap the heat of the relentless \textit{Rifenmist} sun. The imperial elite enjoy spells that cool their homes and shelter them from the sun, but those who toil under the sun pay a heavy price. The Authority's \textbf{hobgoblin} elite are well suited to such a climate, but their \textbf{goblin} and \textbf{bugbear} slaves suffer as they toil—as do the \textit{Orroyen elves}, whose life in the shady, humid boughs of the \textit{Rifenmist Jungle} has left them unprepared for the blistering heat of the plateau.

\section{Verdant Expanse}
East of the \textit{Stormcrest Mountains}, hundreds of miles of massive, unbridled forest shroud the landscape in mystery and shadow. This dense greenwood was for centuries traversed by few but the elves of Tal'Dorei, for it was known to be their domain, and theirs alone. Even in the time of \textit{Zan Tal'Dorei}, when her rebels and their elven allies saved the realm from tyranny, it was understood that the coasts were humanity's domain, and the woods were that of elvenkind. Times have changed, however. The alliances formed to oppose \textit{Thordak the Cinder King} united virtually all of Tal'Dorei's disparate peoples, and many more than simply elves now walk the shaded paths of the \textit{Verdant Expanse}. Even so, much of the Expanse is protected by the watchful eyes of the \textit{Wardens of Syngorn}, as it has been for hundreds of years.


\begin{description}
\item[Majority Faiths:]
The Arch Heart

\item[Minority Faiths:]
The Matron of Ravens

\item[Imports:]
Precious gems, grain, gold

\item[Exports:]
Lumber, elvish goods, weapons, jewelry, meat and vegetables, enchantments

\end{description}

Merchants travel from \textit{Emon} to \textit{Syngorn} along roads sanctioned for trade, hunters both human and elven alike seek glory in tracking the wild beasts that roam the untamed wilderness, and bands of dark elf raiders from \textit{Ruhn-Shak} prowl the forest under cover of night. This enchanted wood holds the decay of the seasons at bay, and its trees are green year-round.

The \textit{Verdant Expanse} is sustained by a massive confluence of ley lines—rivers of magical power that flow through the earth—after the devastation of the \textit{Calamity} shifted the flow of magic across the world. The region is saturated with magic. Magical creatures flock to the forest's supernaturally vibrant boughs, including wayward fey, unicorns, unusual arcane monstrosities like owlbears, and displaced aberrations that call the darker groves their home.

\subsection{The Gladepools}
Surrounding the southern edges of the Expanse lies a cluster of broken ponds and lakes that ever draw fresh water from the \textit{Stormcrest Mountains} to the east, along with saltwater from the Ozmit Sea to the west. The strange mix of habitats across this somewhat marshy grassland has led to unusual ecosystems and odd, dangerous denizens, which in turn has drawn the attention of many fishermen from the \textit{Rifenmist} region, and of \textit{Syngornian} hunters from within the \textit{Verdant Expanse}.

The clay and silt gathered along the shores of the saltwater lakes can be refined into fine ceramics and simple constructs that fetch fair prices in northern Tal'Dorei. \textit{Syngornian} tradition tells of an oracle's spirit that is bound to the lakes, and when given sufficient offering and respect, the oracle will emerge to grant sufficiently reverent supplicants a cryptic prophecy.

\subsubsection{Gladepools Adventures}
Game Masters who set their adventure in the \textit{Gladepools} can use this plot hook for inspiration.

\subparagraph{No Basis for a System of Government (any level)}
While traveling along the edge of the \textit{Gladepools}, an elf character, or a character with elven ancestry, notices a gleaming \textit{longsword} embedded in the bank of the lake. No one but that character can see it until they draw it from the mud. Flowing elven script upon the blade reads, "Whosoever draws this blade bears the mantle of the true scion of Yenlara."

Attempts to appraise the sword and discover if it's a true ancient \textit{Syngornian} artifact plunge the characters into a dark conspiracy that lurks at the heart of \textit{Syngorn}: The High Warden may be wise and just, but she is not the rightful ruler—at least, she is not Yenlara's heir. Centuries ago, someone strove to subvert Yenlara's bloodline. Why? And, if the High Warden is truly a good ruler, does the bloodline deserve to be restored?

\subsection{The Mirescar}
Deep within the western reaches of the \textit{Verdant Expanse} is a region where the trees grow so close together that sunlight is barely permitted to pass through their boughs. These ancient trees cluster and twist into a labyrinth of knotted roots and branches, while a tangled canopy of drooping foliage and gray moss blocks the sky from view. Strange powers dwell within the Mirescar, and beasts and travelers that wander into its depths often emerge with their hearts and flesh touched by a deep and corrupting evil—if they emerge at all. The woodlands surrounding the Mirescar are rife with dire beasts, cruel fey, and many-eyed aberrations.

The Mirescar is the subject of countless \textit{Syngornian} folk tales, many of which purport to tell the true story of how it was formed. Some say it contains the bones of a titan. Others say that a meteor from Ruidus itself landed in the forest. Another, that the The Spider Queen cursed it during the \textit{Calamity}. Any one of the tales may be true—or none of them. Only adventurers of great courage and skill have the power to penetrate deep enough into the Mirescar to learn the truth and retain their humanity.

More often than not, these tales serve to frighten the people of \textit{Syngorn} and \textit{Byroden} away from the region altogether.

\subsubsection{Mirescar Adventures}
Game Masters who set their adventure in the Mirescar can use these plot hooks for inspiration.

\subparagraph{A Sickness Spreads (mid-level)}
A plague has begun to spread from the Mirescar. People living in the woods around this evil realm are suffering as they fall blind, with milky-white eyes; or their bodies are wracked by fevers; or they are struck by a number of other mysterious symptoms. All are aware that the vile Mirescar is the source of the \textit{contagion}. The characters are hired to delve into its boughs and find a cure, or to at least stem the flow of disease.

In ancient times, the Mirescar was called the Heartwood. It was the domain of a wise and kindly \textbf{guardian naga} named \textbf{Sagacitous Erusaire}, who welcomed all with good hearts to share their bounty or receive alms. However, this naga has been corrupted into a cruel shadow of his former self. He and his \textbf{diseased grick} minions have been twisted into hideous, pus-oozing spreaders of disease—and spreading it is now the only thought in their minds. The naga can cast \textit{contagion} at will, and his \textbf{diseased grick} minions can do so once per day. Is it possible to save the mind of this corrupted guardian?

\subparagraph{Tree of Scarred Knowledge (high-level)}
A \textbf{wraithroot tree} dwells at the heart of the Mirescar. It is covered with deep gouges, all of which open and close like eyes and mouths as it watches and, on occasion, speaks in a hissing voice to its surroundings. It stands atop a fell artifact—the very thing that turned the Mirescar to evil in ancient times—and has drunk deep of this artifact's forbidden knowledge and its cruelty.

The tree is excited to bargain with any creature that meets it, trading secrets of unknowable origin for the memories held most dear. It fears no attack, for it is protected by a small legion of invisible stalkers that enact its will throughout the Mirescar.

\subsection{The Shifting Keep}
The Verdant Guard has a secondary base of operations, beyond \textit{Syngorn's} walls, called the \textit{Shifting Keep}. It is built within and around a cluster of massive, living trees that form the heart of the \textit{Verdant Expanse}, and enchanted with the most powerful illusion magic that could be mustered by \textit{Syngorn's} military arcanists and their guild of Spellbenders. Enchantments woven into the keep's foundations allow it to cloak its presence, seemingly vanishing and reappearing elsewhere within the forest at will.

The \textit{Shifting Keep} can't teleport, but it can turn invisible and project an illusion of itself to any point within a mile of its true location. This remarkable enchantment has made it almost impossible for invaders to besiege the keep, and it played a major role in \textit{Syngorn's} victory against Neminar \textit{Drassig} in the \textit{Scattered War}—tales of the unassailable \textit{Shifting Keep} are told in \textit{Syngorn} to this day. Nearly a thousand trained elven warriors, hunters, and spies call the \textit{Shifting Keep} their home, though many spend their time patrolling and exploring the \textit{Verdant Expanse}. These rangers are vigilant, seeking signs of whatever foolish invaders would dare trifle with the might of \textit{Syngorn}.

\subsection{Syngorn}

\begin{description}
\item[City:]
Population 38,540 (85\% elves, 9\% humans, 3\% halflings, 3\% other races)

\end{description}

\noindent Founded by the sorceress Yenlara in the wake of the \textit{Divergence}, \textit{Syngorn} became the heart of elven civilization on Tal'Dorei following its calamitous destruction. The remnants of the fallen Court of Ullusa fled to the Fey Realm for a generation following the end of the \textit{Age of Arcanum}, returning only when Yenlara found the realm safe to return. These four hundred survivors built a new home, taking inspiration from their temporary home among the fey, and formed the foundations of \textit{Syngorn}. Though elves have spread all across Tal'Dorei, even as far as the frigid \textit{Neverfields}, there was once a time where nearly all of them resided in \textit{Syngorn}.

This beautiful city of curving stone is built in harmony with the trees of the \textit{Expanse} and sits against the western base of the \textit{Stormcrest Mountains}. It is all but impervious to external assault—not simply because of its 40-foot-high walls of ivy-covered \textit{jade}, and not only because it's also surrounded by living trees of the \textit{Verdant Expanse} and beacons of detection that keep constant vigil for intruders. No, \textit{Syngorn} is virtually unassailable because every entrance to the city is warded by a series of threshold crests: massive emblems of a crescent moon, flanked by two trees, over a deep cerulean stone. These enchanted stones act as an anchor to the Fey Realm, where \textit{Syngorn's} founders once recovered, allowing the city itself to vanish tracelessly into the Fey Realm as a final act of preservation.

\begin{DndReadAloud}{}
\DndQuote%
{}
{\enquote{All here will say to you that Syngorn is a city without equal on Exandria. Its towers and bowers are impeccable, yes. Its cloisters and arches so graceful and fine they make your heart ache, yes, it has been said. But our dreams and memories are the most precious treasure here. They flicker like fireflies through the night, and they are heard in the laughter of the youngest elven child, and in the reminiscing groans of the eldest of our kind.}}
{Ouestra}
\end{DndReadAloud}

\subsubsection{Government}
The city has been governed since its inception by the \textit{High Warden}, a hereditary monarch who appoints proven individuals to three other offices alongside them as the \textit{Wardens of Syngorn}. Each Warden carries a title and bears a responsibility to guide the city toward safety and prosperity. The Verdant Lord is the head of the city guard, though they delegate this responsibility to a Vice Protector when leading the armies of \textit{Syngorn} to war abroad.

The Guildrunner manages the city's treasury and oversees commerce within \textit{Syngorn's} borders. The \textit{Voice of Memory} is a heralded keeper of history and culture, and is often seen interacting with the elven people, gathering new memories for the archives. The High Warden has always been of Yenlara's bloodline, and is responsible for keeping order within \textit{Syngorn} and its territories, as well as within the \textit{Wardens} themselves.

\subsubsection{Society}
\textit{Syngorn} is steeped in elven tradition that dates back to before the \textit{Divergence}. The arts are lauded and revered, the pursuit of knowledge is respected and encouraged, and some training in refined martial techniques is culturally expected. The idea of trade between other cities and nations is understood to be both healthy and beneficial, but most foreign trade and travel is relegated to outer areas of the city.

There was a time that few foreigners ever saw \textit{Syngorn's} innermost reaches, but that time has passed. \textit{Syngorn's} gates are open to travelers from across Tal'Dorei—though all who visit the city are bound by \textit{Syngorn's} strict cultural laws to respect their millennia-old traditions. Any who disrespect the ways of \textit{Syngorn} are warned but once, and are thrown out of the city thereafter. Delicate crafts by \textit{Syngornian} hands are sought after by collectors around the world, so many \textit{Syngornians} take up the trade for both profit and honor.

In a populace made up of people who are extraordinarily long-lived, having children is a rare and highly-regulated process. Prospective parents must gain approval from the office of the \textit{Voice of Memory} to procreate, and any unapproved children are sent out of the city to be raised in outposts or foreign cities. All elves call such an occasion "regrettable," and truly do treat it with sorrow—but the laws of \textit{Syngorn} are clear on this point, and few elves are willing to speak out against rulings which have stood for a thousand years.

\subsubsection{Crime}
If \textit{Syngorn} appears free of crime, it's because elven criminals have hundreds of years of experience at lurking in the shadows. A thriving market for illegal goods runs invisibly through the city's guilds. Though time often seems to stand still in \textit{Syngorn}, crime moves with the same celerity as it does in other cities—it must, if the criminal underbelly wants to stay one step ahead of the \textit{Wardens of Syngorn}.

The most common contraband to pass through \textit{Syngorn} are shipments of \textit{suude}, stolen goods, and dwarven trinkets. Many come to \textit{Syngorn} to purchase these goods for a low price on the black market, but those unwilling to make a journey and risk being stopped by the Verdant Guard can wait for these goods to make their way up to \textit{Kymal}, where they can be purchased at an outrageous markup of five times their cost in \textit{Syngorn}.

Most people in \textit{Syngorn} have their basic needs taken care of by the city's government, to ensure the wellbeing of their people. Nevertheless, some people still fall into poverty, and these destitute individuals are often preyed upon by pickpockets, swindlers, and smugglers who seek to use them as patsies for a paltry sum of gold.

\subsubsection{Geography and Climate}
The perpetual shade of the thick, green canopy of the forest maintains cooler weather throughout the summer months, while this far south, the ice and snow of winter is present, but minimal.

\textit{Syngorn} is broken up into six main districts and has a variety of neighborhoods:

\subparagraph{Memory Ward}
The mind of \textit{Syngorn} is contained within a ring of marble walls in \textit{Syngorn's} northeast reaches. A winding stair, the One Thousand Steps, leads up to its lofty gate. Standing in the center of its grand court is the mystical Sequoia of Remembrance, a three-hundred-foot-tall redwood within which are stored the lives and memories of every elf in Tal'Dorei. At least, those whose bodies are brought before the sequoia and are granted the Rite of Remembrance, as performed by the \textit{Voice of Memory} and her disciples, have their memories preserved and added to the collective. Nearly all elves in \textit{Syngorn} receive this rite, and a vast number of elves who live elsewhere write into their wills that they must be subject to the rite after their death.

Wisps of violet light flit about the sequoia's ancient boughs, and those who look closely into their centers have been known to catch glimpses of eras past, troubles present, and things that yet may be. The great redwood and its spirits are watched and tended to by the Dreamweavers, elves who have dedicated their waking and sleeping lives to the protection of their heritage. About five hundred Dreamweavers and five thousand scholars, merchants, and other elves call this district their home.

The elves are a culture obsessed with tradition and the past. As such, the Dreamweavers are among the most well-respected factions within \textit{Syngorn}. Their leader, \textit{Ouestra, the Voice of Memory}, is likewise the most revered and socially powerful of the \textit{Wardens}.

\subparagraph{Beryl Keep}
This fortress-district houses the martial might of \textit{Syngorn}. Situated on a hill in the northwest of the city, thick walls of leaf-green beryl separate the district from the rest of the city. Within these walls are rows of barracks, archery ranges, trance chambers, forges, mess halls, and all the other necessities to train an army of long-lived elves. At the far northwest of the district is the Beryl Keep itself. The near-unassailable fortress proves the saying that even the most utilitarian of elven things are beautiful to mortal eyes. Six eucalyptus trees of unknown age mark the boundaries of the keep, and between them grows a thick curtain of ironbark. A dense canopy of impenetrable steelfern forms its roof. Legend says that when \textit{Syngorn's} need is greatest, the six ancient eucalyptus trees that guard the keep will uproot and march as mighty treants to defend the city.

Beryl Keep is commanded by \textit{Verdant Lord Celindar}, \textit{Syngorn's} master strategist. He possesses a \textit{crystal ball} of \textit{true seeing} within his war room, which he can use to \textit{scrying} upon any part of the \textit{Verdant Expanse} and to pierce any illusion. From this vantage point, he commands his forces like a chess master, always one step ahead of the threats within the enchanted forest. So far, \textit{Syngorn} remains safe because its enemies are wild and disorganized. Were they to unify, the might of the Verdant Guard would truly be put to the test.

The armies of \textit{Syngorn} are about five thousand elves strong, though only about three thousand are within this district at any time. The rest are stationed at outposts within the \textit{Verdant Expanse} and the Fey Realm.

\subparagraph{Tarn Ward}
\textit{Syngorn's} central ward is a peaceful commerce district surrounding Lake Ywynnlas and split by the channels that feed it. From sunrise to sunset, no business is conducted here, save for food service at public houses and beds at inns. During the daylight hours, elves socialize, play, sing songs, write poetry, paint great works of art, and meditate here. elves are long-lived, and if half the day is not spent in peace and self-improvement, it is thought that the day was truly wasted, no matter how productive one might have been in matters of business.

At nightfall, however, the Tarn Ward changes. The Tarn Thoroughfare opens at moonrise, lit only by the heavens and by magical, floating lanterns. The streets are silent but for the ethereal elvensong that guides those who know how to listen through the streets. elves and half-elves raised in \textit{Syngorn} know this musical language, but those unfamiliar with it must use a \textit{comprehend languages} spell to follow the melody. When the river of sound guides a customer to a vendor, they speak in hushed Elvish, like an audience whispering at the theatre.

Most shops along the Tarn Thoroughfare are affiliated with one of three guilds—the Spellbenders' Guild, the Elvencraft Alliance, or the Mithral Fellowship—and the sly \textit{Guildrunner Rawndel} sits at the head of all three. The Warden of this district, Rawndel lives in a permanent \textit{Mordenkainen's Magnificent Mansion} with a façade adorned with splendid fey grotesques, and is respected by elven merchants and nobles alike. The elves don't seem to mind his monopoly on Thoroughfare trade, as their lives haven't been affected much, but \textit{High Warden Tirelda} worries of the long-term consequences of Rawndel's power-hungry behavior.

\begin{DndSidebar}[float=!b]{Syngornian Goods}
Any basic goods can be purchased in \textit{Syngorn} at twice their usual cost. These goods, however, are of fine elven make. +1 weapon, horses and other mounts increase their movement speed by 10 feet, and other items are possessed of unearthly beauty and preternatural durability. All common magic items can be easily found here, and any uncommon magic item has a The item is nowhere to be found. chance of being found here in the ward's many market stalls. Finding magic items of greater rarity requires a character to make a DC 21 Charisma (Investigation) check to find a rare item, or a DC 25 (Investigation) check to find a very rare item. These items are rarely for sale, and can only be bartered for—typically by completing a dangerous quest for the person who possesses them.



\end{DndSidebar}

\subparagraph{Emerald Citadel}
The great palace of \textit{Syngorn} looms high over the north of the city—a mighty castle of faded white marble, now covered with climbing ivy and sprouting plants. As the highest point in the gradually sloping city, its emerald-tipped spires can be seen from anywhere below. Visitors to the citadel first climb a grand set of stairs, curved like a flowing river of marble, before reaching its brass gates. The Verdant Guard aggressively protects the castle's main entrance, rejecting any commoner who does not have an invitation to the palace marked with the High Warden's seal. \textit{High Warden Tirelda's} monocled majordomo, a silver-haired elf named Ibbimas, screens all of the High Warden's supplicants.

The palace's high, vaulted ceilings are supported by pillars of gleaming marble carved in the image of tall elm trees, such that the ceilings are a canopy of polished stone. Beautiful portraits and busts of past \textit{Wardens} line the walls, along with opulent tapestries of Yenlara and the creation of \textit{Syngorn}. The citadel is made up of four main levels. The basement levels hold the castle dungeons, as well as grand vaults containing fabulous treasure, historical artifacts, and overwrought gifts—the latter given by emissaries uncertain of how to present a gift fine enough for an elven monarch. The first floor holds lush quarters for ambassadors and visiting dignitaries. The second floor's chambers are dedicated to business of state, including the High Warden's throne room. The towers above the citadel contain studies and quarters for the High Warden and her family.

\subparagraph{Feygrove}
A once-splendid manor house in southern \textit{Syngorn} is now completely overgrown with plant life. Massive mulberry trees sprout at odd angles out of windows and through gables, and fairy lights dance around the house at all hours of the day. When \textit{Syngorn} returned from the Fey Realm after the \textit{Conclave's} defeat, the elves unintentionally brought a little of the Fey Realm with them. A number of fey creatures slipped through a gate from the Fey Realm to Exandria, then wriggled through \textit{Syngorn's} defenses and made their home in the mansion of an elven noblewoman named Lady Ladri Il'shavfa. She has tried for years to reclaim her home, but even the patience of elves wears thin; the fey are simply too numerous and too tenacious to be kicked out for good.

The fey interlopers, however, are having the time of their lives. Far from merely throwing nightly parties—though they do that, too—they have made Il'shavfa Manor a place for them to experiment with wild new magic. Their leader, a \textbf{pixie} prince named Windybranch, has found great pleasure in planting strange fey plants in the house and watching as they integrate with the mansion. The house is now a living, breathing, thinking thing—a new friend for the fey! It keeps the elves out and it's a great conversationalist! To the fey, this arrangement has no downsides! The fairies have named this manor after the Archfey who first allowed them into the city. To them, their new home is Artagan's Lodge (and it's said that the The Traveler occasionally stays within its walls when he visits Tal'Dorei).

Since the Feygrove was first "founded" over twenty years ago, the fey have spilled out of Il'shavfa Manor and made an entire district of the city into their home. By this point, only the most stubborn of elven nobility still gripe about "the fey menace" within their midst. Most of \textit{Syngorn's} residents are delighted to have a district filled with the fey, even if it causes the southern end of the city to be a little more chaotic than usual. The fey here seem to have kind hearts, and their chaos rarely crosses the line from churlish pranks into cruelty.

\subparagraph{Reverie Walks}
The artistic and spiritual center of \textit{Syngorn} is a winding labyrinth of living trees and stones. elves who seek a day-long (or days-long) journey of meditation may wander the Reverie Walks, entering a trance-like state as they do so. Those who seek true enlightenment may dedicate years of their near-immortal lives to wandering the ever-shifting labyrinth, finding peace in isolation or in search of the Stone of the The Arch Heart in the labyrinth's center. An elven monk named Lyssev Sorveline has walked the labyrinth for five hundred years, seeking answers to questions even the gods do not know. Some elves leave gifts of food for the Wanderer within the labyrinth, allowing them to continue their eternal meditation in peace.

No Warden holds sway over the Reverie Walks. Legend holds that the living stones and trees that shift the labyrinth are devotees of the The Arch Heart that swore to be their god's eternal wardens.

\subsubsection{Points of Interest}
\textit{Syngorn} has a number of famous landmarks. They are keyed to the map of \textit{Syngorn}.

\subparagraph{1. Elvencraft Alliance Guildhall}
The Elvencraft Alliance is supposedly the oldest guild in \textit{Syngorn}. They are entrusted with creating works of art that spread the culture of \textit{Syngorn} far and wide across Tal'Dorei. The front rooms of their guildhall are a museum of fine elven crafts from ages past, and they're open to all who wish to appreciate \textit{Syngorn's} cultural heritage.

\subparagraph{2. Guildrunner's Mansion}
\textit{Guildrunner Rawndel}, Warden of the Tarn Ward, makes his home within a permanent \textit{Mordenkainen's Magnificent Mansion}. Its exterior is a magnificent façade of stone and leering fey statues, yet if one were to look inside, they would find it completely empty. The door, which is constantly locked to all but the Guildrunner and his guests, is actually a portal to an extradimensional space containing a manor of supernatural beauty and luxury.

\subparagraph{3. Mithral Fellowship Guildhall}
The master metallurgists and smiths of the Mithral Fellowship are tasked with arming the Verdant Guard and the armies of \textit{Syngorn}. In peacetime, they turn their talents to creating things of beauty from metal and jewels, sometimes in tandem with the artists of the Elvencraft Alliance. Due to their military involvement, none but authorized members of the fellowship are allowed within their guildhall.

\subparagraph{4. One Thousand Steps}
This winding stair is the only way in or out of the \textit{Memory Ward}. Each step is inscribed with an ancient elven idiom, aphorism, or impossible philosophical question. Anyone who enters the \textit{Memory Ward} is told to read the aphorism upon the first step and ponder it the entire time they are climbing the stairway. If they find peace in their answer, they are to read the inscription of the second step on their next journey, and so on. There are few, even among the elves, who have pondered the inscriptions of all of the steps to their satisfaction.

\subparagraph{5. Spires of Yurek}
Roads lined with dormitories and vendors radiate out from the branches of the Sequoia of Remembrance to the walls of the \textit{Memory Ward}, providing essentials for not only the Dreamweavers, but the elven scholars who teach, study, and experiment with magic within the Spires of Yurek. This academy was named for \textit{Yurek Windkeeper}, founder of the \textit{Arcana Pansophical} and personal counselor to the \textit{Voice of Memory} herself. The school is small compared to \textit{Emon's} magnificent \textit{Alabaster Lyceum}, but its seven marble towers are no less awe-inspiring. Characters who visit the Spires of Yurek can learn any historical fact by studying its tomes for one hour and then making a successful DC 20 Intelligence (History) check. They can even uncover hints toward long-lost secrets by completing a day's study and then succeeding on a DC 25 Intelligence (History) check. New spells are constantly being invented here, and countless tomes filled with ancient spells can be found within the libraries.

\subparagraph{6. Tower of Moonlight}
Most visitors to \textit{Syngorn} doubt that \textit{Ouestra, the Voice of Memory}, truly dwells in a tower of moonlight. It is the truth. Her tower does not exist under the light of the sun, and then appears as a semi-solid pillar of Catha's silvery glow, occasionally veined by lines of Ruidus's ruddy light when the smaller moon is full. Its inside has supposedly never been seen by any but the other \textit{Wardens of Syngorn}, and a thousand rumors circulate throughout the realm of the magical and divine marvels that must be contained within.

\subparagraph{7. Spellbenders' Guildhall}
The \textit{Verdant Expanse} is saturated with magic. The ley energies that suffuse the greenwood make it easy for arcanists to create works of spellcraft by themselves, when it might take a half-dozen mages working in concert in other lands. Many of \textit{Syngorn's} most trustworthy magi are members of the Spellbenders, a guild dedicated to creating new spells and using them to create works of art, defend their homeland, and improve their peoples' lives. Their guildhall is only accessible to those with the express permission of \textit{Guildrunner Rawndel}.

\subparagraph{8. Stone of the Arch Heart}
In the heart of the Reverie Walks is a column of pure diamond, carved in the radiant likeness of an androgynous elf. Though the The Arch Heart is worshiped throughout Tal'Dorei as a god of magic and art, they are something more to the elves of \textit{Syngorn}: a progenitor of their people, a deity of fatherhood and motherhood, yet also of neither—a liminal god that by their very nature both embraces and destroys binaries. elves that are "born of the Arch Heart" are honored within \textit{Syngorn}, and many who seek that honor wander the Reverie Ward in search of the Stone, longing for physical transformation, social power, or arcane might. Some search for centuries and never find it, but in times of peril, the The Arch Heart always makes their wisdom known to those who truly need it.

\subsubsection{Syngorn Adventures}
Game Masters who set their adventure in \textit{Syngorn} can use this plot hook for inspiration.

\subparagraph{The River Does Not Run Through It (low-level)}
Word arrives from Verdant Scouts that the Feshun River has stopped emptying into the Ozmit Sea without warning. The powerful flow from the \textit{Tormor Falls} shows no signs of dwindling river flow, so the party is sent to investigate, only to find a sudden sinkhole that has consumed a small section of the \textit{Expanse}, causing the river to pour into this massive, underground cavern. The rock and earth below appears to be crumbling at an accelerated pace, and the air below smells of decay. What could be the source of this rapid destruction, and how can it be remedied?

\subsection{Tormor Falls}
Feeding the mouth of the rushing Feshun River that carves through the body of the \textit{Expanse}, Tormor Falls is an incredible multi-level series of waterfalls that cascade down the eastern side of Orencleft Mountain for hundreds of feet. Swelling with every major rainfall, and beautiful to behold at all times of the year, Tormor Falls is also host to a number of caves that hide beneath the mist and spray, leading beyond the forest and under the \textit{Stormcrest Mountains}.

A squad of Verdant Guard patrols the base of the falls, and many have explored the caverns in the past. Some return with nothing; others discover old relics and trinkets from the earliest days of Tal'Dorei, or even from the time of the \textit{Calamity}. It's rare for a patrol to fail to return, but when they do, no search party has ever been able to make contact with survivors, or to reclaim their remains. The guard has begun to whisper that the cave entrances move, shifting with each rising sun. The superstitions surrounding the caves around Tormor Falls have grown so great that expeditions which fail to return are no longer searched for at all, and no patrols go inside the caves without a critically important reason.

\subsubsection{Tormor Falls Adventures}
Game Masters who set their adventure in the \textit{Tormor Falls} can use this plot hook for inspiration.

\subparagraph{This Cave Was Made for Me! (mid-level)}
Characters that travel the edge of the \textit{Verdant Expanse} hear distressing news that people from small villages, and even from \textit{Syngorn}, have disappeared in the night—and now, a friend of the party has joined the ranks of the missing. What they don't know is that these people are leaving the village themselves, hypnotically drawn to the caves behind \textit{Tormor Falls} by an \textbf{aboleth} living within a pool deep inside the cave system. The beings that emerge from those caves are no longer humanoid, but nightmarish mockeries of life.

\subsection{Vues'dal Waters}

\begin{description}
\item[Scattered Villages:]
Population 2,420 (52\% elves, 28\% goblinkin, 12\% lizardfolk, 5\% humans, 3\% other races)

\end{description}

\noindent The volcano known as Mt. Vues'dal erupted for the last time in the early days of the \textit{Age of Arcanum}. Its explosion shook the earth so terribly that the mountain itself was swallowed, creating the Vues'dal Basin at the edge of the \textit{Stormcrest Mountains}. Today, the basin has been completely filled by the Feshun River, and the land around the Vues'dal Basin is among the most fertile farmland in southern Tal'Dorei. The villages of the Vues'dal Waters are known for their prize produce, exporting tomatoes, cherries, pumpkins, lentils, sweet potatoes, and dozens of other foods, both staples and luxuries, to \textit{Syngorn}—and from there, to the rest of Tal'Dorei.

The people of the Vues'dal Waters are mostly elves, but over the years they've been joined by some of the local \textbf{lizardfolk} that have dwelt in the marshes, as well as goblinkin fleeing the \textit{Iron Authority} to the south. All of these various people have banded together to create something beautiful around the caldera of this long-dead volcano—and to repel bandits and monsters from the nearby marshes. The \textbf{lizardfolk} hate fighting their kin, but human and \textbf{lizardfolk} bandits alike have made the wetlands their home, and prey upon these honest farmers from encampments that shift from week to week.

\subsubsection{Vues'dal Adventures}
Game Masters who set their adventure in the Vues'dal Waters can use this plot hook for inspiration.

\subparagraph{Legions of the Reptile God (mid-level)}
The people of Vues'dal have unwittingly been skirmishing with the vanguard of a force that wishes to bring \textit{Syngorn} to its knees. The \textbf{spirit naga} Maledicta Hexos has come from below the \textit{Stormcrest Mountains} to rally the \textbf{lizardfolk} tribe in the name of the The Cloaked Serpent. Vues'dal will soon be overrun, and the Verdant Guard has called for \textit{Emon's} aid in driving back the legions of this monster with delusions of godhood.

Maledicta's armies contain not just countless \textbf{lizardfolk}, but also several hydras, and a \textbf{night hag} that serves as her lieutenant. Maledicta is also in possession of one of the \textit{Astural Scrolls}, an artifact looted from \textit{Wrettis} that grants her uncanny cosmic powers.

\section{Bladeshimmer Shoreline}
The Bladeshimmer Shoreline, named for the distant glimmering of sun across the Ozmit Sea, stretches across the central western coast of Tal'Dorei. This coast is where humans, halflings, and gnomes first set foot on the continent of Tal'Dorei, and it bears the marks of their first steps on this land. It's also home to the present heart of the Republic of Tal'Dorei.

The shoreline is a hub of international trade, for the western coast of Tal'Dorei is closest to the continents of \textit{Issylra} and \textit{Marquet}. Dozens of tallmasted trading ships set sail to and from its calm shores each day—as well as the majestic skyship that daily soar in and out of the resplendent capital city of \textit{Emon}.

Inland Bladeshimmer is mostly temperate grassland, intercut by cool, winding rivers. The bulk of western Tal'Dorei's produce comes from farms here, blessed with clean water and non-salinated soil, despite the ocean's proximity. \textit{Emon's} presence lends the region stability. Nevertheless, the yeomen living beyond the city's walls struggle to defend their small plots from burrowing ankhegs, corrupting fiends, and the various folkloric spirits and monsters that lurk in the shadows of their minds.


\begin{description}
\item[Majority Faiths:]
The Lawbearer

\item[Minority Faiths:]
The Moonweaver

\item[Imports:]
Precious and industrial metals, lumber

\item[Exports:]
Stone, lumber, ships, fish, grain, produce, cobalt, gold, livestock

\end{description}


\subsection{Daggerbay}
Daggerbay, so named for the jagged Slumber Reef that flanks the bay, was where the first human colonists of Tal'Dorei made landfall. In that bygone time, it was a bustling port for the human city of Port \textit{O'Noa}. It has since fallen into disuse after the \textit{Scattered War} and the destruction of \textit{O'Noa}. Centuries later, the bay is now naught but a haunting reminder of darker days. It is home to hundreds of sunken ships and the bodies of their lost crew, deep below the depths of the reef. Those who know of their history say that the waters are cursed, and the nearby settlers spread rumors of ghost ships and sirens calling looters to their grave should they wander too close to the ominous waves.

\subsubsection{Daggerbay Adventures}
Game Masters who set their adventure in \textit{Daggerbay} can use this plot hook for inspiration.

\subparagraph{Stormbringers (mid-level)}
A society of banished \textbf{merfolk} have claimed \textit{Daggerbay} as their domain. They harass any treasure hunters that seek to loot the wreckage at the bottom of the bay. While the characters are traveling nearby, a terrible storm forces them to take shelter in the ruins of \textit{Daggerbay}.

In truth, the storm was caused by the renegade \textbf{merfolk}. Some weeks ago, they discovered a relic of the The Stormlord with the power to summon and direct storms. The coast is now shrouded in an endless maelstrom. That night, when the characters take shelter in the ruins, they're assailed by the merfolks' land-dwelling humanoid thugs. These brutes attempt to capture the characters and deliver them to their leaders.

\subsection{Daggerbay Mountains}
The pages of history do not remember the Molten Titan, or when this fierce entity of liquid metal was felled by primordial elves at the end of the \textit{Founding}. The entity was sealed within and consumed by the earth, giving birth to a range of mountains that has stood for the ages since. Jutting upward against the northwestern border of the \textit{Verdant Expanse} and stretching to the Ozmit Sea, these stormy peaks functioned as the first boundary between the long-established elves of \textit{Syngorn} and the burgeoning human colonies on Tal'Dorei during their arrival.

For most of history, the \textit{Daggerbay Mountains}, or the lyren'alsi in Elvish, were known to be desolate and devoid of material worth. Prospectors from \textit{O'Noa} came home empty-handed, if they came back at all. The survivors returned with tales of bands of bloodthirsty, one-eyed giants that roamed the peaks: cyclopes. \textit{Issylran} settlers and \textit{Syngornian} elves both learned to leave the lyren'alsi well enough alone.

Decades later, the \textit{Emerald Outpost} was established as a major trading post for the people of \textit{Emon} and \textit{Syngorn}. With such a valuable settlement nearby, the \textit{Daggerbay's} hazardous peaks suddenly garnered the interest of adventurers and fortune-seekers. Small hunts for undiscovered riches grew into vast expeditions. These incursions into the mountains aggravated the mountains' territorial inhabitants; elven scouts reported \textbf{cyclops} raiders and cyclops stormcallers conducting strange rituals within the eyes of thunderstorms at night, and hungry bulettes hunting beyond their normal domain, as if guided by some unknown intelligence.

\subsubsection{Daggerbay Mountains Adventures}
Game Masters who set their adventure in the \textit{Daggerbay Mountains} can use these plot hooks for inspiration.

\subparagraph{Eye of the Storm (mid-level)}
Groups of hill giants are migrating into the mountains, leaving their foothill settlements empty of warriors, civilians, and the infirm alike. A half-elf \textbf{scout} named Thunderchaser has caused a stir in the \textit{Emerald Outpost} by bringing back the lightning-charred corpse of a \textbf{hill giant} with one eye ritually gouged out. What are the cyclops stormcallers planning beneath the concealing stormclouds of the \textit{Daggerbay Mountains}?

\subparagraph{A Metal of Memory (high-level)}
The ancient molten beast that was sundered beneath these rocky peaks left behind pockets, located far below the surface, of an incredibly rare metal with mysterious properties, called orichalcum. The sudden discovery has sent local entrepreneurs into a frenzy, and the ensuing rush has already sparked violence among the bands of treasure hunters. As the adventuring parties dig deeper, they discover that some essence of the ancient creature still lives, and it is seeking vengeance.

\subsection{Emerald Outpost}

\begin{description}
\item[Town:]
Population 1,470 (76\% elves, 19\% humans, 3\% halflings, 2\% other races)

\end{description}

\noindent Originally established by the elves of \textit{Syngorn} as a hidden outpost to protect their homeland from the growing empire of \textit{Drassig}, the \textit{Emerald Outpost} fell out of use after the \textit{Scattered War}. This former military outpost is now an intermediary trading post and a place for merchants traveling between \textit{Emon} and \textit{Syngorn} to rest. The Emerald Chambers, once used as a war room, and more recently used as a diplomatic chamber for the allied governments of \textit{Emon} and \textit{Syngorn}, have these days been converted into a stylish bazaar, in which \textit{Emonian} and \textit{Syngornian} traders can sell their wares without making the full journey to either city.

\subsection{Emon, the City of Fellowship}

\begin{description}
\item[Metropolis:]
Population 365,026 (63\% humans, 8\% dwarves, 7\% elves, 22\% other races)

\end{description}

\noindent \textit{Emon} stands defiantly against all who would threaten Tal'Dorei and its people. It is the cultural and political heart of the Republic of Tal'Dorei, as well as a nexus of commerce, entertainment, travel, education, and adventure.

\textit{Emon} is accessible only through its heavily patrolled gates and by \textit{skyship}. The denizens of the city are generally well protected from outside attackers and sieges.

\subsubsection{Government}
\textit{Emon} is the seat of the \textit{Tal'Dorei Council}, the nation's highest governing body. Though the council originally served under the Sovereign of Tal'Dorei, the nation's final sovereign was killed when the \textit{Chroma Conclave} attacked \textit{Emon}. Following the death of \textit{Thordak the Cinder King} and the rest of the \textit{Conclave}, \textit{Emon} was rebuilt and the council reformed as the backbone of the new Republic of Tal'Dorei. The council and its members are \textit{described in detail starting on page 42.}

\noindent
One unintended consequence of relying on magic to rapidly rebuild \textit{Emon} is that the mages of Tal'Dorei now hold incredible influence over the fledgling council. Some fear that without a sovereign, the council will be unable to keep the arcanists in line, and Tal'Dorei will dissolve into magocracy.

\subsubsection{Current Events}
The \textit{Tal'Dorei Council} is best known for their policies that affect the whole of Tal'Dorei. It is often ignored that the council also governs \textit{Emon} and resolves the day-to-day issues of the city. These are some current events both great and small affecting \textit{Emon} today; these topics can add flavor to the discussions and rumors the characters hear around town, and could even inspire new adventures for them to take on.

\subparagraph{Elemental Disturbances}
Over twenty years have passed since \textit{Thordak the Cinder King} was killed in the caverns beneath \textit{Emon}. His wreckage in the \textit{Cloudtop District}, as well as the \textit{Scar of the Cinder King}, a swath of elemental destruction outside the city walls, have been tended to by members of the Fire \textit{Ashari}, gradually healing the stubborn wounds of his influence. Yet the primordial fires that seeped from his body continue to linger in the caverns beneath the city where he died. Every year or two, an incursion of wild fire and hateful cinderslag elementals burst through the city streets, causing mass chaos, expensive damages, and life-or-death battles for the city guard.

\subparagraph{Mercenary Mages}
A generation ago, the fledgling \textit{League of Miracles} hired hundreds, if not thousands, of mercenary mages to aid in the reconstruction of \textit{Emon} and other cities across Tal'Dorei. These days, Tal'Dorei is rebuilt, and the few infrastructure problems aren't enough to employ all the mercenaries left behind by the league. These out-of-work mercenaries now linger around dismal dive bars and make trouble for the locals. Most people would rather see them gone, but it's hard not to feel sorry for their situation.

\subparagraph{Walls and Slums}
The city's outer walls have fallen under heavy criticism from politicians and civilians alike in recent years. Many people are concerned that the walls' true purpose these days is to separate the slums outside the walls from the "respectable" districts within them. Popular opinion is now that the great \textit{City of Fellowship} should do more to help its less fortunate, but council politicians continue to stall. In private, Master of Commerce \textbf{Vex'ahlia} de Rolo fumes that she wants to help, but can't cut through political red tape with her Fenthras and arrow.

\subsubsection{Society}
While in antiquity \textit{Emon} was the heart of a human empire, the city has grown to encompass people from all over Tal'Dorei, and all over the world. Its citizens include people of nearly all known races from Exandria's many nations, thriving on the new innovations and ideas of its citizenry. It has been this way since the sovereignty of \textit{Zan Tal'Dorei}, when the alliances she formed invited dozens of displaced elves, dwarves, orcs, and \textit{goliaths} into her city. Since then, \textit{Emon} has developed into a city whose people valued innovation and collaboration, especially in small, cohesive groups—virtues that some historians within the \textit{Alabaster Lyceum} believe have given rise to the prominence of the modern adventuring party.

While \textit{Emon} is becoming more commercial and the use of gold as currency has been ubiquitous throughout Tal'Dorei's history, many communities within the city are still close-knit enough to use the barter system. During the reign of \textit{Drassig}, the humans of \textit{Emon} adopted a dwarven oath called rudraz, an intimate promise between two people to repay a deed or trade. Though the rudraz is not a contract, the dwarves believed that an oathbreaker would be forever barred from passing beyond the Brightguard Gates of Hilmaire, the gates that allow dwarves to pass to the afterlife. humans in \textit{Emon} treat the rudraz more lightly, often using it to seal matters of business or politics rather than personal promises, but breaking this oath still carries massive social repercussions. Few look kindly upon a person with the epithet "Oathbreaker".

\subsubsection{Rebirth}
\textit{Emon} has mostly recovered from its destruction at the talons of the \textit{Chroma Conclave} and \textit{Thordak's} subsequent occupation. A number of Tal'Dorei's most significant factions played a role in \textit{Emon's} rebirth, and their members enjoy positions of prominence throughout the city. \textit{Emon} was rapidly rebuilt to its former glory, due to the expensive but effective intervention of mercenary mages serving the \textit{League of Miracles}. Thanks also to the \textit{Clasp's} underground networks and established hierarchies, \textit{Emonian} society and culture was able to bounce back more swiftly than any other major settlement in Tal'Dorei.

And the \textit{Tal'Dorei Council}, formerly an advisory council to the sovereign, is now the ruling council of the land. Folk heroes and career politicians alike sit on the council. Over the past two decades, the council has, with many missteps, enacted the will of the people—allowing Tal'Dorei to transition from a monarchical power to a republic without war or insurrection.

However, the swiftness of \textit{Emon's} rebirth has put the \textit{Tal'Dorei Council} in a precarious position. Not only are they socially indebted to a criminal faction—one whose members now command immense celebrity among the common folk—the council now owes vast amounts of gold to mages of the \textit{League of Miracles}. League conjurers created thousands of tons of stone and steel, and each one of their transmuters rebuilt everything from taverns to townhouses and sewers to stately manors as quickly as a hundred laborers. \textit{Thordak's} hoard of plundered wealth has dwindled to nothing, repaying the league, and they have done the same in cities across the land. Some members of the council have already caved to pressure from factions within the \textit{Clasp} and the \textit{League of Miracles}, and subsequently turn a blind eye to the crime and magical abuse that run rampant throughout the city. Barely a single human generation old, the young republic is already struggling under the weight of ruling in a just and noble manner.

\subsubsection{Defenses}
The Republic of Tal'Dorei has no standing army, and doesn't elect a Master of War—a military commander-in-chief—to its council outside of wartime. In peacetime, the Master of Defense is responsible for the acquisition of mercenaries and adventurers for tasks which require the use of force. They also have the ability to conscript existing militias and town guards into military service and call them to \textit{Emon's} defense at any time, though typically only while the Master of War is being elected.

The only standing forces pledged exclusively to the defense of the realm are the elite forces of \textit{Fort Daxio}. Mobilizing the armies of \textit{Fort Daxio} to \textit{Emon's} aid generally requires the city to withstand a week-long siege. However, in case of internal strife, there are stables positioned every fifteen miles along the Othen Trail, allowing Daxio's legendary outriders to arrive at \textit{Emon} in a mere two days' time by switching out their exhausted steeds.

\textit{Emon} is encircled by 60-foot-high stone walls that stretch from the eastern fields to the western shore—though a number of discussions in the \textit{Tal'Dorei Council} have revolved around the strategic irrelevance of great stone walls. They didn't protect \textit{Emon} against the \textit{Chroma Conclave} over twenty years ago, and they wouldn't have done much against an army of cultists and winged gloomstalkers had The Whispered One led an army against Tal'Dorei. As skyship grow more popular and affordable around the world, even the council's Master of Defense is considering demolishing the city's walls—making the strange case that such fortifications are prized more for their historic and cultural value than for their strategic importance.

These walls are patrolled by lightly armed and armored members of the Arms of Emon, the city's official watch and constabulary. They are divided unofficially into the Arms upon the Wall and the Arms on the Streets, the latter of which are assigned to patrol the city streets to dissuade petty criminals and respond to violent crime.

\subsubsection{Crime}
Hundreds of miles of tunnels of varied provenance run beneath \textit{Emon}. Most are the forgotten remnants of paved-over neighborhoods and secret passageways made by thieves' guilds during the reign of \textit{Drassig}. Others are the lairs of tunneling monsters. The origins of others still are lost forever, for they have already long since been turned into sewer passages to whisk away the refuse of \textit{Emon's} citizens.

The vast majority of the tunnels closer to the surface are known to the \textit{Clasp} and used to rapidly move about the city unseen. Somewhere deep in this network of unlabeled roads is the \textit{Clasp's} secret headquarters, where all manner of thieves, cutthroats, fences, and spies gather beneath their banner. Though the \textit{Clasp} operates in major cities across Tal'Dorei, \textit{Emon} has always been their home, and saving their neighbors from \textit{Thordak's} reign of terror gave the \textit{Clasp} the dubious benefit of a heroic reputation.

However, the \textit{Myriad} crime syndicate from \textit{Wildemount} has spent over two decades relentlessly clawing their way into the \textit{Clasp's} territory across Tal'Dorei—dismantling their networks, poaching their clients, and assassinating their informants. The \textit{Myriad} hasn't managed to completely eliminate the \textit{Clasp's} hold in any of the cities they've targeted, but that hasn't stopped them from bringing the ongoing power struggle to the \textit{Clasp's} doorstep. \textit{Myriad} agents bargain in shadowed alleys, gamble in upscale taverns, and lurk in the very tunnels that the \textit{Clasp} relies on, waiting for their chance to unseat the \textit{Clasp} once and for all.

\subsubsection{Geography and Climate}
\textit{Emon} is not a tourist destination for its climate—the city is known for sudden, pounding rainstorms just as much as for its beautiful, sunny days. The city is blessed with cool summers and warm winters, thanks to its proximity to the cool Ozmit Sea. Though snow rarely falls on the city itself, it relies on springtime snowmelt from the \textit{Cliffkeep Mountains} to fill its reservoirs throughout the year.

\textit{Emon} is made up of the following districts:

\subparagraph{Abdar's Promenade}
The Promenade is an open marketplace district and massive bazaar that dominates eastern \textit{Emon}, named for the legendary spicemonger from \textit{Marquet} who helped fund the construction of \textit{Emon}. To this day, the name of Abdar is synonymous with both generosity and business savvy. Nearly everything a person could want can be found within the tents, carts, warehouses, and shops of the promenade. Its intertwining roads stretch for miles, and it has dozens of tiny neighborhoods, each with their own distinct cultures and identities. The promenade is heavily patrolled by the Arms of Emon, for its sunny reputation is shrouded by an equally significant reputation for literal back-alley criminal deals.

\subparagraph{Central District}
This is the largest residential district within \textit{Emon's} walls. It houses most of the city's merchant class, including traveling traders, ship captains, and guild journeymen. The Central District is a patchwork of thousands of personal homes, tenements, and guildhalls of all shapes and sizes, peppered with small taverns and inns on nearly every street corner.

These neighborhoods lay clustered together among tightly set streets, occasionally broken up by park grasses or the Ozmit Waterways that snake through the region between the promenade and the port. Visitors to the Central District are advised not to go out at night; the wide disparity of wealth from home to home here has seen a recent rise in crime, and the city watch seems reluctant to find a constructive solution. Residents who know the lay of the land have an easier time at night and can help visitors avoid the most dangerous streets.

\subparagraph{Cloudtop District}
Out of all of \textit{Emon's} districts, the Cloudtop has, in some ways, changed the most over the decades. In \textit{Drassig's} time, a monstrous fortress loomed over the manors of his sycophants. In the age of \textit{Zan Tal'Dorei's} descendants, countless mages transfigured that decadent citadel into a sleek, elegant palace. By the time the \textit{Chroma Conclave} came, decadence had once more crept up on \textit{Emon}, and though the Palace of the Sovereign remained much the same, the Cloudtop District was once again a nest of luxury and the throne of the social elite. The latest additions in that era were \textit{skyship} platforms, so that the wealthiest and most connected in Tal'Dorei could travel in exclusive luxury—soaring through the skies rather than enduring arduous journeys over land or sea.

All of that was reduced to slag by the advent of the \textit{Cinder King}. \textit{Thordak} made the Cloudtop District his home, and the palace itself his throne. His very presence caused magma to surge from the caverns beneath \textit{Emon}, where previously no fires burned. \textit{Thordak} ruled there for only a matter of months, but those few months were enough to annihilate unknown millions of gold pieces' worth of art and creations of untold value—though the greedy dragon took care to command his minions to save the most beautiful works for his treasure hoard.

Today, the Cloudtop District is rebuilt, though the most stubborn elemental influences of the \textit{Cinder King} linger in \textit{Thordak's Crater}. A new castle called \textit{Cloudwatch Palace} has been built atop the city of \textit{Emon}, and it contains the chambers of the \textit{Tal'Dorei Council}, as well as grand dining halls and guest quarters for visiting dignitaries. Much of the castle is a museum, displaying what regalia and objets d'art survived the wrath of the \textit{Cinder King}. The private chambers of the palace also contain a number of \textit{teleportation circle} that leading councilmembers can use to return to their homes at night—particularly if they live in far-off cities such as \textit{Whitestone}.

\subparagraph{Temple District}
Some say that this northern district contains more temples than in the rest of Tal'Dorei's cities combined. This is a flagrant exaggeration, but at first glance, most people believe it. \textit{Emon} is a city of many people, from many different places, with many different faiths. Even folk who worship the same gods worship them differently, and have built their own temples. Legend has it that shrines and temples used to be built all across \textit{Emon} (and there are still some small sanctuaries elsewhere in the city), but a cleric of the The Dawnfather declared her temple the greatest in all \textit{Emon}. Her boast was met with scorn from a nearby church of the The Stormlord. Their clerics, always eager to rise to a challenge, built an even greater temple across the street, and soon, people of all faiths were building the most magnificent structures they could, all in one tiny corner of \textit{Emon}.

Today, cathedrals to the The Lawbearer and the The Platinum Dragon stand at apex among smaller temples to the The Dawnfather, the The Stormlord, the The Wildmother, and the The Matron of Ravens. All \textit{Prime Deities} have at least a moderately sized house of worship here, and even some of the \textit{Betrayer Gods} are worshiped in secrecy.

Anyone who seeks shelter here, be they rich or poor, can find sanctuary in the Temple District. Even the faithful of gods of strength, such as the The Stormlord and the The Platinum Dragon, have mercy for those who need the aid of others to help them stand on their own.

Travelers and sailors come to leave tokens and gifts at shrines to bless their journeys, while people of all walks of life seek the graces of the The Changebringer before undertaking risky endeavors. The meek search for wisdom in the great halls of the The Knowing Mentor, as the artists and wistful folk give praise to the The Arch Heart. However, not all folk are drawn toward religion, nor trusting of it, and temples past have been exposed as frauds—or worse, cults who venerate power.

Recently, an off-kilter faith to a lesser idol called the The Traveler has found purchase in the Temple Ward—primarily in the form of puerile graffiti in back alleys. Other quasi-religious spiritual entities have small cult followings in the Temple Ward. They are generally allowed to flourish to whatever degree they can, so long as their beliefs don't advocate for any behavior that would break the laws of \textit{Emon}.

\subparagraph{Erudite Quarter}
This district is home to members of the educated upper-middle class and destitute academics alike. Several institutions of higher learning have their stately campuses here, ranging from well-to-do boarding schools to Tal'Dorei's most premiere academic and arcane institutions.

Beautiful towers in ancient \textit{Syngornian} style stretch to the sky across a cityscape of centuries-old halls of learning and student boardinghouses. Here, the finest schools and colleges draw the wealthy, the gifted, and the brilliant—and none more so than the \textit{Alabaster Lyceum}, often lauded as the largest and most accomplished institute of arcane study on the continent. Though the district's walls were shattered and its libraries crushed during the invasion of the \textit{Chroma Conclave}, the Erudite Quarter has thankfully been mostly reconstructed, thanks to the efforts of the \textit{Lyceum's} staff and students.

Members of the \textit{Lyceum}, under the direction of Headmaster Thurmond Adlam, performed the reconstruction themselves not simply because they sought to project an image of good work ethic and self-sufficiency, but because the heads of their colleges distrust the \textit{League of Miracles}. Their motives and sources of funding are, in the eyes of the \textit{Lyceum}, too opaque to be trusted. Their suspicion has often been derided either as elitism over the diverse array of mercenary mages the league hires, or as sour grapes over losing out on countless lucrative reconstruction contracts.

\subparagraph{Military District}
Since \textit{Emon} has no standing army, the Military District is one of the city's smallest wards. Nevertheless, it has the capacity to house and train a vast number of soldiers if an army needs to be levied from the people of Tal'Dorei. The Military District is mostly filled with several prisons and jails, including the infamous \textit{Black Bastille}, as well as tightly locked barracks and instructional facilities for the Arms of Emon.

\subparagraph{Cemetery District}
This smaller section of the city is filled with graveyards and small religious sanctuaries so that people with means to purchase a gravesite or a mausoleum can bury their dead in whatever luxury they can afford. Countless gravestones line the rolling hills of the district. They are sectioned off from the roads by tall iron fences, with mausoleums of more affluent families dotting the grassy hillsides. As surface space grew limited in centuries past, the city excavated a network of catacombs called the Undervaults. Upkeep, expansion, and general safeguarding of the sites are all overseen by the Gravewatchers, a guild of gravekeepers that have held political and social hegemony over the district for generations.

Their attempts to expand the Undervaults have had mixed results over the centuries. Most new excavations meet without incident (likely because the \textit{Clasp} has learned to keep their tunnels far from this district). However, excavations in recent decades have broken into the \textit{Crystalfen Caverns}, an eerie cave system filled with mind-bending inscriptions and horrific beings. The breach was promptly sealed off.

Yet the sealing of the caverns did not stop some members of the Gravewatchers from seeing an opportunity to gain fame and glory. Most who ventured into the caverns were never seen again, and those who have returned shudder unspeakably of the horrors they escaped. Still, passage is occasionally granted to the \textit{Clasp} or adventurers willing to pay the Gravewatchers for entry to the dangerous realm below.

\subparagraph{Port of Emon}
\textit{Emon} is bordered on the west by the Ozmit Sea and is generally sheltered from the worst of its storms. The vast \textit{Port of Emon} allows hundreds of thousands of gold pieces' worth of goods to flow through the city every single day. Well over a hundred ships fill this port at any given time, and the dock crews are ever working, carrying crates and goods away to \textit{Abdar's Promenade}, or to empty vessels for exportation. The northern sector of the Port is mostly allocated to the Everline, a powerful fisherfolk's guild. They frequently feud with the Onyx Banner, \textit{Emon's} most powerful shipping guild, over which docks they are permitted to use. Both guilds keep their hands clean, but it's an open secret that both hire adventurers and mercenaries to trash the other's ships and warehouses to intimidate them. Adventurers in need of fast cash who don't mind dodging the Arms of \textit{Emon} can find good work as muscle for either of these two guilds.

\subparagraph{Outwall}
During the reign of \textit{Drassig}, a massive city of tents and hovels grew outside of \textit{Emon}, filled with those too poor to pay \textit{Drassig's} extortative taxes. By the end of the \textit{Scattered War}, the slum had grown to a third of the city's size, and nearly half its population. Even during the reign of \textit{Zan Tal'Dorei}, these slums continued to grow, as refugees of the \textit{Scattered War} tried to earn enough money to find stability within \textit{Emon's} walls.

Though their conditions were still meager when the magnanimous \textit{Zan Tal'Dorei} ended \textit{Drassig's} iron rule, the city's denizens chose to remain, having created a community with its own culture of inexpensive living and brotherhood in poverty, away from the bustle of the city's inner streets. A microcosm of \textit{Emonian} society now exists within the city of Outwall, including its own trade square, shrines for worship, and makeshift farms on the outskirts. As early as two decades ago, Outwall was called the Upper Slums—a derogatory name given to it by \textit{Emon's} Cloudtop elite. People within Outwall tried to transform the term "Upper Slums" into a symbol of civic pride, but it never stuck. The name Outwall eventually emerged after community leaders demanded an audience with the \textit{Tal'Dorei Council} in an event that came to be called the Outer Wall Debates—for they demanded the council convene in the slums' own town square.

These days, Outwall is becoming a trendy place to live, thanks to its vibrant culture and relatively inexpensive housing. Its locals are fighting hard to keep young people with money to burn from the \textit{Central District} from gentrifying their home. It hasn't come to violence yet, but the angrier members of the Outwall community clandestinely meet with adventurers to burn down the stylish houses and mansions that are driving out local residents and businesses.

\subparagraph{Southgate Farms}
Inspired by the civic action of the people of \textit{Outwall}, the farmers of Southgate were able to break from the ugly name of "slum" in the past decade. This community of farmers and ranchers, formerly known as the Lower Slum, exists outside \textit{Emon's} southern gates. It first formed when a large refugee band from \textit{Othanzia} was barred entry into the city in the midst of the \textit{Scattered War}. Instead of dispersing, they decided to lay down roots while \textit{Drassig's} armies were out fighting a losing battle against \textit{Zan Tal'Dorei}. These squatters eventually formed a community of farmers who petitioned \textit{Zan Tal'Dorei} for land in exchange for providing produce and grown goods to the people of the city.

Thanks to the aid of the First Sovereign, the descendants of those \textit{Othanzian} refugees have become one of the most significant sources of fresh food for the people of \textit{Emon}. Frustrated by being saddled with the name of "slum," and inspired by the Outer Wall Debates, the leading families of the Lower Slums demanded recognition by the \textit{Tal'Dorei Council} as well. They were met with little debate, and these humble farmers are now a respectable part of \textit{Emonian} society—even if they live outside the city's defensive walls.

This southern region of the city is also home to a famous site of modern history: Greyskull Keep, the former base of operations for Vox Machina in \textit{Emon}. This fortress was gifted to them by the Last Sovereign as thanks for saving his life. The various heroes of Vox Machina have moved on to other homes (and now typically gather in \textit{Whitestone} for holidays), and they allowed the keep to be turned into a museum of their heroic exploits.

\subparagraph{Shoreline Farms}
Less respectable than the \textit{Southgate Farms} are the communities along a stretch of murky shore north of \textit{Emon}. These people of \textit{Outwall} have tried for decades to turn terrible, salty land into effective farmland, to mixed results. They have successfully cultivated the few vegetables that grow in saltwater-soaked soil, leading to bountiful harvests of savory quattet gourds and the hard-to-grow brineroot. However, the farms' distance from the city gives them little protection from dangers that city folk think little of. Adventurers are always wanted to protect the farms from petty criminals and wandering monsters, but the farmers have little to pay. Generally, only utterly novice adventurers find their start in a place like the Shoreline Farms—but all heroes have to start somewhere.

\subparagraph{The Grotto}
This sprawling series of underground chambers composes the \textit{Clasp's} headquarters beneath the city of \textit{Emon}. It's rumored that over a dozen entrances to the Grotto are concealed amidst the labyrinthine tunnels, which are both naturally occurring and a part of \textit{Emon's} sewers. That doesn't even begin to get into the hundreds of other entrances and access shafts abandoned and filled in for fear of discovery. The very nature of the \textit{Clasp's} business is to keep shifting out of sight, and as such, the actual core of the organization moves from base to base between multiple subterranean structures, ever building further underground, or repurposing a long-abandoned hideout whenever necessary.

\subsubsection{Points of Interest}
A number of both popular and historically significant landmarks can be found in \textit{Emon}. They are keyed to the map of \textit{Emon}.

\subparagraph{1. Laughing Lamia Inn}
While the city is home to countless taverns and inns for the weary traveler or practiced performer, few are as grand and well known as the Laughing Lamia. Almost as old as the city itself, this fine establishment has traded hands numerous times over the years, and has expanded with each transition, leaving the massive, four-story institution a hodgepodge mass of odd rooms and themes, all centered around a raucous central tavern room filled with dozens of tables. Currently under the ownership of a brash and cheerful lass called Luthania Wells, this hot spot for traveling traders and adventurers alike rarely sees an hour in the day (or night) where there isn't something interesting going on within.

\subparagraph{2. Azalea Street Park}
Azalea Street, within \textit{Abdar's Promenade}, is one of \textit{Emon's} oldest neighborhoods, and is full to bursting with small businesses, quaint homes, and hole-in-the-wall restaurants, such as the \textit{Laughing Lamia}. Amidst the chaos of daily life, Azalea Street is also home to a peaceful park overlooking the Ozmit Waterways. When the stress of constant battles becomes too much to bear, adventurers in the know take a moment to visit the Azalea Street Park and recover, maybe even swapping stories with other relaxing heroes. Legend has it that heroes who meditate among the flowers and drink in the sea air return invigorated by the spirit of the The Wildmother. A statue of the The Wildmother and The Lawbearer dancing together graces a fountain in the center of the park. Once per week, a character can gain inspiration (inspiration) by taking a resting in Azalea Street Park.

\subparagraph{3. Anvilgate}
Made of stark black stone and adorned with intricate brass metalwork, this grand blacksmithy is one of the most striking buildings in \textit{Abdar's Promenade}. Decades ago, a restless craftsman from \textit{Kraghammer} of the name Grahf Tiltkettle grew weary of solitary pursuits and worried that interest in blacksmithing would give way to the emerging trend of mass-production. Tiltkettle retired in \textit{Emon} and dedicated his respectable fortune to ignite the spark of inspiration in those who wished to learn the ways of the hammer and anvil, hoping to instill the joy of the The All-Hammer in a new generation. He funded the creation of Anvilgate, a massive, public blacksmithy where any folk who wished to try their hand at the smithing of metals could do so under the The All-Hammer guidance.

Cheaper ores and materials were often donated by friends of Tiltkettle to aid the volunteer mentors in teaching their craft to the less fortunate, while folks who wished to bring their own ingots were welcome to do so. The location has become popular enough to begin expanding its facility to include kilns for ceramics and looms for weaving. While the tools are open for anyone to use on the premises, all are warned of the curse that will beset any wanton thief who steals from under the The All-Hammer gaze.

\subparagraph{4. Traverse Junction}
The most magically active structure in Tal'Dorei shines like a pyramid of pure sapphire in the bustling center of the \textit{Erudite Quarter}, its walls thrumming with arcane energy. It is the Traverse Junction, and all visitors to the \textit{Erudite Quarter} stop to see it when they visit.

In nearly every major Exandrian city, there are \textit{teleportation circle} that link mages to other such \textit{teleportation circle} around the world. Each and every one of these major \textit{teleportation circle} has a twin in the Traverse Junction, a travel nexus for approved mages and world leaders. Any characters renowned in \textit{Emon} or honored by the \textit{Alabaster Lyceum} may make use of the Junction (for 50 gp, if they can't cast \textit{teleportation circle}), allowing them to travel to major cities such as \textit{Westruun}, \textit{Syngorn}, \textit{Kraghammer}, and \textit{Whitestone} in Tal'Dorei; Port Damali and Rexxentrum in \textit{Wildemount}; \textit{Ank'Harel} in \textit{Marquet}; and Vasselheim in \textit{Issylra}.

The telemagi who curate the \textit{teleportation circle} are always in search of new teleportation sigils to different lands; they are willing to pay a handsome sum for the services of powerful and magically skilled adventurers willing to seek out and catalog new circle sigils for their records.

\subparagraph{5. Alabaster Lyceum}
The gleaming white halls and elaborate gardens of the \textit{Alabaster Lyceum} signal to all who see them that they gaze upon the greatest institute of higher learning in Tal'Dorei—though numerous scholars in \textit{Kraghammer} and \textit{Syngorn} would heatedly dispute that claim. Regardless, this place is the heart of homegrown academia in the Republic of Tal'Dorei, boldly striving to make a name for itself alongside the ancient scholarly traditions of \textit{Issylra}, \textit{Marquet}, and \textit{Wildemount}. The \textit{Lyceum} (for those who can afford the impressive tuition) is an unparalleled center for study and research of history, economics, alchemy, art, and—most of all—magic and the arcane. Many of Tal'Dorei's greatest scholars and bards paid their dues in these halls, and the history of the \textit{Lyceum} and its graduates are featured in all manner of superstitions and folktales.

\subparagraph{6. Cloudwatch Palace}
In the wake of the \textit{Chroma Conclave's} destruction of \textit{Emon}, the towering palace that stood as home to the sovereign and the center of governance within the city fell to fire and ruin. In the decades since, construction began with the aid of the \textit{League of Miracles} to construct a new palace to help house the \textit{Tal'Dorei Council} and be a symbolic center for all major political affairs within the region. The Cloudwatch Palace now sits atop the once-smoldering hill where \textit{Thordak} once claimed his roost, a symbol of \textit{Emon's} enduring will and optimistic hope for the future. Within these high-arched halls, countless chambers hold space for diplomats, meetings, and in the case of holidays, grandiose celebrations that draw the attention of powerful figures all across Exandria.

\subparagraph{7. Thordak's Crater}
One last unsettling element of the \textit{Cloudtop District} is the elementally scarred remains of \textit{Thordak's Crater}. Even after two decades of healing from \textit{the Ashari} of Pyrah, \textit{Thordak's} fires still burn beneath the Cloudtop. The smoldering entrance to the caverns remains cordoned off from the public because of its unsightly appearance and sulfurous stench, and because fiery creatures occasionally try to claw their way out into the fresh air. The crater is often forgotten, but every few months an adventuring party makes their way inside to seek treasure or glory within the volatile caverns—and such expeditions are a favorite topic of conversation (and clandestine betting) among the politicians, political aides, remaining city nobility, and nouveau riche of \textit{Emon}.

Why \textit{Thordak's} magic lingers enough to create these mindless cinderslag elementals is a mystery; but it seems to be limited to the crater, as the elementals cannot go beyond it without collapsing into inert slurry.

\subparagraph{8. Emon Skyport}
This elevated port is complete with docking platforms solely for the use of skyfaring vessels. The original \textit{Emon Skyport} was the first of its kind in Tal'Dorei, and its successor is every bit as revolutionary as the original.

The \textit{Emon Skyport} was destroyed when the \textit{Cinder King} claimed the \textit{Cloudtop District} as his fiery throne, but like so much of \textit{Emon}, it has since been rebuilt. Unlike other parts of \textit{Emon}, this reconstruction was a joint venture between the \textit{Arcana Pansophical} and the Alsfarin Union, a \textit{Marquesian} guild responsible for the engineering and distribution of all skyship in Exandria.

Booking passage on a \textit{skyship} isn't cheap (see the "Skyships" sidebar), but the privilege is available to everyone that can pay. Generally, one passenger \textit{skyship} leaves \textit{Emon} per week to each of the following locations:


\begin{itemize}
\item Vasselheim, \textit{Issylra} (passenger ticket costs 2,500 gp)
\item \textit{Whitestone}, Tal'Dorei (passenger ticket costs 900 gp)
\item Port Damali, \textit{Wildemount} (layover in \textit{Whitestone}; passenger ticket costs 2,000 gp)
\item Ank'Harel, \textit{Marquet} (passenger ticket costs 2,500 gp)
\end{itemize}

Members of the \textit{Tal'Dorei Council} can also usually pull strings to grant adventurers they hire free passage when on vital business.

\subparagraph{9. Tomb of the Last Sovereign}
Uriel Tal'Dorei II, the Last Sovereign, was slain by the noxious breath of the green dragon Raishan during the attack of the \textit{Chroma Conclave}. Uriel's wife, Salda, survived him and petitioned the reformed \textit{Council of Tal'Dorei} to build their last sovereign a tomb befitting his benevolence and magnanimity. They complied, and the monumental Tomb of the Last Sovereign is now home to not just Uriel's ashes, but to tribute from all who loved him. Those with wealth gave gold, while those whose hearts outweighed their purses gave more personal tribute.

Over twenty years ago, the Gravewatchers closed the tomb to all, supposedly to prevent vandalism, but have not reopened it after more than two decades. The Last Sovereign's eldest daughter, \textit{Odessa Tal'Dorei}, is furious. Despite her position of power on the \textit{Tal'Dorei Council}, she is unable to dispute the act because of the power the Gravewatchers have as keepers of the dead. It has come to the point where she seeks to secretly employ adventurers to uncover the Gravewatchers' true intentions—which involve the psychic influence of an aberrant creature of the \textit{Crystalfen Caverns}.

\subparagraph{10. Godsbrawl Ring}
Though the Temple of the The Stormlord always has a fighting ring in its center, the annual Godsbrawl transforms the earthen, torchlit sanctum into one of Tal'Dorei's most unusual tournament grounds. On the Day of Challenging, the god of athleticism's holy day, the brawny priests of the The Stormlord invite warriors and worshipers of the entire pantheon to the Godsbrawl, asking that each temple offer forth their greatest warrior to act as their god's proxy in the tournament. The clergy of the The Dawnfather and the The Lawbearer send their champions, but rarely take it seriously, for they scorn the storm-priests' notion of "might makes right." Conversely, the champions of the The Stormlord and the The Platinum Dragon have a fierce—though friendly—rivalry, trading the title of Supreme Champion back and forth each year after a bloody (but rarely fatal) final round.

\subparagraph{11. The Black Bastille}
Named for its ash-blackened walls, the ominous Black Bastille is a single-story prison within the \textit{Military District} that holds Tal'Dorei's most dangerous and irredeemable criminals. The compound has no windows, no open courtyards, and only one entrance: two imposing steel doors, flanked by watchtowers. During the reign of the \textit{Chroma Conclave}, the prison was one of the first buildings attacked, setting hundreds of Tal'Dorei's worst criminals loose to sow chaos across \textit{Emon}. The Black Bastille has since been rebuilt, but dozens of its most deplorable inmates still run free, and the Arms of \textit{Emon} are eager to recover them. Some include the human demon-summoner Felrinn Derevar, betrayer of the \textit{Arcana Pansophical}; the half-elf pyromancer Illaman Falconsong, exile of the Fire Ashari; and the vampiric tiefling \textit{Ixrattu Khar}, cultist of The Whispered One and Tal'Dorei's foulest mass murderer.

\subparagraph{12. Walls of Tribute}
Anyone who wishes to uphold justice and keep the peace can apply to join the Arms of \textit{Emon}. But before a recruit is given their armor and takes a vow to protect the city and its people, they must endure a grueling training regimen within the tall, featureless stone barricades of the Walls of Tribute, referred to wryly by recruits as the Quarry. There, they are trained in the physical aspects of peacekeeping, instructed on the habits of the city's criminal underbelly, and drilled in the laws and ideals of \textit{Emon}.

\subparagraph{13. House of Discipline}
\textit{Emon} does not have a standing army of its own. The skills imparted to the Arms of Emon are sufficient to maintain a rowdy populace, but they are carefully trained to use nonlethal force to restrain, rather than kill, all those whom they apprehend. Thus, when war comes to Tal'Dorei, the council appoints a Master of War and levies an army from the Arms of Emon, as well as the militias of its constituent settlements. They must report to the House of Discipline, where they are taught martial combat alongside the harsh realities of warfare. Those who endure their time in the House of Discipline are forged into an army ready to protect \textit{Emon} and the rest of Tal'Dorei from the greater dangers of Exandria.

\subparagraph{14. Gilmore's Glorious Goods}
The broad, single-story façade of Shaun Gilmore's wondrous magical emporium belies the extravagance of its interior. Within, the air is thick with a dozen competing perfumes, each more pungent than the last. The interior is impossibly large—certainly larger than its exterior walls would have you believe—and filled with seemingly endless rows of arcane curiosities and artifacts, many of whose cryptic functions have been lost to time. Everburning candles in an array of unnatural colors light the shop, casting tantalizing shadows over every bubbling phial and mystical orb. Those who enter the shop are first met by Gilmore's long-time assistant Sherri, a half-elf draped in deep purple robes. If they are lucky, or particularly convincing, they may even meet the hero Gilmore himself.

Even rarer a sight than Sherri or Gilmore is the latter's husband, a handsome elf named Darius. He is shy in public, a perfect counterbalance to Gilmore's bombast, but is said to have a singing voice and skill with the \textit{lyre} that can make any mortal weep. Whether or not his husband is present, Gilmore is ever eager to recount his part in the fall of the \textit{Chroma Conclave}, often winkingly embellishing his role in the final battle with \textit{Thordak} himself.

Gilmore's Glorious Goods sells all manner of magic items. Its stock varies, but typically includes magic items of very rare rarity and lower for the price range listed in the fifth edition core rules. Additionally, there are always unidentified items in Gilmore's vast back-room storage that he says he'll "get to eventually;" these may include one or two long-lost legendary items. Artifact-quality items are beyond even Shaun Gilmore's usual fare, but he might know clues to a specific artifact's location, whether through his own extensive knowledge or his rumored relationship with the \textit{Tal'Dorei Council} as an expert consultant on matters of arcane objects. Magic items are expensive, and those strapped for cash may inquire about sponsorships and quests done in Gilmore's service.

\subsubsection{Emon Adventures}
Every district, point of interest, and scrap of current events described earlier in this section could be a plot seed you can turn into an adventure. In addition to those, you can use these adventure hooks to involve your players in an \textit{Emonian} adventure.

\subparagraph{Primordial Children (any level)}
Rumors of terrible creatures of fire and teeth stalking the sewers spread throughout \textit{Emon's} taverns. Such mutterings have gone largely ignored; however, an investigator has recovered what looks to be the broken remains of a large, leathery red egg, not far from a collapsed tunnel beneath the Cloudtop. It seems one of \textit{Thordak's} primordial dragon spawn survived its father's demise and is lurking in the city's tunnels.

\subparagraph{Council Business (any level)}
The \textit{Tal'Dorei Council}—particularly \textit{Tofor Brotoras}, Master of Defense, and \textit{Allura Vysoren}, Master of Arcana—frequently hire adventurers to undertake dangerous quests for the benefit of the realm. Once the characters do a deed great enough to catch the council's attention, they are contacted the next time they're in \textit{Emon} to undertake a quest; choose from any of the other story hooks in this book or make a new one of your own. The council offers a reward equal to 1,000 gp × the average character level of the party.

\subparagraph{Miraculous Escape (high-level)}
Master of Arcana \textit{Allura Vysoren} covertly contacts the characters with a worrisome claim. In the middle of the night, a former \textbf{Remnant cultist} named Hammond Kraith escaped the \textit{Black Bastille}, apparently with the aid of an invisible monster. She provides the party with the magic-inhibiting manacles he wore within the prison, in hopes that they can use them to follow his trail. What she doesn't know is that the \textit{League of Miracles} broke Kraith out of prison using an invisible \textbf{master adranach}. He is hiding out on an island in the \textit{Lucidian Ocean}, aided by a number of \textit{Myriad} mercenaries and a captured \textbf{mage hunter golem} he has reprogrammed to serve him.

\begin{DndSidebar}[float=!b]{The Sisters Grast}
Mythic tales of a pair of ancient hags swirl through learned circles within \textit{Emon}, typically to frighten \textit{Lyceum} students off of wandering forbidden vaults at night. Some folk, however, take these whispers seriously. A number of explorers and treasure seekers wandering the subterranean ruins of \textit{Crystalfen} have encountered one, or both, of these wretched witches, returning with tales of fright, dark dealings, and unholy pacts.



In reality, these creatures are elven scholars who stumbled upon the \textit{Ruins of Salar} long, long ago. Their discovery and research of the profoundly magical location twisted their minds and bodies, infused them with power, extended their lives, and granted a glimmer of knowledge of some immense, unknowable purpose to the lost city. They've become obsessed with the mysteries of the ruins. The elder, Trysta, remains below to continue to excavate and study, while her sister Forscythia travels back and forth from \textit{Emon} under the veil of a half-dozen illusions to acquire goods and gold in exchange for her unsettlingly insightful fortune-telling. Anyone who seeks the sisters with ill intent is rarely seen again, but those who stumble upon them often find themselves temporarily in service to them, a pact made out of fear and self-preservation.



\end{DndSidebar}

\subsection{Crystalfen Caverns}
Deep beneath the \textit{Bladeshimmer Coast} is a sprawling network of natural caverns and rivers that predate the \textit{Age of Arcanum}. These tunnels reach even below the ocean, and have roots in early colonization by denizens of the Realms Beyond.

An ancient society of psychically powerful \textit{aboleths} and other aberrations sprung up from a door between worlds and conquered this subterranean web, drove their slaves to construct \textit{Salar}, the Unseeable City, and slowly pressed upward toward the surface. Their plot was unintentionally ended when the final battle of the \textit{Calamity} sent powerful waves of magical force throughout Exandria, causing much of the underground caverns to collapse and the aboleth civilization to fall under rock and rubble. Fragments of the City Which Eyes Cannot See still remain, however, with surviving terrors slowly rebuilding from the dust within their once-great capital.

The ruins of this sprawling city are now known as the Crystalfen Caverns. These caves are home to scattered veins of \textit{azuremite}, a gorgeous blue crystal that formed from millennia of psychic energies existing near element deposits. Curious explorers who discovered the veins found that, when mined and refined into a fine dust, the \textit{azuremite} powder is a strong mind-altering agent and induces temporary visions and other psychic phenomena.

While a number of small, isolated entrances to this labyrinth of tunnels have been discovered and adventuring parties have attempted to chart the mines, the caverns are so vast and deep that either they gave up for fear of becoming lost, or they were assailed by the terrible denizens of the caverns, never to return. Most who now brave the caverns are foolish treasure hunters or criminals seeking to mine more \textit{azuremite} to sell on the black market.

\subsubsection{Crystalfen Caverns Adventures}
Game Masters who set their adventure in the \textit{Crystalfen Caverns} can use this plot hook for inspiration.

\subparagraph{Mining The Past (any level)}
When the characters finish a resting in the \textit{Crystalfen Caverns}, they must each make a DC 15 Wisdom saving throws. On a failure, a woman's voice fills their minds, saying, "I sense someone...are you truly here, too? Please...I am lost, and alone in the dark." The voice guides them down the twisting caves to a 120-foot-deep mine shaft. The stones at the bottom shimmer with faint azure light. Should they descend into the \textit{azuremite} mine, they travel through an uncharted tunnel and find a lone, dying \textbf{aboleth} with five levels of exhaustion that has wandered too far from its ruined domain.

The \textbf{aboleth} doesn't seek to kill them—yet. Its only priority is self-preservation, and it needs help fending off the kobolds that have made the mine their home in order to reach a subterranean lake where it can recover. Do the characters kill it or save it? If it reaches the lake, or dies, they are rewarded with a haunting premonition of the \textit{Ruins of Salar}, and a murky vision revealing the depths and expanse of the \textit{Crystalfen Caverns}.

\subsection{Ruins of Salar}
In the lightless nadir of the \textit{Crystalfen Caverns}, directly beneath the splendid \textit{Port of Emon}, rest the eternally slumbering ruins of Salar. No mortal has beheld these ruins—once the capital of a mighty \textbf{aboleth} dominion—and returned to the surface with wits to tell the tale. Salar is known only in myth and legend as the Unseeable City, because mortal eyes cannot look upon its alien architecture without destroying their fragile minds. Salar is treated as an urban legend by \textit{Emon's} scholars and archeologists. Nevertheless, many believe in its existence, for numerous scholars have returned from expeditions into the deep caverns with their minds shattered by otherworldly power.

Those who have returned from Salar brought back journals filled with increasingly crazed notes. The most intelligible of their writings indicate that all that survives of the ruined city are emerald marble structures of unfathomable design, seeming to shift and shimmer as the eye tries to comprehend, moving through time as they remain static in space.

The once-renowned archaeologist Karaline von Ethro claimed to have visions of Salar, describing a subterranean city with twisting pillars that stretched higher than any tower, kissing the seat of the surface world from the bedrock of the earth. At the peak of the central pillar floats a giant pearlescent disc. This so-called Moon Disc bathes the entire city in unsettling, milky light that refracts eerily through the emerald spires.

Karaline von Ethro is long since dead, and her journal—by far the most coherent record of the Unseeable City—is an apocryphal artifact. The \textit{Arcana Pansophical} kept it in a closely guarded vault, which only added to its mystique. However, the journal was stolen long ago, and since then, dozens of manuscripts said to be the Von Ethro Journal have been sold in black markets across Tal'Dorei. Anyone lucky—or unlucky—enough to find the genuine journal would learn a terrible secret that the \textit{Pansophical} has held onto in silence for decades: dozens, perhaps hundreds, of \textit{aboleths} survive within Salar. Their master, Durrom of the Emerald Eyes, lives on in secrecy while expanding its army of thralls and learning all it can about the ignorant surface dwellers. Durrom's ultimate goal is to shatter the forgotten Pillars of the Earth beneath \textit{Emon} and cause the city to collapse into its subterranean domain.

\begin{DndSidebar}[float=!b]{Karaline von Ethro}
Archeologist, anthropologist, and explorer extraordinaire; Karaline von Ethro was a human adventurer and inventor. It was her designs that created the Sunbeam Compass, a device that always directs the user to the sun by emitting a small light at the tip of the needle in a glass orb, allowing it to rotate in all directions—a very useful tool when exploring subterranean territory. However, despite her discoveries and inventions, she is most remembered for going mad after allegedly having visions of the Ruins of Salar. Karaline's madness eventually led to her tragic demise at the young age of 37.



\end{DndSidebar}

\subsubsection{Ruins of Salar Adventures}
Game Masters who set their adventure in the \textit{Ruins of Salar} can use this plot hook for inspiration.

\subparagraph{Secrets Within Secrets (mid-level)}
The characters are approached by Jomen Tash, an earnest half-elf mage of the \textit{Arcana Pansophical}, who claims to have obtained pages from the lost Von Ethro Journal. He shares his research with the characters and hires them to accompany him into the ruins, promising that he will guard their minds from any alien influence. But in truth, though Jomen was once a member of the \textit{Pansophical}, he has recently aligned his interests with the mysterious goals of the \textit{League of Miracles}. Is he leading the characters into danger with only counterfeit scrawlings to guide them, or does he actually hold part of Karaline von Ethro's description of \textit{Salar} and its dangers?

\subsection{Scar of the Cinder King}
It was thought that the lands surrounding \textit{Emon} would never fully recover from \textit{Thordak's} overwhelming elemental power. His presence not only incinerated his enemies—it wounded the land itself. After the \textit{Cinder King's} defeat, the \textit{Tal'Dorei Council} beseeched the Fire Ashari to travel from the distant lands of \textit{Issylra} to \textit{Emon} and lend their wisdom.

The frightful \textit{Scar of the Cinder King} and \textit{Thordak's Crater}, both twisted by the red dragon's cancerous, fiery magic, were healed by the druids of the Fire Ashari over the course of about fifteen years—an act that the brightest minds of the \textit{Arcana Pansophical} thought would take an epoch to achieve. Clouds of white ash still occasionally blow across the fields surrounding \textit{Emon}, but the last lingering reminders of the \textit{Cinder King's} terror are beginning to vanish.

Nevertheless, the Fire Ashari left behind a group of four druids, led by a female half-orc named Lorkathar. They dwell in a small stone fortress known as Flamereach Outpost—a defensive fortification used when the scar was still spawning beings of flame and hatred. They vehemently warn all who enter to avoid using elemental magic of any sort—especially fire magic—as these evocations could cause the still tender division between the Material Plane and the Elemental Plane of Fire to resonate and tear.

\subsubsection{Scar of the Cinder King Adventures}
Game Masters who set their adventure in the \textit{Scar of the Cinder King} can use this plot hook for inspiration.

\subparagraph{Reopening Wounds (any level)}
\textit{The Ashari} of Pyrah have claimed that the \textit{Scar of the Cinder King} is healed—yet rumors have spread through \textit{Emon's} \textit{Erudite Quarter} that a representative of Pyrah has been seeking counsel with the \textit{Arcana Pansophical} regarding the scar. While the characters travel through the once-scarred lands east of \textit{Emon} to find their next adventure, they are shocked by a sudden earthquake, as the ground erupts with a serpentine, fiery crack. Magma and cinderslag elementals slowly seep from the crevasse, turning all they touch to ash.

\subsection{Ruins of O'Noa (New O'Noa)}

\begin{description}
\item[Village:]
Population: 310 (83\% humans, 17\% other races)

\end{description}

\noindent The first human colonizers of Tal'Dorei journeyed across the Ozmit Sea, weathering terrible storms and deadly waves, and made landfall in \textit{Daggerbay}. From there, they founded the port city of O'Noa, and though the city quickly grew in wealth and influence, the early settlers were greeted with mistrust by the elves of \textit{Syngorn}. When King Warren \textit{Drassig} took the throne of \textit{Emon}, no city suffered greater devastation than the contested region of O'Noa. It was the site of one of the \textit{Scattered War's} bloodiest battles, and the arcane energies unleashed reduced the city to ruins and permanently scarred the land. Built within a cursed bay, and deprived of fertile land, it was agreed after the war that neither the new nation of Tal'Dorei nor \textit{Syngorn} would attempt to rebuild the city, instead letting it stand as a reminder of the horrors of the \textit{Scattered War}.

Over centuries of dereliction, the ruins became a haven for outlaws and exiles, who named their refuge New O'Noa. They took this place of destitution and dust and made it into a refuge for all rejected by society. Though some of New O'Noa's people were exiled from their homes due to fear and prejudice—like people of ill fate born under the \textit{full light of Ruidus}—most are thieves and delinquents. New O'Noa is ruled by a charismatic tiefling called Waken, who rose to prominence through honeyed words and selective murders.

New O'Noa still appears abandoned to the eyes of outsiders. Its inhabitants live and hide away from prying eyes, amidst ancient stone ruins and their catacombs. Those who venture too deeply into the ruins are rarely seen or heard from again—unless their loved ones can offer up a princely ransom.

\subsubsection{Ruins of O'Noa Adventures}
Game Masters who set their adventure in the Ruins of \textit{O'Noa} can use this plot hook for inspiration.

\subparagraph{This Too Solid Flesh (mid-level)}
New \textit{O'Noa} is a seedy town filled with strange magic. A \textit{suude} dealer named Shef Silverleaf is in major debt to Waken, the leader of the outlaw village. The characters enter the picture when Waken surreptitiously hires them to collect on his debt—perhaps under false pretenses to hide his illegal operation.

When the party tracks down Shef, they find more than they bargained for: his body has been turned to solid silver. A \textbf{Remnant cultist} lurks in the shadows of New \textit{O'Noa}, and has been purchasing \textit{suude} from Shef for nearly three months. Shef uncovered his secret identity and confronted him—but the mage's magic, made chaotic by the \textit{suude}, transformed Shef into solid silver. Waken will take Shef's solid silver corpse as payment—but what if the silver-skinned man were still alive and able to magically walk away? And what of the cultist who transformed him?

\subsection{Seashale Mountains}
A curious range of mountains runs along the \textit{Bladeshimmer Shoreline} north of \textit{Emon}. The Seashale Mountains are named for their unique shape, akin to a tidal wave about to crash on the inland territories, frozen in dynamic stone.

Crashing waves and salty sea breezes have carved miles of twisted valleys between these jagged peaks, creating an unusual home for the creatures lurking within the cavernous rock, shoreside pools, and towering bluffs. High tide leads the briny waters into the Nightwash, a miles-long lagoon that wends its way through the base of the bluffs, carrying hosts of dangerous ocean life with each pass. Conversely, the tidal phenomenon also brings a valuable bevy of deep-sea fish on occasion, which gave birth to the fishing community of the Shalesteps. The mountains' strange peaks are decorated by copses of scraggly trees, stunted by the gravelly soil and stripped nearly bare by salty ocean breezes.

\subsection{The Shalesteps}

\begin{description}
\item[Village:]
Population 770 (40\% humans, 35\% halflings, 32\% dwarves, 3\% other races)

\end{description}

\noindent The Shalesteps are a series of small, allied villages whose inhabitants chose to live a more simple life outside of the bustle of the capital city, but close enough to enjoy imports and protection should the need arise. Most of their food is sourced from the ocean and the Nightwash tide, while a small hunting community wanders the peaks for avian quarry and other sources of meat and natural resources. Those who live here are sturdy and ready to defend themselves against the dangerous wilds surrounding them—and against the occasional monstrosity that washes in with the tide.

\subsubsection{Shalesteps Adventures}
Game Masters who set their adventure in the Shalesteps can use this plot hook for inspiration.

\subparagraph{From Hell's Heart (low-level)}
The Shalestep village of Puddlefoot has never had anything remarkable happen to it. People there are quiet and simple, with no interest in tales of grand heroics. No great adventures have ever come to their rocky shore—until the whale. The dead whale, bloated with corruption, its white hide covered in oozing green lesions, is a source of fear for the nearby townsfolk. What happened to this behemoth? And why is Rill the Hermit ranting about seeing tall shadows spying on Puddlefoot from the mountain peaks?

\subsection{Visa Isle}
Forty miles off the coast of the \textit{Bladeshimmer Shoreline}, west of the \textit{Verdant Expanse}, two islands belch clouds of black smoke into the open sky. \textit{Visa Isle} and its smaller sister island have been surrounded by rumors and sailors' legends since the dawn of civilization. Some say that every night, the island burns to ash—giving it a distinctive orange glow—and that every day new trees grow to maturity by sundown. It is an island of fire and death.

Those who have explored its dense jungles and returned claim only to have found common wild pigs and—at worst—unusually large fire beetles and mantises. They were the lucky ones. They didn't wander unprepared into the Ruins of Vos'sykriss, and did not encounter the serpentine spirits that haunt its elder halls. The only expedition to have ventured into the Ruins of Vos'sykriss and lived to speak of it returned to \textit{Emon} just as the \textit{Chroma Conclave} descended upon the city. Their leader, a half-orc archaeologist and spiritspeaker named Jorlund Vohr, survived the destruction of \textit{Emon}, but all his notes and the artifacts he recovered were stolen. Decades later, the aging professor is still hard at work creating endless proposals for expeditions to \textit{Visa Isle}.

\subsubsection{Visa Isle Adventures}
Game Masters who set their adventure in \textit{Visa Isle} can use this plot hook for inspiration.

\subparagraph{Awakening the Sleepers (mid-level)}
In ancient times, Vos'sykriss was the seat of a serpentfolk empire that dominated southern Tal'Dorei. Before the \textit{Calamity}, their seers had visions of the coming destruction and desperately sought a way to survive it. The serpents constructed a magical stasis field beneath their city and chose their strongest and most powerful to slumber and rebuild the empire when the danger had passed.

They are now awoken. The \textbf{Vos'skyriss serpentfolk} now prowl the island, being prepared by the lingering Vos'skyriss serpentfolk ghosts to claim the world for their own. When Jorlund Vohr's expedition delved into Vos'sykriss's deepest sanctum, they unwittingly ended the serpents' ancient spell and awakened the Sleepers. Soon, a people long thought extinct will return to claim the world as their own.

\part{Wildemount}
\chapter{Welcome to Wildemount}
\begin{DndReadAloud}{}
\DndDropCapLine{T}{}wo armies clash on a fog-shrouded battlefield, heralded by the sound of trumpets and drums on one side, and by an eerie, chirping whistle like countless thousands of wailing crickets on the other. A raven, its ebon feathers spattered with blood, alights on one of the fallen corpses and gently raps its beak on the soldier's helmet. The raven looks up, curious, and flies into the sky. A land embroiled in war expands beneath it, and yet, within the shroud of conflict, the raven sees twinkling pockets of light.


The raven sees souls that shine bright with courage, kindness, and determination. It sees souls lustrous with greed, fear, and anger. And most of all, it sees souls radiating with a desire to discover their place in the world. The raven sees Wildemount.

\end{DndReadAloud}

\section{A New D\&D Setting}
Wildemount is a new setting for your Dungeons \& Dragons campaign, originally created by Matthew Mercer as the setting for the second campaign of the hit series Critical Role. Though it's just one continent in the larger world of Exandria, Wildemount is teeming with varied cultures, locales, peoples, and monsters. It is home to both avatars of evil and paragons of good, but it is defined by the common people who exist in between these two ideals.

This book is for D\&D players and Dungeon Masters of all experience levels. It is for newcomers and for hardcore fans of Critical Role alike. If you've watched every episode of Critical Role to date, this book will still contain surprises for you, as there are corners of Wildemount that even the Mighty Nein weren't able to explore in their adventures.

If you're a Dungeon Master, this book will help you create a D\&D campaign set in a land of war, intrigue, and swashbuckling adventure. You could set your campaign entirely within one of Wildemount's four regions, or you could follow the example set by Critical Role and create a campaign with a constantly evolving tone as the adventurers journey across the continent. You will also find brand-new monsters and introductory adventures that can help you start a campaign in any of Wildemount's four regions.

If you're a player, this book contains new options that will help your character fit in perfectly with this strange new land. You can use the heroic chronicle system to create a backstory that firmly roots your character in the lands of Wildemount. You'll also find new subclasses as well as new magic items that will help you attain greater and greater power as you travel the continent.

\section{Nations of Wildemount}
Wildemount's various regions are described in greater detail in \textit{chapter 3}. Here's a quick overview of the major civilizations of Wildemount.

\subsection{Dwendalian Empire}
\begin{DndReadAloud}{}
The raven soars over a massive city, sprawling outward from a somber, monolithic palace. Soldiers and war machines pour from the city, stretching across the land like innumerable grasping arms, greedily encircling the continent. The raven notes this hungry behemoth and alters course.

\end{DndReadAloud}

The Dwendalian Empire is a young, ambitious state borne from decades of strife in Western Wynandir—the heartland of Wildemount. Its austere lands are dotted with tiny villages, ruins of ancient civilizations, and pockets of uncharted wilderness.

Campaigns set in the Dwendalian Empire are bound to involve elements of war and political intrigue. The empire is at war with the Kryn Dynasty, a nation of surface-dwelling drow which occupies Eastern Wynandir and the wastes of Xhorhas. This war allows various Dwendalian factions, such as the power-hungry Cerberus Assembly, to vie for political dominance.

\subsection{Clovis Concord}
\begin{DndReadAloud}{}
Westward flies the herald of the Raven Queen. It bursts through the clouds above the western mountains and emerges into a land blessed by sun. A sea of lush jungle stretches out beneath the raven, until eventually green gives way to golden sand, and then to a boundless, sapphire sea. Ships swarm on that sea like a cloud of gnats, engulfed by flames and the sound of cannon fire. They fight for freedom.

\end{DndReadAloud}

The Clovis Concord is a republic of city-states lining the Menagerie Coast, the westernmost shore of Wildemount. This tropical region is rich with history and known for a lively, free-spirited culture that contrasts sharply with Wildemount's other, politically narrow-minded nations.

Campaigns set on the Menagerie Coast are often nautical, exploring the conflict between the law of the Clovis Concord and the chaos of the Revelry pirates. They can also investigate the jungles that line the coast and the ancient ruins hidden within their verdant boughs.

\subsection{Tribes of Shadycreek Run}
\begin{DndReadAloud}{}
The raven, battered by ice and snow, sees petty warlords lead bands of slavers across the snowfields and understands that no law but power rules the north. It sees smugglers and thieves flee beyond the grasp of civilization and into the unknown. It sees the vast, roiling wilderness act with a fell will of its own, claiming the lives of those who dare attempt to grasp power beyond their ken. They fight for survival, glory, and greed.

\end{DndReadAloud}

The Tribes of Shadycreek Run are a fragmented coalition of thieves, slavers, and scoundrels who managed to amass power in the frigid backwater of Shadycreek Run. While they struggle to enrich themselves through crime, however, they fail to realize that malign, primordial powers are creeping into the north underneath their noses.

Campaigns set in the Biting North often involve morally gray characters butting heads with equally amoral organizations in a struggle to survive—and then explode into chaos as forces beyond mortal imagination make all these petty conflicts seem insignificant by comparison.

\subsection{Kryn Dynasty}
\begin{DndReadAloud}{}
Finally, the raven flies east. It sees beings that would be slain on sight in the other civilized nations of Wildemount, walking side-by-side with humans and gray-skinned elves. The raven sees that these people have a light within them—and it sees that, far to the east, there is a mighty fortress that contains a light unseen since the dawning of the world. Here, they fight for the revelation of the end of the world.

\end{DndReadAloud}

The Kryn Dynasty is an empire made up of dark elves (drow) and other humanoids typically considered to be monsters by the other nations of Wildemount, such as minotaurs, goblinkin, and gnolls. The Kryn Dynasty rules over the wastes of Xhorhas, a land once blighted by the rule of the Betrayer Gods when immortals still walked the land. The Kryn safeguard a secret; if your players haven't seen any of Critical Role, revealing this secret could mark the turning point of a campaign.

A campaign set in Xhorhas will doubtless involve the war between the Kryn Dynasty and the Dwendalian Empire, but it can also include traditional dungeon crawls, political intrigue within the courts of the dynasty, or even wilderness exploration that could unearth secrets lost since the fall of the Betrayer Gods.

\section{War!}
The continent of Wildemount is a powder keg about to explode. Tensions between the Dwendalian Empire and the Kryn Dynasty have just erupted into open war, and everyone in Wildemount, including the adventurers, must contend with the fallout of the conflict between the continent's two imperial superpowers.

\subsection{The War of Ash and Light}
The war is only as important as you make it. If you place either the battles of the war or their consequences front and center, the war will be of chief concern to your players. If you keep the war and its consequences distant, then the characters will be free to interact with or ignore the war as they see fit. With a little bit of creative tweaking of the setting, you could create a Wildemount campaign where the war has been called off or never even happened!

If you want the war to be an important part of your campaign, you need to give your players a reason for their characters to care about it. The easiest way to get buy-in from your players is to discuss out-of-character how big of a factor they want war to be in this campaign. You can even dial the importance of the war up and down over the course of the campaign, depending on how you and your players feel over time.

\subsubsection{Bring It On!}
If you and your players find war to be dramatic and compelling, urge your players to create characters with strong ties to the conflict. For example, patriotic (or even jingoistic) characters will always fight for king or queen and country. Likewise, greedy characters will fight for profit and glory. Once they've given you their reasons, turn those motivations into plot points in your campaign.

\subsubsection{In the Background}
The adventurers don't have to be active combatants for the war to affect their lives. See below for a list of adventure options that touch on the consequences of war without directly involving mass warfare.

\subsubsection{No War, Please}
If you want the war to be completely absent, simply dial the timeline back by a year, or even just a few months. All the political tensions between the Dwendalian Empire and the Kryn Dynasty still exist, but now they're just simmering instead of at a full boil.

\subsection{Adventuring Options}
Not all campaigns set in Wildemount need to deal with the war directly. That is, the characters aren't expected to conscript with one army or another, or even serve as mercenaries. The effects of war are far-reaching and multifarious, and loose-cannon adventurers are the perfect people to get involved in the countless conflicts sparked by war. These schemes could include:

\subparagraph{Crime}
In times of war, criminals and smugglers take advantage of lax law enforcement. Petty criminals and crime syndicates like the Myriad will need help moving illicit goods in and out of Dwendalian and Xhorhasian cities, and Revelry pirates will need loyal crew members to raid unprotected towns.

\subparagraph{Disaster}
When a town is destroyed by advancing armies, or by a landslide caused by artillery, the common folk need help escaping. This could involve escorting noncombatants through no-man's-land to a nearby town, or carving out a new settlement from the wilds.

\subparagraph{Law}
When officers of the law are conscripted as soldiers, local lawmasters suddenly find themselves in need of mercenary aid. Lawful characters might find great adventure in flushing criminals out of cities and busting their wilderness hideaways.

\subparagraph{Military}
Obviously, characters can choose to join up with one of the sides in this conflict to embark on missions and take orders from a commanding officer. This is a good option if you want to run an episodic campaign, where most game sessions begin with a stated mission that the players must accomplish and ends with a reward and another mission.

\subparagraph{Upheaval}
The explosions of artillery and the pounding of boots against the plains can spook monsters, and invasions can cause even powerful monsters such as giants and dragons to flee their lairs. Whenever an upheaval occurs, soldiers and commoners alike find their lives in danger—and only adventurers can deal with the threat.

\subparagraph{Trade}
Like the crime syndicates, opportunistic traders know that war is a chance to become rich. Adventurers are the perfect pawns for trade barons looking to sell high-value goods, especially if they have to cross dangerous land to reach their buyers.

\subsection{How to Run a War Campaign}
If there's one thing D\&D adventurers hate, it's being told what to do, but serving in an army is all about following orders. It's up to the Dungeon Master to make sure the characters are aware of what they're getting into if they decide to enlist. Similarly, if you force your players to join a military unit and then discover that the characters actually want to go on self-directed adventures without a commanding officer ordering them around, you and your players should talk out-of-game to figure out a way to turn the campaign around.

Since D\&D is primarily a game about a small group of characters going on adventures alone, it can be difficult to simulate massive battles using D\&D combat rules. Because of this, it's generally best to keep the characters away from mass battles. However, huge conflicts with thousands of combatants are a cornerstone of epic fantasy, and your players might be disappointed if your war campaign doesn't have at least one climactic battle.

To solve this problem, you can break down your mass combat into manageable chunks. Find a significant location that the characters can either defend or conquer with minimal reinforcements, like an overrun citadel. Then, have the major battle proceed in waves that guide the characters from one cinematic encounter to another. You can think of these encounters like rooms in a dungeon; some rooms have multiple doors that the characters can choose from, while others only have a single passage.

For example, after capturing the ruined fortress, the characters learn that a garrison of soldiers has arrived to help hold onto the fort, and the characters are needed elsewhere on the battlefield. You could present them with just a single urgent target, such as an enemy warlord, or you could offer them options, such as three different allied battalions under attack.

If you need to depict a lot of both allied and hostile NPCs battling in the same encounter as the players, you can simplify the encounter by ignoring dice rolls entirely and deciding when a creature dies based on your own instinct. This can work once in a while, but try to avoid doing this all the time, since seeing you simply decide the result of combat without dice can damage the believability of your game world.

\begin{DndSidebar}[float=!b]{When Does This Book Take Place?}
This book is set in a certain moment in the history of Wildemount, and thus setting a campaign in a different era—or even in a different decade—could make some information in the gazetteer invalid. The present year is 835 PD ("Post-Divergence"). The history of Wildemount is described in more detail in \textit{chapter 1}.



From a real-world perspective, this book is set at a certain point during the adventures of the Mighty Nein in \textit{Critical Role}'s second campaign. Thus, the point at which the adventures you create using this book diverge from canon is at episode 50 of campaign 2. That's okay! Look at it another way: by creating your own Wildemount adventures, you are boldly creating your own canon, from which all the events of \textit{Critical Role} diverge into myth and rumor.



\end{DndSidebar}

\section{Daily Life in Wildemount}
Certain aspects of daily life in Wildemount remain the same across nearly all cultures. Any character native to Wildemount would be familiar with their nation's currency, languages, and relationship to technology.

\subsection{Currency}
The nations of Wildemount all mint their own coins, and though some slight regional differences exist between currencies, all regions readily accept gold, silver, and copper coins. Spending money of Dwendalian mint in the Kryn Dynasty might turn heads, but a gold coin is still a gold coin.

Platinum and electrum coins are rarely minted by modern nations, but enough of both of these currencies survived the Age of Arcanum that both are considered valid tender across the land.

\subsection{Languages}
All the typical languages found in the \textit{Player's Handbook} and the \textit{Monster Manual} are fair game in Wildemount. While Common is used across the continent and typically associated with human empires, it is not the language that most humans spoke in Wildemount before the rise of Dwendalian Empire three centuries ago.

This setting has three languages unique to the human cultures of Wildemount: Zemnian, Marquesian, and Naush. A human character can learn one of these languages instead of a skill or tool proficiency granted by their background, class, or variant racial traits. Each of these languages is described below:


\subsubsection{Zemnian}
This ancient language was spoken by the people of Zemniaz in the Age of Arcanum. That ancient culture has long since crumbled, but its language and its people live on in the Dwendalian Empire. Many ancient scrolls were written in Zemnian, but it is a language now largely spoken by farmers, as Common is the default language of the empire.


\subsubsection{Marquesian}
The Menagerie Coast was settled by colonists from the arid land of Marquet, and their language now holds an unusual position in the Clovis Concord. It is the language of high society, as many Clovis elite are descended from Marquesians, but it is also the language of piracy, since countless lower-class Marquesians defected from the Concord and formed the Revelry pirates.

\subsubsection{Naush}
Originally spoken by the Ki'Nau islanders native to the Menagerie Coast, Naush is a thriving language within the multicultural cities of the Clovis Concord. Even sailors who only speak Common incorporate dozens of Naush words into their nautical jargon.

\subsection{Technology}
The level of technology in Wildemount is generally consistent with the technologies found in the \textit{Player's Handbook}. However, Dwendalian scientists in Hupperdook and Concordian tinkerers in Port Zoon have made incredible advances in black powder technology. Cannons, mobile war engines, and even handheld pistols and muskets have begun to be used as weapons of war.

Black powder weapons are not common, nor are they available to the general public. Only military engineers and special regiments within the Dwendalian armies and the Concordian navy have access to these powerful weapons, though some aspects of these designs are now finding their way into the Xhorhasian military. Statistics for weapons such as black powder barrels, pistols, and rifles are provided in \textit{chapter 9} of the \textit{Dungeon Master's Guide}. All items listed as "Renaissance Items" on the Firearms and Explosives tables exist in some form in Wildemount. Since these items can't be purchased in a normal store, characters can only obtain them through theft, military service, or a quest of your own devising.

\chapter{Story of Wildemount}
\DndQuote%
{}
{\enquote{Born from nothing, there was something. Light, then thought, then life and fire. We who are born of the will of the gods keep these records to preserve their deeds, their design, and their lessons. We who live by our own will shall keep these records to know where we walk and where not to tread. We who look to our inevitable end keep these records to teach those who succeed our footsteps how to be better than we could.}}
{Archivist Adia Shu of the Cobalt Soul}


\section{Wildemount after the Calamity}
In the wake of the vast destruction of the Calamity and the exile of the pantheon in the Divergence, the survivors began to emerge from the ash and shadows to reclaim and rebuild Exandria. Across the continents, many took the pieces of their cultures and sought a place to forge a new age. This tome focuses on the rebirth of the deeply scarred land now known as Wildemount, the site of the most terrible conflicts of the Calamity, where the echoes of those deadly battles linger to this day.

\subsection{In the Wake of the Gods}
After the Calamity, the lands of Wildemount became wild and perilous, filled with monstrosities and beasts that prowled its broken fields and shattered mountains. The land's new denizens reclaimed the newly untamed realm, hunting any mortals not clever enough to hide. The masterless creations of the Betrayer Gods ruled their own territories, establishing themselves as a new pantheon.

As the dust settled, surviving mortals grasped at what they could from their bygone lives and hid away in fear of what would come next. Those of faith dealt with their sense of failure and abandonment, while enlightened scribes mourned the loss of their research and the great magical secrets that had enabled the previous age to rise to such prominence.

Others picked up what they could from the ruins and chose to start anew, hoping that their descendants would learn from their mistakes. Over those early centuries, the scattered survivors fought back the terrors that stalked the abandoned lands of Wildemount, continuing to endure and build new societies, civilizations that eventually reclaimed Wildemount and brought the world to this modern age.

\subsection{Western Wynandir}
The lands of Western Wynandir were the epicenter of the Calamity's destruction. Rocky ranges of mountains were sundered, entire landscapes burned, and hovering seats of power crashed down from the sky, their floating cities consumed beneath the hungry earth.

\subsubsection{Grimgolir}
The dwarven clan of Grimgol that once helped hold the center of the Ashkeeper Peaks was nearly wiped out during the Calamity. The surviving families and warriors burrowed deep beneath the surface of Exandria to wait out the chaos. Collapsing nearly every established tunnel used in the wars above, the clan found themselves cut off from the rest of Exandria. For centuries, the dwarves of Clan Grimgol endured the darkness alone, defending against the dregs of the Betrayer Gods' forces and adapting themselves into a hardy and stalwart people.

Half-remembered tales and myths of their forgotten history, mingled with worship of Moradin, began to stir an eagerness to return to their ancestral home and build anew. With new purpose, they carved their way back toward the surface, only to find the mountains of their home reduced to dust and pebbles.

Not easily dissuaded, the dwarves returned to the surface for the first time in generations, finding the once-ruined world now recovering with renewed life and color. The resilient clan made their way north along the Brokenveil Bluffs and found a massive mountain to house their new city. They built their stronghold into this monolith among the Dunrock Mountains, and gave their new home the name of Grimgolir.

\subsubsection{Rise of the Julous Dominion}
In the Marrow Valley, the numerous decades following the Divergence left the landscape blackened and inhospitable. The remaining people struggled for food and shelter as they wandered the rocky hills for centuries, avoiding the colder weather of the north. Nomadic sects of humans and halflings defended themselves against the hungry beasts and lingering evils threatening their homesteads, and fought with each other over fertile land.

Small civilizations began to form amid the chaos, but it wasn't until the central township of Zadash was established as a safe outpost that the unification of the Marrow Valley began. Helmed by the Julous family, supposedly of noble blood from before the Calamity, a guarded community with an organized military force was built in the center of the valley around four hundred years ago. It became both a haven for the struggling masses and a center for agricultural expansion.

Increasing their strength with each generation, this movement guided by the Julous family became known as the Julous Dominion, stretching across the south end of Western Wynandir. Not all agreed to join the Julous Dominion, however. Some outliers rejected the Julous interests, founding their own village of Kamordah on the boiling earth of the hills in the southwest Marrow Valley. In time, the people of Kamordah came to worship a being of primordial fire that lives beneath the mountains and speaks to them in dreams and visions.

\subsubsection{Founding of the Dwendalian Empire}
In the frigid fields north of the Marrow Valley, the survivors of the Arcanum city of Zemniaz scraped by in the snow for centuries, clutching their surviving scrolls and fragments of cultural history. They established the dour outposts of Icehaven and Yrrosa, but it wasn't until the Zemnian people wandered east to the Pearlbow Wilderness that they found fields that were fertile enough to nurture a new society.

Founded atop the ruins of an ancient temple to the Dawn Father, the city of Rexxentrum rose to become the hub of Zemnian culture. Now that the Zemnian people no longer had to struggle for survival amid the harsh elements, their leaders argued over how to establish a centralized government. Some wished for a republic, while others sought a theocracy under the priests that shepherded the people into this new age. Eckhardt Dwendal, a well-respected merchant and tremendous force behind the development of Rexxentrum, wished to institute a monarchy and place his son Manfried on the throne. Debate spun wildly for over three years, while Eckhardt privately made deals with the heads of the emerging guilds, using promises or blackmail to secure their cooperation. With the resounding support of the flourishing guild heads, the throne was established. In the year 539 of the post-Divergence era, Manfried became the first emperor of the Dwendalian Empire.

\subsubsection{The Marrow War}
It wasn't long before the expanding borders of the Dwendalian Empire began to push up against the growing Julous Dominion. An agreement was brokered to avoid conflict as long as favorable trade could be established and maintained. Even so, the resource-rich center of Western Wynandir was split between the two states, with both factions eager to discover the other's weaknesses. Rumors of the discovery of a massive platinum mine within the Julous Dominion further increased the empire's desire to conquer their neighbor.

At this time, religious leaders throughout the Dwendalian Empire rebelled against the emperor's religious restrictions. After his coronation, Emperor Manfried Dwendal enacted laws limiting worship and prayer to subjects concerning the power, profit, and expansion of the empire. In the year 544 PD, internal strife and insurgent action turned bloody when the zealous leaders of varying faiths attempted, and failed, to usurp the throne for their council. An angered Emperor Manfried Dwendal had the rebellious priests executed—an event that came to be known as the Admonition.

The emperor spoke of abolishing all worship within his boundaries. This caused further uproar among the populace, but then Emperor Manfried's spies discovered that the seeds of the uprising had been sown by Julous interests. Publicly exposing the plot, the emperor outlawed only select faiths that he felt threatened civilized life. He redirected much of the people's fury by also proclaiming the Julous Dominion an enemy of the empire later that year and declaring war on the neighboring nation.

Meanwhile, the Julous Dominion was dealing with the refusal of the outlying villages near Kamordah to accept their rule in this time of brewing hostility. They eventually drove the independent villagers back into the Cyrios Mountains, just in time to hear the horns of war blaring from the empire to the north. The violent conflict known as the Marrow War lasted over sixteen months and cost the lives of many soldiers and innocents on both sides of the Marrow Valley.

The Julous Dominion lacked the unified military force of the empire. When a major military push left the capital city of Zadash occupied by the empire, Baron Inock Julous and Baroness Tessandra Julous met with Emperor Manfried to negotiate a peace. During this meeting, the baron and baroness were unceremoniously executed. Their bodies were displayed publicly along with the narrative that they had attempted an assassination on the emperor during the peace talks. Disheartened and bereft of leadership, the Julous Dominion immediately came to an end, as did the war.

To prove his benevolence to the people of the newly conquered Julous Dominion, and to help avoid future uprisings, Emperor Manfried declared that all citizens of the dominion would keep their lands and homesteads. No soldiers would stand trial or be punished for their previous allegiance. Life would continue as it had, except for the implementation of imperial taxes, laws, and oversight across the valley under the new emperor. Outlawed religious iconography was destroyed, crops and goods were redistributed according to the needs of the empire, and life slowly returned to normal for the Zemnian people.

\subsubsection{A King Is Crowned}
With citizens of the former Julous Dominion tense beneath their new rule and a growing sentiment of frustration with imperial leadership building, Emperor Manfried sought a way to reform the nation under his banner. The propagandist minds who fought the social war with the Julous Dominion recommended a shift in public image and title, one that could command more respect within Western Wynandir while casting aside the oppressive connotations of the title of emperor. Still riding high on the victory against the Dominion, Emperor Manfried harkened back to the halcyon lore of the Age of Arcanum, recalling the beloved kings and queens of old. It was decided that the people would no longer have an emperor, but a king. Crowning himself the first of the line of Dwendalian kings, King Manfried insisted on maintaining the now-ominous national title of the Dwendalian Empire, famously proclaiming, "I am the Dwendalian king to my people, beloved and open to their wishes. To our enemies, we remain resolute and unstoppable. We are the Dwendalian Empire."

\subsubsection{The Eve of Crimson Midnight}
Not a quarter century after the Marrow War ended and the line of kings was established, a smaller, internal conflict rocked the capital city of Rexxentrum. A number of noble houses with a strong history of studying arcane pursuits began to compete with other high-born magic practitioners from the Julous Dominion. Throughout the lands of the newly expanded empire, these houses escalated their subterfuge and espionage against one another, until the rivalry finally erupted into an all-out magical conflict in the streets of the capital itself.

The events of the Eve of Crimson Midnight destroyed numerous buildings and maimed a number of innocents caught in the crossfire. The struggle ended with all involved shackled and brought before the king. After days of deliberation, an agreement was drawn up that would absolve those involved of the usual punishment in exchange for direct subservience to the Crown and the goals of the empire. Establishing themselves as the Cerberus Assembly, this council of mages became a powerful tool for the empire to maintain its position as the dominant force of Wildemount.

\subsubsection{Modern Dwendalian Empire}
For three hundred years, the empire has expanded and prospered under the rule of the Dwendalian line of kings. King Alfwin Dwendal increased the empire's military power by establishing the stronghold of Bladegarden and assimilating Grimgolir into the empire, seeking to hold back the savage creatures of Xhorhas. King Theoderich Dwendal's benevolent nature brokered the partnership with the Clovis Concord of the Menagerie Coast without bloodshed and ushered in new opportunities for trade with the rest of Exandria. King Willamar Dwendal established the initial prominence of House Truscan and granted them oversight of the Truscan Vale, to the displeasure of the other noble houses.

Other kings have come and gone, but few have made as immediate and lasting a mark as King Bertrand Dwendal. On assuming the throne in 790 PD, King Bertrand immediately demonstrated his love of spectacle and social manipulation by instituting town criers across the empire to disseminate Crown-approved messages. With the aid of Master Ikithon Trent of the Cerberus Assembly, he captured the imagination of the masses with well-crafted tales of his deeds, proclaimed by the criers in every town square.

Seeking to bury any dissent against the Crown, King Dwendal began to use the empire's spy network to seek out sources of unrest in the public and punish those who sowed thoughts of rebellion. Public holidays and celebrations were given greater funding to quell the masses and distract the people from the increase in taxation implemented across the valley. This allowed King Bertrand Dwendal to rise in popularity while also ruling with a far more ruthless and self-serving agenda than most kings before him.

King Bertrand maintains the isolationist world view of his predecessors, even in times of need. By ignoring political alliances from distant nations, preventing the construction of any skyship docks within the empire, and ignoring calls for aid—such as the fight against the rising threat of Vecna two decades before—the King believes he has made the empire stronger than it's ever been.

However, the king's focus on isolationism fades as tensions with Xhorhas grow. Worry about his legacy mingles with the new and deadly threat of the Kryn Dynasty, now emboldened and pushing into his kingdom with strange magics and dangerous warriors. He mistrusts the very Cerberus Assembly that grants him so much influence and arcane might to combat the ever-darker world that barks at the borders of the empire. Public approval of his rule wanes with growing economic disparity, and the whisper of Myriad power returning under his nose only heightens his paranoia.

When the Kryn orchestrated an attack on the Halls of Erudition in Zadash, the Cerberus Assembly pushed for a display of military power to rally the people and bring unity to the national identity of the empire. The king declared war on the Kryn Dynasty and the people of Xhorhas, embroiling the region in this new full-scale conflict, with the hope that, as history has proven, people are easier to unite against a common enemy.

\subsection{The Menagerie Coast}
While destruction came to the farthest reaches of the world during the divine conflicts of the Calamity, the islands of the Lucidian Ocean off the western coast of Wildemount thought they might fare better than most. But though they were far removed from the central fury of the magics unleashed in the final battles, the region's civilizations were instead wiped out by massive floods. Immense tidal surges left the coast drowned and empty, the islands consumed by the raging ocean, leaving behind only the most accomplished of sailors. Eventually, the waters settled once more, and nature swiftly reclaimed the lands with lush green jungles.

\subsubsection{The Ki'Nau}
The few scattered survivors of the floods found shelter among the higher jungles of the Swavain Islands, gifted with the bounty of a more plentiful and tropical environment. These seafaring people rebuilt among the islands and took the name Ki'Nau, meaning "the Water Children" in their language of Naush. Though the terrible floods of the Calamity had passed, the deadly seas and flourishing horrors of the wilds were still a persistent threat to the Ki'Nau. Each generation struggled to survive the elements that sunk their boats and the predators of the jungles and ocean depths.

Over time, an older mind began to call from the waters and entered their dreams. This powerful being, a leviathan creation abandoned by his creator Zehir in the Divergence, called himself Uk'otoa. He spoke of his divinity and promised protection to the Ki'Nau people in exchange for their worship. The Ki'Nau accepted the persuasive entity as an ally and protector, allowing Uk'otoa to rise up as the guardian of the islands and help the Ki'Nau become a proud society of warriors, establishing dominance over the region for centuries.

\subsubsection{The Marquesian Alliance}
Around the year 400 PD, an exploration vessel from the distant land of Marquet was nearing the end of its provisions when it reached the Swavain Islands. Discovered and captured by the ever-watchful Ki'Nau, the desperate explorers made trade offerings of foreign gold, spices, and fine silks. This diplomatic approach won over the leaders of the Ki'Nau, and an alliance was made with the Marquesian travelers to return with trade goods and establish a Marquesian outpost, which was to be named Damali.

During this time, a cruel seafaring culture called the Vukan discovered the Ki'Nau and coveted their lush homesteads and rich, sea-gathered jewelry. These Vukan began a series of coordinated assaults on the Ki'Nau, killing and kidnapping a number of them under the shadow of night. The Marquesian guilds who had sent the exploration ship agreed to offer naval support to combat these attacks, but it was Uk'otoa whose fury ultimately annihilated the Vukan at their home villages and erased them from Exandria.

Having lost a portion of their fleet in these fights, the surviving Marquesian naval ships remained to build small settlements along the coast to aid in their trade with the Ki'Nau. The Ki'Nau leader Tauan kept a tentative peace with the foreign sailors, but Uk'otoa wanted to dominate the new arrivals who refused to renounce their old gods and bow to him. These tensions came to a head when a secret sect of Marquesian worshipers of the Cloaked Serpent Zehir were given a task by their enraged god from beyond the Divine Gate: to destroy Uk'otoa for his insolence. With the aid of a powerful cult leader within an archeological guild known as the Allegiance of Allsight, a series of temples built to Uk'otoa were corrupted with Zehir's influence and used to seal Uk'otoa deep beneath the rock of the ocean floor, robbing the Ki'Nau of their patron.

\subsubsection{Birth of the Clovis Concord}
At this time, the city of Ank'Harel in Marquet had just named the mysterious J'mon Sa Ord as their new emperor. In the transition, the Marquesian colonists now living along the coast demanded independence from Marquet to begin building their own society and identity alongside the Ki'Nau. Ord eventually accepted the request, and the emancipated Marquesians allied themselves with the Ki'Nau to build a new civilization. Over the following centuries, numerous port towns sprung up and grew into city states, trade routes expanded along what is now known as the Menagerie Coast, and the emerging guilds consolidated enough power to become the region's ruling class: the Clovis Concord.

The concord works together to keep the surrounding Lucidian Ocean under their watch and rule, setting trade laws and guiding the fortunes of those who sail on the open waters. The blend of cultures that forms the Menagerie Coast creates a diverse community where the arts are lauded and coin is king. However, the shadow of the Myriad syndicate continues to infiltrate the streets, while the pirate alliance known as the Revelry has risen to harry the trade routes and rebel against the concord. Tensions are further escalated by the threat of expanding war to the northeast between the Dwendalian Empire and the Kryn Dynasty.

\subsection{The Greying Wildlands}
The fierce powers released in the conflicts of the Calamity tore through the heavily forested region now known as the Greying Wildlands, setting ablaze a slow-burning arcane forest fire that smoldered for nearly a century, burning through the hearty trees of the region like a plague of embers. The northern realm was left desolate, ashes and dirt blown into dust and sand as the chill northern winds transformed the rolling hills into a desert of ice and pebbles. Fleeing remnants of the Grimgol dwarves delved into the snow-hooded Flotket Alps, many freezing to death in the wind-battered valleys before finding safety in the natural caverns hiding deep in the mountain range. Within these icy depths, the dwarves cultivated fiery hearths and brought warmth to the heart of the mountain, establishing the new dwarven city of Uthodurn.

\subsubsection{Molaesmyr}
Amid the charred remains of the burning forest that filled the sky with smoke for over a hundred years, one section of the woods defied the flames and held against the embers that battered it. This small patch of idyllic green endured, a sanctuary for the surviving wood elves that held fast and hoped for a miracle. Sensing that the blessings of Melora and Corellon had kept the heart of this forest alive until the fires burned out, the elves began to build a new home among the ancient trees and groves.

This region became the Veluthil Forest, and from this verdant sanctuary, the elves would build the stunningly beautiful city of Molaesmyr. From the surrounding ashes, the fey magics of Veluthil brought forth new life and leaf, expanding the boundaries of the forest over centuries to eventually reclaim much of the Wildlands. Druidic forces and powerful fey enchantments were laid throughout the twisting paths of the forest to protect the elves of Molaesmyr as they expanded their towering city, the central trees reaching magically enhanced heights as homes and halls extended high into the boughs. The elven city stood as the height of reborn civilization on Wildemount.

\subsubsection{The Corruption}
It is still unclear what triggered the event. Whether some knowledge-hungry researcher or arrogant sorcerer was delving in forgotten vaults between the planes, or some terrible power lying dormant beneath the forest was awakened, the source of the corruption is unknown. What is known is that in the year 585 PD, suddenly and without warning, a wave of purple-gray shadow rapidly crept from the center of Molaesmyr to engulf the entire city.

Many elves choked and died on toxic fumes, while others were twisted and driven mad, transforming into terrifying fey abominations. Trees bent and warped, leaves withered and turned a sickly purple, and strange, unknown bramble growths reclaimed the forest floor. This corruption slowly spread throughout the entire Veluthil Forest, decimating the home of the elves. Some survivors fled as fast as they could to the northwest, across the Frigid Depths, eventually settling the township of Bysaes Tyl in the Pearlbow Wilderness. Others escaped to the northeast, seeking refuge in the Flotket Alps, freezing in the snowy mountains until they stumbled on the dwarves of Uthodurn.

Taking pity on the survivors, the dwarves took in the elven refugees and offered them a home. This act of kindness changed the course of Uthodurn's history, eventually forming a unique union of dwarven and elven culture that has endured for nearly three centuries. Craftsmanship within Uthodurn is second to none, but nothing is traded beyond the borders of the city, aside from the odd stolen artifact. Small excavation parties are often sent to the corrupted Veluthil Forest—now called the Savalirwood—to search the ruins of Molaesmyr, attempting to piece together what caused the corruption and determine if the forest can be recovered.

Unfortunately, a township of criminals and outlaws has sprung up on the southern edge of the Savalirwood, just beyond the Dwendalian Empire. In addition to the dangerous denizens of this wretched thicket, which is known as Shadycreek Run, there are numerous raiders and thieves seeking relics from the ruins of the forest, escalating conflict throughout the region.

\subsection{Eastern Wynandir}
The once thickly forested landscape of Xhorhas was the site of the final battleground of the Calamity. It was subjugated to such intense destruction at the hands of the warring gods that when the dust of the Divergence cleared, Eastern Wynandir was naught but a cracked, ashen wasteland. Ghor Dranas, the citadel of the Betrayer Gods, lay broken, cursed, and abandoned, while lingering demons and monstrosities fled to the surrounding jagged cliffs and blackened marshlands. Discarded abominations seeped from beneath the ruined earth of Xhorhas's eastern shores, climbing past the Penumbra Range where the forgotten creations of the Betrayers were locked away.

Throughout this wasteland, scattered bands of aimless men, beasts, and goblinkin battled for shelter and sustenance. These tribes were confined to their desolate home by the Ki'Nau leviathan who guarded the waters past the southern shores and the surviving dwarves of Grimgolir, who prevented any Xhorhasians from moving farther west. The nomadic, warring clans fought against the harsh conditions of the valley for centuries, only able to develop small bastions of civilization within the harsh wilderness.

\subsubsection{The Kryn Dynasty}
Beneath the ruins of Ghor Dranas, a group of scavenging drow turned from Lolth, their primary deity for centuries, finding renewed faith in an esoteric entity they called the Luxon. Guided by the light of their new deity, the drow survived their ascent from deep within the Underdark and reclaimed the halls of Ghor Dranas for themselves, naming their new home Rosohna, or "Rebirth." The leading house of Kryn helped build a new nation of dark elves who sought the surface as part of their worship, and slowly began to reach out to the nearby nomads to unite under the light of the Luxon.

As the first nation of post-Calamity Xhorhas, the Kryn Dynasty worked to expand its reach throughout the wasteland and bring unity to the ravaging hordes. Dormant evils continued to rise from beneath the marshes and badlands of Xhorhas, spreading chaos and battling with the warriors who called the desolate valley their home. The struggle to convert the denizens of Eastern Wynandir to the light of the Luxon continues to this day, while the agents of Lolth seek to spite and destroy her once-chosen for their arrogance.

Meanwhile, the dragonborn have begun to reach out to Exandria's sprawling political world, wishing to exchange knowledge and trade their gold and jewels for powerful magics to better their standing.

\subsubsection{Draconia}
After the Calamity, a shard of the destroyed dragonkin city of Kethesk came to rest within the southern Dreemoth Ravine. The surviving dragonborn took advantage of the ravine's defensible layout to build a new city and made an alliance with the cave-dwelling, indigenous dragonborn they called ravenites, beginning the reformation of civilized dragonborn society.

With the ravenites as a labor force, the dragonborn harvested the massive brumestone crystals that had once held their city aloft, and the new city of Draconia began to take shape as a series of floating islands within the heart of the chasm. Calling themselves draconbloods, they spent hundreds of years slowly reconstructing their nation in seclusion, focusing entirely on their own enlightenment and pursuits. They eventually enslaved the ravenites to mine the ravine's gold and gemstones, and to ensure draconblood dominance in the region, leaving their minds free to focus on internal politics and arcane advancement.

\subsubsection{The Chroma Conclave}
In the year 815 PD, an unexpected assault on Draconia came in the form of a group of ancient chromatic dragons calling themselves the Chroma Conclave. The mighty city of Draconia was immediately destroyed and left to crumble in the Dreemoth Ravine. A powerful white dragon called Vorugal, the Frigid Doom, claimed the ravine as its new domain, while the rest of the dragons sought conquest elsewhere in Exandria. These events shook the nations of Wildemount, sending the Clovis Concord and the Dwendalian Empire into a state of alert, while the Kryn Dynasty temporarily withdrew to Rosohna, fearing another attack.

As time passed, Vorugal and the entire Chroma Conclave were slain and the danger seemed to have dissipated. With Draconia sundered, the no-longer enslaved ravenites greatly outnumbered their previous masters; they claimed their freedom and punished the surviving draconbloods. A handful of draconian nobles fled to the Menagerie Coast or the Dwendalian Empire to seek asylum, leaving the ravenites to rebuild Wildemount dragonborn society on the ruins of their oppressors.

\subsubsection{The War of the Ash and Light}
In the past century, the Dwendalian Empire has clashed with the Kryn Dynasty. Small skirmishes and slain soldiers on both sides deter expansion beyond the Ashkeeper Peaks, with each faction keeping a close eye on the other.

Recently, the empire's interest in growing Kryn influence has led to increasing curiosity about the source of the Kryn's mysterious arcane capabilities. Dwendalian spies successfully stole two of four artifacts known as "luxon beacon," which seemed to be powerful sources of the unique magical energy called dunamis. The Bright Queen Leylas Kryn began an initiative steal back the beacons and assault Dwendalian strongholds until these sacred relics are returned. Now, in 835 PD, the resulting skirmishes have escalated into all-out war. A fifth beacon has been recently uncovered by the Cerberus Assembly, and the violence threatens to grow past the Ashkeeper Peaks and embroil other denizens of Wildemount in the conflict.

\chapter{Factions and Societies}
\DndQuote%
{}
{\enquote{Queens, kings, emperors, nobles... beneath the rhetoric, it's all the same. They say they fight to protect your families. They say they fight to enforce your beliefs, your honor. Fight for a better tomorrow. Fight for virtue, for prosperity. Bottom line, it's all about fighting, and it's you that dies fighting for their cause.}}
{Unknown soldier, Collected Tales of the Marrow War}
From the ashes of a land devastated by the Calamity grew the civilizations of a new age. The survivors claimed land across the continent, forged alliances, and made enemies as they fought for power and influence. Wildemount is shaped by the ideals and interests of its factions, from powerful governments to covert societies. This war-torn continent is a fertile landscape for adventurers on a quest for greatness.

\section{Dwendalian Empire}
The Dwendalian Empire has reigned over Western Wynandir for over twelve generations. Imperial rule is maintained through its powerful military, its policy of religious restriction, and the support of the secretive Cerberus Assembly. Though these institutions nominally exist to ensure the safety of the populace from foreign threats, terrible dangers that wander the wilds, and lingering evils that seep from the scars of the Calamity, they also allow the empire to maintain strict social order within its territories.

King Bertrand Dwendal is now two months past his sixty-eighth birthday, and his iron will has not wavered with age. A tall, powerful man, King Dwendal rewards loyalty and swiftly punishes insolence and failure. His family has developed a totalitarian rule that reaps heavy taxes and looms over society with a military presence that only grows stronger with proximity to the capital city of Rexxentrum. While civic unrest does rise beyond mere mutterings from time to time, reminders of the dangerous creatures that live beyond the guarded borders of each city and the protections offered within the empire generally keep the populace in line. Those that rebel often find themselves in chains or forgotten in a dungeon.

Dwendalian society has a rigid caste system, in which denizens who have mastery over the arcane are considered elite members of society. Many Dwendalian ladies and lords came from a sorcerous bloodline or studied the arcane arts within the walls of the Soltryce Academy. Those who especially excel in their magics or were born into prominent families can come to hold power as a member of the Cerberus Assembly, a council of extremely powerful mages that works directly with King Dwendal and his councilors to develop and enforce laws throughout Western Wynandir. However, distrust between the king and the Cerberus Assembly runs deep, for they both wish to eliminate the other and enjoy uncontested power.

Most cities and townships throughout the empire are ruled by a local government under the leadership of a starosta, usually a native baron appointed by the Crown and the Cerberus Assembly. Reporting directly to the king's council, the starosta is given absolute control over local government, appointing and maintaining responsibility for the heads of the city guard (watchmaster), the regional military (warmaster), local commerce (coinmaster), infrastructure (pithmaster), religious practices (idolmaster), and the courts of law (lawmaster). Twice a year, tithe collectors (nicknamed "Reapers") visit each imperial city to inspect households for records and proof of income, collecting a percentage of their earnings in coin as determined by the tithe collector.

The empire's military force currently stands at around 35,300 Righteous Brand soldiers across all regiments, and 5,200 Crownsguard across all cities. The majority of the Righteous Brand forces were deployed to fortifications along the Xhorhasian border at the onset of the current war.

Soon after the Dwendalian bloodline came into power, there was a short-lived rebellion helmed by religious civilians. It was swiftly quelled and the rebel leaders were executed in an event known as the Admonition. The ruling class has regarded religion with suspicion ever since, believing that faith inspires the unworthy to rise above their station and spark fruitless rebellion. Instead of abolishing all religious practices, the empire regulates worship within its boundaries. All temples are owned and run by the government, and must be dedicated to deities approved by the Crown. Priests and clerics are on the Crown's payroll to preach only approved religious texts. Keeping private shrines or worshiping outside official temples is punishable by fines or imprisonment. The only deities that are approved by the empire are Erathis the Law Bearer, Bahamut the Platinum Dragon, Moradin the All-Hammer, Pelor the Dawn Father, Ioun the Knowing Mentor, and the Raven Queen, Matron of Death.

\subsection{Goals}
The empire uses the plentiful resources of the Marrow Valley and sizable tithes gathered from citizens to protect their people and maintain control over Western Wynandir, while keeping a keen eye out for opportunities to expand their influence across the continent.

The empire's religious proscriptions are proving more and more difficult to enforce in this time of conflict. Theolocrat Kon Bruda is helming a new directive to align religious texts with less oppressive terminology, in an attempt to redirect discontent away from the Crown and onto the heretics beyond the borders.

Now that tensions with Xhorhas have erupted into all-out war with the Kryn Dynasty, the Crown is spreading propaganda through all aspects of daily life to unite the people of the empire against their common enemy. They hope to garner support for the war by reinforcing rumors about the kidnapping of Dwendalian children and the Kryn's violent and cannibalistic nature.

\subsection{Relationships}
Though many citizens gripe about the empire's restrictive laws, the majority have grown accustomed to a comfortable lifestyle under imperial protection. Fear of the evils that stalk the edges of civilization generally outweighs any frustrations about civil liberties, and the relative peace provided within the confines of Western Wynandir is welcomed. However, dissidents have begun to speak loudly of their discontent, and the growing trend of uproar is making King Dwendal quite irritable.

Landlocked and reliant on the trade avenues of the Menagerie Coast, the Dwendalian Empire maintains a respectful and mutually beneficial relationship with the Clovis Concord. Agents of both nations watch the borders, and an occasional political gesture of good faith satisfies those in power. Still, the empire doesn't trust the concord by any stretch of the imagination, and a secret world of subterfuge is alive and well.

Historically, the Cerberus Assembly has been responsible for much of the empire's acquisition and retention of power, but the royal bloodline has always harbored a paranoid dislike of the council of mages. Handshakes and smiles hide distrust, and the king correctly assumes that the assembly is merely biding their time to attempt to challenge his power.

The King and his chamber ignore most foreign powers. Requests for alliances from the Council of Tal'Dorei and the distant city of Ank'Harel have been met with indifference, and efforts to install a skyport in Rexxentrum to enable easier foreign travel have been continuously shut down.

\subsection{Figures of Interest}
A number of personages maintain the power and influence of the empire, some working to expand its power, while others seek to redeem their country.

\subsubsection{King Bertrand Dwendal}
\textit{Lawful neutral, male human}

\noindent
King Bertrand was raised with the strict beliefs of the Dwendalian royal family, and his mounting responsibilities as the leader of one of Exandria's greatest nations weighs heavily on him. Bertrand is a dark, clever, and often paranoid king—and a lonely one since his wife Duvia passed away nearly ten years ago. Fiercely nationalistic and proud of his heritage, the king harbors a disdain for outsiders and their challenges to imperial rule; he often has emissaries thrown from the borders or imprisoned until a replacement is sent to barter for their release. Bertrand does genuinely wish to protect the lives of his people and be remembered as the best of his bloodline. However, civil unrest is brewing, and his frustrations with the Cerberus Assembly only intensify his forceful demeanor. His only living son, Eidys, has been unable to produce an heir, and King Bertrand fears that his family line will come to an end, an obsession which drives him to seek out a way to prolong his life.

\subsubsection{Chamberlain Queselia Vir}
\textit{Neutral good, female elf}

\noindent
Chamberlain to the previous king and now to Bertrand, Queselia handles all affairs regarding the management of the house of Dwendal. Trusted deeply by the king, Queselia refuses to give into the brooding tendencies of her monarch and hopes to guide him toward a more just rule. Although she is unable to protect him from his own unreasonable impulses, she hopes to find a way to curb his self-destructive behavior.

\subsubsection{Prince Eidys Dwendal}
\textit{Neutral good, male human}

\noindent
Heir to the throne of the empire, Eidys has shown great strength in physical combat and war games, but little cunning in social or political matters. Though he was once eager to please his father, the consistent lack of approval has left Eidys jaded. Happy to live an extravagant life free from crushing responsibility, he's been eager to produce a child to whom he could abdicate should his father pass. However, he and his wife, Princess Suria Dwendal, have been unable to conceive. The king has all but given up on his son, leaving the prince dejected and bitter.

\subsubsection{Princess Suria Dwendal}
\textit{Chaotic neutral, female human}

\noindent
Wed into the royal family from her own noble house of Saugiss, Suria was filled with pride to bring honor to her family. Unfortunately, the unexpected difficulty in producing an heir has put her and her husband out of favor with the king. Unwilling to give up, Suria is seeking to remedy the situation by any means necessary, no matter how unnatural.

\subsubsection{Prime Arbiter Sydnock Truscan}
\textit{Lawful evil, male human}

\noindent
The Truscan family are longtime allies of the Dwendalian bloodline, and Sydnock managed to advance his family's influence within the empire above most other noble houses. Rising to become the Prime Arbiter of Rexxentrum, Sydnock has final say on all matters of law that do not require the king's attention. A brutal man with a grudge against religious empowerment, Sydnock is constantly attempting to sniff out corruption within the approved structure of worship within the empire. His estranged daughter Lydia disagrees with his ruthless methods and now lives in Zadash as an advocate for civilian rights.

\subsubsection{Crown Marshal Hayden Damurag}
\textit{Lawful neutral, nonbinary human}

\noindent
As a decorated veteran of the imperial military who has served for nearly twenty years, Hayden is a highly respected soldier and clever tactician who is instrumental in the campaign against the Kryn Dynasty. Known to be stern but adaptable in the face of adversity, they've risen through the ranks and now directly advise the king on all major military actions. Hayden prefers to avoid court life, especially as their practicality often clashes with Sydnock Truscan's cold cynicism. Currently assigned to oversee the defense of the besieged eastern front at the Rockguard Garrison, the Crown Marshal is doing their best to keep their soldiers safe and the enemy at bay.

\subsubsection{Exchequer Aethia Drooze}
\textit{Lawful neutral, female dwarf}

\noindent
Aethia is in charge of the tithes from the various cities within the empire and is the overseer of the biannual collections by her tithe collectors. She absolutely loves numbers and wealth, taking joy in squeezing every copper she can from those who are legally obliged to pay. Obsessed with efficiency and intolerant of sloppy work, Aethia will even comb through the ledgers of her subordinates late into the night in order to achieve numerical perfection.

\subsubsection{Guildmaster Kai Arness}
\textit{Lawful evil, male human}

\noindent
Kai came to prominence within Rexxentrum through the Arness family's illustrious connections within the merchant guilds of the empire. He earned the nickname "Goldfang" for his penchant for sneaking missing profit from the pockets of businesses unable to pay membership fees and taxes. Respected and feared for ruthlessly governing the many guilds and organizations that drive the Dwendalian economy, his private worship of Asmodeus guides his ambition to mingle with dangerous figures in the empire. The guildmaster often works with Exchequer Drooze to maintain the proper flow of gold into the coffers of the empire.

\subsubsection{Argonomist June Wiston}
\textit{Chaotic good, female half-elf}

\noindent
June is seen as a studious spitfire, and she has unexpectedly proven to also be a masterful bureaucrat with a talent for exploiting legal loopholes. Driven by a youth in poverty, she swore to help repair a broken system of agricultural distribution. By making the right moves and the right friends—and making a loud fuss when necessary—she eventually became a folk hero to the underprivileged. Hoping to watch her eat her words and collapse under the pressure of imperial politics, the Crown offered her chief management over all agriculture in the position of Argonomist. Much to their chagrin, she's proven both driven and capable, pushing for advanced change in the empire. Unable to discard someone so high profile during a time of general unrest, the Crown instead endures her work while attempting to bury her in red tape.

\subsubsection{Theolocrat Kon Bruda}
\textit{Lawful neutral, male human}

\noindent
Once a monk of the Cobalt Soul, Kon Bruda's faith in the Knowing Mentor has waned since accepting the office of Theolocrat. Kon is a specialist in the historical and social workings of modern religion. When the previous Theolocrat passed away, he accepted the post out of a sense of duty, then became accustomed to the lavish lifestyle afforded by his position. He helps the Crown alter existing religious texts and write approved sermons and prayers for distribution throughout the empire.

\subsubsection{Emissary Lord Zeddan Graf}
\textit{Neutral good, male human}

\noindent
A beloved noble and socialite from the House of Graf, Zeddan showed prowess in gathering rumors for his own entertainment. The cunning Sydnock Truscan saw promise in Zeddan's talents and recommended him to the king to become the empire's lead emissary. Now acting as the diplomatic contact to the world beyond the empire's boundaries, Zeddan is eager to do the Crown's bidding while rooting out gossip that may prove politically expedient. In his travels, Lord Zeddan often crosses paths with Lord Athesias Uludan of the Cerberus Assembly, and the pair loathe each other deeply. Zeddan is desperately trying to convince the king to install a Skyship Tower in Rexxentrum to make his travels easier.

\subsubsection{Oliver Schreiber, Retainer to the Crown}
\textit{Lawful neutral, male gnome}

\noindent
Most know Oliver as the immediate head servant to the king and the royal family for nearly three generations, but Oliver is also the secret head of the Augen Trust, the Crown's elite spy network. The Augen Trust uses extreme secrecy, false identities, and infiltration into allied and opposing societies to acquire information and inform the Crown of any threats to its goals. As one of the most trusted members of King Bertrand Dwendal's circle, Oliver holds more power than his compatriots would believe.

\begin{DndSidebar}[float=!b]{Laws of the Dwendalian Empire}
As judged by the eyes of the local lawmaster, appointed hands of that lawmaster, or the Law Bearer herself, any criminals discovered and/or convicted of crimes within local halls of judgment throughout the empire are subject to punishments befitting the crime. All fines and punishments for crimes against a Crownsguard member, Righteous Brand soldier, or noble are doubled. All fines and punishments regarding crimes against an imperial official or associate of the Cerberus Assembly are tripled.


\begin{description}
    
    
    \item[Theft]: 5–10 days incarceration and/or a fine of twice the worth of the stolen goods.
    
    
    
    \item[Worship of non-empire gods:] 30 days incarceration and a fine of 250 gp.
    
    
    
    \item[Ownership of another intelligent humanoid creature:] 12–20 months incarceration and a fine of 1,800 gp.
    
    
    
    \item[Trespassing]: 12 days incarceration and/or a fine of 215 gp.
    
    
    
    \item[Failure to pay taxes to the Crown:] 5–30 days incarceration and a repossession of property equal to the amount owed.
    
    
    
    \item[Damage to private property:] 5–15 days incarceration or a fine of twice the worth of property destroyed.
    
    
    
    \item[Damage to Crown property:] 15–60 days incarceration and a fine of twice the worth of property destroyed.
    
    
    
    \item[Kidnapping]: 60 days incarceration and/or a fine of 800 gp.
    
    
    
    \item[Assault with intent to injure:] 30 days incarceration and/or a fine of 350 gp.
    
    
    
    \item[Assault with intent to kill:] 250 days incarceration and/or a fine of 1,800 gp, or execution.
    
    
    
    \item[Murder]: 8–15 years incarceration and/or a fine of 6,000 gp, or execution.
    
    
    
    \item[Treason]: execution.
    
\end{description}


\end{DndSidebar}

\section{Kryn Dynasty}
The Kryn Dynasty rules the northern wastes of Xhorhas, and has only revealed itself to the wider world within the past century. Centered around—and under—the ruins of Ghor Dranas in Eastern Wynandir, the dynasty governs many townships and small villages where the Kryn have helped establish a modicum of civilized living among the nomadic wastefolk. They believe that beings not yet beyond redemption can be turned to the light of the Luxon, and often struggle with the unruly denizens of the wastes, hoping to gain converts to their cause and faith. Those that do accept the dynasty often do so because they seek protection from the dregs of the Calamity that wander the wastes, as well as the scouts of the Dwendalian Empire.

Dark elves are the most populous race within the Kryn Dynasty. The drow were the first people to uncover the first buried Luxon beacon many ages ago, in a time before they turned from Lolth and escaped enslavement from the Betrayer Gods. The dark elves claimed the dread citadel of Ghor Dranas as their own, as well as the massive network of underground caverns that make up the Underdark of Wildemount. The Kryn are stealthily becoming a force to be reckoned with.

The Kryn drow who emerged from the shadowed depths of the caverns beneath Xhorhas now endure periods of sunlight as part of their worship, though their cities are shaded by umbral magic during daylight hours. When hunting, or in times of conflict, the warriors of the Kryn don their signature chitinous armor. A series of small airways woven through this armor can be opened to emit a loud, clicking buzz to intimidate their foes as they charge. This noise has inspired outsiders to call the Kryn "cricks," referring to this cricket-like din.

Artifacts left behind by the Luxon that hold its essence are known as Luxon beacons and are coveted by the Kryn Dynasty as central to their civilization. If a proven soul is ritualistically bound to a Luxon beacon, they can enter a process called the consecution: when a mortal bound to the Luxon dies within a hundred miles of a beacon, the beacon ensnares their soul from its intended afterlife and prepares it to be reborn in a new child within that same vicinity. When that child then matures into adolescence, they begin to recall subtle knowledge from their past life's experiences, a process called anamnesis. Through meditation with a guide, they can unlock the memories of their past lives and become a more complete being. When a soul has cycled enough to believe they are ready to greet the Luxon, they are called an Umavi, or "perfect soul," and are revered within the Kryn society. Some beings that naturally reincarnate, such as the phoenix, are considered holy creatures that carry a shard of the Luxon within themselves.

If a life continues long enough to achieve a state of clarity called achess, where they believe that they've learned all they can learn from their current life, they can choose to partake in a ritualistic termination of their current life, assisted by an Umavi in proximity to a beacon, thereby setting them on the path to a new birth. While considered an honor within the Kryn Dynasty, this practice is misunderstood and seen as barbaric outside their culture.

The Kryn Dynasty is a society restructured from the Lolthite history they rejected. The people of the dynasty stand under the first Umavi, known as Bright Queen Leylas Kryn. Beneath the first Umavi, there are twelve noble Dens that mark families of souls who have lived together through generations: Dens Kryn, Mirimm, Thelyss, Hythenos, Icozrin, Biylan, Tasithar, Duendalos, Daev'yana, Omrifar, Beltune, and Zolaed. When someone is believed to be experiencing anamnesis, they submit to a deep meditation with a guide trained in the ways of consecution. This process helps restore memories of the past life, which then mingle with current memories.

The awakened soul is then reunited with their Den. Dens do grow as new souls enter the consecution and join their ranks, and new Dens are established with each generation. The oldest Dens, especially those helmed by an Umavi, are held in the highest esteem, and generally answer directly to the Bright Queen. The older the soul, the more prestige it holds in the dynasty. Beneath the Dens of Rosohna, guilds and Den representatives govern the other cities of the dynasty that stretch across Xhorhas, spreading the faith of the Luxon and watching over all Xhorhasians who seek the safety of civilization.

This seemingly benevolent societal structure does have a dark side. A closely guarded secret of the highest Umavi is a challenge in the consecution called typhros: some souls that have undergone multiple cycles of the consecution are driven mad when their minds cannot reconcile the memories of many lifetimes. Those who fall to typhros are quietly removed from the dynasty and either left in the wastes or mercifully slain away from the beacons. Signs of madness are beginning to show in the Bright Queen herself, but those closest to her do all they can to shield others from such portents.

It is believed that the Luxon cherishes unity in community, but also values free choice and the power of the individual mind. For each soul to carve a unique path is the best way for a soul to learn about the world. Casting off the teachings of Fate and Destiny taught by some of the current pantheon, the followers of the Luxon are taught that the real power of the universe comes from the power of choice. Many possible futures all await your path, and the universe trembles with anticipatory power as you approach every decision. This power is called dunamis, the energy of potentiality, and the coalescing architecture of the multiple futures leading into each moment is released when an outcome is brought to reality. It is the energy that suffuses every Luxon beacon to prepare the lives to come, and it is the basis of the Kryn's unique divine crafts and arcane abilities.

The Kryn Dynasty has uncovered four beacons and is certain that more remain to be discovered. It is believed that once all the beacons are brought together, the Luxon will be summoned from their slumber to ask their children the great question and impart the truth. It is said that at this time, the Luxon will take those who entered the consecution and abandon this lesser world to start a new world elsewhere.

The Aurora Watch legionnaires of the Kryn currently number around 21,600 in total. Around 13,200 Aurora Watch are divided among the cities under the dynasty's protection, with another 3,300 Aurora Watch stationed in and around the capital of Rosohna—formerly known as Ghor Dranas—to keep the peace and defend the region. Approximately 5,100 soldiers are scattered throughout the Underdark beneath Wildemount at different locations of interest, many employing magically controlled purple worms to allow elite units to navigate the caverns quickly. The remainder of the watch is stationed along the border of Western Wynandir to prevent the empire's incursions—and to seize land from the empire as they advance.

\subsection{Goals}
The primary objective of the Kryn Dynasty is to seek out and recover the Luxon beacons presumed to be scattered across Exandria. Not only does each beacon bring them closer to the eventual rebirth of the Luxon, they also help broaden the boundaries of the Kryn society as a new center of the consecution, allowing expansion into new lands and the continued conversion of others. With two of their four beacons now missing, the Bright Queen and her people are desperate to discover more of these sacred artifacts.

The Bright Queen and her fellow Umavi see the strength dormant in the roving beastfolk and goblinoids of Xhorhas, and have taken to teaching these nomadic peoples their ways, helping them build villages under the banner of the dynasty, and converting them to their faith. Even so, not all ruling Dens like the idea of introducing what they consider lesser creatures into the consecution, leading to tensions between the people of Rosohna and the scattered villages in the Xhorhasian wilds.

With some beacons stolen by operatives of the Dwendalian Empire, the entire Kryn Dynasty has mobilized to retaliate. They are gathering armies to assault Western Wynandir while sending assassins and spies through the subterranean paths of the Underdark, hoping to recover the beacons and simultaneously weaken the evil and blasphemous empire.

The Kryn Dynasty has spent a generation training and preparing a select group of volunteers to gather information on their enemies and find routes through both political and physical borders. These agents are collectively known as the Lens, a network of spies and infiltrators who can bypass boundaries and assimilate into foreign nations to gather intelligence, help the Bright Queen protect their people, and seek out beacons across Exandria.

\subsection{Relationships}
Open war has recently broken out between the Kryn Dynasty and the Dwendalian Empire—a conflict that neither nation can afford to lose. Those bearing the colors of the enemy are subdued and questioned before being punished and sometimes executed. Judgment is swift against the perceived imperial oppressors, and the forces of Rosohna will continue to push until their foes yield the stolen relics.

Though the Kryn have remained fairly isolated throughout their history, they are now reaching out to other factions in the hopes of making alliances against the Dwendalian Empire. The Kryn Dynasty has begun hiring mercenaries through the Tribes of Shadycreek Run and purchasing dangerous weapons of war from the Myriad. Rumors of strange magics and powerful relics in northern Eiselcross have reached the dynasty's ears, and they are willing to seek out any potential advantage over their Dwendalian foes.

There are many humanoid tribes that roam the southern reaches of Xhorhas that have declined to ally themselves with the Kryn—or to be assimilated by them. While the Bright Queen has allowed them to continue their neutrality for the time being, she watches them closely, lest they be courted by the empire's coin and empty promises.

\begin{DndSidebar}[float=!b]{Laws of the Kryn Dynasty}
Within the sight of the Luxon, the Umavi, and the Dens who protect the lands of Xhorhas and its people, any who are presented for judgement before the presiding officials entrusted with the wisdom of the Bright Queen are to have their sentence carried our accordingly. All fines and punishments for crimes against an Aurora Watch soldier or noble are doubled. All fines and punishments regarding crimes against a Den official or member of the Bright Queen's council are tripled.




\begin{description}
    \item[Theft:] 2–5 days incarceration and a fine equal to the worth stolen.
    
    \item[Ownership of another intelligent humanoid creature:] 10–20 months incarceration and a fine of 1,500 gp.
    
    \item[Trespassing:] 5 days incarceration and a fine of 120 gp.
    
    \item[Damage to private property:] 5–20 days incarceration or a fine of twice the worth of property destroyed.
    
    \item[Damage to dynasty property:] 12–50 days incarceration and a fine of twice the worth of property destroyed.
    
    \item[Kidnapping:] 60 days incarceration and/or a fine of 800 gp.
    
    \item[Worship of Betrayer Gods:] 150 days incarceration, or execution.
    
    \item[Assault with intent to injure:] 12 days incarceration and a fine of 200 gp.
    
    \item[Assault with intent to kill:] 200 days incarceration and a fine of 1,200 gp, or execution.
    
    \item[Murder:] 5–10 years incarceration and/or a fine of 4,000 gp, or execution.
    
    \item[Treason]: execution.
\end{description}



\end{DndSidebar}

\subsection{Figures of Interest}
There are many powerful figures that guide the faith of the Luxon and the expansion of the Kryn Dynasty's power and influence.

\subsubsection{Bright Queen Leylas Kryn}
\textit{Lawful neutral, female dark elf}

\noindent
Leylas was present when the first beacon was found a millennium ago, and has lived eight lifetimes of intense struggle and self-discovery since. She built the dynasty around her journey, spearheaded the recovery of the four beacons known to the Kryn, and died numerous times in the defense of her people. Leylas Kryn is deeply revered and respected as a mighty warrior, an empathic leader, and the elder Umavi who guides her people in the light of the Luxon. Her behavior has grown erratic at times, worrying her trusted inner circle, but she is far too driven to acknowledge such weakness in a time of war.

\subsubsection{Dusk Captain Quana Kryn}
\textit{Neutral good, female dark elf}

\noindent
An extremely powerful echo knight and general of the forces of Rosohna, Quana has been Leylas's partner for three of her lifetimes, and five of Leylas's. She is wholly devoted to her people—and her queen. Quana holds a private hope that the Kryn might one day leave the inhospitable wastes of Xhorhas and find a proper place in the world among the other peoples of Exandria. The Bright Queen was once receptive to her ideas, but Leylas Kryn has no time for idealism in this age of war. Quana is disappointed, but she is more worried by the strange shifts in her partner's usually wise and controlled demeanor.

\subsubsection{Shadowhand Essek Thelyss}
\textit{Neutral evil, male dark elf}

\noindent
Just over a hundred and twenty years old and still within his first life, Essek is a prodigy who displays an unprecedented talent for the manipulation of dunamis energy. Essek is both respected and feared for his intelligence and cunning, and he carefully plots each step he takes to further his and his Den's climb up the hierarchy of the dynasty. Essek is more interested in dunamancy magic than worshiping the Luxon, and he is strangely excited for the war with the empire. He is eager to use the conflict as an excuse to practice the deadlier aspects of dunamancy, and also curious to see what the powerful minds of the Cerberus Assembly may have gleaned from their research into the beacon they stole.

\subsubsection{Sunbreaker Olomon}
\textit{Neutral, male minotaur}

\noindent
Having lived a number of lifetimes as an early follower of Leylas, Olomon has settled into this recent life as a massive minotaur, embracing the abilities granted to him on the battlefield and taking his place as a champion on the front lines. His presence inspires his fellow warriors, and his cunning has led the armies of the Kryn to many victories against the empire.

\subsubsection{Taskhand Verin Thelyss}
\textit{Lawful good, male dark elf}

\noindent
An acclaimed soldier and talented leader, Verin Thelyss is the younger brother of Essek and the youngest soul of Den Thelyss. Verin lacks his older brother's ambition but possesses a keen sense of justice and honor. Recently given command of Kryn forces in Bazzoxan, Verin hopes to bring about a victory for the dynasty that will grant him respect and acceptance from his accomplished family.

\subsubsection{Skysybil Abrianna Mirimm}
\textit{Lawful neutral, female goblin}

\noindent
Seen by many within the dynasty as a talented diviner and voice of reason, Abrianna is an Umavi nearly as old as the Bright Queen herself. The beloved Skysybil has spent her current lifetime as a goblin, a fact that has drawn both ire and jests from the Bright Queen's court for decades. This experience has given Abrianna a new perspective on the wastefolk under the dynasty's banner; she is now one of the few councilors in the Lucid Bastion who bother to vouch for the goblinkin.

\subsubsection{Professor Tuss Waccoh}
\textit{Neutral, female half-orc}

\noindent
A scatterbrained scholar whose subtle brilliance has increased the effectiveness of the dynasty warmachines, Professor Waccoh usually oversees a number of academic programs and private research projects at the Marble Tomes Conservatory within the Firmaments of Rosohna. Recently obsessed with the desolate, cursed lands of Blightshore and the mysteries that lie hidden within, Tuss is actively seeking allies and mercenaries to send into the dangerous lands to recover information, relics, and specimens for further study.

\section{Cerberus Assembly}
The Cerberus Assembly was established nearly three hundred years ago, when strife and competition between warring mage lords threatened to unravel the social order of the Dwendalian Empire. The house leaders of this conflict nearly destroyed one another in what was called the Eve of Crimson Midnight, an arcane battle within the capital of Rexxentrum that leveled two wards of the city, killing hundreds of civilians. At the peak of the violence, a parlay was called, and the mage lords decided to put aside their differences and avoid mutually assured destruction, instead joining as a unified conclave. When brought before the king for judgment of their crimes, they offered an accord to instead officially instate the conclave and submit entirely to king's rule. Taking the name the Cerberus Assembly, a reference to a powerful, historical union of mages during the Age of Arcanum, the mages took the empire to new heights of power and expansion.

In the centuries since, the Cerberus Assembly has subtly ingrained itself within most facets of the empire's government, becoming a force of political and military power that many say could rival the king himself. This reality is not lost on the royal council, and tensions often simmer between the assembly and the Crown. The mages of the assembly are known for their extravagant lifestyles and intimidating history of political and economic control, and much of the populace mistrusts them, though no one dares to cross them or the king. There are eight ruling positions within the assembly, and the families of each member are appointed as lesser members without responsibility, called "annexes." Mages of the assembly wander throughout the empire with little fear of argument or resistance.

Members are elevated into one of the eight major offices on the death of a member that leaves an office vacant, or when a standing member abdicates their position. The remaining members then recommend replacements and vote on the mage to fill the newly open office.

\subsection{Goals}
The assembly plays the long game, working to subtly increase its power across the empire and beyond. Popular decisions are often publicly presented as the will of the assembly, while conflicted issues are tailored to seem purely the "will of the Crown," driving a wedge between the two bodies. The assembly seeks to write history to favor its members as the architects of modern law and culture, with the eventual goal of subverting the power of the king, turning him into yet another puppet of the assembly's goals. Assembly mages believe that only they have the knowledge and talents to truly guide the empire into greatness.

Aside from political matters, the assembly is focused on locating and excavating powerful relics and magic lost to the Age of Arcanum. Its members believe that their original namesake was the apex of society across all history, and they seek to reclaim their place as paragons of mortal rule. Thus, all salvaged texts and artifacts of power from that era are to be recovered at all costs.

The curious discovery of the power behind the Kryn Dynasty has recently drawn much of the assembly's attention. The mysterious, primal force known as dunamis could change the world's understanding of magic itself—and could even be the key to placing the assembly above all other rulers and nations. In pursuit of this goal, agents of the assembly stole two Luxon beacons from the Kryn Dynasty, though they subsequently lost one to a band of adventurers called the Mighty Nein. The assembly has dedicated immense effort to studying and exploiting this new form of magic, knowing full well that stealing the beacons is what caused the current war in the first place.

\subsection{Relationships}
The assembly has worked directly with the king and his chamber for so long that the two factions have become deeply intertwined, frustrating both of them. Though they rely heavily on each other, every smile and agreement barely conceals an undercurrent of venomous loathing. Some members of the assembly feel that it is far too soon to overtly challenge the Crown, while others believe it would be foolish to miss the opportunity the war is creating.

The assembly has established itself as the more forward-thinking face of the empire to the outside world, with Lord Athesias Uludan going out of his way to upstage the Crown at every diplomatic opportunity, all "in the name of the king." Most foreign officials happily deal with the assembly instead of the Crown, especially those of the empire's closest trade partner, the Clovis Concord. The assembly ensures that such meetings are recounted to the king with embellished tales of their diplomatic success.

The acquisition of a Luxon beacon and the hunger among assembly mages to unlock its powers have drawn the animosity of the Kryn Dynasty, triggering their attacks on the empire.

\subsection{Figures of Interest}
The mysterious minds of the Cerberus Assembly often quarrel over goals and methods, but ultimately these eight mages are united by their heritage and a shared thirst for knowledge.

\subsubsection{Master Trent Ikithon, Archmage of Civic Influence}
\textit{Chaotic evil, male human}

\noindent
Trent is respected as the acclaimed Propagandist of the empire and the third oldest member of the assembly. Once an instructor at the Soltryce Academy, he only returns every few years to collect young students for his experiments in the mental conditioning that he calls "awakening." Many of these students go mad and are locked away, but those who endure become zealots for the assembly and join the Volstrucker, an elite group of arcane thugs commonly known as Scourgers, who perform the assembly's dirtiest work under Trent's direction.

\subsubsection{Martinet Ludinus Da'leth, Archmage of Domestic Protections}
\textit{Lawful evil, male elf}

\noindent
Ludinus is the oldest and only original member of the assembly, as well as the master of warfare and conflict. Charged with overhauling the military structure of the Dwendalian Empire, Ludinus directed the construction of the garrisons on the Xhorhasian border and often oversees their maintenance. He was one of the mages who survived the destruction of Molaesmyr and fled to Bysaes Tyl, but he saw the opportunity to achieve greatness within the empire and left his culture behind to continue his arcane pursuits. Wise, if emotionless, he bears a deep hatred for the Kryn Dynasty and spares no effort gathering information on their weaknesses and secrets. Ludinus spends most of his time developing arcane weapons of war and shoring up the military might of the empire, while subtly challenging the leadership of Crown Marshal Damurag.

\subsubsection{Baroness Jenna Iresor, Archmage of Industry}
\textit{Neutral doppelganger}

\noindent
One of the younger members of the assembly, Jenna is known for her business acumen and her extravagant lifestyle. By hiding her nature as a doppelganger and using memory-altering magics at a young age to fabricate a false past, Jenna constructed her human persona from the ground up, leveraging her powers of deception to essentially write herself into history as a Clovis Concord expatriate. She helps oversee central guild business in Rexxentrum under Guildmaster Kai Arness, and helps Exchequer Aethia Drooze organize the collection of tithes through starostas across the empire.

\subsubsection{Master Doolan Tversky, Archmage of Dysology}
\textit{Chaotic neutral, female gnome}

\noindent
The second-oldest member of the assembly, Doolan is in charge of the study and understanding of abnormal creatures and deviants of arcane creation that might threaten the empire's way of life. She is an absentminded yet brilliant gnome who is obsessed with all beasts, aberrations, and creatures of legend. Doolan imports creatures from around the world to study, disassemble, and use in her attempts to revolutionize magical practices. She resents the Library of the Cobalt Soul, as her reputation has caused them to bar her from their facilities. She wishes to catalog the unstudied horrors of Xhorhas and has covertly obtained the services of the Myriad to retrieve new specimens.

\subsubsection{Lord Athesias Uludan, Archmage of Diplomatic Union}
\textit{Neutral good, male human}

\noindent
Athesias's charm and bombastic personality serve him well as a diplomat. His duty is to foster a positive relationship with people of power both within and beyond the borders of the empire. He was originally one of the most effective instructors at the Soltryce Academy, but his penchant for spectacle and his rampant narcissism made him a difficult ally to trust with state secrets. When the office of Diplomatic Union opened, he was quickly and quietly reassigned. Athesias finds great pleasure in ruining or usurping the plans of his counterpart in the Crown's employ, Emissary Lord Zeddan Graf.

\subsubsection{Lady Vess de Rogna, Archmage of Antiquity}
\textit{Neutral evil, female half-elf}

\noindent
A public recluse for most of her life, Vess is both a brilliant mage and dedicated historian. She assumed this post after replacing her criminal predecessor, Lady Delilah Briarwood. As an instructor at the Soltryce Academy for over two decades, Vess has studied and unraveled a number of historical mysteries and pre-Calamity riddles—and hoarded some of the spoils for herself. Always eager to pursue forgotten lore and artifacts of eons past, Vess has been known to quietly vanish to Xhorhas for weeks at a time, returning with fewer guards and more uncovered secrets.

\subsubsection{Headmaster Zivan Margolin, Archmage of Conscription}
\textit{Lawful neutral, male human}

\noindent
Zivan Margolin inherited the position of headmaster from his father, the late Jorma Margolin. Zivan has been the headmaster of the Soltryce Academy in Rexxentrum for nearly twenty years. Calm, patient, and quietly imposing, Zivan walks the halls of the Academy with a keen eye for talent. He is in charge of the curriculum and also watches for any latent powers that may be worth grooming as future allies of the assembly, dangers to be monitored, or prospective minds for Ikithon to conscribe into the Volstrucker. Zivan has rarely had the opportunity to demonstrate his full power, for he is typically busied with keeping the peace between the feuding members of the assembly. Those who have witnessed his true might, however, now know that his words are backed by some of the most powerful magics within the Cerberus Assembly.

\subsubsection{Headmaster Oremid Hass, Archmage of Cultivation}
\textit{Lawful neutral, male earth genasi}

\noindent
The current headmaster of the Hall of Erudition in Zadash, Oremid is tasked with watching and grooming the next generation of mages and arcane specialists outside Rexxentrum. While he himself is a gentle soul who adores animals, he puts on the façade of a strict man with no sense of humor, which is further enhanced by the elemental influence of his earth genasi blood. He teaches students that failure is not an option, and that emotion is a barrier to one's true ability. Equally feared, respected, and privately loathed by the students (and some instructors), Oremid personally dismisses those who break under his school's curriculum and heaps joyous praise on those who endure their training.

\section{Clovis Concord}
Nearly four hundred years ago, an alliance of the Ki'Nau islanders and foreign Marquesian traders founded a new nation on the tropical Menagerie Coast. The Clovis Concord is a democratic nation that enforces law and order, and regulates commerce along the length of the coast. The nation is composed of eight independent city-states, each ruled by its own marquis, which operate as one union under the banner of the Clovis Concord. This arrangement ensures uniform laws, regulation of trade, and mutual protection between the city-states. These eight cities are Port Damali, Port Zoon, Gwardan, Tussoa, Othe, Feolinn, Nicodranas, and Brokenbank.

Open shipping lanes and inviting tropical scenery have made the domain of the Clovis Concord a cultural melting pot and a popular place to live or visit. Knowing this, the respective marquises uphold a long tradition of welcoming all weary travelers and providing an atmosphere of color and delight. If there is gold to be spent, the concord is eager for you to spend it within their cities. Between the venues for entertainment and vice, and the ample business opportunities, the concord hopes to convince those with wealth to spend it here, outside the gates of the Dwendalian Empire.

Internally, each marquis of the Clovis Concord keeps a tight grip on the shipping and trade that comes though their cities, overseeing all guilds that organize imports and exports and managing a handful of guildmasters that form the backbone of local commerce. When a marquis dies or is impeached by the other members of the concord, the remaining members choose a successor.

Warriors who prove their might through gladiatorial bouts during times of celebration and ceremony are asked to join the Zhelezo, a well-paid guard force that works for the governing marquis to enforce the laws. Cities have one or more magistrates that oversee judgment on legal matters and criminal punishment.

While the cities do work together for the good of the concord, they are also each in silent competition with each other, flaunting their affluence and clientele at every opportunity. When a prominent figure from the far reaches of Exandria finds their way to the Menagerie Coast, it's not uncommon for them to be courted by multiple marquises in an attempt to convince the personage to stay as a guest of their city.

In recent times, the Myriad has subtly infiltrated the Clovis Concord, often with the unwitting aid of greedy local politicians. Rumors of criminal collusion have raised suspicion amid members of the concord. Sooner or later, the tension will have to break—and the fallout won't be pretty.

Despite the concord's cordial relationship with the Dwendalian Empire, their proximity to the imperial power is a lingering source of unease. Currently, the concord maintains a standing military of 7,800 Shore Warden soldiers across the coast, a nautical fleet of 370 ships, and 5,200 Zhelezo divided between the cities. Some of the Shore Wardens are posted inland, scattered around and within the Cyrios Mountains, where they maintain forts along the border with Western Wynandir and the Dwendalian Empire.

\subsection{Goals}
The Clovis Concord is concerned with maintaining the trade and commerce that ensures their continued influence, both in Wildemount and the rest of Exandria. Anything that might jeopardize business is viewed as a major threat to the sanctity of life along the Menagerie Coast. Faced with the hazards of dangerous creatures that lurk in the nearby jungles and beneath the waves, as well as piracy on the open waters, the concord is employing substantial numbers of adventurers and mercenaries.

Forty years ago, when the concord worked to enforce rising taxation on independent merchant companies operating within the Menagerie Coast, the outrage led a small union of angry enterprises to turn to piracy, calling themselves the Revelry. These pirates have claimed the island of Darktow and terrorized the shipping lanes for four decades, all while deftly avoiding pursuit and naval retribution. The concord is eager to destroy the Revelry in its entirety, and have put massive bounties on the heads of the so-called scourges of the ocean.

\subsection{Relationships}
After the fall of the Julous Dominion and the rise of the Dwendalian Empire, the concord was eager to reach a diplomatic arrangement with the victorious empire. The empire recognized the difficulty of conquering the Clovis Concord and the challenge of managing their immense network of international trade, and acquiesced to a treaty. This loose alliance has been mutually beneficial for many generations, and the concord city of Nicodranas acts as a common ground for people of both nations. Both imperial Crownsguard and local Zhelezo keep the peace in Nicodranas, and both imperial and concord laws hold sway within the city. While this arrangement occasionally causes strain between the two governments, the marquises of the concord are happy to have the buffer of the empire between them and the ravenous wastes of Xhorhas. As the war between the empire and Xhorhas continues to grow, the concord is determined to stay neutral. This has caused some tension with the empire, but nothing the concord worries will overshadow their long-standing treaty.

\subsection{Figures of Interest}
Though the marquises of the Clovis Concord outwardly embrace the spirit of cooperation on which their nation was founded, most scheme and conspire against one another to attain ever-greater power.

\subsubsection{Olesya Lapidus, Marquis of Port Damali}
\textit{Neutral good, female human}

\noindent
The Lapidus family's work in protecting and continuing the prosperity of the city for four generations has kept a Lapidus as Marquis of Port Damali for just as long. Olesya was raised as a boy, but quickly discovered that she identified as female and transitioned in her youth. As the fourth and current leader of Port Damali, Olesya takes pride in her charge and is respected by the rest of the concord. However, with the growing distraction of the Revelry, Olesya failed to notice the infestation of Myriad agents throughout her city and the guilds beneath her. Now realizing the dangerous web surrounding her and the interests of her city, she is seeking a way to turn the tide.

\subsubsection{Zhafe Uludan, Marquis of Nicodranas}
\textit{Lawful neutral, male human}

\noindent
Zhafe is a distant relative of Lord Athesias Uludan of the Cerberus Assembly; his placement within the concord was primarily a diplomatic concession to maintain a positive relationship with the empire. However, Zhafe shuns this notion, insisting that he earned his role on his own merits, and eagerly organizing small galas and parades to ensure his popularity among the citizens of Nicodranas. Even so, he secretly fears the assembly and bows to Athesias's requests whenever they meet, making him a useful pawn.

\subsubsection{Alamads Haddou, Marquis of Port Zoon}
\textit{Neutral good, male dwarf}

\noindent
The old, crotchety, and honorable leader of the industrial city of Port Zoon, and head of the Silverbone Smithing Guild that operates all along the coast, Alamads is widely known for his insurmountable will and strong sense of duty to his people. As a former Shore Warden and a legendary ironsmith, his skill with a hammer—both on and off the battlefield—has earned him great respect, as well as the whispered nickname of "Ironpuss."

\subsubsection{Vasan Atrith, Marquis of Gwardan}
\textit{Lawful good, female elf}

\noindent
Unflappable and driven to protect the people of her city, Vasan is a bit of an outsider among the other marquises of the Clovis Conchord for her standoffish nature and the way she unapologetically prioritizes the wellbeing of her city over that of the whole coast. Though Vasan presents herself as aloof and uninformed, much of her political persona is intended to sow underestimation of how informed and clever she may be.

\subsubsection{Portcaller Mazin "Fat Fish" Fahreed}
\textit{Neutral, male tabaxi}

\noindent
Tough as nails and quick to bite back, Mazin is an honorable tabaxi and dependable leader in times of trouble. This led to his rise as the Portcaller of the island village of Brokenbank, the first major trade junction along the Menagerie Coast. Though the village is small, Mazin deals with a perpetual string of smugglers and pirates while also organizing the incoming and outgoing shipping lanes for much of the coast. Even so, Mazin is rarely included in major concord business, to his obvious and public frustration.

\begin{DndSidebar}[float=!b]{Laws of the Clovis Concord}
All who are found by the esteemed Zhelezo to have broken the fair laws of the Clovis Concord are judged and sentenced by the local magistrate and are subject to punishments befitting the crime. All fines and punishments for crimes against a Zhelezo official or a noble are doubled. All fines and punishments for crimes against a Clovis Concord official are tripled.



Permits must be purchased and owned to do business for more than 15 days in Menagerie Coast cities. Permits are approved and purchased from the Compass Lodge installations.



All religions are legal to practice in private, but public worship of Betrayer Gods is considered dangerous to public discourse and is strictly forbidden.



\begin{description}
    
    \item[Theft:] 3–10 days incarceration and/or a fine of twice the worth stolen.
    
    \item[Public worship of Betrayer Gods:] 10 days incarceration and/or a fine of 300 gp.
    
    \item[Ownership of another intelligent humanoid creature:] 6–12 months incarceration and/or a fine of 1,000 gp.
    
    
    
    \item[Trespassing:] 7 days incarceration, and/or a fine of 100 gp.
    
    
    
    \item[Failure to pay taxes to the Clovis Concord:] 3–15 days incarceration and a repossession of property equal to the amount owed.
    
    
    
    \item[Damage to private property:] 5–15 days incarceration or a fine of twice the worth of property destroyed.
    
    
    
    \item[Damage to concord property:] 30–60 days incarceration and a fine of twice the worth of property destroyed.
    
    
    
    \item[Kidnapping:] 30 days incarceration and/or a fine of 500 gp.
    
    
    
    \item[Assault with intent to injure:] 15 days incarceration and/or a fine of 250 gp.
    
    
    
    \item[Assault with intent to kill:] 170 days incarceration and/or a fine of 1,000 gp.
    
    
    
    \item[Murder:] 5–10 years incarceration and/or a fine of 5,000 gp.
\end{description}



\end{DndSidebar}

\section{The Myriad}
Any civilized society will turn the people it has failed into criminals, and the societies of Wildemount are no exception. The continent's largest criminal organization took shape about eighty years ago, when a shipping company in the Dwendalian city of Yrrosa turned to smuggling contraband to make ends meet. This tightknit group of clever smugglers soon began to bargain with their competitors, employing blackmail and offering membership to their organization as an alternative to elimination.

The syndicate grew with alarming speed, infiltrating the criminal underbelly of every major city in Western Wynandir. Masquerading as purveyors of antiquities and foreign textiles, the Myriad focuses on providing their clients with exotic goods, such as illicit substances and magical beasts, or supplying hired muscle to intimidate their clients' rivals. The worst of the Myriad even deal in human merchandise.

Originally centered in Yrrosa, the Myriad now operates across the continent as a loose network of gang bosses who run their own local sects without direct oversight from the mysterious heads of the syndicate. Each satellite group is expected to regularly deliver information and a cut of their profits to the leadership. Those who fail to pay up receive quiet threats of enslavement or assassination—and the Myriad always makes good on its threats.

Members of the Myriad are sworn to keep their syndicate secret by pretending that their chapter is still just another local gang. This has given the Myriad a sinister and enigmatic reputation, which further obfuscates its activities from the authorities. Since members of the the Myriad prefer coercion to public displays of violence, officers of the law often fail to recognize Myriad activity until it's too late.

The empire scored a major victory fifteen years ago when they discovered and raided the central Myriad stronghold in Yrrosa, forcing the surviving leadership to scatter across the continent. The Myriad has emphasized the narrative of its own fall in order to keep a low profile, though in fact it has retained most of its power. Its influence has stretched to the distant corners of Wildemount as it bide its time and waits to reinstate its control within the empire. Shifting tactics, the Myriad has now begun seeding major factions with double agents, who turn the most foolish or corrupt of their colleagues into the Myriad's unwitting pawns.

\subsection{Goals}
The loss of their headquarters taught the surviving Myriad leaders to temper their arrogance with caution. Now a decentralized network of allies and informants, the Myriad is reconnecting with their temporarily dormant contacts, not all of whom are happy to see their cruel masters return.

Beyond the borders of the Dwendalian Empire, the Myriad has deftly utilized the chaos of the Revelry's decades-long reign of piracy to infiltrate the guilds of the Clovis Concord. Between the war in the east and the distractions on the western waters, there has never been a better time for the Myriad to strike.

\subsection{Relationships}
Merely mentioning Myriad activity within the Dwendalian Empire or the Menagerie Coast is enough to draw unwanted attention; members of the local law enforcement might investigate the rumors—and the person who started them. Those involved with the Myriad often vanish into the darkest prisons, or worse. Though the Myriad has apparently scattered since the raid on their base in Yrrosa, the authorities still keep an eye open for cell activity.

Twenty years ago, the Myriad established satellite cells in the distant lands of Tal'Dorei and Marquet. Since their headquarters in Wildemount was destroyed, some have suggested sending agents across the oceans to reinforce their foreign operatives. If the network can quickly regain its former strength, there will eventually be little chance of stopping the spread of its power and influence.

In the supposed absence of the Myriad, the Tribes of Shadycreek Run have begun to push into the empire to stake their claim in Myriad territory. The leaders of the Myriad watch from the shadows and plot ways to conscript (or, if necessary, to destroy) the prideful families of Shadycreek Run.

\subsection{Figures of Interest}
The dangerous aspirations of the Myriad are driven by a curiously scattered collective of intelligent liars.

\subsubsection{Iris Bethelas}
\textit{Lawful evil, female human}

\noindent
This lovely, elderly woman is deeply respected and beloved for her wise, spiritual guidance and seemingly gentle demeanor. This persona kept her safe in Yrrosa, and now in Rexxentrum, where she is one of the masterminds planning the return of the Myriad. Insightful, clever, and unhesitating in her cruelty, Iris has long been considered the driving force behind the Myriad's success.

\subsubsection{Doctor Waldorf Perifeather}
\textit{Lawful evil, male halfling}

\noindent
A talented physician and healer, Waldorf often travels between Port Damali and Zadash, tending to affluent patients with rare diseases. His proximity to powerful figures in their most vulnerable moments grants the good doctor a source of highly sensitive information and a reputation as a savior. This combination of influence and access makes him a key information broker for the Myriad.

\subsubsection{The Gentleman}
\textit{Neutral evil, male water genasi}

\noindent
The Gentleman has held a position of influence in the shadows of Zadash for nearly two decades, using his charm and connections throughout the city to carve a sizeable niche for himself in the Underworks beneath the city. These days, he's respected as a suave crime lord with dozens of reliable cutthroats and thieves at his beck and call.

\subsubsection{Korshad Seddiki}
\textit{Neutral evil, male human}

\noindent
A patient and intimidating figure in the underground of Deastok for over forty years, Korshad is often known as the "debtkeeper." He sends agents and missives to individuals throughout the empire who fall behind on their payments to the Myriad. Korshad is a mysterious figure by his own design, and he shrouds himself in illusion magic for his own protection and prestige, feeding the whispered myths about him that are legendary in unscrupulous circles.

\subsubsection{Chessia Wakiam}
\textit{Lawful neutral, female elf}

\noindent
A talented songstress and dancer renowned throughout the Menagerie Coast for her sweet voice and enchanting allure, Chessia is a popular performer in many opulent social circles across Port Damali. Chessia uses her charming persona to deflect suspicion as she weaves her way into the lives of politicians and guild masters to steer them toward Myriad interests.

\section{Children of Malice}
In the echoing stillness that followed the Divergence, the ruined lands of the surface retained remnants of civilizations trying to rebuild their world from the ashes. Below, however, the surviving children of Lolth could feel her presence fade with her banishment. The chains she held over her drow were weakened, and those who had found renewed faith and purpose in the discovery of the Luxon abandoned their terrible mistress and slowly climbed their way to the surface to reclaim Ghor Dranas and be rid of their spiteful god for eternity.

Still reeling from her banishment to the Abyss, Lolth felt a rapid shift in her influence as these drow turned from her. Her fury burned with a poisonous hate, and she seethed in her webs as she plotted vengeance on her ungrateful and treasonous progeny. As time passed and her dark divinity gradually recovered, the Spider Queen set aside her pride in favor of her wish to punish; she began seducing the minds of goblinoids, beastfolk, and the wayward people of the wastes, creeping into their dreams and convincing them to love her, bending them to join in her violent hate for the rising Kryn Dynasty. She found drow who felt spurned by the Dens of the dynasty and preyed on their fears and insecurities, amplifying their xenophobia as the dynasty recruited more and more outsiders to their cause. Gradually, Lolth developed a new network of scattered worshipers, spies, and assassins. She calls them her Children of Malice, and they will gladly serve as the instruments of her revenge.

The Children of Malice aren't numerous, but they are clever, driven, and directed by a god of pure hatred. They established hidden cabals throughout the fringes of Xhorhasian society, forming motley bands of Lolthite bugbears and hobgoblins who raid outlying Kryn settlements. Drow agents of Lolth have infiltrated even the highest echelons of Kryn society, slowly poisoning the dynasty both metaphorically and literally. The true extent of this infiltration is only now being uncovered within Rosohna, spreading fear throughout the dynasty as the war against the Dwendalian Empire continues to divide the attention of the Bright Queen between the enemy outside and the enemy within.

\subsection{Goals}
The primary goal of the Children of Malice is to subjugate and dismantle the flourishing, Luxon-worshiping dark elf society in Xhorhas. They wish to raze everything the Kryn have built, crush their will, ruin the holy implements of the Luxon, and drag these broken, traitorous drow back under Lolth's influence and command.

To achieve this goal, Lolth and her children know well the necessary tenets of secrecy, lies, and patience. Most of their spies are sleeper agents who wait in the shadows, quietly gaining converts while worming their way into the heart of Rosohna and awaiting the order to strike.

\subsection{Relationships}
The Children of Malice have no formal alliances; they prefer to infiltrate other factions as part of their design to undermine the dynasty. Keeping their true nature hidden, members of the Children trade with outlying camps and townships across the Wastes for supplies and information. Thus, their informants are usually unaware they are informants as they trade gossip for goods and conversation, unwittingly feeding the enemy.

\subsection{Figures of Interest}
The Children of Malice vary dramatically in both appearance and methods.

\subsubsection{Zyn Daev'yana}
\textit{Lawful evil, male dark elf}

\noindent
Zyn was born in squalor in the Coronas of Rosohna. Lolth found him at his nadir and showed him how to refine himself into a true liar. He's since risen to become one of the more even-tempered members of Den Daev'yana, climbing the ranks as a respected strategist and writer. Using his position to sew discord and doubt amongst the Dens, he privately communes with the Spider Queen as one of her most trusted operatives.

\subsubsection{Ereldra Icozrin}
\textit{Neutral evil, gender-fluid dark elf}

\noindent
A talented harpist lauded throughout Rosohna for their musical gifts, Ereldra worshiped the Luxon before they were betrayed by their own brother, who sabotaged their marriage to their wife. Lolth found Ereldra in their despair and molded them into a valuable spy. Ereldra weakens in the social foundation of the Dens by spreading rumors and telling lies.

\subsubsection{Ironeye Jaglord Brukk}
\textit{Lawful evil, male hobgoblin}

\noindent
A feared warlord who guides a nomadic band of Lolthite hobgoblin warriors, Jaglord honed the skills of his marauders in the broken lands of the Barbed Fields, where they brazenly attacked Kryn trade caravans and supply wagons before vanishing into the deadly fields. As his war band grew larger and more capable, he boldly led them farther south, eventually coming to rest within the Ghostlands outside Rosohna, where they now patiently await further orders from the Spider Queen.

\subsubsection{Head Matron Vivurk Tonn}
\textit{Neutral evil, female dark elf}

\noindent
Vivurk commands the drow houses and the hobgoblin horde that claim the mountainside stronghold of Dumaran. She rules with a bloody, iron fist in the image of her god to impose the same love and fear that characterized the drow societies of old. Ruthless and demanding, Head Matron Vivurk builds, trains, and prepares her forces to one day help Lolth strike down the Kryn Dynasty with delicious brutality.

\section{Diarchy of Uthodurn}
When the elven city of Molaesmyr unearthed a terrible curse that ruined the Savalirwood, many survivors fled to the north. They endured the blizzards and dangers of the Flotket Alps, clinging to the hope that the withdrawn dwarves of Uthodurn would find enough mercy to aid them. Though many perished in the harsh elements, the elven refugees finally located the iron doors of the subterranean stronghold and requested asylum. It was with only mild hesitation that the dwarves accepted their desperate guests, taking them into the warm halls of their city and giving them a place to stay.

Many years passed, with little information on the cause of the cataclysm that shattered the elven homeland, and everyone gradually realized that their new situation was permanent. Social tensions rose as cultural differences and forced proximity placed the denizens of Uthodurn into rising discomfort. This growing divide worried Judessa Fruunast, the queen of the dwarves of Uthodurn. She met with the throneless elven king, Imathan Talviel, and unexpectedly offered to instate a diarchy within the city. Months of deliberation followed as the monarchs discussed how best to establish social harmony, where new homes would be built, and how to expand the discs of the city to fit a more permanent populace. Eventually, the two came to an agreement that the majority of both dwarves and elves felt was fair and tolerable.

The process of restructuring the stubborn societies of this new Uthodurn forged a long partnership and deep friendship between the two leaders until Queen Judessa's death in 794 PD. Judessa's daughter, Simone Fruunast, took up her mother's mantle, and the bonds of friendship forged in her mother's time flourish under her rule. The benefits that were uncovered through the unity of these disparate cultures birthed new prosperity throughout the city, and the new dual society has grown into a synergistic civilization over the centuries.

Now spreading beyond the underground halls of Uthodurn, the people under the Dual Monarchs seek to build a brighter future, and perhaps even resolve the mysteries of their past.

\subsection{Goals}
Many historians and prominent members of both halves of Uthodurn's society are investigating the source of the corruption that persists in the Savalirwood, hoping to one day heal the forest and reclaim the homeland of the elves. Shreds of evidence have been recovered by scouting parties sent into the deadly thicket surrounding the ruins of Molaesmyr to uncover signs of the curse and its origins. However, looters and thieves from Shadycreek Run have discovered the ruins as well, stealing relics and selling them to private interests. These thefts, coupled with recent research indicating that the corruption is slowly spreading, have accelerated the need to find answers to the mystery of the Savalirwood.

The growing population of the city has a constant need for more resources, so the diarchy agreed on the need to expand beyond the mountain pass. The village of Uraliss was founded using experimental magic designed to thaw snowfields, with the goal of developing sections of the Rime Plains into farmland. Unfortunately, the ground beneath the snow turned out to be untillable rock, so the settlement was abandoned. The dual monarchs then established the northern fishing village of Palebank, which supplies new goods and a coastal outpost in spite of the grueling cold.

Small vessels built to withstand the icy waters of the Frigid Depths began to explore the shores surrounding the region, discovering the lands of Eiselcross. In these frozen wastes, the settlers of Uthodurn uncovered signs of scattered arcana from a pre-Calamity era buried beneath the snow. These discoveries have piqued the interest of Molaesmyr researchers and Uthodurn's monarchs alike. As the conflict in the south escalates, the leaders of Uthodurn are eager to gain whatever power or protection might be gained from the secrets hidden in the treacherous tundra of Eiselcross.

\subsection{Relationships}
Uthodurn has little affiliation with other factions across Wildemount. Aside from the occasional explorers and adventurers who pose no threat, most people who come to Uthodurn are met with skepticism and avoided. The locals want nothing to do with the war in the south.

Anyone who might be from Shadycreek Run is treated as a lowlife brigand and immediately captured and interrogated. Any information about the powers behind the criminal civilization is considered useful, and once non-violent prisoners expend their usefulness, they are unceremoniously stripped of their equipment and set loose in the snowy fields of the Rime Plains.

\subsection{Figures of Interest}
Distanced from the political bustle of Western and Eastern Wynandir, these figures within the Uthodurn diarchy are beginning to seek their place in the wider world of Exandria.

\subsubsection{Queen Simone Fruunast}
\textit{Lawful good, female dwarf}

\noindent
Daughter of the previous queen, Judessa Fruunast, and the current ruler of the dwarven people of Uthodurn, Simone is a respected, charismatic leader known for her endless compassion and warrior spirit. Often driven more by her heart than her wits, she has come to rely on Imathan for guidance and clarity. While a diplomat by nature, Simone has certainly proven her mettle with an axe in times of conflict.

\subsubsection{King Imathan Talviel}
\textit{Lawful good, male elf}

\noindent
The elder king of the surviving elves of Molaesmyr, Imathan is as candid as he is graceful in his old age. Patient but stern, he guides the elven people of Uthodurn with a calm strength and wisdom that helps maintain the continuing peace between the dual communities of Uthodurn. Imathan was a priest of Corellon in his younger days; now he maintains the faith among his people and shares the light of the Arch Heart with the dwarves of Uthodurn.

\subsubsection{Professor Gulrim Shalebrow}
\textit{Lawful good, male dwarf}

\noindent
Gulrim is praised as the brazen—though perhaps crazy—explorer who convinced Queen Fruunast to allow his team to cross the Frigid Depths. He and his surveyors found the mysterious banks of Eiselcross, established the Syrinlya outpost, and recovered fragments of ancient relics previously lost in the ice and snow. Professor Gulrim now helms the Uthodurnian efforts to expand Syrinlya and learn more about the mysteries hidden among the glacial fields.

\section{Tribes of Shadycreek Run}
The collection of rival criminal families known as the Tribes ran Shadycreek Run for nearly a century before losing much of their dominance to violent infighting and internal power struggles. Comprised of the Mardoon family, the Uttolot family, the Trebain family, and the Jagentoth family, the Tribes still control most of the business that transpires in the city as they work against each other in secret. While each family center is usually comprised of relatives by blood or marriage, their employees are also considered to be "within the family," so long as they remain loyal. The Tribes operate without a unified symbol or banner, and their interests lie in fencing stolen goods, importing and providing vice, and controlling the populace through confusion, intimidation, and fear.

The Mardoon family is a powerhouse of crime within the Graying Wildlands, thanks to its vast network of foreign contacts as well as its prosperous alliance with the Myriad. Known for relentless business acumen and a love of enforceable contracts, the Mardoons would rather shake hands than swing blades. Even so, when contracts are broken, examples must be made, and such circumstances demand a bloody show. The family fortune is made from brothels, jewels, and the slave trade, often specializing in exotic finds from Xhorhas and beyond. Leadership is shared between Lord Anselm Mardoon, Kriemhilde Mardoon, Ophelia Mardoon, Demis Mardoon, and Rufus Mardoon.

The Uttolot family rose to prominence through a network of familial trade alliances, making it the foremost importer of goods difficult to find this far north of the empire. It controls the gates through the Quannah Breach, and anyone who passes into or out of the region must pay tolls to the Uttolots, making them a necessary partner to the other families whether they like it or not. They dabble in poaching dangerous beasts and monsters, and often clash with the hunters of the Greytrader Union, a collective of citizens who want to improve life for the people of Shadycreek Run. Unlike the other families, the Uttolots rely on the Grudge Gang of Shadycreek for most of their muscle. Leadership is shared between Lord Milos Uttolot, Mira Uttolot, and Jagoda Uttolot.

The Trebain family is the most covert of the houses, preferring to do its business through intermediaries and hired muscle. It lives beyond the boundaries of the city, deep within the Savalirwood, in a network of homesteads surrounding a stronghold. The Trebain family focuses on controlling hunting grounds and harvesting relics from the ruins of Molaesmyr. The family's interests put it in constant conflict with the Uttolot family and the Greytrader Union. The triangle of animosity between the Uttolots, Trebains, and Greytraders sometimes makes for strange alliances—after all, the enemy of my enemy is my friend, no matter how temporary that bond might be. Leadership of the family is shared between Lord Darko Trebain, Juraj Trebain, Tatia Trebain, and Zorka Trebain.

The Jagentoth family has earned a reputation for being especially brutal and unforgiving. Its members publicly execute perceived traitors and high-profile competitors, but promise great rewards in exchange for absolute loyalty. The family's most successful business ventures include assassination, cultivating and distributing illicit substances, and slavery (which puts them in direct competition with the Mardoon family). Leadership is shared between Lady Sulia Jagentoth, Reese Jagentoth, Numa Jagentoth, and Ionos Jagentoth.

\subsection{Goals}
Conflict between the families is a constant in Shadycreek Run, but the apparent absence of Myriad activity has given the Tribes a unique opportunity to expand their horizons—and their influence. Lately, the Tribes have occasionally endured each other long enough to embark on joint ventures into the empire to rekindle old alliances with the Myriad.

Some see this age of tenuous cooperation as a chance to escape from the long shadow of the Myriad and stand alongside them as equals, but others fear that forging new bonds with the Myriad will just give it the freedom to destroy the Tribes from within.

\subsection{Relationships}
The Tribes endure each other in order to maintain a level of civility within the otherwise cutthroat streets of Shadycreek Run, but they are each on the lookout for ways to subvert or weaken their opposition and become the city's dominant family.

All members of the Tribes are documented criminals within the empire, with bounties to match. Travel within the borders of Western Wynandir is dangerous for these criminal overlords, but each family has their own connections within the major cities of the empire that enable business.

In the past, the Myriad has oft attempted to subjugate the Tribes, but it has been rebuked time after time. The Tribes are too proud of their well-oiled engine of vice to permit the Myriad to take over. Even though the Tribes are pursuing new alliances with the Myriad, the people of Shadycreek Run fear that conflict between these shady powers is inevitable.

\subsection{Figures of Interest}
Life in Shadycreek is perpetually dangerous for the members of the Tribes, where every smile hides ill intent.

\subsubsection{Ophelia Mardoon}
\textit{Neutral evil, female tiefling}

\noindent
One of the most prominent members of the Mardoon family, Ophelia commands the Estate Sybaritic to the north of Shadycreek Run. Hedonistic and forceful, she does not enjoy being denied and has the skills and connections to ensure that she gets whatever she desires. One of the more ambitious members of the house, Ophelia has made a number of contacts in the Wildemount criminal world outside the Greying Wildlands.

\subsubsection{Jagoda Uttolot}
\textit{Chaotic evil, male half-orc}

\noindent
Jagoda is a master hunter and expert on the terrifying beasts that wander through the Savalirwood, using his skills to track and capture rare creatures for his family to sell. Gleeful in his brutality and prone to ignore the graces of diplomatic relations, Jagoda has occasionally gotten himself into trouble that the rest of his family has had to dig him out of.

\subsubsection{Lord Darko Trebain}
\textit{Lawful neutral, male human}

\noindent
The elder of the Trebain family, Darko helms the efforts to excavate the ruins of Molaesmyr from within his museum-like mansion. He is obsessed with the secrets of the blight on the Savalirwood, and he's convinced that it was caused by the elves of Molaesmyr tampering with misunderstood magics. He hopes to discover the source of the corruption and learn how to weaponize it. With such power at his disposal, the rise of the Trebains would be unstoppable.

\subsubsection{Ionos Jagentoth}
\textit{Lawful evil, male half-elf}

\noindent
Known for his skills in negotiation and deceit, Ionos has proven himself the master of the family's distribution network. Inspired to expand the business alliances of the Jagentoth family far beyond the northern edges of the Dwendalian Empire, Ionos has been making arrangements with Myriad operatives while simultaneously rooting out Myriad influence in those same circles. Either the noose is about to tighten, or he's designed a brilliant ploy to eliminate the competition.

\section{Library of the Cobalt Soul}
Under the enlightening scriptures of Ioun, the Knowing Mentor, teachers, priests, and monks who have been drawn to the calling of truth and knowledge spend their lives training within the Library of the Cobalt Soul. Rather than a single physical building, the Library of the Cobalt Soul is a collective term for the universal knowledge and philosophies upheld by those who follow Ioun's teachings. The Cobalt Soul is guided by a central belief that true strength is found in understanding the world around you. Despite their idealism, the pursuit of truth is hampered by the realities of life in the empire: politics, propaganda, and the dangers of the wilderness.

The Cobalt Soul is based in Rexxentrum but operates throughout the empire and across the whole of Exandria. Temples to Ioun under the management of the Cobalt Soul act as massive libraries called archives, usually located in larger cities and cultural centers. Its members come from all walks of life and are expected to assist in the maintenance, organization, and protection of the archives. Its satellite archives collect artifacts and research information both historical and contemporary along the Menagerie Coast and in the distant lands of Tal'Dorei, Issylra, and beyond.

Archivists act as administrators at each Archive, delineating tasks, overseeing the training of new members, and even negotiating for or purchasing artifacts and records—often from those unaware of their true value.

The monks of the Cobalt Soul are the enlightened knowledge-seekers of their order. They research places where ancient knowledge could be hidden and lead large-scale expeditions to these places.

Expositors are the covert agents of the Cobalt Soul. These enlightened infiltrators extract information that others would keep secret and use their newfound knowledge to better the world. As masters of acquisitions, expositors answer only their High Curator.

At the top of each Archive is a High Curator, who dictates the Archive's goals to subordinates and uses this power to assign marks to expositors, outline allies and enemies to their archivists, and approve the expeditions proposed by monks.

The Cobalt Soul was founded within the Julous Dominion before the nation fell to the armies of the Dwendalian Empire. The Dwendalian Crown permitted the Cobalt Soul to be assimilated into the social hierarchy of the empire rather than destroying it. This was not an act of mercy, but a political maneuver; the Cobalt Soul was allowed to survive only to provide a sense of normalcy and appease the restless populace.

The Library has continued to work with and under the empire out of necessity, but its archivists are savvy political players—whenever possible, they avoid making moves that could put them in danger of being exploited by the Crown or the Cerberus Assembly. Nowhere is this tension more apparent than in the very existence of the militant truth-seekers of the Cobalt Soul: the expositors, rigorously trained to root out corruption and falsehoods. Behind the scenes, a silent war is brewing as expositors bring the dark secrets of the empire's social elite out into the light.

\subsection{Goals}
Untold knowledge was lost when the Calamity brought the Age of Arcanum to its fiery end, and so too was the case when Ioun was mortally wounded in that terrible war. Since that day, the Cobalt Soul and their progenitors have sought to recover what was lost. They hope that, in time, their tireless effort will eventually usher all societies to an enlightened future.

Ioun is still wounded, but her wisdom and grace has returned to the world—at least in part. Under the guidance and wisdom of their patron god, members of the Cobalt Soul seek to enlighten themselves through study and research. They devote their lives to studying the mysteries of the world, and turning that knowledge to protect it from another Calamity. To this end, the Library makes the immense breadth of knowledge within their archives available to anyone and everyone who seeks to educate themselves—though some secrets are so potent that they must remain hidden from the public until the proper time.

Empires thrive on misinformation and propaganda, yet it is necessary for the Cobalt Soul to operate within the Dwendalian Empire. As such, the Cobalt Soul works to provide the truth to those who will listen, to covertly update revisionist history, and to remove false information before it can spread too far. Unfortunately, imperial agents and even spies from foreign powers have successfully infiltrated the Library over the centuries, and some of their attempts to rewrite fact and history in favor of their own interests have succeeded. In response, the Library has had to construct a powerful, hidden force of expositors to defend their texts from manipulation and uncover truths that remain guarded.

\subsection{Relationships}
The Library tries to remain autonomous, but their agents within the empire are under constant siege by Theolocrat Kon Bruda, a once-respected monk of the Cobalt Soul who is now just another corrupt lackey of the Crown. Bruda and the archivists wage a quiet battle of wills and subterfuge, each trying to outsmart and undermine the other. This situation, combined with the empire's efforts to prevent the education of their people, has created palpable unrest within the Library.

The perceptive ears of the expositors often hear rumors of unwholesome practices behind closed doors within the Soltryce Academy and other places under the Cerberus Assembly's oversight. The Library is concerned about the assembly's secretive cabal but has thus far been unable to infiltrate deep enough to uncover anything damning.

Meanwhile, a small number of volunteers from within the Library have agreed to travel deep into Xhorhas in the hopes of communicating with the Kryn Dynasty to learn more about their mysterious society, and to understand the level of threat they pose to the people of Western Wynandir and beyond. None of the volunteers have yet returned.

\subsection{Figures of Interest}
Operating under the will of Ioun, and driven by curiosity and the desire to further universal understanding, the diverse agents of the Cobalt Soul seek to change the world for the better.

\subsubsection{High Curator Loman Turray, Herald of the Zadash Archive}
\textit{Lawful neutral, male human}

\noindent
Curator Turray has held the position of High Curator of the Valley Archive in Zadash and the title of Herald of the Archive under the Crown for over forty years. He faces increasing pressure to instill his subordinates with loyalty to the Crown, and he is slowly giving in to imperial indoctrination. He fears that enforcing the will of the Crown when it contradicts the philosophies of the Library will cause his own colleagues to depose him—which only leads him further into the Crown's protective embrace.

\subsubsection{High Curator Yudala Fon, Herald of the Rexxentrum Archive}
\textit{Lawful good, agender half-elf}

\noindent
Yudala is a wise and respected figure who has spent most of their life fighting to protect and maintain the Rexxentrum Archive. They have mastered the art of parrying the Crown's and the assembly's attempts to subjugate the Archive and bring it under state control. For nearly two decades, they have been engaged in a battle of wills with the infuriated Theolocrat Kon Bruda—and they fear that Kon's patience is at an end. His vengeance will be elaborate, exacting, and swift, or so Yudala fears. Their network of spies and expositors is frantically seeking to find and publicize evidence of Kon's unlawful dealings and end the feud before it escalates further.

\subsubsection{Expositor Sia Kresh}
\textit{Neutral, female halfling}

\noindent
As Yudala Fon's secret right-hand, Sia organizes the expositors of the Rexxentrum Archive. She is determined to root out the infiltrators who would seek to undermine the leadership of the Library and grab power for the Crown, and she is swift and brutal in her work. Those that know of her fear her, and many wonder if her methods are crossing a line.

\section{Scars of Scale and Tooth}
For hundreds of years, the small, cloistered dragonborn nation of Draconia reigned over the Dreemoth Ravine to the south of Xhorhas, removed from the troubles of the wastes. The erudite draconbloods of the ruling houses, marked by their intellect, jeweled tails, and ostentatious wealth, looked down from the floating chain of islands that formed their capitol while the droves of tailless slave dragonborn, known as ravenites, toiled away in the ravine below, gathering resources and serving their masters for generations. This structure collapsed around twenty years ago, when the Chroma Conclave, a union of ancient dragons, destroyed Draconia and left the city ruined in the ravine below.

When the Conclave was destroyed, the surviving dragonborn ravenites outnumbered their previous masters, who were bereft of their former privilege and power. In the uprising that followed, many draconbloods were executed for their crimes, and others fled to the empire for asylum. Thus freed from their shackles, the ravenites built a new society adjacent to the ruins of the old nation of Draconia.

The past two decades have seen the rise of a new society, known as the Scars of Scale and Tooth, and it is wracked by the growing pains inherent to the dawn of a new nation. The Scars' leadership is tenuous, as many seek to claim the prestige of leading their people, and disagreements have led to bloodshed between the newfound councilors. This has made the building of the new city of Xarzith Kitril slow and fraught. Some heroic figures are seeking to establish a better direction for the ravenites, while others fled to the empire to seek better fortunes away from the infighting in their homeland. The war between the Kryn and the empire has closed the border pass into Western Wynandir, so the ravenites must find a way to coexist and thrive before the strife of warfare engulfs them.

\subsection{Goals}
The conflicting goals of the Scars' leaders have placed their nascent settlement in peril. Every would-be ruler wants something different for their people and has sufficient support within their faction to cause the streets of Xarzith Kitril to run red with ravenite blood, if they so desired.

Some amid the Kitrilians are outraged over news that a massive number of ravenite and draconblood refugees have been relocated to a shantytown called Talonstadt. Not only is the treatment of their kin unfair and inhumane, but this makeshift settlement practically straddles the Dwendalian Empire's eastern border, and they fear that Talonstadt will soon become the first casualty of a Kryn invasion.

Meanwhile, the monstrosities of the Xhorhasian wastes can sense the power vacuum left in the wake of Draconia's fall. Xarzith Kitril has no formal military, just a loose militia of hunters and trappers. This leaderless force is stretched to its limits, and more hunters are lost every day to opportunistic monsters. Outside aid is urgently required, but the squabbling councilors of Xarzith Kitril can't agree on whose aid they should request.

\subsection{Relationships}
Rumors of the Dwendalian Empire mistreating dragonborn ravenites have bred a heavy distaste for King Dwendal and his nation. The empire has recently sent emissaries to request an alliance against the dynasty, but Kryn diplomats have also arrived in Xarzith Kitril, seeking friendship and offering to aid the dragonborn in their efforts to rebuild. The Scars have heatedly debated which power they should side with—or if they can afford to aid either.

Ships from the Clovis Concord have docked on Kitrilian shores in recent years, eager to resume trade with the dragonborn people. In particular, Nicodranas has welcomed the emerging nation, and a few emissaries from Marquis Zhafe Uludan currently reside within Xarzith Kitril to keep an eye on the developing city-state and to aid to whatever limited extent they can.

\subsection{Figures of Interest}
The former slaves of Draconia are eager to forge their own destiny in the world. Some have risen to the challenge with glee, while others have had leadership suddenly thrust on them, but all the bearers of this charge hope to create a new beginning for the dragonborn people.

\subsubsection{Druvis Koothalok}
\textit{Lawful good, male ravenite dragonborn}

\noindent
Once a trained combatant in draconblood gladiatorial games, Druvis gained notoriety when he refused to strike his foes dead. His skill and creed garnered attention, and he was quickly sold to a wealthy noble house atop the floating islands of Draconia. In the chaos after Draconia's fall, Druvis became widely respected for his courage and sense of morality, and has hesitantly become a champion for the downtrodden people of Xarzith Kitril.

\subsubsection{Vemosi Runek}
\textit{Chaotic neutral, female dragonborn}

\noindent
Vemosi is a zealous and militant priest of the Storm Lord, and her tempestuous fervor and blistering sermons in the aftermath of Draconia's fall made her a popular voice of the people, though her detractors call her little more than a demagogue. Vemosi's brutal sense of justice has earned her—and the Storm Lord—many followers, all of whom are determined to hunt down and exact retribution from the surviving draconbloods. Vemosi has become a powerful figure within the Scars, and now she seeks to cement her place as a leader of her people.

\subsubsection{Troka, the Scoundrel King}
\textit{Neutral good, male dragonborn}

\noindent
Druvis first noticed Troka for his silver tongue, and tasked him with infiltrating the empire and appraising the situation in Talonstadt. Frustrated by the meager living conditions and demeaning treatment of his people, Troka has retaliated by sowing chaos throughout the surrounding Dwendalian territories. Troka and his troupe of slippery thieves steal from the Crown, confuse the local military, and keep a watchful eye on the dragonborn refugees, making him a folk hero among their community.

\section{Claret Orders}
Centuries ago, the high cleric of the Julous Dominion made secret pacts with a devil of the Nine Hells, hoping to gain enough power to repel the aggressive Dwendalian Empire. However, he failed to read the fine print of the devilish contract—he succeeded in his mission, but his people paid a terrible cost. The empire was indeed repelled, but it was because the lands of the Marrow Valley were brimming with devils and undead. The common folks hid in their homes, and the valley grew unhallowed and gray. In this brief era of terror, a humble priest of the Raven Queen named Trence Orman prayed for a way to protect his flock, his people, and the lands they called home. The Matron of Death blessed him with the ancient secrets of blood magic. Trence trained his most trusted friends in the techniques of hemocraft, and they sacrificed a portion of their humanity in exchange for the power to defend the dominion. This marked the origin of the Claret Orders.

Using their forbidden skills, the Claret Orders managed to seal away the source of this dark deal, free the high cleric from his torturous form, and clear the countryside of the evil that pervaded it. Unfortunately, before the Claret Orders could be celebrated as saviors, the armies of the empire swept into the freshly purged Marrow Valley and crushed the Julous Dominion.

The newly installed Dwendalian lords blamed the Claret Orders for inviting fiends into the Marrow Valley. While some knew the truth and tried to clear their name, it quickly became clear the Claret Orders were no longer welcome in their own homeland. Trence and his troupe of blood-tithed warriors went into hiding. Even in the face of such opposition, they never stopped protecting their lands from the shadows, seeking and destroying evils that emerged from the scars of the Calamity.

The orders have spread to the outskirts of Western Wynandir and beyond. They maintain small keeps to train new generations of protectors who are willing to take the Hunter's Bane and embrace their inner darkness to combat the outer darkness in the world.

\subsection{Goals}
The creed of the Claret Orders asks its followers to commit their lives to hunting monsters that threaten the sanctity of life and joy. While the sacrifice is great, the reward is the continued existence of purity and good in the world. The orders strive to protect those who cannot protect themselves from the terrors that go unseen by their noble lords and ladies. Most take little credit for their services, beyond means to live and travel. Some have strayed, using their skills to amass a fortune, but the heads of the orders deliver swift justice unto those who exploit the people they are sworn to protect.

\subsection{Relationships}
The existence of the Claret Orders is largely rumor and myth in most places, as the orders have no alliance with or allegiance to any standing government or group. The empire is well aware of their presence and makes a point to ignore their heroic deeds while publicizing misdeeds of those few hunters that succumb to the corrupting influence of their powers.

Some splinter orders have traveled beyond the boundaries of the empire, seeking more welcoming folk and more unique challenges in lands across Exandria. Still, no matter where the guardians of the Claret Orders travel, they know they will never be free from the hatred of those that misunderstand them.

\subsection{Figures of Interest}
Weeding out the terrors that feed at the edges of civilization, the members of the orders work for little more than the satisfaction of a deadly job well done.

\subsubsection{Claret Director Jette Brashan}
\textit{Neutral, male half-orc}

\noindent
An elderly priest of the Raven Queen who's as revered as he is feared, Jette has a powerful will and a tactician's mind. Jette never makes direct eye contact, and it seems that his attention is scattered among a dozen points at any given time, yet his ability to read patterns and calculate the hidden nature of evil leaves his compatriots in awe. Absolutely dedicated to the eradication of fiendish influence on the mortal plane, Jette has sacrificed much in his pursuit—all who serve him suspect that, one day, he will gladly make the ultimate sacrifice.

\subsubsection{Elias de Corvo}
\textit{Neutral good, male human}

\noindent
A famed blood hunter across the orders, Elias has successfully hunted some of the vilest creatures in all the planes and lived to tell the tale, though his scars show the cost of his victories. In his later years, Elias is now focused on training and mentoring new members of the orders, trying to instill in them morals that anchor them to their better nature.

\subsubsection{Lilyana Cayd}
\textit{Chaotic good, female werewolf (half-elf)}

\noindent
A sarcastic spitfire who joined the orders to help gain control over her lycanthropy, Lilyana has become a masterful blood hunter in an impressively short period of time. She leads many hunts with her keen, lycanthropic senses, and is known for being more personable than most who undergo the Hunter's Bane. Lilyana is spoken of with equal parts irritation and awe within the circles of the orders.

\section{Golden Grin}
When the lands of Tal'Dorei were ruled by the tyrant Drassig, a hidden society of muses, storytellers, bards, and folk heroes assembled to oppose their despotic king. The Golden Grin has endured since those dark years, and it has secretly guided the just, inspired the common folk, and planted the seeds of discontent and rebellion amid tyranny ever since.

Each member, known as a grinner, pledges to uphold the ideal that every individual has the power to change the world. They seek out places where hope is faint and the soul grows dim without inspiration, lifting the spirits of the downtrodden and heartening those who feel forgotten. Little is known about the Golden Grin aside from myths and rumors, and its members prefer it this way. The grinners listen for whispers and information using fabricated personas they create when they join the faction.

The Golden Grin arrived in Wildemount by following the trade routes from Tal'Dorei to the Menagerie Coast, where its members discovered a land plagued by warring despots. A sect of the Golden Grin has spent the past generation spreading its roots throughout the Clovis Concord. Now firmly established, the Grin has set about undermining corruption within the concord and has begun working their way into the Dwendalian Empire.

\subsection{Goals}
The Golden Grin—or simply the Grin—exists to enlighten and inspire, keeping ears and eyes out for the murmurs of tyranny and cruelty and deposing those responsible through social upheaval, or occasionally direct intervention. The idea that each person is capable of great things permeates its mantras and beliefs, and the wandering storytellers of the Grin are known to pay handsomely for tales of local heroes and majestic deeds of courage to add to their repertoire of stories and songs. If the people's hearts grow cold enough to reject words and poems, the noblest of grinners will often rise to the occasion and get their hands dirty to display such an act of heroism themselves.

\subsection{Relationships}
The Grin claims no formal allies, since its members never reveal themselves to others except in the direst of circumstances. They uplift and bolster other groups that work with the people's best interests at heart. Some grinners occasionally find their way into government positions, granting them a platform to disseminate the ideals behind their cause. While some may not agree with their rather anarchistic views, grinners are generally regarded as idealistic paragons of individual virtue.

The Grin's mischief in Wildemount began relatively recently, so it is largely unknown and able to conduct its work unimpeded; however, those versed in historical lore might know of its deeds and read the signs of its meddling. The Golden Grin currently is seeking an alliance with the Library of the Cobalt Soul to root out the untrustworthy within the Cobalt Soul and align their interests to aid the people of Wildemount.

\subsection{Figures of Interest}
In a world where strife and struggle often define one's life, it takes a certain quality of hero to step up and answer the call of the Grin.

\subsubsection{Sir Linus Denwallop}
\textit{Neutral good, male gnome}

\noindent
A wandering minstrel who claims a sourceless knighthood, Sir Linus travels all across the Menagerie Coast with his lyre and bag of illusions to astound and excite the children and poor folk of the countryside. Sir Linus intends to make up for his years of villainy as a smuggler and thief by subtly undermining the same sects of the Myriad with whom he once did business.

\subsubsection{Shakäste (Hush)}
\textit{Chaotic good, male human}

\noindent
A warm, gentle soul with blind eyes and a heart overflowing with patience, Shakäste was quick to learn the harsh realities of the world around him. After finding faith in the storms that shake the coast where he grew up, he was gifted with a hummingbird companion that serves as his eyes. Armed with a talent for going unnoticed and powers granted by his willful love of freedom, Shakäste has operated under the name "Hush" since joining the Golden Grin, quietly seeking out those in need of a little guidance.

\subsubsection{Master Duasad Keef}
\textit{Chaotic good, male human}

\noindent
An affluent noble among the upper crust of the Clovis Concord, Duasad is an inheritor of the Keef family's ancient, generational wealth. Duasad was raised as a woman, but realized as a young adult that something wasn't right, and sought a mage to help him transition—a fairly simple task in the concord, at least for those with affluence.

These days, Duasad has become well known as a powerful empath who uses his psychic talents to help others deal with trauma or seek a more fulfilling path in life. He has become one of the most influential and effective members of the Golden Grin in Wildemount. However, the radical shifts in his clients' life choices have not gone unnoticed, and Duasad is aware that the watchful gaze of the Clovis Concord is now firmly fixed on him.

\section{The Revelry}
The Revelry's allies within the Clovis Concord will tell you that they turned to piracy because of the concord's excessive taxes and unjust laws. Perhaps it began that way, but over time, a movement that began as public dissent turned to lawful protest, which then became violent action. Eventually, a union of disenfranchised salvage ships and trade vessels fought their way into the Dragshallow Reef and violently overthrew the small sanctuary port of Darktow.

After establishing a lawless base of operations on the isle of Darktow, these pirates began stalking the well-known shipping routes of the Menagerie Coast to harry, intimidate, steal from, and destroy passing ships that flew under the concord's banner. Calling themselves The Revelry, these "free folk of the sea" have now spent over forty years terrorizing the Lucidian Ocean. With loose morals and a twisted code of honor, they take what they want from less protected targets while living a life of debauchery in the well-guarded haven of Darktow. The Revelry lacks a uniform symbol, instead changing the colors and symbols of its banner regularly to avoid counterfeiting and infiltration.

The Revelry is governed by a collective of captains, all of whom defer to the final word of the Plank King, a capricious monarch who never leaves Darktow. The original Plank King, Hunnis Breeah, ruled for twenty-five years before a goliath named Wyatt Maranoss challenged him for his increasingly lackadaisical attitude—and won. Now Breeah is dead, and Wyatt claims the title of Plank King.

\subsection{Goals}
While the Revelry's original goal was to scare the concord into revoking their soaring taxes, its charter has changed. The Clovis Concord refused to change its laws, and the pirates found a certain dark joy in theft and murder on the waves. The original members of the Revelry (and now their children) have no plans to return to their former lives as legitimate traders. They intend to maintain and grow their operations across the Lucidian Ocean, while continuing to defend and improve their fortress in Darktow, finding intricate ways of reinforcing the reefs with dangerous traps and impassable barriers.

The Revelry has since expanded their interests beyond mere piracy. Relics and secrets dating back to the Age of Arcanum and beyond are rumored to be scattered all across the coast and the Swavain Islands. Some scavengers have returned to Darktow with proof of such artifacts, and a new kind of aspiring treasure seeker has emerged within the Revelry: one who seeks to plunder from the dead rather than the living. The drive to outdo their fellow sailors has already placed several ships in serious danger.

\subsection{Relationships}
Each member of the Revelry known to the Clovis Concord is considered a major criminal, amassing an ever-growing bounty with each infraction they commit and every ship they raid. Those arrested are generally executed and strapped to the rocks outside Brokenbank as a warning.

Some of the more affluent Marquesian merchant guilds that have little investment in concord shipments use back-channels to pay off the Revelry and fly specific colors when crossing their territory to avoid being boarded.

\subsection{Figures of Interest}
Whether it be for glory, for gold, or for freedom from law on the ocean waves, the banner of the Revelry calls many to join the cause.

\subsubsection{Wyatt Maranoss, the Plank King}
\textit{Chaotic evil, male goliath}

\noindent
After seizing the mantle of Plank King, Wyatt quickly established his position with displays of tactical brilliance, fair judgment, and brutal retribution. Wyatt is eager to maintain the approval of his captains and often throws bawdy, drunken celebrations for his subordinates and the people of Darktow to offset periods of unrest or frustration. Initially friendly and welcoming to an almost unsettling degree, Wyatt has a reputation for abruptly turning into a cruel and violent dictator when his trust and hospitality are violated.

\subsubsection{Captain Alyson Paij}
\textit{Chaotic good, female human}

\noindent
Known to many as the "Screaming Knuckle," Alyson is the loud and joyful captain of the ship The Ill Hand who enjoys a good brawl as much as she appreciates a stiff drink. Paij took to piracy because of her hatred for the Clovis Concord, but she is driven by an uncompromising sense of moral duty to the downtrodden. She never raids civilian vessels, and she does her best to rein in the Revelry's cruelest captains and focus their efforts on those who deserve to have their goods stolen.

\begin{DndSidebar}[float=!b]{Cults: Dubious Worship from the Shadows}
While this chapter outlines the major players across the continent operating in ways that directly affect the political landscape of Wildemount, there are numerous smaller factions that can be introduced in your own Wildemount campaigns. Among these lesser powers are countless cabals, cults, and leaderless groups of people who fall under the influence of the numerous Betrayer Gods and other entities who corrupt and manipulate mortals. These may operate on their own or within larger factions, providing challenges, villains, and an element of chaos that can make your Wildemount campaign unique.



Below are a few examples of cults deeply rooted throughout Wildemount, and their respective goals.



\textbf{The Remnants}. When Vecna's physical form was destroyed during the Age of Arcanum, his most devoted followers founded the Remnants, a collection of secretive sects dedicated to realizing Vecna's plan to ascend to godhood, despite his death. The cult succeeded in aiding his resurrection and ascension, but they were scattered when the heroes of Vox Machina banished and sealed Vecna beyond the Divine Gate. The Remnants now seek new ways to restore their god to his promised position of absolute dominance over creation.



\textbf{The Caustic Heart}. The leaders of the Caustic Heart are dedicated to Tiamat and eager to free her from the hellscape of Avernus. The Caustic Heart uses social leverage and binding contracts, as well as seductive promises of power and vice, to ensnare the weak-willed and greedy within their plots and expand their influence.



\textbf{The Harbingers of the Core.} These aberrant cultists worship the terrors born from the dreams of the Elder Evils: impossibly ancient, malevolent entities who exist within the Far Realm and the spaces beyond the known planes. Touched by madness, their bodies often subtly mutated by their contact with the horrors to whom they pray, the Harbingers of the Core gather in the broken lands of Blightshore to summon and unleash the terrible children of their masters.



\end{DndSidebar}

\chapter{Wildemount Gazetteer}
\DndQuote%
{}
{\enquote{I've certainly seen much of it, and it is very beautiful. Beautiful, and terrifying, and cold, and humid, and deadly, and joyous, and full of strife, struggle, and hope. It's home, and that's enough for me. Doty, did you get all that?"}}
{Taryon Darrington, The Daring Trials and Tribulations of Ser Taryon Darrington}
Along the western shores, turbulent ocean currents batter trade ships as they ride the winds toward port cities that shine like diamonds on the verdant coastline. Monstrous behemoths plod across the cracked and blighted barrens of the eastern valley. The desolate northern tundra is a lawless wilderness where power-hungry scavengers hunt for ancient artifacts buried beneath the ice. Wildemount is a vast and varied landscape brimming with possibilities for any adventure.

This chapter presents the various areas of Wildemount, broken into seven major regions: the Menagerie Coast; the Zemni Fields and Marrow Valley, which together form Western Wynandir; the Greying Wildlands and Eiselcross, which shape what is known as the Biting North; and lastly the Wastes of Xhorhas and Blightshore, which comprise Eastern Wynandir. Each section describes the region's important locations and includes plot hooks for adventurers who travel there.

Adventure hooks presented within this chapter are broken down into recommended ranges of adventuring party level:

\begin{itemize}
    \item For low-level characters: 1st to 4th level
    \item For mid-level characters: 5th to 10th level
    \item For high-level characters: 11th to 16th level
    \item For epic-level characters: 17th to 20th level
\end{itemize}

Adventure hooks labeled as "any level" can be presented at any point in your party's journey, with a little planning. You are welcome to adjust the difficulty of any of these adventure hooks to fit your campaign.

\section{Menagerie Coast}
Spanning the length of Wildemount's southwest coastline along the Lucidian Ocean, the Menagerie Coast is a lush, rain-swept, tropical region filled with scattered forests and jungles, rocky seaside cliffs and beaches, overgrown islands, and bustling trade routes that traverse the perilous seas. The coast is ruled by a nation comprised of seven city-states whose cultural and political leaders have united to form the Clovis Concord. This land is known far and wide as a place for lively trade, vibrant art, performance, and excellent food—as well as the indulgence of vice, piracy, and illicit business.

Five of the seven cities have open shipping ports, supporting a robust economy of goods, services, and contraband. The Myriad crime syndicate has had its hands in the affairs of the concord for decades, and while officials publicly deny that the Myriad is still active along the Menagerie Coast, local guildmasters and politicians fear that the Myriad is stronger than ever.

The major cities that comprise the Menagerie Coast are Port Damali, Feolinn, Othe, Nicodranas, Gwardan, Port Zoon, and Tussoa. Beyond the coastal cities, the Swavain Islands are scattered across the Lucidian Ocean, bearing shipping posts, dangerous reefs, and ancient secrets hidden in the bellies of hungry jungles. The pirates of the Revelry stalk the oceanic trade lanes, hassling merchants, marines, and adventurers alike. As any adventurer will attest, all this danger guarantees that travel along the Menagerie Coast is an anxious, thrilling, and profitable experience—if you know where to look.

\subsection{Areas (A-C)}
\subsubsection{Barren Shores}
The northern edges of the Cyrios Mountains drop off into dark pebble beaches known as the Barren Shores. Many believe that an ancient being of immense strength pulverized the edge of the mountains with a series of mighty blows. Others speculate that the shore is somehow related to the Pallid Grove. Whatever the case, these stories of forgotten magic exert an undeniable pull on the curious minds of adventurers and historians alike. Traversing the icy waters of the Frigid Depths is nearly as dangerous as climbing the uncharted expanse of the Cyrios Mountains, and those who survive the journey are greeted by monstrous, twisted crustaceans and seaside cults eager to torment trespassers. Some divine historians theorize that something ancient and cruel might have been released into the nearby sea when the mountains were sundered; they speculate that the strange ring of small islands surrounding the nearby inlet to the west may be connected to the being's whereabouts and purpose.

\subparagraph{Barren Shores Adventures}
Adventurers might seek out the secrets of the Barren Shores through archaeological curiosity or find themselves stranded during an unfortunate sea voyage.

\subparagraph{Be Our Guest (Low Level)}
The characters discover a wrecked smuggler's ship, its ruined structure washed on shore. Inside, they find the long-dead bodies of the smugglers who once manned it, as well as the mysterious animated cargo that slew the previous sailors. Flying swords, rugs of smothering, and \textbf{animated armor} haunt the interior, while the baleful spirit that grants them life remains tethered to a golden chair hidden within the depths of the hull, seeking release from the curse that binds it there.

\subsubsection{Bisaft Isle}
The largest island in an archipelago southwest of Port Damali, Bisaft Isle is a lush expanse of jungle, nut and fruit trees, exotic creatures, and thick mangroves. Swarming with denizens that call the dense foliage home, Bisaft has become a favored location to hide contraband or conduct unlawful activity beyond the sight of the Clovis Concord. Bisaft is dotted with abandoned ruins of the Ki'Nau who once ruled this jungle, and some believe that these villages hide knowledge and gifts granted to the Ki'Nau by their elemental patrons.

Along the southern shore of the island lies a single bastion of civilization, known as Bisaft Post. A small village of no more than eight hundred people, the harbor town is a tropical escape for concord citizens and a surprisingly busy place of business for ships in need of resources and repairs while passing through the trade lanes of the Swavain Islands.

\subsubsection{Brokenbank}
\begin{description}
\item[Population.]
1,520 (74\% humans, 9\% tabaxi, 6\% dwarves, 11\% other races)

\item[Government.]
The settlement is run by the portcaller and his personal Zhelezo soldiers. Minor power struggles occur with local business leaders.

\item[Defense.]
Beyond the Zhelezo, there is a small volunteer militia of fishers and ex-military locals.

\item[Commerce.]
Basic supplies are available, though sources are limited to a few inns and places of trade. Some merchant sailors bring rarer goods on occasion.

\item[Organizations.]
Aside from scattered shrines, no full temples exist in Brokenbank. The settlement is too small to support independent guilds; large businesses control internal commerce.
\end{description}
\noindent
Brokenbank is a coastal settlement built along the beautiful northern beaches of the island of the same name. Its bustling docks that are filled with foreign ships coming and going at a moment's notice. Despite its muddy streets, Brokenbank plays an important role in the trade business of the coast. It is the first major entry port for the Menagerie Coast; most trade ships and shipping vessels check in at Brokenbank to resupply and submit to official inspection before entering any of the coastal ports. However, since Brokenbank lacks the strong cultural identity of the mainland city-states, the concord doesn't recognize it as an official city with a marquis.

The constantly shifting clientele, who often arrive after many weeks at sea, leads to fruitful and rowdy nights in Brokenbank's crowded taverns and brothels.

\subparagraph{Government}
Since Brokenbank has no marquis, its de facto leader is Portcaller Mazin Fahreed (described in \textit{chapter 2}). Mazin performs the duties of a marquis without the title or respect, and he often gets ornery when dealing with other members of the Clovis Concord. He has been known to severely punish those who question his authority. He retains a garrison of hardened Zhelezo who maintain the peace among sailors and peddlers, while aiding the port inspectors to handle any fishy business among the ships in port.

\subparagraph{Crime}
Sailors carrying goods across the ocean often attempt to skim a few goods off the top of their cargo to make an extra silver at market. This is considered normal unless one gets too greedy. Less acceptable is outright stealing from other ships in port before setting sail. The local Zhelezo mete out stern punishment to petty thieves, cross-referencing all documentation to ensure that the sterling reputations of the shipping guilds and Brokenbank itself aren't tarnished.

\subparagraph{Geography}
The busy port consists of a curved main street that follows the shape of the inlet, forming a stretched half-circle facing the northern ocean. It's not uncommon to see the port full of docked vessels while other ships weigh anchor in the harbor, eager for spots to open up. The strong sea breeze keeps the denizens cool under the glaring sun that bakes the dirt and sand along the banks of the settlement.

\subparagraph{Brokenbank Adventures}
The characters could come to Brokenbank to buy or sell a ship or exotic items, or to find a merchant hiring adventurers for protection.

\subparagraph{For the Myriad (Low Level)}
Agents of the Myriad hire the characters to find information they can use to blackmail Portcaller Mazin Fahreed. What evidence might the characters find of the portcaller's dirty deeds?

\subsubsection{Bwualli}
Bwualli is the largest of the Swavain Islands, and its volcanic peak, Itaa, is said to have been one of the many thrones of Kord the Storm Lord while he walked Exandria. The active volcano is a smoldering monolith of black rock and gray smoke that shakes the surrounding lands with sporadic eruptions. Where the glass and rock has cooled, a crowded jungle teeming with aggressive wildlife and lurking pitfalls awaits those who brave the untamed isle. The smugglers and pirates of the Revelry take advantage of Bwualli's dangerous and unstable reputation to hide their plunder on the island, carving out lairs to hide their gains. The recent tales of a nameless horror that stalks the volcanic crags only embolden the pirates. They jeer that the Clovis Concord's officers lack the guts to investigate.

\subsubsection{Copia Lake}
North of Tussoa, the sparkling waters of Copia Lake are nestled within the Copia Wildwood, a thousand feet up in the Cyrios Mountains. The lake is considered a sacred pool where Ki'Nau warriors and priests would ascend to purify themselves of terrible dreams and spirits. As such, the concord has outlawed fishing or diving except for of ceremonial purposes. However, the lake is isolated and only sparsely guarded by Zhelezo, so the area has become a favored drop point for contraband, and passing smugglers have decided to pick up some extra coin by poaching the sacred fish unique to the lake.

\subsubsection{Copia Wildwood}
The Copia Wildwood winds through the Cyrios Valley north of the city of Tussoa, and wraps around the base of the mountains like a coiling verdant serpent. While the outer edges of the forest provide strong timber for the surrounding cities, deep beyond the welcoming tree line lie the swaths of heavy spruce and fir trees infested with all manner of large, hungry predators. Locals spread tales of an intelligent behemoth who lives in the nearby mountains and devours or enslaves any who wander too far into the woods.

\subparagraph{Copia Wildwood Adventures}
This dangerous forest is the perfect spot for characters to face powerful predators.

\subparagraph{Something Wicked This Way Comes (Mid Level)}
Something is forcing the giants of the mountains to migrate southward into the wildwood. The party is hired by the foreman of the Shekis Timber Company to eliminate these unwelcome brutes, offering gold for each giant head delivered. However, when the characters face one of these towering giants, they might notice their sallow figures and the look of fear in their eyes—they're not attacking, they're running from something.

\subsubsection{Cyrios End Forest}
Along the southern edge of the Cyrios Mountains is a dense timberland of softwood trees and jungle vegetation. The Cyrios End Forest provides lumber used to create quality paper, but those who harvest from the woodland generally keep to the perimeter. Deep within the twisting paths and fern-cluttered floor, deadly creatures hunt in the shadows, while abandoned huts and lodges crumble around a toppled temple, now mostly reclaimed by the forest. Adventurers and historians have helmed expeditions into the remnants of this Ki'Nau village and the temple, but have often returned wounded or diseased, or have never returned at all.

\subsubsection{Cyrios Mountains}
The largest mountain range in Wildemount, the Cyrios Mountains wind across the western side of the continent, dividing the Menagerie Coast from Western Wynandir and the Dwendalian Empire. The jagged mountains make overland travel between the neighboring continents of Wildemount and Tal'Dorei all but impossible—to say nothing of the nightmarish beasts that hunt amid the craggy cliffs and gnarled trees.

\subparagraph{Cyrios Mountains Adventures}
The vast mountain range harbors all kinds of challenges for those who value treasure more than their own lives.

\subparagraph{The Unlucky Few (Low Level)}
The characters are hired by a minor official from the Dwendalian Empire to find a small community of Righteous Brand deserters who have set up an encampment deep into the Cyrios Mountains and bring them to justice. When they locate these fugitives, the adventurers discover that the entire community is afflicted with lycanthropy, the reason for their desertion. The characters can attempt to capture or slay the former soldiers to complete their contract, make allies with them and keep their secret, or even search for a cure and help them reintegrate into society.

\subparagraph{Many Eyes for Many a Prize (Mid Level)}
Rumors reach the coast that a hidden mountain valley holds a mysterious collection of statues. The characters are hired to accompany a curator of the Cobalt Soul to investigate the nature of this strange menagerie, only to discover that numerous curators have already been added to this display by an eccentric \textbf{beholder} called Enxanash who wants to build a beautiful museum of zoology. The characters must defeat the beholder, convince it to relinquish its collection, or steal the petrified victims.

\subsection{Areas (D-F)}
\subsubsection{Darktow Isle}
The volcanic rock that forms the jagged cliffs of Darktow Isle is wildly varied in height and shape, creating the impression of a massive, toothy maw. These rocks hide many caverns and tunnels that have inspired numerous local myths of secret hideouts and legendary beasts. The island is barren of greenery, save for the verdant basin in its center, known to locals as the Pit. This small pocket of jungle is the largest source of natural resources on the isle, and its fruit, game, and timber sustain the pirate haven in lean times.

The Dragshallow Reefs that surround Darktow Isle have claimed many ships over the years. Their wrecked skeletons have been dragged atop the rocks as an additional barrier to intruders and as a warning to trespassers. A \textbf{dragon turtle} named Idok stalks the reefs, but an accord between the cunning creature and the Plank King himself keeps Idok away from pirate vessels. In return, the dragon turtle can do whatever it wishes with any other ships that enter the island's twisted waterways. The denizens of Darktow have further fortified the reefs with sentry posts, hidden weaponry, and deployable pitfalls to prevent unauthorized entry through the southern channel. The defenders of the channel have orders to fire without warning on any ships that do not fly the colors of the Revelry.

\subsubsection{Darktow}
\begin{description}
\item[Population.]
1,306 (61\% humans, 12\% elves, 9\% dwarves, 18\% other races)

\item[Government.]
The island is ruled by the Plank King, whose power is maintained through the pirate lords loyal to him.

\item[Defense.]
Numerous hidden ballistae, siege weapons, explosives, and magical wards guard the waterway to the city. The harbor holds a powerful navy of pirate vessels, and most Darktow denizens are accomplished cutthroats, ready to fight any threat to their city.

\item[Commerce.]
Most goods and services can be found here, as well as all kinds of rare, exotic, and dangerous supplies that find their way here from far-off raids.

\item[Organizations.]
Most homesteads have small personal shrines, but only two rickety temples to Kord and Melora stand along the shoreline. Each ship is its own independent organization, but all work as part of Darktow's community.
\end{description}
\noindent
This eerie, fog-shrouded town is built into the southern cliff of Darktow Isle. Over a century ago, Darktow was established as a sanctuary for sailors traveling to Wildemount from Marquet or Tal'Dorei. About forty years ago, a group of aggrieved sailors and traders within the Clovis Concord rebelled against their government's excessive taxation of independent nautical businesses. The rebellious pirates conquered Darktow and claimed it as their headquarters, making this town the beating heart of piracy in the Lucidian Ocean. This union of thieves and pirates, known as the Revelry, plunders ships on the surrounding waters then retreats to their reef-encircled base, amassing wealth and influence through intimidation and cleverness.

\subparagraph{By Land or Sea}
While most Revelry pirates are sailors by trade, many of the isle's residents never take to sea. Instead, they maintain the town's defenses, repair ships in port, entertain pirates on shore leave, or sell goods and services. A large number of these landlocked folk were denizens at the time of the town's capture or are their descendants, having been granted safety in exchange for their allegiance to the Revelry.

\subparagraph{Government}
Within the Revelry, each captain is the lord of their own ship, but all business must be reported to the Plank King, the leader of the Revelry as a whole. Though the Plank King rarely leaves the Throne Roost in the town of Darktow, he oversees operations, organizes raids, protects the island, and punishes traitors.

\subparagraph{Crime}
One would think a landscape of law-breakers would be rife with crime, but the rule of the Plank King heavily punishes thievery within Darktow. A number of bird-picked skulls adorn the rocky shore aside his lair as a warning to play nice with one's fellow pirates. Even so, some folks do push those boundaries should they feel confident in their ability to get away with it. Theft and conflict beyond the boundaries of the Dragshallow Reef is reviewed and judged by the Plank King.

\subparagraph{Geography}
The town's many ramshackle structures sprawl across numerous rock shelves and caverns, connected by bridges, tunnels, and rickety platforms. A vast shipyard surrounds the southern cliffs with a winding network of docks and bridges, while other fortifications bristle with cannons. Other siege weapons dot the jagged edges of the isle.

A series of watchtowers keeps vigil over the inlet, their torches flickering within the midnight fog. A precarious, winding cliffside stair leads to the Throne Roost, the cavern where the Plank King makes his home. Atop that same cliff, a mysterious arcane device is under construction; some believe that it is intended to defend the Revelry's stronghold in case some fool from the Clovis Concord attempts to bring a skyship over their heads.

\subparagraph{Darktow Adventures}
Darktow is a pirate settlement, making it the perfect starting place for an adventure on the high seas. The characters might come to the town to join a crew, challenge a pirate captain, or hunt down criminals.

\subparagraph{Bloody Maggie's Crew (Mid Level)}
The characters are hired by Bloody Maggie (chaotic neutral, female human \textbf{bandit captain}) of the pirate ship Dragon Turtle to find out who or what keeps murdering members of her crew. The characters eventually catch the culprit, Blue Martita, a \textbf{revenant} who wants revenge on Bloody Maggie and her crew for leaving her behind during a Clovis Concord raid on one of their safe houses. The revenant tries to recruit the characters to her cause.

\subparagraph{Heroes on a Dragon Shell (High Level)}
The arrangement between the Plank King and the \textbf{dragon turtle} Idok of the Dragshallow Reefs has broken down, so the creature is now also sinking Revelry ships that attempt to enter or leave Darktow. The Plank King hires the characters to parley with Idok and find a way to meet the creature's demands, convince it to end its assaults, or slay it in its lair for a share of its treasure hoard.

\begin{DndSidebar}{Song of Dashilla the Dreadful}

{
    \small
    \itshape
    \noindent
    \enquote{Where once the waters rolled calmly by,\\
    Where ocean's heart did sing,\\
    There stalked a lass with a hungry eye,\\
    And claws that grasp and wring.\\
    With lantern lure of a glimmer pot,\\
    With treasure in their eyes,\\
    Below they swam to the gold they thought,\\
    Would await them as their prize.\\
    Yet each of them went to surely die,\\
    Their graves to not be found,\\
    For Dashilla and her hungry eye,\\
    Collects the lost and drowned.
    }
    \vfill\textemdash Brokenbank folk ballad
    
}

\end{DndSidebar}

\subsubsection{Diver's Grave}
Deep beneath the waves of the open Lucidian Ocean lies a terrifying graveyard of wrecked, merrow-infested ships scattered across the ocean floor. This sunken cemetery of broken vessels is strewn over miles and miles, and each salvage ship that tries to reclaim cargo from the Diver's Grave only adds another corpse to its hungry sands. The Grave seems to sense ships as they approach and sends out a storm to consume them. At the heart of the Grave is an underwater maze of carnivorous seaweed, and within this dread labyrinth lives Dashilla the Dreadful, a terrible \textbf{sea fury} who hungrily awaits her next shipment of wayward sailors.

\subsubsection{Dolithil Mine}
Northeast of the elven city of Gwardan lies the Dolithil Mine, owned and operated by the Dolithil Orekeepers Guild from within the city. The entrance to the mine remains hidden from view thanks to simple illusions put in place to deter thieves. While the mine yields a respectable amount of gold and platinum, it is notably rich in iron, tin, copper, zinc, and other more common metals that are used in the secret trade of steamcasting (see "\textit{Gwardan}" later in this chapter).

\subparagraph{Dolithil Mine Adventures}
Dolithil Mine is such a critical location for the Clovis Concord that adventurers could be hired to find the mine for the organization's enemies, or they might be paid to protect the place from threats outside and within.

\subparagraph{Revelry Razing (Mid Level)}
The Revelry knows that finding and destroying Dolithil Mine would be an excellent way to hinder their enemies in the Clovis Concord. The characters could be hired by the Revelry to find and destroy the mine, or by the concord to defend the mine from Revelry invaders.

\subsubsection{Doralle Woods}
The Doralle Woods are tucked away at the base of the Cyrios Mountains. The woodlands are a major source of game, timber, and herbal medicines for the nearby elven city of Gwardan and the surrounding coast, fed by the watershed of the nearby mountain range. Fishers often follow the path alongside the Doralle River into the northern lake to catch rare fish and silt-born shellfish only found within those crystal waters. Travelers and hunters that traverse the woods are advised to bring offerings of fresh bread, tropical fruits, and even pearls for the woodland spirits that are rumored to call the forest home. A community of fey beings watch over these woods, stepping from the veil of the Feywild to meddle with trespassers who show disrespect, or terrorize those who seek to rob the wilds of their boons. Local stories speak of a crew of unsanctioned loggers from Tussoa who vanished naught but a week into their venture, only to return a month later aged thirty years, with no memory of what occurred in the interim.

\subsubsection{Everplume}
Those who brave the waters surrounding the Torrid Reef, south of the Shearing Channel, find their view of the horizon broken by an endless column of jet-black smoke rising from beneath the ocean. Hidden under the rolling waves is the monstrous volcano Everplume, which fills the surrounding waters with murky soot, toxic fumes, and boiling heat. Elemental energy swirls around this ominous peak, leading many experts to believe that a rift to the Elemental Plane of Fire might sit at its heart, unknown to the Ashari people, who safeguard breaches between the Material Plane and the Elemental Planes elsewhere in Exandria. Natural creatures of the deep have been slowly twisted into fiery mockeries of their original selves, and regular tremors suggest that the mountain may be growing.

\subsubsection{Feolinn}
\begin{description}
\item[Population.]
12,700 (73\% humans, 10\% elves, 5\% gnomes, 12\% other races)

\item[Government.]
Marquis Hector Soliva runs the city, with wealthy, non-competitive businesses operating as an unofficial council.

\item[Defense.]
Local Zhelezo keep the peace and protect the city perimeter. Affluent members of society also keep a hired guard.

\item[Commerce.]
Most basic supplies are readily available in taverns, inns, and the marketplace.

\item[Organizations.]
Shrines and temples to the Prime Deities are present, if downplayed in prominence. Local merchant families vie for financial and political dominance, leaving room for growing Myriad corruption.
\end{description}
\noindent
Smaller than its nearest coastal neighbors, Feolinn is widely referred to as the Garden City for its rural aesthetic, proximity to the Plumgroves, and well-kept parks and gardens. Feolinn acts as a harbor for the plentiful fishing around the Mother's Sigh Reef, as well as a destination for travelers looking to escape the chaos of Exandria's woes. Feolinn is a popular destination for wealthy socialites, and it is filled with businesses that specialize in indulgence and self-pampering.

\subparagraph{Tree and Vine}
The abundance of fruit and sweet roots harvested from the Plumgroves has given rise to a competitive wine and spirits industry. Feolinn wines and ciders are renowned throughout Exandria. The vintners of Feolinn are known for their intense pride—and for the disdain they have for other vineyards, such as those in the Dwendalian city of Kamordah. Classic vintages auction well in Port Damali, while high-end liquor is presented as a fine gift in noble circles. Competition between the local wineries and distilleries is cutthroat, with major organizations such as the Soliva Winery, the Laffath \& Legrand Vineyards, the Leoleoa Cellars, and the Purple Breast Company all vying for local and international renown.

\subparagraph{Government}
Feolinn's reputation for vanity is well-deserved; even its leaders believe in their city's superiority. Marquis Hector Soliva of the Soliva Winery family is the third marquis of his bloodline to maintain the city's luxurious image, while also doing his best to elevate the prestige of his family's business. Local Zhelezo are often tipped by businesses to pay particular attention to their safety, or to pry into the affairs of competitors. This culture of bribery has caused Feolinn's Zhelezo to neglect the concerns of their city's underclass in favor of lining their pockets with "tips" from merchants.

\subparagraph{Crime}
Visible, violent, and unbecoming criminal activity is swiftly quelled to avoid tarnishing the reputation of the fair city. However, quiet lawbreaking occurs within back rooms of wineries and the mansions of the upper class. A number of powerful local businesses have even fallen in with the Myriad to unseat their competitors, and though few are aware of it, Feolinn is now nearly as corrupt as Port Damali.

\subparagraph{Geography}
Buttressed against the southern tree line of the Plumgroves and surrounded by a number of small family farms, Feolinn is replete with rural charm. The city itself is comprised of four main roads that form a large rectangle, with each corner ending in a promenade filled with shops and pubs, while the city center is mostly made up of residential homes and affluent estates. Suamalie Way specializes in food and fine dining, while Alt'afi Avenue sports theaters, bathhouses, gambling halls, and other venues for private entertainment. Minceli Road contains administration buildings and Zhelezo barracks, and Soyer Street is home to many of the local harvesters and working class folk.

\subparagraph{Feolinn Adventures}
Feolinn is a good place for adventurers to get caught up in shady business deals and cutthroat competition.

\subparagraph{Loose Lips (Low Level)}
The characters are invited to a citywide alcohol fair, where Feolinn's wineries and distilleries enter their products to compete for prizes. The fair is secretly sponsored by the Myriad, which plans to get the city's most influential people drunk and learn their secrets to use as blackmail. The characters get wind of this plan and must decide whether to stop the Myriad or get in on the plan.

\subparagraph{Celebrity Endorsement (Any Level)}
After the characters gain recognition for their deeds, the competing wineries and distilleries in Feolinn attempt to convince them to become official spokespeople of their brands, offering money and free booze. If the characters accept an offer, the rejected companies turn on them, trying to embarrass the characters at every turn and hiring rival adventurers to take credit for their accomplishments.

\begin{DndSidebar}[float=!b]{Wines of Feolinn}
Tropical weather and fertile lands have enabled Feolinn to develop a prosperous wine trade across the Menagerie Coast and beyond. There is fierce competition between local families as well as the rival wineries of Kamordah in the empire, and the numerous wines exported from Feolinn can claim a respectable bit of coin, though the prices might vary depending on the market.




\begin{description}
    \item[Solvia Groves Diamond Plum Wine:] 60gp per bottle.
    \item[Oveso Family Dark Red:] 35 gp per bottle.
    \item[Rustlecall's Evening Dream Plum Wine:] 40 gp per bottle.
    \item[Solvia Groves Sunset White Wine:] 30 gp per bottle.
    \item[Cuthras Silver Vineyard Premium Red Wine:] 50 gp per bottle.
\end{description}


\end{DndSidebar}

\subsubsection{Festerwill Pool}
Festerwill Pool's bubbling waters emit a powerful sulfurous smell that fills the valley, while the nearby rock is discolored by mineral buildup. The pool itself is a green lake choked by algae and weeds at the base of a ragged mountain valley, surrounded by moss-covered, leafless trees. This lake and the surrounding rocky cliffs are home to the grotesque and ravenous gorefeather harpies. Screeching through the air, the harpy troupes bring their prey here to devour, the bloody remnants falling into the brackish waters below. Under their matron, Uada, the gorefeathers claim this region as their territory, hiding baubles and trinkets among their filthy nests.

\subsection{Areas (G-O)}
\subsubsection{Gilded Roadway}
Running nearly the entire length of the Menagerie Coast, the Gilded Roadway is a ribbon of civilization that links every major city in the Clovis Concord. This well-traveled highway is too long to be constantly patrolled, but small bands of Zhelezo do move from city to city, giving travelers at least the illusion of security. The Gilded Roadway is a symbol of unity between the coastal cities of the Clovis Concord and is even represented on the concord's official seal.

\subsubsection{Gravid Archipelago}
These three islands are comprised of steep, volcanic peaks with sheer cliff faces, pock-marked with numerous caves housing blood hawks, harpies, and (rumor has it) an impossibly big \textbf{roc}. There are only small pockets of trees and jungle, so resources are scarce and few ships have a legitimate reason to stop here. However, this does not stop the Revelry from using the region for storage, ambushes, and meetings outside Darktow.

\subsubsection{Gwardan}
\begin{description}
\item[Population.]
18,900 (63\% elves, 16\% humans, 11\% gnomes, 10\% other races)

\item[Government.]
The city is governed by the four Noldar, one of whom must act as marquis.

\item[Defense.]
Gwardan has a respectable force of Zhelezo, mercenaries, and trained citizen militia. Elders and prominent figures keep a small personal guard.

\item[Commerce.]
Most basic goods and facilities can be found within the city, while rarer commodities are scarce. Inns and taverns are limited compared to other coastal cities.

\item[Organizations.]
The Houses of Sagacity (see "Geography" below) are the center of worship, though small shrines dot the city. The Clovis Concord is involved in the city's affairs, though kept peripheral by the Noldar. Myriad interests are slowly invading.
\end{description}
\noindent
After the sea guardian of the Ki'Nau was banished, the Menagerie Coast received an influx of travelers from Marquet. Many Ki'Nau of elven heritage joined the desert elves of Marquet to seek a new home. They traveled inland toward the protection of the Cyrios Mountains, where precious metals and stones lay unclaimed, and paid their respects to the fey folk of the Doralle, taking only what lumber they felt was needed. Together, they built a new center for elven culture among the fields outside the forest.

When the Clovis Concord was established, Gwardan was the force that mandated autonomy between each city-state. The elegantly carved wooden arches of Gwardan rise out of the fields like the waving grass writ large, with neighborhoods of domed yurts rising into tall spires of brass and silver. The Houses of Sagacity allow a limited number of students annually to join their halls to seek higher knowledge.

\subparagraph{Steamcasting}
The elves of Gwardan built sprawling mines throughout the nearby mountains, becoming the most renowned metallurgists and smiths across the Menagerie Coast. The mingling of elven cultures has occasioned the guarded development of a unique style of metalworking. Blending organic designs with super-heated caustic mists, the "steamcasting" process imbues metal with a prismatic shimmer and enables cheap, brittle materials to become as workable and as strong as steel. Gwardan is known across Exandria for its unique steamcast goods, and no one else has managed to get the style quite right, making Gwardan's decor, jewelry, and weapons coveted throughout the world.

\subparagraph{Government}
General city affairs are overseen by the Noldar, a council of four figures who represent the larger facets of Gwardan society. Though the Noldar were established in the interest of fairness, they were forced to select one of their number as their marquis when Gwardan entered into the Clovis Concord. Currently, Marquis Vasan Atrith has led the Gwardan people for just under a century, and is seen as both just and forthright.

\subparagraph{Crime}
While Marquis Atrith has been willing to allow outsiders to visit and live within the city, rising crime and rumors of the Myriad's return to the Menagerie Coast have set her on edge. Thieves have tried numerous times to uncover the secret process of steamcasting, causing her to increase the number of Zhelezo patrolling the city streets. These Zhelezo were drawn from Port Damali, and their presence has given Port Damali's leaders the power to interfere in local affairs, to the Noldar's consternation.

\subparagraph{Geography}
Gwardan is a beautifully rustic yet elegantly refined city of stained wood and gold that rises from the windswept tallgrass fields, cultivated farm land, and clustered trees of the chilly northern plains. The outer ring of the city, called the Anorans Ring, is comprised of numerous domed abodes of bowed lumber and carved bark where families make their homes. The Houses of Sagacity form a prestigious region where three temples built for worship also provide a variety of advanced studies to a limited student body. These are the House of Enlightenment under Ioun, the House of Invention under Moradin, and the House of Enchantment under Corellon.

The Korrwa Ring is the inner ring of Gwardan, where business is conducted between peddlers, merchants, and craftsmen. Within that ring stands the Ornatora, a massive tower of latticed wood spiraling upward like a crafted, leafless tree. Disc-like homes and chambers protrude from the tower at varying intervals, granting the wealthy folk of the city a magnificent view. The peaks of the tower house the Noldar, as well as the central meeting chambers where local political discussion takes place.

\subparagraph{Gwardan Adventures}
Characters could come to Gwardan seeking steamcast goods or any of the city's strange crafting secrets.

\subparagraph{Bring 'em Back (Low Level)}
Talamin Raanan (chaotic neutral, female, high elf \textbf{bandit captain}) once worked for the tribes of Shadycreek Run as an assassin. The elf disappeared five years ago, but spies report that she resides in Gwardan under the name Valemis Tirdra. One of the Shadycreek Tribes, or a faction that wants to see Talamin pay for her crimes, hires the characters to capture the elf. When the characters find Talamin, they learn that all she desires is to leave her violent past behind and live out her days with her new family, including her two young children.

\subsubsection{Hearthstar Peaks}
Visible even from the ocean, the Hearthstar Peaks are a compact, dizzyingly tall cluster of mountains near Port Damali. The highest peak, known as the Hearthstar, has a giant column of ever-burning arcane fire. The Hearthstar is a beacon that warns nearby ships of the rough and unwelcome rocky coast in the region, guiding them eastward to the Edgepoint Lighthouse at the tip of the peninsula and onward to the safe harbor of Port Damali. The sheer slopes of the mountains are incredibly dangerous to climb, but some do enjoy the challenge, and there are occasional rumors of arcanists who float to the highest peaks to hide away their secrets in the topmost caverns.

\subsubsection{Lushgut Forest}
Choking the valley floor between the Stonecage Cliffs and the entire northwest region of the Menagerie Coast, the Lushgut Forest is a monument to the fury and anger of Melora. It is said that a terrible battle between Melora and Zehir took place during the Calamity, when the Cloaked Serpent plotted an ambush by a shrine to the Wild Mother. This ambush poisoned and slaughtered an innumerable number of Melora's chosen. In her anguish and rage, Melora's scream shook the heavens and birthed a twisting, thorny woodland that crushed and consumed the serpentfolk of Zehir, sending the dark god fleeing the valley. The Lushgut Forest remains a physical manifestation of the Wild Mother's vengeful nature; her wrath is embodied in the tall, curling boughs that still ooze toxins from the poison left in Zehir's wake. The poisonous denizens of the labyrinthine wilderness of shadow and fumes are rumored to protect something that lingers near its center, and scholars still debate if these slithering sentinels are born from the will of Melora or the spite of Zehir.

The deep woods of the Lushgut Forest are also home to a civilization of goblins and bugbears known as the Vinewreathed Enclave. This druidic order is led by bugbears who have used their natural magic to break the control which Bane, the Strife Emperor, holds over their westward goblin kin. They are devotees of the Wild Mother, whom they know as \textit{an-Melzidanye}.

\subparagraph{Lushgut Adventures}
Whether traveling from Tal'Dorei or wandering far north of the Menagerie Coast, adventuring parties can find unique challenges within these deep woods.

\subparagraph{Venomous Request (Low Level)}
Wandering deep into the shady heart of the forest, the characters discover a hidden society of yuan-ti who worship the Wild Mother instead of Zehir. However, they bear a cancerous curse that cuts their lives short, which they believe stems from the poisonous sunken temple they choose to guard in Melora's name, slaying the twisted creatures that spill from its corruption. The village leader asks the characters for aid in venturing deep into the deadly temple to seek the source of these terrors, and perhaps their own curse, in hopes of destroying it once and for all.

\subsubsection{Nicodranas}
\begin{description}
\item[Population.]
31,900 (68\% humans, 13\% halflings, 8\% dwarves, 10\% other races)

\item[Government.]
The city is ruled by Marquis Zhafe Uludan, who maintains his position through his family's ties with the Dwendalian Empire.

\item[Defense.]
A respectable army remains stationed within Nicodranas to dissuade empire foolishness. Zhelezo mingle with small squads of imperial Crownsguard. Powerful nobles and diplomats also keep hired swords and guards.

\item[Commerce.]
The Dwendalian Empire is Nicodranas's primary foreign trade partner, and Nicodranas is renowned as a city where foreign travelers can find respite and entertainment.

\item[Organizations.]
The port is home to many places of worship, though temples to the gods approved by the empire are kept to the north side of the city, separated from those of banned gods. The Dwendalian Empire has a strong presence among high society. The Myriad is very quiet but unquestionably present.
\end{description}
\noindent
The easternmost city of the Menagerie Coast, Nicodranas (meaning "gathering of colors" in Draconic) sits close to the Dwendalian Empire's southern border. The proximity to the empire has made Nicodranas as an embassy state between the two powers, and the city hosts most diplomatic relations between the Clovis Concord and the Dwendalian Empire. The second-largest city of the coast, Nicodranas is similar to its larger sister city, Port Damali. Although it lacks the bustle of the unofficial capital, Nicodranas is preferred by travelers who dislike the crowds and rapid pace of Port Damali.

A lighthouse of the Wild Mother overlooks the docks of the Restless Wharf, where traders and sailors maneuver their goods in and out of the harbor. Many of the local businesses are connected with either the ruling Uludan family or the entrepreneurial mage Yussa Errenis in the Tidepeak Tower, creating an additional level of competition. Those with coin to spend eventually find themselves drawn to the delights of the Opal Archways, where dens of indulgence are plentiful, though they are always under the supervision of the local Zhelezo guard.

\subparagraph{The Nearest Escape}
Sitting just beyond the southern borders of the empire, Nicodranas has been nicknamed "the Escape" for many who seek to enjoy the simple pleasures of the tropics—but also because other folk come here to flee from the dreary realities of daily life in the city's plentiful gambling houses. Some successful Dwendalian merchants have even gained the Cerberus Assembly's approval to relocate to Nicodranas.

\subparagraph{Government}
Nicodranas is currently led by Marquis Zhafe Uludan, a powerful guildmaster who is related to Lord Athesias Uludan, a respected member of the Cerberus Assembly. This connection has enabled Zhafe's family line to maintain the tenuous truce between the empire and the Clovis Concord, including the controlled trade and travel between borders. The Zhelezo who guard the city intermingle with assigned Dwendalian sentries, though the overlapping jurisdiction occasionally causes conflict.

\subparagraph{Crime}
The awkward intersection of Dwendalian and concord law makes Nicodranas an especially dangerous place to delve into criminal activity. Both factions are competitive in their enforcement of the laws, so the local criminals are extremely cautious to avoid attention. However, many things considered unlawful within the empire are carefully regulated but nonetheless legal here, giving Nicodranas the dubious honor of being the safest place in Wildemount to indulge in one's vices.

The Myriad is quiet here but unquestionably present, manipulating the black market to maneuver their contraband throughout both the Clovis Concord and the Dwendalian Empire.

\subparagraph{Geography}
This vivid, lively port city is built below the mouth of the Wuyun Valley, where the Cyrios Mountains form the southeastern border of the Dwendalian Empire. The districts within the city vary in architecture and atmosphere, meshing colorful Menagerie pavilions with the gothic, brooding Dwendalian aesthetics as the cultures clash and mingle. Divided into four loose regions, each section is further divided by the lifestyles of the denizens based on which side of the valley they hail from.

The Restless Wharf marks the eastern end of the city's shoreline, with a number of docks, warehouses, and a massive shipping bay run by the Glisteners Guild, a privately owned subsidiary of the Uludan family. Many ships are guided into port by the Mother's Lighthouse, a colossal lighthouse carved in the image of Melora. In perpetual competition with the Open Quay docks, the Restless Wharf is filled with cheap entertainment, pamphlets for services, and false promises of competitive prices. The dock yard in the wharf has managed to keep a tight hold on repair supplies coming into the city, leaving the Open Quay starved for materials.

On the western side of the city, the Open Quay is privately owned by Yussa Errenis, a wealthy merchant mage who purchased the land two centuries ago after most of it burned down. He ignored offers to sell the land to the Uludan Family, and instead built an affordable bastion of rentable business space out of spite for the monopolizing family. Considerably seedier than the Restless Wharf, the Open Quay is known for the merchants' loyalty to Yussa and the diversity of work found throughout the area.

The Opal Archways, a loop of eleven stone arches placed seemingly at random, were created when Nicodranas was founded centuries ago to mark the edges of the city, before it expanded far beyond its original boundaries. Now the Archways are markers for the lively social center of the city, glowing with multicolored torches and incense at dusk when music and evening commerce fill the streets. Housing is quite expensive in this vibrant district, so only nobles, well-off merchants, and political figures can afford to call the Archways home. It's also common for foreign diplomats and moneylenders to own a part-time leisure home within the area.

The Skew encompasses the residential sprawl of the city, a colorful collection of mismatched homesteads whose varying heights and sizes create a jagged skyline. The streets are safer the closer to the Archways they are, but at night the lack of well-lit walkways makes it dangerous to walk through this region alone.

\subparagraph{Nicodranas Adventures}
Characters looking to relax in Nicodranas might stumble into an adventure involving the Myriad, the Clovis Concord, or the Cerberus Assembly.

\subparagraph{Vine Shine (Low Level)}
Overnight, heavy vines grow over the lighthouse of the Wild Mother, and anyone who cuts their way through to investigate the sacred place doesn't return. The characters are hired to investigate and find an infestation of vine blights in the temple. How did the blights come to be? Are they part of a plan to corrupt the lighthouse for dark purposes, or are they a priest's experiment gone terribly wrong?

\subparagraph{Chaos Companions (High Level)}
While the characters visit the city, they witness a noble drop dead as a \textbf{slaad tadpole} bursts from her chest. More nobles die in the same way as the characters investigate. They eventually discover that the victims all visited the same brothel, an establishment run by a \textbf{death slaad} in disguise with several slaad bouncers.

\subsubsection{Othe}
\begin{description}
\item[Population.]
8,320 (64\% humans, 16\% halflings, 11\% half-orcs, 9\% other races)

\item[Government.]
Matters of state are managed by Othe's marquis, while local industrial and religious leaders govern on a smaller scale.

\item[Defense.]
Zhelezo guards and volunteer citizen warriors maintain the law and protect the city.

\item[Commerce.]
Basic needs can be readily tended to through a trading post or inn. Some traveling merchants bring outside goods, but rarer commodities are sparse. Othe exports herbs with natural medicinal properties at high cost.

\item[Organizations.]
A collection of small temples resides in the Apa-feids. Merchant guilds are small but work in conjunction with one another. Myriad agents are embedded in numerous elements of the city.
\end{description}
\noindent
Established around 410 PD by the Ki'Nau as a shamanistic site for meditation and medicinal ritual, the location expanded into the village of Othe when Marquesian explorers discovered the unique properties of the Othemoor and the flora it sustains. Othe is the smallest of the coastal Clovis Concord domains, and the local leaders are generally left to their own devices. This makes the social environment of Othe extremely insular and less hospitable to outsiders who don't go out of their way to engage with the locals.

\subparagraph{A Spiritual Elite}
While the bogs surrounding Othe aren't the most picturesque of settings, the city is a regional center for spiritual guidance. Within the Apa-feids of Othe lie many humble temples, where numerous holy guides welcome travelers who seek a comforting word or a moment of clarity. The spiritual traditions of the native Ki'Nau people are alive and well here, and have melded with worship of the gods to create a unique religious practice. Here, connections are forged between the modern tenets of the gods and the animistic faiths of the Ki'Nau people in search of universal truth.

\subparagraph{Government}
The Marquis Mata Laau-afu is a respected spiritual shaman who has followed in his family's tradition for generations. While he acts the part of the wise and mysterious guide to the people of the marsh, Mata is also a shrewd negotiator and clever bureaucrat when dealing with the other marquises and the politics of the Menagerie Coast. Bolstering his Zhelezo with a handful of faithful druidic guardians of the Othemoor, Mata remains vigilant in the protection of his people and their interests.

\subparagraph{Crime}
As watchful as the marquis intends to be, the growing demand for Othe's unique medicinal herbs has drawn the attention of opportunists and swindlers. Myriad agents masquerading as potential investors and partners in small family businesses have tried to abscond with herbal cuttings, but haven't yet been able to break into the highly secretive society that cultivates them. Those that are caught trying to steal or counterfeit Othenian herbs are usually thrown into the Bogpit prison, and weightier crimes come with severe penalties, such as being left tethered and submerged in the most putrid swamps of the Othemoor.

\subparagraph{Geography}
Shaped like a large crescent, this settlement is built across a series of five small land masses that rise above the waterline of the Othemoor. These landmasses are all connected by bridges, and each of the city's districts is built atop one or more of these islands. First, the Rudwashes are a residential district whose houses have been built using bricks made of red clay from the marsh. The port, known as the Stark Docks, is where small boats and storage buildings are used by the laborers of the city to gather resources. The Apa-feids form the central bulk of the city, with more modern buildings and the elevated hall where the marquis resides. Here, a circle of small temples represents the prime pantheon.

While the winds are often earthy and mingled with marsh air, the simpler, pastoral lifestyle of Othe is more relaxed and spiritual than other locales along the coast.

\subparagraph{Othe Adventures}
The characters may come to Othe to seek guidance from one of the city's many temples.

\subparagraph{Seeking Misguidance (Any Level)}
A ranking official of the Clovis Concord is in Othe seeking guidance on handling the Myriad. The criminal organization hires the characters to pose as holy guides and convince the official to work with the Myriad instead of against them.

\subsubsection{Othemoor}
Thick with humid air stinking of stale water and mildewed wood, the marshland of the Othemoor is considered by most outsiders to be an unpleasant place filled with nasty creatures and sinister trees. As unwelcoming as the landscape might be, the minerals of the bog mud are known to have restorative properties, and a number of extremely rare plants and roots that only grow in the Othemoor have become popular in modern medicine, as well as meals within affluent social circles. The marsh is large and ancient, and has swallowed many secrets over the millennia, some of which have been known to rise to the top of the muck unexpectedly.

\begin{DndSidebar}[float=!b]{Herbs of the Othemoor}
The rich, primal nature of the Othemoor marsh nurtures a number of rare and highly sought-after grasses and herbs. Herbalists, alchemists, and spice traders often visit Othe to bolster their stores. Below are a number of the herbal goods often found for sale within Othe:



\textit{Theki Root}. This thick marsh root tastes bitter but is thought to aid digestive health. When you use an action to consume a dose, you gain advantage on saving throws against the effects of poisonous or toxic substances for 8 hours. \textit{Cost}: 3 gp per dose.



\textit{Muroosa Balm}. This paste made from the muroosa bush is known to help prevent sunburn, but it is also a fire retardant. After spending 1 minute applying a quarter pint of muroosa balm to your skin, you gain resistance against fire damage for 1 hour. \textit{Cost}: 1 gp per dose for sunburn, 100 gp for a full application against fire.



\textit{Olisuba Leaf}. These dried leaves of the Olisuba tree, when steeped to make a tea, can help a body recover from strenuous activity. If you drink a dose of Olisuba tea during a long rest, your exhaustion level is reduced by 2 instead of 1 at the end of that long rest. \textit{Cost}: 50 gp per dose.



\textit{Willowshade Oil}. A dark blue oil can be extracted from the rare fruit of the willowshade plant. A creature can use its action to apply the oil to another creature that has been petrified for less than 1 minute, causing the petrified condition on that creature to end at the start of what would be that creature's next turn. \textit{Cost}: 30 gp per dose.



\end{DndSidebar}

\subparagraph{Othemoor Adventures}
The muck of Othemoor is filled with bodies, treasure, and even more mysterious things.

\subparagraph{Blood Diamonds (High Level)}
When several diamonds are revealed by the shifting of Othemoor's muck, many adventurers, including the characters, head for the swamp, hoping to strike it rich. The diamonds were actually planted by Pillia Ravenosa, a \textbf{vampire}, to bring prey into the Othemoor for her \textbf{vampire spawn}.

\subsection{Areas (P-R)}
\subsubsection{Pallid Grove}
Hidden near the northern heart of the Cyrios Mountains, the valley known as the Pallid Grove is cut off from the rest of Wildemount by deadly peaks. Once a beautiful, untouched woodland atop a mountain plateau, the grove was home to reclusive elves who worshiped Sehanine, protected by the illusory magics of their moon priests.

During the Calamity, Torog and his loathsome minions pierced the elves' obfuscating barriers and burrowed beneath the grove, causing the plateau to sink into a shadowed valley. His servants lapped at the roots of the trees from below, sucking the life and color from them, while many of the denizens of the wood were dragged beneath the earth to be tortured in the bladed halls of the Crawling King. Now, nearly a millennium later, the bleached, broken, petrified forest remains an eerie specter of the natural world, and its cursed history keeps travelers at bay.

\subparagraph{Pale Survivors}
It was thought the elves of the wood had been annihilated during the destruction of what is now known as the Pallid Grove, but a number of them survived among the sunken roots. Living in fear of the Crawling King's cruelty and hiding from the destruction of the Calamity, these enduring elven clans kept to the shadows of the valley and persisted through the ages. The blessings of the Moon Weaver shielded them from harm, and over time, these blessings mingled with their natural adaption to the blighted forest and gave rise to a new subrace of elves known as the pallid elves. Only in recent years have these timid survivors begun to step out into Exandria, curious and confused about the new world that has developed without them.

\subsubsection{Palma Flora}
\begin{description}
\item[Population.]
1,780 (68\% humans, 13\% halflings, 11\% dwarves, 8\% other races)

\item[Government.]
An elder manages settlement-wide issues, but most people tend to their own affairs.

\item[Defense.]
An untrained militia composed mostly of shark hunters defends the settlement from pirates and sahuagin.

\item[Commerce.]
The settlement exports some fish, fruits, and alcohol, but makes most of its money through tourism.

\item[Organizations.]
Local farmers and brewers established the Palma Flora Traders to collectively barter with larger, more powerful ports.
\end{description}
\noindent
The community of Palma Flora occupies the southern tip of the Vezdali Peninsula, and is a favorite vacation destination for merchants and Zhelezo from nearby Port Damali. While this settlement is small, it's famous for its annual shark hunting tournament and a popular cocktail that bears the settlement's name.

Most tourists are familiar with Flora Isle, which is home to the port and a shimmering beach. Most vacationers stay in the Riptide Inn and Tavern, sipping on fruity drinks while watching the shark hunters practice on the southern shore. The island is separated from the rest of the peninsula by a wide channel spanned by a wooden bridge. Most residents of Palma Flora live and work in the northern half of the settlement and retire to the southern island after a hard day's work.

\subparagraph{Palma Flora Adventures}
Though Palma Flora is the perfect place for a vacation, it's not without trouble.

\subparagraph{Tide of Retribution (Low Level)}
Palma Flora's shark-hunting tradition has earned its people the ire of a \textbf{sahuagin} clan living in the Sharkfeather Abyss, a nearby undersea trench. The sahuagin tend to their flock of sharks peacefully and view the shark-hunters as cruel poachers. One sahuagin has made a pact with the leviathan Uk'otoa and sworn vengeance on the people of Palma Flora. This adventure hook is expanded into a full introductory adventure in chapter 5.

\subsubsection{Plumgroves}
The Plumgroves are a luscious woodland filled with soft grass, swaying tropical trees with white bark and purple flowers, and numerous varieties of delicious, celebrated fruit. Nowhere else is there such a condensed collection of sought-after berries and fruit throughout the continent, and the people of Feolinn have based their identity and lifestyle around this abundance.

Local laws limit the hunting of wild game to ensure a controlled ecosystem, and Zhelezo patrol the forest pathways to prevent poaching. The Plumgroves are both beautiful and safe, and it's not uncommon for families and children to be seen strolling through the public trails.

\subsubsection{Port Damali}
\begin{description}
\item[Population.]
82,110 (51\% humans, 16\% halflings, 15\% elves, 18\% other races)

\item[Government.]
The city is helmed by Marquis Olesya Lapidus, who works with local guildmasters and nobles, bolstered by the distant support of the other Clovis Concord members.

\item[Defense.]
Port Damali maintains a sizable army alongside the many Zhelezo who guard the city. Nobles and wealthy merchants keep bodyguards.

\item[Commerce.]
Nearly any sort of goods or services can be found here, legal and otherwise. The high volume of travelers means that there are plentiful inns, taverns, underground markets, and rare auctions.

\item[Organizations.]
Temples to the Prime Deities are common, and some notable shrines to the Betrayer Gods are secreted away within this large city. The numerous merchant guilds run most businesses, but the Myriad is the real power behind the guilds themselves.
\end{description}
\noindent
The jewel of the Menagerie Coast, the original Marquesian colony, and the unofficial capital of free trade outside the empire, Port Damali is a bustling city filled with music, color, and commerce. The massive port hosts ships from all around the world, creating a mishmash of cultures and trades that mingle to form the society of Port Damali. Trade routes splinter from all other sides of the city to ensure rapid transport of goods to destinations throughout Wildemount and beyond.

The city's architecture is varied and inconsistent; the inner wards display their own historical roots or lavish opulence, while the outer wards attempt to combat the slum-like atmosphere with cheap, colorful banners and street buskers. The city has a roiling underbelly of crime, murder, and corruption that many say is regrettable but necessary for the city to function.

\subparagraph{Guided through the Darker Waters}
The waters along the Menagerie Coast are rampant with piracy and roving ocean creatures, threatening trade and travel routes throughout the region. The concord has established a formal navy of nearly four hundred warships to maintain the sanctity of the surrounding ocean. These ships are scattered along the coast, with many stationed in Port Zoon or situated in the waters south of Nicodranas to prevent the current war from bleeding into concord territory, but a quarter of the entire navy is maintained within and around Port Damali.

\subparagraph{To the Highest Bidder}
A popular destination within Port Damali for wealthy merchants, nobles, and collectors from all across Exandria is the Exalted Collection Auction House within the Gilded Esplanade. Run by the esteemed Lord Gabriel Rymmer of the influential Rymmer family, this grandiose, if gaudy, establishment has all manner of rare items, collectibles, and historical artifacts brought from around Exandria to be sold to whomever offers the most gold, with a small portion of the sale filling the family's coffers. Socialites, nobility, celebrities, and the fabulously wealthy of Exandria travel here from all over the world, forging bonds of rivalry with each purchase. When particularly powerful relics find their way to the auction house, it's rumored that beings from the Outer Planes will masquerade as merchants to snatch up such prizes.

\subparagraph{A Steamy Shift}
Until recently, many of the businesses in Port Damali were heated and powered by a system of steam-powered engines, which were installed beneath the city's streets. Unknown to most citizens, a marid had been shackled within the sewers to power the steam engines for decades. Recently, another marid in Nicodranas was freed from a similar situation, inspiring an anonymous adventurer to free the marid in Port Damali from its shackles, leaving the engines inert and many businesses without power. A large group of hired druids and mages are maintaining what functionality they can, but Marquis Olesya Lapidus is desperately looking for an alternative energy source or a replacement for the steam engines.

\subparagraph{Port of Sky and Sea}
Port Damali is the central city for most seafaring business and shipping throughout the Lucidian Ocean and a popular destination for those with gold, but it is also the only city in Wildemount with a skyport. While the empire refuses to establish a skyport within Rexxentrum, in what many believe to be a symbol of their dedication to their isolation, Port Damali occasionally receives skyships from Tal'Dorei or Marquet, drawing crowds that rush to see the majestic display.

\begin{DndSidebar}[float=!b]{Skyship}
Luxurious passage through the clouds themselves—it's a romantic idea cultivated to appeal to the wealthiest citizens of Exandria. As an added bonus, skyship travel avoids the inherent dangers of a journey by land or sea, and is more readily available than teleportation magic. These beautiful airships range in size but generally are the same size as a merchant coaster. Skyships are held aloft and guided with the help of three or more enchanted crystals called brumestones, which are installed along the hull or topside of the ship. The enchantments required to produce and maintain these brumestones are closely guarded by the Alsfarin Union within the distant city of Ank'Harel. While the Alsfarin Union occasionally sells skyships to foreign businesses, they retain absolute control over the maintenance of skyships in operation.



Skyships are reserved for goods and individuals who want to circumvent the slow speed or danger of an ocean voyage and can afford the exorbitant cost of air travel. No commoner can afford to charter a skyship, and the lower classes of Exandria generally see air travel as just another way for the wealthy elite to lord their superiority over the poor. Others, however, see the majestic craft as aspirational and long to become rich enough to afford one, or skilled enough to pilot one.



Port Damali is the only city in Wildemount with an open skyport. The others are in the cities of Ank'Harel, Vasselheim, Emon, and Whitestone. Few skyships in service will deviate from the dedicated shipping routes between these major cities.



\textbf{Purchase Cost}: 100,000 gp



\textbf{Speed}: 10 mph



\textbf{Maximum Cargo}: 10 tons



\textbf{Crew}: 6–10



\textbf{Maximum Passenger Occupancy}: 30



\textbf{Personal Travel Cost}: 2 gp per mile per passenger



\textbf{Shipping Cost}: 1 gp per 100 pounds per mile



\end{DndSidebar}

\subparagraph{Government}
The laws of the Clovis Concord are enforced with varying stringency within each city. Nowhere is this more true than in Port Damali, where the law can be bought like any other commodity. Anything can be purchased here, even the ear of Marquis Olesya Lapidus, the celebrated leader of Port Damali.

The Lapidus family rose to prominence over a century ago, when they shifted from employing merchants to deploying mercenaries. In time, anyone who sought safety on the waves turned to Lapidus family privateers. As the fifth in her family to lead Port Damali, Marquis Olesya steadfastly collaborates with the other leaders of the Menagerie Coast to keep the concord an independent nation, free from imperial control.

However, Marquis Lapidus faces threats on all sides. When she was distracted by the rise of the Revelry, the long-dormant Myriad swooped in and conscripted many of the guilds she was supposed to oversee. Now she struggles to excise the Myriad claws embedded within the belly of her city.

\subparagraph{Crime}
Port Damali has been a pivotal nexus of Myriad activity for some time, and nearly all elements of major commerce within the city are tied to the crime syndicate in some way. Many of the powerful merchants and company owners are either part of the Myriad or are indentured to it. As such, the city is the beating heart of Exandria's slave trade, with slavers passing through often to discreetly move their bounty.

Those who fail the Myriad are shackled and shipped off as living contraband. Myriad operatives are plentiful within the law enforcement of Port Damali, so most illegal activity is covered up or allowed to slide, unless such crime interferes with Myriad interests.

\subparagraph{Geography}
The largest city on the western coast of Wildemount, the resplendent seaside metropolis of Port Damali is nestled against the Tyodan River where it meets the Lucidian Ocean. The city is built atop a cluster of leveled hills, and slight changes of elevation physically mark the variation between wards. Past the farmsteads, beyond the tents and ramshackle homes that mark the boundaries of the city, lies the Prism Path, a series of connecting main roads made of smooth, multicolored stones reclaimed from the ocean shores. These roads carve through the interior wards: the Tumbledowns, the Beaded Alley, the Crescents, the Gilded Esplanade, and the Larboard Light.

\subparagraph{The Tumbledowns}
Within the Tumbledowns, the majority of citizens live in cramped structures of varied shape and design, from clay domes to thatched huts to warped wooden shacks. The architecture resembles a mishmash of cultures from all over Exandria, which, along with the Prism Path, creates a pleasant patchwork atmosphere even in the more run-down regions. A number of small temples are also situated in the Tumbledowns.

\subparagraph{The Beaded Alley}
The Beaded Alley is the center of general commerce and trade within Port Damali, named for the decorative cords that hang over each intersection bearing all manner of streamers and dangling beads. Businesses both young and old all mingle here, competing for the attention of anyone with coin. Many bards and performers bring an air of lively irreverence to the bustling atmosphere.

\subparagraph{The Crescents}
The ward known as the Crescents was named in reference to the seven large, moon-shaped statues that were a gift from the Alkamar family when the city was founded. This affluent area houses competing guild families and well-connected nobles who wish to live near the politics of the city or the delights of the port. The beautiful Pearl Shrine temple to Avandra and the Sunfall Sanctuary temple to Pelor reside within the Crescents.

\subparagraph{The Gilded Esplanade}
The Gilded Esplanade is definitely the most extravagant locale within Port Damali. Each street is lined with the most lauded of artists and performers, with markets lush with expensive wares and high-end imports from around the world. Home to the Auction House and venues for recurring festivities, as well as the spire-like skyport, the Esplanade is known to be the sparkling locale that comes to mind whenever people mention Port Damali.

\subparagraph{The Larboard Light}
The massive harbor that forms the southern edge of the city is known as the Larboard Light. This upscale dockside region is as lavish as docks come. Dozens of ships come into port every day, carrying goods and gold from all over, and many visitors coming to briefly live like a king with the many amusements the city has to offer. The office of the marquis is located on the northern end of this ward.

\subparagraph{Port Damali Adventures}
Between the Clovis Concord, the Myriad, the Revelry, skyships, auctions, and more, adventure abounds within the colorful city of Port Damali.

\subparagraph{Mephit Madness (High Level)}
The \textbf{marid} who was freed beneath the city isn't pleased about its incarceration. Every day, it summons steam mephits to terrorize the people of Port Damali with pranks and mischief. The characters are hired to discover the source of the mephit problem and solve it. They must face the marid in its lair, where they might learn the truth of the city's steam power.

\subparagraph{Hot Purchase (Epic Level)}
The characters aim to legally or illegally acquire a powerful item from the Exalted Collection Auction House. They could want this item for themselves or be hired to acquire it for the Revelry, the Myriad, or another faction. While the characters attempt to claim the item, they are confronted by a \textbf{balor} and a crew of demons who were charged by a demon lord to claim the same item.

\begin{DndSidebar}[float=!b]{The Golden Chain Mercenary Company}
The sheer volume of trade and travel along the Menagerie Coast creates plentiful opportunities for sellswords, bodyguards, and adventurers seeking to make some coin and a name for themselves. One crew in particular established itself within Port Damali nearly sixty years ago and has become the most sought-after, expensive, and efficient mercenary band on this side of the Lucidian Ocean. Known as the Golden Chain, it prides itself on having no affiliations, no morals, and no tolerance for failure. Those who can afford its services earn the temporary loyalty of the most ruthless hunters and enforcers on the coast, no questions asked. The Golden Chain is organized and led by Duos Sabado, son of the famed bounty hunter Quan Sabado who founded the company. Few ever dare to cross the Golden Chain, let alone fail to pay it, for those who have challenged these mercenaries are remembered mostly for their bloody demise.



\end{DndSidebar}

\subsubsection{Port Zoon}
\begin{description}
\item[Population.]
19,120 (80\% humans, 7\% halflings, 13\% other races)

\item[Government.]
Marquis Alamads Haddou collaborates closely with his trusted Zhelezo captains. Numerous crafting guilds guide major facets of industry and influence.

\item[Defense.]
A strong force of dedicated Zhelezo and paid militia keep order in and around the city. Artisans and guild heads keep bodyguards, while newly designed defense weapons are employed along watchtowers.

\item[Commerce.]
Common supplies and businesses are abundant, as are inns and taverns. Esoteric goods can be hard to find, as many of the rare materials imported are already spoken for by the crafting guilds.

\item[Organizations.]
Temples across the pantheon are present, but the crafting guilds are tied to worship of Moradin, and are second in power only to the marquis. Myriad activity is present and steadily growing.
\end{description}
\noindent
Settled between Feolinn and Nicodranas, Port Zoon is the industrial center of the Clovis Concord. Wealthy merchant guilds employ dozens of arcane and scientific experts, and the cityscape is dotted with strange towers, pillars of smoke, and exhaust from raging furnaces. Though it's less vibrant than the other coastal cities and jokingly referred to as "New Othe" by outsiders, what Port Zoon lacks in liveliness and splendor it more than makes up for in manufacturing and cutting-edge science. Entire districts dedicated to metallurgy and smithing produce some of the best anvil masters in the region, and the majority of the goods and tools created here are sold along the Gilded Roadway to support the rest of the Menagerie Coast. A gargantuan temple to Moradin looms over the city, with large cables that tether his towering statue to other high points across the city. Dangling from these cables are glowing lanterns that light the evening sky with flickering flames. A quarter of the docks here are home to the Clovis Concord's reserve fleet, while a sizable portion of the concord's soldiers remain within the northern barracks in times of peace, often taking up small trades of their own.

\subparagraph{A City of Training}
Most artisan crafts and skills can be learned in Port Zoon under various schools and halls run by the guilds of the Clovis Concord. Many consider Port Zoon to be the best place for common folk to learn a trade. Training programs generally have moderate to high fees, depending on the promise of the student, or an agreement that a percentage of the profits from the student's future work will belong to the guild. Many students end up finding work within the guilds that trained them, increasing the appeal of apprenticeships in Port Zoon while also strengthening the stranglehold that the guilds have on production and business along the coast.

\subparagraph{Government}
Often referred to as "Ironpuss" by those not in his presence, Marquis Alamads Haddou is an older dwarf, a master of ironworking, and the head of the Silverbone Smithing Guild. Ornery but passionate about the city and its people, Alamads rules with clear intent, and swiftly punishes anyone who crosses his laws or endangers the people of Port Zoon. His Zhelezo are trained to be naturally mistrusting of idle hands. Convicted criminals are chained and forced into labor for the various guilds, doing the harshest jobs in each industry.

\subparagraph{Crime}
The sheer volume of equipment and goods produced within the city draws the attention of opportunists and counterfeiters from all over. An entire industry has developed to differentiate true work from fakes, and the laws of Port Zoon are especially strict when it comes to forgeries. Unfortunately, Marquis Haddou's attention has been divided across issues of production and supply with the growing war to the northeast, diverting his attention from the growing Myriad influence over the city's businesses and political leaders.

\subparagraph{Geography}
Port Zoon is a compact metropolis of tall stone structures, metal beams and towers, layers of drifting exhaust dust, and rising plumes of industrial soot. The city is separated into three smaller "spheres," or districts, connected to the central Genesis Sphere through main thoroughfares lined by homes and warehouses.

The Harbor Sphere is an oblong district where the docks, seafaring businesses, and warehouses reside, as well as naval housing and weapon storage. The Assay Sphere is home to smelters, coal refinement factories, and the most mysterious of experimental towers, called the Timuafa, where experiments in controlling the nearby weather patterns have led to varying—and sometimes undesirable—results.

The eastern Sentry Sphere is divided between impoverished homes, humble temples, and the barracks of the Sentry Keep, where many Shore Wardens of the Clovis Concord's army remain until called.

The central Hammer Sphere is the largest district and the nexus for the others. Crammed hovels and student housing form the outer ring, while the most industrious factories, forges, and schools cluster around the center Temple Genesis of the All-Hammer.

\subparagraph{Port Zoon Adventures}
The characters could come to Port Zoon seeking incredible arms and armor, then get caught up in adventures throughout the city.

\subparagraph{Criminal Uprising (Low Level)}
While the characters visit Port Zoon, captured bandits forced to work for the guilds break out of their bonds, overpower their captors, and free other brigands working in the city. These criminals take control of vital industries in Port Zoon, threatening to destroy industrial sites unless given their freedom. The characters are hired to quell the uprising. Who is behind this revolt? Maybe the criminals planned it themselves—or perhaps an outside organization engineered it for their own purposes.

\subparagraph{Smith Contest (Any Level)}
The characters get wind of a smithing contest to be held in Port Zoon. The winner gets a set of magic smith's tools and free tutelage from a master in the craft. If any characters participate, they should be wary: the Myriad wants those tools, and the competition is literally cutthroat.

\subsubsection{Quoraska Jungle}
This incredibly dense tropical jungle wraps around the southernmost reaches of the Ashkeeper Peaks, engulfing a massive portion of the southeastern Menagerie Coast. The region is filled with treacherous vine-covered trails stalked by all manner of jungle monsters and venomous vermin. The mist-filled air smells of ancient earth and fresh rain, while carnivorous plant life hides among the roots and brush, waiting for a foolish meal to wander by. The regions surrounding Raska Lake contain a number of small fishing villages populated mostly by the jungle-dwelling tabaxi.

Numerous secrets lurk within the ravenous rainforest, from lost clans of the Ki'Nau to smuggler dens abandoned by brigands who were unprepared to survive in the hungry jungle. Growing interest in the Quoraska Jungle stems from recently recovered writings of the lost people of Quoraska that seem to hint that they might have known of the long-lost city of Obarra and where it vanished to after the Age of Arcanum.

\subsubsection{Ruins of Sepesca}
In the centuries that followed the Calamity, the Ki'Nau people established Sepesca, meaning "Nature's End," as their holiest of burial temples, built over a lush grove where the blood of the Wild Mother is said to have spilled before the Divergence. Every Ki'Nau priest and king that passed was interred within Sepesca's massive stone structure.

An ancient feud between the goblinkin of the Lushgut Forest and the Ki'Nau came to a head a few centuries ago, when the goblins stormed the region and desecrated the Wild Mother's temple. The Ki'Nau protectors fled from the worshipers of Gruumsh, and the abandoned holy site fell into ruin. Vengeful spirits and gnarled monstrosities are said to stalk its sunken halls, guarding the seeds of darkness left by the Ruiner's faithful.

A few descendants of the Ki'Nau across the Menagerie Coast still seek the means to cleanse this holy location and finally bring peace to the forgotten line of kings.

\subsubsection{Rumblecusp}
The isle of Rumblecusp is a beautiful paradise of sparse trees and white sands, topped by the single mountain that gives the island its name. Only a few dozen indigenous folk call the island home. The shallow shores of the island often fall beneath the sea during high tides, and a small trail of smoke frequently rises from the volcano. Be that as it may, the waters surrounding the island make for fine fishing, and it's not uncommon for fishing vessels from Nicodranas to travel to Rumblecusp during unfavorable fishing seasons.

\subsection{Areas (S-Z)}
\subsubsection{Siltbasin Pass}
A muddy forest of coastal swamp trees and mangroves that occasionally dips just below sea level, the Siltbasin Pass is a well-traveled portion of the Gilded Roadway between Gwardan and Port Damali, as well as one of the major coastal locations to harvest seaside fruit and clay. Small clusters of huts and hovels line the edges of the Siltbasin every few miles where the locals gather their yields and prepare them for travel to any of the major Clovis Concord cities. Bad weather can often make the pass extremely difficult to traverse, so some caravans have to camp for several nights while waiting for the pass to open once more. Their travels are further threatened by the many swamp creatures and night predators that also call the Siltbasin home.

\subsubsection{Stonecage Cliffs}
A craggy, nigh-impassable range of coastal mountains and bluffs that wraps around the entire northeast peninsula of Wildemount, the Stonecage Cliffs have been a natural barrier between Wildemount and Tal'Dorei since before the societies of each continent first began to communicate. The cliffs enclose the dense and perilous Lushgut Forest, and this combination of deadly terrain has rendered most of the cliffs unfit for the construction of towns or cities. As such, the region has remained a monster-infested wilderness since time immemorial.

The frothing waters of the Shearing Channel churn against the towering cliffs, and thousands of years of tempestuous wind and water have pockmarked them with caves that harbor terrible flying and climbing beasts. It is said that only the bravest—or stupidest—explorers attempt to seek what secrets might lie among the rock and tides of the Stonecage Cliffs.

\subsubsection{Tussoa}
\begin{description}
\item[Population.]
15,110 (74\% humans, 10\% elves, 8\% halflings, 8\% other races)

\item[Government.]
Beneath the just rule of Marquis Tusila Latu, community leaders and wealthy families helm guild business.

\item[Defense.]
A small force is maintained for the defense of the city alongside the protection of the Zhelezo. City leaders often keep bodyguards.

\item[Commerce.]
This rural city can provide most common services and products, though exotic goods might prove difficult to obtain. Outside exporting lumber, the biggest businesses in the city are food, lodging, and entertainment.

\item[Organizations.]
Humble temples are found near and around the Lotoe Plaza. Family businesses work beneath the marquis and the local guild masters.
\end{description}
\noindent
Tussoa began as a Ki'Nau freshwater fishing village along the Tyodan River, but it has grown into a respectable city in the centuries following the creation of the Clovis Concord. The central district where the two main streets of the city meet has larger, more modern structures and bustling streets, but the outer regions are still dominated by humble homes of river reeds and straw.

Tussoa boasts a more relaxed way of life than most of the other coastal cities. Its people try to live a simple life, unburdened by business or politics, but they still grapple with issues of cultural identity. Tussoa's most immediate problems are the threats of the natural world that menace this frontier city on a daily basis.

\subparagraph{Logging a Profit}
The competing Majest and Regver timber guilds are the leading suppliers of lumber to Port Damali. They work along the Copia Wildwood and send their trunks down the Tyodan to their mills at the port, usually accompanied by hired hands to protect their goods. Because the prime logging sites of the forest are hotly contested and not overseen by the Zhelezo, both guilds hire mercenaries to guard their work and clear out nearby monsters stalking the shadows of the wood.

\subparagraph{Government}
Marquis Tusila Latu currently directs the rule of law in Tussoa, taking particular pride in her Ki'Nau heritage—and the Ki'Nau heritage of her people and her city. The marquis treats her Zhelezo and constituents well, eager to build loyalty in a time of creeping Myriad interests. Marquis Tusila insists on being involved in all trials where life and death are concerned and is renowned as a champion of justice, unlike her corrupt predecessor, Marquis Maros du Simmas.

\subparagraph{Crime}
Try as the marquis might, the comparatively rural location and understaffed Zhelezo of Tussoa make the city a breeding ground for petty thieves, charlatans, and forgers. Independent crime rings operate in small bursts, then halt their work for weeks or months to cover their tracks. There are some thugs within the city with connections to the Myriad, but their current intention is to watch and wait until the time is right.

\subparagraph{Geography}
Laid out in the shape of a large V, this bustling fishing city of lively docks and a hard-working populace is spread out over two main roads that meet at the northeast edge, with numerous smaller roads connecting the major thoroughfares. The northern road is Pulum Road, and most of the city's shops, inns, and storage houses are kept to the road's eastern edge, interspersed with neighborhoods made up of simple shacks. The southwestern road is Vaigvasa Road, built along the edge of the Tyodan River. Here, more homes and businesses fill the inner regions, while fishing huts and river logging stations line the southern, outer side. The two roads meet at the Lotoe Plaza, home to the Olo, the repurposed temple to Uk'otoa that now acts as the court of the marquis. All holidays and celebrations are centered in the Lotoe Plaza, which is known to become crowded enough with people on such days that the streets can barely be seen.

\subparagraph{Tussoa Adventures}
Tussoa's bustling logging industry, minimal presence of Zhelezo, and location near the Copia Wildwood make it an excellent place for urban adventures to overlap with wilderness expeditions.

\subparagraph{Log Wars (Low Level)}
A business war turns physical when a camp of Regver loggers is murdered and evidence points to a nearby camp of Majest loggers as culprits. Regver responds by burning Majest vessels, and the violence continues to escalate. The marquis fears that the war might spill into Tussoa's streets, so she hires the characters to negotiate a peace. They might find that the Revelry committed the inciting murders as a way to cause chaos in the city.

\subsubsection{Twinward Isles}
A sacred location to the native people of the coast, this pair of surprisingly tall islands are all cliffs, with no beaches. While the surrounding waters are often choppy from the waves that crash against the nearby Mother's Sigh Reef, the seaside cliffs are filled with healthy fruits and game to hunt. At the inception of the Clovis Concord, the Ki'Nau people demanded that the islands' natural resources were to always remain untouched. This agreement has been dutifully upheld over the centuries.

While open to common use, the dangers of the trek still limit those who choose to make use of the land here. Tales tell of a \textbf{couatl} spirit bound to the inner reaches of the jungles, willing to bestow wisdom and insight to those who bring a mighty offering.

\subsubsection{Urukayxl}
The island of Urukayxl was once the site of the temple of the leviathan Uk'otoa. Many centuries ago, the spiteful followers of Zehir murdered Uk'otoa's Ki'Nau adherents and used the temple as the seal preventing Uk'otoa from returning to the Material Plane.

Forgotten in the years since, the fallen temple is surrounded by jungles that now harbor feuding clans of lizardfolk and yuan-ti. These yuan-ti are servants of the Cloaked Serpent and still maintain the crooked temple in the image of their banished god. Between the disease-carrying insects, the deadly jungle, and the treacherous Inkclaw Reef, many consider the region cursed and avoid it at all costs.

\subsubsection{Vezdaweald}
The Vezdaweald is a bright and sunny tropical forest composed of softwood trees that loggers from Port Damali commonly harvest for smaller crafts and goods. The beautiful woods are full of vibrant tropical flowers and birds that fill the air with a symphony of mingling music and calls—but it's at night that the forest truly comes to life. Millions of glowing insects emerge and transform the canopy into a dancing starscape.

Weddings and religious ceremonies of all faiths take place amid the stars of the forest—but the lights are not always a friendly omen. The glowing insects also summon forth the nocturnal dinosaurs, amphibious Swavain basilisks, and other night beasts that call the forest home.

\subparagraph{Vezdaweald Adventures}
The softwood trees and beautiful scenery of the Vezdaweald attract loggers in need of the characters' protection from the forest's denizens.

\subparagraph{Dinosaur Ivory (Mid Level)}
A wealthy artist in Port Damali wants to use dinosaur bones for her next sculpture, so she sponsors a contest. The adventurers that bring her the most dinosaur bones by weight will be gifted 5,000 gp and immortalized in the sculpture.

\subsubsection{Vide Cay}
The isle of Vide Cay is completely surrounded by vicious coastal rocks and choppy waters, so most ships avoid it. The island itself is a series of large, craggy hills that surround a small valley where a small society of fey live within an enchanted tropical forest. Revelry captains who try to use the island complain of a faerie witch that protects the isle, scrambling the minds of those who attempt to rob or enslave the denizens of the trees.

\subsubsection{Wuyun Gorge}
The dry, rocky valley that divides the southern pass between the Cyrios Mountains and the Ashkeeper Peaks is known as the Wuyun Gorge, home to the Wuyun Gate that acts as the transition point between the Dwendalian Empire and the Menagerie Coast. The road through the gorge is well-traveled, seeing carts and caravans throughout the year. Wandering creatures, violent giantkin, and hungry monsters occasionally descend from the mountains to prey on travelers. As such, it's common practice for people of means to hire mercenaries in Nicodranas or Trostenwald—the closest settlements on either side of the Wuyun Gorge.

\section{Marrow Valley}
The southern slopes between the Ashkeeper Peaks and the Cyrios Mountains are filled with the beautiful green grasslands of the Marrow Valley. This rural expanse of lush farmsteads and formative industry welcomes the majority of travel and trade from outside the empire. Fields of grains and grasses stretch into verdant forests and crystal lakes within this southern region of Western Wynandir. The central metropolis of Zadash is a major nexus for business and travel within the region, calling people from all walks of life to find work, bargain, and carouse among the many vendors and guilds that run the city. The imperial military is headquartered in Bladegarden, on the eastern edges of the Marrow Valley, where the bulk of the standing legions await orders and train the next generation of imperial soldiers to join the conflict against Xhorhas. While the war keeps many eyes turned eastward, the Marrow Valley has its own host of unspoken threats and dark secrets hiding beneath the veneer of everyday rural life.

\subsection{Areas (A-C)}
\subsubsection{Alfield}
\begin{description}
\item[Population.]
3,410 (61\% humans, 22\% halflings, 11\% gnomes, 6\% other races)

\item[Government.]
Starosta Kosh Clearbarker cravenly delegates the responsibilities of governance to the local watchmaster and lawmaster.

\item[Defense.]
A small cadre of Crownsguard keep the peace as the empire focuses on larger settlements.

\item[Commerce.]
Alfield is a minor stop for traders, with little beyond basic supplies and a couple of inns.

\item[Organizations.]
Minor public shrines offer solace to the faithful. Businesses are small and independent, requiring no guild oversight.
\end{description}
\noindent
Alfield is a rural town nestled into a scattered patch of trees near the Amber Road. The old mines winding beneath the nearby hills yield a fair amount of tin, copper, and quartz that supplement the town's economy. Small in size compared to the surrounding townships, Alfield is regarded by most travelers as a stop on the route to more important business elsewhere. This frustrates some proud residents and local entrepreneurs, but most prefer the small-town lifestyle and are happy to see the passing traders and travelers move on. Alfield is currently recovering from a recent assault by a roving band of violent gnolls, so a number of buildings are still being repaired or rebuilt.

\subparagraph{Government}
The starosta, a middle-aged gnome named Kosh Clearbarker, leans heavily on his dutiful watchmaster, the half-elf Bryce Feelid. Bryce personally shoulders the town's burdens while the starosta continues to oversee the town from within his quarters, excusing himself from any hands-on governance.

\subparagraph{Crime}
Since Alfield is both overlooked and under-protected by the Crown, charlatans and thieves often sweep through town in an attempt to relieve the townsfolk of any loose coin. The relatively small Crownsguard presence is perpetually frustrated with the lack of support from the empire; they must split the attention of their small force between policing the borders and watching for hooligans who seek to abuse the civilians within the town.

\subparagraph{Geography}
The welcoming farming community and beautiful countryside surrounding Alfield often call to those who want a simpler life in the empire. Flanked to the north and east by the Alfield Wood, the town is shaped like a small diamond in shape, only a couple of miles across on all sides. The Dustway and Harvest Path roads divide the town into four quarters, meandering lazily through the ramshackle wooden buildings until they meet at the recently rebuilt Candleglow Inn at the center of town.

\subparagraph{Alfield Adventures}
Between the recent gnoll attack, overburdened town leaders, and numerous criminals, the people of Alfield need the aid of adventurers.

\subparagraph{Yeenoghu's Army (High Level)}
While the characters are staying in Alfield, they learn the truth of the recent \textbf{gnoll} attack. The first battle was just a way to test the town's defenses for the real war: one fought by demons! The gnolls sacrificed the townsfolk captured during their raid to open a portal to the Abyss, letting demons loyal to Yeenoghu, the Gnoll Lord, invade the Material Plane. With so few defenses in the settlement, the characters are the only people who can truly defend Alfield from the threat.

\subparagraph{Myriad Takeover (Any Level)}
Overworked leadership and weak law enforcement mean that Alfield is ripe for Myriad takeover. The characters could be hired to lead the Myriad forces, bringing the town's leadership and petty criminals to heel, or they could fight against the takeover, waging battles both physical and political to keep the town out of the syndicate's grasp.

\subsubsection{Ashguard Garrison}
\begin{description}
\item[Population.]
5,720 (74\% dark elves, 26\% other races)

\item[Government.]
The garrison is currently helmed by the acting general of the Kryn forces holding the stronghold.

\item[Defense.]
The vast majority of the occupants are frontline soldiers and deadly combatants. Mobile siege engines stand at the ready.

\item[Commerce.]
Beyond the rations needed to survive and small-scale trade between the soldiers, little commerce is available.

\item[Organizations.]
The outpost is currently held by the Kryn Dynasty. Temporary shrines to the Luxon have been installed.
\end{description}
\noindent
The Ashguard Garrison is the partner fortress to the newer Rockguard Garrison, and the two strongholds flank the city of Bladegarden. On the south side of the Brokenveil Bluffs, Ashguard sits atop a steep mountain pass that slides into the Brokenveil Marsh of Xhorhas, making it near impossible to assault from the front. Although the garrison successfully repelled the dangers beyond for many years, it fell to a series of clever assaults from the Kryn—they used magic to darken the sky and cause confusion while their forces burrowed beneath the battlements. Over a thousand Dwendalian soldiers were slain, and the army was pushed back to Bladegarden.

Now occupied by the forces of the Kryn Dynasty, this garrison has been refitted to defend against the very people who built it, its walls darkened by the shadows of perpetual twilight.

\subparagraph{Ashguard Garrison Adventures}
If the characters are at Ashguard Garrison, they're likely either reclaiming the garrison for the empire or working for the Kryn Dynasty on the front lines of the war.

\subparagraph{Take Back Ashguard Garrison (Mid Level)}
The characters are tasked with taking back the stronghold for the empire. Since a direct assault on the stronghold would likely end in the characters' demise, they are encouraged to use magical disguises, infiltrate the outpost, and form a strategy over a period of months while they gather intelligence. During this time, the characters might learn that the Kryn Dynasty isn't the evil threat it was made out to be, prompting them to question their loyalties.

\subsubsection{Ashkeeper Peaks}
This grand mountain range extends northward from the Lucidian Ocean through the center of Wildemount, marking the boundary between the Dwendalian Empire and the wastelands of Xhorhas. The Ashkeeper Peaks reach skyward in a jagged array of steep mountaintops, like the pikes of a shadowy army looming over the Marrow Valley. The mountains house the treacherous paths to the Gehen Basin and the hidden valley of the Verstglade—secret lands that only the truly courageous or truly desperate attempt to reach.

The hostile landscape of the Ashkeeper Peaks has prevented any attempts to install an imperial presence, and the peaks are teeming with all kinds of monsters and terrible beasts. Most of the cities bordering the Ashkeepers request a reinforced assignment of Crownsguard to protect their citizens from the many roving mountain predators.

Beneath the surface of the mountains, huge networks of caverns and tunnels weave between subterranean civilizations and the deadly shadows of the Underdark. Unknown to the empire, a massive fire giant citadel rests at the northern heart of the mountain range.

\subsubsection{Berleben}
\begin{description}
\item[Population.]
3,230 (63\% humans, 17\% halflings, 12\% gnomes, 8\% other races)

\item[Government.]
Starosta Carliyle Vesperwind has left most of the governing to idolmaster Julienne Droogba and other community figures.

\item[Defense.]
Though a bare-bones garrison of Crownsguard guards the town, sellswords protect travelers within the Labenda Swamp.

\item[Commerce.]
Supplies beyond basic goods and tools are uncommon. A handful of meager inns provide limited lodging.

\item[Organizations.]
The idolmaster maintains the small shrines within the city. Local business is thriving, but few outsiders are willing to invest in this lopsided backwater.
\end{description}
\noindent
Near the green waters of the Ounterloch, the marsh town of Berleben is nestled deep in the Labenda Swamp. Berleben was settled early in the expansion of the Dwendalian Empire, and the town initially showed promise with a booming industry in the cultivation of unique swamp plants sought after by healers and mages alike. However, the watershed from the Silberquel Mountains runs into the Labenda Swamp, and the land slowly grew unstable, causing half of the town to begin sinking into the surrounding marsh. Those who remained in the town have had to adapt over the years, as much of the settlement now sits uncomfortably under inches of swamp water.

Other regions of the empire regard Berleben with distaste and ridicule, colored by exaggerated stories of rampant illness, a perpetual plague of buzzing insects, and the failing landscape it was built on. Berleben is sometimes call "the Sog Sticks," a moniker used by outsiders to ridicule the town and by its denizens with sarcastic pride.

The ragged swampfolk that call this place home have remained due to familial ties, impoverished necessity, or the promise of work harvesting what natural commodities the swamp offers. The dangers of the swamp and the isolated location have bred a tight-knit and nosy community. Neighbors frequently pry in each other's business, aggressive curiosity taking hold whenever something new or out-of-the-ordinary occurs within the murky streets.

\subparagraph{Government}
Atop the Stilts region of the settlement lies the Summit House, home to Starosta Carliyle Vesperwind. Vesperwind greatly despises his assignment and quietly resents his inability to protest his appointment. He spends most of his time locked within his abode or traveling to neighboring locales to avoid responsibility. The governance of the town is mostly left in the hands of the Mossmire Temple's idolmaster, Julienne Droogba.

\subparagraph{Crime}
Without an active starosta or a well-staffed Crownsguard, crime in Berleben is fairly common and rarely punished. Most folk who live within the town carry their own means of self-defense and are slow to trust outsiders.

\subparagraph{Geography}
Built over unsteady, stinking swampland, the somewhat miserable and half-sunken town is divided into thirds. The three districts correspond to different ages of construction in the empire and visually differ based on whatever construction materials the townsfolk had access to at the time.

The Puddles are the oldest and poorest region of Berleben, with numerous old stone and wood buildings that now sit in several inches of marsh water. The Midway Docks are the central region of Berleben, where the workers and fishmongers do business and where a majority of locals socialize. The Stilts are, for what little it's worth, the newest and most affluent part of town, marked by sturdier houses atop heavy load-bearing stilts that raise the buildings a few feet above the swamp in hopes of outlasting the rest of the sinking settlement.

\subparagraph{Berleben Adventures}
Berleben is an all-but-lawless settlement that needs heroes to help keep the peace. Likewise, it is a wonderful place for unscrupulous sellswords, since no one takes notice of petty crime.

\subparagraph{A New Drink (Low Level)}
People in Berleben are obsessed with swamp swill, a new whiskey that gets folk drunk more quickly and for longer. Some claim to have wonderful reveries while intoxicated by swamp swill, and many have become addicted to the stuff, resorting to violence to get their next drink. If the characters investigate, they find that a tribe of \textbf{lizardfolk} is making the swill and giving it to a merchant in Berleben to sell. The lizardfolk plan to take over Berleben once enough people are dependent on the swill, but only because trolls threaten the lizardfolk's territory.

\subsubsection{Bladegarden}
\begin{description}
\item[Population.]
9,910 (32\% half-orcs, 21\% orcs, 25\% humans, 22\% other races)

\item[Government.]
General Kruth Doorav and Starosta Leman TuFinn are often at odds, but work together to manage Bladegarden.

\item[Defense.]
The city, which houses a large number of Crownsguard, is perpetually host to a large legion of Righteous Brand soldiers and soldiers-in-training.

\item[Commerce.]
There are many modest inns and taverns. Armaments and basic supplies are readily available, but more luxurious services and goods are extremely difficult to find.

\item[Organizations.]
The Righteous Brand maintains a strong presence within Bladegarden. Small temples dot the knotted streets.
\end{description}
\noindent
When the Julous Dominion was defeated, the expanding empire saw the nomadic Odakar orc clan, formerly under Julous rule, as an opportunity to secure their hold on the eastern border. The empire reached out to the more approachable members of the orc settlement of Bladegarden, offering them wages, resources, and access to imperial goods and services. A number of them became the cornerstone of the eastern branch of the Righteous Brand and emboldened the empire's military techniques with the prowess of the border-hardened orcs. Other orc clans who rejected these ideals were forced eastward into the Brokenveil Marsh and forgotten.

From this accord, the new Righteous Brand training facility in the central Marrow Valley was established. Bladegarden was the site of one of the bloodiest battles of the Calamity; its name referred to the volume of weaponry scavenged by orcs from the war-ravaged earth, and has since been adapted to the idea that warriors are grown within its walls. Bladegarden is both a training complex and the core of orcish society in Western Wynandir. Contests of skill and might are commonplace among the recruits and their trainers. Even the merchants come from sturdier stock and often partake. A sense of honor and pride fills the ranks here, stoked by the rivalry between Bladegarden and Grimgolir.

As scuffles with the Kryn escalated and a growing number of dragonborn refugees left Talonstadt to join the military, a certain belligerent mirth filled the warriors of Bladegarden. This oddly jovial wartime atmosphere was shaken when the Xhorhasian assault on the Ashguard Garrison proved far more swift and brutal than anyone predicted. However, now that the surviving Ashguard forces have retreated to Bladegarden and the scattered armies of the empire have moved east to join them, the city is as crowded and lively as ever—but morale is tenuous. One more defeat like the one at Ashguard could crack Bladegarden like a nut.

\subparagraph{Government}
While Bladegarden is technically a military settlement under the watch of General Kruth Doorav, the Crown also assigned a starosta, Leman TuFinn, to oversee the day-to-day minutiae of imperial governance.

\subparagraph{Crime}
As the chaos of the war grows more frenzied with each passing week, the overworked Crownsguard pay little attention, and the distracted populace make easy marks for petty criminals. However, those who are caught find themselves among many well-trained soldiers who often deal out punishment themselves in times of crisis, so pursuing crime in Bladegarden is a very risky, if sometimes rewarding, endeavor.

\subparagraph{Geography}
Aside from the impressive halls of the sprawling Righteous Brand barracks that loom over the eastern edge of the city, Bladegarden is considered by many a remarkably drab place. The oval settlement is threaded by wandering dirt roads that seemingly stop and start without reason, and there is little color in the streets. Beyond the western walls, there are a number of farmsteads where local produce is gathered to aid in feeding the populace, but the majority of food is brought via cart from other sources across the Marrow Valley.

\subparagraph{Bladegarden Adventures}
Bladegarden is an excellent base of operations for characters fighting in the war on the side of the empire, or for Xhorhasian spies who wish to infiltrate their enemies' ranks. With the Crownsguard overworked, there are also plenty of adventures fighting crime within the city.

\subparagraph{Morale Pals (Any Level)}
The Crown sends the Morale Pals, a famous comedic halfling duo, to keep spirits up in Bladegarden. But when the curtain goes up on opening night, it reveals the decapitated bodies of the Morale Pals! The characters are hired to answer the question everyone is asking: who killed the halflings? It might be Xhorhasian spies determined to keep morale low, thieves who wanted the halflings' purses, or perhaps a rogue killer who wanted to send a message to the whole city that this is just the beginning.

\subsubsection{Brokenveil Bluffs}
Long ago, the Brokenveil Bluffs were once a section of the Ashkeeper Peaks that linked to the nearby Dunrock Mountains. The mountain ranges were sundered in the final battles of the Calamity, creating two mountain passes—the only land route that connects Eastern and Western Wynandir.

The bluffs are known for their sheer cliff faces and brisk winds that loudly whistle through the network of caves and channels between the cracked rock. This region is home to many winged monstrosities, sure-footed terrors, and swamp creatures who wait at the boundaries of the Marrow Valley to swoop in and snatch the wayward traveler or soldier. The stories of the Brokenveil Bluffs are nearly as ominous as those told about the Xhorhasian landscape beyond.

\subsubsection{Crispvale Thicket}
To the east of the Silberquel Ridge, surrounding the juncture between the Glory Run Road and the Crispvale Roadway, the shady Crispvale Thicket marks the center of the eastern valley. A haven from the bustling lifestyles of the empire, the dense elms and tall underbrush shelter a handful of firbolg clans living away from the eyes of the Crown. As long as one sticks to the roads, it's normally a peaceful place to pass through, but the deeper areas of the thicket are stalked by hungry behemoths and bandits.

\subsubsection{Cyrengreen Forest}
At the southwestern end of Wynandir, the Cyrengreen Forest hugs the base of the Cyrios Mountains. The forest is the second-largest source of lumber in the empire, with logging companies from Deastok and Trostenwald working the north and south sides, respectively. Most logging is limited to the outer tree line due to dangers in the shady interior of the Cyrengreen. The lush heart of the forest is home to some wilder families who brave the dangers of the forest to live a freer life, as the tax reapers are reticent to visit too often. There are also scores of abandoned cabins scattered throughout the deep woods where it is believed folks were either forced to flee or were killed by the wild denizens of the forest. To stumble on such a site is considered terrible luck; local tales speak of hungry fey spirits and shapeshifting hunters that stalk the dark shadows of the forest in search of prey to toy with. The Cyrengreen Forest is also one of the wild sanctuaries of those cursed with lycanthropy, though the sects of those shackled with this terrible affliction are divided between those who struggle to contain their feral hunger and those who embrace their predatory instincts. The conflict between these factions builds with each passing moon cycle, threatening to spill into the nearby cities.

\subparagraph{A Superstitious Shade}
While the dangers of the Cyrengreen include the more common threats such as subterranean ankheg nests and packs of displacer beasts, mischievous fey also call the forest home, wading through errant moonbeams on bright evenings in search of distraction. Also hidden within the forest is the obscured garrison of the Claret Order of the Lycan, where outcasts cursed with lycanthropy train to control and harness their feral burden away from the judging eyes of society.

\subparagraph{Cyrengreen Forest Adventures}
Crossing through the Cyrengreen Forest is an adventure that has travelers looking over their shoulders at every stray sound. The creatures in the shady depths of Cyrengreen love to toy with travelers.

\subparagraph{Lumberjack Madness (Low Level)}
Lumberjacks from Deastok and Trostenwald working in the Cyrengreen have reported strange events. They fall asleep randomly on the job, witness twinkling lights, see horrific nightmarish creatures for mere moments—and some even turn into beasts such as frogs and rabbits for short periods of time. The loggers are spooked and refuse to return to their work until the characters deal with the problem. The problem is pixies—they're upset that the loggers are tearing their forest apart! The good fey are stubborn but might cut a deal with the loggers if the characters agree to negotiate.

\subparagraph{More Perfect People (Mid Level)}
Reports of wolves attacking travelers on the road from Deastok, Kamordah, and Trostenwald have become more frequent. All travelers tell the same story: the beasts tear through caravans, biting as many people as possible before moving on. There is never a death toll. The reason for the strange attacks becomes obvious during the next full moon when the bite victims turn into werewolves. The characters are hired to find and stop the lycanthropes spreading the curse.

The instigating werewolves live in the dark heart of the Cyrengreen. They call themselves the Brotherhood of the Moon and aim to spread their curse throughout the Dwendalian Empire. They believe their mission will make humanoids evolve into a stronger, more perfect people.

\subsection{Areas (D-K)}
\subsubsection{Deastok}
\begin{description}
\item[Population.]
10,090 (60\% humans, 26\% dwarves, 10\% halflings, 4\% other races)

\item[Government.]
Starosta Deter McGinneas is nothing more than a puppet who answers to Prime Arbiter Truscan in Rexxentrum.

\item[Defense.]
Deastok enjoys the protection of the Crownsguard. The city supplements its defenses with a small army of sellswords and private bodyguards.

\item[Commerce.]
Basic goods and services can be found alongside comfortable lodging. Some imported, rare goods can be sought by those with black market connections.

\item[Organizations.]
A single temple to Bahamut is the only official religious locale. The Truscan family and their allies hold tight to their power, while Myriad influence grows in secret.
\end{description}
\noindent
Travelers in the Truscan Vale find the city of Deastok at the southwestern edge of the Dwendalian Empire. The sister city of Kamordah to the north, Deastok is a borough that appears unusually affluent for its size, largely because the Truscan family has called the city home for many generations. Fine cobblestone streets connect numerous modern buildings and charming country shops, while elegant entertainment venues and picturesque views make the city a quaint holiday location for many wealthy travelers.

\subparagraph{Government}
While Starosta McGinneas is the Crown-appointed leader of the city, it is no secret that the Truscan family is the real guiding force in Deastok. The starosta and all titled subordinates within the city were assigned at the behest of Prime Arbiter Sydnock Truscan, and it is widely assumed that anyone in power is somehow related to the Truscans by marriage or blood, or is deeply indebted to them.

\subparagraph{Crime}
The clean, well-maintained streets are surprisingly free of crime for a city so socially stratified. There's a good reason that precious few petty criminals exist in Deastok: they've all been recruited by the Myriad!

Under the guidance of local crime boss Corshad Dumon, this city has been a major Myriad hub for decades, with secret subterranean bases and pathways beneath the city that enable the organization to go about its business unimpeded. Those that attempt to blow the whistle on Myriad activities vanish or meet a terrible end, so most city folk turn a blind eye to organized crime.

\subparagraph{Geography}
Deastok is a tidy city bordering the Cyrengreen Forest, with a beautiful collection of quaint homes, fine goods, and affluent citizens. The central metropolitan portion of the city is divided into three areas. Most of the residential homes are located in the Garden Grounds, framed by curling parks and greenery. The Bursar Plaza marks the industrial district where shops, markets, and busking spots are carefully arranged for both commerce and atmosphere. The perimeter of the city is known as the Jeweled Gates, where the Crownsguard barracks and small military reserve are located. The soldiers guard the city and take shifts atop the Cyrios Mountain watchtower.

Most of the city is built atop a large, leveled hill, surrounded by farms that continue for hundreds of acres on the western, northern, and eastern sides. At the southern end of the city, where the hills slope down toward the Cyrengreen Forest, a number of lumber mills bustling with gnomish technology harvest timber from the forest with impressive efficiency.

\begin{DndSidebar}[float=!b]{The Darrington Brigade}
Following his incredible journey as a member of the legendary adventuring group Vox Machina, famed artificer and writer Taryon Darrington returned to his home of Deastok and founded the Darrington Brigade, a nonprofit adventuring group modeled after Vox Machina. Under Taryon and his husband Lawrence Darrington, the brigade has gained notoriety in the past twenty years, both for the deeds accomplished and the series of mostly-non-fiction books published by Taryon himself.



The Darrington Brigade welcomes eager wanderers and adventurers with open arms, and they sponsor a number of small adventuring parties across Western Wynandir and beyond. Even so, Taryon can often be found helming an adventure or two himself, with his trusty automaton Doty and Lawrence by his side.



\end{DndSidebar}

\subparagraph{Deastok Adventures}
Elements of intrigue, noble finery, and the Myriad define adventures in Deastok. Characters might have adventures centered around grand parties, social gatherings, and whispered rumors.

\subparagraph{Merry Spouses of Deastok (Mid Level)}
A member of the nobility in Deastok is murdered, and the characters are hired to find the killer. During the investigation, more nobles and prominent members of society are slain. It seems the victims cannot be connected, and each murder employs a different grisly method. The truth is that a group of incubi and succubi married into the Truscan family over the past decade. These fiends are committing murders for the thrill of it.

\subparagraph{Kryn-filtration (Any Level)}
The Kryn Dynasty wants to infiltrate the nobility of Deastok to manipulate and steal information from wealthy Dwendalian citizens. The Kryn could hire the characters to accomplish this, or the characters could uncover the Kryn plot and try to find the spies. In either case, the Myriad's involvement in the city adds another dimension to the conflict between the Dwendalian Empire and Kryn Dynasty, making the characters wondering who—if anyone—they can trust.

\subsubsection{Egelin Quarry}
This centuries-old strip mine was the primary source of stone used in the construction of Zadash by the Julous Dominion, but it was largely abandoned when the platinum-rich Herathis Mine was discovered to the southeast. Now this gaping chasm is only worked by a small group of miners who still collect errant stone and rocks to sell to masons. Under the watchful gaze of the gnome prospector Lason de Weeler, the self-appointed foreman of the quarry, are a number of rock-breakers who toil in the mines in hopes of discovering some hidden cache of gems or precious ore.

\subsubsection{Felderwin}
\begin{description}
\item[Population.]
8,180 (58\% halflings, 21\% humans, 6\% dragonborn, 15\% other races)

\item[Government.]
The locally elected starosta, Theanor Whisperthin, works with the Crown-appointed lawmaster to organize and distribute the harvests.

\item[Defense.]
A sizable regiment of Crownsguard is now bolstered by a patrolling legion of Righteous Brand soldiers, while guildmasters keep personal bodyguards.

\item[Commerce.]
Basic supplies and general goods are produced locally, and humble rooms and amenities are available to travelers.

\item[Organizations.]
Shrines to the Dawn Father bless the outer fields, while most worship takes place in homesteads within the city. The farming families and guilds aren't independent, but are operated directly by the Crown.
\end{description}
\noindent
Felderwin is the heart of halfling society in Wynandir. The settlement itself is surrounded by fertile farmland—the sprawling Felderwin Tillage that has been cultivated by the halfling community for over seven centuries. In the wake of a sudden and vicious attack from the Kryn Dynasty, the Crownsguard presence has recently increased, and paranoia has gripped the citizens of Felderwin. Taverns and shops are filled with hushed whispers about the fate of the empire.

\subparagraph{A Tilling Tradition}
This area was originally settled by the halfling Felder clan many centuries ago, and the community has always valued the simple, grounded life of a farmer. The exceedingly fertile soil surrounding Felderwin is known throughout the empire, and the best farmers and students of agriculture in the land have settled here. As such, the agricultural skills taught in Felderwin are second to none.

Some family bloodlines have tilled the earth for as long as Felderwin has existed and fiercely guard their plots against any incursion—be it from other farmers, local goblins, or invading dark elves. Nomadic goblin camps moving through the nearby Cyrios Mountains often encroach in the dark of night during harvest time.

\subparagraph{Government}
Without Felderwin, the Dwendalian Empire—and the imperial military—would surely starve. Recognizing the importance of Felderwin's fertile fields, the Crown has allowed the settlement's starosta to be selected by the local community instead of appointing a leader from Rexxentrum. When the starosta dies or is deposed, the community votes on their replacement. The current starosta, Theanor Whisperthin, has served the community for over thirty years and does his best to maintain order. However, Theanor is under pressure by his assigned overseer on the Cerberus Assembly, Marquis Jenna Iresor, to expand the settlement's contributions to the war effort—or be replaced.

\subparagraph{Crime}
Felderwin is not a haven for unlawful activity, but the importance of the local crops have prompted the empire to assign a large number of Crownsguard to prevent any danger to agricultural production. The crime that does exist within Felderwin tends to involve small scale larceny, burglary, and the occasional attempt to forge official documents to avoid taxes and tithes to the Crown.

\subparagraph{Geography}
This rustic settlement is home to the beating heart of the Marrow Valley's farming community. Felderwin is built around three small hubs, each surrounding a communal gathering site called a Hall, which give the settlement a roughly triangular shape. Each section is filled with intersecting dirt roads, large barns, storehouses, lively taverns, and livestock pens that give the streets their unique, shifting smell of fresh grass, sweet pollen, and animal dung. The Green Hall calls forth trade meetings and day celebrations, while the Sunset Hall acts as a massive tavern for evening festivities and diplomatic events. The Lake Hall is built at the edge of the Scymir River and helms the small fishing community within the city.

\subparagraph{Felderwin Adventures}
Felderwin is an excellent location for new adventurers. The war has visited the settlement, and there are criminals and goblins aplenty for the characters to defeat.

\subparagraph{Short Uprising (Any Level)}
The characters find themselves in Felderwin when Starosta Theanor gives in to the Cerberus Assembly and raises taxes. The people of Felderwin, still trying to recoup their losses from the Kryn attack, begin to riot, demanding that the starosta step down so they can elect a replacement. Both sides pressure the characters to join them in the conflict.

\subparagraph{A Price Paid in Blood (Low Level)}
A series of macabre murders have shaken the denizens of Felderwin. The watchmaster has heard rumors that a \textbf{blood hunter} recently arrived in the city and hires the characters to hunt down the possible murderer. When they finally find the blood hunter, she tries to convince them that the murders are the fault of a rakshasa that's secretly living within the settlement as the watchmaster. The characters must either help the blood hunter take down the watchmaster or capture the blood hunter to claim their reward.

\subsubsection{Felderwin Tillage}
The Tillage, as it's called by locals, is the expanse of fertile farmland surrounding Felderwin, west of the Ashkeeper Peaks. This sprawling network of farms cultivates most of the empire's crops and livestock used and consumed across Western Wynandir. Crops are cultivated and harvested in the Tillage, while most business and shipping is handled within Felderwin itself. Farming families frequently travel between their acreage and the settlement of Felderwin to sell their goods and distribute their tithes to the empire.

While most conflicts throughout these rural homesteads involve disputes over land, the open and remote nature of the farms makes them easy targets for burglary and theft. Stories of such dangers continue to grow as financial strife in the empire pushes working folk to extremes. Increasing mistrust of outsiders and even some neighbors has led to some families becoming quite hostile. For the most party, however, there are few communities known to be as welcoming and hospitable as those of the Felderwin Tillage.

A stealth unit of Kryn warriors recently sneaked into the Marrow Valley from below the earth and set fire to a number of fields before being repelled. With large plots in ashes, many farming families are trying to rebuild, while the leaders of Felderwin are seeking ways to better defend their populace.

\subparagraph{Felderwin Tillage Adventures}
The characters could aid or stop Kryn raids on Felderwin Tillage, or defend the farms against burglars and monsters.

\subparagraph{Kryn Spy (Low Level)}
The characters are hired by a farming family in the Tillage to get to the bottom of recent arsons that burned crops in Felderwin Tillage. The culprit behind the fires is difficult to find because they are a \textbf{doppelganger} working for the Kryn, masquerading as multiple absent members of the community. Moreover, once the characters ask too many questions, the doppelganger begins hunting them as well.

\subparagraph{A Friend in Need (Mid Level)}
Stumbling on a secret gathering of cultists of Tiamat preparing to sacrifice a woman to their evil goddess, the adventurers must stop the ritual and free the captive, who aids them in the battle. She introduces herself as Calianna, a half-elf sorcerer with a draconic lineage that covers half of her body in black scales. She breathlessly describes this terrible cabal, known as the Cult of the Caustic Heart, and requests the characters' aid in vanquishing it once and for all.

\subsubsection{Grimgolir}
\begin{description}
\item[Population.]
19,090 (81\% dwarves, 8\% humans, 6\% halflings, 5\% other races)

\item[Government.]
Underbaron Brock Dummaroff rules the city with a council of dwarf lords and ladies. The Crown-appointed starosta is effectively just another councilor.

\item[Defense.]
A fair Crownsguard presence is overshadowed by the impressive force of elite dwarf soldiers known as Monoliths that protect the city. Private Monolith warriors are hired to protect merchant families. Most civilians are trained in basic combat.

\item[Commerce.]
A majority of services can be found, as well as raw ore, materials, gems, and goods. Taverns and inns are plentiful.

\item[Organizations.]
Local merchant guilds and conglomerates control the flow of business. Temples to Moradin and Bahamut are prominent, followed by the other Dwendalian gods. All answer to the Crown, but it holds less political influence here than in other cities.
\end{description}
\noindent
The dwarven stronghold of Grimgolir sits at the core of an immense network of tunnels carved throughout the southern Dunrock Mountains, perfectly situated to watch over the eastern border with Xhorhas. Though the Dwendalian Empire and the Grimgol dwarves have clashed in the past, the empire now counts the disciplined armies of Grimgolir as one of its staunchest assets in the war against the Kryn Dynasty. The predominantly dwarven society of Grimgolir prizes fine metals, finer craftsmanship, and proven skill on the battlefield, and their culture drills these virtues into every young dwarf, human, and halfling born in their city. Since days long past, the people of Grimgolir have waged eternal war on the monstrous denizens of the lightless caverns, in order to mine the precious mineral veins that suffuse the Dunrock Mountains.

All soldiers from Grimgolir have stories of their first patrol along the Brokenveil Bluffs. Every new recruit's first station is atop the dizzying cliffs, as an introduction to the hostile terrain of Xhorhas and to the monstrosities that wander out of the marsh below. The Grimgolir reputation for martial excellence soon reached the ears of the Dwendalian military, and the Crown made overtures to the ruling elders of the Grimgol dwarves, inviting them to peacefully become citizens of the empire. The negotiations were nearly spoiled by the Crown's simultaneous courtship of the orcs of Bladegarden, but a concordance was eventually reached. The empire granted Grimgolir an extremely generous trade agreement and all the benefits of imperial citizenship, along with unprecedented autonomy for an imperial vassal state, in exchange for the ability to marshal the Grimgol warriors to fight for the empire. As part of the empire's unified eastern front against Xhorhas, Grimgolir has aided in the training of the soldiers of Bladegarden, as well as the Rockguard and Ashguard Garrisons.

\subparagraph{A Society of Service}
Life under the mountain has never been easy. Every bountiful cavern is home to a deadly beast, and ancient tunnels once used by the armies of the Betrayer Gods now are haunted by lingering evils. This has shaped the core value of Grimgol society: to succeed is to struggle. Worth in oneself is measurable by what one has overcome. This has built a philosophy that breeds generation after generation of dwarves seeking ways to challenge themselves and prove their capability. Everyone born within the city learns the basics of armed and unarmed combat as a child—often making outsiders, merchants, and travelers somewhat uneasy with the knowledge that any native Grimgoliran could take them in a fight.

Early combat training combined with a cultural reverence for wars of ages past creates immense pressure for native-born Grimgolirans to join the city-state's military: the Monoliths. This is especially true for the Grimgol dwarves, for most can trace their family lineage back to a famous warrior or two. However, other Grimgolirans find their calling as storytellers and artisans. Those unable to make a living as a warrior or as an artisan often find themselves as laborers in the service of the Bleeding Stone Syndicate, a massive business enterprise that controls the entire pipeline of ore excavation to weapon manufacture. The Bleeding Stone maintains mining facilities beneath the Dunrock Mountains, a network of railcart lines between the mines and Grimgolir, a cutting-edge blast furnace for ore smelting, and an army of smiths hammering ingots into blades, shields, and mail.

\subparagraph{Government}
The agreement between the Crown and Grimgolir is the only treaty that has ever guaranteed a city-state autonomous rule while under the empire's banner. While Grimgolir does indeed have an appointed starosta, currently Starosta Ichabod Ashelon, the treaty between Grimgolir and the empire relegates the starosta to little more than an advisor of the ruling Underbaron of Grimgolir. As such, Underbaron Brock Dummaroff oversees the city and its goals with his own council and appointed officials. This arrangement has caused a few frustrations on both sides, and with the war against the Kryn now in full swing, these disagreements continue to grow larger and more frequent.

\subparagraph{Crime}
Culturally, denizens of Grimgolir genuinely believe that crime is a waste and a dark mark on one's life's work. If someone struggles, others lift them up until they can walk on their own. As such, it is very rare for an individual raised in Grimgol society to turn to a life of crime unless absolutely necessary. There is some criminal activity within and around the city, but most of it is helmed by foreign, outside forces.

\subparagraph{Geography}
From outside, Grimgolir is an impressive series of battlements and crenelations built into the side of the Dunrock Mountains, with granite staircases and bridges adjoining them. Within the mountain, the city is divided into two massive, cavernous wards, one sitting at an angle above the other in a gradual grade.

The interior cavern of the top ward, known as the Cage Ward, resembles a giant, angled ribcage. The Cage Ward contains numerous levels of structures that house civilians, businesses, and guard posts for Monoliths to watch over the area. Farther down the sloping tunnels lies the Heart Ward, a gargantuan, spherical cavern with dozens of elevated walkways crossing at multiple heights, connecting many of the hundreds of open tunnels. The Heart Ward is the hub of all mining and gathering operations throughout the mountain, surrounded by Monolith barracks embedded in the outer edges of the ward. Many barracks adjoin wide tunnels that breach the surface at the base of the mountain, with the exit point heavily guarded and shrouded in illusory stone to prevent discovery.

\subparagraph{Grimgolir Adventures}
With monsters aplenty and war at its doorstep, Grimgolir has no shortage of adventures.

\subparagraph{Mimics in the Mine (Low Level)}
Monsters are killing Grimgolir's miners. The typically tough and unbeatable Grimgolir soldiers can't seem to find the creatures. Looking for a fresh perspective, they hire the characters to hunt the monsters in the abandoned mine while operations are shutdown. The monsters are mimics that turn into pickaxes, railcars, and other mining equipment. They strike at the characters with hit-and-run tactics, retreating into darkness and using their unique powers to blend in with the environment.

\subparagraph{Rise of the Worms (Epic Level)}
While the characters are in Grimgolir, several purple worms tunnel into the city, causing mass destruction. After the worms are killed or driven off, Kryn forces come through the tunnels created by the monstrosities, razing strategic military outposts while the city is already in disarray. The Kryn Dynasty knows how critical the city is to training Dwendalian troops, so they mobilized the worms and followed behind to deal Grimgolir a decisive blow.

\subsubsection{Herethis Mines}
Platinum is as rare as it is beautiful, so when a very bountiful vein of the precious metal was discovered east of Deastok within the bounds of the Julous Dominion, the greedy lust for such a financial wellspring finally pushed the empire to break their agreements and conquer the dominion. Within the Crown's control, a vast mine has since expanded to extract the extremely valuable metal. However, the mine's proximity to the tenacious ankheg population that infests the Cyrengreen Forest continues to be a challenge. In addition, the Myriad has slowly infiltrated the operations of the mines over the past fifty years, regularly skimming from the platinum haul to fund its endeavors.

\subsubsection{Hupperdook}
\begin{description}
\item[Population.]
12,090 (76\% gnomes, 10\% dwarves, 8\% humans, 6\% other races)

\item[Government.]
Starosta Emma Zezbok, who has ruled for two centuries, works with elected officials and the major houses of industry.

\item[Defense.]
A small number of Crownsguard are backed by a larger civilian militia of volunteer soldiers. Many war machines and engines also defend the city.

\item[Commerce.]
A strong industry around metals, mining, tinkering, and entertainment makes most basic supplies and some rarer goods available. Upper Hupperdook has many taverns, inns, and places of revelry.

\item[Organizations.]
Industrial families are heavily invested in the city and its business. Temples to empire-approved gods, especially Moradin, are easy to find.
\end{description}
\noindent
The gnomish community of Hupperdook is built along the eastern base of the Silberquel Ridge, marked by the ever-rising columns of dark smoke and steam. Predating the empire, Hupperdook remained autonomous until the Crown absorbed the community through diplomatic agreements over two centuries ago. This smoke-stacked city is now the heart of industry in the empire, manufacturing technological marvels and large-scale weapons of war for the Dwendalian Empire.

The gnomish culture within Hupperdook prizes hard work and hard play. Family trades vary from coal mining and clockwork development to practical enchantment and the brewing of zesty libations. Within the bounds of the city, a general air of upbeat anticipation frames the day-to-day activities of residents as they bound their way toward dusk.

\subparagraph{Carpe Noctis}
For as long as the people of Hupperdook have lived on the mountain, they've been pushing to better their respective trades. The ten-hour work days are long and sometimes arduous, as sparks fly and clangs of metal ring through the city, but when dusk comes, the whistles blow and the carousing begins. Nearly every evening is a wild party celebrating the end of a productive day. Music, fireworks, dancing, and lots and lots of alcohol have given Hupperdook a reputation as one of the rowdiest locales within the empire. Such a lifestyle can be taxing on folk not used to such grueling labor and formidable merriment, but as the saying goes, "ya come here for an early grave worth earning."

The traditional nighttime drink is called fuse, a mixture of strong liquor and an energy-boosting herbal concoction that is considered by the hardest workers as a necessity to properly embrace an evening of revelry, though some visitors complain of the rather unpleasant aftermath of imbibing the powerful substance. A jet-black coffee with the consistency (and some say, the taste) of lantern oil is customary the following morning.

\subparagraph{The Core of Warfare}
The Crown understands the benefit of keeping Hupperdook and its citizens happy through relative autonomy, as no other city can develop and produce weapons like the ones found here. A sizable portion of Hupperdook's industry is allocated to the production and development of defense weaponry, siege engines, and firearms. Finished implements of war are shipped across the valley to Bladegarden or Rockguard Garrison, with a few left stockpiled within the city.

\subparagraph{Government}
As per the accord enabling it to become part of the Dwendalian Empire, Hupperdook is one of the only imperial cities allowed to choose its own leadership. Starosta Emma Zezbok is a bombastic gnome who negotiated the absorption of the city into the empire and has held her seat for two centuries. Emma has sharp mind for problem solving and a love of debauchery—and cares little for politics outside the city walls or for the short-lived "Kings of Men," so long as they leave her and her people alone.

Due to the independent nature of the city, the Crownsguard presence is relatively small, and the volunteers of the Citizens' Watch pick up the slack. This has led to some tension between the two law enforcement factions, but many of the Crownsguard forgive any slights when the evening drinks begin to flow.

\subparagraph{Crime}
The diligent days are almost completely free of crime as every nook and cranny of the city buzzes with business, but the nights lend themselves to lackluster oversight and intoxicated Crownsguard. Lawful folk tend to travel in pairs at night, lest their drunkenness be exploited by a cutpurse. While most criminals native to Hupperdook disdain pickpocketing people on the street, burglary is seen as a noble and skilled pursuit among thieves.

\subparagraph{Geography}
Hupperdook is divided into four distinct areas across two levels built against the mountain base.

Lower Hupperdook sits at ground level and is divided into two sections. The soot-stained, dusty Ironlot is where the majority of mining, metalwork, and small-scale production takes place amid residential neighborhoods peppered with a dozen ash-piping smoke stacks. The larger war machines and siege engines are constructed in the Assembly Yard, while more volatile projects are tested within the pock-marked turf known as the Craterfield.

The more charming, residential region of Upper Hupperdook sits about two hundred feet above the ground atop a platform of leveled rock. This level of the city has two regions: Silver Falls Grove, where the massive Silver Falls descend from atop the mountainside and spill into a small lake surrounded by residential living towers, and the Idleworks Shelf, a recreational region of sprawling commerce, colorful markets, and social gatherings.

\subparagraph{Hupperdook Adventures}
The nearby Silberquel Ridge is a ready source of dangerous beasts and unnatural dangers, while the evening celebrations of the city conceal all sorts of shenanigans.

\subparagraph{Life Is a Drag (Any Level)}
The annual Iron Runway celebration is seeking entertainers to wow the nobility with show-stopping, gender-bending performances and delightfully bawdy humor. The characters can choose to participate for an impressive top prize, or be hired to guard the event, which is known to get a little rowdy.

\subparagraph{A Hungry Problem (Low Level)}
An illegally transported \textbf{rust monster} escaped from its cage in the night, nested within the Silberquel Mines, and gave birth to a small brood. The hungry creatures now threaten to consume the veins of ore needed for the empire's war effort, unless the characters manage to slay or capture them.

\subsubsection{Kamordah}
\begin{description}
\item[Population.]
7,440 (58\% humans, 23\% halflings, 13\% dwarves, 6\% other races)

\item[Government.]
Starosta Brennan Keef oversees his network of appointees, but his true loyalty is to the Truscan family.

\item[Defense.]
Crownsguard and hired mercenaries protect independent business interests.

\item[Commerce.]
Basic supplies are available, but imported goods are more difficult to obtain. Rare plants, herbs, and wines support a profitable and well-guarded local industry.

\item[Organizations.]
A few small, family-run temples are open to travelers, but the local businesses are fearful of outsiders.
\end{description}
\noindent
Kamordah was built within the Bromkiln Hills by a collection of zealous followers of Desirat the Twilight Phoenix (described in \textit{chapter 1}). They were eventually driven out by the Julous Dominion, who claimed the city and attempted to harness the intense geothermal activity in the area. When the dominion was conquered by the Dwendalian Empire, it was discovered that this phenomenon infused the local earth with minerals that support a number of rare plants, ones that normally would not survive in Wildemount. This has led to a very profitable, carefully guarded industry of farming exotic herbs and fauna, as well as local vineyards that produce exquisite wines with flavors not found elsewhere in Exandria. Unfortunately, the ground water here is undrinkable, forcing denizens of the city to survive off captured rainfall and fresh water imported from outside the city.

\subparagraph{A Fire Below}
While the majority of Western Wynandir and the Marrow Valley is filled with rolling fields of green grasses and tall, vibrant trees, the Bromkiln Hills that surround the city of Kamordah are gray and rocky, unable to sustain the vegetation. Patches of gray, bubbling mud form between the hills, where natural vents perpetually release lazy plumes of sulfurous gases that give the region its unmistakable smell. Mineral shelves surrounding areas of intense seismic or thermal activity sometimes produce rings of vibrant colors, and many locals study these minerals for their medicinal and arcane properties. Others worry that the source of this activity is a portent of terrible things to come, whether it be an imminent eruption or some slumbering elemental power.

\subparagraph{Government}
Starosta Brennan Keef is the Crown-appointed governor of Kamordah. Though his close ties to the Truscan family grant him good standing with the Crown and local traders in the Truscan Vale, locals grouse that he puts Truscan interests before those of his own people.

Coinmaster Dela Folger has a long-standing feud with Exchequer Aethia Drooze. Folger, a staunch supporter of Kamordan wineries and their local businesses, resists Drooze's attempts to squeeze every last copper out of Kamordah and into the Crown's coffers.

\subparagraph{Crime}
The businesses in Kamordah carefully guard their rare goods. Crownsguard mingle with hired mercenaries to keep a watchful eye on business deals, making the few criminals within the city extremely cautious.

\subparagraph{Geography}
Kamordah is a strange settlement built from pale stone mined from nearby Mount Mentiri and petrified wood from the surrounding forest. Constructed over the prismatic mud of the Bromkiln Hills, it resembles a collection of pearls against a smear of colored paints. The settlement is loosely broken up into two regions based on industry. The Rainbow Vineyards wind along the northern end, twisting up the base of the mountain paths, while the massive central Mudfields harbor greenhouses and research facilities alongside the scattered homes that comprise the remainder of the city.

\begin{DndSidebar}[float=!b]{Wines of Kamordah}
The natural phenomenon that churns the mud and earth surrounding Kamordah has helped its people cultivate a valuable collection of vineyards. Wines produced by these families are some of the best in Exandria, often presented as extravagant gifts or a display of wealth. These bottles fetch a healthy sum of gold, with some of the finest listed below, though the prices vary depending on the market.



\textbf{Stassman's Thistle Branch Dark Blood Red Wine:} 85 gp per bottle.



\textbf{Wagner \& Waltz Emerald Embrace Green Wine:} 65 gp per bottle.



\textbf{Stassman's Truscan Sunrise White Wine:} 45 gp per bottle.



\textbf{Lionett Vineyard Purple Cinder Red Wine:} 50 gp per bottle.



\textbf{Errenath Family Private Wish Red Wine:} 75 gp per bottle.



\end{DndSidebar}

\subparagraph{Kamordah Adventures}
Adventurers in Kamordah might become embroiled in squabbles over local goods, politics, or speculation concerning the source of Kamordah's geothermal activity.

\subparagraph{Broken Earth (Low Level)}
While the characters are in Kamordah, an earthquake topples buildings and releases lava from beneath the earth. The broken earth also opens the way for magma mephits that cause mischief around the settlement. If the characters manage to save themselves from the chaos, there are plenty of other citizens in need of rescue.

\subparagraph{Twilight Phoenix (Epic Level)}
The truth of Kamordah's geothermal activity is far more sinister and fantastical than many know. The ancient and mad twilight phoenix Desirat is chained beneath Mount Mentiri. A cabal of archmages who worship Asmodeus seek to free the Lord of the Nine Hells' steed, secretly gathering components for a ritual to sunder the mountain and unleash the crazed entity to incinerate the countryside. If the characters cannot stop the ritual, they must combat the twisted, primal phoenix herself.

\subsection{Areas (L-Z)}
\subsubsection{Lebenda Swamp}
The sprawling marsh known as the Labenda Swamp encompasses the valley along the southern side of the Silberquel Ridge and is fed by frequent rainfall. Surrounding the Ounterloch, this immense swampland is humid for much of the year, making travel uncomfortable. The incessant hum of buzzing insects mingles with the musky stench of decaying plants and stagnant water, and thick fog obscures much of the muck-laden region. The marsh is home to trolls, carnivorous plants, and a host of other deadly creatures. Only the most experienced travelers dare wander too deep into the soggy reaches of the inner swamp, where it is said that rare plants and ancient ruins hide behind the mists.

\subsubsection{Mount Mentiri}
Mount Mentiri is a sky-piercing crimson monolith and an ever-present visual landmark within the Truscan Vale. Thin wisps of steam rise from numerous vents scattered throughout the massive mountain. The nearby Bromkiln Hills and the northern lands around Pride's Call have been altered by the elemental power that radiates from the rock beneath Mentiri.

The source of this subterranean heat is a mystery to the people of the empire. Scholars have proposed many theories, suggesting that Mentiri may be a young volcano—or that it sits atop a rift to the Elemental Plane of Fire.

None of their theories are correct. The truth is that Desirat the Twilight Phoenix, prized steed of Asmodeus, is imprisoned deep within the curling caverns beneath the mountain, and her prison is surrounded by the fire giant city of Vuthos. The presence of Desirat has affected the people of Kamordah since the city was founded, and her messages call to the weak-minded in dreams, urging them to delve beneath the mountains to worship her as a god.

\subsubsection{Ounterloch}
The mossy green waters of the Ounterloch lie brackish and foul in the center of the Labenda Swamp. Home to swamp beasts and malevolent forces drawn to the mystic secrets buried beneath the algae and stones, this lake is a source of myths for folk all across the Marrow Valley. Many speak of the temple that once sat at its northern shore before the quaking of the Calamity sent it sinking down into the dark water. Others tell tales of the Marsh Siren who calls to people in their dreams and draws them into the swamp with her song. These unfortunate souls are lost to the marsh, or they drown in her embrace beneath the waters. Most who discuss the Ounterloch, however, just speak of how miserable it is.

\subsubsection{Talonstadt}
\begin{description}
\item[Population.]
1,810 (82\% dragonborn, 11\% humans, 4\% halflings, 3\% other races)

\item[Government.]
Talonstadt's refugees are organized by community leaders who manage to work together, despite their bickering.

\item[Defense.]
A small but functional band of volunteer fighters keeps watch alongside a minimal Crownsguard presence.

\item[Commerce.]
General goods and arms are available, if somewhat scarce beyond what the populace requires. Trade is minimal, and lodging is hard to find.

\item[Organizations.]
There are makeshift shrines to Bahamut in public gathering places. The Scars of Scale and Tooth have representatives present to help protect the interests of ravenite refugees.
\end{description}
\noindent
The fall of Draconia at the hands of the deadly chromatic dragon sect known as the Chroma Conclave sent the region into chaos. The ravenite resistance rose up against the surviving draconbloods in a bloody rebellion, while the dangerous creatures that stalked the Dreemoth Ravine closed in when the arcane defenses of Draconia fell. While many remained to rebuild from the ruins of the fallen city, dozens of ravenite refugees and survivors fled across the dangerous lands of Xhorhas, climbing over the mountains into the Marrow Valley. The Crown wasn't eager to offer clemency to a number of foreign revolutionaries who overthrew an allied nation, but the Cerberus Assembly managed the development of a ramshackle shantytown for the dragonborn at the base of the Ashkeeper Peaks.

Although it is part of the Dwendalian Empire, Talonstadt receives very little support from the Crown. Instead, the Cerberus Assembly pays meager wages to the refugees while promising that their hard work will one day allow them to earn imperial citizenship and a better quality of life for their families. With tensions between the refugees and the empire escalating, and a Crownsguard that consistently turns a blind eye to crimes within the community, very few residents believe these empty promises.

Amid the chaos, a handful of Draconian sympathizers within the refugees actively undermine the ravenites' community leaders.

\subparagraph{Talonstadt Adventures}
The people of Talonstadt have great need of adventurers to protect them from crime, monsters, and each other. The community can't offer much in the way of monetary rewards, but they can provide forgotten stories and describe (relatively) safe routes through Xhorhas to the characters.

\subparagraph{Chroma Council (Low Level)}
A secret group of Draconian sympathizers, who call themselves the Chroma Council, tempts the young adults of Talonstadt with promises of wealth and a better life if they adhere to the commandments of Tiamat. Their recruits disappear into the Ashkeeper Peaks to meet with the cult leaders, who plan to raid passing caravans for money and supplies. The characters are hired by the recruits' families to steal them back from this seditious cult.

\subsubsection{Trostenwald}
\begin{description}
\item[Population.]
8,900 (66\% humans, 13\% halflings, 8\% half-elves, 13\% other races)

\item[Government.]
Starosta Fryda Langer works alongside resident nobles and small local guilds to maintain the affairs of the city.

\item[Defense.]
A respectable number of Crownsguard protect Trostenwald, occasionally bolstering their defenses with hired mercenaries.

\item[Commerce.]
Most basic supplies and services are available, with occasional rare goods from the Menagerie Coast. The city's brewing industry is prominent and celebrated.

\item[Organizations.]
Small temples and shrines to empire-approved gods are present. The brewing families compete for dominance while Myriad activity simmers in secret.
\end{description}
\noindent
The settlement of Trostenwald lies at the southern outskirts of the Dwendalian Empire, on the banks of the Ustaloch. The seeds of civilization were planted here nearly two centuries ago when the fair weather and stretching grain fields gave rise to a booming brewing economy. Three major families of brewers and alewives earned fame throughout the empire by producing some of the best and most unique ales Dwendalians had ever tasted.

\subparagraph{A Finer, Sweeter Grain}
Ordering a "trost" has become common terminology in many taverns that serve the unique, sweeter ales of the settlement's breweries. Many attribute the distinctive flavor to the waters of the nearby Ustaloch, while others attribute the flavor to an enchantment on the nearby soils hearkening back to the Age of Arcanum. Regardless, the ale has become the flagship export for the settlement, causing tension between competing families. The reigning three breweries fight for dominance over farmland and promotion of their products, while occasionally allying to squash out any upstarts who wish to encroach on the already crowded market. The known family brands that produce Trostenwald ale are Husseldorf, Bambauch, and von Brandt.

\subparagraph{Government}
The city is headed by Starosta Fryda Langer, daughter of the last starosta who fell to illness two years ago. Her family has strong ties to the Bambauch family, leading some to mutter of favoritism, though they have no power to challenge her. Starosta Langer enjoys her position, often making public shows of her power by holding small celebrations and contests in her honor. Those she did not favor before taking office are rumored to have disappeared or been driven into poverty.

\subparagraph{Crime}
Crime in Trostenwald usually involves petty thefts of shipments and burglaries rather than any major conflicts. The Crownsguard have their pockets padded by the three brewing families to pay special mind to their wares. The Myriad has agents in Trostenwald, but the settlement is too sleepy to become an effective center of criminal activity. Myriad operatives do their part to push contraband smuggled in from the Clovis Concord north toward Zadash and Rexxentrum.

\subparagraph{Geography}
Trostenwald is a rural borough of simple homes, competing breweries, and fishing huts that form a kidney-shaped layout, wrapping around the western end of the Ustaloch. Travelers from the south consider Trostenwald a simple and folksy welcome to the empire.

The city is broken into three main sections. The Northward borders the wind-kissed grain fields, and houses the privately owned breweries and other local industries. The Lochward forms the eastern portion of the settlement, comprised of residential neighborhoods and lakeside fishing communities. The southwest region is the Hillward, where the affluent and noble-blooded live among the guild halls, easily protected by the nearby Crownsguard barracks.

\subparagraph{Trostenwald Adventures}
Characters in Trostenwald could get caught up in adventures with the three brewing families or the Myriad. They might suddenly get called to action while deep in their cups.

\subparagraph{Grain Hijackers (Low Level)}
The characters are hired by one of the brewing families to find out who tore into their silo and stole their grain. The characters might suspect one of the other brewing families, but that notion is quickly dashed when a second family's silo is similarly vandalized. The true culprits behind these attacks are two owlbears that meet each night and devour grain from a new silo as part of a courtship ritual.

\subparagraph{Find the Formula (Any Level)}
The characters are hired by an anonymous entity to uncover the secrets of Trostenwald's ale. This means investigating all the rumors, infiltrating the brewing families to find recipes, and examining all the ingredients. The brewing families keep their secrets under lock and key. The plot thickens when the characters' anonymous patron wants to make sure the recipe doesn't get out, so they attempt to kill the characters when the mission is complete. This patron could be a member of the Myriad, a family from Shadycreek Run, or a rival brewery in another settlement.

\subsubsection{Ustaloch}
This beautiful lake is the largest body of fresh water in the southern reaches of the empire and the largest source of fish in the Marrow Valley. The city of Trostenwald rests firmly against the western banks, tending to its waters and the bounties they provide. Clear waters aid fisherman in their trade, and the purity of the lake water is thought to be a key ingredient in Trostenwald's sweet ales. Crownsguard patrol the shores of the lake at all hours to prevent anyone from tampering with the ecosystem that the locals depend on.

\subparagraph{Myths and Murmurs}
The Ustaloch is not without dangerous creatures that stalk the deeper waters, and it has developed a number of myths, the most popular of which stems from the larger of its two islands, known as Crooked Stone. This island contains a small grove where the abandoned hovel of a feared witch still stands in disrepair. Many youth paddle out to the island to conduct rites of passage or pranks, though others warn that the witch still returns to her home for the colder seasons.

\subparagraph{Ustaloch Adventures}
From fishing for food to hunting a witch, there are adventures aplenty for the characters to find in and around Ustaloch.

\subparagraph{Witch's Amulet (Mid Level)}
Something in Ustaloch is turning the fish and crabs in the lake into skeletal abominations that attack boats and people near the shore. The characters are hired to uncover the source of this curse, which happens to be an amulet from the witch's island grove that a group of children discovered and threw in the water. If the characters can find and retrieve the amulet, the witch (a \textbf{night hag}) might soon come looking for it.

To create the skeletal creatures in this adventure, take any aquatic beast and make the following changes to it:


\begin{itemize}
\item The creature is undead and doesn't require air, food, drink, or sleep.
\item It has immunity to poison damage and exhaustion, and it can't be poisoned.
\item It has vulnerability to bludgeoning damage.
\end{itemize}

\subsubsection{Vol'antim}
\begin{description}
\item[Population.]
3,890 (96\% aarakocra, 4\% other races)

\item[Government.]
Vol'antim is ruled by a council of twelve elders known as the Horizon.

\item[Defense.]
The settlement's perilous surroundings make it easily defensible, and trained aarakocra warriors vigilantly protect it from above and within.

\item[Commerce.]
General arms, simple supplies, and natural goods can be purchased or traded here. Lodging for guests is sparse.

\item[Organizations.]
Dual temples to both Melora and Erathis mark the cliff face and interior caverns. No major organizations make their home here.
\end{description}
\noindent
The hidden cliffside settlement of Vol'antim is nestled deep within the heart of the Cyrios Mountains. Tucked beneath an unreachable mountain precipice that soars over a thousand feet in the air above jagged rocks and brush, the community is home to aarakocra who live in relative seclusion. Few ever see these aarakocra, as they keep to the Cyrios Mountains and watch intently for uninvited guests.

The denizens of Vol'antim are proud of their isolationism and their adherence to the ancient ways of Melora the Wild Mother. For generations, they have carved their homes into the mountain cliffs. There are over a hundred caves adorned with colorful banners that mark the entrances to tunnels. These tunnels twist into a vast cavern, which houses the town proper.

\subparagraph{Aerial Enemies}
The skies above the Cyrios Mountains are a battleground between the aarakocra of Vol'antim and their ancient enemies, the gorefeather harpies. Despite the birdfolk's best efforts, the gorefeather clan has only grown larger over time. Some fear that their matron has forged a pact with an evil power that has granted her people a stronger brood with each passing generation.

\subparagraph{Mythical Guardians}
While the more learned people of Wildemount know of and understand the nature of the winged citizens of Vol'antim, the common folk often only catch occasional glimpses of them in the sky, and few people can claim to have met one. A handful of rural communities in the Truscan Vale refer to the flying creatures as their protectors and leave them small offerings before making dangerous treks in the region. The rare arrival of an aarakocra in one of these communities is a major event for the locals, and they shower the winged herald with gifts, much to the aarakocra's embarrassment.

\subparagraph{Government}
The community is led by a dozen elders called the Horizon. Each elder chooses their successor on their deathbed; if no successor is named before the elder expires, one is chosen by the remaining Horizon members. When there is disagreement within the Horizon, the eldest member makes the final decision and all others respectfully follow. The community's laws are structured around the tenets of Melora and the respect that she has for her lover, Erathis the Law Bearer. Their foremost creed is that natural life is sacred within the ways of nature and must be protected.

\subparagraph{Crime}
The culture of Vol'antim eschews criminal behavior and acting outside the interests of the community. However, there are tales of past aarakocra who fell from the grace of Melora and followed a selfish path, and these figures were often punished with exile. Visiting outsiders have also been known to bring criminal activity to the city, contributing to the mistrust of those who come to the cliffs unbidden.

\subparagraph{Geography}
Standing high above a rocky valley, the mountain spire curls at its peak, forming a protective lip over the cliff face where the numerous cave entrances to the city are visible. Only accessible via flight, each entrance bears a small shelf or lip for arrival, often bearing banners or thatched decor to signify what is found within. Beyond these caves lies a network of small tunnels that open into a shared interior cavern, in the middle of which a one-hundred-foot-tall stalagmite merges with a matching stalactite to form a great pillar riddled with caves that connect to homesteads in the cavern walls via bridges and ropes.

\subparagraph{Vol'antim Adventures}
If the characters earn the respect of the aarakocra of Vol'antim, they could be asked to help keep the Cyrios Mountains free from evil. If the characters instead earn their ire, they would do well to steer clear of the mountains.

\subparagraph{Break a Few Eggs (Low Level)}
The characters are asked by the \textbf{aarakocra} of Vol'antim to help destroy a nest of gorefeather harpies. During the quest, the characters uncover a clutch of harpy eggs, which the aarakocra believe should be destroyed, since legends say that each generation of harpies is stronger than the last. If the characters destroy the eggs, they will drive the harpies to greater acts of destruction. But they refuse, how might the aarakocra react?

\subsubsection{Wuyun Gate}
The southernmost boundary of the Dwendalian Empire, the Wuyun Gate rests within the Wuyun Gorge where the Ashkeeper Peaks meet the Cyrios Mountains. This imposing gate stands roughly twenty feet tall and fifteen feet wide, opening when passage is given to travelers seeking to enter or exit the empire. A wall of stone twenty-five feet tall continues from each side of the gate, stretching to meet the nearby mountains. These walls are topped by patrolling Crownsguard whose weapons are always at the ready. Originally built as a military fortification to prevent the Clovis Concord from assaulting the empire, it now acts as a waypoint on the trade route, and occasionally as a Righteous Brand outpost when dangerous creatures begin to wander too far from the nearby mountain ranges.

\subparagraph{Wuyun Gate Adventures}
The characters could help defend the Wuyun Gate outpost when monsters attack and hold the gate until reinforcements arrive, or they could reclaim the gate for the empire after it falls.

\subparagraph{Take Back the Gate (Mid Level)}
As the characters travel through the region, they need to pass through the gate. When they get to the outpost, they learn that monsters from the mountains (gnolls, hobgoblins, orcs, or trolls) murdered the Crownsguard and have claimed the gate as their own. From this secure position, they launch more raids into the Marrow Valley. Someone has to stop them and take back the gate.

\subsubsection{Zadash}
\begin{description}
\item[Population.]
89,210 (70\% humans, 11\% halflings, 9\% dwarves, 10\% other races)

\item[Government.]
Starosta Wyatt Fedar is the city's nominal leader, but nothing gets done without the assent of the lawmaster, other Crown-appointed officials, and the wealthy elite.

\item[Defense.]
An impressive host of Crownsguard is integrated throughout the city, while a sizable legion of Righteous Brand soldiers is stationed within the Signet Ramparts. Political figures and nobles also employ personal guards.

\item[Commerce.]
Nearly all goods and services are available in Zadash, including magic items. Illegal imports and contraband are found in underground circles.

\item[Organizations.]
Major temples to empire-approved gods are scattered throughout the city. The Cerberus Assembly holds great sway, as do the high nobility and the guildmasters. The Myriad, once thought banished, is still deeply entrenched in the city's businesses.
\end{description}
\noindent
Zadash was the capital of the Julous Dominion before it was conquered by the Dwendalian Empire nearly two centuries ago. Now the city is the central crossroads of the southern region of the empire and a smaller sister city to the capital of Rexxentrum. Zadash accommodates heavy traffic through its winding streets, serving noble diplomats, countless trade caravans, and rushing messengers from all across Western Wynandir and the Menagerie Coast. The city's ruling class is comprised of powerful merchant lords and ladies who oversee the guilds that control the markets throughout the city.

Zadash is a bastion of law, civilization, and military might in the Marrow Valley, whose small townships often escape imperial scrutiny. The city is home to the Hall of Erudition, an extension of the Soltryce Academy and foothold of power for the Cerberus Assembly. As a common waypoint for travelers to and from Rexxentrum, the city offers many forms of entertainment and recreation. It also boasts the largest military force in the southern end of the empire.

\subparagraph{Government}
The majority of day-to-day business and politics are overseen by Starosta Wyatt Fedar, Lawmaster Orentha Stonegrasp, and High-Richter Dolan Thyme, who must work closely to manage the large, sprawling city. However, it is common knowledge that the starosta is deep in the pockets of the city's wealthiest families, and is often manipulated by those powers without the people's best interests at heart. Community leaders are rightfully furious about the starosta's corruption, but Starosta Wyatt is finding it difficult to gain support from the capital, which is focused on the war with Xhorhas. The growing hum of unrest makes him increasingly nervous.

\subparagraph{Crime}
Criminal elements are part of the very fabric of the city. Independent swindlers make their way through the poorer districts, pickpocketing freely behind the backs of overworked Crownsguard. The Myriad has patsies in every industry and most political offices, too.

A crime boss known in some circles as the Gentleman runs a thieves' guild based in a network of underground tunnels called the Underworks. Crime is extremely profitable in Zadash, and has been fine-tuned for generations to be efficient, hard to trace, and largely invisible to the untrained eye.

\subparagraph{Geography}
Zadash has a mix of older Julous architecture and modern Dwendalian construction, giving the picturesque, triangular city a varied and mismatched look. Zadash is broken into five major districts: the Outersteads, the Pentamarket, the Innerstead Sprawl, the Tri-Spires, and the Signet Wall.

\subparagraph{The Outersteads}
Divided into the East and West Outersteads, this area is home to the poor and powerless of the city, a collection of slums and neighborhoods of hovels that cluster around and beyond the city walls. The majority of folks who live within these areas are humble crafters hoping to catch the eye of traders passing through, workers who handle the less desirable careers that keep the city functioning, or the destitute who are trying to regain control of their lives. Many sections are humble but welcoming, with family communities who work together to keep the peace when the Crownsguard are too sparse or too uncaring to watch over them. The regions beyond the walls are equal parts farming communities and tent cities. However, some areas of the Outersteads have become hotbeds for criminal activity, with locals disappearing often enough to cause discontent with the empire's indifference.

\subparagraph{The Pentamarket}
The Pentamarket is the bustling heart of commerce, a massive bazaar that rests at the junction of the city's five main roads. Here, all manner of guild shops and businesses provide both wares and lodging for droves of weary travelers. During many calendar holidays, this region of the city is flamboyantly decorated, and music and games fill the streets surrounding the crafters' temple, Hearth of the All-Hammer, and the massive gladiatorial arena known as the Victory Pit.

\subparagraph{The Innerstead Sprawl}
A labyrinth of houses owned by the working-class people of Zadash, the Sprawl bends between the other districts, acting as the boundary between them. Except for a few of the temples, such as the Valley Archive of the Cobalt Soul, the Platinum House, and the King's Hall, much of the architecture is simple and built for function.

\subparagraph{The Tri-Spires}
The Tri-Spires are known throughout the walls of Zadash as the neighborhood of the noble elite and the wealthy. Built on the ruins of the old Pious District after the fall of the Julous Dominion, the Tri-Spires now stand as the center of aristocracy and higher learning within the city—and, indeed, throughout the empire, second only to Rexxentrum. The region takes its name from the three large, uniquely constructed towers that lord over the skyline: the Zauberspire, the Constellation Bridge, and the Triumph Chime. The Zauberspire is the central tower of the Halls of Erudition, a beautiful ivory campus that is an extension of the Soltryce Academy. The Constellation Bridge looms over the Gilded Willows neighborhood as an extremely tall, thin watchtower that offers an incredible view of the entire city and the Marrow Valley beyond. Once a temple to Avandra, The Triumph Chime is now a squat, seven-story belfry outfitted with sleeping chambers, bars, and a raucous gambling hall.

\subparagraph{The Signet Wall}
The Signet Wall encloses the eastern end of the city and houses Zadash's reserve force of Righteous Brand soldiers. The Signet Ramparts are a collection of four simple stone and wood buildings that act as barracks and living quarters for all local Crownsguard. The clangs of training blades and blacksmiths' hammers fill the air for half the day, while drinking songs and bawdy rhymes fill the other half. Beyond those grounds lies the dark stone entrance to the impressive, eight-floor, subterranean prison known as the Umber Dungeon, where criminals from the streets of Zadash are held and punished, with the most dangerous imprisoned in the deepest cells. The Shrine of the Dawn, a temple of pale stone and stained glass dedicated to Pelor the Dawn Father, where both civilians and soldiers can worship, marks the center of the district. The eastern wall ends in the Marrow Keeps of the Righteous Brand, a massive stronghold surrounded by a half-dozen smaller fortifications that make up the southern bastion of the Crown's military power.

\begin{DndSidebar}[float=!b]{The Invulnerable Vagrant}
There are few locations in the empire where enchantments and arcane protections are readily available for sale, as the cost of such devices is prohibitive for the majority of the populace. Most recovered magic items find their way to black markets, are sold through auction houses, get taken by deft thieves, or wind up in the collection of some affluent family.



This makes the success of the Invulnerable Vagrant, one of the few magic item retailers in Wildemount, unique. The Vagrant, established in 814 PD, is owned and operated by a city-bred firbolg called Pumat Sol and his team of identical quadruplets—that is, his three arcane simulacra. Many visit this establishment simply to marvel at the curious atmosphere and equally curious staff.



Pumat Sol was trained by Headmaster Oremid Hass of the Halls of Erudition, but he tends to keep quiet about the specifics of his ties to the Cerberus Assembly.



\end{DndSidebar}

\subparagraph{Zadash Adventures}
As a large city full of travelers, criminals, and politicians, Zadash has plenty of opportunities for adventure.

\subparagraph{Send a Message (Mid Level)}
The characters are hired—or blackmailed—by the Myriad to convince Starosta Wyatt Fedar to let the Myriad quell unrest in the city in exchange for favors from the city leader. If the starosta accepts, the Myriad asks the characters to crack down on non-Myriad crime in the city. If the starosta refuses, the Myriad asks the characters to run amok, causing and inspiring as much chaos as possible until the starosta relents.

\subparagraph{Victory Breakout (High Level)}
The characters are in Zadash when several monsters are freed from their confines in the Victory Pit. These monsters—an \textbf{abominable yeti}, a \textbf{behir}, a \textbf{hydra}, a \textbf{tyrannosaurus rex}, and others, tear through the streets, leaving carnage in their wake. Did these monsters escape by themselves, or did someone else free them to distract the Crownsguard and characters while another crime took place?

\section{Zemni Fields}
Along the frost-dappled fields of the northern plains, the heart of the Dwendalian Empire is formed by a sprawl of dirt-packed roads, dense forests, and austere townships that fight back the dangerous creatures of the wilderness. Divided from the southern Marrow Valley by the mountains of the Silberquel Ridge, this region is known for its gray skies, fits of cold rain, and occasional snowfall, keeping the feet of the denizens muddy and their crops healthy. The Dunrock Mountains shield the people from the threats that loom beyond the boundaries of the king's rule, with the Rockguard Garrison to the east holding firm against the deadly armies of the Kryn Dynasty and the dangers that brew beyond the Ashkeeper Peaks. Home to the mighty imperial capital of Rexxentrum, these fields are the seat of power and commerce for the empire. Diplomats often find themselves traveling to and from the Zemni Fields, while traders and merchants from around the world traverse the long roads in hopes of making their fortunes in the capital city. Even so, the growing ferocity from beasts and creatures that share this region mingle with the whispers of a recovering criminal element and the escalating war with the Kryn to create a perpetual air of tension and anxiety.

\subsection{Areas (A-K)}
\subsubsection{Amber Crossroads}
The Amber Road and the Havenpath Road merge with the Cloakjaunt Path and the Glory Run road to form a great crossroads, where a small makeshift marketplace has sprung up. Mercenaries flock here to find work among traveling caravans, and merchants exchange goods in the absence of imperial oversight before heading into the nearby cities to set up shop. Wandering performers bring color and life to the crossroads, while charlatans of all walks of life take advantage of the fact that the market has no Crownsguard presence.

\subparagraph{Amber Crossroads Adventures}
The conjunction of trade and travel creates a location rife with mystery, danger, and the need for adventurers.

\subparagraph{Rob from the Rich (Low Level)}
Arriving just in time for a semiannual bazaar known as the Carnival of Coin, the characters begin to notice a strange, trance-like state that temporarily overtakes a number of wealthy merchants, who then sell off their most expensive wares for next to nothing. These merchants hire the characters to find the source of this enchantment, which turns out to be a wandering band of charlatans who seek to redistribute wealth to the poor of the Marrow Valley. The characters must decide whether they wish to turn in these thieves for rewards from the merchants, or let them go and risk becoming hunted alongside them.

\subparagraph{War Plans (Mid Level)}
The characters are hired by the Myriad or the Kryn Dynasty to intercept imperial documents as they travel from Rexxentrum to Zadash. The first courier takes the documents as far as the crossroads before handling them off to a second courier. Both couriers are knights of the Crownsguard mounted on warhorses. The documents are war plans bearing magical wards and the seal of the king.

\subsubsection{Blumenthal}
\begin{description}
\item[Population.]
3,850 (71\% humans, 12\% dwarves, 11\% elves, 6\% other races)

\item[Government.]
Starosta Judus Olessia is the appointed leader of Blumenthal, but the countless directives from Rexxentrum often bypass him entirely.

\item[Defense.]
A respectable cadre of Crownsguard and a small body of Righteous Brand soldiers keep the farmsteads of Blumenthal safe.

\item[Commerce.]
While food, produce, and basic supplies are available, little else can be found since most trade heads straight to the capital.

\item[Organizations.]
A temple to Pelor stands vigil over the heart of the farmland. The farming guilds all work under the strict expectations of the Crown.
\end{description}
\noindent
A sprawling collection of rural farmsteads connected by administration buildings, Blumenthal is Rexxentrum's main source of produce and meat. Numerous government-owned fields and orchards harvest cold-weather crops and livestock, and families chosen to tend to the land are paid by the Crown to gather their goods and deliver them to the gates of the capital. Some who till the land are frustrated that they don't own their plots outright, but with generational incentives for those who keep to the family plot like their ancestors, others see it as investment toward eventually living in greater comfort than the typical farmer.

As Rexxentrum's survival depends on Blumenthal's fields, the Crown has erected dozens of Crownsguard watchtowers at the corner farmers' plots to watch for thieves, scavengers, and monstrous predators. The people of Blumenthal can't help but feel the watchful eye of the Crown on them at all times, looming over their work and dampening their spirits. Yet even in such troubled times, Starosta Judus Olessia works directly with Argonomist June Wiston to improve the lives of the hard-working farmers who support the success of Blumenthal—and thus of the empire.

\subparagraph{Blumenthal Adventures}
Since Blumenthal's fields are critical to the survival of the people of Rexxentrum, adventurers might be hired by the empire to investigate disappearing or dying animals or crops, or to stop monster raids from the Silberquel Ridge. Adventurers who work for the Kryn Dynasty could be asked to destroy crops or poison goods headed for Rexxentrum.

\subparagraph{Bad Beef (Low Level)}
Tainted beef from three different cattle farms in Blumenthal has killed several people in Rexxentrum. Crownsguard investigators concluded that the meat was intentionally poisoned and hire the characters to find the culprits. Why hire outsiders? The Crownsguard in Blumenthal is on the list of suspects, since they escorted the meat to Rexxentrum. Of course, the list of suspects doesn't stop there. Did the farmers work together because they abhor living under imperial rule? Is there an agent of the Kryn within Blumenthal? Is someone using enchantment magic from afar to make farmhands poison beef? One thing is for sure: the culprits are willing to kill and won't come quietly.

\subsubsection{Bysaes Tyl}
\begin{description}
\item[Population.]
19,090 (83\% elves, 7\% humans, 10\% other races)

\item[Government.]
The Ring of Three reluctantly cooperates with the appointed starosta, Uther Prent, to govern the city. Exalted elven houses also work with the Ring of Three to preserve their culture under imperial rule.

\item[Defense.]
Both local and assigned Crownsguard stand ready to defend the city.

\item[Commerce.]
Nearly any equipment, goods, or services are available, though rarer supplies are generally kept from non-elven travelers.

\item[Organizations.]
While smaller temples dedicated to the empire-approved gods are present, shrines to Corellon, Melora, and Sehanine hide in plain sight. The beginnings of Myriad influence are taking hold.
\end{description}
\noindent
The rustic elven city of Bysaes Tyl is tucked away in the northern region of the Pearlbow Wilderness. Built synergistically within the tangled boreal forests, this hidden settlement of tall wooden arches and warm hearths sits below a massive network of treehouses and bridges that form the canopy and upper levels of the city. The streets are made of smooth stones set in packed earth, gently enchanted to warm the pathways and melt the frequent snow. Large, bowed tutoring halls hold lectures for the masses, while cleared patches of forest facilitate the training of warriors in the defense of both the city and the empire.

Bysaes Tyl was founded by elves who fled the destruction of Molaesmyr in the Savalirwood and quietly began to build a new home away from the eyes of the expanding empire. The elves of Bysaes Tyl struggled through the harsh winters for many years as they constructed their homestead, leaving the past behind to look to the future of their society. Adversity bred a hard, guarded, and tight-knit people.

When the Crown discovered that an independent city-state had formed within imperial territory, an imperial envoy was dispatched to present the elves with an ultimatum: accept Dwendalian rule or prepare for the empire to claim sovereignty over Bysaes Tyl by force. The Ring of Three, the elven elders, convened and agreed not to fight. Many of their kind would die in a senseless war—and perhaps the elves would outlive their imperial rulers many years hence. They agreed to submit and follow imperial laws and pay imperial taxes, under the condition that they could maintain some level of cultural autonomy within their city.

The elves guard the secrets of their culture as best they can, maintaining a tenuous relationship with the Crown that holds to this day.

\subparagraph{Leveraging a Legacy}
Bysaes Tyl remains the only city of elven creation within the Dwendalian Empire. The denizens feel an intense need to protect and maintain their cultural legacy, and to search for clues surrounding the mysterious destruction of their ancestral home. As per the agreement between Bysaes Tyl and the empire, the Ring of Three has incorporated some elements of imperial society into their own and shared a number of arcane secrets and recovered artifacts from Molaesmyr with the empire in exchange for their city's cultural autonomy. The elves wish to preserve their unique art and music, their crafting skills and techniques, and knowledge gleaned from the ruins of Molaesmyr that might decide the future of their people. The rumors that secrets and relics within Bysaes Tyl still remain hidden from the Crown and the Cerberus Assembly heighten the tension between the starosta and the Ring of Three, in addition to the slowly increasing imperial oversight.

\subparagraph{Government}
As part of the agreement that brought the city into the fold of the Crown, the Ring of Three, comprised of the elven nobles Allusus Dawyne, Gwey Yethneiros, and Amra Balleei, still oversee most aspects of city life. They must, however, work very closely with Starosta Uther Prent and his appointed counselors when governing the infrastructure and enforcement of law and order. Uther's prejudice against elves is well-known throughout the empire, and many believe this is why he was chosen for his position. The Ring of Three continues to deftly undermine Uther's leadership while simultaneously stroking his ego, a dance that they have thus far executed with grace but which cannot continue forever.

\subparagraph{Crime}
The elves of the city rarely engage in criminal activity, but escalating theft and brutality from non-elven imperial citizens living within—or passing through—Bysaes Tyl is a problem that the starosta largely ignores. Agents of the scattered Myriad have offered to help the Ring of Three manipulate the starosta and his council. The Ring of Three are weighing the benefits of letting an unlawful organization into their midst to help their oppressed people against the long-term consequences of working with such untrustworthy allies.

\subparagraph{Geography}
The city exists partially on the forest floor and partially within and around the trees, giving the settlement an air of rustic nostalgia that tempers its elven mystique. The walled surface streets are known as the Rough, or \textit{Talan}, where structures weave between rocks and wide-trunked trees, giving a slanted octagonal shape to the tall, wooden exterior walls of the city. Nearly all structures here are built from pine and other taiga trees, favoring designs that resemble a bowed arch rising to an apex. Spiraling staircases known as the Shrouds, or \textit{Halya}, twist around and within trees, climbing up to the elevated walkways and buildings amid the forest canopy.

\subparagraph{Bysaes Tyl Adventures}
The tension between the Crown and the elves in Bysaes Tyl plays a part in any of the city's adventures. Elves trying to maintain their cultural connection with nature and their lost homeland of Molaesmyr could look to adventurers for aid, the Crown could use help keeping the peace, and the Myriad could hire the characters to infiltrate the city.

\subparagraph{Elves, Arise (Mid Level)}
A group of elf assassins unhappy with the Crown's leadership hatch a plot to simultaneously assassinate the Ring of Three, Starosta Uther Prent, and other important figures in Bysaes Tyl. They believe that with the Crown distracted by the war, now is a perfect time to rise up and take back their city. The characters get wind of the plot and must decide whether to ignore, stop, or support the assassins.

\subparagraph{Evil, Arise (Epic Level)}
A powerful artifact of the long-deceased Archmage Oleyahs has been recovered from the ruins of Molaesmyr and is secretly in transit to Bysaes Tyl. Researchers hope that this object holds a key to understanding the destruction of their ancestral home, but when the object arrives, it is revealed to instead be a deceptive casing for the archmage's \textbf{demilich} form. A plague of arcane terror rains on Bysaes Tyl, and the characters are tasked by the Ring of Three to destroy or banish the demilich Oleyahs from the burning city.

\subsubsection{Druvenlode}
\begin{description}
\item[Population.]
12,110 (70\% humans, 14\% dwarves, 9\% elves, 7\% other races)

\item[Government.]
The city is run by Starosta Uvilia Whek, who oversees the mining guilds. Her rule is supported by Watchmaster Quin Theramast and her Crownsguard.

\item[Defense.]
A regiment of Crownsguard help protect the city, as well as the mines. Guildmasters and nobles also maintain personal guards.

\item[Commerce.]
Basic supplies, inns, and taverns are easy to find here. Ore and ingots are readily available, and rough, rare metals are traded on occasion.

\item[Organizations.]
Several temples can be found across the city. While smaller organizations exist within the city, the mining guilds are by far the most powerful and wealthy merchant groups.
\end{description}
\noindent
Established in the early years following the Calamity, Druvenlode is the primary source of precious metal in Western Wynandir. Located at the western base of the Silberquel Ridge by the mouth of the largest silver vein ever found on the continent. Dusty from wind-swept mountain dirt that turns to mud in the wetter seasons, Druvenlode is a far cry from beautiful Rexxentrum or any of the pastoral villages of the empire. Though the city is wealthy, buildings here were designed for function and storage while the talented architects focused on the noble homesteads and wards of the capital, leaving Druvenlode seemingly rough and rustic by comparison.

\subparagraph{Government}
Starosta Uvilia Whek oversees all major mining operations, and the larger mining guilds report directly to her. Maintenance of these mines also requires constantly protecting the workers and battling the dangerous creatures of the Silberquel Ridge, creating a thriving market for mercenary work when the Crownsguard are too busy—or too cowardly. The majority of the population works directly in or adjacent to the mining industry, laboring within the mines and tunnels, categorizing and smelting recovered minerals and metals, or forging materials into useful tools for distribution across Western Wynandir.

\subparagraph{Crime}
The local Crownsguard do their best to keep the peace, but with most eyes looking east to the war, the streets are ripe for criminal activity. Theft of goods and ore is increasing, and a call for additional, outside aid grows with it.

\subparagraph{Geography}
Druvenlode is divided into three main areas. Most citizens live in the ramshackle Dustbellows, where the clustered mass of simple hovels and retrofitted warehouses house those living hand-to-mouth. In these dreary streets, the need for distraction leads to a plethora of experimental music and dance at street corners and simple taverns. Silverstreet is a long thoroughfare where the talented folk of Druvenlode hawk their wares, advertise their skills, and seek legal ways to earn an income under the watch of the empire. The street is bookended by two massive gambling taverns that vie for dominance: the Sterling Satchel and the Cabin of Fortune. The Pillars is the industrial part of the city, where dark columns of smoke and exhaust belch forth from the massive smelting operations that dot the landscape. Most miners and physical laborers work in this area, and many of them live in the same soot-dusted factories where they toil.

\subparagraph{Druvenlode Adventures}
Adventure can be found in Druvenlode's mines, which occasionally fill up with burrowing beasts, as well as in the city proper, where shady characters are often found indulging in gambling or petty theft.

\subparagraph{Fire in the Mines (Mid Level)}
A being made of fire attacked miners from the Surgoth Guild who were working a vein of silver ore. The miners survived but are too scared to go back into the mine. Instead, they gamble their savings away at the Sterling Satchel. If the characters investigate, they find a path of destruction in the mine that leads to a \textbf{fire elemental} under the control of Casper Windmouth (neutral evil, male, human \textbf{mage}). Casper is a gambler with many debts to pay, so he took a job from an anonymous source to scare the miners out. Was it a rival miner's guild? Was it the proprietors of the Sterling Satchel, looking to hang on to their customers day and night? Or was it the owner of the Cabin of Fortune, hoping to bankrupt their competitor's clients?

\subsubsection{Dunrock Mountains}
The snow-dusted peaks of the Dunrock Mountains encircle the northern end of Western Wynandir, separating the Dwendalian Empire's territory from the scarred northern region of the Greying Wildlands. This majestic chain of rocky mountaintops is a snow-dappled bulwark against the corruption that torments the Savalirwood. The steep grades and icy terrain make the Dunrock Mountains perilous to traverse, and few beyond smugglers and gutsy hunters find reason to climb far into these sierras. Dappled with clusters of dense pine trees and cold streams, these beautiful mountains would be picturesque were it not for the roaming giants and wild griffons that hunt among the peaks.

Beneath the surface, the Dunrock Mountains contain a vast network of tunnels and caverns carved by the dwarves of Grimgolir before the Calamity, as well as hidden, collapsible passages utilized by the Kryn Dynasty to infiltrate the empire during this time of war.

\subparagraph{Dunrock Mountains Adventures}
These rocky peaks are home to many creatures and secrets beyond the empire's influence.

\subparagraph{The Crawling Caverns (Mid to High Level)}
Hidden deep in the heart of the Dunrocks are the forgotten depths of the Bastille of Torment, a massive, multilevel dungeon where Torog tortured the enemies of the Betrayer Gods during the Age of Arcanum. This horrifying abandoned prison is now run by a cabal of spirit nagas, who guard the many arcane secrets and treasures left by the Betrayers.

\subsubsection{Erdeloch}
The Erdeloch is the largest lake in Western Wynandir, often used as a landmark for travelers making their way toward the capital. Its chilled waters are the most prolific site for freshwater fishing in the empire, and the south side of the Erdeloch sees quite a number of large and small fishing vessels moving in and out of Odessloe daily. The many rivers that intersect with the lake make it a nexus of river transport within the empire, especially for loggers. The north shores, nestled within the Pearlbow Wilderness, are known to be filled with local wildlife and indigenous monsters, so few dare make landfall near the treeline at night.

\subsubsection{Icehaven}
\begin{description}
\item[Population.]
5,090 (71\% humans, 12\% elves, 10\% dwarves, 7\% other races)

\item[Government.]
Though technically run by Ernath Noack, the Crown-appointed starosta, the town is mostly ruled by a committee of locals who pay little mind to the will of the Dwendalian Empire.

\item[Defense.]
A small company of Crownsguard keeps the peace, while a fresh group of Righteous Brand soldiers have been assigned to the docks.

\item[Commerce.]
Resources are limited, so mostly basic supplies and standard arms are available. Relics recovered from Eiselcross can occasionally pass through, providing strange and exotic curiosities for acquisition.

\item[Organizations.]
Simple temples mark city squares, while hidden shrines to Crown-forbidden gods lurk on the outskirts of the town. The empire continues to push for stronger influence, and Myriad operatives interested in Eiselcross's mysteries are also beginning to work their way into Icehaven.
\end{description}
\noindent
Along the icy northeastern shores of the Zemni Fields stands the only port in Western Wynandir. Icehaven is home to the dozen or so ships in the region outfitted to wander the ice floes of the Frigid Depths, and is the central trade post and only open port for hunters and adventurers who brave the Saltwallow Bog or the frozen north toward Eiselcross. The inhabitants of Icehaven earn their living ice fishing within the oft-frozen Kaltenloch, or are trappers who stalk the Velvin Thicket or edges of the Pearlbow Wilderness and return to peddle their furs and meats. Some more confident individuals might try their hand at netting some Eiselcross crab, a rare delicacy throughout the empire, but the shifting, freezing weather makes such endeavors dangerous.

\subparagraph{A Sudden Concern}
As the war with the Kryn Dynasty escalated far to the east, few eyes were on the backwater port of Icehaven until fishers and explorers began to return from expeditions to Eiselcross with stories of odd ruins and even odder creatures. Rumors of pre-Calamity artifacts began to spread among the nobles of Rexxentrum, until the king and the Cerberus Assembly suddenly took an interest. With an influx of soldiers, sellswords, and fortune-seekers looking to make their way to the mysterious landmass, Icehaven and its residents are trying to adapt to the unexpected attention.

\subparagraph{Government}
While Icehaven is rarely of concern to the Dwendalian Empire, the town has a starosta. Currently the office is held by Ernath Noack, an older man who survived an extensive forty-year career as an Eiselcross crab fisherman. The local government is mostly left well enough alone except when it comes time for the Crown's tithe collectors to exact taxes. This yearly event has become something of an occasion to the lower-class fisher folk, who harangue the tithe collectors with fusillades of insults and endless pranks, each more humiliating than the last.

\subparagraph{Crime}
Due to the relative lack of imperial oversight and the harsh weather, many of the Crownsguard assigned to Icehaven put minimal effort into their watch, leaving room for laws to be bent or broken. The general sense of being at the end of the world during the long winter makes many within the town feel disconnected from the empire and its laws, which is both liberating and disconcerting.

\subparagraph{Geography}
Surrounding the small inlet harbor, Icehaven is a cramped, crescent-shaped town of snow-coated log cabins and lantern-lined street corners. The city is divided into two main areas: the Harborline, where the docks house the fishing vessels and the few ice-ready ships, and the Haven, where the crowded buildings seem to be clustered together for warmth. The streets shift between snowy mud and slicks of ice, curving through twisting alleys in a haphazard layout.

\subparagraph{Icehaven Adventures}
As a remote port town, Icehaven has need of adventurers to enforce the law, escort caravans, and find supplies in the wilderness.

\subparagraph{Snowed In (Any Level)}
While the characters are in Icehaven, a terrible blizzard strikes, leaving many residents to fend for themselves inside their homes. As the snowdrifts rise and wind howls, a gang of criminals or monsters comes in to cause mayhem and take whatever they like. If the characters don't brave the snow to stop these villains, no one will.

\subsubsection{Kaltenloch}
The harsh winters in the region cause the Kaltenloch to remain frozen over for three to four months out of the year, but the cold lake provides fish to the community of Icehaven year-round. With the ice comes a mysterious low mist that many say is made up of the spirits of fallen Calamity soldiers, looking for warmth to steal.

\subsection{Areas (N-R)}
\subsubsection{Nogvurot}
\begin{description}
\item[Population.]
15,270 (70\% humans, 17\% dwarves, 8\% elves, 5\% other races)

\item[Government.]
Starosta Brumae Luddie handles the affairs of the city with the support of the lawmaster's regiment of Crownsguard. Nobles under the starosta handle in local administration.

\item[Defense.]
A number of Crownsguard share protective duties with a regiment of Righteous Brand soldiers called the Last Aegis.

\item[Commerce.]
Well-stocked trading posts and inns make most supplies readily available, though rare or exotic goods are difficult to locate.

\item[Organizations.]
There are temples to empire-approved gods, and most businesses are independently run without the oversight of a guild.
\end{description}
\noindent
A major stop along the Glory Run Road heading eastward from the capital, Nogvurot hosts travelers making their way through the empire, many of whom need provisions and like to buy locally made trinkets as souvenirs of their journey. There is some business to be had in a small mine at the base of the Dunrock Mountains to the north of the city, but most Nogvurot folk earn their keep by crafting, farming, or joining the local military.

\subparagraph{A Healthy Defense}
Originally founded as a guard post near the Quannah Breach to watch for deadly creatures wandering from the Savalirwood, this settlement has since grown to be a fair-sized military reserve for the Righteous Brand. The regiment of the Righteous Brand stationed within the city would be the defending army against any assault headed for the capital, granting them the title of the last Aegis. Nearly half of the city's population is military reserve, with just over six thousand troops present at any given time. The bored or restless soldiers, seeking distraction and entertainment, developed the barracks of the Aegis Block into a place of sport and competitive game, often with impressive amounts of coin wagered on the outcomes. However, between reports of the Kryn taking the Ashguard Garrison and rumors of unexplained Kryn attacks within the Marrow Valley itself, the Aegis is very much on edge.

\subparagraph{Government}
Starosta Brumae Luddie, a retired general of the Righteous Brand, has watched over the city for fifteen years. Brumae is driven by a strong sense of duty and pride in the empire, never questioning the words of the king. Because Brumae is regarded by most as a fair and just starosta, some fear that his advanced age will soon bring an end to his tenure and allow less-scrupulous figures to vie for his position.

\subparagraph{Crime}
The strong military presence within Nogvurot keeps petty crime within the city to a minimum. Unfortunately, the Crownsguard often puff out their chests and try to one-up the soldiers of the Righteous Brand—and some local criminals have taken advantage of this petty rivalry to keep the soldiers distracted and continue their criminal behavior unimpeded.

Corruption is also becoming a problem, as Last Aegis soldiers use their sway within the city to make easy coin by turning a blind eye to smuggling to and from Shadycreek Run and other criminal hotspots in the Greying Wildlands.

\subparagraph{Geography}
The oblong city is divided between the eastern Aegis Block, which houses the bulk of the stationed soldiers, and the western region of the city that is broken into two sectors: the northern Scatters and the southern Peddler's Row. The Scatters is the larger region, housing dive taverns, simple temples, and a majority of the populace, who live in numerous neighborhoods of rickety domiciles. In contrast, Peddler's Row is a bustling home to local craft guilds and a colorful, if muddy, marketplace where the Glory Run Road passes through the city.

\subparagraph{Nogvurot Adventures}
Characters could go to Nogvurot to find someone in the city or to join the Righteous Brand. The city is also a good place for adventurers looking to test their mettle in competitive games against trained soldiers.

\subparagraph{Traitors in the Midst (Low Level)}
The characters are approached by Telva Cornwrath (neutral, female, human \textbf{spy}), who claims to work for the Augen Trust. She hires the characters to infiltrate the ranks of the Righteous Brand in Nogvurot and discover which soldiers are accepting bribes from the Tribes of Shadycreek Run. She wants to know which criminal family pays each crooked official. She could be who she claims, or she might be a spy for the Myriad or one of the Tribes in Shadycreek Run looking to remove competition. In the latter case, she might not want the characters sticking around as loose ends.

\subsubsection{Odessloe}
\begin{description}
\item[Population.]
6,970 (73\% humans, 12\% elves, 8\% dwarves, 7\% other races)

\item[Government.]
Starosta Vesdan Wothassic is eager to please the Crown and demands huge harvests from the fishers of Odessloe. As resentment builds, the few noble families who have business interests here work with Watchmaster Lein Untgaf to maintain order.

\item[Defense.]
A small cadre of Crownsguard watch over the city under the direction of the watchmaster.

\item[Commerce.]
As a popular trade stop on the way to the capital, Odessloe has many taverns and hostels. General stores provide basic supplies and occasionally offer strange relics pawned by cash-strapped adventurers.

\item[Organizations.]
The city harbors many small shrines to empire-approved gods. Most businesses are run by locals but owned by absentee nobles. The Myriad has a quiet presence here.
\end{description}
\noindent
Adjacent to the clear waters of the Erdeloch, Odessloe is a prominent fishing and logging settlement, supplying Rexxentrum and the surrounding townships with fresh fish and lumber. The only major stop on the Cloakjaunt Path for those traveling between Bysaes Tyl and the capital, Odessloe also acts as a supply stop for traders and drifters moving through the shaded trees of the Pearlbow Wilderness.

\subparagraph{The Waters Abide}
The bounty provided by the Erdeloch is a source of local pride as well as the backbone of the local economy. The people of Odessloe are a superstitious lot, and many believe that the lake is home to a spirit that must be kept happy if they are to continue to receive its bountiful blessings. They leave offerings at the shores of the lake to appease this unseen entity. Imperial officials, including the starosta, find this practice childish and actively dissuade citizens from indulging in "unlawful idol worship."

\subparagraph{Government}
Starosta Vesdan Wothassic has held his position for a year, since the passing of the beloved Starosta Ruell. Vesdan has big shoes to fill and is eager to prove his worth to the Crown. He has demanded resource yields from his citizens and enforced penalties on those who fail to meet his standards. Unsurprisingly, the people are already voicing their dissatisfaction with the new starosta. The nobles who profit from the local businesses grow worried that their investments are in danger—and are working with the watchmaster to try and suppress any talk of outright insurrection.

\subparagraph{Crime}
The city sees its small share of petty thefts and solitary criminal activity, but lacks the size and infrastructure to house an organized criminal element. This works to the advantage of the Myriad, which has been using the rivers that connect to the Erdeloch as smuggling routes, with Odessloe acting as a hub for transactions and storage, all without the government's knowledge.

\subparagraph{Geography}
Built along the eastern shore of the Erdeloch, this kidney-shaped city is comprised of three major areas. The muddy Lakeside region on the western end houses the docks, river-based storage buildings, and scattered housing for denizens of the district. Verdant and clean, the Greenset region on the eastern side of the settlement holds numerous logging camps and most of the residential homesteads for working families. The Crossmark region is the smaller center of the city, where the other regions meet at the juncture of the two largest thoroughfares, Cedarline Street and Hogmarshal Road. The smell of starchy foods and fresh fish suffuses the inns and market stalls of this sleepy intersection.

\subparagraph{Odessloe Adventures}
With the Myriad lurking in Odessloe, there are plenty of opportunities for characters who want to work for or against the organization. The settlement also makes an excellent base of operations for adventurers exploring the Pearlbow Wilderness.

\subparagraph{Feeling Crabby (Low Level)}
The residents of Odessloe fear that they somehow angered the spirit of the Erdeloch. Starving giant crabs from the lake have terrorized the Greenset region of the settlement every night for the past week. The cityfolk have left expensive offerings for the spirit on the lake's banks. These offerings are taken, but the crabs continue to attack. If the characters investigate, they find a group of bandits who are capturing and starving the crabs before releasing them into the settlement, and stealing the offerings for themselves. Are these ruffians connected to the Myriad? Could the stolen offerings anger the actual spirit of the lake enough to punish the settlement for the disrespect?

\subsubsection{Pride's Call}
\begin{description}
\item[Population.]
16,090 (81\% dwarves, 8\% humans, 6\% halflings, 5\% other races)

\item[Government.]
Starosta Ghuldi Demosk walks the fine line of honoring imperial rule without undermining dwarven tradition. Small guilds run by dwarf nobles support his rule.

\item[Defense.]
Watchmaster Hildas Cliffhollow oversees the local Crownsguard, which is sometimes bolstered by mercenaries.

\item[Commerce.]
Most goods and services are available here, including room and board for travelers. Pride silk is a popular commodity.

\item[Organizations.]
Empire-approved gods have humble temples alongside the local guilds and businesses, while Myriad agents lurk within the shadows of the city.
\end{description}
\noindent
Within the rolling hills of the southern Zemni Fields lies the dwarven city of Pride's Call. Established in the ashes of the Calamity and slowly rebuilt as the center of hill dwarf culture in Wildemount, Pride's Call chose to ally with the Dwendalian Empire during the empire's expansion nearly two centuries ago. Their favor as the empire's first allies eventually waned when the dwarven city of Grimgolir was absorbed a century later, leaving Pride's Call overshadowed by the new, impressive, and strategically significant dwarf city. However, Pride's Call is nonetheless respectable in martial might and historically rich in culture, music, and their coveted pride silk. The city is a bright spot in the otherwise cold and dour northwestern Wynandir.

\subparagraph{Even Favor Fades}
During the invasion of the Julous Dominion, the empire saw many great deeds done under the banner of Pride's Call. After helping the empire achieve victory, the people of Pride's Call enjoyed a period of celebratory attention. Their fame subsided, however, as Dwendalian attention turned to the mounting threat of Xhorhasian terrors boldly crossing the Brokenveil. Talks had already been privately underway to bring the dwarven stronghold of Grimgolir under imperial rule, given their military might and advantageous proximity to the brewing danger. With the combination of Grimgolir's power and the empire's plentiful resources, victory was ensured and a strong sentinel force at the Xhorhasian border was established. Grimgolir became the new favored home of the dwarves in Wynandir, and many Pride's Call dwarves who lived through those years still hold a grudge—a fact that the Crown deftly wields to exploit both cities.

\subparagraph{Squeezing Silk from a Stone}
The dwarves of Pride's Call produce a number of useful exports, but none are as unique nor as sought after as pride silk. This remarkable fabric can only be woven from materials found in the unique ecosystem of the geothermically active hills around Pride's Call. A special breed of silkworms has developed in these strange conditions. These fist-sized silkworms produce abnormally thin and incredibly resilient strands, and a secret process of reinforcing the threads creates a material similar in texture and appearance to thick silk, but with the strength of treated leather. Affluent folk across Exandria pay handsomely for garments made of pride silk, but rumors are beginning to spread that outsiders have uncovered the secret, and the tradition might be slipping from the grasp of Pride's Call.

\subparagraph{Government}
Pride's Call is governed by the dwarf who signed the original pact with the empire. Starosta Ghuldi Demosk spends his twilight years doing what he can to maintain his culture in the face of an increasingly invasive rule. Idealistic and unable to see any error in his judgment, Ghuldi is seen by many as hardheaded among even the most stubborn of dwarves. Many of his people are loudly voicing their discontent, and younger politicians such as Usula Starsburn are rising in popularity to succeed as starosta when Ghuldi eventually passes away. Fearing that he will have to choose between his legacy and the opinions of younger generations, Ghuldi is beginning to exhibit more totalitarian behaviors, slowly becoming a petulant dictator.

\subparagraph{Crime}
While the city's colorful veneer makes it seem a place of plenty, the economy of Pride's Call has been in slow decline for some time. The landscape's strange mineral composition makes its produce a healthy but unsavory and unpopular export, while other goods that previously brought the city prosperity are now more readily and quickly available from traders on the Menagerie Coast. This economic downturn has led to a rise in crime in the impoverished regions of the city, and also opened the door for bolder Myriad activity.

\subparagraph{Geography}
The city is comprised of four boroughs built atop and within a cluster of large hills that wrap around the base of the Cyrios Mountains. The two largest boroughs, known as the Tillers and Crafters, are twin sections of mixed residential neighborhoods with areas of concentrated commerce. The Tillers cultivates produce and meats for the city, though most local food has a bland and bitter taste, making spices a popular import. The Crafters contains numerous forges, tanneries, and the well-guarded subterranean silk-production warehouses. Nearest to the mountains is the Vigil borough, where the locally trained Crownsguard mingle with those assigned from Bladegarden to watch over the hillside. The western branch of the Righteous Brand are also stationed in the Vigil borough. Between Vigil and Crafters lies the fourth and smallest borough, called Moot, which is marked by a circular road surrounded by the wealthier houses. The Moot itself is an open, stone forum for discussion and debate that long predates the empire's rule.

\begin{DndSidebar}[float=!b]{Pride Silk}
The hearty silk unique to the region has long been a point of pride with the locals. When processed and woven properly, the resulting cloth is \textit{pride silk}: an extremely flexible material with the durability of tanned leather. Pride silk has become a favorite cloth for use on military ship sails, war banners, and other utilitarian applications, but has also gained popularity in affluent circles. High-end clothing made from pride silk is both fashionable and protective, leading many nobles and wealthy merchants to seek full outfits and formal gowns made with it. The tailors skilled enough to create such outfits are few and far between, so while pride silk might be expensive as a raw material, outfits made of the cloth are prohibitively expensive for most common folk.



One square yard of \textit{pride silk} weighs 1 pound and costs 100 gp. An Pride Silk Outfit weighs 4 pounds and costs 500 gp. If you aren't wearing armor, your base Armor Class is 11 + your Dexterity modifier while wearing it.



\end{DndSidebar}

\subparagraph{Pride's Call Adventures}
If the characters seek pride silk or master artisans, Pride's Call is where they should go. In a city experiencing political and economic turmoil, there is no end to the intriguing adventures characters might experience.

\subparagraph{Worm Your Way In (Low Level)}
The characters encounter a traveling textile merchant named Cardorn Brentahill (neutral evil, male, human \textbf{noble}), who lives in Yrrosa. He's seeking bodyguards to help him escort his cart back to his home, paying 100 gp per hired hand. Along the journey, the characters and their employer are attacked by a crew of mercenaries representing the Durrodan family, claiming that Cardorn stole silk worms from their silk farm. Cardorn offers to increase the characters' pay if they defend him. If they side with the Durrodan family, Cardorn is taken prisoner but the Durrodans grant them a small reward and the promise of future work. If they side with Cardorn, fight off the mercenaries, and escort Cardorn home, he introduces the characters to Gothad Miskal (neutral evil, male, halfling \textbf{bandit captain}), a Myriad taskmaster from Yrrosa who is interested in offering better pay for less legal work.

\subsubsection{Rastum Den}
Hidden among the Dunrock Mountains, north of the Pearlbow Wilderness and Rexxentrum, is an encampment for breeding and training military griffons. These mighty beasts are used by the Dwendalian griffon riders of the Righteous Brand, some of the best-respected soldiers in the empire.

Known as Rastum Den, this moderately sized encampment is nestled in a small valley four thousand feet up in the mountains. The area's dizzying heights and tumultuous weather make it all but inaccessible to those without a griffon rider's prowess. This location was originally home to the largest griffon population in Wildemount, making it natural terrain to quickly condition any young mounts bred for the Righteous Brand.

\subsubsection{Rexxentrum}
\begin{description}
\item[Population.]
205,200 (81\% humans, 8\% dwarves, 6\% halflings, 5\% other races)

\item[Government.]
King Bertrand Dwendal rules the city and the Dwendalian Empire from Rexxentrum. He delegates lesser matters to Arbiter Sydnock Truscan, while citywide guild business falls under the purview of Guildmaster Kai Arness.

\item[Defense.]
Crownsguard patrol all major streets, while the bulk of the Righteous Brand soldiers within the empire are stationed in the King's Bastion to defend the capital city.

\item[Commerce.]
Rexxentrum is the largest city in Wildemount and the seat of the Dwendalian Empire's power, enabling massive volumes of trade in and out of the city. Most any goods, services, and places of business can be found within the city walls.

\item[Organizations.]
Grand temples to empire-approved gods can be found throughout the city, while the most powerful guilds in the empire vie for power and influence. The Cerberus Assembly takes an active role in political affairs, while the deep-rooted Myriad weaves tapestries of crime and corruption throughout the city.
\end{description}
\noindent
The majestic capital of the Dwendalian Empire was founded in the year 514 PD. The central location and plentiful surrounding resources of the region made Rexxentrum the most favorable place to oversee the rule and expansion of the freshly established Dwendalian Empire. Over the next three centuries, Rexxentrum grew to become the largest city in Wildemount, and the cultural and military center of Western Wynandir.

At the core of all imperial politics are the nobles, politicians, and wealthy elite who live within the city or own property there. Grandiose displays of wealth and power are common among the inner wards, while sprawling poverty consumes the outer districts and slums beyond the protective walls of the city. Cobblestone streets twist and mingle through the bustling neighborhoods, cart merchants and street performers vying for the attention and patronage of passersby throughout the cold summer months and the even colder winter.

Glowing orbs of arcane fire flicker at the tops of dark iron lamp poles at major intersections, lighting the evening air, while colorful flowers adorn the inner archways of each neighborhood entryway. Buildings cluster close and reach toward the sky, often snow-dappled with long wooden slats placed between rooftops for riffraff to skitter across the city unseen. Although the city teems with life and activity, it is also full of lawkeepers and military watch, with griffon riders occasionally circling overhead above the ever-patrolling Crownsguard and the nearby Righteous Brand headquarters.

The massive and beautiful center for arcane study known as the Soltryce Academy is located in the Candles, a thin crescent of a neighborhood marked by the tall and imposing tower homes of all eight leading members of the Cerberus Assembly, each unique and imposing in its own way. The beating heart of the city, and of the empire, lies positioned within the arc of the Candles in the terrace of Castle Ungebroch. From within the castle halls, King Bertrand Dwendal and his council of chosen authorities oversee all matters of state.

\subparagraph{Magic Masters}
Much of the empire's prosperity over the centuries stems from the Cerberus Assembly, a union of the nation's most powerful mages working directly for the Crown. These arcane practitioners and scholars aid the king in engaging foreign interests, in developing magical technology to better oversee the land and enforce his will, and making use of the rare relics that the Crown confiscates from unearthed sites of pre-Calamity might. This arrangement affords the assembly homes in the Candles and across the major cities of the empire, as well as access to every Dwendalian political figure. Thus, members of the assembly are extremely influential and, by the nature of their arrangement, very carefully watched by the king. Most civilians swap rumors of demonic pacts and demigod-level magics at the assembly's disposal, adding to their aura of mystique among the common folk. When an assembly mage walks the streets of Rexxentrum or attends the halls of an establishment within the city, it is an event that everyone will be talking about for days to come.

\subparagraph{Rule with a Golden Fist}
The resource-rich economy of Rexxentrum supports a powerful guild and merchant class, who work directly with the Crown under the watch of Guildmaster Kai Arness. Any permanent place of business is required to join a corresponding guild, paying an annual fee on top of their taxes to the Crown, and to follow the guild's requirements for conditions of operations. Some complain that these conditions are often manipulated on a case-by-case basis to remove competition between similar interests, but cries of corruption within the guilds fall on deaf ears. This has led to a rapid growth of mobile businesses on street corners and shops that move between districts to avoid inspection by Crownsguard or guild members.

\subparagraph{Politics and Persecution}
Historically, public displays of unease under the king's rule were often met with brutal punishment. However, the previous king, Everard Dwendal, loosened the laws applying to public discourse. This gave rise to a number of political figures who gathered the unhappy populace behind their philosophies. In response, King Everard established the Augen Trust, a network of spies working for the king to blackmail or corrupt these politicians. With these high-profile dissenters under the thumb of the Crown, controlling political discourse became easier, and politics within the empire became a tense game of cat and mouse. Forthright nobles and frustrated guild heads unswayed by propaganda, gold, or threats still dissent publicly, testing the boundaries of their influence. King Bertrand's patience with public objections to his policies grows ever thinner, and he is enacting even more stringent laws against public criticism of the Crown.

\subparagraph{A Growing Unrest}
In the wake of frustrated citizens, neglected concerns, and a clear caste system, the common folk of Rexxentrum are disenfranchised. The rallying call of war has brought some renewed imperial pride, but whispers of discontent have continued, coalescing into a small group of rebels within the capital that call themselves the Locksmiths. Wary of the history of failed rebellions in the empire, they gather information and gradually inspire the common folk to question authority, waiting for the right moment to ignite a revolution with a chance of success. With the Crown distracted by war and the populace brimming with poverty and discontentment, that time might come sooner than anyone thinks.

\subparagraph{Government}
King Bertrand Dwendal rules the capital himself. No starosta holds office within Rexxentrum. Prime Arbiter Sydnock Truscan oversees all major matters of law that do not require the attention of the king, while Guildmaster Kai Arness works with the various guilds of the city to ensure smooth trade both within and beyond the walls of Rexxentrum. Beneath Sydnock is the head of Rexxentrum's Crownsguard, Watchmaster Brunda Gleshin, who ensures that the laws of the empire are enforced in the streets of the capital.

\subparagraph{Crime}
The capital city harbors a diverse criminal underground both within and beneath the streets. Surviving Myriad leaders remain hidden in plain sight throughout the city, including the beloved Iris Bethelas, who is organizing the syndicate's return to power. Independent crime rings plague trade centers and busy districts where there are easy gold and gullible cityfolk. The populace is simply too large for the Crownsguard to oversee every corner of cobblestone where laws are bent or broken. Now, with the collective attention of the law focused on the war in the east, thieves and cutpurses are upping their game, and crime is on the rise within Rexxentrum.

\subparagraph{Geography}
The city spreads across the tallgrass fields south of the Pearlbow Wilderness, its walls forming a distended pentagon. Farmland around the north and east sides of the city are dotted with farmsteads and small rural communities that provide food for the capital. The northeast side of the wall curves around a trio of immense, fortified towers surrounded by barracks known as the King's Bastion, where the central force of the Righteous Brand military trains and resides. Within the outer walls of the city are seven major districts: the Mosaic Ward, the Mudtop Ward, the Tangles, the Shimmer Ward, the Academy Grounds, the Candles, and Castle Ungebroch Terrace.

\subparagraph{Mosaic Ward}
A sprawl of impoverished neighborhoods are hidden behind colorful banners and small fences that segregate this ratty residential district from the central thoroughfares and avenues. The region, which encompasses most of the southern and southwestern regions of the city along the outer walls, is a hotbed for crime and violence. Crownsguard interest in policing the region wavers, and opportunists often use the lack of military presence to victimize poorer citizens. This has led to a growing collection of vigilantes and self-appointed protectors of the Mosaics.

\subparagraph{Mudtop Ward}
The northern end of the city houses its own large slum, known as the Mudtop Ward. The cobblestone streets of the city were never quite built up to the northern wall, causing this area to be rife with mud and muck year round, leaving the denizens cold, ornery, and prone to illness. Many of the sick and diseased are relegated to the Mudtop to be forgotten. However, some healers have come to their aid, as well as some priests who wish to remedy the unjust conditions of the poor, becoming local heroes for their selfless work.

\subparagraph{The Tangles}
This western region of Rexxentrum houses the working class folk of the city. It is here that most of the city's general commerce takes place. Dozens of smaller streets branch off major roads to create a labyrinth of twisting neighborhoods and industrial districts. The network of major roads converges to form three major plazas across the district, where shops, inns, and various forms of entertainment can be found. The Platinum Veranda plaza lies at the northern end of the Tangles, where the more expensive homes and luxurious entertainment can be found, as well as the central temple to Bahamut, called the Dragon Seat.

Near the southern end of the Tangles is the Vigil's Circle, a marketplace built around a loop of road that surrounds the Tower of Writ, the temple to Erathis where the Prime Arbiter delivers his judgment. All manner of commerce takes place here, while the nearby Claykeep prison looms as an ominous reminder to any who might seek to break the laws of the land.

The Court of Colors, on the western edge of the Tangles, is the largest of the three main plazas. The Court is a winding knot of streets, each of which is wholly dedicated to traders who come from far-off places to sell their exotic wares, while artists and performers bring an air of lively celebration despite the gray and dreary skies of the region. It's within this lively area that the Rexxentrum Archive of the Cobalt Soul can be found. Near the southern wall of the Shimmer Ward is the Chantry of the Dawn, the ancient temple to the Dawn Father on which the city was originally founded, now restored and expanded into an awe-inspiring cathedral.

\subparagraph{Shimmer Ward}
Within the inner walls of the city, the Shimmer Ward is where the noble houses and guild leaders of Rexxentrum dwell. Here, beautiful gardens and museums celebrate Zemnian culture, and specialist merchants promote their unique wares. The debate forum is rarely used aside from holidays, where intellectuals will challenge each other to arguments about philosophy and history. Cottages and tasteful estates sit behind tall fences and walls, and a pair of open squares host occasional galas when the local elite are in need of socialization.

\subparagraph{Academy Grounds}
To the northwest of the castle lie the Academy Grounds, a well-kept and heavily guarded campus where the halls and towers of the Soltryce Academy reside. A college of higher learning for all pursuits related to the arcane arts, the Soltryce Academy is all but exclusively a school for wealthy scions, though a number of rare talents found in the wild are offered entry every year. Each subject's curriculum is tailored so that the young mind is indoctrinated to support the interests of the empire, and the Cerberus Assembly is always looking for prodigies with exemplary talent to take under its wing.

\subparagraph{The Candles}
East of the castle stands the collection of eight tall spires that house the members of the Cerberus Assembly. Each tower has its own unique appearance. The original Candles were built to the specifications of the original leading members of the Cerberus Assembly. In the years since, each tower has been magically altered to the taste of their successors. This creates a jagged and subtly unsettling sight visible from any point in the city.

\subparagraph{Castle Ungebroch Terrace}
Standing vigilant in the center of the inner walls is the Castle Ungebroch Terrace, a small district of compact military outposts that surround Castle Ungebroch. The dark stone of the castle splits into numerous towers, while high-arched stained glass windows dress the exterior of each floor of the citadel. Within the shadowy halls of the castle, lit by flickering torchlight, the king rules with his family and councilors at his beck and call.

\subparagraph{Rexxentrum Adventures}
A noble house in Rexxentrum that isn't scheming for more power is a noble house in decline, or so the gentry believe. The city's officials and nobles have schemes to grab more money and power, and more than a few of the common folk have their own plots to launch themselves into positions of influence. Most of these social climbers need the help of adventurers willing to take risks, while others often need protection from their fellow schemers.

\subparagraph{Call a Locksmith (Low Level)}
The exploits of the characters might draw the attention of the Locksmiths, a resistance group seeking to undermine the totalitarian rule of the king. A friendly face the characters have come to know within the city invites them to meet Lady Quana Seledo (chaotic good, female, human \textbf{noble}), a well-connected figure who reveals herself to be a key figure in the Locksmiths. She asks the characters to join the cause, offering allies, rare goods, and well-paid work. If they agree, the first mission begins: help the Locksmiths infiltrate the office of Exchequer Aethia Drooze and steal the current year's tithe ledgers.

\subparagraph{King Killers (Mid Level)}
The characters intercept a cryptic note which indicates that there will be an attempt on King Bertrand Dwendal's life the next time the monarch makes a public appearance. The characters are too inexperienced to get an audience with the king or anyone close to him, so it's up to the characters to thwart the plot. These assassins could be sent by the Kryn Dynasty, the Cerberus Assembly, the Myriad, the cult of a Betrayer god, or by the Crown itself to garner sympathy among the people. If the characters manage to stop the assassins, they earn the notice of the king, who will likely ask more of them.

\subsubsection{Rockguard Garrison}
\begin{description}
\item[Population.]
6,800 (71\% humans, 12\% halflings, 12\% dwarves, 5\% other races)

\item[Government.]
The garrison is commanded by Marshal Garad Velious, who oversees regiment captains that hold sway in the area.

\item[Defense.]
Three regiments of Righteous Brand soldiers defend and patrol the garrison and the surrounding mountains.

\item[Commerce.]
Most supplies within are military rations and stationed equipment. Trade with the local soldiers can yield goods and services.

\item[Organizations.]
Small shrines to empire-approved gods have been constructed by the soldiers. Aside from the Righteous Brand, the Crown, and their allies in Grimgolir, no other organizations have authority over the garrison.
\end{description}
\noindent
The Rockguard Garrison was constructed in a joint effort between the Crown and Grimgolir as a symbol of their political union. The sister-fortress to the Ashguard Garrison is the foremost defense post against the ever-encroaching forces of the Kryn Dynasty and the monstrous hordes of the wastelands to the east. It boasts a formidable army and a near-impenetrable perimeter that expands into a series of defensible weapon posts along the sides of the valley to further dissuade any force from attempting an invasion. This fortress has held strong for generations without being breached, but the loss of the Ashguard Garrison is still fresh in the minds of the soldiers stationed here. A cold intensity hangs over each passing day as the Dwendalian forces plot to reclaim their captured outpost and remain vigilant against a similar assault.

\subparagraph{Rockguard Garrison Adventures}
Rockguard Garrison adventures usually relate directly to the war. Defending the outpost from attacks, planning strikes against the Kryn, and other military missions for the Crown are carried out here on a daily basis.

\subparagraph{Steal the Plans (Low Level)}
Agents of the Kryn hire the characters to infiltrate the Rockguard Garrison and uncover the Crown's plans to retake the Ashguard Garrison. The Kryn want any physical plans the characters find to remain in place to ensure the Crown doesn't realize the plot, so all written documents containing key information need to be copied. The garrison's guards, scouts, and veterans patrol the place and question characters who obviously sneak around or don't belong.

\subsubsection{Ruins of Shattengrod}
This sinkhole is a rocky pit measuring roughly thirty feet across and sixty feet deep. Littered with ancient bones, the pit was long considered just a macabre curiosity, referred to by the dwarves of Pride's Call as the Grave Hole. When tales of the strange location reached inquiring minds within the empire, an investigation was launched. Beneath the pit, explorers discovered the sundered ruins of a prominent city from the Age of Arcanum called Shattengrod. Heavy excavation expanded the pit into a spiraling hole that delves into the subterranean city, unveiling broken relics and forgotten secrets—as well as some enduring shades and hungry entities now awakened. The promise of beneficial discoveries pushes both interests of the Crown and the Cerberus Assembly to delve into the ruins, but the deaths of mercenaries are already drawing attention as rumors of the dangers below spread and scare off many possible hires.

\subsection{Areas (S-Z)}
\subsubsection{Saltwallow Bog}
Born from the sunken muck that merges with the northern treeline of the Velvin Thicket, the Saltwallow Bog is a marsh of briny, slightly acidic waters. Myths speak of a terrible servant of the Crawling King that was slain during the battles of the Calamity, its corrosive corpse sinking into the floor of the forest and corrupting the thicket around it until this toxic bog was left as a final cry of defiance. Strange and twisted creatures have adapted to live in this grotesque place, delving into the Velvin Thicket to hunt and drag their prey back to their lairs in the bog. As dangerous as they are, some of these creatures are now highly prized by tanners and alchemists for their leather and harvested parts, and hunters venture into the depths of the bog year-round, risking life and limb for coin.

\begin{DndSidebar}[float=!b]{Dyolet Fruit}
A round, blue fruit roughly the size of a plum, a dyolet is known for its sweet flavor and the mild surge of energy granted to whoever eats it. Dyolets can also be imbued with minor healing abilities, a property kept secret by the gnomes of the Velvin Thicket until recent years. Herbalists and apothecaries can imbue healing within the fruit for a fraction of the time and gold cost of brewing a potion. A dyolet can be imbued with one of the following properties by a character proficient with a healer's kit:




\begin{itemize}
\item You can saturate the fruit with carefully brewed curative draughts to imbue a dyolet with the properties of a \textit{potion of healing}. This process takes 12 hours and costs 15 gp.
\item You can saturate the fruit with strengthened restorative properties to imbue a dyolet with the properties of a \textit{potion of greater healing}. This process takes 4 days and costs 70 gp.
\end{itemize}



You can consume a dyolet as an action to gain its healing effect. A dyolet lasts up to 20 days after being harvested before it begins to rot. A rotting dyolet loses its healing properties.



\end{DndSidebar}

\subsubsection{Silberquel Ridge}
The imposing mountain range known as the Silberquel Ridge extends across the middle of Western Wynandir, dividing the Marrow Valley from the Pearlbow Wilderness. A range of jagged mountain tops visible from most cities in the empire, the Silberquel Ridge is an important landmark, but it is best known for its seemingly endless cache of rich metals and deep caverns. The ore and gems mined here provided Rexxentrum with most of its wealth and funded the rise of the Dwendalian Empire. The surface veins are the largest source of silver in Wildemount, but the deeper paths reveal other precious metals, gems, and forgotten remnants of pre-Divergence curiosities—and many deadly beasts that have come to call those ancient corners home.

\subparagraph{Silberquel Ridge Adventures}
Traveling across the Silberquel Ridge is its own adventure, given the number of dangerous creatures that inhabit the region's mountains and tunnels. Caravans, miners, and prospectors hire adventurers to guard them as they travel and work in the mountains.

\subparagraph{Hell Hound Purge (Mid Level)}
Marta Denizreitch (chaotic neutral, female, human \textbf{mage}) lives in a remote tower where she experiments with conjuration magic to develop new spells. The Cerberus Assembly funds Marta's research. Years ago, Marta had a problem with trolls attacking the tower, so she summoned hell hounds to roam the wilds near her tower and keep the trolls at bay. She does not want the hell hounds to cause problems down the line, so she hires the characters to hunt a few to keep the population down. The characters discover they aren't the only ones killing hell hounds: a \textbf{hill giant} and a small gang of trolls are killing the hounds in preparation for an attack on Marta's tower!

\subsubsection{Velvin Thicket}
\begin{description}
\item[Population.]
850 (98\% gnomes, 2\% other races)

\item[Government.]
The people of the Thicket are governed by a loose union between the leaders of the Velvins, the collection of local nomadic clans who convene only when threatened.

\item[Defense.]
While few are specifically trained for combat, many denizens of the Thicket are well-versed in illusion magic for self-defense.

\item[Commerce.]
Those that have the trust of the Velvins can acquire basic supplies, dyolet fruit, and the occasional unusual relic.

\item[Organizations.]
Scattered altars to the Wild Mother exist throughout the Thicket. Each of the seven nomadic clans handles its own internal affairs. Representatives from Pride's Call often visit for a time to maintain their relationship with the Velvin leaders.
\end{description}
\noindent
Bordering the soggy fields of the Saltwallow Bog is the strange forest known as the Velvin Thicket. The nomadic gnomish clans that call these dense woods home have a loose network of allied families rather than a central township, and they defend their shaded sanctuary with illusions and enchantments. Because of the Velvins' trade agreement with Pride's Call, the dwarven city was able to request that the empire keep the thicket from becoming formally assimilated into Dwendalian society so long as the Velvins provided a small but consistent crop of dyolets. The dyolet fruit is a rare delicacy receptive to minor enchantments and only grows within the Velvin Thicket. The empire's interest in permitting Velvin autonomy has waned over the years, and the Cerberus Assembly is keen to take the Thicket under its jurisdiction and seize control of the dyolet crop.

\subparagraph{Velvin Thicket Adventures}
The gnome families that live in the Velvin Thicket are suspicious of outsiders. The Cerberus Assembly's interest in the area might be enough to bring the characters to this strange place.

\subparagraph{Dyolet Mounds (Mid Level)}
Herwin Melvot (chaotic neutral, male, forest gnome \textbf{mage}) attempted to grow a new kind of dyolet with magic. His experiment resulted in shambling mounds made of dyolet plants that terrorize the Thicket. The gnomes allow the characters into the forest if the adventurers slay the plant monsters.

\subsubsection{Vergesson Sanatorium}
During the war between the Dwendalian Empire and the Julous Dominion, a vast prison known as the Archevault was made as a jail for prisoners of war. After the war, the Archevault was repurposed by the Cerberus Assembly as the largest imperial sanatorium and hospital. Named after Lord Symon Vergesson, the late assembly member who founded it, the massive mansion is surrounded by vast gardens and filled with beautiful artwork designed to please the primarily wealthy patients.

The sanatorium's patients are treated for physical and mental ailments severe enough to require alchemical, magical, or experimental therapy. Helmed by Gertrude Wagner, the Sanatorium Director, the faculty is monitored by the Crown. The facility offers treatment to those who require assistance and can pay handsomely, but beneath the pleasant mansion lies a torturous oubliette for political prisoners.

\subparagraph{A Dark Utility}
While the reputation of the Vergesson Sanatorium is sterling, it hides dark secrets. The original purpose of the Archevault was never fully discarded, and a number of vocal critics of the Crown have disappeared from the streets of Rexxentrum. Deep beneath the floors of the estate lie hidden chambers and cells that house enemies of the king, foreign spies, and traitors to the Crown privately sentenced by Prime Arbiter Sydnock Truscan. There, they are subjected to experiments conducted in secret under the supervision of Master Trent Ikithon of the Cerberus Assembly. A tower on the northern outskirts of the grounds acts as Trent's summer home, where he houses and trains the prospective members of his Volstruckers, an elite group of spies and assassins who use magic to do the assembly's dirty work.

\subparagraph{Vergesson Sanatorium Adventures}
Most who adventure in the Vergesson Sanatorium are trying to free a prisoner. Just finding a captive hidden in the massive complex is a challenge, but if the characters become imprisoned in the Archevault themselves, they must survive and form a plan to escape.

\subparagraph{Killing Regin (Mid Level)}
Regin Kavla (lawful evil, female, human \textbf{spy}), a former spy for the Myriad, is trapped in the bowels of the Vergesson Sanatorium. Worried that she might share important secrets, the Myriad hires the characters to get into the prison, capture Regin, and delivered her to them, offering a hefty sum or some other reward—whatever it takes to get the job done. As the characters search for Regin, they uncover the hospital's dark secrets and must decide what to do with this disturbing information.

\subparagraph{Patient Zero (High Level)}
One of the most trusted allies of Prime Arbiter Sydnock Truscan, Doctor Reynold Brueska, has been privately developing an arcane disease to be used in the war against the Kryn Dynasty, with the sanatorium's patients as his test subjects. When all contact with the sanatorium ceases, the Prime Arbiter hires the characters to investigate the premises and uncover the doctor, his work, and the events that transpired. The characters discover that a majority of the denizens of the Vergesson Sanatorium have been mutated and driven mad by the disease. The adventurers must fight their way through the horror-filled halls while trying to avoid contamination, or discover a cure before the disease completely corrupts them.

\subsubsection{Yrrosa}
\begin{description}
\item[Population.]
3,220 (54\% humans, 28\% dwarves, 14\% elves, 4\% other races)

\item[Government.]
The Crown's will is imposed by Starosta Zulla Sauer, in tandem with the local noble families.

\item[Defense.]
Yrrosa is protected by a battalion of Crownsguard supplemented by mercenaries.

\item[Commerce.]
An older and smaller trade post, the town has basic equipment and supplies available for trade or purchase. Rooms and taverns are humble, but available.

\item[Organizations.]
A temple to Bahamut looms over smaller shrines to the other empire-approved gods. The Myriad was once headquartered here; in its absence, the local businesses are entirely independent.
\end{description}
\noindent
Once a snowy outpost between Icehaven and Rexxentrum, Yrrosa has slowly grown into a sizable township where travelers and traders warm themselves as they pass through the frigid fields. The town is comprised of humble homesteads and numerous inns that promote themselves as "the hearth before the hoarfrost."

Recently, adventurers and fortune-seekers have flocked to the area to aid in the excavation of the nearby Ruins of Shattengrod (described earlier in this chapter).

\subparagraph{Siring a Syndicate}
Eighty years ago, a trapper guild turned smuggling gang was ousted from Icehaven and came to Yrrosa to seek better fortune. Life was no easier, and the crew of hungry rogues was forced to become creative. Noticing the lack of oversight from the empire, they began to build a network of like-minded individuals throughout Western Wynandir. Taking the name "the Myriad," they helmed illegal activity from the comfort of this small outpost town until fifteen years ago, when their base of operations was discovered and raided by imperial forces. The remaining Myriad leaders have scattered, but the shadow of their operation still looms over Yrrosa like a dark cloud.

\subparagraph{Government}
The previous starosta, Stefan Weberrat, was arrested fifteen years ago for protecting the Myriad. His successor, Zulla Sauer, is the willful and respected daughter of Petrov Sauer, a mining baron in Druvenlode. She was chosen after she volunteered for the position, but unknown to most, she was already involved with the Myriad long before she became a starosta. She now helps inform their hidden network from a position of safety.

\subparagraph{Crime}
Petty crimes occur, but most stem from disagreements or competition over parties using Yrrosa as a base of operations while delving into the Ruins of Shattengrod. Unbeknownst to all but their allies, the Myriad has been slowly carving out new bases, and in time, they might grow bold enough to reclaim their old hideout.

\subparagraph{Geography}
Ramshackle wooden domiciles with snow-covered roofs line muddy streets throughout the town. The main roads were more improvised than planned, and most of the older buildings are in varying states of disrepair. Where Pride's Walk Road intersects with Havenpath Road is a cluster of taverns, inns, liveries, and stores.

\subparagraph{Yrrosa Adventures}
The shadow of the Myriad hanging over Yrrosa makes it a fine place for intrigue, and the many travelers moving through the town have rumors and jobs to share with adventurers.

\subparagraph{Blood Pact (High Level)}
A human blacksmith named Wilhelm Arrisgon never returned from his latest journey to collect a wagon of raw materials. Wilhelm's family hires the characters to find the smith, and they trace him to Yrrosa. During their investigation, the characters discover that several other solitary travelers have disappeared while passing through the city. Velima Shanglian, a \textbf{vampire} who lives in a hidden lair outside Yrrosa, turned the travelers into her \textbf{vampire spawn}. The Myriad has a deal with the vampire and sends victims her way in exchange for help with various unsavory tasks. If the characters destroy the vampire, they make an enemy of the Myriad.

\section{Greying Wildlands}
Pressed against the northeastern coast of Wynandir is a desert of ice and snow, separated from the Dwendalian Empire by the jagged Dunrock Mountains and the Quannah Breach. This landscape of dense, desiccated forests, rolling mountains, and icy tundra is known as the Greying Wildlands, the southern part of the Biting North that encompasses this region as well as the islands of Eiselcross. The empire has struggled to expand into the region, leaving its denizens to focus on their own survival in the harsh territory, where wandering outlanders and barbarous hunters stalk the unchecked wilderness. Tucked beyond the edge of the empire's reach, past the shadowed dangers of the dense, frozen foliage, lies the dwarven citadel of Uthodurn, which has become the home of many elves who survived the destruction of Molaesmyr.

\subsection{Areas (A-M)}
\subsubsection{Blooming Grove}
Nearly ten miles northwest of Shadycreek Run, deep in the Savalirwood, a centuries-old willow grove and small temple have held back the curse that darkens these woods. The Clay family, heralds to the Wild Mother and the Arch Heart, built this holy place as a funeral grove of renewal, designed to return the bodies of the fallen to nature while letting their spirits move on. Each bloodline that is interred among the swampy ponds and vibrant leaves nurtures a wholly unique flower, giving the location the macabre nickname of the Bone Orchard.

Time does not stand still, however. The curse of the wood has eaten away at the holy boundaries of the grove. Dark vines and thorns reclaim layer on layer of fences and grave markers, pushing the once-expansive grove into a small patch of defiant green at the base of the temple. Only one of the keepers remained to watch over the decay of their grove, and now this priest, named Caduceus, has left in the hopes of discovering a way to reclaim the orchard and the ancestry it protects in the name of Melora.

\subsubsection{Boroftkrah}
\begin{description}
\item[Population.]
3,060 (70\% orcs, 17\% half-orcs, 13\% other races)

\item[Government.]
The town is guided by a circle of elder warriors known as the Runemarked.

\item[Defense.]
A force of hunter-warriors defends the settlement, but every able-bodied citizen is trained in combat to defend Boroftkrah.

\item[Commerce.]
Basic everyday supplies are available, as well as simple crafts and weapons.

\item[Organizations.]
A shrine to Kord sits in the center of the town, but no organizations, religious or otherwise, exist in Boroftkrah.
\end{description}
\noindent
After the orc's creator god Gruumsh was banished in the Divergence, many of the war clans of his bloody army fled and went into hiding. Splintered from the same throng of Odakar orcs that held the Ashkeeper Peaks, the leaderless droves migrated northward to the Rime Plains. Although they were severed from the direct influence of the Ruiner, the directionless orcs began to succumb to the Ruiner's "gift" of bloodlust. A schism tore the clans asunder, and from the bloodshed emerged a new clan of orcish marauders known as the Jez-Araz.

The remaining champions of the Odakar clan, determined to survive, found inspiration in Kord, the Storm Lord they had once fought against. They formed an organized clan of warrior orcs who seek to overcome the Ruiner's murderous influence and find strength within themselves. From this new path, the Boroftkrah orc clan and the settlement of Boroftkrah was born, creating a strong community that values deeds done in service to the clan. The orcs of Boroftkrah wander the surrounding plains in heavy furs and leathers, seeking ways to commune with Kord through challenges of survival and contests of hunting and tracking skill. While the clan is overwhelmingly of orcish descent, on rare occasions they welcome outsiders who prove themselves strong in the eyes of Kord.

\subparagraph{Surmounting the Curse}
Like many of their southern brethren, the orcs of Boroftkrah continue to eschew the hunger for violence that Gruumsh supposedly instilled in the orcish people when he created them. They take pride in victory and the success of the hunt, making great offerings of powerful prey to the Storm Lord. They channel their desire for war into battling the harsh lands that surround them to ensure the survival of their community. While most Boroftkrah orcs choose to shed their shackles of murderous rage, they remain territorial, often fighting their ferocious Jez-Araz cousins or coming into conflict with hunting parties from Uthodurn.

\subparagraph{Government}
The settlement is guided by a circle of elder warriors who have proven themselves the strongest and most dedicated to the survival of the clan. If the circle has a vacancy, volunteers prove their worthiness in a series of contests before the central idol of Kord, ending with a rune-casting ceremony that invokes Kord's will toward who will become an elder warrior. Under these elder warriors, all others are equal and trained to take up arms if the settlement is in danger.

\subparagraph{Crime}
Crime is not common within Boroftkrah. A strong sense of honor often prevents the urge for criminal activity. It is rare for anyone within the clan to steal or take from others, unless both parties agree to a challenge to see who is worthy. Those that do steal are shamed as unworthy and exiled from the clan to wander the Greying Wildlands alone.

\subparagraph{Geography}
The settlement of Boroftkrah is a loose collection of tents and shelters surrounded by a tall fence of wooden pikes. The town is created to be mobile if serious danger threatens its current location. Neighborhoods are formed in large circles of homesteads called lods. Giant boars and powerful worgs of the plains are trained as mounts and kept in scattered holding pens.

\subparagraph{Boroftkrah Adventures}
Navigating Boroftkrah is its own adventure. To gain the trust of the orcs, the characters must prove themselves in contests of strength.

\subparagraph{Storm Celebration (Low Level)}
When lightning splits the sky, the Boroftkrah orcs honor Kord with competitive games that last until the storm breaks. Characters who have earned the orcs' trust might find themselves volunteered as participants or judges during such a time. Meanwhile, an \textbf{orc eye of Gruumsh} leads a band of \textbf{orc} and \textbf{orog} warriors who have secretly begun worshiping the Ruiner. They plan to display their faith during the contest by turning the games into a bloodbath, hoping to convert other orcs with the promise of glorious bloodshed, and murder any who resist.

\subsubsection{Cinderrest Sanctum}
Hidden along the eastern slopes of Kravaraad is a cavern of onyx, where a pool of uncooling lava sits at the heart of the black stone grotto. Within this spacious chamber, black iron adornments and supports bear tales etched in runes. Carved reliefs depict dwarven interpretations of Melora granting stone and ore from below the ground to Moradin the All-Hammer, for him to craft into great halls and mighty guardians.

This shrine has been watched over by the Dust family for generations, maintained by the offerings delivered by the reverent of Uthodurn. Those who make the pilgrimage to Cinderrest Sanctum are invited to pray to Melora and Moradin, while perhaps gaining access to the Underforge, a great forge built over many decades as tribute to the partnership between the two deities. With their divine blessings, it is said that powerful artifacts with great destinies are crafted.

\subsubsection{Crystalsands Tundra}
During the Age of Arcanum, this land was a vibrant continuation of the Veluthil Forest. However, stories speak of a great battle during the Calamity that violently dashed this region into rock and sand. Winter quickly reclaimed the new desert, the sands freezing over with northern ice and snow. Tall dunes glitter like piles of diamonds in the daylight, shifting with high winds and the movements of the remorhazes and frost worms that burrow and hunt beneath the sands. Nomadic clans of northern goblins form hunting parties and ride winter wolves, stalking the furred beasts that call these frosted dunes home. Other beings stay hidden to avoid the giant hawks and giant owls that fly over from the Flotket Alps in search of prey. In times of fresh snowfall, travel is extremely dangerous, since the soft powder obscures the shifting sands and ice-marked tunnels that lead to the gnashing jaws of the subterranean frost worms. Some travelers have speculated that, depending on the winds, great structures of forgotten epochs become partially revealed as the sands and frost are temporarily scattered, only to be swallowed up once more as the dunes shift in the next storm.

\subparagraph{Worshipers of the Worm}
A tribe of wildfolk who worship Quajath are scattered across the northern fields of the Crystalsands, calling themselves the Mawcotters. They hunt by digging tunnels and caves in the sands and nearby rock, then lying in wait to spring when their quarry falls into the trap. These fanatical people make an annual journey across the treacherous Frigid Depths to the Eiselcross island of Gelier to mine flesh from the frozen body of Quajath. This meat is then brought back to the Crystalsands and ritualistically devoured by the leaders of these worm-cults, and by the young as a rite of adulthood. Fortunately for them, the flesh travels long enough before consumption to avoid the full wormkin transformation (see "\textit{Tomb of the Worm}" later in this chapter). The Mawcotters eagerly seek a way to one day free their patron from its icy prison.

\subparagraph{Crystalsands Tundra Adventures}
Traversing these endless fields of frozen dust is a thrilling and dangerous journey for any adventurer.

\subparagraph{Fair-Weather Friends (Mid Level)}
While crossing the wind-swept snow dunes, the characters encounter three young remorhazes that do not attack. They are instead following Shwells (chaotic neutral, female, gnome \textbf{scout}), who wields a recently scavenged pale, blue orb that keeps the creatures docile. She will not part with the orb and has grown attached to her monstrous companions, but characters with arcane knowledge will notice the orb is drawing more remorhazes to it from all across the tundra, and the magic that subdues the creatures appears to be rapidly fading.

\subsubsection{Flotket Alps}
Dominating the central region between the Crystalsands Tundra and the Rime Plains, the range of mountains known as the Flotket Alps is a majestic, snow-covered land of scattered pine trees and rough slopes that rise to dizzyingly sharp peaks. Nomadic trolls and roving wyverns claim the caves and alcoves of the mountainsides, hunting other inhabitants who've also adapted to the arctic weather. This alpine expanse is often battered by harsh blizzards, pushing some denizens underground into the maze of ice caverns that wind throughout the range, abandoned by some long-forgotten society and turned into the home of several warring frost giant clans. Among the natural chaos of the alps, the subterranean city of Uthodurn and its people seek to tame the wild crags around them and spread beyond the boundaries of this frozen highland.

\subparagraph{Flotket Alps Adventures}
Icy spires and snowy paths wind though these mountains, leading travelers who survive the Alps to new challenges and rewards.

\subparagraph{A Deserving Dare (High Level)}
Hidden among the shallow chasms of ice is the lair of Yondis the Winter Star, an \textbf{androsphinx} with pure white fur. It watches over the tomb of a heroic group of champions who sacrificed themselves in the final battles of the Calamity. Those who wish to claim the arms and instruments of these legendary figures must prove their worth to Yondis through numerous quests or challenges—and the sphinx has no mercy for those who fail.

\subsubsection{Ivory Lake}
Fed by the annual snow melt and heavy rains that wash across the Flotket Alps, Ivory Lake is the largest source of standing fresh water in the region. Many creatures throughout the Rime Plains make their way to the lake during the summer when the waters are thawed and the fish are plentiful. The banks of the lake are dangerous for smaller beasts and humanoids alike, as wyverns and other predators lie in wait to gorge on unsuspecting travelers.

\subsubsection{Kravaraad}
Black, burning, and defiant against the frost and cold of the north, the volcano of Kravaraad looms in the heart of the Flotket Alps. Perpetual streaks of dark smoke rise up to blend with heavy steam from the trails of molten rock that cool against the surrounding ice. The mountain is known for a continuous, gradual flow of glowing lava from multiple vents, but massive eruptions are extremely rare. The heat of the lava allows a pair of adjacent freshwater lakes, known as the Mistpools, to provide liquid water year-round. Its dramatic presence among the pale, snow-powdered trails of the Flotkets makes Kravaraad an important landmark for travelers in the snowy region. Hardier creatures live in the surrounding snowmelt, often competing for territory against the fire giants who claim the northern side and interior of the volcano. Many such battles echo through the mountains like distant thunder, calling other denizens to come witness the struggle, or to wait and finish off the survivors.

\subparagraph{Kravaraad Adventures}
The volcano attracts creatures dangerous enough to make traveling near Kravaraad an adventure of its own.

\subparagraph{Vault of the Fire Giants (High Level)}
Kravaraad's fire giants have greedily amassed treasure since the time before the Calamity. To reach their subterranean vault, the characters must sneak by or carve through the latest creatures who have challenged the fire giants (often frost giants or other creatures that can endure the cold). But making it to the fire giants is only half the battle. The fire giants have hewn a vast dungeon under the mountain, and characters must navigate its perils while maintaining their balance whenever Kravaraad causes the entire complex to tremble. The vault itself is guarded by an enslaved \textbf{adult red dragon} that longs to be free.

\subsubsection{Mythburrow}
In the farthest reaches of the northeast Greying Wildlands, tucked low near the base of the Flotket Alps, the ominous, icicle-wreathed mouth of a deep mountain cavern plunges into the heart of the rocky range. The locals who call this place Mythburrow tell stories of a monstrous, ancient being that claims dominion over the surrounding lands. Tales of a bellowing roar that shakes the valley have been told for centuries, and some even claim to have seen the frightful wingspan of the immense white-scaled dragon that calls Mythburrow home. While many joke about the siren call of a dragon's hoard, few have the mettle to venture into this cavern in search of wealth and glory—and those who did enter the cave never returned.

\subparagraph{Mythburrow Adventures}
Deep within the Mythburrow, something old and terrible shifts in seclusion, plotting its next hunt in the lands above.

\subparagraph{Gelidon (High Level)}
All trade and communication from Palebank Village (described below) halts suddenly and without warning. When the characters are sent to investigate, they find the entire area transformed into a frozen nightmare, the village leveled and the remaining corpses locked within jagged ice. Uthodurnian scouts report sightings of the long-dormant \textbf{ancient white dragon} Gelidon returning to its lair in Mythburrow.

The Diarchy of Uthodurn puts out the call to all monster hunters and heroes, setting a handsome bounty on the head of Gelidon. The characters must now clamber down the twisting tunnels of the Mythburrow and destroy the powerful wyrm before their competition can claim the gold and glory.

\subsection{Areas (P-Z)}
\subsubsection{Palebank Village}
\begin{description}
\item[Population.]
690 (61\% dwarves, 32\% elves, 3\% gnomes, 4\% other races)

\item[Government.]
The elf ranger Elro Aldataur is the appointed leader of the village. He delegates tasks to proven friends and trustworthy allies.

\item[Defense.]
Elro commands a small force of errant Glassblades. Some civilians are willing to take up arms to defend the city in moments of crisis.

\item[Commerce.]
Basic supplies can be purchased or traded at the small inns or supply shacks.

\item[Organizations.]
There are two humble temples to Moradin and Corellon, respectively. Uthodurn governs from afar, while the Tribes of Shadycreek send agents to seek more information about Eiselcross.
\end{description}
\noindent
Along the cold shore of the Frigid Depths, Palebank Village marks the first surviving Uthodurn outpost beyond the mountains. The settlement is situated against the low cliff boundary of the Crystalsands Tundra, just north of the Flotket Alps. The denizens of Palebank often take to ice fishing, trapping, or hunting, returning in small caravans to Uthodurn to sell their wares and restock their supplies. Territorial monsters wander close to the village, so rotating squads of Glassblades are assigned to protect the people. The small docks have been recently expanded to accommodate the increased interest in heading northwest toward Eiselcross. Palebank has become a launching point for northern expeditions, leading to growing trade and a pressing need for more inns.

\subparagraph{Government}
The village is guided by Elro Aldataur, a weathered elf ranger who helped found Palebank roughly sixty years ago. Stoic and humorless, Elro showed a knack for leadership when the fledgling outpost was beset with danger on a near-constant basis. Elro earned the trust and respect of the residents, and they asked him to govern Palebank. He begrudgingly accepted, and has since overseen the expansion and protection of the village. Elro generally has sixty to ninety Glassblades at his command.

\subparagraph{Crime}
The people of Palebank must work together to survive in this frigid land, and there is little room in the small community to conduct any worthwhile larceny or immoral activity. What crimes do occur are born out of necessity and punished accordingly by Elro.

\subparagraph{Geography}
Surrounded by four watchtowers and a ten-foot-tall palisade of sharpened logs, Palebank is a ramshackle village locked in eternal winter. The village is made up of over a hundred cabins and shacks. It's too small to have distinct neighborhoods, so most businesses and residences were established wherever they could fit at the time, creating a meandering layout. Off the northern cliff, about fifteen feet down to the sand-and-snow beach, the docks hold around a dozen ships, ranging from fishing vessels to small cargo ships that carry expeditions to and from Eiselcross.

\subparagraph{Palebank Village Adventures}
Palebank Village would be a sleepy place without much excitement were it not for the presence of eager explorers and the settlement's importance as a waystation between Uthodurn and the islands of Eiselcross.

\subparagraph{Troll Night (Mid Level)}
While the characters stay in Palebank Village, the only survivor of an unlucky Glassblades patrol staggers into the village. An organized band of trolls is coming down from the Flotket Alps toward the settlement, which needs every able-bodied warrior to keep the trolls at bay. How did these trolls get so organized, and what do they seek in the village?

\subsubsection{Rime Plains}
These fields of rolling tundra and bursts of boreal forest stretch on and on, covering the southeastern Greying Wildlands. Wild muskox and snow foxes wander across the rocky grassland, as wolves, saber-toothed tigers, and other beasts hunt to survive the freezing nights. Humans have not yet encroached on this region, and the people of Uthodurn abandoned their expansion into the plains, so wayfaring orcs and goblinoid bands have claimed areas of the region unchallenged. The Rime Plains are a mostly open, untamed wilderness where survival depends on one's ability to endure the elements and navigate the wilds skillfully.

\subparagraph{Roving Savagery}
While the Boroftkrah orcs managed to temper their fury through their worship of the Storm Lord, a number of their kin embraced their violent past and abandoned the rest to seek greater conquests. Now known as the Jez-Araz, or "Breakers of Spirits," these violent orcs live in scattered, nomadic groups that hunt around the plains and the lower Flotket Alps. Praising Gruumsh through acts of brutality, the Jez-Araz hunting parties are a fearsome sight for any wanderer to behold. of Molaesmyr

Just north of the Boreal Omen River, toward the center of the warped Savalirwood forest, looms the ravaged elven city of Molaesmyr. Once a beautiful city filled with erudite scholars and designed to blend in with the towering trees by masterful elven architects, the metropolis now lies abandoned. The canopy of trees above the ruins is thick enough to block most of the sun, and massive, twisting bramble vines choke the streets and rooftops, spreading into the surrounding woods like veins from a corrupted heart. Crumbling spires lean against trees swollen with disease, halls of learning are smothered by toxic fungus and moss, and a perpetual heavy fog obscures the dangers that lurk in unseen alcoves. All manner of terrible fey creatures now call this city home. The rest of the woods are filled with natural forest denizens mutated by errant magics as well as many elves who were caught in the horrifying catastrophe and who now stalk the shadows as terrible monstrosities, mockeries of their former selves.

The ruins are considered the source of the corrupt magics that still spread throughout the Savalirwood, twisting its inhabitants and keeping the elves of Molaesmyr from their home for over two hundred and fifty years. Many factions throughout the Greying Wildlands have their eyes on the ancient city, eager to uncover its secrets and its treasures. Expeditions from Uthodurn into the blighted city delve beyond the outer edges, seeking the central tower of Caes Mosor and the labyrinth beneath, where it is believed that the dark magic was unleashed. The Uthodurnian expeditions hope to uncover relics and clues that could be the key to reversing this desolation, but these dungeoneers now contend with not only the hungry monsters, but also hired mercenaries and greedy thieves from Shadycreek Run who seek to steal these same secrets for profit. The blood of both factions has been spilled in the streets of Molaesmyr, and the dark spirits of the city seem to relish the violence.

\subparagraph{Ruins of Molaesmyr Adventures}
The ruins of Molaesmyr can tempt the characters into their dangers with the promise of forgotten treasures and magic.

\subparagraph{Cackles in the Dark (Low Level)}
Auntie Rottongue, Nanny Filthtouch, and Granny Pustreat are a coven of green hags living in the labyrinth beneath the ruins of Molaesmyr. The hags send their minions to steal children from Shadycreek Run, then use them as test subjects for new, dark rituals. The minions are supposed to take children who will not be missed, but when they kidnap a Mardoon child, the characters are hired to delve into the ruins and recover the victim while facing the hags' corrupted fey minions.

\subsubsection{Savalirwood}
This area was once a beautiful thicket known as the Veluthil Forest, enchanted with fey magics and guarded by elven outposts. Nearly three centuries ago, the forest suffered a terrible curse wrought from the disaster that felled Molaesmyr. In the centuries since, the fey magics have grown twisted and dark, the denizens corrupt and feral. Spirits and monsters crawl up from under the cold, rotted forest floor, while shadowy whispers bring nightmares to those who sleep beneath the sinister trees. The forest has been renamed the Savalirwood, for the elven word \textit{savalir}, meaning "guilt." It has become a dangerous obstacle for travel, playing a major role in the Dwendalian Empire's inability to expand northward to plunder the ruins of Molaesmyr or conquer Uthodurn. In this most recent century, the less-dangerous outskirts of the forest have become havens for criminals, leading to the establishment of Shadycreek Run.

Throughout the Savalirwood, a handful of small groves and clearings combat the magical miasma that infects the rest of the wood. Within many of these relatively protected locales, wild magic thrives and gentler (if slightly mad) fey guard the surrounding glades, often enthralling wandering humanoids and capturing them to keep the fey company—forever. A community of reserved, thick-furred tabaxi called the Nightback Clan lives among the dangers of the Savalirwood, occasionally coming to the aid of others in need, seeking allies in their struggle to survive.

\subparagraph{The Nightback Clan}
Many denizens of the Veluthil were slain following the destruction of Molaesmyr and the spread of the mysterious curse that transformed the wood, but one group of tabaxi from the southern reaches of the forest survived and refused to flee. Knowing that the spirits of their ancestors were bound to the trees, they chose to stay and adapt, learning to live a nomadic life among the perpetual threats of the forest. Devoted to Melora, the Nightback Clan patrols the forest for food and locations to briefly make home. Unpleasant encounters with folk from Shadycreek Run have left some of the Nightback Clan eager to terrorize any unlucky thieves found in the Savalirwood. The Nightback Clan also blames the elves of Molaesmyr for the tragedy that claimed the forest and refuses to speak to or work with any representatives of Uthodurn.

\subparagraph{Savalirwood Adventures}
Every journey through the dangerous, weird, and corrupt Savalirwood is an adventure.

\subparagraph{Nightback Negotiation (Low Level)}
The Nightback Clan captured a merchant caravan traveling through the Savalirwood, suspecting that the travelers are thieves. The characters are hired by the merchant's family to rescue the caravan. If they can track the tabaxi through the woods, they must either convince the Nightback Clan that the travelers are innocent or fight to free the captives. Things grow more complicated when the tabaxi present compelling evidence that this merchant caravan is more than just a group of simple travelers.

\subparagraph{Treantwatch (Mid Level)}
While traveling through the forest, the characters encounter a \textbf{treant} who insists on guiding them safely through the Savalirwood. The treant claims to know every nook and cranny in the forest, as well as its dangers. On the surface, it seems very helpful. However, the curse afflicting the forest has driven the treant mad and made it delusional. All the dangers it describes are inventions of its twisted mind, and it has a knack for leading the characters into one perilous situation after another.

The treant's refusal to divulge its name is the first clue that something is amiss. In truth, the treant no longer remembers its name.

\subparagraph{Mapping the Gates (High Level)}
Erma Schnieb (chaotic neutral lightfoot halfling \textbf{mage}) believes that the curse infecting the Savalirwood creates portals to the Lower Planes. She hires the characters to explore the forest with her to find and map these gates. In addition to the usual nightmares that stalk the forest paths, fiends patrol the gates and attack anyone who tries to close the planar doors.

\subsubsection{Shadycreek Run}
\begin{description}
\item[Population.]
14,770 (56\% humans, 15\% elves, 14\% dwarves, 15\% other races)

\item[Government.]
Four organized crime families and the citizen-run Greytrader Union vie for control of the city, while the mercenary peacekeepers of the Grudge Gang attempt to prevent complete chaos.

\item[Defense.]
The Tribes of Shadycreek Run each employ personal soldiers and bodyguards, while the Grudge Gang defends the city at large.

\item[Commerce.]
Most general goods are available, and rare imports and stolen commodities are sold and auctioned in secret. Many taverns, inns, and places of vice are open to the public.

\item[Organizations.]
Most business is run by one of the Tribes or the Union, even if alliances are hidden. The city has scattered, humble temples.
\end{description}
\noindent
Past the northern outposts of Wynandir, a scattered line of shantytowns and hovels chokes the pathway into the trees of the wildlands alongside a small creek. Originally comprised of a patchwork of makeshift, impoverished societies run by cruel slumlords, Shadycreek Run is now known as a hive of criminal activity and savage humanity given free reign. The degenerate dregs of civilization see this place as a chance to gain untold fame and power. Most debates end in a brawl, and most brawls end in bloodshed.

The harsh, unguarded environment breeds a tough people, and despite the ravaging beasts that pick off easy prey from the outskirts, proud groups of hunters and skinners have risen to become prominent traders. Some intelligent figures have begun to organize and civilize life in the Run, hoping to bring a modicum of order and respect to the collection of reprobates, while the veterans of the Run fight to keep things unpredictable, undermining these attempts with violence. The volatile society underscores the age-old adage, "they call it the Run for a reason."

\subparagraph{Dangerous Vocations}
Many who choose to live within Shadycreek Run do so to avoid the persecution (or prosecution) that awaits within the empire, while others have a taste for the brutal lifestyle afforded by a land of overt moral flexibility. Professions here range from artisans and craftsman to farmers, mercenaries, and entertainers. However, the most profitable entrepreneurs in the region are thieves, hunters, and scavengers who brave the deadly Savalirwood and the ruins of Molaesmyr. Powerful beasts yield hearty meat and furs, while recovered relics and lost magics fetch quite the price in distant markets. The struggle over these prizes often grows bloody, as the Tribes and the Greytrader Union battle to gain control of these trades, while also competing with the people of Uthodurn who wish to reclaim their heritage and lost belongings.

\subparagraph{Adversarial Rulership}
Currently, Shadycreek Run is divided between three competitive factions: the Greytrader Union, the Tribes, and the Grudge Gang.

The Greytrader Union is a faction of reformed outcasts and criminals who want to establish firm trade relationships with outside societies and help salvage Shadycreek Run by making it a safer, more functional community. The members of the union don't abide imperial rule, and they enjoy the freedom of a lawless lifestyle, but they also wish to be free of the legacy of violence for which Shadycreek Run is known. Made up of trappers, hunters, builders, and craftsmen from all walks of life, the Greytrader Union wishes to institute self-appointed marshals to keep the peace. They aim to increase the safety of civilians, turn the Run into a respectable destination for trade, and eventually diminish the influence of the Tribes and the Grudge Gang.

The Tribes are a collective of rival criminal families who ran the Run for nearly a century before they lost much of their dominance to violent infighting and internal power struggles. Comprised of the Mardoon, Uttolot, Trebain, and Jagentoth families, the Tribes work to maintain control over most of the business that transpires in the city, and conspire against each other in secret. While each family center is usually composed of relatives by blood or marriage, their employees are also considered "within the family," so long as they remain loyal. The Tribes fence stolen goods, provide vice and bawdy entertainment, and control the populace through confusion, intimidation, and fear.

The Grudge Gang is a hardened company of mercenaries who act as the muscle to enforce local interests within the Run. Founded by a band of deserters from the Dwendalian army, they've grown into a battalion of warriors loyal to coin and each other. Members of the gang will happily work for the Greytrader Union or any of the families of the Tribes, so long as their business isn't bloody enough to harm innocents or put the safety of the Run at stake. Since the other factions rely on the Grudge Gang for protection from the terrors of the Savalirwood, its members consider themselves neutral instruments in the political chess game. Some members of the gang have assigned themselves to maintain order in a land without law, calling themselves Taskers and doing their best to keep the innocent from being crushed by the city's predators. Their priorities revolve around lending their skills to whoever pays the best, maintaining livable streets, and keeping the dangers of the Savalirwood and the Dwendalian Empire out of the city.

\subparagraph{Government}
The city has no central government or laws, so much of it is held together by a vague sense of "live and let live," broken by occasional bursts of "take what I want." If violence spirals out of control, the Grudge Gang's Taskers step in and beat everyone involved to a pulp. It's not uncommon for careless folk to wind up dead in a snow-covered alley, stripped of their goods and left to be eaten by starving dogs. Denizens of the Run learn to watch their backs and warily wield blades.

\subparagraph{Geography}
What began as an extended string of shanties and shacks built along the southern treeline of the Savalirwood has turned into a three-mile-long, crooked main street called Clover Alley, lined with hovels, buildings, and mudholes. The neighborhoods on the southern side of the Alley, called the South Clover, are the most impoverished and least maintained areas of the city, known for extremely desperate folk and even more dangerous business. The northern side, called the North Clover, is partially built into the Savalirwood itself, and houses the wealthy and powerful denizens of the Run, including most of the families of the Tribes, who live in various mansions and small strongholds.

The western end of the Clover Alley, known as the Wickerskids, is home to the scavengers and farmers who do their best to coax the harsh lands of the cold north to produce food for the people of Shadycreek Run. The central square of Clover Alley is named Clover Plaza and marks the heart of commerce, conflict, and vice within the city. In the plaza, all factions mingle in broad daylight, disguising their hidden agendas with a bawdy and loud lifestyle of indulgence and excess. The eastern end of Clover Alley, called the Gruff Slum, houses a cluster of peasants and cutpurses who live among drug dens and pig pens.

\subparagraph{Shadycreek Run Adventures}
Slack laws and the promise of easy money draw many opportunists to these dangerous streets.

\subparagraph{A Stolen Heart (Low Level)}
At the Trench Tavern, the characters join a drinking game with some of the local Grudge Gang. Win or lose, an \textbf{orc} member of the gang, Dratto, reveals their connection with the Mardoon family. Dratto mentions that they need some hands to steal a Molaesmyr relic currently held by the Trebain family. The characters can turn Dratto into the Trebains to gain their favor or agree to aid in the theft.

If the characters agree to help Dratto, a dangerous burglary ensues. They must infiltrate a guarded shed outside the Wickerskids, only to discover the relic is actually a slumbering \textbf{pixie}. This creature wakes in the presence of the characters, when a mysterious aspect of one character's nature draws the pixie from out of its torpor and binds its affection to the character.

\subparagraph{King's Spy (Any Level)}
Nella Durlon (lawful neutral, nonbinary, human \textbf{spy}) lives in Shadycreek Run as a spy for the empire. They gather information about the settlement's defenses to take back to Oliver Schrieber (see "\textit{Dwendalian Empire}" in chapter 8). The characters learn of Nella's job from a faction that hires them to find the spy. Depending on the faction, the characters must either escort Nella safely out of the settlement so they can report to Oliver, or ask that Nella be killed for giving away the settlement's secrets. Either way, finding Nella is no small task. The spy is a master of their craft and blends in with the criminals and ruffians of the Run.

\subsubsection{Uraliss}
The small town of Uraliss was Uthodurn's first attempt to expand beyond its mountain halls. Uraliss was settled with the intent of becoming a farming village to aid in crop production. As the town began construction, arcane wards were developed and placed along a patch of snow-covered field at the base of the Flotket Alps. The wards were designed to slowly melt the snow and hopefully help develop the area into farmable land. However, the town was beset by orcs and other threats from the Rime Plains, and the town's hopes were further crushed when the land beneath the snow was found to be rocky and untillable. The settlement was quickly abandoned and is now a crumbling ghost town left to be vandalized and repurposed by the Boroftkrah orcs of the Rime Plains.

\subparagraph{Uraliss Adventures}
Characters might go to Uraliss to retrieve something left behind, study its arcane wards, or hunt a beast that lurks in the ghost town.

\subparagraph{Haunted Temple of the All-Hammer (Low Level)}
The Boroftkrah orcs avoid Uraliss's temple of Moradin at all costs. The temple is filled with specters of dwarves and elves captured by a high priest who went mad and locked her congregation in the temple during the final orc raid. All her victims starved to death, as did the priest herself, who became a \textbf{wraith}. The characters could come to this building to find a lost treasure of Moradin, or to escape other threats as they pass through the ruined town.

\subsubsection{Uthodurn}
\begin{description}
\item[Population.]
26,240 (56\% dwarves, 36\% elves, 4\% gnomes, 4\% other races)

\item[Government.]
A diarchy between the dwarven and elven monarchs rules over the populace, each leader supported by a council of nobles who share the responsibilities of overseeing the city and its functions.

\item[Defense.]
Uthodurn maintains a well-trained army of Glassblades to keep the peace and stations scouts and skirmishers across the Flotket Alps. Nobles also employ bodyguards.

\item[Commerce.]
Most basic goods and services are available. Uthodurn's isolation means that imports are limited and inns for outsiders are scarce.

\item[Organizations.]
Numerous temples dot the Volition Disc area of the city, and competing guilds do business under the oversight of the councils. Hidden agents from the Tribes of Shadycreek Run have begun to find their way into the city.
\end{description}
\noindent
Hidden among the icy peaks of the Flotket Alps, the northernmost dwarven stronghold of Wildemount is ensconced in the base of a steep mountain valley. Uthodurn is a bastion of dwarven culture, home to a secluded people who've endured the freezing temperatures and continuous snowfall across the mountains and surrounding lands since they founded the stronghold nearly eight hundred years ago. The dwarves of Uthodurn are sharp of wit, hairier than their southern brethren, and as unyielding as the rock they live in. The Queen of Uthodurn broke from stubborn dwarven traditions and opened the city's doors to thousands of elven refugees following the destruction of Molaesmyr nearly three centuries ago. While there have been moments of conflict between the dwarves and elves as ancient grudges are briefly revived, the society within these fire-warmed stone halls has grown into a unique union of both cultures.

\subparagraph{A Harsh Environment}
Beyond the iron doors of the city, the wildlands surrounding the underground stronghold are brutally cold and battered by cutting winds. The people of Uthodurn hunt and track in bursts, like sudden flurries of snow, stalking the wild white bison and cliffneck goats that call the Flotket Alps home. The dwarves and elves of Uthodurn have adapted to their environment, becoming skilled hunters and trackers as well as hearty travelers used to persisting through extreme conditions. As the hunter's saying goes, "squeezing water from a stone might be hard, but snaring a rabbit from the Flots is harder." A strong martial force is necessary to combat dangers such as frost giant clans and snow-faring goblinoid bands that hunt throughout the slopes and caverns of the mountain range, and the masters of Uthodurn have trained fierce warriors and spellslingers.

\subparagraph{Dual Societies, Dual Skills}
Uthodurn has the unique distinction of being the only place in Wildemount where dwarven and elven lifestyles and traditions come together within one location. The resulting collaborations have created incredible new crafts and means of circumventing the challenges of the region. Dwarven carving techniques blend with elven arcane knowledge to elevate architecture and excavation capabilities. Historical knowledge of the local vegetation meets refined herbalism to cultivate and grow cold-resistant crops. The crafting and smithing techniques of each lineage entwine and complement one another, making Uthodurnian goods that pass beyond the Greying Wildlands valuable indeed.

\subparagraph{Love Knows No Bounds}
Dwarves and elves have lived together under the same roof and banner for centuries within Uthodurn. Over those years, they have forged ties of mutual respect, and romances and unions between elves and dwarves have also become increasingly common as the years progress. Though such partnerships were initially considered scandalous, they came to be met with tolerance, and eventually celebration. While an elf and dwarf may indeed produce offspring, the curiously stubborn nature of both ancestries produces a child that is either an elf or a dwarf, though a child of either race often inherits minor physical attributes from the other parent. Such children are affectionately referred to as attalwen, or "two-hearted."

\subparagraph{Strike It Rich}
The Diarchy of Uthodurn rewards those who discover magic items in the wastes of Foren, driving Uthodurn's most experienced—and most desperate—explorers to Syrinlya. Uthodurn doesn't have the funds to support a full expedition in Foren, so they provide incentives for citizens who gather such items. The city provides free transportation to and from Syrinlya, in addition to rewarding discoveries, because they are desperate for items that could provide them with more resources, guard them from an invasion from the Dwendalian Empire or the Kryn Dynasty, or help cure the corruption in the Savalirwood. A handful of scholars believe the source of the Savalirwood's corruption might even be tied to something within the ruins of Aeor.

\subparagraph{Government}
The stronghold of Uthodurn is ruled by the dual monarchs, the dwarf Queen Simone Fruunast and the elf King Imathan Talviel, who each oversee their own small council. Most matters of the city and its provinces are handled by these councils, and well-trained soldiers and rangers called Glassblades—named for their signature weapon design, which resembles jagged glass—act as enforcers of order and law.

\subparagraph{Crime}
Uthodurn society is far from perfect, and a seedy criminal element exists in the shadows under the mountain. Guildmasters who seek to surmount their rivals are known to resort to theft and bribery, while some honorless nobles have occasionally embraced the effectiveness of blackmail, and even murder, to gain power and wealth. Curious spies and thieves from Shadycreek Run have also sneaked or lied their way into the hearth-lit halls, and a small surge of crime might be imminent.

\subparagraph{Geography}
Above the entrance to Uthodurn stands a series of small watchtowers, which surround a cluster of livery stables that house thick-coated, hardy horses and cliffneck goats. The entire city is built around a vertical pit that descends straight down from a half-valley, with numerous tunnel doorways locked with cap-like iron doors, creating a wide, downward spiral of carved stone, ruddy iron, polished marble, and cold subterranean waterfalls lit by ever-burning pyres. The city's sections are divided and marked by sets of wide pillars that hold aloft steel gates that can be dropped to seal off entrances and defend the deeper regions of Uthodurn.

The top level of the spiral, called the Volition Disc, is the widest and oldest of the regions. Many citizens live along the outer walls of the Volition Disc in a unique mishmash of old, angular stone abodes and lofty elven homesteads. The district is bisected by a central road called the Auger Trail, which is both a thoroughfare and a forum, lined by tables, benches, and peddlers' carts. Many construction businesses have small building yards here, while the upper, or Topper, mines work hard to pull iron from the surrounding rock. The main Glassblade barracks are centered around the entrance tunnels that descend into the city.

Below the Volition Disc, in a smaller spiral, lies the Deliberation Disc, or Liber Disc as it's often called. The Liber Disc is lined floor-to-ceiling with homes, emporiums, forges, and taverns that fill the air with the scents of cooked meats and boiling tubers. Many Uthodurnians come here to shop, enjoy the bustling atmosphere, or unwind with a hearty ale or piquant wine. Foreign travelers may be tickled to come across dwarves dancing to uncharacteristically whimsical elven music or elves enjoying dirty barroom games.

At the bottom of the city, the Grand Disc contains the Luddenbrock Mines that pull precious metals from the deeper veins of the mountain, as well as the Grandcast Citadel where the dual monarchs rule their people. The citadel itself is an impressive palace carved into the central stalactite pillars to resemble an inverse tower, with marble walls that make it look like a frozen, underground tornado.

\subparagraph{Uthodurn Adventures}
Travel to and from Uthodurn is its own adventure, given the severity of the surrounding environment. The Glassblades might also turn to adventurers to help stem the rising tide of crime.

\subparagraph{Awful Wails (Low Level)}
The characters are hired by the Glassblades to investigate a series of murders. Ancient elves, old enough to remember the destruction of Molaesmyr, are withering and dying in the streets at night. Witnesses all report hearing the same thing before finding the bodies: a horrible wail. The culprit is Myra Quirn, an elf who was turned into a \textbf{banshee} when she was left behind as her people fled Molaesmyr, and turned into a vengeful spirit by the dark magic of the ruins. After many years of searching, Myra found her still-living kin in Uthodurn, and she plans to exact revenge on every last elf that left her behind.

\subparagraph{Barren Shelves (Mid Level)}
Uthodurn's stock of food is running out as winter approaches. The monarchs sponsor a mammoth-hunting contest, offering a prize of 1,000 gp to the team that brings back the largest \textbf{mammoth}. If the characters wish to partake, they must hunt and kill the largest mammoth possible, then deal with a rival crew of adventurers that tries to take their kill by force.

\subparagraph{A Dreadful Bore (High Level)}
The Kryn mount a daring incursion into Uthodurn, using purple worms to bore into the subterranean city. After breaching the outer wall of the Grand Disc, Kryn invaders storm the Grandcast Citadel and hold both monarchs hostage after they're pinned down and unable to escape with their captives. The Glassblades turn to the adventurers for help freeing the monarchs and bringing the Kryn invaders to justice.

\section{Eiselcross}
Far north of the rest of Wildemount is a collection of imposing islands of ice, danger, and powerful, forgotten relics. Eiselcross and the Greying Wildlands are collectively known as the Biting North. For centuries, most of the world dismissed Eiselcross as an uninhabitable glacial land of blizzards and yetis. Only fishing vessels ever got close to the frozen shores—until one of those ships found an Aeorian item.

The flying city-state of Aeor, a place of inventors, magic, and progress, was hurled into Eiselcross by gods both good and evil during the Calamity. The Prime and Betrayer deities agreed on a momentary armistice to achieve this mammoth task, for Aeor's people crafted weapons made to kill the gods and had to be stopped. The gods tried to wholly destroy Aeor, but its protections against divine entities were carefully tuned and unexpectedly strong. It took the most of the pantheon to crash the city.

Aeor was long believed to have been obliterated, and the gods' silence on the matter only furthered this conviction, so the books of lore left nothing more to be studied or discussed on the once-powerful nation. However, recent discoveries of Aeorian artifacts revealed the truth of Aeor's resting place and spurred rumors of arcane secrets frozen beneath the ice of Eiselcross. Now Dwendalian, Uthodurnian, Xhorhasian, and mercenary expeditions search the hazardous terrain for materials that could change the tide of the war between the Dwendalian Empire and the Kryn Dynasty.

\subsection{Encounters in Eiselcross}
Eiselcross is a land of strange magic, bizarre creatures, and fascinating people. Each day that the characters travel the land, roll a d20. On a result of 15 or higher, an encounter chosen or rolled on the Eiselcross Encounters table of the appropriate level occurs. Not all of these encounters are meant to result in combat, and the ideas on the tables can inspire encounters of your own creation.

\begin{DndTable}[header=Eiselcross Encounters: Levels 1–4]{cX}

d6 & Encounter\\
1 & The characters encounter 1d4 wildfolk scouts hunting a \textbf{saber-toothed tiger}. The scouts are out of arrows and willing to trade equipment and information to replenish their ammunition.\\
2 & The characters encounter a lost Cerberus Assembly \textbf{mage} trying to return to Balenpost. The rest of the mage's team died in the ruins of Aeor, and there is a 50 percent chance the mage carries a recently recovered Aeorian item.\\
3 & The characters find a \textbf{griffon} with an arrow in its wing that prevents it from flying. A character who succeeds on a DC 15 Wisdom (Animal Handling) check can approach the griffon without being attacked. Such a character can remove the arrow, restoring the griffon's ability to fly, after which the griffon allows the character to ride it as a mount for the next 1d10 days.\\
4 & The characters encounter the \textbf{revenant} of a dead explorer from Balenpost who was murdered by one of its team and is looking for vengeance.\\
5 & The characters find a 3-foot-tall, 50-pound egg. If they keep the egg in a cold environment for 1d10 weeks, a \textbf{white dragon wyrmling} hatches from it.\\
6 & A \textbf{young remorhaz} ambushes the characters.\\
\end{DndTable}

\begin{DndTable}[header=Eiselcross Encounters: Levels 5–10]{cX}

d6 & Encounter\\
1 & The characters encounter 1d4 + 1 \textbf{mammoths} being hunted by 2d6 \textbf{wildfolk scouts}.\\
2 & An \textbf{abominable yeti} and 1d4 yetis stalk the characters.\\
3 & The characters encounter 2d4 \textbf{yetis} from Allowak's Sanctuary (see "\textit{Allowak's Sanctuary Yeti Statistics}" later in this chapter). The yetis attempt to scare the characters away but do not wish to do harm.\\
4 & The characters pass an Aeorian building that crashed nearby. The building contains 1d4 + 4 \textbf{wights}.\\
5 & An \textbf{Aeorian absorber} ambushes the characters.\\
6 & The characters encounter a large herd of elk. If the characters scare the elk, 1d6 \textbf{elk} attack while the rest of the herd flees. During the combat, a \textbf{remorhaz} bursts out of the snow and attacks indiscriminately.\\
\end{DndTable}

\begin{DndTable}[header=Eiselcross Encounters: Levels 11–16]{cX}

d6 & Encounter\\
1 & An \textbf{adult white dragon} confronts the characters, demanding treasure for passage through its territory. The dragon accepts 1,000 gp worth of treasure or a rare magic item as acceptable tribute.\\
2 & The characters encounter an \textbf{iron golem} that keeps walking in circles. The golem attacks any creatures within 10 feet of it but otherwise continues walking in circles.\\
3 & The characters find a camp of 2d4 \textbf{minotaurs} and 4d6 \textbf{orcs} looking for wildfolk to recruit for the Kryn. These explorers want to keep their operation covert. If they aren't convinced the characters can keep their secret, the Xhorhasians attack.\\
4 & A \textbf{roc} flies overhead and attempts to scoop up a character of your choice as its next meal.\\
5 & The characters encounter a coven of three \textbf{night hags} searching for Aeorian items. In their humanoid guises, these hags pretend to be lost explorers and request the characters' escort. When the characters are in a vulnerable spot, the hags reveal their true nature and attack.\\
6 & The characters are ambushed by 1d8 + 1 \textbf{frost giants}.\\
\end{DndTable}

\begin{DndTable}[header=Eiselcross Encounters: Levels 17–20]{cX}

d8 & Encounter\\
1 & The characters find the buried top of an Aeorian tower that crashed into the ice. Ice fills the lower levels of the tower, and 1d4 + 1 \textbf{Aeorian nullifiers} are frozen in the ice.\\
2 & The characters encounter 2d6 \textbf{gladiator wormkin} (see "\textit{Tomb of the Worm}" later in this chapter) attempting to force-feed a group of 3d6 \textbf{wildfolk commoners} the meat of the Undermaw.\\
3 & The characters find a humanoid skull frozen in the ice. If they free the skull, it comes alive, revealing its true form as a \textbf{demilich}.\\
4 & The characters encounter two mages from Balenpost planning to steal an Aeorian item from a camp of 3d6 frost giants and their tamed \textbf{frost worm} The mages ask for the characters' help with the heist.\\
5 & The characters encounter the \textbf{ghost} of an Aeorian citizen. If the characters are kind to the ghost, the undead might be able to share the history of the crashed city with the characters.\\
6 & The characters encounter 2d6 \textbf{frost giant zombies}. These zombies did not come from the Fortress of the Dead Jarl. Their origin is up to you.\\
7 & The characters cross over an area of thin ice. A \textbf{kraken} obsessed with Aeor floats beneath the ice and orders the characters to share any knowledge they have of the place or face its wrath.\\
8 & An \textbf{ancient white dragon} ambushes the characters.\\
\end{DndTable}

\subsection{Islands of Eiselcross}
Eiselcross is a collection of eight icy islands, some of which have small, snow-covered mountain ranges:

\subparagraph{Beissel}
Eiselcross's third largest island is so far north that even prepared explorers have trouble surviving the cold. Most don't dare brave an expedition, and those who do turn back once they see the white dragons that rule Beissel.

\subparagraph{Dampfkan}
This island is an active volcano that steams into the sky. Dampfkan's few explorers have seen a large structure in the lava, but nobody knows what this mysterious monolith is.

\subparagraph{Foren}
The largest island of Eiselcross is also its most desirable destination for treasure-seekers, since Aeor's main crash site is on Foren.

\subparagraph{Frostbogen}
Given its proximity to Foren, one might think that this island is also a popular destination for explorers. But ever since a Cerberus Assembly delegation of over one hundred people, their ship, and all their belongings vanished during an expedition to Frostbogen, no one has dared to explore the island.

\subparagraph{Gelier}
Eiselcross's second-largest island holds the body of Quajath, the Undermaw, and its wormkin (see "\textit{Tomb of the Worm}" later in this chapter).

\subparagraph{Kaltsel}
The smallest of Eiselcross's islands floats off the southeast coast of Foren and is said to be the home of Mynarc Furdahl, an undead warlock of Asmodeus willing to make deals in exchange for souls.

\subparagraph{Schneescel}
A permanent magical blizzard rages across Schneescel, a small island to the west of Foren, making it nigh impossible to explore. Many believe that a fallen Aeorian relic causes this storm.

\subparagraph{Taergoss}
Eiselcross's second smallest island is a favorite spot for fishing vessels to anchor overnight. The crews tend to stay on their boats and never set foot on shore, since faint wails can be heard at all times of day coming from Taergoss's center.

\subsection{Wildfolk of Eiselcross}
While most of the world did not know about Eiselcross's secrets for centuries, the islands are home to nomadic clans of humans and halflings. These people move in small groups, following herds of elk and mammoths. They have no gods but instead revere the animals and natural phenomena such as the wind, snow, and waves. The wildfolk avoid Aeor's ruins and the terrible dangers that lurk there.

Eiselcross's wildfolk live simple, harsh lives. They hunt, gather, and support each other, taking shelter in tents made of furs. If they have disputes with one another, they reconcile or simply go separate ways. Taking the life of another Eiselcross native for any reason is an even greater transgression than murder normally is in the other areas of Wildemount. The penalty for such a crime is a facial branding in the shape of a knife and permanent exile from any wildfolk clan.

For all their peace with each other, the wildfolk of Eiselcross are fierce warriors. Their cold environment shapes them, and they must defend against saber-toothed tigers, yetis, and even more dangerous predators such as white dragons. The wildfolk tend to be distrustful of outsiders, especially after seeing explorers from across the sea going into locations the native people dare not enter and using arcane magic to uncover the frozen secrets of their land.

\subsection{Areas (A-M)}
\subsubsection{Aeor Crash Site}
While pieces of Aeor were thrown all over Eiselcross and the Frigid Depths, the city's main crash initially created a three-mile-diameter crater in the center of Foren. In the centuries since, ice covered the crater. Ruins poke through the frosty surface, creating tunnels, pits, and chasms. The buried ruins are filled with dangerous, wonderful creations and pockets of people frozen in time.

\subparagraph{Stasis Bubbles}
Most of Aeor's streets, buildings, and people were smashed to pieces when the city collided with Foren. Amid the ruins are magic fields of translucent blue force surrounding moments frozen in time. Many of these bubbles hold diverse crowds of humanoids huddled in the street, their bodies, clothes, belongings, and the ground itself unharmed by the crash. Others contain families holding hands as they brace for impact around dining tables, mages waving wands in the air around laboratory benches, soldiers with glowing swords raised to the sky, siege crews climbing atop weapons meant to fire massive projectiles, and more. Each magic field is a 10-foot-radius sphere, and nothing within appears damaged or moved by the crash.

While these preserved scenes are sometimes useful for gathering information, the items and people within are the real prize. Totally intact Aeorian technology is a rare find, but having even one of Aeor's inventors or mages to demonstrate, repair, and rebuild their creations is a reward beyond any of the other treasures the ruins hold. However, all attempts to dispel, break, or otherwise remove the impenetrable magic barriers have failed. The people within appear frozen in time, but no evidence suggests this stasis would change if the barriers dropped, or that the change would be for the better. Some say the barriers dropping could cause the creatures within to age rapidly, turn to dust, or go mad. Despite these warnings, explorers search Aeor's ruins for an item or spell that can break these bubbles.

\subparagraph{Breaking the Bubbles}
If and how the stasis bubbles of Aeor can be broken is up to you, but here are a few suggestions for solutions:


\begin{itemize}
\item An encoded ritual hidden in the bowels of Aeor's crash site can break the bubbles. It requires one powdered unicorn horn and one dried flumph tentacle per bubble.
\item The crystal in Allowak's Sanctuary powers the stasis bubbles. If it is destroyed, the stasis bubbles break.
\item In their exalted state, the Arms of the Betrayers are powerful enough to break the bubbles with an attack.
\end{itemize}

\subparagraph{Arms and Armor}
Aeorian technology was created to battle the gods and their servants. Explorers seek weapons and armor, most of which is crafted to pierce divine defenses and guard against deific assaults. Weapons as small as daggers and as large as cannons were crafted in Aeor. Although much was destroyed in the crash, there are still items in the rubble. The pieces are valuable, and the blueprints for creating such powerful items would be priceless, though none have been found so far.

\subparagraph{Disease Laboratory}
A large facility in the city developed diseases meant to weaken and kill the gods and their servants. Aeor's mages tested the diseases on beasts and prisoners, creating horrific monsters, some of which still stalk the ruins.

\subparagraph{Frigid Woe}
Frigid woe is a special disease developed by Aeor's mages that cannot be cured by conventional treatment or magic. The only way a creature infected with the disease can be cured is by finding and drinking the manufactured antidote, a milky liquid stored in gold vials found in Eiselcross's ruins. This disease was created to slow down the forces of the gods and get around the healing power of their clerics and angels.

The disease is transmitted by breathing in blue spores that Aeor's mages created long ago. When a creature comes into contact with these spores, it must succeed on a DC 11 Constitution saving throw or become infected with frigid woe. It takes 1d4 days for the symptoms to manifest in an infected creature. These symptoms include fatigue, chills, and visible blue veins that appear on the creature's body. The infected creature's speed is reduced by 5 feet as long as it remains infected. Every 10 days after symptoms appear, an infected creature must succeed on a DC 11 Constitution saving throw, or its speed is reduced by another 5 feet. If a creature's speed is reduced to 0 as a result of this disease, the creature dies and its body turns into a statue made of ice.

A creature can drink the antidote as an action, ending all symptoms and effects of the disease instantly.

\subparagraph{Aeorian Hunters}
Aeor's people didn't just develop weapons that could kill the gods; they created soldiers. Strange experiments turned captive beasts, prisoners, and volunteers into hulking dreadnoughts called Aeorian hunters (described in \textit{chapter 7}). Those that survived the crash still stalk the ruins, driven mad by isolation and unsated bloodlust.

\subparagraph{Strange Obelisk}
A mile west of the crater that defines the crash site is a rocky hill that thaws any snow or ice that rests atop it. The mouth of a cavern can be found on the hill's northern side, which leads to a 60-foot-wide, domed chamber. In the center of this cave stands a fifteen-foot-tall obelisk made of dark, reflective metal that resembles smoothed obsidian. The stone quietly hums with an aura of abjuration from some mysterious enchantment. Travelers that complete a long rest within proximity of the stone are visited with terrible dreams of strange, alien places and visions difficult to understand.

\subparagraph{Aeor Crash Site Adventures}
The main reason for characters to explore Aeor's crash site is to uncover Aeorian magic items and other treasures, but any time spent in the region is sure to bring excitement and danger.

\subparagraph{Wild Child (Mid Level)}
An eight-year-old wildfolk child named Essora was playing with others near the crash site when a friend dared her to stand on the ice over the ruins. Essora fell through the ice and has not come out of the city's remains. None of the wildfolk are willing to risk themselves in the ruins to save Essora. The child's grandmother, Leriea, begs the characters to bring Essora home. While looking for a way out of the ruins, Essora found the enhancement workshop and drank a potion that gave her frigid woe. There is a vial of antidote deeper in the ruin.

\subparagraph{Weapons of War (High Level)}
The characters are hired by a representative of the Dwendalian Empire, the Kryn Dynasty, the Myriad, or some other faction to travel to the Aeor Crash Site and search for gold vials containing frigid woe, a milky liquid. A previous expedition found one such vial, and the characters' benefactors want more of it. The party's benefactor wants to keep the vials from falling into the wrong hands or wants to sell them to the highest bidder. To reach a hidden cache of vials in the depths of the crash site, the characters must content with Aeorian hunters and a rival search party of high-level adventurers.

\subsubsection{Allowak's Sanctuary}
\begin{description}
\item[Population.]
532 (100\% yetis)

\item[Government.]
Brightheart, the settlement's smartest and strongest abominable yeti, rules Allowak's Sanctuary as a dictator.

\item[Defense.]
Every resident of Allowak's Sanctuary is capable of defending the settlement.

\item[Commerce.]
Aside from trading for simple goods and services, there is no major commerce.

\item[Organizations.]
The yetis have no organizational ties.
\end{description}
\noindent
When Aeor fell into Foren, the impact threw a monolithic, purple crystal miles away from the crash site. The stone pierced the body of an abominable yeti named Allowak, but the crystal's magic preserved the yeti's corpse and absorbed Allowak's soul. When other yetis touched the crystal, their intellect increased and their innate desire to slaughter faded away. A society of the monstrosities living in huts and keeping mammoths for food sprung up around the crystal. These peaceful yetis abhor their more savage kin.

\subparagraph{A Secret Place}
The yetis of Allowak's Sanctuary keep their home a secret, for they worry that outsiders might try to steal the crystal. Although the crystal's effects seem permanent, the yetis fear that if it were removed from Allowak's remains, they would lose their newfound intelligence. If outsiders get close to the village, the yetis patrolling the village's borders approach and act like the more common and savage members of their species. The yetis of Allowak's Sanctuary do not kill unless forced to, but they will cause injury to guard their secret.

In recent years, the yetis have acquired books and journals from explorers who died close to their camp. The tomes gave the yetis an understanding of the world beyond Eiselcross.

\subparagraph{Brightheart's Intentions}
The yetis of Allowak's Sanctuary are led by an abominable yeti named Brightheart. She organizes adventuring parties of her people to search Foren for more of their kin to awaken with the crystal. Brightheart claims this recruitment is to increase the size and security of the community, but some whisper that she is building an army to take over Eiselcross, and possibly even more places beyond the Frigid Depths.

\subparagraph{Allowak's Sanctuary Yeti Statistics}
The yetis of Allowak's Sanctuary are yetis and abominable yetis, with these changes:


\begin{itemize}
\item Their alignments are neutral.
\item Each has an Intelligence score of 16 (+3) and a Charisma score of 10 (+0).
\item They can speak, read, and write Common and Yeti.
\end{itemize}

\subparagraph{Allowak's Sanctuary Adventures}
The characters could stumble across Allowak's Sanctuary or fight their way through the perimeter guards to uncover the village.

\subparagraph{A Matter of Trust (Mid Level)}
As the characters travel Foren, they find Dawnmind, a female Allowak yeti (see "Allowak's Sanctuary Yeti Statistics") holding the tatters of a net as she lies in the snow, covered in claw wounds from others of her kind, with 1 hit point remaining. Three other yeti corpses, Dawnmind's companions, lie in the bloody snow nearby. Dawnmind begs the characters to escort her back to Allowak's Sanctuary along with the corpses. Her group failed to capture six wild yetis, but those yetis are nearby and stalk the characters if they help Dawnmind.

If the characters travel to Allowak's Sanctuary with Dawnmind, Brightheart welcomes them as heroes and thanks them with a place to stay, provided that they swear to keep Allowak's Sanctuary a secret—or the \textbf{abominable yeti} could order the characters to prove their loyalty by capturing a group of wild yetis or mammoths for her.

\subsubsection{Balenpost}
\begin{description}
\item[Population.]
205 (69\% humans, 11\% halflings, 7\% gnomes, 7\% dragonborn, 6\% other races)

\item[Government.]
Mages of the Cerberus Assembly, led by Wardlow Akron, run Balenpost.

\item[Defense.]
A 10-foot-high wooden wall surrounds Balenpost. All the inhabitants are mages or adventurers capable of defending the outpost.

\item[Commerce.]
The residents here are outfitted with all they need to survive by the Cerberus Assembly. Visitors can trade goods or services for lodging, rations, and other survival supplies.

\item[Organizations.]
The Cerberus Assembly runs everything in the camp, but a few of King Bertrand Dwendal's spies are among the explorers' ranks.
\end{description}
\noindent
A fort of frozen logs stands on the southwestern edge of Foren, housing intrepid explorers overseen by the Cerberus Assembly. Every waking moment of life in Balenpost is dedicated to planning and executing expeditions to uncover Aeorian items on the island. Lady Vess de Rogna (see "\textit{Cerberus Assembly}" in chapter 8) appointed Wardlow Akron (lawful evil, male, human \textbf{mage}) to oversee day-to-day operations of the fort. The explorers say that Wardlow is harsher than the winds when it comes to rationing supplies and driving expeditions to go farther and faster in the hazardous tundra.

\subparagraph{Icy Expeditions}
The expeditions that leave Balenpost are run by mages that work for the Cerberus Assembly. These mages and their crews are given honors, accolades, and special privileges, such as more rations, when they uncover Aeorian items. Teams that return to Balenpost to resupply have their wages docked until they return with a find.

Dwendalians with nothing to lose come to Balenpost to earn good pay, but many of the other explorers are prisoners quietly approached by the Cerberus Assembly and offered freedom in exchange for service. Prisoners that don't contribute to the expeditions are sent back to the mainland in chains, so many are willing to cheat, steal from, or kill other crews.

\subparagraph{Horswell's Lab}
Balenpost is not just a place of discovery but also one of arcane experimentation. Wardlow Akron placed Korberta Horswell (chaotic neutral, female, rock gnome \textbf{mage}) in charge of the team that studies and tries to duplicate Aeorian items retrieved from the wilds of Eiselcross. The reckless mage and her team perform all manner of experiments to recreate items brought back from the wastes. These tests often go dangerously awry. Death and magical destruction are not uncommon results, and a few items have warped or twisted subjects into horrible monstrosities that are locked away somewhere secret in the lab.

\subparagraph{Dwendalian Spies}
While the Cerberus Assembly is in charge of exploring Eiselcross, the Crown wants to make sure its interests are protected. The assembly reports enough discoveries to King Bertrand to make it seem like the sovereign is informed, but he knows better. King Bertrand planted spies among the explorers because he doesn't want the assembly keeping the most powerful Aeorian items secret, lest they be used to overthrow him. The leader of this spy ring is Karta Meck, a neutral human spy who also happens to be one of the camp's most accomplished treasure hunters.

Wardlow suspects that the king planted spies in Balenpost and regularly tries to read the minds of explorers and look at communications sent out of the camp to find the king's agents. His attempts thus far have proven fruitless.

\subparagraph{Resupply}
New supplies, letters, and people come to Balenpost from Icehaven by boat every month. These ships then take people and missives from Eiselcross back to the mainland of Wildemount. The journey through the Frigid Depths is perilous and can delay or destroy a resupply boat, an event which forces the outpost to restrict rations.

\subparagraph{Balenpost Adventures}
The main reason for many characters to come to Balenpost is to seek out Aeorian items for themselves or the Cerberus Assembly.

\subparagraph{Conviction (Low Level)}
Nathrow Arple (neutral evil, male, stout halfling \textbf{bandit captain}), a former enforcer for the Myriad, came to Balenpost with his old crew of four thugs. The group was captured fifteen years ago during the raid on the Myriad's central facility in Yrrosa. Nathrow got his imprisoned crew on the same exploration team, then they murdered their Cerberus Assembly leader and made off into the wilds of Eiselcross, looking to steal a fishing vessel. Wardlow Akron offers the characters 100 gp to hunt down the criminals and throws in a \textit{gem of brightness} if the characters bring back Nathrow alive (so the assembly can make an example of him). If the characters brave the dangers of Eiselcross and catch up to the criminals, Nathrow offers them anything within his power to avoid arrest and turns to combat if all else fails.

\subparagraph{Cold as Ice (High Level)}
An entire team of explorers was torn apart in the wastes just outside Balenpost. The tracks nearby indicate that only humanoids were involved in the struggle. Agents of the Dwendalian royal family send the characters to investigate, suspecting that one of the other teams ambushed the victims. Wardlow Akron is loath to cooperate with the characters since they are sent by the empire, and because he knows the responsibility for the attacks lies partly on his shoulders. The person behind the attacks is Mera Vacross, a female human transformed into a \textbf{vampire} by one of Korberta Horswell's experiments. The vampire escaped the lab and makes her home in a cave just outside Balenpost. If the characters take care of the vampire, Wardlow tries to bribe or threaten the characters to keep quiet about the monster's origins.

\subparagraph{Spy Hunters (Any Level)}
The Cerberus Assembly learns that King Dwendal's spies have infiltrated Balenpost. They hire the characters to find and report the spies to Wardlow Akron. The characters might uncover Karta Meck's network of four spies within the camp or find a lead-lined box among Karta's supplies which holds a Sending Stones paired with another held by Oliver Schreiber (see "\textit{Dwendalian Empire}" in chapter 8). If discovered, the spies attempt to win over the characters, promising an Aeorian item in exchange for helping to discover what secrets are hidden in Horswell's lab.

\subsubsection{Fortress of the Dead Jarl}
From the outside, the Fortress of the Dead Jarl is an enormous fairy-tale castle made entirely of ice. The horrors inside the place tell a different story. The frost giant leader, Jarl Conessa Berg, was obsessed with gathering objects from Aeor. Two centuries ago, one of those items turned Conessa and every giant within her castle into an undead zombie. These beings are trapped inside the castle along with a hoard of Aeorian items.

\subparagraph{Fortress of the Dead Jarl Adventures}
The Aeorian items within the fortress are enough to tempt most adventurers, but there are other reasons to go to the fortress.

\subparagraph{They're Free (Mid Level)}
When the characters arrive in Balenpost or Syrinlya, a \textbf{frost giant zombie} attacks the camp. If the characters help defeat the giant, the leader of the camp tells the characters it is the third attack in as many days and asks them to find the source of the zombies. Characters can track the zombies back to the fortress. Each day the characters follow this path, roll any die. On an even result, the characters encounter 1d4 frost giant zombies. When they reach the fortress, they find that one of the walls is breached and must repair it to stop the zombies from leaving the castle. This breach could be natural or could be caused by a faction with a score to settle against the attacked outpost.

\subsubsection{Mutalos}
In the northern part of Foren, a strange patch of wilderness exactly 30 miles in diameter is avoided by the wildfolk at all costs. They call this region Mutalos. Every morning at dawn, the area's environment suddenly changes. Mutalos shifts from desert to forest, from mountainous to swamp, or from any one habitat to another without any discernible pattern. The area can even be a seemingly normal icy tundra, so unsuspecting explorers traveling through Mutalos could suddenly find themselves in the middle of the sea when the sun rises.

The origins of these environmental changes are unknown, but scholars believe an Aeorian object fell from the city and created the region. All Wildemount's factions are interested in an item of such incredible transmutation power.

\subparagraph{Denizens of Mutalos}
When the environment of Mutalos changes, so do its inhabitants. At dawn, a mountainous region filled with kobolds, goats, and ogres might become a swamp inhabited by lizardfolk, rats, and yuan-ti. It is unknown where exactly these creatures come from and where they go when the regions change. They could be transported from other places in Exandria, teleported from elsewhere in the multiverse, or created and destroyed by this strange magic. It is also unknown if the same creatures reappear with the same habitats or if the creatures are always new to Mutalos. The origin of the region's changes likely holds the answers to these questions.

Sometimes a creature that appears in Mutalos leaves the borders of the region. When this occurs, the creature becomes untethered from the magic that brought it to Foren and remains in Eiselcross even after Mutalos's environment changes. Many creatures from other habitats don't last for long in the arctic wastes, but a few are smart or strong enough to survive.

\subparagraph{Environment Changes}
Each day at dawn, the habitat of Mutalos changes. You can choose or roll this new environment on the Mutalos Habitats table. If you roll the result of Mutalos's current environment, roll again until you get a new result.

\begin{DndTable}[header=Mutalos Habitats]{cX}

d8 & Habitat\\
1 & Arctic\\
2 & Desert\\
3 & Forest\\
4 & Grassland\\
5 & Hill\\
6 & Mountain\\
7 & Ocean\\
8 & Swamp\\
\end{DndTable}

\subparagraph{Mutalos Adventures}
Trying to find the origin of Mutalos's magic or simply traveling through the region is an adventure, but there are other reasons to explore the area.

\subparagraph{Lost Staff (High Level)}
Farrhan Yost (lawful evil, female, human \textbf{archmage}) was a part of the initial Cerberus Assembly team that explored Eiselcross. Now back on the mainland, she tells the characters she lost her \textit{staff of power} in Mutalos when the region was an ocean habitat, hoping it tempts them into exploring the region to find her old staff. As the characters search, she uses the \textit{scrying} spell to keep tabs on them. Should the characters discover her staff or any other interesting treasures, Farrhan teleports to their location with her shield guardian and offers a meager payment for the find. If the characters refuse, she attacks. Farrhan intends to get her staff back no matter what and wants to claim the characters' discoveries as her own.

\subsection{Areas (R-Z)}
\subsubsection{River Inferno}
This odd, mile-wide river of lava flows south-to-north across the length of Foren. The River Inferno disappears at both its ends into the sea, springing up seemingly from nowhere. Most scholars believe that the river was created by the crash of Aeor, though this theory has yet to be confirmed. It is clear that the river has ties to the Elemental Plane of Fire, for fire elementals, magma mephits, and salamanders can be occasionally found along its banks.

The River Inferno's lava has the destructive power one expects, though there is an exception: ice and snow that fall into the lava remain intact. Bold explorers traveling north sometimes ride the river's floating icebergs to get to their destinations more quickly. The river flows at an impressive rate of seven miles per hour.

\subparagraph{River Inferno Adventures}
The river itself is best avoided, but characters can have interesting encounters if they follow it.

\subparagraph{Haughty Efreet (High Level)}
A pair of efreet recently emerged from the River Inferno, and the characters are the first intelligent creatures they find. Using \textit{tongues} spells, the efreet demand that the characters tell them everything they know about this strange land, then demand that the characters point them in the direction of the nearest settlement. The efreet plan to enslave the settlement's inhabitants and put them to work building an obsidian tower on the shore of the River Inferno, where the efreet can stay whenever they deign to visit this dismally cold corner of the multiverse.

\subsubsection{Syrinlya}
\begin{description}
\item[Population.]
182 (72\% dwarves, 24\% elves, 4\% other races)

\item[Government.]
Professor Gulrim Shalebrow is in charge of Syrinlya, appointed to run the outpost by Queen Fruunast of Uthodurn.

\item[Defense.]
Syrinlya's inhabitants are explorers capable of defending the settlement.

\item[Commerce.]
Syrinlya's explorers are willing to trade for meager belongings. Aoerian items can be sold to the Buyer (see "The Buyer's Hut" below).

\item[Organizations.]
Elves, dwarves, and a few others from Uthodurn are the only residents of Syrinlya. Most have no organizational ties.
\end{description}
\noindent
Syrinlya is a collection of yurts filled with mercenary explorers from Uthodurn looking to get rich. Professor Gulrim Shalebrow (see "\textit{Diarchy of Uthodurn}" in chapter 8) runs this outpost with few rules: just an unspoken agreement that all explorers have each other's backs against the harsh environment, just as they would back in their mountain home.

\subparagraph{A Lawless, Moral Land}
Syrinlya has few laws, but its explorers look out for each other. This is a courtesy they learned at home in Uthodurn, where the dwarves and elves are interdependent. Most folks in Syrinlya share food when their hunts are good so that others return the favor when their luck runs out.

Syrinlya doesn't suffer bullies. The honor system is enforced by the toughest explorers, who don't hesitate to take justice into their own hands. One such individual is the dwarf Morgo Delwur (neutral good, male \textbf{werebear}), a staunch friend and an even more staunch holder of grudges. Gulgrim is content to allow the explorers to keep order, as it saves him from having to govern and allows him to go off on expeditions of his own.

\subparagraph{The Buyer's Hut}
The largest yurt in the outpost is occupied by an elf called the Buyer (lawful neutral, nonbinary, high elf \textbf{mage}). The Buyer was brought to Eiselcross by Professor Gulgrim Shalebrow to run the day-to-day operations of the camp while the dwarf explores. The Buyer keeps to themselves and refuses to share their real name, which is Elloway Tenvidas. When the Buyer acquires an object, they place the item into a large stone chest, which transports the object back to Uthodurn or Palebank Village. Some of the explorers think the chest could also be used to transport people, but none are brave enough to ask the mysterious Buyer or risk experimenting with the chest themselves.

\subparagraph{Syrinlya Adventures}
The main reason for many characters to come to Syrinlya is to seek out Aoerian items for Uthodurn, but there are plenty of other adventures to be had in the outpost.

\subparagraph{The Arcane Tree (Mid Level)}
While the characters are in Syrinlya, Morgo Delwur returns from an expedition with a bag of five green seeds from one of Eiselcross's Aeorian ruins. The \textbf{werebear} lost a seed in the snow when showing them off, and a 15-foot-tall, green tree springs up overnight in the middle of the camp. Whenever a creature casts a spell that restores hit points or deals radiant damage to another creature within 90 feet of the tree, the tree uproots itself from the ground and attacks the creature. After killing the spellcasting creature and other threats, the tree replants its roots. Every 1d10 days the tree remains rooted, a new tree of the same type grows up from the ground within 30 feet. If the trees become a problem, Professor Gulgrim Shalebrow asks the characters to deal with them.

The \textbf{animated tree} uses the \textbf{treant} stat block, with these changes:


\begin{itemize}
\item It is unaligned.
\item It has resistance against necrotic damage and radiant damage.
\item It understands Common, Druidic, Elvish, and Sylvan but cannot speak.
\end{itemize}

\subsubsection{Tomb of the Worm}
\begin{description}
\item[Population.]
345 (83\% humans, 17\% halflings)

\item[Government.]
Quajath commands the settlement's inhabitants telepathically, though a cult fanatic named Galgonos serves as the Tomb of the Worm's leader for day-to-day operations.

\item[Defense.]
All adult inhabitants in the Tomb of the Worm are capable combatants.

\item[Commerce.]
The Tomb of the Worm has no commerce, though Quajath's meat is offered freely to travelers.

\item[Organizations.]
The wildfolk that live here serve only Quajath, but they are friendly toward explorers from the Kryn Dynasty.
\end{description}
\noindent
Three decades ago, a clan of wildfolk on Gelier found the exposed body of Quajath the Undermaw poking through the wall of an icy cave (see "\textit{Lesser Idols}" in chapter 1). These accessible sections of the worm peek through the cave's ceiling and walls, and are still attached to the worm. The wildfolk who discovered this place consumed the Undermaw's exposed meat and found that it regrew. They now live in the cave, enjoying the easy, renewable source of food, and are now fully under the sway of the Undermaw. They invite others to join them, calling themselves wormkin.

The people living in the Tomb of the Worm work together in harmony. Galgonos (chaotic neutral, male, human \textbf{cult fanatic} with the wormkin template described below) is the leader of the settlement, and he organizes parties of wormkin that search Gelier and the rest of Eiselcross for more wildfolk to join their ranks. Outsiders often find the settlement charming because of how well the villagers get along, though a character who succeeds on a DC 15 Wisdom (Insight) check can tell that the wormkin are more dangerous than they appear. The villagers offer the Undermaw's meat to visitors and force-feed it to those who refuse but would make worthy additions to the Undermaw's army.

The wormkin have two goals: grow their numbers and free Quajath from the ice. What Quajath plans to have its followers do once it leaves the ice is known only to the Undermaw.

\subparagraph{Wormkin Template}
Humanoids that consume the Undermaw's flesh ten or fewer days after it has been removed from Quajath must make a DC 15 Constitution saving throw or become wormkin 1d6 hours later. Wormkin can be identified by the black veins that cover their appendages and neck. Characters who become wormkin might become NPCs at the DM's discretion. A \textit{greater restoration} spell or similar magic removes all the wormkin features from a creature, restoring it to its original state. Wormkin gain the following features:

\subparagraph{Alignment}
The creature's alignment becomes chaotic neutral.

\subparagraph{Connected to Quajath}
No matter where the creature is, Quajath can read the creature's every thought and memory and speak with the creature telepathically. Quajath can also access the creature's senses and perceive what the creature perceives. As a bonus action, Quajath can take direct control of the creature, choosing the actions, bonus action, reactions, and movement the creature takes until Quajath ends this control as a bonus action.

\subparagraph{Magic Resistance}
The creature has advantage on saving throws against spells and other magical effects.

\subparagraph{Regeneration}
If the creature has at least 1 hit point, it regains 5 hit points at the start of its turn.

\subparagraph{Tomb of the Worm Adventures}
The characters might find adventure, allies, and adversaries in the Tomb of the Worm.

\subparagraph{Meat Thieves (Low Level)}
As the characters explore Gelier, they cross paths with a group of four wildfolk commoners led by Jarinmine (chaotic good, female, lightfoot halfling \textbf{scout}). The halfling begs them to take a 300-pound package of meat to the sea and dump it in the icy water. She warns the characters not to eat it and offers them her ivory necklace (worth 100 gp) if they agree to the task. If pressed, she says the meat is tainted and that it turns people into slaves for Galgonos. Jarinmine stole the meat from the wormkin because she believes the diet poisons their minds. She needs to get back to the Tomb of the Worm before Galgonos realizes what she's done.

Before the characters reach the sea, six wormkin catch up to them. Their leader is Redjaw (chaotic neutral, female, human \textbf{cult fanatic}), and the others are cultists, all with the wormkin template applied. Redjaw demands that the characters give back the food that was stolen and forces a fight if the characters refuse. If the characters return the meat, Redjaw invites them back to the Tomb of the Worm. If the characters go to the village, they find that Jarinmine and her friends have been captured by the wormkin and are about to be executed for their crimes.

\subsubsection{Vurmas}
\begin{description}
\item[Population.]
193 (51\% dark elves, 23\% orcs, 14\% gnolls, 12\% other races)

\item[Government.]
Captain Hallwas Denalor has authority over this floating outpost, but Norca Brighttusk leads the expeditions.

\item[Defense.]
Every inhabitant is an experienced explorer capable of defending the outpost.

\item[Commerce.]
The outpost has no official economy, but people might be willing to trade for goods or services.

\item[Organizations.]
The people of Vurmas are not associated with any organizations beyond the Kryn Dynasty.
\end{description}
\noindent
Vurmas is a floating outpost for the Kryn Dynasty near the north shore of Gelier, and the dynasty keeps Vurmas secret from other factions. The Kryn established this base near Eiselcross's second-largest island not to find Aeorian items but to make allies of Eiselcross's wildfolk.

\subparagraph{Floating Outpost}
Vurmas is not a true settlement: it's actually three sailing ships tied together and anchored near the north shore of Gelier. These boats are the \textit{New Radiance}, the \textit{Candlefire}, and the \textit{Soulsailor}. At any given moment, one of the ships could be traveling to Xhorhas with a hold full of new Luxon-worshiping wildfolk converts or on the way back to Eiselcross with more supplies.

The Kryn's explorers stay on the boats when they can, bringing supplies to shore when they go out in search of allies. These ships are maintained and run by Captain Hallwas Denalor (lawful neutral, female, dark elf \textbf{veteran}), while Norca Brighttusk (neutral, female, orc \textbf{priest}) leads the expedition teams.

\subparagraph{Ally Recruitment}
The Kryn Dynasty believes that the wildfolk of Eiselcross are a better find and more reliable than most Aeorian objects, especially since the items tend to interfere with the Kryn's divine magic. With the Cerberus Assembly and the Uthodurnians busy in Foren, the Kryn stay clear of the largest island and search for allies in the rest of Eiselcross's lands. Many wildfolk are easy converts of the Luxon since they respect divine magic but do not worship any deity.

The Kryn approach each wildfolk clan with caution and kindness, healing the sick and injured while preaching the faith of the Luxon. They take only willing wildfolk away from the clan, knowing many others will follow because they miss their friends and family.

Choosing the waters near Gelier as a base of operations was no accident. This island is home to the wormkin, wildfolk that devour the frozen from of Quajath (see "\textit{Tomb of the Worm}" earlier in this chapter). These wildfolk have strange minds, but they display regenerative abilities because of their diet. Norca Brighttusk regularly preaches to the wormkin, telling them that the beast they feed on is a creation of the Luxon. The wormkin have agreed to convert and join the Kryn once Quajath is freed from the ice—an effort that is already underway. If the Undermaw and its people can be turned toward the light of the Luxon, the Kryn Dynasty will have gained a considerable advantage in its war against the Dwendalian Empire. None of the Kryn, not even Norca Brighttusk, knows that Quajath has no intention of serving anyone and that the wormkin are completely under its sway.

\subparagraph{Vurmas Adventures}
The main reason for characters allied with the Kryn Dynasty to come to Vurmas is to join Norca Brighttusk's recruitment efforts, but there are plenty of other adventures to be had in the area.

\subparagraph{Keep Us Secret (Low Level)}
Kryn spies report that the Bold Adventurer, a scouting ship from the Dwendalian Empire, is planning to visit the island of Gelier, near Vurmas. Captain Hallwas Denalor and Norca Brighttusk don't want the Bold Adventurer's crew to see Vurmas or reach Gelier, ensuring that the Kryn's operation in Eiselcross and the settlement of the wormkin remain secret. The pair offer the characters 500 gp to stop the Bold Adventurer's approach. The Bold Adventurer is a keelboat (see the Airborne and Waterborne Vehicles table in Owning a Ship; Airborne and Waterborne Vehicles of the \textit{Dungeon Master's Guide}) that carries six scouts. The captain of the Bold Adventurer, Sally Brizzberm, has a \textit{sending stone} that she uses to communicate with the captain of a Dwendalian warship called the Bright Queen's Doom. If her ship is attacked, Sally uses the stone to send a warning to the warship, which might threaten Vurmas in the future.

\subparagraph{Save Our Ships (Mid Level)}
While the characters are aboard one of Vurmas's boats, three wildfolk druids, each riding on the back of a \textbf{killer whale}, attack the boats. The druids are angry that wildfolk are leaving Eiselcross in droves to go to Xhorhas. They believe that the Kryn are using enchantment magic against these wildfolk, tricking them into leaving their friends and family. The druids use hit-and-run tactics, dealing as much damage as possible before fleeing. Captain Hallwas Denalor asks the characters to deal with the attackers. The druids could be open to negotiation, or they might fight with their dying breath for their beliefs.

\begin{DndSidebar}[float=!b]{Moorbounders}
Beasts of burden are common throughout Xhorhas, and often horses and other riding animals are the best or only option. For those who require speed and have a little more coin to spend, the best option is often a \textbf{moorbounder}. However, moorbounders that haven't undergone proper training or established bonds of trust with their masters tend to attack and even eat their riders.



Moorbounders can cost anywhere from 300 to 500 gp each. They're often cheaper and more abundant in towns on the wastes than in cities like Rosohna.



\end{DndSidebar}

\section{Wastes of Xhorhas}
Stretching across the lands of Eastern Wynandir is a valley that still bears the lingering scars of the final battles of the Calamity. Fields of broken earth and muddy paths weave around muggy swamps and corrupted forests, home to creatures twisted by the echoing shadows of the Betrayer Gods who once ruled this region. Fierce storms rule the sky as hungry predators stalk the open, desolate plains.

All manner of beastfolk, goblinkin, and nomadic outlanders claim the marshes and badlands of this region, honing their survival skills and warring ways against the elements and each other. Meanwhile, the dark elves of Rosohna have emerged as the core of civilization in Xhorhas, reaching out to the wayward tribes and offering them enlightenment if they join the Kryn Dynasty.

The dark elves of the Kryn Dynasty worship the Luxon, a mysterious deity of light, after rejecting the evil designs of Lolth long ago. Nevertheless, the Spider Queen's venomous will still seeps into the dynasty, as she and her devotees plot revenge.

\subsection{Areas (A-D)}
\subsubsection{Asarius, the City of Beasts}
\begin{description}
\item[Population.]
48,025 (41\% goblinoids, 32\% gnolls, 10\% dark elves, 17\% other races)

\item[Government.]
Den Olios governs the city, the laws of which are enforced and monitored by captains of the Aurora Watch.

\item[Defense.]
An impressive army is stationed on the outskirts of the city, comprised of Aurora Watch soldiers and trained local warriors, as well as mercenaries.

\item[Commerce.]
Basic supplies and equipment are available, as well as quality arms and some rare goods. Inns and taverns offer austere accommodations.

\item[Organizations.]
There are several temples to the Luxon, as well as some shrines to other deities. Merchant unions are beholden to the will of Den Olios. Meanwhile, Lolth's Children of Malice lurk in the shadows.
\end{description}
\noindent
Encircling and covering a massive, cracked hillside in the center of northern Xhorhas, Asarius has grown into a cluttered sprawl of roughshod homesteads, mud huts, and military fortifications that are home to a bustling blend of the many peoples of the Kryn Dynasty. The city is a cultural hub established by the Kryn for the wilder folk of the surrounding wastes and nomadic peoples who joined the dynasty. The Kryn convinced these wastefolk to unite under their banner by promising them protection, resources, and, in many cases, enlightenment (followed by conversion to the faith of the Luxon).

Although the establishment of Asarius was initially fraught, the city quickly grew into a symbol of unity for the scattered cultures of northern Xhorhas, who eschewed the dangerous, competitive fight for survival to instead focus on building a community. In the generations since, the city has blossomed into a tangled web of quirky neighborhoods, filled with industry and musky scents. The muddy plains outside are tilled for hearty roots and produce, monstrous livestock are herded and kept in large pens and pits, and a growing force of Aurora Watch soldiers train and prepare to battle the forces of the Dwendalian Empire.

\subparagraph{Forged in Frenzy}
The wastes have always bred strong peoples who could survive despite the harsh environment, deadly wildlife, and scarce resources. Many of the wastefolk fought each other over hunting grounds and territory until the Kryn rebuffed attempts at violence and instead offered to bring together the tribes for the common goal of mutually assured survival. Some were intrigued by the opportunity, while others did not fully trust the strange elves who chose to wield words instead of swords. The early establishment of Asarius was marked by conflict, and while the community has since flourished, there are still many who live beneath the dynasty's protection but mistrust the dark elves' intentions.

\subparagraph{A Martial Pillar}
While many see the establishment of Asarius as pure altruism on the part of the Kryn Dynasty, one cannot deny the importance of the city's central location among the wastes and strategic proximity to the border of the empire. The Bright Queen has transformed the tiny township into a gargantuan military hub for the dynasty, enlisting many thousands of wastefolk to form the backbone of its standing army. Wide swaths of military tents cover the western outskirts of the city as countless soldiers prepare for the coming conflicts in the west, training atop the terrifying war beasts and towering war machines brought from Rosohna or built alongside the towering walls of Asarius.

\subparagraph{Government}
The rule of law is tenuously imposed throughout the chaotic city, where the squabbles of argumentative civilians often lead to physical conflict. The laws of the dynasty, which stem from the faith of the Luxon, are enforced by a network of trusted Aurora Watch captains and soldiers who report to Den Olios within the Aurora Hold. Local affairs are the jurisdiction of Lady Zethriss Olios, who works with representatives of each major neighborhood to manage the city and its people.

\subparagraph{Crime}
Blood often runs hot in Asarius, and folk here often come to blows, simply to deal with the stress and frustration of life on the plains. The Kryn know this; as such, their laws regarding physical assaults and fights are far laxer in this city than anywhere else in Kryn territory. Even so, disagreements can escalate to a point where theft, brutality, and murder sometimes occur with worrying regularity.

\subparagraph{Geography}
The majority of this rough, boisterous city is draped across a wide, muddy hill that overlooks the surrounding plains. The hill is surrounded by fifteen-foot-tall walls of purple vermaloc wood, harvested from the eponymous Vermaloc Wildwood. The interior of the city is a tangled mass of wayward, mud-slick streets and a hodgepodge of homes and structures that architecturally mirror the social mixture of dark elf society and the many nomadic lifestyles of the wastefolk. The center of the city is marked by the ominous spire known as the Aurora Hold, home to Den Olios.

The southeast region of Asarius is called the Pits, named for the multitude of breeding pits for war beasts and other beasts of burden. The Pits are responsible for the city's ever-present smell of wet animals and dung, and the area is also home to many poorer citizens seeking work wherever it can be found. The southwest portion of the city is known as the Creets, a rustic tangle of goblinoid neighborhoods and houses of farming folk who till the harsh environment beyond the eastern walls. The entire northern side of the city is simply known as the Northroads. Within this area, the Scowl Square is the center of commerce, and the North Center forms the industrial core of Asarius. An average of four thousand Aurora Watch are stationed outside the western walls of the city, with additional reinforcements during times of looming conflict.

\subparagraph{Asarius Adventures}
Characters might come to the City of Beasts to find a reliable war beast, receive military training, or spy on the Kryn military preparations.

\subparagraph{War Beast Escape (Mid Level)}
While the characters stay in Asarius, spies from the Dwendalian Empire infiltrate the city and simultaneously release war beasts from their pens after enraging them with whips and hot irons. As the characters try to stop the rampage, they learn that the spies are using the confusion to steal the city's secrets.

\subparagraph{Echo of the Past (Mid Level)}
A chaotic good, male, dranassar survivor named Jibu-Corryn has been living as a dark elf in Asarius under the name of Alyxikon, a minor government secretary. A cult of hobgoblins and bugbears who call themselves the Lords of Strife are hunting Jibu after hearing a rumor that one of the last of the dranassar are hiding in Asarius. Their leader, a hobgoblin captain named Gorthux, possesses the power to transform into a \textbf{bone devil} once per day, and plans to offer Jibu as a sacrifice to Bane.

Jibu has heard tales of the characters' prowess and contacts them with a cryptic note. He begs for protection while he settles his affairs in Asarius and travels to a hermitage in Blightshore.

\subsubsection{Barbed Fields}
The valley north of Rosohna is covered by this broken and marred landscape of jagged rock, sunken craters, and scabrous, hooked spires of glass and rock jutting from the ground. It is said that the final battles of the Calamity were fought on this valley floor. The earth is so scarred here that few question the veracity of this legend. Treacherously unpredictable earth crumbles into sinkholes surrounded by thick, caustic tar pits that swallow careless wanderers, while the storms common to the region bring rain-slicked mud traps and flash floods.

However, powerful relics and armaments have been recovered from these ancient battlefields, so the strong and the stupid occasionally venture into the dangerous fields, often finding themselves hunted by the monsters that lurk in this ruined domain. It is rumored that toward the center of the barren fields grows a beautiful, towering tree that grew from the final seed left by Melora before she vanished to her realm beyond the Divine Gate. Powerful followers of the Wild Mother sometimes make a pilgrimage to this tree for guidance and communion, but the fields are so fraught with evil that few ever return.

\subparagraph{Barbed Fields Adventures}
The deeply wounded lands of the Barbed Fields are rife with planar entities bleeding through where the veils between worlds are thin.

\subparagraph{Dark Days (Epic Level)}
The extremely rare double eclipse, when the moons Catha and Ruidis overlap before the sun, is approaching, and it will occur over Xhorhas. Seers and oracles warn of portents that speak of an uncontrollable planar crossover, where multiple planes might briefly collide. The Kryn Dynasty hires the characters to bolster their ranks in preparation—but no one is prepared for the army from the Nine Hells that invades the fields when the realms merge. Leading this evil army is a \textbf{pit fiend} named Vanatoth.

\subsubsection{Bazzoxan}
\begin{description}
\item[Population.]
2,610 (81\% dark elves, 8\% goblinoids, 11\% other races)

\item[Government.]
Martial law is enforced by Taskhand Verin Thelyss and maintained by the local military.

\item[Defense.]
The Aurora Watch defends the town and keeps the demonic incursions at bay. The few citizens learn combat skills out of necessity.

\item[Commerce.]
Very basic supplies are available, if scarce, due to the dangerous circumstances surrounding the town. Most trade involves provisions and goods for survival and defense.

\item[Organizations.]
There are a number of partially destroyed temples and a handful of small ones that are still intact.
\end{description}
\noindent
Bazzoxan was once a dark temple instrumental in the machinations of the Betrayer Gods during the Calamity, but the remaining ruin was left to gather dust like many relics of those terrible times. After Ghor Dranas was claimed by the Kryn, settlers roamed north through the Barbed Fields to discover this site with hopes of renovating it as well. A community began to flourish within the ruins of Bazzoxan over a few generations, building the foundations of a new city around the looming stone structures of the abandoned sanctum—until demonic forces from within the old temple threw the burgeoning society into chaos. Bazzoxan is now perpetually locked in a stalemate between the dynasty's forces and the abyssal incursion from within the temple. This unholy site is the subject of constant research and of growing worry for the dynasty.

\subparagraph{Never Left Alone}
About forty years ago, a group of curious local dungeoneers uncovered a long-dormant gateway to the Abyss. Its demonic influence drove them mad, and they accidentally reactivated the gateway of Bazzoxan—an immense rift between the Abyss and Exandria. This enabled demonic forces to pour out into the ruins and surrounding city. The incursion was eventually stifled, though at a great loss of life. The civilians that remain in Bazzoxan support the warriors who keep the forces of the Abyss at bay, as well as the arcanists who work to reseal the rift and reclaim the city.

\subparagraph{Government}
The town is under martial law, with Aurora Watch soldiers keeping the peace under the current governance of Taskhand Verin Thelyss. The taskhand's lieutenants serve as minor officials and bureaucrats, and even the civilians who continue to stay are trained to defend themselves and fall into rank if malevolent forces rise from beneath their home.

\subparagraph{Crime}
With the tension of living in constant fear of demonic threat, the focus on survival in Bazzoxan leaves little room to break the law for personal gain. Even so, the rare grifter sometimes profits from gouging the price of provisions or swindling the fearful. Criminals who are caught by the Watch are sent south to Rosohna for punishment or simply executed, depending on the severity of their crimes.

\subparagraph{Geography}
The towering ruined temple sits against the rocky southern base of the Penumbra Range, carved directly from the natural stone of the mountains. The ominous onyx doors that lead into the heart of the temple remain barred and under constant watch. Surrounding the temple are over a hundred abodes and structures, as well as a partially reconstructed town square. Construction equipment lies abandoned in the streets, and vacant homes sit untouched since their denizens were slain or forced to flee. Makeshift Aurora Watch barracks are continuously expanded to make room for the reinforcements sent to push back the fiends that assail the city from beneath.

\subparagraph{Bazzoxan Adventures}
Adventurers loyal to the Kryn Dynasty are welcome in Bazzoxan and will almost certainly be asked to aid with keeping Abyssal forces at bay or gathering spell components for the mages trying to seal the rift.

\subparagraph{Heir Aberrant (Mid Level)}
One of the dungeoneers who first discovered the gateway beneath Bazzoxan, a lunatic who nonetheless lived out his final days in luxury, has just passed away. The inheritor of his estate is offering a modest reward to any investigators who can locate certain relics of abyssal origin rumored to be hidden in the deceased madman's trapped abode. But what are the intentions of the so-called inheritor, and what danger do these relics pose to the people of Xhorhas?

\subsubsection{Brokenveil Marsh}
The sheer cliffs and ravines of the Brokenveil Bluffs were created when the Dunrock Mountains and the Ashkeeper Peaks were sundered in the Calamity. Between these precipices, a dense morass of thick, stinking marshland claims the valley floor. Gnarled trees and tangle vines clutter the shaded canyons, while sucking mud pits and buzzing insects make travel a slow and arduous experience. All kinds of swamp beasts and sinking muck traps threaten to swallow those who trespass in the marsh, and filthy bluff harpies are eager to pluck away stragglers.

The marsh is a comfortable home for horizonback tortoises, so the challenge often lies in perceiving which dry mound in the Brokenveil is an actual hill and which is the shell of a territorial monster.

\subparagraph{A Buffer of Battle}
The marsh helps discourage the eastern expansion of the Dwendalian Empire and deters the Kryn Dynasty from pushing farther west. However, nearly two decades of skirmishes between these nations have left many soldiers on both sides buried in the marsh, their bodies forever lost to the dark muck. The Kryn's recent victory at the Ashguard Garrison has given them control over the marsh and a foothold in the enemy territory beyond it.

\subsubsection{Charis}
\begin{description}
\item[Population.]
950 (53\% dragonborn, 41\% halflings, 6\% other races)

\item[Government.]
Two elders preside over this village, one representing its dragonborn population and the other representing its halfling population.

\item[Defense.]
Numerous bands of experienced fighters, druids, and hunters patrol the area. Civilians know basic combat techniques to defend the village if necessary.

\item[Commerce.]
Travelers can find basic supplies, most likely at a small inn. Visits from traveling merchants are infrequent, so rare goods are special finds.

\item[Organizations.]
A mighty treetop temple to Melora shades the center of the village.
\end{description}
\noindent
A wandering community of halflings survived in the fiend-twisted Lotusden Greenwood for generations through their cunning and the grace of their faith in Melora. When they found themselves ambushed by a group of starving ravenites fleeing from the destruction of Draconia, the halflings shunned conflict in favor of kindness. They offered to feed the ravenites, and a bond between peoples was born.

This burgeoning community of halflings and dragonborn was able to carve out a permanent home within a safer region of the Lotusden and, with the blessing of the Wild Mother, ward off the evils of the forest. They gave their woodland village the name Charis, from the Draconic word for "hope." This hidden hamlet of hollowed trees, moss-covered grottoes, and mushroom fields harbors a nature-worshiping commune of vegetarian farmers, forest guardians, and druids who work to live in harmony with the surrounding forest. They hope to cleanse the Lotusden of the evils that have pervaded its shadows since before the Divergence.

\subparagraph{Ignorance Is Bliss}
Despite the chaos and turmoil that rock the lands of Wildemount beyond the Lotusden Greenwood, the people of Charis are happy in their isolation. Lacking any real interest in outside affairs, they focus entirely on protecting the surrounding forest and their fellow villagers. Outsiders occasionally visit Charis, and while the villagers are welcoming to those who offer no sign of aggression, they usually find a way to quickly usher them along, preferring to be left to their simple lives and pursuits. The people of this village would never consider leaving their tranquil home, unless so commanded by the Wild Mother herself.

\subparagraph{Government}
The community is small enough that most intricacies of formal government are unnecessary. Volunteers take up responsibilities in shifts, from keeping watch to gathering food. Leaderships is shared between two elders, one for each of the halfling and dragonborn populations, and their authority is mostly required for giving spiritual guidance and resolving minor disputes.

\subparagraph{Crime}
The society within the village operates without money, gold, or emphasis on precious stones. Everything is provided by the community and shared equally, in theory. In practice, however, disputes do occasionally arise, and ravenites who remember their times in the Dreemoth Ravine sometimes feel compelled to secretly hoard things once considered valuable. While there are no formal laws, if enough citizens feel wronged by an action, they bring the issue before the elders, who will judge the situation and choose a fitting punishment.

\subparagraph{Geography}
Nestled within a sunken meadow surrounded by dense trees and hills, this hidden sanctuary is a beautiful pocket of protected nature within the haunted shadows of the Lotusden. Colorful flowers and fungi line the pathways between stacking treehollow homes and hillside domiciles, and vine-wrapped bridges hang between treetops. Sunbeams pierce the canopy of the forest, bringing vibrant warmth during the day, while moonlight calls faint whispers of mischievous fey into the village to mingle with its denizens.

\subparagraph{Charis Adventures}
Finding Charis is an adventure the characters might undertake to locate a person with a specific item or piece of information that could aid them in another quest.

\subparagraph{Overlord of the Overgrowth (Mid Level)}
A demon has made the forest bordering Charis its lair. Every day, the \textbf{hezrou} Garlagrag attempts to summon more and more demons into the forest, and the woods are now flush with hezrous and dretches. The people of Charis seek aid to fight off the demons or to end Garlagrag's plot once and for all.

\subsubsection{Deepriver Mine and Icozrin Steelworks}
Established around a plentiful source of iron and copper within a dry, subterranean riverbed, the Deepriver Mine is one of the dynasty's most productive mining operations. Workers toil away within the mines beneath the red canopy of the Vermaloc, or work the smelters that litter the courtyard, preparing to send their goods to Rosohna or farther into the west of Xhorhas.

Since the conflict with the Dwendalian Empire has escalated in recent decades, Den Icozrin established an impressive steelworks alongside the Deepriver Mine to expedite the forging of armaments. The Icozrin Steelworks has since grown to become a construction yard for war machines and siege weaponry to be sent to the front lines in the west.

\subsubsection{Dreemoth Ravine}
Along the southeastern range of the Ashkeeper Peaks, an impressively deep gorge cuts into the mountains for nearly eighty miles. Winding and widening, the enormous ravine reveals a beautiful, rocky valley of sparse trees, natural caverns, and incredible waterfalls tumbling from the Ashkeeper runoff. Strong winds blow through the gorge, and the numerous natural cliffside outcroppings provide shelter for wild griffons, wyverns, and other flying beasts. The central bend of the Dreemoth Ravine still contains the ruins of Draconia, as well as the countless abandoned caves and quarries that once held the enslaved ravenites, now freed.

\subparagraph{Dreemoth Ravine Adventures}
Characters might come to Dreemoth Ravine to capture and train a flying mount, or to explore the ruins of Draconia.

\subparagraph{Bones of Frigid Doom (Epic Level)}
A \textbf{death knight} named Pentrakath lurks in a cave in the Dreemoth Ravine, and he has uncovered the bones of Vorugal, the \textbf{ancient white dragon} that destroyed Draconia twenty years ago. He gathered a host of profane relics and stole the souls of hundreds of dead dragonborn in an attempt to stitch together a soul powerful enough to resurrect Vorugal as an ancient white dracolich (see \textit{Monster Manual}).

\subsubsection{Dumaran}
The hobgoblin fortress of Dumaran is built into the western base of the Penumbra Range, hidden within the south end of the web-infested Vermaloc Wildwood. Dozens of defensible portholes and arrow slits look down over the walkways that wind across the face of the impressive fortress, while curling paths lead up into the mountains. Within, countless halls, chambers, and tunnels are carved through the rock.

Following the Calamity, the hobgoblin survivors were scattered, fighting the clans of the Iothia Moorland for scarce resources. They drew the attention of Lolth, the Spider Queen, who seduced and empowered the hobgoblins of Dumaran. Although they were not drow, the hobgoblins joined the ranks of the Children of Malice, giving them dominion over the spiders throughout the Vermaloc and honing them into a ferocious army to topple the Kryn as punishment for their insolence.

\subparagraph{Changed by Hatred}
Twisted by the dark gifts of their murderous goddess, the hobgoblins of Dumaran have onyx black fur, pale yellow eyes, and a hunger for slaughter. Their nocturnal hordes strike out in the dark of night, commanding the massive, venomous arachnids of the Vermaloc to aid their raiding parties. Their rituals emphasize their loathing of the Kryn drow, the same who abandoned their Spider Queen, and they work to slowly build paths through the Penumbra Range toward Rosohna to eventually corrupt the city from the shadows.

\subparagraph{A Jealous Mistress}
Like an equally dark mirror of the Lolth-worshiping drow in distant regions of Exandria, the warrior matrons of the Children of Malice hold power within Dumaran, leading their family houses while serving under Head Matron Vivurk Tonn. Though many drow societies are filled with political scheming and infighting to win Lolth's favor, their queen's seething hatred of the dynasty focuses these drow against a common enemy, making the forces of Dumaran a ruthlessly efficient engine of vengeance.

While crimes occur from time to time, the threat of swift and brutal punishment is enough to keep the inhabitants of the fortress in line, unless they are absolutely certain they can get away with a misdeed or prove that their actions further Lolth's goals.

\subparagraph{Dumaran Adventures}
Those who serve the Spider Queen in Xhorhas often pass through Dumaran, and heroes who seek to end her poisonous influence might attempt to destroy the stronghold.

\subparagraph{Matron of Malice (Mid Level)}
The Kryn Dynasty aims to create upheaval within Dumaran by hiring the characters to infiltrate the fortress and assassinate Head Matron Vivurk Tonn. To get to her, the characters are given potions that can polymorph them into dark elves for 24 hours.

\subparagraph{A Looming Visit (High Level)}
Intercepted messages between the Children of Malice speak of a ritual that is being prepared in the heart of Dumaran in the coming weeks. Victims have been captured for a mass sacrifice, and the messages mention a "great arrival." The Kryn Dynasty hires the characters to uncover more information about this ritual, attempt to release any prisoners, and, should the opportunity present itself, stop this ritual or destroy whatever arrives at its completion.

\subsection{Areas (F-R)}
\subsubsection{Fevergulf Lake}
Fevergulf Lake is a massive body of brackish water known for its warm temperature and sulfurous smell. Depending on the season, the waters can become almost scalding in places, causing many of the fish and other creatures that live beneath the surface to migrate from one end of the lake to the other. The banks of the lake collect layers of strange minerals and stones of all colors, creating a rainbow of odd pastel hues among the pebbles, sands, and muck on the lake shore. The immensity of the lake gives plenty of room for all manner of strange terrors to thrive beneath the waters, making it extremely dangerous to cross via boat.

\subparagraph{Fevergulf Lake Adventures}
The characters could attempt to cross Fevergulf Lake without waking its monsters, or they could hunt the denizens beneath.

\subparagraph{Eyes of the Lake (Any Level)}
The two islands near the center of the Fevergulf, known as the Eyes of the Lake, are home to rocky shores and twin forests of brown, moss-covered, dead trees. Tall, twisted, and leafless, these woods are ominous enough from a distance, but the clans of the Iothia Moorland still tell fables of a strange, isolated cult from long ago that lives there to this day, mourning the loss of some long-forgotten god.

\subsubsection{Iothia Moorland}
The Iothia Moorland fills much of the central valley of Xhorhas, a landscape of hilly grasslands and muddy fields that stretches from horizon to horizon, framed by soggy swamps and massive lakes. The Moorland is surprisingly fertile, supporting a robust ecosystem of game and beasts that call the valley home. The roving Dolorav and Koshtask clans claim the Moorlands as their hunting grounds, fighting back the deadly predators that stalk the marshes while also fighting intruders and each other.

\subparagraph{The Dolorav: Wayfaring Warriors}
The Dolorav clan is a fierce and disciplined warrior tribe, as well as one of the largest human societies in Eastern Wynandir. The Dolorav are descended from the surviving members of the Vukan, who were defeated by the Ki'Nau and their patron Uk'otoa in centuries past. They've now grown into a matriarchal society ruled by a champion called the Skyspear, who is chosen through feats of strength and mental acuity. Tribal warriors gain glory by surviving dangerous trials and rituals, and the majority of the Dolorav hunters and fighters are women.

All who belong to the clan take a blood oath called the Marking as a teenager (or on being accepted into the clan, if they joined later in life). This vow demands celibacy until an appointed mate is chosen by the Skyspear. The clan is a tight-knit community that understands the perilous nature of the land, and is ever vigilant to protect its own and rebuff the attempts of the dynasty to indoctrinate members into the Luxon's worship. On rare occasions, the Dolorav clan has been known to ally with the Koshtask clan against a common foe, and recently there have been talks of renewing such a truce to fight the Many Hosts of Igrathad.

At present, the Dolorav clan is recovering from the massacre of their previous leadership by one of their own only a year ago.

\subparagraph{The Koshtask: Willful Brutishness}
The Koshtask clan is a powerful union of ornery orc marauders and abandoned wastefolk who were left behind when they chose not to follow the clans that joined the Kryn Dynasty and established Igrathad. As worshipers of Gruumsh, the Koshtask value efficiency in the kill and take what they want. They prefer fearmongering to diplomacy, and most interactions with the Koshtask end in violence, unless a greater enemy is present against whom they must ally. The current chosen of Gruumsh, a fire-scarred orc called Arakki, leads with cunning and foresight, seeking convenient alliances to ensure his clan's survival, fight off the Kryn missionaries that seek to convert his people, and eventually betray and destroy those foolish enough to trust them.

\subparagraph{Iothia Moorland Adventures}
These muck-laden fields of swamp and grass hide many lost secrets that have sunk beneath the mud—and deadly denizens intent on keeping those secrets hidden.

\subparagraph{Lost Resting Place (Mid Level)}
A disorganized clan of cyclopes wander the moorland, searching for the grave of their founding champion who died hundreds of years ago. They might parlay with the characters and are willing to trade fine, rare items for aid in seeking out their champion's resting place. When the grave is discovered, the cyclopes immediately begin a ceremony that raises the bones of their founder from the muck into a new existence as a giant. The characters can decide to halt the ritual, angering the cyclopes, or learn later that the clan is terrorizing the tribes of the moorland.

\subsubsection{Jigow}
\begin{description}
\item[Population.]
13,210 (56\% goblinkin, 31\% orcs, 13\% other races)

\item[Government.]
The assigned taskhand, Durth Mirimm, has a strained partnership with the two elder leaders of the goblin and orc clans.

\item[Defense.]
Taskhand Mirimm's small platoon of Aurora Watch soldiers bolsters the sizable ranks of the local warriors.

\item[Commerce.]
Most basic goods and services can be found here, though their quality reflects their rough and rural origins. Hostels and lodges offer rooms and provisions for travelers.

\item[Organizations.]
A central temple to the Luxon presides over a district called the Jumble (see "Geography" later in this section). Each clan maintains their own network of artisans and merchants.
\end{description}
\noindent
A tenuous alliance of goblin clans and territorial orc clans once scavenged along the northern shores of Xhorhas, warring and competing for survival. Over the past century, the Kryn Dynasty has worked to bring the two roving factions together. The negotiations were long and arduous, but the Kryn Dynasty's introduction of the faith of the Luxon ultimately helped the clans to form a unified society.

From these talks, the coastal city of Jigow rose to become an important source of seafood within the dynasty. A booming fishing society that spreads over both the banks of the Emerald Gulch Sea and the northern Ifolon River, this community of powerful hunters, cunning fishers, and tireless crafters has become a valuable symbol of the dynasty's growth and influence.

\subparagraph{Always at Odds}
Despite the peace between the peoples of Jigow, their long-standing rivalries require other outlets wherein they can express their need to conquer and gloat, which is why Jigow celebrates competition in any form. Whether the rivalry of the day is a contest between hunting parties, a cooking competition, or just children playing sporting games, the people of Jigow love turning amicable rivalries into a loud and colorful spectacle—especially when they can make bets or help their chosen side cheat to achieve victory.

\subparagraph{A Stubborn Culture}
The beliefs and superstitions of the original Gakthash, Uvuroh, and Zetek goblin clans permeate Jigow's goblin society. Similarly, the honorable traditions of the Norgab and Shoomat orcs continue to guide their people. While some have adopted the beliefs of the Kryn, many resist the teachings of the drow, as they feel that the presence of drow is suffocating their own culture. The mounting tensions between the dynasty and the people of Jigow are only further exacerbated by the worry that the dynasty's war will spill into Jigow territory.

\subparagraph{Government}
Taskhand Durth Mirimm represents the dynasty's interests first and foremost, putting the needs of Jigow's residents a distant second. His controversial judgments are often challenged by the most vocal of the clan elders: Elder Cubu-Ka, an elderly goblin who enjoys frustrating the Kryn representatives, and Elder Ushru, an aging orc hunter eager to maintain the cultural heritage of his people. To the dynasty, maintaining order resembles a handful of adults failing to keep unruly children under control. The native people of Jigow, however, see their willfulness as a way to establish boundaries against a clueless enforcer.

\subparagraph{Crime}
While members of the Aurora Watch do their best to keep things under control, a constant hum of small-scale thefts can't be quelled. Between watching their own pockets, the occasional violent confrontation from a contest gone sour, and illegal attempts to undermine a competitor's business, the Watch is perpetually overworked, which only invites more shenanigans from the people of Jigow.

\subparagraph{Geography}
Jigow is more of a collection of coastal villages than a singular, central city. The settlement is comprised of a chain of settlements clustered among the trees of the shoreline. The main docks, known as the Meatwaters, lie on the westernmost shore, where the fishing industry thrives due to the healthy marine ecosystem. The central villages, called the Jumble, sit slightly farther from the coast, built among the massive mangrove trees and nomadic horizonback tortoises that serve as roving homesteads. The eastern edges, known as the Wetwalks, are filled with huts and shacks built on stilts over the low wetland patches where grains are cultivated in the paddies amid the reeds.

\subparagraph{Jigow Adventures}
Characters in Jigow might get pulled into local competitions, or they might be keeping an eye on the populace for the Kryn, who are trying to avoid a revolt.

\subparagraph{Battle of the Braggarts (Low Level)}
At a tavern, the adventurers witness two braggarts engage in a boast-off. The two boasters, a battle-scarred old \textbf{orc} named Temgor and a young \textbf{goblin} named Gerb, are both prideful and refuse to back down. It seems like the conflict is going to come to blows, and both Temgor and Gerb's friends separately ask the adventurers to defuse the tensions before the Aurora Watch gets involved.

\subsubsection{Lotusden Greenwood}
The dark, verdant forest of the Lotusden Greenwood envelops the southern end of Eastern Wynandir in eerie mists, murderous trolls, and hungry trees. Incredibly lush and filled with life compared to the barren and desolate wastes to the north, this unchecked forest is home to many creatures and beasts, some of whom have grown to a nightmarish size because of the residual power of the long-buried druidic fane that empowers and protects the Greenwood.

Giant predators and dire monstrosities stalk the undergrowth, while vicious, carnivorous plants patiently await their next meal. Cold nights bring icy mists to the wood, drawing the spirits of those claimed by the Lotusden out into the moonlight to siphon the warmth from the living. The Kryn Dynasty considers exploiting the mysterious and deadly forest to be an unnecessary risk, so the Lotusden Greenwood remains one of the few places on this side of the Penumbra Range untouched by Kryn interests.

\subparagraph{A Bite as Bad as Its Bark}
No one knows the source of the wild curses that plague the Lotusden, but there are a number of stories that speak of whole portions of the forest suddenly coming to life. The trees uproot themselves and march as one, going berserk and assaulting any they perceive as outsiders—until mysteriously, and without warning, they take root again and return to their docile, unmoving state. While academics maintain that there is no definitive proof of this behavior, some explorers have noticed the shape of the forest's edges changing dramatically, sometimes overnight.

\subparagraph{Lotusden Greenwood Adventures}
The characters might enter the forest to retrieve someone who is lost or to escape Kryn forces in pursuit. Any plant in Lotusden Greenwood could be a carnivorous monster waiting to attack.

\subparagraph{Mother's Milk (Mid Level)}
Someone or something has been kidnapping the infants and children of Charis. A terrible evil lurks at the heart of the matter: an unholy coven of hags, led by a \textbf{night hag} of prodigious influence and power. The hag preside over a dark corner of the Greenwood, where they lair beneath the roots of a gargantuan tree. Here, the nightmarish coven and its minions manufacture the slow doom of the Greenwood itself as they siphon the lifeblood of the Charis children into a loathsome artifact known as the Gloom Lathe.

\subsubsection{Many Hosts of Igrathad}
\begin{description}
\item[Population.]
6,750 (41\% goblinkin, 19\% orcs, 15\% humanoids, 25\% other races)

\item[Government.]
A local leader governs independently alongside a Kryn representative in each of the seven villages. The village leaders occasionally convene as a council.

\item[Defense.]
A small force of Aurora Watch soldiers weave through the villages, and each village maintains its own army of warriors.

\item[Commerce.]
Most goods or services can be found within the villages, though it might take some searching, and strange or rare commodities are occasionally available from merchants or scavengers. Inns and taverns are scattered throughout the villages.

\item[Organizations.]
Small shrines and temples can be found in each village, and they each maintain their own marketplace and merchant unions. The Children of Malice have begun to infiltrate many of the villages.
\end{description}
\noindent
South of the Sorrowseep Waters, where the fields are more fertile, lies a scattered network of seven villages with robust communities of goblinkin, orcs, and giants, known as the Many Hosts of Igrathad. The seven villages are Kranad, Broonbobah, Aruuth, Oshinik-Ka, Wrathfall, Unyeethi, and Dodafir.

All seven villages were built on a central battleground where the factions of this region feuded and warred over resources and glory. The Kryn Dynasty reached out to the village leaders nearly a century ago and managed to broker a peace to halt the bloodshed. This disparate community is a uniquely powerful, if quarrelsome, ally to the dynasty.

\subparagraph{A Warrior Culture}
Over the generations of fighting each other as well as the roving clans of nomads that wander the Iothia Moorland, the many people of Igrathad have come to prize displays of combat prowess and hunting skills. Children are trained with arms from an early age, and the Many Hosts are respected as the finest warriors in the region, making them feared soldiers within the dynasty's military forces.

\subparagraph{Tentatively United}
While the villages of Igrathad have enjoyed the fruits of a somewhat collaborative society under the guidance of the dynasty, the centuries of turmoil have not been forgotten. Neighbors argue over territory, and youths challenge each other for social dominance. The people take pride in the unique history of their own village, and that pride can win out over the dynasty's message of unity. Interactions between the villages often grow tense when the Kryn aren't around to keep the peace. This isn't helped by the fact that a record number of able-bodied villagers from Igrathad have been sent north to fight on the front lines of the war. The center of the Kryn's government is far away from Igrathad, and many of the Dens are concerned that the tenuous union between the Many Hosts is crumbling.

\subparagraph{Government}
Each of the seven villages has its own leader, who is usually chosen for their lineage or with a contest of strength. Each leader is assigned a representative from Rosohna to help enforce the will of the Bright Queen, though the leaders often disagree with various facets of the dynasty's decrees, making total unity challenging. The various leaders meet in a public space to make most major decisions and judgments: they debate and sometimes duel to resolve the matters before them. While this unconventional court is not consistent with the dynasty's intended order, few of the Kryn representatives find the will to argue.

\subparagraph{Crime}
Some Aurora Watch soldiers keep the peace and manage the day-to-day quarrels within Igrathad. However, the majority of the Hosts' people are expertly trained warriors, so they both handle their own problems and make it difficult to enforce any laws on them. The idea of attempting to punish or incapacitate a criminal ogre or angered giant keeps all but the strongest from taking a post in Igrathad.

\subparagraph{Geography}
The seven villages are connected by muddy paths that cut through the rocky grassland. Each village has its own culture, complete with distinct customs and religious practices:

\textbf{Kranad} began as a bugbear clan of the same name and is made up of wooden shanties built around stands of trees.

\textbf{Broonbobah} is a village of goblins, where clusters of huts and rickety wood structures climb high to oversee the surrounding plains.

\textbf{Aruuth} is home to the enduring orc people of the Aruuth-Uk clan, once the dominant clan of the area. The stone village is built around a prominent, deep-set vein of iron that provides ore to the villages.

\textbf{Oshinik-Ka} is a village of the hobgoblin folk, who reside in stone bunkers behind barricades that markedly push the militant air of life within.

\textbf{Wrathfall} is the village of humanoid outlanders who never joined with the Odarami or the clans of the Iothia Moorland. It is little more than a few scattered clusters of yurts.

\textbf{Unyeethi} is home to the gnoll clan of the same name. It is the least social of the villages. The gnolls live in primitive huts made from bark, furs, and bones.

\textbf{Dodafir} is the fortress village of hill giants that marks the southernmost borough of Igrathad. The fortress is made of great logs crudely fitted together, with roofs made of bark tiles and thatch.

\subparagraph{Igrathad Adventures}
Characters could come to Igrathad seeking help from a powerful giant, or to help keep the villages united, or, conversely, to try to drive them apart.

\subparagraph{Necromancer's Army (Mid Level)}
An exiled Dolorav necromancer has allied himself with a clan of bloodthirsty minotaurs to manufacture an army of undead, whose slow expansion poses a threat to Igrathad. Each of Igrathad's seven villages appoints a champion to march with Dolorav and Koshtask warriors against this cruel usurper and his legion of unholy monstrosities. The characters must join the fight if they want the villages to help them with some other matter.

\subsubsection{Penumbra Range}
Dividing the marshes and badlands of Xhorhas from the ghastly boneyards of the Blightshore, the Penumbra Range is an unusually steep series of black rock mountains that rise against the horizon like a huge jaw of pointed teeth. Hiking across the sharp, glassy rock is an extremely dangerous feat that would challenge even the most experienced travelers.

\subsubsection{Rosohna (Ghor Dranas)}
\begin{description}
\item[Population.]
113,190 (66\% dark elves, 9\% goblins, 7\% duergar, 18\% other races)

\item[Government.]
The twelve ruling dark elf Dens oversee the city, all at the command of Den Kryn and the Bright Queen herself.

\item[Defense.]
A mighty army of Aurora Watch protects the city, while mounted weapons and ballistae stand at the ready. Each Den employs personal guards, as well as highly trained bodyguards for their nobles.

\item[Commerce.]
Most any goods or services can be found within the city, with distant imports available on the black market.

\item[Organizations.]
The majority of temples here are dedicated to the Luxon, but smaller shrines to the rest of the pantheon are scattered throughout in the outer districts. The merchant guilds are organized by the Dens and closely tied to the faith of the Luxon. Agents of the Children of Malice are sown like seeds of chaos throughout the city.
\end{description}
\noindent
The realm of Xhorhas was home to the strongholds of the Betrayer Gods as they walked the realm of mortals, but the beating heart of their mercurial alliance was the dark city of Ghor Dranas (meaning "gathering of shadows" in Draconic). A monument to the union of tyranny and chaos, the shaded skies framed the dozens of crooked towers and war barracks that surrounded the central citadel of Virkwzual, or "Dark Bastion" in the Infernal language. The final battles of the Calamity were waged around Ghor Dranas, and much of the city was reduced to ash and dust. Towers crumbled, the ground opened to swallow the legions of the Betrayers, and the citadel was left sundered and abandoned. It wasn't until the surviving Kryn drow emerged from the Underdark that the city rose from its ashes.

The Kryn slowly cleared the ruins of the corrupt creatures that had claimed it for so long. They built new structures on the foundations of the ancient city while expanding the underground levels beneath, bringing new life to the city of Ghor Dranas, which they named Rosohna, or "Rebirth," in their tongue.

Rosohna is the center of the powerful Kryn society, a symbol of dark history reclaimed. Its tall, dusky spires and beautiful buildings rise from the broken shards of the Calamity. Breaks in the constant cloud cover offer glimpses of sunlight, which is seen as a holy gift shining on the many monuments to the Luxon. Crystalline halls of glass and quartz glow with interior light in the traditional Kryn hues of green and purple.

The streets are lively with city-dwelling wastefolk, the Luxon-worshiping drow who live beneath the open sky to spite Lolth's design, and even a few duergar and other beings from the Underdark. Local businesses thrive as the Xhorhasian community grows, and goblinoid merchants have found profit and comfort within this dark elf civilization. Even so, most people beyond the borders of Xhorhas only see Rosohna as an unbroken link in the legacy of Ghor Dranas, a city born of malevolent gods and their murderous spawn.

\subparagraph{A Collection of Characters}
As the Kryn discovered their cycle of reincarnation, many began to find beings outside the drow lineage undergoing anamnesis, the process in which a creature reincarnated by a \textit{Luxon beacon} begins to recall memories from its past lives. Unexpectedly, the spirits of heroic Kryn were reborn in various intelligent creatures outside the dark elf society. With this knowledge, the ruling dens of the Kryn Dynasty began to reevaluate their treatment of non-drow, instead seeking to bring the light of the Luxon to the wastefolk of Xhorhas. While many beings across Xhorhas and beyond still cultivate mutual animosity toward the Kryn, there is a growing number of goblins, ogres, gnolls, orcs, and many other beings who have found a new home and purpose under the rule of the dynasty. Thanks to the miracle of anamnesis, Kryn society has become a diverse and thriving community that spreads across Xhorhas.

\subparagraph{The Dusk City}
The dark elf and duergar bloodlines are averse to sunlight, which presents a challenge to a culture that journeyed to the surface out of love for their light-bearing god. Periods of worship involve occasional prayer and meditation in the warmth of daylight, but twelve hours of daylight each day can become painful and physically detrimental to those beings adapted to subterranean life. To mitigate this problem, powerful dunamancers have woven arcane shields to temporarily block out the sun above Rosohna. Periods of days, sometimes weeks, will transpire in a state of perpetual evening, enabling the darkness-bound denizens to go about their work. Such periods briefly come to an end to usher in periods of scheduled, mass worship under the sunlight.

\subparagraph{Above and Below}
Rosohna is built on the remains of Ghor Dranas, and an entire, unseen level of the city lurks beneath the surface. These caves and tunnels were used as torture dungeons by worshipers of Torog, and the Dens of the Kryn have spent centuries reclaiming the space for their own people. Much of the dungeon remains abandoned, but a great many layers have been cleared and transformed into underground neighborhoods, hubs for subterranean travel across Wildemount, and farming caverns for crops of edible fungi. However, the Kryn Dynasty has secretly allowed some hidden chambers to fulfill their original grim purpose as cells for vile criminals and political prisoners.

\subparagraph{Government}
The ruling class of the dynasty is seated in Rosohna, spread across twelve noble Dens—all of whom seek to one day rule the dynasty. They play an intricate political game by expanding their families of timeless souls and covertly undermining the authority of other Dens.

Each Den traditionally oversees a portion of local industry, but the more distinguished Dens also vie for control over aspects of spiritual or military leadership. Den Masters keep a careful eye out for weakness among other Dens, seeking ways to sully the reputations of their rivals while displaying their own worthiness to the Bright Queen. Such social sabotage is reminiscent of the old ways of the drow, but assassinations and acts of outright violence are rare. The three ruling Dens—Den Kryn, Den Thelyss, and Den Mirimm—are each helmed by an Umavi. Of those families, Den Kryn holds supreme power, for its head is Bright Queen Leylas Kryn herself.

\subparagraph{Crime}
Conflict and struggle are as common in Rosohna as in any other metropolis, and though the unifying faith of the Luxon helps maintain the peace between disparate and sometimes violent residents, old habits are hard to break. Den Mirimm has enacted a tough, sometimes brutal approach to law enforcement that helps discourage violent behavior.

Blackmail is common between political figures, and those that indulge in this vice believe that they can absolve any sins through pious prayer, reasoning that their actions are only illegal if they're caught.

Those who have been found guilty of major crimes against the dynasty are locked away in the Dungeon of Penance (see "Geography" below), or taken out beyond the reach of the beacons and executed, with no hope of rebirth.

\subparagraph{Geography}
Nestled protectively within a deep valley alongside the Penumbra Range, the majestic, ominous, and beautiful city long known as Ghor Dranas stands as a symbol of renewal and change. A vast medley of bright and wondrous structures, grandiose towers, and crystalline cathedrals soar into the sky, thriving among the broken remnants of a ruined metropolis that was once nothing but ash and haunted memories. Districts are laid out in the husk of the ancient capital like a series of layers, each ward growing more refined and enlightened the closer one gets to the center of the city. These layers are the Ghostlands, the Coronas, the Gallimaufry, the Firmaments, the Shadowshire, the Lucid Bastion, and the Dungeon of Penance.

\subparagraph{The Ghostlands}
Large sections of the foreboding obsidian walls that once surrounded the city have fallen, but they are in the process of being rebuilt as the war with the empire grows more violent. Within these walls, the outlying Ghostlands remain empty, haunted by ghosts and shadows, the razed structures and barracks of the larger, ancient city left to crumble to dust. Two functional roads known as the Hallowed Paths allow caravans, travelers, and military forces to reach the city.

\subparagraph{Coronas}
Most of the city's poorer citizens are scattered throughout this sprawling, foggy district of salvaged structures and newer homesteads. Patches of magically maintained farmland are tilled and tended by nearby denizens, while industrial factories and bellows release stacks of black smoke and faintly colored steam. Aurora Watch captains coordinate law enforcement from three fortresses, while mercenaries, volunteers, and desperate folk often strike out into the Ghostlands with hopes of exorcising new neighborhoods to shelter the continuously growing population of Rosohna.

\subparagraph{Gallimaufry}
Deeper into the heart of Rosohna is the massive region called the Gallimaufry. The rowdiest and most vibrant ward of the city, the Gallimaufry marks the intersection of culture and commerce where citizens come to share ideas, peddle wares and services, find employment, or celebrate personal victories large and small in the company of other folk from the wastes. The area is lit with ever-burning lanterns of green light strung across street corners and roadways, and the percussive music of the wastefolk mingles with the haunting melodies of drow musicians to produce an ever-shifting atmosphere that many find at once both welcoming and claustrophobic.

\subparagraph{Firmaments}
Near the center of Rosohna, travelers will find the Firmaments, where the haunted memory of Ghor Dranas has been erased with a beautiful district of glowing drow gardens, luxurious chapels to the Luxon, and the spacious halls of Marble Tomes Conservatory, the city's center of higher learning. Wealthy estates house Den members and figures of import to the dynasty, and the impressive towers of the Nimbus Keep loom over the region. The Nimbus Keep is the center of the Aurora Watch and the first of two paths to the Dungeon of Penance beneath the city.

\subparagraph{Lucid Bastion}
Located at the very center of Rosohna is the awe-inspiring and towering citadel of the Lucid Bastion. This impressive cathedral is built from the ruins of Virkwzual and now stands as the symbol of Kryn rebirth and their faith in the light of the Luxon. The Lucid Bastion's spires of gray quartz pierce the sky, and the soft interior glow that suffuses the interior chambers never darkens, shining over the city day and night. Major religious events are centered around the Bastion, for it houses the "heart of consecution," the first \textit{Luxon beacon} to have been discovered. Within these halls, the Bright Queen, her trusted Umavi, and their Dens work to govern and protect all citizens of the Kryn Dynasty.

\subparagraph{Shadowshire}
Several well-guarded stairways curl beneath the Firmaments, descending under the surface of Xhorhas into the Shadowshire, a vast subterranean district built throughout the reclaimed prisons of Torog. Less affluent dark elves live among the stone and glass abodes built into the rock and sediment. Caverns amass bioluminescent fungi to be harvested for food and medicine, while enchanted lanterns light the various levels with a soft and many-colored luminescence. A magical underground farmland has been established here, where sunless crops are grown to feed the city. The second entrance of the Dungeon of Penance is located within the Shadowshire.

\subparagraph{Dungeon of Penance}
The many levels of the Dungeon of Penance delve under the city. This terrifying oubliette bears the marks of the Kryn Dynasty's renovations on its upper levels, but become less explored and less changed from the days of Ghor Dranas the deeper one delves.

As the largest prison within the dynasty, a microcosm of the dynasty itself can be found in the complex bureaucracy of jailors, supervisors, and elite Aurora Watch soldiers who guard the frigid halls of the dungeon.

\subparagraph{Rosohna Adventures}
To outsiders, Rosohna is a city of bizarre wonders, but to its residents, it is simply home—a place of business, crime, worship, and more. All kinds of adventures await in this diverse city.

\subparagraph{An Early Dawn (Any Level)}
While the characters are visiting Rosohna, the shadows above the city suddenly disappear, letting in harsh light that hinders the \textbf{drow} and \textbf{duergar} residents. Is this phenomenon part of an imminent attack, or did the Kryn Dynasty anger the Luxon and must figure out a way to regain their god's favor?

\subparagraph{A Bloody Dawn (Mid Level)}
During a particularly clear day, the priests of the Luxon dispel the darkness above the city for a moment of prayer. While the denizens bask in the light, agents of the Children of Malice strike, attacking and trying to kidnap influential members of the dynasty's clergy. The characters must stop the assault, root out the hidden traitors behind this plan, and save any surviving prisoners.

\subparagraph{Ghor Dranas Returns (Epic Level)}
The Cerberus Assembly sends several disguised archmages to reopen the sealed rifts in Rosohna, hoping to destroy the dynasty by filling the city with otherworldly threats. The characters might be part of the plot, or they might work to stop it.

\subsubsection{Ruins of Draconia}
At the bottom of the deep Dreemoth Ravine lie the crumbled ruins of the once-proud nation of Draconia. The floating islands of the city-state were the original home of dragonborn society within Wildemount, though the ruling class installed a strict caste system and forced the tailless ravenite dragonborn to labor in the ravine below their magically elevated island city. When the city was assaulted by the Chroma Conclave two decades before, the islands were sundered and fell into the ravine, killing many of the inhabitants. However, this disaster gave the ravenites a chance to rebel against their cruel masters.

The plundered ruins of the former city are now abandoned heaps of rock and rubble, a dismal memorial to the ignominious end of a glorious, cruel civilization.

\subparagraph{Not Forgotten}
The ravenites did have allies among the ruling houses. Many Draconian citizens disagreed with the oppression of their brethren and fought for generations to overturn the laws that kept them enslaved. Others aided the ravenites in their emancipation after the fall of Draconia; those folk came out on top in the end, for they're considered honored members of the emerging ravenite society.

Those champions who were lost in the struggle are immortalized in statues left among the ruins, the most famous of whom is the hero Tiberius Stormwind of Draconia. Tiberius was one of the legendary heroes of Vox Machina, and he gave his life to defend ravenites from the ire of Vorugal, the white dragon of the Chroma Conclave. His statue and surrounding shrine are a place of pilgrimage for ravenite historians.

\subparagraph{Ruins of Draconia Adventures}
Those brave enough to delve into the ruins of Draconia can find treasures and a few ghosts of the destroyed city.

\subparagraph{Former Glory (Low Level)}
A cabal of disaffected draconblood dragonborn have been secretly assembling in the ruins of their former metropolis. They are guided by the whispers of a \textbf{death knight} who promises them the power to reassert control over the ravenites. They have begun learning necromancy and even making minor sacrifices in his name. The death knight, Pentrakath, is currently gathering artifacts elsewhere in Wildemount, but will soon use his cultists and cult fanatics to carry out his twisted plans, if no one intervenes (see "\textit{Dreemoth Ravine Adventures}" earlier in this chapter).

\subsection{Areas (S-Z)}
\subsubsection{Sorrowseep Waters}
The Sorrowseep Waters are one of the twin lakes that dominate the central wastes of Xhorhas. These dark waters are home to a curious ecosystem of swamp beasts, amphibious behemoths, moorbounders, and apocryphal creatures said to live within the lake itself. Fishers, hunters, and trappers often brave the surrounding knotted boughs and stinking mud of Sorrowseep Marsh in search of unique prey, while healers or their hirelings also risk the journey to collect the waters of the lake, which are believed to have restorative qualities. The crisp waters of the lake branch out across the wastelands of Xhorhas, feeding rivers and sustaining life throughout the scarred plains.

\subparagraph{The Tanglebed}
The true depth of the Sorrowseep has yet to be discovered due to the density of vegetation near the bottom, called the Tanglebed. Scholars believe that something ancient and powerful sleeps at the bottom of the lake, and that such a creature could be the source of the lake's unique plant growth, unnaturally clean waters, and possible healing properties.

\subparagraph{Sorrowseep Waters Adventures}
Characters exploring the Sorrowseep Waters encounter all manner of strange creatures and eccentric denizens.

\subparagraph{Something Borrowed (Low Level)}
One of the characters is contacted by a close childhood friend or family member to announce their sudden betrothal to a passionate love interest. This union is taboo, however, for the love interest is actually a \textbf{green hag}. The character is invited to be a member of the wedding party—but something about this marriage seems fishy.

\subparagraph{Entangled (High Level)}
A gnome inventor claims to have created a submersible vessel that can survive the lightless journey through the Tanglebed, past the known depths of the Sorrowseep. He hires the characters to help him helm the machine. Together, they discover a mysterious location beneath the lake, where a \textbf{marid} is gathering a cult to purge the nearby rivers of civilization.

\subsubsection{Urzin}
\begin{description}
\item[Population.]
5,930 (62\% goblinoids, 21\% gnolls, 10\% orcs, 7\% other races)

\item[Government.]
The reigning champion of Urzin, Sunbreaker Olomon, guides the townsfolk, who follow with zeal and dedication.

\item[Defense.]
A legion of Aurora Watch defends the town, while the hordes of Urzin hunters scout the marsh. The horizonback tortoises that bear the town on their backs are more than capable of defending themselves.

\item[Commerce.]
Trade in Urzin is fairly basic and geared toward survival and warfare. Many common goods can be bought or traded for, but rarer items are coveted and appear infrequently. Inns are very rare.

\item[Organizations.]
Some small Luxon temples travel on tortoise-back alongside shrines to older goblinkin patrons.
\end{description}
\noindent
A uniquely mobile settlement with its permanent dwellings built on the backs of dozens of horizonback tortoises, Urzin is home to swamp hunters, scouts, and trappers that know the Brokenveil Marsh like no one else. The most recent society to come under the wing of the Kryn Dynasty, Urzin acts as a powerful, mobile fortress for the Kryn military throughout the Brokenveil Marsh. The settlement rarely stays in one location for more than a month, and its people are skilled at the art of camouflage and ambush. Urzin has become a major boon to the dynasty as the conflict with the Dwendalian Empire escalates.

\subparagraph{Respect for Strength}
The swampfolk of Urzin have always prized strength and guile: their leader is simply the reigning champion of one-on-one combat. The previous champion, Ogre Lord Buhfal, ruled the settlement for decades, assaulting the encroaching empire and fighting off the dangers of the marsh. Roughly forty years ago, Buhfal was challenged by the Kryn paragon Sunbreaker Olomon. In what is still retold as legend, a mighty bout between the two ended with Buhfal defeated and Olomon elevated to rule the settlement.

With his guidance, Olomon brought the teachings of the Luxon, in the name of the Bright Queen, and promised that with the hearts of the Urzin people behind him and his drow, the terrible soldiers of the empire would no longer terrorize the marshfolk. Renewed by this promise, Urzin follows the lead of the dynasty with an impressive fervor, working like clockwork to achieve any goal dictated by Olomon. In his absence, \textbf{Ogre Lord Buhfal II} proudly fills the role of leader and protector.

\subparagraph{A Shifty Settlement}
A large number of the denizens of Urzin make their homes atop the shells of the horizonbacks, but many also carry temporary tents and makeshift abodes that allow a home to be made wherever they find themselves. Whole packs of tortoises reconfigure daily, and often the settlement breaks up and scatters across the marsh for periods of time, leaving satellite outposts all over the region before reforming elsewhere among the muck. Because of this, the Dwendalian military finds the movements of Urzin hard to track and even harder to predict.

\subparagraph{Government}
While there is a comparatively thin presence of authoritarian Aurora Watch soldiers among the ever-shifting boundaries of Urzin, Sunbreaker Olomon and his appointed council are treated with reverence and their orders are followed with impressive zeal. The culture of Urzin prizes unity and adherence to the wishes of their champion, so Urzin's governance at times seems strangely smooth compared to other locations in the dynasty. While Olomon is away on other business, the Urzinians can become restless or slower to follow instruction, but this insubordination lasts only until his return.

\subparagraph{Crime}
Criminal activity within the settlement is rare, given the zealous nature of the marshfolk. Dissent within the ranks is often punished with exile, so the few that do disobey find themselves ostracized and left to survive alone among the foggy dangers of the marsh.

\subparagraph{Geography}
The town of Urzin consists of around forty horizonback tortoises and many hundreds of mobile tents and huts that shift with every relocation. The backs of these tortoises are covered front to back in rickety shacks and barracks, sometimes stacked in wild towers, with long ropes and vines anchoring them in place. The horizonback shells are reinforced at the edges with jagged pikes and sharpened bones to help defend those who live atop, while bolt throwers and short-range catapults can be found throughout the neighborhoods. Every tortoise contains its own microcosm of a village, and denizens all take up skills to make the miniature community nearly self-sufficient, complete with moss farmers, blacksmiths, and physicians.

Note that any specific indicator of Urzin's location on a map is an approximation of its location, due to the roving nature of the community.

\subparagraph{Urzin Adventures}
Just finding the settlement and winning the trust of the people of Urzin is an adventure. Characters on the run could also try to hide in Urzin, since the settlement does the running for them.

\subparagraph{Horizonback Horror (Mid Level)}
A \textbf{horizonback tortoise} has broken away from Urzin's fleet and is rampaging through Jigow. The orcs and goblinkin fought to contain it, but every wound they inflict on the beast seems to stitch together before their eyes! An \textbf{oni} that once led a clan of trolls the orcs defeated several months ago is out for revenge and devised an alchemical infusion created from troll blood. The oni administered the toxin to the tortoise, granting it the Regeneration feature of a troll and causing it to go berserk, terrifying not only the denizens of Jigow, but also the people clinging to their homes atop the frenzied tortoise.

\subsubsection{Vermaloc Wildwood}
The Vermaloc Wildwood is an enormous timberland of vibrant trees with deep purple and maroon trunks and a bright canopy of red-orange leaves. This curious and beautiful burst of color stands in stark contrast to the often dreary landscape of Xhorhas. Many denizens of Xhorhas have camps or small settlements along the treeline, as the forest is one of the major sources of wood and logging throughout the wastes. Those who live here, however, remain alert for the many deadly beasts and creatures that call this colorful expanse home. Monstrous predators twisted long ago by the Betrayer Gods hunt from the trees and brush, while the deeper woods are overrun with the webs of the malicious giant spiders that claim those shaded regions. Only the most experienced and well-armed hunters dare wander the depths of the Vermaloc.

\subparagraph{Vermaloc Wildwood Adventures}
The characters might hunt in the Vermaloc Wildwood to prove their mettle or be forced to survive a journey through the wood.

\subparagraph{Ettercap's Web (Low Level)}
Aramue, a \textbf{dryad} of the Vermaloc Wildwood and devout follower of the Wild Mother, possesses a unique gift: she can send dreams to other worshipers of Melora. Using this power, she reaches out to such individuals. Aramue needs help ridding her forest of an \textbf{ettercap} that has taken control of many of the forest's giant spiders. The spiders attack the forest's outlying settlements and bring prisoners to the ettercap, which it sacrifices to gain Lolth's favor. The dryad is kept awake at night by the ettercap's screaming victims.

\subsubsection{Verstglade}
A small valley of crooked, leafless trees is hidden in the tallest peaks of the southern Ashkeeper Peaks, though few know of its existence. During the Age of Arcanum, this dark and menacing wood once was an experimental playground for the terrible lich Vecna, where he would manipulate and twist all manner of life, arcana, and undead to his sick whims. The countless years of misuse left the withered forest corrupt and overrun with nightmarish abominations and cursed, undying minions who were abandoned to their madness. Mercifully, Vecna was banished before he could unleash his creations on Wildemount, and his aspirations for godhood left his intended plots with Verstglade forgotten.

\subparagraph{Verstglade Adventures}
Verstglade is filled with forgotten experiments that characters might be asked to retrieve or destroy.

\subparagraph{The Whispered Study (High Level)}
Cloaked by illusion and filled with deadly traps, one of the primary laboratories of Vecna remains hidden within the Verstglade. It is home to many powerful artifacts and arcane secrets, and the promise of such mighty rewards draws a bold few to attempt to locate it. Of the few mortal minds that even know of Verstglade and what lies hidden within it, only brave adventurers or the minions of evil are foolish enough to enter the corrupted woods.

\subsubsection{Xarzith Kitril}
\begin{description}
\item[Population.]
9,110 (93\% dragonborn, 7\% other races)

\item[Government.]
A council of six oversees the city, with the Wardens of the Ravine keeping order under their guidance.

\item[Defense.]
An army of Wardens of the Ravine are ever at the ready to defend the city. Many self-trained citizens can form a militia if need be.

\item[Commerce.]
Trade is mostly limited to the city, making rare goods hard to find. Most basic equipment, arms, and services can be found here, as well as taverns and lodging.

\item[Organizations.]
Small, rough sanctums of worship can be found within the city. Unions of artisans and merchants work together on a small scale.
\end{description}
\noindent
Shortly after the fall of Draconia, when the lords of the ruined city-state were deposed and scattered, the remaining ravenite dragonborn took the reins of their new freedom and began to develop a new dragonborn society in Wildemount. The ravenites have spent nearly two decades building a new city of hope for the dragonborn survivors within the cliffs of the Dreemoth Ravine. The fortified stone city of Xarzith Kitril continues to grow, and its borders are constantly being rebuilt to accommodate its exploding centers of residence and industry amid the ore-bearing mines and tall smoke pillars from ever-burning forges. The leaders of the city, members of the newly formed Scars of Scale and Tooth, work to bring unity to the fractured, feuding ravenites.

\subparagraph{Freedom from Two Terrors}
When the city of Draconia fell, the surviving ravenites enjoyed a brief moment of freedom before the terrifying, ancient white dragon Vorugal the Frigid Doom claimed the ruins and the surrounding ravine as his new lair. The ravenites were freshly bound by fear under their new overlord, whose presence plunged the ravine into an endless winter. It wasn't until Vorugal was slain by the heroes of Vox Machina that the ravenites' chains were truly shattered. As the heavy ice melted away, the new city's construction began, and the city took the name Xarzith Kitril, Draconic for "ice born."

\subparagraph{Growing Pains}
Xarzith Kitril represents a new dawn for the dragonborn people, and many of the citizens place personal pride in what their city stands for and what they believe it will eventually become. However, the council of new leadership, known as the Scars of Scale and Tooth, continues to debate over the direction and development of the city, unable to agree on a clear vision of the future. Some wish to see the city rise in prominence and respect. Others wish to keep the city isolated. The populace grows restless and frustrated without united leadership.

\subparagraph{Government}
With the support of the ravenites behind them, a council of six agreed to form the Scars and govern the city. Ruling by committee, however, sometimes leads to disagreements and slow progress. In particular, Vermosi Runek and Druvis Koothalok differ greatly on the perceived priorities of the city. Vermosi rules with religious passion and wants to remain in isolation to spite the worldly ways of the previous rulers of Draconia, while Druvis seeks the benefits from forging alliances with neighboring powers. As this and similar arguments continue to escalate, governing the growing city is becoming challenging.

\subparagraph{Crime}
Local law enforcement is comprised entirely of a force of volunteer soldiers called the Wardens of the Ravine. The Wardens truly believe in the future of the city and the Scars of Scale and Tooth, and they vigilantly maintain the peace and prevent criminal activity. Unfortunately, resources are relatively scarce, and the recent introduction of trade with the Menagerie Coast has brought an influx of coin into the previously barter-based society. As such, theft and racketeering are on the rise.

\subparagraph{Geography}
Constructed within the thinner, southern end of the Dreemoth Ravine, this city of sturdy stone structures has expanded into the walls of the ravine on either side. Abodes are stacked along the cliffs and entwined with winding stairways and gondola platforms that glide across the open crevasse, elevated with recovered chunks of brumestone that once held the floating islands of Draconia aloft.

The central, base region of the city is referred to as the Roots, where the majority of commerce and industry meets the political core of the Scars of Scale and Tooth. The West Wall region is a tall, nearly vertical neighborhood of homes built into the rocky western cliff face of the ravine above the Dreemoth Quarry. The East Wall region looks similar to its western counterpart, though at a slightly less extreme grade, and houses three separate major mine shafts that are peppered with residential caves and abodes. Dozens of cords and chains stretch from ravine wall to ravine wall, allowing people and goods to move between them on the floating brumestone platforms.

\subparagraph{Xarzith Kitril Adventures}
All manner of adventures and challenges await within this recovering dragonborn society.

\subparagraph{Best Left Alone (Mid Level)}
 One of the prominent mining companies within the city has recently uncovered a mysterious purple stone of unknown origin, about fifteen feet in diameter. Dozens of superstitious miners have fled this site, so the company hires the characters to investigate the stone and uncover its nature. When any of the stone's smooth, reflective surface takes damages, this buried relic of the Shadowfell begins to summon shadowghasts and gloomstalkers at an increasingly fast rate. The characters must discover a way to repair and reseal the stone, destroy it, or banish it from Exandria before the city is overwhelmed.

\section{Blightshore}
The blackened seaboard that forms the easternmost edge of Wildemount is a savage wasteland struggling to recover from the arcane havoc wrought by the Calamity so many years ago. The Miskath Strand was once a temperate expanse of arboreal pulchritude and coastal tranquility, but it has since become a lonely stretch of hazardous terrain and maddening isolation. This coastline of cursed earth and the surrounding area have come to be known as Blightshore, a fitting moniker for the once-proud seaside paradise defiled by the Betrayer Gods. Here, among the Miskath Strand's smoldering remains, explorers vie for the opportunities offered up by Blightshore's weird new varieties of flora and fauna—not to mention the hidden wealth that lurks within the lightless vaults of the Betrayers themselves. Life in Blightshore is difficult at best; for many, it represents a true test of spirit as they weather the persistent chaos of Blightshore's terrain.

While the Dwendalian Empire and its constituents seek to establish a foothold in the more hospitable regions of the Miskath Strand, these efforts are consistently hampered by Blightshore's less scrupulous denizens. Remote operatives of the Myriad and the Revelry, intrepid pilgrims from the Kryn Dynasty, outcasts from the fallen city of Draconia, and goblinoid nomads from the Xhorhasian wastes form the motley civilization of Blightshore. Moreover, this land is home to many undead horrors and wandering hollow ones.

\subsection{Bizarre Terrain of Blightshore}
The strange landscape of the Miskath Strand has been savagely reshaped by the Calamity, infusing the earth with chaotic new elements and especially inhospitable terrain. As a result, much of this region of Eastern Wynandir remains relatively unmapped by modern scholars.

The terrain and weather features described in the sections that follow can be encountered anywhere within Blightshore's defiled expanse. Many of these areas present hazardous (and potentially fatal) conditions for travelers and explorers. As your players traverse the desolate badlands of the Miskath Strand, you can present Blightshore's weird and wondrous terrain as a character in its own right. And perhaps there are even stranger features of this accursed frontier that wait to be discovered, tamed, or avoided at all costs.

\subsubsection{Earth Motes}
These chunks of broken land hover and swirl in the air, held aloft by unknown arcane forces. A typical earth mote is 1d6 × 10 feet in diameter and hovers 2d10 × 10 feet above the ground. Earth motes are capable of supporting the weight of modest structures as well as creatures.

\subsubsection{Electric Wind}
The galvanic breezes of this supernatural gale are latticed with currents of lightning, which stab through the air with primordial fury.


\begin{itemize}
\item Ranged attacks through electric wind are made with disadvantage.
\item Any creature that starts its turn surrounded by electric wind has a 5 percent chance of being struck by lightning. A creature struck by lightning must make a DC 15 Dexterity saving throw, taking 4d6 lightning damage on a failed save, or half as much damage on a successful one.
\end{itemize}

\subsubsection{Ember Thicket}
A dense grove of blackened, leafless trees perpetually smolders with the dull glow of arcane embers. This copse is choked by a cloud of ash and cinders, which constantly rain on the cremated forest floor.


\begin{itemize}
\item An ember thicket is a lightly obscured area.
\item Creatures gain one level of exhaustion for every 10 minutes spent inside an ember thicket without suitable protection from the heat, smoke, and ash.
\item Ranged attacks through an ember thicket are made with disadvantage. A ranged weapon attack made through or within an ember thicket automatically misses a target beyond the weapon's normal range.
\end{itemize}

\subsubsection{Etheric Rift}
An etheric rift appears as a light fog or spectral haze. The interplanar tracts of this anomaly reach in and out of the Border Ethereal. Unseen trenches and invisible paths overlap throughout its vaporous expanse.


\begin{itemize}
\item An etheric rift is a lightly obscured area.
\item Any creature that enters an etheric rift or starts its turn there has a 33 percent chance of gaining the effect of the \textit{etherealness} spell until the end of its next turn.
\end{itemize}

\subsubsection{Gravity Funnel}
This cylindrical area is affected by the unseen pull of heightened gravitational force. A typical gravity funnel is 3d10 × 10 feet tall and 5d10 × 10 feet in diameter.


\begin{itemize}
\item Strength checks made within a gravity funnel have disadvantage.
\item A ranged weapon attack made through or within a gravity funnel automatically misses a target beyond the weapon's normal range.
\item Creatures and objects that are fully contained within a gravity funnel have resistance against force damage.
\item The gravity funnel counts as difficult terrain.
\end{itemize}

\subsubsection{Psychic Miasma}
A virtually imperceptible mist of preternatural psychotropic vapor hovers above the ground. This coruscating miasma corrupts the sanity of the living and fortifies the grim resolve of the undead.


\begin{itemize}
\item Creatures that are fully immersed in psychic miasma have vulnerability to psychic damage.
\item Any creature that enters a cloud of psychic miasma or starts its turn there must succeed on a DC 15 Constitution saving throw or become afflicted with a random form of long-term madness (see "\textit{Madness}" in chapter 8 of the \textit{Dungeon Master's Guide}). Creatures that can't be poisoned automatically succeed on the saving throw. A creature that successfully saves against the psychic miasma is immune to its effects for 24 hours.
\item An undead creature that spends at least 1 minute in a cloud of psychic miasma gains advantage on ability checks and saving throws for 24 hours.
\end{itemize}

\subsubsection{Rust Rain}
This crimson precipitation is corrosive to ferrous metal and caustic to creatures.


\begin{itemize}
\item Rust rain corrodes nonmagical, ferrous metal objects. If the object isn't being worn or carried, the rain destroys 1 cubic foot of it for every 10 minutes of exposure. An object must be covered to avoid exposure. If the exposed object is nonmagical metal armor, a nonmagical metal weapon, or a nonmagical metal shield being worn or carried, it takes a permanent and cumulative −1 penalty to damage rolls or the AC it offers for every 10 minutes of exposure to rust rain. A weapon that drops to a −5 penalty, armor reduced to an AC of 10, or a shield that drops to a +0 bonus is destroyed.
\item Any creature that starts it turn in rust rain and isn't covered by an umbrella, a shield, or some other form of rain protection takes 1d4 acid damage.
\end{itemize}

\subsubsection{Toxic Fungus Colony}
A virulent expanse of malignant fungal growths blankets the landscape and fills the air above it with spores.


\begin{itemize}
\item Any non-plant creature within 20 feet of a toxic fungus colony is poisoned. A creature poisoned by the colony for 1 hour has a 20 percent chance of gaining a random form of short-term madness (see "\textit{Madness}" in chapter 8 of the \textit{Dungeon Master's Guide}).
\item A toxic fungus colony can be bioluminescent. A bioluminescent colony sheds dim light within its area and for an additional 20 feet beyond its outer edges.
\end{itemize}

\subsubsection{Uncanny Phenomena}
The arcane fallout from the Calamity is so widespread throughout Blightshore that the very land, sea, and air have been subtly altered in strange and mysterious ways. As characters explore the post-Calamity Miskath Strand, roll on the Uncanny Phenomena table as frequently as you like to augment the world they encounter. If you roll something that doesn't make sense, roll again, choose a more appropriate entry, or use the result as inspiration to make up your own phenomenon. The magical effects of these phenomena can't be dispelled.

\begin{DndTable}[header=Uncanny Phenomena]{cX}

d10 & Phenomenon\\
1 & It smells of death here. This area is permanently ridden with the stench of rot and decay.\\
2 & Incorporeal phantoms move in and out of reality here. These apparitions ignore any interaction with creatures.\\
3 & Everything is a single color here. All physical matter in this area adopts the pigmentation of one of the seven colors of the rainbow or the color black.\\
4 & Nonmagical fire can't burn here.\\
5 & Sources of light here appear kaleidoscopic. Empty air refracts light in this area, creating dazzling displays of prismatic colors.\\
6 & Any character who sleeps here dreams of a different life on a different world in a different body.\\
7 & There is no sound here. This area is permanently under the effects of the \textit{silence} spell.\\
8 & This place is supernaturally unlucky. Attack rolls, ability checks, and saving throws cannot be rerolled in this area.\\
9 & Water turns to acid here. Water introduced to this area becomes acid in 1d4 days.\\
10 & This place is haunted by fragments of the Far Realm. Whenever the target of a Wisdom saving throw fails its save, that creature has a 5 percent chance of becoming afflicted by a random form of long-term madness (see "\textit{Madness}" in chapter 8 of the \textit{Dungeon Master's Guide}).\\
\end{DndTable}

\subsection{Cauldron Sea}
A perpetual storm chokes the stony shores of this dark and brackish expanse, where aspects of the Elemental Plane of Water and the Abyss collide in constant entropic turmoil. Aquatic terrors lurk in its murky fathoms, which are rumored to hold everything from the ruined battlements of a sunken city to the haunted graveyard of a forgotten fleet.

Although the waters of the Cauldron Sea present challenges to those who would tame its tides, a handful of unmapped (and often unnamed) villages are scattered along its savage shores. Some of these settlements offer safe havens to weary travelers of the northern Strand, prospectors and pirates alike, but most treat outsiders with disdain or even outright hostility.

\subsubsection{Oblivion Vortex}
Soothsayers claim that Tharizdun, the Chained Oblivion, created the Cauldron Sea to loosen the veil between worlds and warn that one day he will rend the very fabric of reality from the depths of its accursed waters. There are those who would see this mad prophecy fulfilled, but many more would do anything to prevent the Betrayer God's malevolent return.

\subsubsection{Cauldron Sea Adventures}
The tumultuous waters of the Cauldron Sea breed all manner of creatures and crises for intrepid characters to overcome.

\subparagraph{Presage in a Bottle (Mid Level)}
Unlike the restorative essences of the Sorrowseep Waters, a draught of liquid from the Cauldron Sea is said to induce madness along with various other arcane effects. A world-class alchemist and chirurgeon, the eccentric Brannon Ven Brava, seeks to analyze the primordial liquid's supernatural essence firsthand and will pay handsomely for an escort through the wilds of Blightshore to secure it.

\begin{DndSidebar}[float=!b]{Blightshore Contraband}
The unique environs of the post-Calamity Miskath Strand have produced a number of strange new substances that can be cultivated for both noble and illicit purposes. The regulation of these bizarre drugs proves difficult, however, as their exotic qualities often confound those who would subvert their questionable usage. As such, supply and demand remain steady despite the inevitable hazards of production.



\textit{Black Sap}. This tarry substance harvested from the dark boughs of the death's head willow is a powerful intoxicant. It can be smoked as a concentrate or injected directly into the bloodstream. A creature subjected to a dose of black sap cannot be charmed or frightened for 1d6 hours. For each dose of black sap consumed, a creature must succeed on a DC 15 Constitution saving throw or become poisoned for 2d4 hours—an effect that is cumulative with multiple doses. Black sap is valued at 300 gp per dose.



\textit{Blight Ichor}. This bitter chartreuse concoction is distilled from a fungus native to the Blightshore badlands. The sickly green liqueur harbors potent psychedelic properties. Provided it is neither a construct nor undead, a creature subjected to a dose of blight ichor gains advantage on Intelligence and Wisdom checks, as well as vulnerability to psychic damage, for 1 hour. For each dose of blight ichor consumed, the creature must succeed on a DC 15 Constitution saving throw or become poisoned for 1d6 hours and suffer the effects of a \textit{confusion} spell for 1 minute. An undead creature subjected to a dose of blight ichor gains advantage on all Dexterity checks and is immune to the frightened condition for 1 hour. Blight ichor is valued at 200 gp per dose.



\textit{Soothsalts}. Soothsalts are derived from a naturally occurring crystalline substance discovered throughout the wilds of the Miskath Strand. The crimson crystals have been mined from cavernous veins like those in the mouth of the Miskath Pit and found within smaller geode formations near sites ravaged by the Calamity. Soothsalts are consumed orally in lozenge-sized doses, and frequent users can be identified by the telltale crimson stain around their mouths. A creature subjected to a dose of soothsalts gains advantage on all Intelligence checks for 1d4 hours. For each dose of soothsalts consumed, the creature must succeed on a DC 15 Constitution saving throw or gain one level of exhaustion—an effect which is cumulative with multiple doses. Soothsalts are valued at 150 gp per dose.



\end{DndSidebar}

\subsection{Ebonglass Massif}
The sharp obsidian peaks of the Ebonglass Massif spell doom to those who would attempt to tame them, but many still try. This jagged range of lustrous volcanic rock divides northeastern Blightshore into two large peninsulas: Far Hharom to the north, and the Tooth of Zehir to the east. Expeditions to the Ebonglass Massif from the nations of Western Wynandir and beyond have been fraught with peril, and the craggy mountains are considered all but impassable by cautious outlanders, curious prospectors, and others who would explore their mysteries.

Like many places in the Miskath Strand, the Massif is haunted by a dark past. Historians have unearthed tales of ancient dwarven strongholds that were utterly destroyed by the upheaval during the Calamity. Now, the only known denizens of the Ebonglass Massif are the birds of prey that scour its peaks for sustenance and a maniacal clan of gray-scaled kobolds that employs simple technology—hooks, cables, and gliders—to traverse the Massif's treacherous terrain.

\subsubsection{Fang of the Worm}
The kobolds of the Massif live in a subterranean sprawl on the southwestern edge of the Ebonglass foothills. Known as the Fangs of the Worm, they are led by a kobold queen whose eldritch magic is sustained by the clan's central religious totem—a powerful artifact from ancient Eiselcross, which has no doubt played a large part in the wayward sect's rapid evolution and questionable sanity. The kobolds of the Fangs of the Worm gather around this massive cylinder of metal and light to pay fealty to their queen and honor the visions she receives from beyond. This columnar artifact is actually a cryogenic containment unit that houses and preserves a fetid tissue sample from the festering body of Quajath, the Undermaw, and what the kobolds of the Ebonglass Massif mistake for divine visions are actually the abhorrent murmurs of Quajath, who telepathically guides the clan's erratic actions.

\subsubsection{Phantom Wardens}
The dwarves of the Ebonglass Massif persist as incorporeal caretakers in their hidden subterranean halls, where they carve the volcanic glass with psionic powers like an artist moves paint on canvas. These forlorn specters preside over an empire of impeccable craft and staggering natural wealth, but their physical tombs have no entrances from the outside world, much to the chagrin of would-be burglars.

\subsubsection{Ebonglass Massif Adventures}
Adventurers who find themselves near the peaks of the Ebonglass Massif have strayed very far indeed from the political turmoil of Wildemount. The nigh impassable mountains hide strange creatures and even stranger secrets.

\subparagraph{Lost Xagonstar (Low Level)}
The adventurers travel to Ebonglass Massif by airship to join a well-funded archaeological expedition mounted by Duurek Kronst, a wealthy dwarf who is determined to find the dwarven stronghold of Xagonstar, lost and buried since the Calamity. An unnatural storm sweeping across Blightshore causes the airship to crash, but the characters survive. Once they reach the dig site, they see that great progress has already been made. The first great hall of Xagonstar has been discovered! Unfortunately, the archaeologists are missing and the site has been abandoned. Were the diggers captured by the Fang of the Worm or carried off by a hungry \textbf{gloomstalker} (as described in \textit{chapter 7}), or did something slip out of Xagonstar to devour them?

\subsection{Far Hharom}
Beyond the Ebonglass Massif's hyaline peaks, a lonely expanse of moorland stretches northward toward the icy waters of the Frigid Depths. Though the landscape is drab and featureless, the peninsula known as Far Hharom is littered with the barrows of a long-forgotten barbarian people, who vanished centuries before the Calamity ravaged the Miskath Strand. Here and there, gelid bogs home to bizarre flora and fauna dot the landscape. And all around, an obstinate air of doom lingers among the peninsula's perpetual howling winds.

Far Hharom is a dangerous realm, even by Blightshore's standards, and many suggest that the impasse formed by the Ebonglass Massif is actually a blessing in disguise. The only charted settlement in Far Hharom is an abandoned village on the rocky island west of the peninsula, known as Olc Skerry (described later in this chapter). Travelers have been known to camp amid the stony ruins—but these interlopers leave as soon as their affairs allow, typically preferring the haunted expanse of Far Hharom's open fells over the eerie isolation of Olc Skerry's craggy shores.

\subsubsection{Barrows of the Betrayers}
The ancient burial mounds scattered across Far Hharom are rumored to be haunted by restless undead that were animated just as the arcane meddling of the Betrayer Gods reached its abominable zenith. Though the antiquities buried within Far Hharom's ancient barrows promise considerable wealth to those who would plunder their lightless depths, few excavators of these accursed places actually live to tell the tale.

\subsubsection{Domains of the Deep Ones}
There are rumors of a frightening clan of amphibious humanoids who dwell along the icy northern shores of Far Hharom. The stories speak of a hoary clan of pale, degenerate folk who pay fealty to dark, arctic demigods in the form of warm-blooded sacrifices. Some believe the Ebonglass Massif was erected by the Betrayer Gods to keep these strange and sinister people at bay, while others report to have seen raiding parties of scaly Far Hharom monstrosities as far south as the Cauldron Sea.

\subsubsection{Far Hharom Adventures}
The northern peninsula of Blightshore is a remote and unwelcoming place, but characters in search of secret histories and treasured lore will find it most inviting.

\subparagraph{Deep Ones of Far Hharom (Low Level)}
A coastal storm causes the characters' ship to run aground on the northern shore of Far Hharom. As they take stock of the damage and make plans to head inland, waves of \textbf{kuo-toa} converge on the ship and attack them.

\subparagraph{It Kills at Night (Mid Level)}
While exploring a Far Hharom barrow, the characters are surprised to discover a group of spooked travelers who have taken shelter within. Instead of the subterranean wights and poltergeists of legend, these survivors fear an even greater threat that stalks the moorland at night—something so supernaturally horrible, they must huddle in darkness every evening to evade it as they wearily cross the accursed landscape.

\subsection{Ghor Veles}
The ruined keep at the heart of this crumbling city once served as an arcane repository and research facility for the Betrayer Gods during their pre-Calamity ascent to power. Artifacts and wondrous items from the Age of Arcanum are rumored to linger within its shadow-haunted vaults, lying in wait for those brave or foolish enough to challenge the citadel's timeless guardians. Now, its decaying battlements stand as a stark monument to the Calamity's unstoppable devastation.

The name Ghor Veles means "Vault of Shadows," and a vast dungeon indeed lurks below the ruins of the city. These catacombs were home to the aberrant and undead minions of the dracolich \textbf{Oracs the Enduring}. Countless adventuring parties and would-be tomb robbers have attempted to lay siege to the vaults below Ghor Veles, only to perish at the clawed hands and creeping doom of the keep's fearsome inhabitants. But a destiny worse than death awaits them: most of these ill-fated interlopers become undead minions themselves, bound to Oracs in unending servitude. As this army grows, so does Oracs's hold over the Miskath Strand.

\subsubsection{Land of Refugees and Hermits}
Ghor Veles is pockmarked by crumbling houses and dilapidated fortresses, and countless people displaced from elsewhere in Xhorhas by war or exile have made this blighted battlefield their home. A notable resident of the ruins surrounding Ghor Veles is a dranassar survivor named Liv-Burrath. She uses magic to disguise herself as a dark elf to avoid attracting attention, but is covertly collecting arcane relics from Ghor Veles with the hope that she can one day reunite with other dranassar and retake their ancient home of Ghor Dranas. She pays a fair sum for any artifacts that adventurers gather from the ruins.

\subsubsection{Ghor Veles Adventures}
Only the undead horrors beneath Ghor Veles rival the cutthroat grave robbers that characters might encounter among the crumbling ruins.

\subparagraph{Shadow Essence (Mid Level)}
The characters are hired by Cerberus Assembly sages in New Haxon to gather a rare magical ingredient—shadow essence—that can be collected from dead gloomstalkers. As gloomstalkers are known to circle above Ghor Veles like vultures, the sages ask the characters to travel there, slay a gloomstalker, and collect its shadow essence in a specially made flask. The sages warn the characters not to linger near Ghor Veles, lest they fall prey to the terrible creatures lurking in its haunted fortress.

\subparagraph{Of Vaults and Vengeance (High Level)}
One or more of the characters are kidnapped by minions of \textbf{Oracs the Enduring} (chaotic evil, male, ancient black dracolich) and are carried away to the vaults below Ghor Veles for indoctrination and necromantic experimentation. With the help of a \textbf{core spawn seer} from the Miskath Pit, the dracolich and his minions are weaponizing their mortal enemies to gruesome effect. As one half of the party struggles to survive captivity, the rest of the characters must plot their harrowing rescue.

\subsection{Incanter's Rest}
This Calamity-era battlefield is now a massive graveyard haunted by the undead spirits of unnameable horrors and their cruel, unending vengeance. Named after the unknown mage who brought this region's struggle to its devastating conclusion, Incanter's Rest has few landmarks other than the gigantic skeletal remains of three monstrous titans, which reverently tower above the area's dusty steppes. Howling winds perpetually sweep the Rest with ghostly ardor, while ramshackle monuments of rusted weapons and armor sporadically mark the graves of the fallen.

Some historians speculate that Incanter's Rest is actually the amalgamated site of multiple conflicts, while others claim that it was the most massive battlefield Wildemount has ever seen. No matter the cause, the grave-strewn fields of Incanter's Rest are seeded with powerful necromantic energies that many unscrupulous mages seek to exploit.

\subsection{Miskath Pit}
This geographical anomaly is an arcane sinkhole of unknown depth and origin. Its oblong opening measures roughly 120 feet in diameter and plunges endlessly into the planet's surface. Few survivors have returned from the pit's staggering subterranean expanse, and those who have appear maddened and deranged with fear. Whatever precipitated the advent of the pit is ultimately inconsequential compared to the sheer devastation it causes in the surrounding Blightshore wilds.

Despite its astounding depth, the Miskath Pit's twisting tunnels are far from lightless. A melange of psychedelic sights and sounds are found throughout its bizarre subterranean landscapes, where strange, sporadic lights can be just as terrifying as the darkness. The cloying presence of something ancient and evil is unmistakable to anyone who seeks firsthand to understand the pit's harrowing phenomena. This area is an appalling fusion of the Underdark and the alien quasi-geometries of the Far Realm.

\subsubsection{Spawn of the Core}
Travelers have recently reported sightings of a new variety of horrific aberrations in the vicinity of the Miskath Pit. Known as core spawn, these creatures embody the chaotic nature of Blightshore itself with every tenebrous fiber of their eldritch being. Hordes of lesser core spawn are presided over by core spawn seers—corrupted magic users whose loathsome agendas are one of Blightshore's greatest (and most troubling) mysteries. See \textit{chapter 7} for more information about core spawn.

\subsubsection{Salt of the Arcane Earth}
Despite its unique perils, the Miskath Pit is a darkly alluring place that seems to beckon travelers with a proclivity for eldritch magic or a taste for trouble. The bizarre arcane phenomena that plague the Miskath Pit are also responsible for an array of enchanted minerals and precious elements. Profiteers from Rotthold and beyond risk madness, and even death, in an attempt to mine the pit's supernatural resources.

\subsection{New Haxon}
\begin{description}
\item[Population.]
1,250 (70\% humans, 20\% dragonborn, 10\% other races)

\item[Government.]
The Cerberus Assembly works in tandem with allied mercenaries to maintain a stratocracy of lawful military power.

\item[Defense.]
Satyana Vorvexis, an adult copper dragon, commands a garrison of over two hundred active duty soldiers, and each citizen of New Haxon is enlisted for reserve military duty.

\item[Commerce.]
Common goods and services are available via the city's vast network of enterprising military merchants. Meanwhile, the Cerberus Assembly strives to subvert New Haxon's thriving and profitable black market for rare arcane antiquities.

\item[Organizations.]
Although the Cerberus Assembly's sages hold power here, various mercenary guilds from remote corners of Wildemount have also made New Haxon their base of operations. The city has become an unlikely haven for esoteric cults of the Betrayer Gods and other dark entities.
\end{description}
\noindent
Whatever disaster befell Old Haxon is a tale lost to time and the Calamity. In its place, a fair-weather confederation of humans and dragonborn have erected this lone bastion. The stalwart outpost of New Haxon is garrisoned by agents of the Cerberus Assembly, along with a contingent of allied mercenaries from nearby Draconia. An amalgam of crumbling ruins and modern battlements, New Haxon serves as the Dwendalian Empire's only trusted waypoint for "civilized" exploration throughout Blightshore's savage badlands.

The adult copper dragon Satyana Vorvexis commands New Haxon's threadbare legion of soldiers with a healthy measure of cunning and inventiveness. In addition to her military prowess, Satyana's draconic predilections and experience in the field have earned her a reputation as one of Wynandir's foremost curators and collectors of arcane relics. Some members of the assembly (and their imperial patrons) question the chaotic dragon's true intentions, but the people of New Haxon would be hard pressed to find a more competent leader to navigate the region's tumultuous history and enduring threats.

\subsubsection{Lost Vaults}
Among its other mysteries, Old Haxon is rumored to be the site of one of the lost vaults of Vecna, where the lich hid countless magic items of awesome power and ill intent. Whether or not this arcane storehouse is merely the stuff of legend remains to be seen, but its presence would certainly explain the Cerberus Assembly's keen interest in Haxon's wayward history.

\subsubsection{New Haxon Adventures}
Characters who adventure in New Haxon are likely conscripts, volunteers, or opponents of the local stratocracy's oppressive military force.

\subparagraph{Right Side of Wrong (Low Level)}
As dutiful servants of the Cerberus Assembly, the characters are tasked with ridding a local safe house of cult activity. But the operation doesn't go according to plan: one of the cultists is the teenage son of a prominent assembly official, and the characters must decide whether they will shine a light on his indiscretion or bury it at the potential cost of becoming accomplices themselves.

\subparagraph{Plunder the Crypt (Mid Level)}
A high elf mastermind, Drahir Baiegn, hires the characters to masquerade as Cerberus Assembly agents and plunder a notorious vault below New Haxon: the Amethyst Crypt. When they infiltrate the vault, however, the characters discover they've been duped by Baiegn, who sacrifices his disposable agents to feed an abhorrent arcane construct deep within the inescapable tomb.

\subparagraph{Eater of Worlds (High Level)}
The characters are summoned to a meeting with Satyana Vorvexis. One of the dragon's spies managed to infiltrate a cult operating within the city and determined that its leaders are preparing a ritual to summon Shothotugg, the Eater of Worlds. Shothotugg's arrival would almost certainly place New Haxon in jeopardy. The cult controls a crumbling dungeon complex below the ruins of Old Haxon. The cult's leaders are a pair of vampires aided by a small conclave of cloakers, but most of the cult members are nihilistic humanoids and \textbf{vampire spawn}.

\subsection{Olc Skerry}
Legends say the people of Olc Skerry were either slaughtered by outlanders or claimed by a virulent contagion. Learned sages, however, suggest that the inhabitants of this island hamlet abandoned it long ago and journeyed into the mainlands of the Far Hharom peninsula. The nature of this exodus remains a mystery, but relics found scattered among Olc Skerry's crumbling wooden hovels and tenebrous seaside caves suggest the people of the village revered an unnamed pantheon of forgotten entities from the Far Realm. Modern explorers who return from expeditions to Olc Skerry tell stories of pernicious phantoms who poison the mind with madness, and blood-hungry wolves that lurk among the island's dark and craggy woodlands.

\subsubsection{Olc Skerry Adventures}
Olc Skerry is a spooky, abandoned settlement sure to raise the hackles of even the most jaded explorer.

\subparagraph{Ol' Bucket Helm (Low Level)}
Flocks of black gulls haunt the lonely north shores of the island, congregating around a seaside cave that attracts ordinary crabs and a few giant crabs as well. Just outside the cave mouth, someone has planted a grim scarecrow with a rusted bucket helm for a head, and anyone who tampers with the helm causes the black gulls to form into three hostile swarms (use the \textbf{swarm of ravens} stat block in appendix A of the \textit{Monster Manual}). The gulls also gather into swarms at the command of the cave's reclusive occupant, Peggy Silva, a sea hag with snails in her hair who wears an old fishing net as a shawl. A deranged kuo-toa whip named Plookoolp is the only company Peggy Silva tolerates.

\subparagraph{Black Sap (Mid Level)}
Visitors to Olc Skerry can't help but notice an old stone windmill perched atop a rocky hill just south of town, its four spindly vanes stripped to their wooden bones. Behind the crumbling, desolate structure grows a large death's head willow, from which characters can extract 20 doses of black sap (see the "\textit{Blightshore Contraband}" sidebar earlier in this chapter). The black sap is worth a small fortune in the black markets of Rotthold. However, the windmill is home to an oni that doesn't take kindly to the sap's theft.

\subsection{Rotthold}
\begin{description}
\item[Population.]
8,800 (30\% humans, 15\% dark elves, 15\% goblins, 15\% tieflings, 5\% hollow ones, 20\% other races)

\item[Government.]
Rotthold has no official governing body but is held together by a consortium of local guild leaders, merchants, and veteran adventurers who meet in secret once a week to discuss current affairs, shared interests, and potential threats.

\item[Defense.]
The small but precious amount of sanctuary Rotthold provides from the Blightshore wastes is upheld by a militia of battle-hardened locals and dutifully concerned citizens. Other fortune-seekers often join this militia simply for the opportunities and information the alliance provides.

\item[Commerce.]
Almost anything can be bought or sold in Rotthold, where the merchants are just as unscrupulous as in the surrounding Blightshore wilderness. Rotthold is the center of black market trade in Eastern Wynandir, and perhaps in all of Wildemount.

\item[Organizations.]
In addition to the city's handful of merchants and thieves' guilds, a confederacy of agents representing the mercantile interests of the Myriad and the Revelry operate in public throughout Rotthold, while a covert sect of the Cerberus Assembly carries out their own clandestine business.
\end{description}
\noindent
Some say this lawless harbor city is as dangerous as the arcane wilds of Blightshore that surround it. But whether or not visitors enjoy its chaotic disposition, Rotthold is an unmistakable center of travel and commerce in eastern Wildemount. The decayed stone dwellings and wooden tenements that line Rotthold's ramshackle streets are enveloped by a luminous, rust-colored moss that bathes the city's evenings in a sick crimson glow. Many outlanders view this crime-ridden settlement as little more than a collection of degenerate safe houses and slums, while others see it as a raw jewel—full of potential—that shines bright among the dusky shadows of the poisoned east.

\subsubsection{Government}
Multiple factions operate in and around Rotthold, including covert agents of the Cerberus Assembly and a tenuous coalition between members of the Myriad and the Revelry. While the Plank King's remote forces maintain the southern waterways and oversee the export of Blightshore's esoteric resources, Myriad spies monitor and manipulate the exchange of goods and services inside the city limits, and the fragile alliance between these two is the closest thing this unruly corner of the Miskath Strand has to a governing body.

\subsubsection{Crime}
The open trade of illicit goods and services is the lifeblood of Rotthold's mossy streets. All manner of pleasures are available for those willing enough to pay the price, which can be quite steep. Many of Rotthold's denizens owe the Myriad one favor or another, and the constant exchange of these debts keeps the city full of new faces. Crime is a way of life in Rotthold, where the Dwendalian Empire and Kryn Dynasty have no sway over its opportunistic, unethical inhabitants.

\subsubsection{Geography}
No walls protect Rotthold from the surrounding Blightshore wastes, but eleven watchtowers serve as a line of defense. The city is roughly divided into four quarters: North Rotthold, East Rotthold, Dockside, and Cairn Hill. North Rotthold houses the city's most impoverished denizens, who encounter the Miskath Strand's impinging monstrous wildlife on a regular basis. Most of the city's trade is conducted in East Rotthold, where a bazaar known as the Red Market takes up a majority of the quarter. On the south side of the city, Dockside is the purview of the Revelry and houses Rotthold's most affluent citizens, who no doubt believe they are safest when located as far away from Blightshore's savage interior as possible. On the west end of the city, the quarter known as Cairn Hill rests above a mysterious lattice of pre-Calamity catacombs. Here in Rotthold's graveyard, furtive agents meet in cemeteries to carry out their agendas.

\subsubsection{Rotthold Adventures}
The moss-lit streets of Rotthold are ripe with adventure for characters with a taste for small-city mystery or black market mayhem.

\subparagraph{Rough Times in Rotthold (Low Level)}
A moment of respite goes south in one of North Rotthold's taverns, where the characters—along with the proprietors and other patrons—become targets of an impromptu shakedown by an upstart gang with ties to a prominent thieves' guild. Successfully thwarting the thugs could very well earn the characters the ire (or the respect) of the guild, which operates under the clandestine guidance of a mysterious cabal of enlightened nothics.

\subparagraph{The Cairn Hill Horror (Mid Level)}
A serial killer \textbf{cultist} stalks the streets of western Rotthold from his catacomb lair beneath Cairn Hill, where the viscera of a hundred victims have been offered to summon his otherworldly patron: a \textbf{glabrezu} named Vrath. This demon is a cruel, immortal deceiver that promises its acolytes unrivaled power in the wake of its awful arrival—and only the characters can stop it.

\subparagraph{The Krakenwhacker (High Level)}
Captain Veluna Samarshan of the galley Krakenwhacker has decided to retire in Rotthold. Having heard about the adventurers' exploits, she offers to sell the Krakenwhacker to them for one gold piece, though they'll need to find their own crew. The ship is in excellent condition and has a magical figurehead carved in the likeness of Bahamut that grants the vessel immunity to fire, lightning, and thunder damage. The ship is the envy of several other Revelry captains, and the characters might have to fend off their advances to keep it. Captain Samarshan also warns them about an \textbf{adult bronze dragon} named Walarkus. Corrupted by the magic of Blightshore, this once-noble dragon has become a coastal menace consumed by greed. It would like nothing more than to add the ship's magical figurehead to its undersea trove.

Chapter 5 of the \textit{Dungeon Master's Guide} has rules and statistics for ships, as well as information about hiring a crew and making ship repairs.

\subsection{Tooth of Zehir}
This eastern arm of Blightshore is a mist-shrouded peninsula, the northernmost swath of which is covered by the stone forest of Strathfell (described below).

The lower region of the Tooth is marked by craggy basalt beaches covered in black sand, and the ocean is full of hazardous, fossilized coral reefs. Exotic fey creatures are said to congregate among the Tooth's darker dominions, which are rumored to skirt the Shadowfell via the Border Ethereal.

\subsubsection{Strathfell}
This petrified woodland is the domain of an ancient silver shadow dragon named \textbf{Karkethzerethzerus, the Sable Despoiler} (Karketh for short). Before his malign reign, this region of the Miskath Strand was a verdant haven for a long-lost nation of wood elves and centaurs. For reasons unknown, Karkethzerethzerus plotted to annihilate the woodland nation, killing it off with energies from the Shadowfell. Since doing exactly that, the Sable Despoiler has ruled Strathfell with equal measures of gleeful arrogance and baleful suspicion. The shadow dragon's lair, like his motives, remain hidden. Etheric rifts run throughout Strathfell's gloomy expanse, allowing denizens such as darklings and meenlocks to travel to and from the Shadowfell with unsettling regularity. Most of these interplanar paths are invisible to the naked eye—but minions of Karkethzerethzerus know all too well where the veil grows thin.

The future of Strathfell is not without hope, for good-aligned centaurs continue to patrol the woods as their ancestors did in ages past, battling the evil forces of the shadow dragon at every turn. The centaurs could use some help, however, as they are few in number.

\subsubsection{Chosen of Zehir}
The southern tip of the peninsula abuts the Mordant Isle, a smaller land mass separated from the mainland by the reef-laden waters of the Venom Strait. On this island, an aberrant enclave of winged serpentkin dwell within the ruins of Sariss, the Jade City, where they have erected foul monuments to the Cloaked Serpent himself. These abhorrent snake folk plot a terrible doom for denizens of the Miskath Strand from their horrid lairs. No warm-blooded creature has ever been to Sariss and survived—or so the story goes.

\part{Marquet}

\chapter{Marquet Gazetteer}

Following the winds and waves to the southwest of Tal'Dorei across the Ozmit Sea eventually leads one to the chain of islands called the Hespet Archipelago and to the shores of the continent of \textit{Marquet}. The pebble beaches of the Bay of Gifts greet travelers with the charming sights of untouched nature, and the port city of Shamal invites all to engage in the pleasures of civilization, with beautiful ballrooms and gilded gambling halls.

Beyond the port of Shamal lies the rocky edge of the Aggrad Mountains, whose canyons and caverns house untold numbers of predatory creatures, before giving way to the immense Marquesian Desert. This sandy wasteland is rumored to have once been verdant jungle brought to ruin by the spear of the The Ruiner in the \textit{Calamity}. Some stories tell of a great hero that sacrificed himself to save \textit{Marquet} from utter annihilation, but only scattered fragments of this tale endure.

Centuries after the Calamity, Marquet is known as a continent of deserts and mountains, with few bastions of civilization to be found. But only one-third of the continent is desert.

The Beryl Sea lies south of the continent. The Lapis Channel which leads from Lake Koron and the Gloomed Jungles of Aeshanadoor feeds into it.

\section{Ank'Harel}
Near the center of this intense desert, surrounded by a sprawling network of small villages, is the grand city of \textit{Ank'Harel}, known as the Jewel of Hope. This city, built atop the Drowned City of Cael Morrow, is the center of culture, history, and power for all of \textit{Marquet}. Ruled by the mysteriously undying J'mon Sa Ord, \textit{Ank'Harel} is a magnificent haven of music and culture. It's not without its dangers, but it is a paradise compared to the monster-infested sands that surround it.

The Hands of Ord, the city's guards, also protect nearby villages and townships within the desert, but they can only reach so far. Further south, beyond the reach of \textit{Ank'Harel} and the Hands of Ord, the sands give way to scrubland and marshes tangled with the volatile mountains surrounding the smoking Suuthan Volcano and the dangerous clans that worship it. Westward, beyond the arid desert, lies a vibrant landscape of verdant jungles, craggy mountains, and magical mysteries that the people of \textit{Marquet} have spent eons unraveling. Within the jungles of \textit{Marquet} are dozens of cultures, some nomadic, some seafaring, and some with great cities of their own. One such city, founded by orcish survivors of the The Ruiner wrath during the \textit{Calamity}, is known across Exandria for its mathematically immaculate architecture. This realm welcomes scholars and travelers from across Exandria, who are eager to learn from the minds of some of the greatest architects the world has ever seen.

Ank'Harel translates to "Jewel of Hope" in the northern dialect of the Marquesian language. The city was established as an oasis trading post. A rare source of underground fresh water made this one of the few habitable locations in the vast desert.

As more travelers arrived to take part in building this metropolis, unpaved roads were upgraded to stone streets, and walls were bolstered with towers and parapets. Ank'Harel flourished, becoming a cultural melting pot of people from across the continent and some from farther away.

\subsection{Ank'Harel at a Glance}

\begin{description}
\item[Population.]
293,500 (predominantly humans, orcs, and elves)
\item[Leader.]
J'mon Sa Ord, a lawful neutral \textbf{ancient brass dragon}, helped found the city over 400 years ago and has ruled it ever since. Depending on which supposition one subscribes to, J'mon is either an archfey of the Feywild, a demigod in disguise, or multiple individuals pretending to be a single entity. J'mon is rarely seen in public, and it is extremely difficult to gain an audience with them. Because of the dragon's reclusive nature, an elderly female human named Gemeshega is Ank'Harel's Grandmaven, or acting chancellor.
\item[Laws.]
The laws of Ank'Harel are based on the Code of Ord. Theft, exploitation, and murder are the worst crimes, punishments for which range from life in prison to permanent exile from the city. Laws are enforced by the Hands of Ord, the city guard. Offenses for which culpability is in question or that involve particularly heinous crimes are ruled on by the judges in the Cerulean Palace and, in some cases, by J'mon Sa Ord.
\item[Religion.]
Worship of all Prime Deities is welcomed in Ank'Harel. Worship of the Betrayer Gods, when such activity is discovered, is punished on a case-by-case basis depending on the nature of the reverence. Someone in possession of forbidden iconography is typically sentenced to a few weeks of imprisonment, but someone who conducts rites dedicated to a Betrayer God is treated more severely.
\item[Neighborhoods.]
The city is divided into eight districts: the Alluvium District, the Circlet Walk, the Guided District, the Ridge, the River District, the Sand-Herald District, the Sigil District, and the Suncut Bazaar.
\item[Geography.]
Ank'Harel lies in the center of a vast desert, nestled among clay-rich mountains. Three of these peaks have been leveled off at different elevations to form the Ridge, the Guided District, and the Cerulean Palace. Canals filled with fresh water circulate through the River District, branching outward from the nexus of the oasis.
\end{description}

\subsection{Lodgings}

\begin{DndTable}[header=Ank'Harel Lodgings]{Xll}

Lodging & Type & Lifestyle\\
Back alley & Side Street & Wretched\\
Ajir's Whetstone (in the Alluvium District) & Hostel & Squalid\\
Golden Chip (in the Suncut Bazaar) & Casino & Poor\\
Harried Mongoose (in the River District) & Tavern & Modest\\
Step Aside (in the Circlet Walk) & Tavern/Inn & Comfortable\\
Lap of the Gods (in the Sand-Herald District) & Resort & Wealthy\\
\end{DndTable}

\subsection{Major Factions}
Numerous political, academic, and criminal factions are active in Ank'Harel. Driven by ambitions both apparent and unseen, these organizations and their rivalries shape the city's social and political landscape, and they offer opportunities for adventurers to rub shoulders with powerful individuals and thereby enhance their reputations.

Three of the groups have a strong interest in the Apotheon and what he represents: the Allegiance of Allsight, the Consortium of the Vermilion Dream, and the Library of the Cobalt Soul. Each of these factions has a set of missions the adventurers can undertake, which are detailed later in the chapter.

\subsubsection{Allegiance of Allsight}
\textit{"To peer into the past is to illuminate the future."}


\begin{description}
\item[Headquarters.]
The Crystal Chateau, a university in the Sigil District, is the Allegiance's headquarters.
\item[Leader.]
The Allegiance is led by two headmasters: the dour \textbf{James Cryon}, a lawful neutral elf, who is also the head of the Crystal Chateau's program for the arcane arts; and the kindly \textbf{Gryz Alakritos}, a neutral good \textbf{goblin}, whose infectious enthusiasm makes him the more popular of the pair. Both use the \textbf{scholarly mastermind} stat block (see appendix A).
\item[Allies.]
The Allegiance and the Library of the Cobalt Soul are staunch allies.
\item[Opponents.]
The Allegiance opposes the Consortium of the Vermilion Dream and the Sentinels of Memory (see "Minor Factions" later in the chapter).
\end{description}

\noindent The Allegiance of Allsight is an academic collective known across Marquet. Its original members came to Ank'Harel from across the continent to rebuild the robust academic tradition that defined the elves and orcs of ancient Cael Morrow.

The Allegiance of Allsight is embodied in two of Ank'Harel's academic institutions: Teres Schoolhouse, the city's largest center of education, and the Crystal Chateau, an elite school of arcane principles and the center of the Allegiance's operations. Students from anywhere in Ank'Harel are welcome at Teres Schoolhouse, but the Crystal Chateau's programs are extremely competitive, and all pupils there are required to join the Allegiance of Allsight.

In recent years, the faction's leadership has been focused on uncovering the secrets of Cael Morrow, the sunken ruins beneath the city. For decades previously, the headmasters of the Allegiance petitioned J'mon Sa Ord for permission to unseal the long-lost ruins, and finally, the city's mysterious leader relented. The opening of an entrance to the ruins infuriated the Sentinels of Memory, a minor faction whose members fanatically believe the sunken city is best left forgotten, lest some unknown terror left over from the Calamity be disturbed.

\subparagraph{How to Join}
Joining the Allegiance of Allsight requires one to be referred by an active member, then be reviewed and audited by the Allegiance's council of leadership prior to the issuance of a judgment on the candidate.

\subsubsection{Consortium of the Vermilion Dream}
\textit{"Seek the stories others fear to know."}


\begin{description}
\item[Headquarters.]
Consortium members meet in First Eclipse, a tavern in the Suncut Bazaar, and hold private conversations in a secret storeroom there.
\item[Leader.]
The Consortium is led by a council of five masters who all use the \textbf{occult silvertongue} stat block (see appendix A). They are \textbf{Aradrine the Owl}, a lawful neutral goliath (see the "Goliaths of Exandria" sidebar later in the chapter); Dendarron the Sun Bear, a chaotic neutral halfling; Larthul the Wolf, a chaotic evil human; Khelkur the Gull, a neutral evil dwarf; and Vrill the Moth, a lawful evil elf.
\item[Allies.]
The Consortium maintains a tenuous alliance with the Sentinels of Memory (see "Minor Factions" later in the chapter).
\item[Opponents.]
The Consortium opposes the Allegiance of Allsight and the Library of the Cobalt Soul.
\end{description}

\noindent This group of occult mystics has assembled a membership of like-minded souls from across Exandria. The Consortium of the Vermilion Dream focuses on studying the magic described in folklore that has been dismissed as fanciful by more "proper" institutions. The group takes its name from its foremost object of fascination: Ruidus, the ruddy moon of Exandria that is fabled to bestow ill fortune and suffering upon those who are born, enter into contracts, or begin new ventures under its baleful light.

The Consortium was founded by five influential individuals, all renowned in occult circles from across Exandria, who sought to gain wealth and prestige from their studies. The Consortium is still young, and thus all five of its founders still live. Older academic institutions such as the Allegiance of Allsight and the Library of the Cobalt Soul express concern and sometimes disdain toward the upstart Consortium and its aggressive, profit-driven attitude toward scholarship.

Due to the expensive nature of its research, the Consortium is in constant need of funding. It leans heavily on its members for such support, selling the services of its magically skilled devotees to government officials, wealthy diplomats, and other persons of means. A large portion of its income goes to support the Sentinels of Memory, the Consortium's tentative ally against the Allegiance of Allsight.

\subparagraph{How to Join}
The Consortium gladly welcomes any enthusiast of the strange and occult to its ranks—as long as the faction's leaders can be reasonably certain that the individual has no ulterior motives that would undermine their efforts.

\subsubsection{Library of the Cobalt Soul}
\textit{"Reason. Knowledge. Truth."}


\begin{description}
\item[Headquarters.]
The Marquesian branch of the Library of the Cobalt Soul is based at the Temple of the Mentor—a temple of Ioun the Knowing Mentor—in the Guided District.
\item[Leader.]
At 29 years old, \textbf{Jamil A'alithiya}, a chaotic good human, is the youngest \textbf{monastic high curator} (see appendix A) in the history of the Cobalt Soul. He is often underestimated by his contemporaries and uses that fact to his advantage.
\item[Allies.]
The Cobalt Soul maintains alliances with the Allegiance of Allsight and the Hands of Ord (see "Minor Factions" below).
\item[Opponents.]
The Cobalt Soul opposes the Consortium of the Vermilion Dream.
\end{description}


\noindent The Library of the Cobalt Soul was founded in the heart of Wildemount centuries ago, before any of the nations that currently wage war over that land were born. Its archivists have stood strong against conflict, propaganda, and upheaval during all this time in their dogged hunt for the truth. Their mission has spread across Exandria, with advocates in various locales defying local governments in the pursuit of truth.

The Cobalt Soul founded the Temple of the Mentor in Ank'Harel nearly 150 years ago. The archivists here focus on rooting out evil in the city's civilian factions—principally the Consortium of the Vermilion Dream, whose members reportedly wield magic weapons made of some strange, red mineral the likes of which no Cobalt Soul archivist has ever seen.

The faction is wary of the Allegiance of Allsight's fascination with Cael Morrow, but as fellow academic institutions, the Allegiance of Allsight and the Cobalt Soul work together most of the time, sharing resources from their libraries.

\subparagraph{How to Join}
The Library of the Cobalt Soul accepts for consideration any volunteers who express interest in gathering and preserving knowledge. These volunteers are put to work as librarians, performing menial tasks such as organizing records and transcribing notes from expeditions—drudgery designed to weed out those who want to join the organization for the sake of fame and adventure. Typically, this duty lasts about six months, after which a librarian is formally inducted.

\subsection{Minor Factions}
Other, smaller factions are important to life in Ank'Harel.

\subsubsection{Hands of Ord}
\textit{"By the code of Ord, we stand vigilant."}


\begin{description}
\item[Headquarters.]
The city guard is based in Ord Bastion in the Sand-Herald District.
\item[Leader.]
Ironhand Sem, a lawful good \textbf{minotaur}, commands the Hands of Ord and oversees all military operations in Ank'Harel.
\item[Allies.]
The Hands of Ord and the Library of the Cobalt Soul are staunch allies.
\item[Opponents.]
The Hands of Ord aim to root out and destroy the Veil (see below).
\end{description}

\noindent Established four hundred years ago by J'mon Sa Ord, the Hands of Ord are a peacekeeping order of desert warriors who have long watched over Ank'Harel and kept the city safe and stable.

The order is governed by the code of Ord, which is the backbone of the city's legal system. Members of the Hands dedicate their lives to the city and the code, and many come from a long, respected bloodline of previous Hands. Well-appointed barracks in the Sand-Herald District provide members of the Hands with a comfortable home; places of worship; and sources of food, armor, and weapons.

Despite their best efforts, the Hands of Ord have failed to thwart the Veil, the crime syndicate that pervades Ank'Harel society.

\subsubsection{Scarbearers}
\textit{"We do not break."}


\begin{description}
\item[Headquarters.]
The Bowl of Judgment, an arena in the Sand-Herald District used for combat tournaments, is the Scarbearers' headquarters.
\item[Leader.]
Once a mercenary in his own right, Quartermaster Croog Lynn, a chaotic neutral, goliath \textbf{gladiator}, has now taken to running the guild in his later years (see the accompanying "Goliaths of Exandria" sidebar).
\item[Allies.]
None.
\item[Opponents.]
None.
\end{description}

\noindent The Scarbearers are the most famous mercenary company in Ank'Harel. Named for the scars their founders displayed to J'mon Sa Ord as proof of their service to the realm, the Scarbearers occupy an honorable position in Marquesian society. They often take high-profile jobs, such as serving as bodyguards for diplomats and protecting caravans and expeditions into the desert. The members promote themselves as a guild that respects honor and guidelines, but they often come across as brash, thuggish fighters who are out for a paycheck.

In addition to their mercenary exploits, the Scarbearers also supervise the tournaments that take place in the Bowl of Judgment.

\subsubsection{Sentinels of Memory}
\textit{"Let the past die in peace."}


\begin{description}
\item[Headquarters.]
The Tower of Memory in the Guided District is the meeting place for the Sentinels of Memory.
\item[Leader.]
Watcher Trast, a neutral, elf \textbf{priest} and a former cleric of Ioun, founded the Sentinels of Memory and maintains command over the faction to this day.
\item[Allies.]
The Sentinels are allied with the Consortium of the Vermilion Dream.
\item[Opponents.]
The Sentinels oppose the Allegiance of Allsight.
\end{description}

\noindent The Sentinels of Memory are a group of fanatics who believe the sunken city beneath Ank'Harel isn't some grand ancient wonder—it's a prison that must be kept sealed at all costs. Citing the lack of recorded history about Cael Morrow as proof that powerful forces wanted to expunge the city from memory, the Sentinels of Memory are willing to do whatever it takes to ensure that no one speaks of that ill-fated city ever again.

Since they oppose the excavation of Cael Morrow, the Sentinels of Memory are constantly at odds with the Allegiance of Allsight. Though their efforts to date have mostly been political, some of their members have been connected to violent acts intended to thwart the excavation efforts.

Recently, the Consortium of the Vermilion Dream has been supplying the Sentinels of Memory with funds and agents to support the Sentinels' conflict against the Allegiance of Allsight. The Sentinels are grateful for the aid but treat the Consortium with caution, suspecting the faction has ulterior motives.

\subsubsection{The Veil}
\textit{"Welcome to the other side."}


\begin{description}
\item[Headquarters.]
None.
\item[Leader.]
Twenty-year-old Ilena Hapayhari, a chaotic evil, human \textbf{assassin} and ruthless pirate from the western coast of Marquet, recently disposed of the Veil's old leadership. She now rules the Veil from the shadows, though only high-ranking members are privy to this knowledge.
\item[Allies.]
None.
\item[Opponents.]
The Veil opposes the Hands of Ord.
\end{description}

\noindent A secret network of spies, thieves, and assassins, the Veil is the largest crime syndicate in Ank'Harel. It is the chaotic counterpoint to the order of the Hands of Ord.

Not much is publicly known about the Veil; the identities and personal details of its members are kept secret even among the group, and those who let such secrets slip are often found murdered before they can be apprehended. The Veil has no permanent headquarters in the city, constantly changing bases to cover its trail.

Although the Veil doesn't play a role in this adventure, you can use this faction as an adversary in your own adventures set in Ank'Harel.

\subsection{Alluvium District}
The Alluvium District is the northernmost neighborhood of Ank'Harel. Named for the rich clay mined in the cliffs that border the district, this area is a mixture of dense residential sections, sites of clay strip mining, and spacious plazas used for oral performances and artisan markets.

\subsubsection{Alluvium Gardens}
\textit{Arboretum}

\noindent The arboretum known as the Alluvium Gardens is famous for the meticulous, geometric layout of its hedges and planters, mimicking the tile art of the district's ceramic masters. After a long day of hauling clay, miners often retreat to the shade of the gardens to unwind and chat with friends. At the center of the gardens is a grove of fig, apricot, and pomegranate trees, and visitors to the gardens are welcome to help themselves to any ripe fruit they can reach. The gardens also contain kurrak, a hard-to-cultivate fruit native to the jungles of Marquet that has a tough, spiny shell and soft, sweet flesh like that of a mango.

\subsubsection{Opalite Forum}
\textit{Entertainment Square}

\noindent Eight pillars of iridescent glass, created from discarded sand and debris from the clay mines, ring the outside of this wide stone courtyard in the heart of the Alluvium District. The Opalite Forum is a central gathering place for district residents, often used to host festivals and holy celebrations.

The forum is also a place where storytellers from across Exandria tell their tales. During the day, children gather around the elderly chroniclers, who relate the legends of Exandria's founding and stories of the Calamity in the oral tradition of their forebears. At night, the forum is overtaken by crowds of burgeoning writers, reading their latest works aloud and, on occasion, engaging in friendly duels of poeticism and wit.

\subsubsection{Sa'irah Mines}
\textit{Clay Strip Mines}

\noindent The angular ridges of the Sa'irah strip mines loom over the rest of the Alluvium District, and the cacophony of hammering and drilling fills the nearby streets almost constantly during the workday. Most of the clay harvested here is used for the construction and repair of buildings. Some of it is worked by artisans who have perfected the art of blending and firing the clay, creating vibrantly colored mosaic tiles and other ceramic works.

The quarry is run by a gruff, middle-aged miner named Zala Keencutter, a neutral, tiefling \textbf{thug}. Though he puts up an abrasive facade, Zala has a shy and empathetic soul, and he wants dearly to work up the courage to read his prose writings at the Opalite Forum some evening.

\subsection{Cerulean Palace}
The purple and blue spires of the Cerulean Palace reflect the sun's light off their edges like a sun-dappled waterfall. Most citizens are familiar only with the palace's ornate exterior, covered with glittering designs of lapis lazuli and turquoise with brass accents. The inside of the palace, where only those with official business are typically allowed, contains chambers used for deliberation, debate, and judgment regarding legal matters. Beneath the palace lies a labyrinthine dungeon called the Oblag, also known as the Scarlet Oubliette.

At the center of the palace complex is the Chamber of Judgment. Intricate tilework covers the outer surface of the chamber. The inside is dimly illuminated by sunlight that shines through the translucent panes of bluish crystal that make up the roof. On a raised stone dais sits an intricately carved bronze throne studded with sapphires—the Seat of J'mon Sa Ord, said to be a duplicate of the one the leader uses in their personal quarters at the top of the palace's highest tower. In the Chamber of Judgment, perpetrators of the most terrible crimes are brought to justice. J'mon Sa Ord descends to occupy the seat while presiding over these hearings and then personally delivers the verdict.

\subsection{The Circlet Walk}
The Circlet Walk is primarily a residential district, though taverns and other places of commerce can be found here as well. The district is so packed with buildings that it's customary for people to travel through the district by hopping from rooftop to rooftop rather than navigating the narrow roads.

\subsubsection{Lyrean Linen}
\textit{Front for the Black Market}

\noindent At first glance, Lyrean Linen looks like any other cloth-dyeing establishment, with long sheets of fabric hanging off every rafter and vats of colored liquid arrayed across the front of the shop. Asking for cloth dyed "whitestone green" gives a customer access to the store's other wares: a dangerous narcotic called suude and other forbidden concoctions sold on the black market.

The shop is run by Bas and Kovra Lyrean, twin sisters who are neutral evil, half-elf bandits. Originally alchemists by trade, they use their knowledge to brew magic potions deemed by the Hands of Ord to be too dangerous for public consumption.

\subsubsection{Step Aside}
\textit{Inn and Tavern}

\noindent his cozy establishment is a local favorite, known for its delicious food and a cheery atmosphere that invites folks to "step aside" from their daily troubles and rest a while. Desert flowers flourish in planters that decorate the tavern's windowsills and bar, and the upstairs rooms are redolent with the scents of garlic rice and fried meats wafting up from the kitchen. On the most sweltering days, the tavern offers its signature dessert: haluh, a dish from the southwestern Marquesian coast consisting of shaved ice and sweet milk topped with jackfruit or kurrak, coconut strips, and sweet purple yam.

The place is run by two former mercenaries: Irawan, a lawful good, elf \textbf{veteran}, and his husband Calinao, a neutral good, elf \textbf{druid}. Unbeknownst to many, Irawan and Calinao were among the original Scarbearers who came to Ank'Harel during the city's nascent years and defiantly showed to J'mon Sa Ord the scars they earned during their defense of the realm. The couple has long since retired from mercenary life, but they serve the people of the city in other ways. If the phrase "Pardon my steps as I pass through" is whispered to either Irawan or Calinao, they'll lead the individual to a safe house beneath the tavern, where the beneficiary is welcome to stay until other arrangements can be made.

\subsection{Guided District}
The Guided District contains several houses of worship, with each of the Prime Deities represented. The district's architecture is more ornamental than in other areas of the city; many temples are adorned with intricate mosaics and crowned by minarets with gilded spires.

The few vagabonds to be found in Ank'Harel frequent the alleys in this district, where many temples send representatives to feed and clothe the downtrodden. The Hands of Ord also maintain a notable presence here.

\subsubsection{Blessing Well}
\textit{Holy Site}

\noindent At the heart of the Guided District is a concave stone courtyard known as the Blessing Well. A line of alabaster tiles spirals down to the center of the plaza, ending at the well itself, whose deep blue bricks are speckled with silver in a representation of the night sky.

The Blessing Well is not dedicated to a deity and instead is a public space for anyone in Ank'Harel who seeks spiritual guidance. On many occasions, an individual has been lowered into the well, away from the bright lights and bustle of the city, to calm their thoughts and commune with the divine; they emerge hours later, eyes gleaming with renewed purpose as they savor the divine vision they experienced while in the well's depths.

The water of the well is separate from that in the city canals and is used only in holy ceremonies. Drinking the well's water bestows a terrible divine curse upon the drinker that even the most potent mortal magic can't undo.

\subsubsection{Crossings of Eventide}
\textit{Necropolis}

\noindent Over two hundred feet high, the Crossings of Eventide are a series of tombs carved into the cliffs that border the Guided District. Reliefs depicting the most important deeds of the tombs' occupants adorn the spaces between the stone slabs that seal the graves. At sunset, crystal braziers that hang throughout the necropolis catch the fading light, painting the nearby cliffside in vibrant hues of pink, orange, and purple.

At the base of the cliff rests a small temple dedicated to the Matron of Death. In Ank'Harel (and in much of Marquet at large), the god of death and fate is perceived differently from how she is regarded in the lands of Wildemount, Tal'Dorei, and Issylra. Bereft of her moniker "the Raven Queen," here the god is associated with twilight, inevitability, and the passage of time, and she presides over the transition between life and death. Some Marquesian sects refer to her as the "Duskmaven," portraying her as a keen-eyed vulture with a golden funerary mask and feathers whose colors range from sunset hues to midnight black.

\subsubsection{Temple of the Mentor}
\textit{Temple to Ioun and Headquarters of the Cobalt Soul}

\noindent A ceramic mosaic of three open eyes framed in cobalt blue crowns the entrance to the Temple of the Mentor. Dedicated to Ioun, the temple is more a library than a hall of worship, and it doubles as the base of operations for the Marquesian branch of the Library of the Cobalt Soul. Shelves line the walls of the temple's sanctuary, all filled near to capacity with meticulously catalogued tomes, scrolls, and periodicals that are available for public inspection.

Staircases lead up from the sanctuary to the residents' apartments and the scriptorium, and up farther to the quarters of the Cobalt Soul's higher-ranking members. Stairs lead down from the sanctuary to a locked door opened with a key held only by select individuals and the Cobalt Soul's high curator. Behind this door is the Cobalt Soul's vast subterranean archive, containing all sorts of esoteric knowledge that is purposely kept away from the public. The entire expanse is lit by \textit{continual flame} spells.

At the high curator's behest, a sardonic monk named Iwo Zalarre, a chaotic good, half-orc \textbf{monastic operative} (see appendix A), oversees all visitors to the public archives and registers all requests for access to the private archive. Though just as dedicated to Ioun's service as any other member of the Cobalt Soul, Iwo is eager to get out of the temple and take part in a challenging mission.

\subsubsection{Tower of Memory}
\textit{Headquarters of the Sentinels of Memory}

\noindent Situated on the far western side of the Guided District, the octagonal Tower of Memory juts out of the ground at a slightly crooked angle. The building was slated for demolition when the founder of the Sentinels of Memory, Watcher Trast, purchased it to use as a home for his small but zealous group of followers. The Sentinels have since spent little time repairing the structure, so cracks still show between the ivory bricks, and the wooden floorboards creak with every step.

Watcher Byron, a lawful evil, gnome \textbf{veteran}, oversees most of the sentinels' operations for Watcher Trast and rarely leaves the Tower of Memory. Gentle of voice but volatile of temper, Byron has grown increasingly impatient with his faction's political strategy and has begun to plot more aggressive ways of hampering the Allegiance of Allsight's excavation of Cael Morrow. Though he would never admit it, Byron secretly longs to see the wonders of Cael Morrow for himself.

\subsection{The Ridge}
Sitting atop one of the city's three plateaus, the Ridge is Ank'Harel's industrial district. Its streets are crowded with smithies, foundries, shipyards, and warehouses.

\subsubsection{Alsfarin Union Shipyard}
\textit{Skyship Conglomerate}

\noindent The foremost producer of skyships in Exandria, the Alsfarin Union, has its base of operations in Ank'Harel. Its shipyard, thundering with the peal of hammers and glowing with arcane energy when work is under way, is off limits to anyone except employees. The Alsfarin Union protects its supply of brumestone—the enchanted crystal that magically holds skyships aloft—with an iron fist.

Since all skyships in operation across Exandria must come to the shipyard once a year for maintenance and repair, the area near the shipyard has become a hodgepodge of cultures from around the world. Curious folks can purchase iconic headwear from Whitestone in Tal'Dorei at the Lord Percy Haberdashery (a name not endorsed by Lord Percival de Rolo III of Whitestone), then patronize a cart selling pastries from Port Damali in Wildemount and indulge in a mug of ale brewed in Vasselheim in distant Issylra, all in the same city block.

\subsubsection{Sunfire Forge}
\textit{Smithy}

\noindent Built partially into the side of the mountain that borders the north side of the district, the Sunfire Forge is the premier blacksmith and weapons manufacturer in Ank'Harel. Its name comes from the forge's unique power source: liquid sunlight, harvested from the iridescent crystals lining its roof. The consistent heat of this concentrated form of energy allows for the creation of fine metalwork that incorporates fragile components such as arcane minerals and gem dust. Two glowing streams of sunfire cascade down the sides of the forge's pitched roof before coming to rest in stone reservoirs, where the liquid is stored before being funneled into the forge.

The Sunfire Forge is under the control of Rohaya Tak, a neutral good, half-elf \textbf{commoner}, who inherited the business from her mother two decades ago. An innovative metalworker, Rohaya combines her smithcraft with shrewd business intuition. In addition to crafting blades and tools, Rohaya uses the forge to produce cannons for use on sailing ships and skyships, modeled after the bronze cannons of her family's homeland in southern Marquet.

\subsection{River District}
Its atmosphere permeated by the smell of wet earth, the River District is where most of the city's farming is done. Strips of irrigated cropland are intermingled with newer hydroponic farms, and in place of roads, canals connect the subdivisions of the district. Most folks traverse the canals of the River District in skiffs and lightweight canoes.

\subsubsection{Life Dome}
\textit{Nexus of the Canals}

\noindent This glass dome crowns the source of water in Ank'Harel: the oasis around which the city was built. All the canals in the River District originate from the Life Dome, and the structure is a center of celebration and a monument to the city's history. At night, the inside of the dome is magically lit by drifting orbs of blue and green luminescence, painting the nearby area in a display of colored light.

The Life Dome also serves an important agricultural purpose. Because the water that fills Ank'Harel's oasis flows not from a natural source but from underground caverns created by the acts of the gods, the water level frequently fluctuates. The water can surge forth or slow to a trickle, threatening to bring about floods or dry spells, respectively. To counteract this problem, the Life Dome instantly transmutes any water that exceeds the accepted level into vapor and stores it in the seams of the dome. When the water level drops, the stored vapor is transmuted back into liquid and fed into the irrigation canals.

The Allegiance of Allsight scholar Carliale Kroogan, a neutral good, gnome \textbf{scholarly excavator} (see appendix A), is currently the chief supervisor for the Life Dome. Lonely and a bit of a chatterbox, Carliale is always ready to chat about arcane theory with any individual who comes his way.

\subsubsection{Old Man Kruuk's}
\textit{Bakery and Fence for Stolen Goods}

\noindent Old Man Kruuk, a neutral, half-orc \textbf{bandit captain}, is a mediocre baker; his pitas are tough and chewy, and the caramel of his baklava is bitter and burned. Kruuk's talent lies in his true vocation as a fence for stolen goods. After an individual gains his trust, he drops his bumbling baker facade and reveals himself to be an excellent businessperson, able to sweet-talk even the most uncompromising thieves into a favorable deal.

The scope of Old Man Kruuk's operation far exceeds what anyone might suspect—or even what he should be capable of as a single individual. Some rumormongers suggest that Kruuk is a high-ranking member of the Veil, but no one has verified this supposition.

\subsubsection{Steam Gardens}
\textit{Arboretum and Bathhouse}

The most luxurious of Ank'Harel's arboretums, the Steam Gardens are popular because they include a bathhouse and spa. Water from the canals is fed into the gardens' magically fueled hot springs, which bubble as they release massaging jets of water.

An hour's visit at the most exclusive spa costs 5 gp. Characters who spend 8 hours eating, luxuriating, and relaxing in the spa awaken the next morning with 1d10 \textit{temporary hit points} that last for the next 24 hours, in addition to the usual benefits of finishing a long rest.

\subsection{Sand-Herald District}
The center of military activity and law enforcement in Ank'Harel, the Sand-Herald District houses the headquarters of the Hands of Ord. The organization's stronghold is a major feature of the district, dwarfed only by the Cerulean Palace.

The residences in the Sand-Herald District are occupied mostly by wealthy folk: powerful merchants, decorated mercenaries who have earned riches by winning tournaments in the Bowl of Judgment, and local celebrities of song and stage. The district is also the site of Ank'Harel's skyport.

\subsubsection{Bowl of Judgment}
\textit{Combat Tournament Arena}

\noindent On the north side of the Sand-Herald District, an immense circular depression has been carved into the ground. The Bowl of Judgment is where warriors are made into legends. Sporting tournaments are held here on a regular basis, highlighted by a yearly event known as the Grand Tournament. At each competition, spectators wager on who will be the winner of the upcoming match. Anyone who can pay the entry fee of 20 gp can sign up to compete. Although most bouts end with one participant yielding or falling unconscious, death is always a possible outcome in the Bowl of Judgment.

Sheed Caltor, a chaotic neutral, orc noble, is the suave ringmaster of the Bowl. Sheed is mostly concerned with his image, his income, and his status among the upper crust of Ank'Harel.

\subsubsection{Indala Skyport}
\textit{Skyship Port}

\noindent Skyship travel is undertaken primarily by those who want to avoid the time and danger of an ocean voyage or a long trek through the desert—and who can afford the cost of such transportation. The Indala Skyport dates to the earliest days of skyship travel. Its central terminal is world-famous for its elaborate stained-glass ceiling, which depicts a mosaic of the surface of Exandria with glowing dots moving across it, representing each skyship the Alsfarin Union currently has in flight.

Outbound flights from the Indala Skyport connect to other major skyports in Exandria, including Vasselheim in Issylra, Port Damali in Wildemount, and Emon and Whitestone in Tal'Dorei.

\subsubsection{Ord Bastion}
\textit{Headquarters of the Hands of Ord}

\noindent In contrast to the heavily embellished architecture of the Cerulean Palace, the towering walls of Ord Bastion are built from sturdy mud bricks, undecorated except for simple, angular carvings that adorn the tops of the building's watchtowers.

Ord Bastion serves principally as the barracks and training center for the members of the Hands of Ord, but since membership in the Hands is a life-long commitment, many members have brought their families to live in the bastion, creating a small community bound together by love and loyalty.

\subsection{Sigil District}
The Sigil District is one of the smaller districts in Ank'Harel but also one of the busiest. People from across the city commute here every day to study at the district's various centers of learning. From basic skills and trades to the most esoteric arcane ideals, various intellectual pursuits are represented in the Sigil District.

The Allegiance of Allsight is a pervasive presence in this district; many of the faction's members and affiliates are on the faculties and staffs of these schools. Institutions of learning constructed in Ank'Harel's infancy abut structures of more modern architecture, and buildings representing cultures from across Exandria are a counterpoint to the grand, Marquesian-style domes of the Crystal Chateau. At the heart of the district is the excavation that leads to Cael Morrow, the mysterious sunken city (see "Maw of Cael Morrow" below).

The Streets of the Sigil District map shows a small segment of the district associated with the faction missions later in the chapter. It includes the following locations:


\begin{description}
\item[S1: Crystal Chateau Entrance.]
This elite university is described below.
\item[S2: Teleportation Atrium.]
A permanent teleportation circle stands in the center of this austere, domed building. Most Teleportation Tablet (see appendix B) made by the Allegiance of Allsight are keyed to this circle.
\item[S3: The Frog and Kebab.]
No actual frogs are served at this popular tavern, which is frequented by students of the Crystal Chateau.
\item[S4: Emerald Ibis Relic Hall.]
Watched over by statues made of green glass and shaped like giant ibises, this hall is where the Allegiance of Allsight examines newly discovered items. See "Consortium Mission 3: Elephant in the Room" in the Consortium Story Track.
\item[S5: Scroll's Alcove.]
This apartment building is used mostly by professors at the Crystal Chateau and visiting guest lecturers. See "Cobalt Soul Mission 1: The Cultist of Zehir" in the Cobalt Soul Story Track.
\item[S6: Lawbearer's Librams.]
One of dozens of libraries in the Sigil District, this is the most extensive legal library in the city.
\item[S7: Boughs of the Wild Mother.]
Students can be found studying in the shade of this courtyard's giant sycamore tree during the day. Some say a ghost lurks here at night (see "Consortium Mission 1: A Ghost in Our Midst" in the Consortium Story Track).
\end{description}

\subsubsection{Crystal Chateau}
\textit{Elite University and Headquarters of the Allegiance of Allsight}

\noindent An immense structure of smoky quartz and white marble, the Crystal Chateau is a major landmark in the Sigil District. A small stream winds through the gardens around the building before flowing under an opening into the main hall, where the bubbling sound of its water is said to help inspire the highest of thoughts. A white iron fence closes off the premises to all except members of the Allegiance of Allsight.

The Crystal Chateau serves double duty as a university for the teaching of spellcraft and arcane theory and as the center for the Allegiance's operations. Though not all members of the Allegiance are students or professors at the university, every would-be pupil must first become a member of the faction. Tuition is free, since the Chateau's expenses are subsidized by the city in the expectation that the university's graduates will use their gifts to protect Ank'Harel and enhance its prestige across the world.

Jor Raashid, a lawful neutral, halfling \textbf{scholarly agent} (see appendix A), is the gatekeeper for the Crystal Chateau. Suspicious and condescending, Jor sternly prohibits nonmembers from entering the Chateau.

\subsubsection{Maw of Cael Morrow}
\textit{Excavation Site Entrance}

\noindent What was once a public square has been converted into an excavation site, from where the Allegiance of Allsight conducts its exploration into the ruins of Cael Morrow. The "maw" is a 150-foot-diameter, roped-off hole, admission into which is prohibited to all except for holders of Allegiance of Allsight badges. A 10-foot-wide path spirals down along the edge of the maw, descending hundreds of feet into Cael Morrow (see chapter 5).

\subparagraph{Protection}
A contingent of five lawful neutral Scarbearer veterans patrols the rim of the maw, on the lookout for unauthorized visitors. They are led by a loyal, lawful good, orc \textbf{scholarly excavator} (see appendix A) named Hakzorne. He permits only people who display Allegiance badges to enter the maw; Hakzorne gives only one warning before ordering the mercenaries to attack.

\subsubsection{Omnival Library}
\textit{Public Library}

\noindent Tan stone columns and pointed archways form the circumference of this enormous tower on the northeast side of the Sigil District. Inside, texts from across Exandria fill shelves that line the wall space between the floor and the library's vaulted gold ceiling. Staircases and magical lifts provide access to the highest shelves, and on the floor sits a mixture of weathered wooden desks and colorful cushions.

Any individual is welcome to borrow from the library, though access to the most esoteric and delicate materials requires the approval of an Allegiance of Allsight headmaster. Recent additions to the library include copies of notes from the latest Cael Morrow expeditions.

The Omnival Library is managed by Bookkeeper Khime, a neutral good, orc \textbf{mage}. He's convinced that shadowy forces are constantly trying to undermine Ank'Harel, and this penchant for conspiracy theories makes him incredibly observant. He's an excellent source of information about happenings in the city, despite his wild conclusions.

\subsubsection{Teres Schoolhouse}
\textit{University}

\noindent The larger and less stringent of the city's two universities, Teres Schoolhouse offers a curriculum that features practical applications of magic as well as a wide range of trade skills. Its physical presence, though not as imposing as that of the Crystal Chateau, occupies a large portion of the Sigil District, made up of both older buildings of weathered brick and newer structures built from stone quarried from northern Marquet.

Academic schedules at Teres Schoolhouse are typically divided into two semesters, each five months long. The institution also puts on single-session educational seminars that anyone in the city is welcome to attend. As with the Crystal Chateau, tuition at Teres Schoolhouse is free for all residents of Ank'Harel.

\subsection{Suncut Bazaar}
Crimson banners, hanging flags of all colors, and cloth drapery that provides shade over the crowded streets herald the entrance to the Suncut Bazaar, the center of commerce in Ank'Harel. Awnings stretch above open-air food stalls that offer treats ranging from fresh fruit to sweet semolina cakes to warm meat carved directly off a roasting spit. Most of the densely packed stone buildings are stores that sell artisan goods and shops that buy and sell magic items and nonmagical trinkets. In the center of the district lies the palatial Luck's Run casino.

The Streets of the Suncut Bazaar map shows a small segment of the district that comes into play in the faction missions later in the chapter. It includes the following locations:


\begin{description}
\item[T1: Welcome Arch.]
This vaulted stone archway draped with crimson banners welcomes everyone to the bazaar.
\item[T2: Battlements.]
The walls of the Suncut Bazaar are made of sandstone and patrolled day and night by Hands of Ord guards and veterans.
\item[T3: Central Street.]
This dusty red road that runs throughout the district is so crowded with shoppers, pickpockets, tourists, and locals that it can be hard to see the paving stones.
\item[T4: Market Stalls.]
Market stalls shaded by awnings line the streets and alleys of the bazaar.
\item[T5: Luck's Run.]
This casino is described below.
\item[T6: Bone Garden.]
The bazaar's most infamous pawn shop is described below.
\item[T7: Mystic Pursuits.]
This arcanist's abode is described below.
\item[T8: First Eclipse.]
Both a tavern and the secret headquarters of the Consortium of the Vermilion Dream, this venue is described below. Its roof is where the characters appear if they used one of Aloysia's Teleportation Tablet at the end of chapter 3.
\item[T9: Courtyard.]
During business hours, hungry folk line up in front of the food stalls on the west side of this courtyard while others sit around the edge of the courtyard's fountain and share gossip.
\end{description}

\subsubsection{Bone Garden}
\textit{Purveyor of Oddities and Exceptionalities}

\noindent The clatter of fish-bone wind chimes announces one's entrance into the Bone Garden: a combination pawn shop, antique store, and repository of strange and misplaced items. The shop, which stands near the main entrance of the bazaar, has become a running joke among the residents of Ank'Harel. Fruitless searches for a missing key or a lost earring often end with the lighthearted comment "It'll probably end up among the Bones within the week."

The Bone Garden is run by Rerosha, a chaotic good, human \textbf{acolyte} who's a nonverbal young woman with a brazen attitude toward life. She uses facial expressions, gestures, and a sassy gray parrot named Akil to communicate.

\subsubsection{First Eclipse}
\textit{Tavern and Headquarters of the Consortium of the Vermilion Dream}

\noindent In sharp contrast to the glimmering Luck's Run casino across the street, First Eclipse is notable for its maroon stone walls and ash-black door. Inside, old ale barrels have been turned into tables, and the back wall is covered with painted lids from ales, wine, and liquor casks from across Exandria.

One of these cask lids bears a symbol not associated with any known brewer: a red crescent moon. Servers, if asked about it, are told to explain that this lid represents the "house ale." In truth, knocking three times on the cask lid reveals a secret door that leads to a hidden storeroom. Members of the Consortium of the Vermilion Dream use this room as a base of operations. (The Consortium changes its headquarters frequently to avoid detection.)

The secret room is little more than a cellar with a long table in it, with a few locked cabinets containing fist-sized chunks of ruidium. It is used only for private mission briefings and rare gatherings of all five of the faction's masters.

First Eclipse is owned and operated by Satzrak Runestrider, a neutral evil, blue dragonborn \textbf{occult extollant} (see appendix A). He is a loyal member of the faction, and he longs to keep the Consortium headquartered at his place permanently; the group has already used the tavern for three months longer than originally planned, so his hopes are up.

\subsubsection{Luck's Run}
\textit{World-Renowned Casino}

\noindent The Luck's Run casino stands three stories tall and is made of speckled red porphyry and topped with metallic domes in the palatial Marquesian style. After dark, faint purplish light escapes from the slit-like windows of the casino, and the sounds of laughter and music carry into the nearby streets.

Inside the casino, tables draped in crimson cloth fills the floor space, each location crowded with patrons trying their hand at one of the games of chance. Employees dressed in uniforms of dark gray and gold push through the throng while carrying trays of food, drinks, and coin between tables.

The casino is owned and run by Adima Shemsilver, a chaotic neutral, halfling \textbf{noble} who walks with the confidence of someone twice his height. Adima secretly uses half of the casino's earnings to fund the Veil, as his family has done for generations.

Three of the most popular games at Luck's Run are described below.

\subparagraph{Avandra's Favor}
This dice game has a 25 gp minimum bet. Each participant rolls 2d6 and wins back twice their bet on a roll of 7 or 12. A participant who doesn't win can roll another d6 and add it to their total by wagering an additional 25 gp. A participant can do this as many times as they want, until they "bust" by accumulating a total higher than 12.

\subparagraph{Gambit of Ord}
This card game requires a 50 gp initial bet, which goes into a pot in the center of the table. Each participant draws a card, then a second, and then a third. After each draw, each participant has a chance to raise the bet, stand pat, or fold. The participants reveal their hands after the third card is drawn. Whoever has the highest-ranking hand takes the pot.

To simulate the play of this game, have each participant roll dice in secret. At the first draw, all participants roll a d8. The second draw is a d6 roll, and the third is a d4 roll. Then, all reveal their dice, and whoever has the highest total wins.

\subparagraph{Quon a Drensal}
Also known as "Run of Luck," this game is a contest of lizard racing. Each player can bet 10 gp or more on a single racer in a field of three, and anyone who backs the winner gets back double the amount wagered. Those who bet on the second-place lizard get back half their wager.

To simulate the play of this game, roll 3d4 for each lizard. The winning lizard is the one with the highest total (ties are possible).

\subsubsection{Mystic Pursuits}
\textit{Arcane Shop}

\noindent Colorful tapestries frame the entrance to this establishment. Beyond the threshold, silk curtains cover the walls, and magical baubles cover low wooden tables illuminated by scented candles. A soothsayer sits in the back of the room, her deep blue hood pulled down over her eyes and a deck of cards spread out before her—offering an invitation to peer into one's future, if one so dares.

The "soothsayer" is a charlatan named Amkezne, a chaotic neutral, tiefling \textbf{mage}. She is a Crystal Chateau dropout who makes a living defrauding wide-eyed tourists with her mystical displays and fake enchanted trinkets. Customers who see through her act but treat her with respect can purchase any spell component worth 1,000 gp or less from Amkezne. She also has a collection of fine jewelry for sale, including five gold bracelets (25 gp each), three gem-studded gold rings (250 gp each), and two jewel-encrusted necklaces and tiaras (750 gp each). An \textbf{invisible stalker} watches over this collection and attacks anyone who tries to steal from it.


\end{document}
